
%%%%%%%%%%%%%%%%%%%%%%%%%%%%%%%% SCRIPT FOR E-RARA AND MDZ FILES     %%%%%%%%%%%%%%%%%%%%%%%%%%%%%%%%%%%%%%%%%%%%%%%%
%%%%%%%%%%%%%%%%%%%%%%%%%%% fini le 30.04.2025 par F. GOY            %%%%%%%%%%%%%%%%%%%%%%%%%%%%%%%%%%%%%%%%%%%%%%%%
% !TeX TS-program = lualatex
\documentclass{article}
\usepackage[T1]{fontenc}
\usepackage{microtype}
\usepackage[pdfusetitle,hidelinks]{hyperref}

\usepackage{polyglossia}
\setmainlanguage{english}
\setotherlanguages{latin,greek}
\usepackage[series={},nocritical,noend,noeledsec,nofamiliar,noledgroup]{reledmac}
\usepackage{reledpar}

\usepackage{fontspec}
\setmainfont{TeX Gyre Termes}

\usepackage{sectsty}
\usepackage{xcolor}

\usepackage{fancyhdr}
\pagestyle{fancy}
\fancyhf{}
\fancyhead[LE,RO]{\nouppercase{\leftmark}}  
\cfoot{\thepage}
\renewcommand{\headrulewidth}{0.4pt}

% Redefine \section to remove numbering
\usepackage{titlesec}
\titleformat{\section}[block]{\normalfont\scshape\color{gray}}{}{0pt}{} % no number in heading
\titleformat{\subsection}[hang]{\normalfont}{}{0pt}{} % also remove subsection number
\titleformat{\subsubsection}[hang]{\normalfont\footnotesize\color{black}}{}{0pt}{}

% Modify how section marks are stored to exclude numbers
\makeatletter
\renewcommand{\sectionmark}[1]{%
	\markboth{#1}{}} % Only store the section title, without number
\renewcommand{\subsectionmark}[1]{%
	\markright{#1}} % Only store the subsection title, without number
\renewcommand{\numberline}[1]{} % Hide the section number in TOC
\makeatother

\begin{document}

\date{}
        \title{In D. Pauli priorem Epistolam ad Timotheum commentarius [Geneva] : [Lambertus Dannaeus], [1577]}
\maketitle
\tableofcontents
\clearpage
\begin{pages} 
\beginnumbering
        IN PRIOREM EPISTOD.  PAVLI lam ad Timotheum comIN QVO NON SOLVM IPmentarius. sius Epistolae doctrina, et artificium singulorumque argumentorum loci explicantur: sed etiam vera Disciplinae Ecelesiasticae forma, tum ex dei verbo, atque ex ipso Paulo, tum ex veteribus synodis repetita atque restituta est, vt legitima et Apostolica regendae Dei Ecclesiae ratio et norma hodie in vsum reuocari possit. PER LAMBERTVM DANAEVM Accessit operi duplex index. APVD EVSTATHIVM VIGNON GENEVÆ M. D. LXXVII. \pstart AD LECTOREM. Quod saepius, Lector optime, iuris Pontificii canones, et Gratiani Decretum citauimus, ea solùm mente fecimus, ut ex his tanquam ex quadam illius temporis historia, quae veteris Disciplinae Ecelesiasticae vestigia tunc adhuc (et si tamen corruptioribus iam Ecclesiae temporibus ) restarent, intelligatur:non vt ex istis scriptoribus quicquam dubium et anceps decidatur. Nam etiam ex hoc ipso perspicient omnes, quantùm posteriori aetate Romana Ecelesia etiam ab eo tantillos quod veteris disciplinae reliquum erat) discesserit: et in quam miseram horribilemque tyrannidem et confusionem iam, iusto Dei iudicio, inciderit, vt ab ea nos summo Dei Optimi Maximi beneficio liberatos esse laetemur, Deoque ipsi perpetuo gratias agamus.  \pend
\phantomsection
\addcontentsline{toc}{subsection}{\textit{ILLVSTRISSIMO ET POTENTISSIMO PRINCIPI ac domino D. Gulielmo Principi Auracio Comiti Nassauao etc.  Christi gloriae assertori in Hollandia et Zelandia fortissimo, Gratiam et pacem a Domino. }}
\subsection*{\textit{ILLVSTRISSIMO ET POTENTISSIMO PRINCIPI ac domino D. Gulielmo Principi Auracio Comiti Nassauao etc.  Christi gloriae assertori in Hollandia et Zelandia fortissimo, Gratiam et pacem a Domino. }}\pstart \huge\textbf{V}\normalsize Etus est illud, Princeps Illustrissime, etiam ab Aristotele vsurpatum, Nihil esse in ipsa rerum vniuersitate pulchrius ordine. Quid enim vel oculis iucundum, vel menti etiam ipsi et animo hominis gratum obuersari potest, quod sit confusum, et nulla partium apta separatione distinctum? Certe quocunque non tantum oculos, sed omnes animisensus conuerterimus, si quae in eos incurrunt, neque ordine digesta, neque apto siru inter se cohaerentia, neque conuenienti loco et modo collocata a nobis apparebunt: ea neque vtilitatem, neque venustatem aliquam habere statim pronuntiabimus, tantûmque ab illis oculorum, animique intuitum auer\pend
\section*{EPISTOLA. }\pstart temus, quantum ea nos ad se rapiunt, quae commoda, propriaque ratione, et dispositione distinguuntur. Hoc in Regno, hoc in Rep. hoc in oppidis, hoc in pagis, hoc in priuatorum aedibus, hoc in hortis, et cultis sedibus, hoc in solitudine, hoc in rebus que natura gignuntur, hoc in artificiis quae hominum industria efficiuntur, verumesse ipsa rerum experientia, et publica mortalium omnium, non tantum piorum, sed etiam profanorum hominum vox testatur: denique Mundus ipse, pulcherrimum Dei opus, ab ordine κόσμος nominatur. Quo magis detestanda est eorum sententia, qui rerum confusione non tantum apud se delectantur: sed omnia publica priuataque euertere, et perturbationem (iusto ordine sublato) in res omnes inducere conantur, quos homines ἀτακτους et Ἀναρχικους, tanquam monstra, et generis humani ἀλάςτορας, iam pridem veteres appellarunt. Caeterum si qua in re alia ordinis obseruatio diligens non tantum vtilis, sed etiam omnino necessaria est, maxime in Ecclesia Dei locum habere ex Pauli Apostoli praecepto debet. Est enim Ecclesia domus Dei, id est summi illius regum Regis, in qua idcirco omnia circumspecte, apte, decenter, et, vt Paulus ipse lo\pend
\section*{EPISTOLA. }\pstart quitur, ἀσχημόνως καί κατα τάξιν fieri oportet: nihil autein confuse, et pertubate, ne ex ista domus turpitudine ad ipsum patrem familias, id est, ad Deum Optimum Maximum aliqua labes quodammodo redundet. Itaque ipse eadem benignitate qua in sacratissimo suo verbo omnia, quae sunt ad salutem aeternam, ipsiusque Ecclesiae suae fundamentum necessaria, praescripsit: it a et iam quae sunt ad eiusdem Ecclesiae ordinem, et ornamentum vtilia, demonstrauit, nequa deesset pars ad verae Ecclesiae constitutionem necessaria, quam non ab ipsius ore et verbo abunde disceremus. Vt vero hae praeceptiones, quae Ecclesiae necessarium ordinem continent etiam nomine ipso ab iis distinguerentur, quae doctrinam fidei complectuntur, generaliter disciplinae Ecclesiasticae vocabulo vulgo sunt appellatae. Quemadmodum enim ea verbi Dei pars, et ea veritatis cognitio, quae quid de Deo Patre, de Christo, eiusque naturis et officio, de Spiritu Sancto, de causis nostrae iustificationis, de vita aeterna, nobis credendum sit, ostendit ossa et neruos, ipsamque adeo verae Ecclesiae, tanquam corporis οὐσίαν, et constitutionem continet et aedificat.  \pend
\section*{EPISTOLA. }\pstart Sic illa eiusdem verbi Dei pars, quae disciplinae Ecclesiasticae praecepta tradit, huius Ecclesiae, velut corporis iam suis membris constantis, sed nudi vestimentum, cultum, et ornamentum compingit et consarcinat, quo decore vestiri possit, vt ex eo quanta sit vtriusque, id est, doctrinae et disciplinae Ecclesiae inter se consensio coniunctióque, satis iam intelligi possit. Erit enim coetus hominum, etiam eorum, qui puram Dei praedicationem amplexi sunt, velut nudum quoddam, et suis ornamentis carens, et destitutum corpus, nisi disciplinam Ecclesiasticam ad veritatis lucem et professionem adiunxerit. Denique illa eadem disciplina recte custos ipsius fidei, doctrinae, mater virtutum, vinculum honestatis, neruus societatis Christianae a Patribus Ecclesiasticis appellatur: cuius certe tanta vis est, vt sola illud conciliet Ecclesiae, quod decorum, venustatem, splendorem, et ornamentum appellamus, adeóque speciosam eam reddat, vt ab haereticis (qui istius disciplinae fuerunt senper hostes maximi et infestissimi ) inuidioso nomine lenocinium vt ait Tertull. Ecclesiastica disciplina diceretur. Cùm sit igitur haec disciplina illa iustitiae  \pend
\section*{EPISTOLA }\pstart pars et regula, per quam externus ordo, et totius Ecclesiae legitima administratio praescripta est, quis eam spernendam negligendamue vnquam existimauerit? Cuius institutio non tantum sub nouo testamento fieri coepta est: sed et iam in lege veteri a Deo ipso iam olim tradita fuit, vt in libello de Politia Iudaica doctissime iam pridem docuit Cornelius Bertramus Hebraicae linguae professor eruditissimus, et collega meus. Quid enim aliud sibi volunt illa Dei praecepta, quae Leuiticam tribum a reliquis distinguuntvt ex ea sola ministri templi eligantur: quae in ipsa tribu Leuitica iura Sacerdotum et Leuitarum discernunt, quae quanta summi Sacerdotis, quanta aliorum sit potestas describunt? quae qui ritus cuique  sacrificiorum generi, quae libamina, quae oblatio sit necessaria deffiniunt? denique tota ipsa lex ceremonialis vetus a Deo ipso dictata, quid aliud, obsecro, est, quam amplissima et luculentissima disciplinae Ecclesiasticae descriptio, non tantum autem quaedam illius delineatio? Nec vero minus benigne Deus per Apostolos sub nouo testamento eandem disciplinam nobis instituit vel potius collapsam instaurauit et restituit: vt pote omnino ad re\pend
\section*{EPISTOLA. }\pstart cte constituendam regendamque suam Ecclesiam necessariam, quam sub veteri descripserat. Id quod cùm ex variis aliis noui Testamenti locis apparet: tum maxime ex libro Actorum, ex prima Pauli epistola ad Corinthios, ex hac ad Timotheum, et vna illa, quae est ad Titum scripta:in quibus omnibus quicquid ad legitimam Ecclesiae gubernationem, et disciplinam Ecclesiasticam maxime pertinet, siue totam ipsam Ecclesiam, siue priuatos regendos spectes in qua illud totum accurate et copiose descriptum est. Vnde mirum et absurdum saepe mihi videri solet, quod tot postea canones, tot synodorum decreta, tot Ecclesiae constitutiones, tum generales, tum particulares ad constituendam disciplinam Ecclesiasticam fuerunt conscriptae. Quid enim his fuit opus? Sed magno contemptu et ignorantia eorum pre ceptorum, quae ab Apostolis ipsis, id est, a Dei spiritu antea tradita fuerant id certe factum fuit, quod haec diuina praecepta spernerent homines, vt sua commenta, vt traditiones, saepe etiam ineptissimas, at disciplinae ipso Dei verbo praescriptae contrarias in Ecclesiam Deil inueherent. Quid igitur, dicet quis, estne verbo Dei, et in no\pend
\section*{EPISTOLA. }\pstart uo Testamento disciplina haec Ecclesiastica tam accurate tradita, vt nihil iam addi possit: aut ea authoritate constituta, vt ab ea sine piaculo, discedi non debeat? Certe vt huic quaestioni satisfaciam, distinguendum dico. Sunt enim plures istius Ecclesiastice disciplinae partes, quarum alia generalem tantum Ecclesiae (quacumque orbis terra rum parte esse potest) et fundandae et retinende, et regendae modum continet, et describit, veluti in Dei Ecclesia vocem Domini nostri Iesu Christi audiendam esse eamque solam: Sacramenta piis tantùm et fidelibus danda esse et caet. huiusmodi, quae generaliter sunt conseruandae Ecclesiae praecepta. Alia magis particulares illius administran dae normas et praecepta tradit et persequitur. Haec vero posterior disciplinae Ecclesiasticae pars haec 4. capita complecti videtur: nempe. I. Personarum, quae ad Ecclesiasticum ali quod munus vocandae sunt, electionem, 2. Muneris et officii Personarum iam electarum descriptionem:3. Cen surae, morumque correctionis, quae tum in consistorio tum extra consistorium fieri et applicari singulis debet, rationem. 4. Eorum qui censura Ecclesiastica affecti sunt, et Ecclesiae iudicio abstenti, vel excom\pend
\section*{EPISTOLA. }\pstart municati dicuntur, cum eadem Ecclesiareconciliationem. Cùm igitur de hoc quaerimus, vtrum omnium horum capitum singula sic praecepta tradita sint, vt his nihil addi possit, respondeo, in quaque re quaedam esse οὐσιώδη, quaedam vero συμβεβηκότα. Sunt οὐσιώδη id est essentialia, quae rei ipsius vel tota natura sunt, aut illius pars: ve Iuti in hoc ipso argumento, Si quis neget in pastorum aut eligendorum vocatione Iegitima eorum vitam et doctrinam examinandam, vel approbationem a populo expectandam esse, ille ea negat, quae sunt Electionis legitimae οὐσιώδη. Si quis item hoc examen non ab Ecclesia, non a pluri bus, sed ab vno quopiam solo fieri et mero vnius arbitratu stare debere doceat, is οὐσιώδη electionis tollit. Sed si quis hoc vel illo loco, hoc vel illo tempore, hoc vel illo modo examen a pluribus  fieri debere disputet, is περί συμβεβηκότων tantum disputat, veluti, exempli gratia vetus electio et examen pastorum olim a tribus vicinioribus episcopis seu pastoribus conuocatis cum presbyterio et Ecclesia (cui dandus erat Pastor) fiebat: item in metropolitana tantum vrbe, certo tempore, nempe ieiunii publici tempore. Denique manuum eleuatione (quae χειροτονία  \pend
\section*{EPISTOLA. }\pstart dicebatur) non solo consensu populi, aut solo vocis suffragio, sane ista omnia sunt accidentaria, non essentialia in electione. Sic igitur statuimus, quae verbo Dei fieri et obseruari ad regendam Ecclesiam iubentur, siue illa de personis eligendis, siue de iam electarum munere, siue de reIiquis disciplinae Ecclesiasticae capitibus agant, illa, inquam, esse οὐσιώδη. Quae autem in quaque vel prouincia, vel Ecclesia ad haec ipsa praestanda, et ad praxin deducenda commodiora iudicantur, et constituuntur, ea esse συμβεβηκότα quaedam, veluti vt hoc potius quam illo tempore fiat ver bi Dei concio: hoc potius quam illo presbyterorum vel diaconorum numero contenta sit haec Ecclesia, hoc est συμβεβηκoς minimeque οὐσιῶδες. Iam igitur ad quaestionem propositam ex hac distinctione sic respondemus, vt quae sunt disciplinae Ecclesiasticae οὐσιώδη augeri minuiue nouis praeceptionibus, vel constitutionibus hominum posse negemus: quae sunt autem σομβεβηκότα, et augeri et minui et temperari a quaque non tantùm prouincia, sed etiam coetu et Ecclesia pro sua vel commoditate vel incommoditate posse: pro varia denique locorum, rerum, temporum,  \pend
\section*{EPISTOLA. }\pstart et personarum circunstantia. Eodemque modo respondemus ad alteram illam quorundam petitionem, in qua quaerunt vtrum Ecclesia a constituta semel per Apostolos disciplina recedere possit, ac vtrum sit vna et eadem disciplinae Ecclesiasticae in omnibus euangelicis Ecclesiis forma retinenda? Nam si de essentialibus disciplinae Ecclesiasticae agunt, dicimus illa ab omnibus Ecclesiis aequaliter eademque esse et seruanda et retinenda, veluti vt sit regimen Ecclesiae non Monarchicum, sed Aristocraticum: vt sit in quaque Ecclesia sena tus et coetus aliquis Presbyterorum ab Ecclesia electorum, sine quo nemo vel Pastor, vel alius quisquam negotia Eccle siae gerat : sit aliqua scelerum et morum vitiosorum censura, eaque non vnius alicuius arbitratu, sed ἡγουμένων, et vt ait Hiero. Presbyterorum, aestimatione et iudicio fiat: sit denique excommunicatio, non ex quorundam Iurisconsultorum et pragmaticorum, sed ex Dei verbo iudicantium hominum sensu et iudicio, lata in non poenitentem, approbante et collecta Ecclesia. Itaque οὐσίαν disciplinae Ecclesiasticae in omnibus Ecclesiis reformatis et Euangelicis eandem esse oportere ingenue, et a\pend
\section*{EPISTOLA. }\pstart perte respondemus. Sin autem de externis et accidentalibus huic disciplinae rebus quaeritur, in quibus commoditas incommoditasque praxeos, vsus, et executionis disciplinae istius versatur, dicimus illas variari posse, atque etiam saepe oportere propter diuersam diuersorum temporum, rerum, et personarum rationem : nedum vt loci et regiones diuersae eandem semper disciplinae formam patiantur, aut requirant. Itaque quae hoc loco, et in his regionibus commode geruntur (vti in Gallia Coenae Dominicae administratio per singula anni quadrimestria) in aliis fortasse regionibus commode iisdemque ue temporibus non fierent aut celebrarentur. Sic in aliqua Ecclesia excommunicatio ex scripto fiet, in alia vero potius ore ipso Pastoris pronuntiabitur. Sic in hac Ecclesia eius qui post professionem Christi fuit idololatra et transfuga, posteaque ad Ecclesiam rediit confessio die tantum dominica fiet, in illa quotidie: hic ex scripto, alibi ore ipsius poenitentis. Haec igitur omnia diuersam formam habere possunt, modo tamen coercitio morum, agnitio culpae, separatio hoedorum ab agnis in domo Domini semper retineatur, fiatque quod fieri prae\pend
\section*{EPISTOLA. }\pstart cipit Deus, Denique modo sit Ecclesia et Dei domus εὔτακτος, non autem sentina improborum, sed recte administrata familia. Substantia igitur disciplinae vbique in Ecclesiis Dei eadem esse debet: accidens autem illius mutari potest. Ergo diligenter quid essentiale, quid accidentale sit, distinguendum esse respondemus. Porro nemini, ac ne quidem ipsi summo magistratui, vel Ecclesiae toti concedendum arbitramur, vt illa, quae sunt οὐσιώδη disciplinae Ecclesiasticae, tollere, vel immutare possit. Etsi enim in rebus politicis Rex, summusque magistratus nouas leges ferre, et veteres abrogare iure potestatis sibi a Deo concessae potest, in disci plinae tamen Ecclesiasticae fundamento vel veteri, et a Deo iam posito tollendo: vel nouo reperiendo constit uendóque idem ius non habet. Cuius rei haec ratio est, et quidem manifestissima, quod cùm disciplinae Ecclesiasticae scopus et finis hic sit, primùm vt Pastores, et qui Ecclesiae Dei prae esse debent legitimam vocationem habeant. Postea vt hominum conscientiae gubernentur, et ad poenitentiam conuersionemque ad Deum adducantur, ad solum Deum disciplinae Ecclesiasticae definitio  \pend
\section*{EPISTOLA. }\pstart legislatioque sane quidem pertineat necesse est, quia ille solus Ecclesiae, et coetus istius Rex est, dominus, moderator et legislator, idemque solus conscientiarum nostrarum legitimus gubernator, arbiter, atque rector. Itaque solus ille potest leges, tum legitimae vocationis eorum, qui in sua domo praeesse debent, praescribere: tum etiam mouendae hominum conscientiae, ingignendǽque poenitentiae modum nobis demonstrare. In quibus duobus capitibus laus et vis summa vniuersae Ecclesiastice disciplinae, quemadmodum diximus, finisque positus est. De quibus rebus leges alias, et ab iis, quas Christus praescripsit, diuersas ferre neque Reges vlli, ac ne vniuersa quidem ipsa Dei Ecclesia (quae est in his terris) potest. Nam et ipsa ecclesia est duntaxat serua Dei, et non domina, et reges quoque ipsi sunt in hoc coetu membra tantùm non caput: membra quidem honorabilia, non tamen sic membra, vt in hoc coetu regem do imperent pro suo arbitratu, sed tantùm vt Dei mandata a piis et inter pios obseruari sedulo curent. Vocationum enim ecclesiasticarum, et conscientiarum nostrarum procul dubio solus Deus et scrutator et gubernator et Rex et autor esse debet  \pend
\section*{EPISTOLA. }\pstart et potest. Itaque cùm hominum conscientias immutare, inflectere, ad Deum adduce. re, huius disciplinae ecclesiasticae propositum sit et finis, quomodo, quaue via id fieri recte possit, a nemine, praeterquam ab vno Deo praescribi, definiriue potest, quia solus ille eas fabricatus est: solus quomodo eae regendae, flectendae, ac immutandae sint, nouit, solus earum dominus est. Ipsius enim solius templa et sedes sunt nostrae conscientiae. Vnde velle ab alio, et mortali quidem homine, leges inflectendae, et conuertendae nostrae conscientiae propter Dei verbum habere, est hallucinari velle, est operam perdere, est summam iniuriam Deo (qui sibi vni hoc totum ius reseruauit )inferre. Possunt igitur Reges, regumque consiliarii, viri prudentes, leges de publica inter ipsos ciues pace tuenda dictare et commentari: leges autem de conscientiae hominum regendae, aut inflectendae ratione et modo (praeter eas, quae Dei verbo praescriptae sunt) nouas cogitare et comminisci, denique de legitima eorum qui ecclesiam (id est domum, Dei non ipsorum Regum gubernant) vocatione praeter Dei. i. ipsius Domini domus voluntatem leges inquam et regulas condere et sancire  \pend
\section*{EPISTOLA. }\pstart certe quidem illi nullo iure possunt. Ac, vti dixi nec Ecclesia quidem ipsa vniuersa et collecta in Oecumeniis, quae vocant, conciliis hoc quoque potest, adeo vt mera fuerit tyramnis, et in ius ipsum Dei inuasio improbissima quicquid Synodi illae Oecumenicae contra praescriptam verbo Dei disciplinae Ecclesiasticae οὐσίαν induxerunt, et sanciuerunt: quanquam se aliquid melius et sanctius moliri existimarent. Quod verum esse, vel ipsa Ecclesiae confusio, ruinaque, quae inde con secuta est, manifeste declarat. Nam si illa Apostolicis scriptis praescripta disciplina retenta fuisset, numquam in eam perturbationem Ecclesia incidisset, a qua miserrime tandem pessundata est. Quid igitur, damnamûsne omnes eas Euangelicas Ecclesias in quibus vel nulla disciplina Ecclesiastica omnino obseruatur, vel Apostolica non retinetur: sed in qui bus ea vel corruptaest: vel etiam alia, pro vera substituta ? Ad hanc odiosam qui dem quaestionem, et quae tamen saepe nobis obiicitur, ingenue quidem, sed tamen placide et modeste respondemus. Nos propter carentiam et εἰπουσίαν? disciplinae Ecclesiasticae minime damnare Ecclesias et coetus illos hominum, in quibus vera pu\pend
\section*{EPISTOLA. }\pstart raque Euangelii Christi praedicatio retino tur, summota omni idololatria: ergo non diffitemur et eas etiam veras esse Dei Ecclesias. Vnica enim est, eaque fundamentalis verae Ecclesiae nota in his terris, pura nimirum verbi Dei predicatio a legitimo pa store facta, quae vbicumque est, illic quoque est Grex Domini, iuxta illud Christi, Oues meae vocem meam audiunt: et si ab eo coetu saepe Sacramentorum administratio, sae pe etiam morum correctio, Ecclesiasticǽque disciplinae vsus et praxis exulat, atque abest. Quod enim hae duae posteriores notae, tanquam necessariae, ad verae Ecclesiae designationem a nomnullis pronuntiantur, ostendunt illi tantum, quid quantumque requiratur, vt omnibus suis numeris absoluta sit verae Ecclesiae facies. non autem illud volunt, omnino Dei Ecclesiam dici eum coetum non posse, in quo pure et vere verbum diuinum praedicatur, si nulla sit disciplinae Ecclesiasticae praxis. Sunt igitir huiusmodi Ecclesiae quidem, sed non omnibus suis partibus et numeris absolutae: quin potius claudicantes. Sic ille a nobis dicitur homo, cui brachium vtrumque abscissum est, si modo idem ille ratiocinatur. Sed hoc conari, hoc studere, hoc ef\pend
\section*{EPISTOLA }\pstart ficere omnibus modis debent illae Ecclesiae, quae summo Dei beneficio idololatriam abiecerunt, verum Dei verbum et vitae aeternae fundamentum receperunt, et adhuc tamen ista disciplina carent, vt hoc tam praeclarum ornamentum. i. disciplinam hanc Ecclesiasticam ( qualis ab Apostolis praescripta est) ad corpus illud Ecclesiae suae adiungant, et addant, vt non modo neruos et ossa Ecclesiae habeant, sed etiam vt sanguinem, colorem, venustatem, vires denique Ecclesiae integras retineant. Attendendum est enim diligenter ad illud Spiritus sancti dictum rectas faciendas esse vias nostras, ne quod claudicat tandem prorsus abducatur a via. Nam si quae necessaria sunt per incuriam in Dei Ecclesia negligantur, temporis successu illa penitus a Deo tandem abalienantur. Quia vero haec discipli na maxime hac Pauli Epistola, quae prior est ad Timotheum, est explicata, idcirco eam potissimum delegi, in cuius interpretatione meam de toto hoc argumento sententiam breuiter, more meo, exponerem. Scripserunt in hanc ipsam Pauli Epistolam varii, tum veteres, tum recentiores interpretes, eruditi illi quidem: sed qui se in hoe de disciplinae Ecclesiasticae ar\pend
\section*{EPISTOLA. }\pstart gumento maxime ostentare voluerit fuit Claudius Spensaeus. Is enim praeter magnos in ipsum Pauli contextum commentarios etiam tres libros Digressionum, quas appellat, et ipsos quoque prolixos edidit in hanc Pauli Epistolam. Tantum autem abest, vt, meo quidem iudicio, quisquam iis adiuuari possit, vt nihil sit et scriptore isto obscurius, et illius disputatione magis perturbatum atque confusum. Sunt enim isti scholastici, e quorum numero fuit Spensaeus, iusto Dei iudicio occaecati: et tam tenebriones, vt in plerisque ne seipsos qui dem intelligant, quia sunt in tenebris et mataeologia educati. Certe cum in eos commentarios incidissem, nihilque in iis proficerem, mature abistis digressionibus et errore ad ipsum Paulum et verum iter redii, fecique illud, quod habet prouerbium vetus bibi et fugi. Quantum autem hoc meo labore ipse sim consecutus, et rem caeteroquin propter temporis vetustatem, diuturnanque desuetudinem, iam obsoletam ac difficilem illustrarim, aliorum esto iudicium: certe magnis laboribus et lucubrationibus haec vetustatis rudera nobis eruenda fuerunt. Precor autem Dominum Deum vt hos meos labores Ecclesiae suae vtiles,  \pend
\section*{EPISTOLA. }\pstart piisque viris suaues esse velit. T.vero E. Princeps Illustriss. cui eos animo demisissimo ac lubentissimo dono, dicóque, si probabuntur, certe operae meae recte collocatae fructum hunc maximum iam ero consecutus. Quamobrem autem hoc, quicquid est operae meae, T. E. dedicarem cùm variae et multae sunt, eaeque graues causae, tum hoc imprimis, quod cum summo Dei Optimi Maximi bene ficio, tuaque illa plus heroica virtute et constantia Zelandiae et Hollandiae, totiusque illius tractus Ecclesiae Dei pace recreentur, nihil vel ad eas, iam praeclare instituendas: vel quae erunt constitutae retinendas hac disputatione aptius atque vtilius, adeo etiam magis necessarium esse existimem. Quemadmodum enim Ecclesiasticae disciplinae contemptus, atque praetermissio magnas in Ecclesiam ruinas semper induxit:ita diligens illius obseruatio et ad Apostolicam normam reuocatio non tantum caelestem ipsam doctrinam sartam tectam puranque inter mortales conseruauit, sed etiam mores ipsos hominum nomen Christi proffitentium, si qua in parte ad vitia delabebantur ad meliorem frugem conuertit, et  \pend
\section*{EPISTOLA. }\pstart In Dei metu officióque continuit: est denique haec eadem aduersus enascentes haereses sepes et munimentum validissimum, aduersus iam natas et infelicissime pullulantes, tanquam aduersus magica venena, remedium praesentissimum. Quanquam autem T. E. doctissimorum virorum, imprimis autem generosissimi praestantissimi que illius viri Ph. Marnixii Domini montis Sancti Aldegondii, et fidelissimi verbi Dei Ministri Taffini, (a quibus haec ista discet) eruditissimis sermonibus non caret, haec tamem in commune vestrarum Ecclesiarum commodum, tanquam symbolum, et ipse pro virili volui conferre, vt illae intelligant quantum pro T. E. et illis ipsis bene ex animo velim, ac felicissima atque optima quaeque a Deo ardentibus votis exoptem. Tuae vero virtuti, Princeps Illustrisime, quid non debeatur, cùm ab omnibus fidelibus et Christianis viris, tum vero a nostris Gallis, quorum salus et incolumitas semper summae curae fuit tibi, fratique tuo Illustriss. Comiti Ludouico ( cuius felix memoria nunquam caneset, ac ne seclis quidem innumerabilibus, cuius vtriusque et armis et prudentia Ecclesiae Gallicanae ab aduersariorum atrocissimis  \pend
\section*{EPISTOLA. }\pstart Iniuriis contra leges patrias, moremque maiorum illatis tandiu sunt defensae: denique cùm tua sit hodie in orbe terrarum gloria, tantaque rerum pro Dei gloria asserenda gestarum laus, vt parem tibi tanque constantem principem nulla adhuc hominum aetas viderit. Tu Prouinciam tibi commissam ab Hispanorum ac extraneorum tyrannide eripuisti, tu Belgium intestinis morbis profligatum ac iam deploratum inexpectata ratione et medicina sanasti, ac in iura maiorum pristinumque splendorem restituisti. Tu hostibus io triumphum iam canent ibus fugam iniecisti, et victor ipse triumphatórque Christi vexilla praeferens rediisti, Tu pietatis assertor, tu Iustitiae publicaeque salutis vindex fortissimus ore omnium mortalium celebraris, adeo quidem, vt qui tanta, tamque praeclara et pro Domini nostri Iesu Christi regno amplificando et pro patria libertate tutanda bene geris, tu quoque  bene inter omnes mortales mereris audire Βροτοῖς γὰρ vt ait Eurip. χρηστὰ πρα- γματα τα Χρηστῶν ἀφορμὰς ἐνδιδωσ᾽ ἀεί λογων. Ego vero Princeps Clementiss. cùm tot admirandas virtutes tuas commendari vbique gentium intelligam, et hoc quoque  \pend
\section*{EPISTOLA. }\pstart ue in T.E. obseruantia testimonium et tuarum laudum praedicationis conatum et symbolum extare volui quod vt grato benignóque animo T. E. accipiat obnixe precor. Dominus Deus solus, Æternus Sapiens, Bonus, Iustus, Misericors, te velut alterum Danielem, Spiritu suo corroboret ac confirmet, et ad Ecclesiae suae presidium diutissime felicissimeque nobis, et toti Belgio incolumem conseruet. Geneuae Cal. Augusti Anno temporis vltimi. CIↄ Iↄ LXXVII.  \pend
\phantomsection
\addcontentsline{toc}{subsection}{\textit{IN D. PAVLI APOSTOLI priorem ad Timotheum E pistolam. }}
\subsection*{\textit{IN D. PAVLI APOSTOLI priorem ad Timotheum E pistolam. }}
\section{L. Danaei Commentarius. CAP. PRIMVM.}
\phantomsection
\addcontentsline{toc}{subsection}{\textit{\huge\textbf{P}\normalsize AVLUS Apostolus Iesu  Christi ex ordinatione Dei seruatoris nostri, et Domini Iesu Christi, spei nostrae. }}
\subsection*{\textit{\huge\textbf{P}\normalsize AVLUS Apostolus Iesu  Christi ex ordinatione Dei seruatoris nostri, et Domini Iesu Christi, spei nostrae. }}\pstart Paulinarum Epistolarum genus est duplex, vnum earum, quae ad Ecclesias totas scriptae sunt. Alterum, earum quae ad Particulares quasdam per sonas quae erant, vel Ecclesiasticae, vti vtraque ad Timoth. et vna ad Titum. vel Priuatae, vt ea quae est ad Philemonem. Nam quae circunferuntur Pauli Epistolae ad Senecam et contra, sunt fictae et ementitae. Quanto vero est dignior, qui Ecclesiae praeest priuato homini, ita vberiora et locupletiora sunt earum Epistolarum argumenta, quae ad praefectos Ecclesiae scriptae sunt (quales sunt illae quae ad Timotheum, et Titum, quae pene sunt paris argumenti) quam earum, quae ad  \pend
\section*{AD I. PAVL. AD TIM. }
\marginpar{[ p.2 ]}\pstart priuatos. Quae fuerit autem huius Epistolae ad Timotheum scribendae ratio, et occasio facile ex ipso Paulo verss.3.intelligitur, nimirum, quod iam quidam sanam Euangelii doctrinam, ab ipso Paulo et aliis Fidelibus Christi Ministris traditam, et explicatam varie corrumperent tum in Substantia ipsa et doctrina, tum in Modo et forma illius tradendae. Cui malo vt mature occurreret Paulus, Timotheum monet quid opus sit facto, non tam ipsius quidem gratia, quam aliorum, id est, istorum doctrinam veram et sanam deprauantium. Quando scripta fuerit infra cap. hoc ipso vers. 3. Vnde scripta sit in fine huius Epistolae dicetur, Domino fauente. Diuiditur autem totum corpus, et σύνταγμα huius Epistolae in duas tantùm partes, nempe, Salutationem, et argumenti suscepti tractationem. Epilogo enim, id est, postrema illa parte caret, in qua variae salutationes esse solent, quae in aliis tamen Epistolis extant, quae hoc loco non erant necessariae. Salutatio autem haec tria continet vbique, nimirum (Quis scribat (Ad quem scribatur (Summam precationis Ergo quaeritur, quis seribat? Scribit ( Plaulus, personam designat ( Apostolus, munus suum describit (Plus est enim dixisse Apostolum, quam seruum Dei simpliciter, quod alibi facit. Chrysosto. hoc annotat Apostolum autem se vocat, et Iesu Christi:non  \pend
\section*{CAPVT I. }
\marginpar{[ p.3 ]}\pstart hominum: in quo est commendatio muneris. Nihil enim nisi apprime salutare habet legatio et munus huiusmodi et a Christo demandatum et Dei ipsius ordinatione siue decreto. In quo inest confirmatio huius muneris. In hac autem statim epigraphỹ apparet tum summum amoris Pauli erga Timotheum et luculentissimum testimonium: tum illius muneris et Apostolatus firmamentum maximum. Quod non ipsius quidem Timothei gratia fit, sed aliorum, quos admoniturus erat Timotheus, nec quidem tam suo, quam Pauli nomine et authoritate. Itaque hanc authoritatem assertam prius confirmatamque esse oportuit. Deinde ipsius quoque Timothei et confirmandi et commendandi apud Ecclesiam causa, qui a Paulo missus erat, hoc additum est, quo iustior ipsius vocatio appareat, Pauli vocatione semel stabilita, etsi de ea Ephesii incerti esse non potuerunt, inter quos vixit per biennium Paulus. Nihil autem hoc loco fit a Paulo ostentationis et iactantiae animo, sed quod muneris authoritatem defendi aduersus quorundam impostorum calumnias oportuit, vidit que pro temporis circunstantia esse necesse. Caeterum docet hic locus a quo mitti de beamus et vocari, nempe a Deo. Nemo enim vsurpat sibi honorem Hebr. 5.V.4. etsi non omnes extraordinarie (vti Paulus) vocantur.  \pend\pstart Finem autem, et caussam tanti Dei beneficii erga homines ostendit, dum appellat Deum seruatorem. Quia enim Deusnos seruare vult, etiam mitti ad nos, qui ipsius Euangelium et voluntatem annuntient. Psal.74.Luc 13.vers.34. id est,  \pend
\section*{AD I. PAVL. AD TIM }
\marginpar{[ p.4 ]}\pstart qui istius viam salutis exponant nobis. Vocat autem Patrem seruatorem vti Tit I.vers.3 et 2.vers 10. et Christum spem nostram. Vti enim Pater dedit Filium, per quem seruamur. Ioan.3.V. I6.ita Filius est fundamentum spei et salutis nostrae, vt cùm haec Patri et Filio communiter tribuuntur, intelligamus esse indiuisa Trinitatis opera, et tres personas paris, non autem dissimilis maiestatis et potestatis. Hec breuiùs persequor, quia a multis copiose tractata sunt.  \pend
\phantomsection
\addcontentsline{toc}{subsection}{\textit{ Timotheo germano filio in fide. Gratia sit tibi misericordia et pax a Deo Patre nostro, et Christo Iesu Domino nostro. }}
\subsection*{\textit{ Timotheo germano filio in fide. Gratia sit tibi misericordia et pax a Deo Patre nostro, et Christo Iesu Domino nostro. }}\pstart Timotheus quis fuerit apparet ex Act.16. et variis Epistolarum Pauli locis: Hunc autem vocat et Filium in fide, vt distinguat a carnali filio: et Germanum. i. verum et referentem veros mores Patris Pauli: non igitur fuit hypocrita et fictitius, vt hodie multi qui magnorum Dei seruorum et filii sunt et discipuli. Quod vtrunque pertinet ad Timothei laudem et imitationem.  \pend\pstart Sic Titum appellat Filium Tit.I.v3 et Corinthios quoque I Corinth. 4. vers.14. Tamen obstat quod ait Christus Math.23.vers.9. et Hebr.12.V.9 Vnus est Pater animorum, nempe, Deus. Resp. Verum est per se et proprie vnicum esse animorum nostrorum Patrem, nimirum Deum ipsum: sed communicatione quadam, et quatenus homines sunt huius nostrae in fide et Ecclesia Dei generationis, aut potius regenerationis instrumenta, vox ipsa Patris benigne illis a Deo communicatur, eti et Patribus carnalibus eadem vox Patris. Plus  \pend
\section*{CAPVT I. }
\marginpar{[ p.5 ]}\pstart autem debemus Patribus qui nos in Dei metu instruunt, quam qui nos tantùm corpore gignunt. Nec quod pastores dicantur patres nostri in fide excusat patres ipsos corporum et carnales, quo minus instituant liberos suos in vero Dei metu et pietate Deut.4. vers.II.et II.V.I9. sed differentia tamen ostenditur inter vtrunque genus patrum, quorum alii sunt nobis tantùm spirituales Patres, alii sunt carnales.  \pend\pstart Filii autem dicuntur, ait Augustinus contra Adimantum cap.5. (Natura aut ( Doctrina Imitatione Hic et doctrinae et imitationis ratione Timotheus filius dicitur, non tantùm aetatis, vt tamen putat Erasmus. Salutatio est in his verbis posita Gratia et haec (Deo Patre nostro Misericordia (omnia ( Pax a (Domino Iesu Christo. Misericordia addita est hic, vti in Epistola ad Titum, praeter stylum et communem morem Pauli, quo magis suum erga Timotheum amorem testaretur, et se miro erga eum affectu impelli et affici ostenderet. Recto autem ordine et naturali vtitur in his enumerandis Paulus, et progreditur a causis ad effecta. Est enim gratia causa misericordiae, vt misericordia ipsa causa pacis nostrae. Est igitur gratia gratuita nostri acceptatio, quae fit a Deo, quam sequitur misericordia, id est, peccatorum nostrorum remissio. Itaque vocem misericordiae hic passiue pro effecto, quod in nobis apparet: non autem actiue in Deo accipimus  \pend
\section*{AD I. PAVL. AD TIM. }
\marginpar{[ p.6 ]}\pstart Author autem tot et tantorum donorum est Deus Pater noster, in et per Christum Dominum nostrum, in quo omnes Dei promissiones et beneficia sunt, Amen et Etiam  \pend
\phantomsection
\addcontentsline{toc}{subsection}{\textit{3 Quemadmodum te sum precatus vt permaneres Ephesi, quum proficiscerer in Macedoniam, permaneto, vt denunties quibusdam ne diuersam doctrinam doceant: }}
\subsection*{\textit{3 Quemadmodum te sum precatus vt permaneres Ephesi, quum proficiscerer in Macedoniam, permaneto, vt denunties quibusdam ne diuersam doctrinam doceant: }}\pstart Secunda pars huius Epistolae sequitur, in qua de re ipsa et argumento est agendum. Primum autem occasionem scribendi narrat, deinde rem ipsam aggreditur Paulus. Ex ipsa vero disputatione apparet magnam partem esse hanc epistolam ἐπανωρ θωτικήν et δειδακτικήν. Constat autem sex capitibus, quorum summa breuiter hac nostra tabula comprehensa est. In summa in hac Epistola plena ministerii Euangelici descriptio continetur, et disciplinae Ecclesiasticae institutio explicatióque. Nam quae ad particulares quaestiones pertinent, suntque  scitu necessaria videntur a Paulo tradita in Epistola ad Corinthios.  \pend
\section*{AD I. PAVL. AD TIM. }
\marginpar{[ p.8 ]}\pstart Precatus) blanda et modesta ad monitio: quae eo acutius pungit et magis serio, quo maior est Pauli in Timotheum authoritas, id est, praeceptoris et patris in filium, et discipulum ius et affectus. Vtrunque autem potest Pastor, Denuntiare vt praefectus et pastor: et Precari vt frater2 Tessal. 3.V.I2. Nam quam hic sibi non tribuit, dat Timotheo postea potestatem:nimirum, Denuntiandi et quasi iubendi. Hic autem nihil de Papali authoritate et Petri iure agitur in Timotheo Epheso praeficiendo, cuius consensus in Timotheo vocando non interuenit, neque a Paulo est vel quaesitus, vel expectatus.  \pend\pstart Annotauit Eras. lucidiora esse Graeca, quam Latina, quod verum est, et repetendum est aliquid ἀπό τοῦ κοινοῦ, sic Quemadmodum te rogaui, vt permaneres Ephesi, maneto, vt denunties, et c. Sic enim facile fluit sensus.  \pend\pstart Proficiscerer in Macedoniam, Haec est tertia profectio (quemad modum obseruauit D.Th. Beza praeceptor meus) quae contigisse videtur annoa Pauli conuersione 28. et a nato Christo 62. Huius sit mentio Act. I9. et 2O. Ergo scripta est haec Epistola ante eas, quae sunt ad Ephesios, Philippenses, Colossenses. Post eam tamen vtranque quae est ad Corinthios, et ad Titum. Post vtranque eam, quae ad Thessalonicenses scripta est, vt quidam putant.  \pend\pstart Apparet etiam Paulum habuisse Ephesiorum Ecclesiam commendatissimam, cui tam sollicite prospicit.  \pend\pstart Ne aliter doceant) Cur hoc timuerit Paulus ratio satis intelligi potest ex Apocalyp.2.vers. 6.  \pend
\section*{CAPVT I. }
\marginpar{[ p.9 ]}\pstart et quod iam Galatae, quae erant vicinae Ecclesiae, abducti fuerant a synceritate et puritate doctrinae Euangelii per huiusmodi pseudofratres concionatores et nebulones, qui aliter, quam Apostoli, Euangelii praetextu docebant. Quid sit autem ἑτεροδιδασκαλεῖν hic et inf cap. 6.V.3. varia est interpretum sententia: sed ex legitimi pastoris definitione facile tanquam ex contrario perspicietur. Illa autem definitio extat apud Matthaeum 22.vers.I6. vbi de Christo agitur. Est igitur verus Pastor qui docet viam Domini in veritate, idque  sine personarum acceptione Malach.2. Ergo peccatur vel in Re et doctrina, vel in Forma et modo docendi, quam ὑποτύπωσιν appellat Paulus  Timoth I.vers.13. vel in Applicatione et vsu doctrinae 2 Timot.2.vers.15. Quibus omnibus vitiis et modis aliquis ἑτεροδιδασκαλεῖ, id est peccat in docendo.  \pend
\phantomsection
\addcontentsline{toc}{subsection}{\textit{4 Nec attendant fabulis et genealogiis infinitis, quae potius quaestiones praebent, quam aedificationem Dei quae est per fidem. }}
\subsection*{\textit{4 Nec attendant fabulis et genealogiis infinitis, quae potius quaestiones praebent, quam aedificationem Dei quae est per fidem. }}\pstart Εʹξήγησις. Illustrat exemplo quod sup. versic. praecepit. Quid sit nimirum ἑτεροδιδασκαλεῖν. Ex hoc autem loco apparet non eos tantum in docendo peccare, qui falsa tradunt, sed etiam eos qui vana et inutilia. Denique qui non necessaria curiose et diu inquirunt et perscrutantur. Cuiusmodi olim fuit genealogiarum percensio, vel inuestigatio. Hodie infinitae quaestiones quae sunt in scholastica. Primùm de hypothesi agamus, deinde veniamus ad thesin. Ac falsa est eorum inter\pend
\section*{AD I. PAVL. AD TIM. }
\marginpar{[ p.1O ]}\pstart pretatio et sententia qui volunt nomine geneaIogiae hic intelligi, vel Mathematicas descriptiones natiuitatum, quas tamen damnatas esse agnoscimus. vel longas illas series et disputationes de Æonibus, quas Valentinus ipsiusque discipuli etiam post Pauli mortem, nimirum circa annum a Christo passo cx. inuexerunt in Ecclesiam, tanquam fidei doctrinam necessariam.  \pend\pstart Quod ad horoscoporum inspectionem et natiuitates quae ab Astrologis fiunt, illae non γενεα- λογίαι appellantur, sed γενέσεις et γενεθλιαλογίαι, et qui eas tractant et docent, non γενεαλόγοι, sed γενεθλιακοί nominantur. Hanc tamen interpretationem videtur amplexus Chrisostom. quae huic loco minime conuenit.  \pend\pstart Quanquam autem inter Ephesios, qui Iudaei non erant, versabatur Timotheus, tamen inter Iudaeos maxime huiusmodi quaestiones, quas hic damnat Paulus, fiebant, vti etiam Tit. I.vers.14 eas Iudaicas fabulas vocat Paulus. Sed vel iam huiusmodi fermentum in Ecclesiam Dei irrepserat, vel ex ea communicatione, quae fidelibus cum Iudaeis conuersis intercedebat hoc periculum iam Ecclesiae imminere potuit, et praeuideri a Paulo.  \pend\pstart Fuerunt autem huiusmodi de Genealogiis disputationes. Quis ex qua familia ortus esset, siue de se Iudaei agerent: siue de gentibus. Id quod partim ostentationis plenum est, partim inanis cognitionis.  \pend\pstart Μύθους, igitur vocat huiusmodi quaestiones, quod etsi saepe non sint prorsus res falsae et fictae pro animi libito:tamen sunt inutiles, vti fabulae et ficte narrationes Μύθος autem vnde nostri Gal\pend
\section*{CAPVT I. }
\marginpar{[ p.11 ]}\pstart li duxisse videntur suum illud Mot et Latine Mutire, est vox aliquid significans hominis ore egressa et prolata, et generalis illius significatio est et acceptio. Sermonem enim quemlibet notat etiam apud Homerum vox Μῦθος. Iubet igitur sobrie potius eas delibare Paulus, quam iis inhaerere et immorari. Ait enim Ne attendant. Nam quasdam tradit etiam Dei verbum nec inanes, nec infructuosas: nec tantùm vetus testamentum, sed etiam nouum, vti Genealogiam Christi ipsius. Id quod fit ad fidei nostrae confirmationem, et vere humanitatis Christi approbationem, in quibus eatenus nobis versari licet quatenus et modum teneamus, et scopum a Spiritu sancto propositum non transiliamus. Sic fecerunt quoque Patres, vt Augustinus lib 2.3.7. contra Faustum, et liber 2. de Consens. Euangelist.  \pend\pstart Hae disputationes et inuestigationes nimiae, videntur ex eo cepisse, quod ante Christum et sub lege studiose inter Iudaeos cuiusque familiae descriptio fiebat et tum publicis archiuis, tum priuatis monimentis religiose et diligenter seruabatur, adeo vt reiecti sint et expulsi ex sacrificatorum numero qui suam Genealogiam ostendere non possent. Num. 1.vers.18, et Nehem. 7. V.62. Igitur iussi sunt eas habere tum in archiuis publicis tum priuatis scriniis Iudaei Numer. 25.vers.5. adeo vt Prophetae ipsi libros totos de iis ediderint vt est 2 Chron. 12. vers. 15. Sed in eo ratio diuersa fuit illis et nobis, illi expectabant Christum ex vna quadam et certa suae gentis familia, et genere nasciturum, non ex quauis tribu. Idcirco distinctio familiarum diligenter retinenda ab iis  \pend
\section*{AD I. PAVL. AD TIM. }
\marginpar{[ p.12 ]}\pstart fuit vt melius postea Christum agnoscerent. Nos iam ex ea familia natum scimus et habemus vt hac inuestigatione nobis hodie non sit opus. Deinde praedicatur Christus, vt non vnus tantùm familiae vel gentis: sed omnium populorum, omnium nationum Mediator, Deus et Redemptor, vt in sola Iudaeorum gente non sit hodie nobis Christus quaerendus. Dixit, et quidem recte, August. etiam veteris Testamenti Genealogias ad Christum pertinuisse liber 19. contra Faust. cap.351. Vnde apud illos ante natum Christum fuit haec quaestio tolerabilior, quae hodie cessat nisi si sumus infideles, et Christum vere natum secundum Prophetas dubitemus.  \pend\pstart Duplici nomine condemnat eas Paulus, quod et sint Infinitae, et Inutiles ad vitae pietatem in nobis serendam et gignendam. Vtrunque pugnat cum vero omnis studii et inquisitionis fine. Quaerimus enim vt inueniamus: et vt tandem ad aliquam cognitionem, atque etiam eam nobis vtilem perueniamus et insistamus ei acquiescentes. Quem fructum aut vsum non praestant hae quaestiones. Edificationem igitur eam dicit quae est in fide, id est, quae secundum ipsius Dei voluntatem, verbum, et praeceptum est, et solida. Nam finis Theologiae verae, est aedificatio conscientiarum nostrarum, non vana et turgida scientia variarum disputationum et quaestionum. Ædificant enim, sed ad pugnas et clamores tantùm eae quaestiones, quae sunt de huiusmodi rebus, ac propterea vt sunt ita merito appellantur inutiles 2 Timoth.2.v. 16 et 23.  \pend
\section*{CAPVT I. }
\marginpar{[ p.13 ]}
\phantomsection
\addcontentsline{toc}{subsection}{\textit{5 Porro finis mandati est charitas ex puro corde, et conscientia bona, et fide non simulata. }}
\subsection*{\textit{5 Porro finis mandati est charitas ex puro corde, et conscientia bona, et fide non simulata. }}\pstart Αʹντίθεσις . Opponit ex aduerso inutilibus istis quaestionibus veram Theologiae tractationem, quae charitatem pro fine habet, Verbum Dei profundamento, πρᾶξιν autem, eamque legitimam, pro laboris fructu et pro iusto finis, ad quem tendit testimonio, quae tria breuissime quidem, sed tamen grauiter complectitur hoc loco Paulus.  \pend\pstart Finis praecepti) id est legis Dei, quam suae curiositati praetexebant, est charitas. Ergo non ineptae et inutiles quaestiones istae. Legis nomine Scripturam sacram intelligit, quae Legis nomine appellatur, quod tunc temporis alia nulla sacra scriptio adhuc esset, nisi de Lege, etsi Legem ipsam a Prophetarum scriptis nonnunquam distinctam legimus Luc 24. Sed et Legis nomen latissime patere certum est apud Paulum. Vnde Lex fidei et Lex Spiritus a Paulo dicitur Romanor.3, v. 8. Hic autem pro Verbi diuini tractatione. Nam docet Paulus quis sit verae Theologiae finis. Is autem est non theoriticus sed practicus, non inanis sed fructuosus consolatióque  animarum, quia charitas actionem producit, nec in sola contemplatione nec in verbis solis et lingua versatur a Ioan.3.v 18. Scholastici aliud de sua Theologia sentiunt, cuius finem constituunt contemplationem et Theoriam. Aiunt enim praestantissimae scientiae est praestantissimus finis. At Theologia est ars praestantissima. Haec concedimus. Sed quod aiunt,  \pend\pstart Contemplatio est nobilior actione. Resp.  \pend
\section*{AD I. PAVL. AD TIM. }
\marginpar{[ p.14 ]}\pstart Actionem hic non accipi a nobis pro opere mechanico: sed pro praestatione aliqua, cuius fructus sit Dei gloria et aedificatio conscientiarum. Contemplatio autem nuda et mera imperfectus actus et humanae imbecillitatis indicium potius est quam laus potentiae, si desit scilicet istius contemplationis aliquis effectus et fructus qui cognitionem nostram consequi possit.  \pend\pstart Aiunt isti Dixit Christus Luc. 1o. meliorem partem elegit Maria sibi. Respond. Versabatur in actione, non in nuda disputatione, cuius exercendae gratia mundo ipsi renuntiabat Maria.  \pend\pstart Aiunt Tertiae praestantissimae hominis partis actio est praestantissima. Animus autem, cuius propria est contemplatio, est corpore excellentior. Ergo et contemplatio est nobilior actione. Resp. Vera esset conclusio, si actio, de qua loquimur, solius corporis viribus fieret, non autem praecipue operatione, id est, intellectu et voluntate ipsius animi nostri, ad cuius obsequium postea for mantur corporis membra. Ergo quum dicimus actionem esse finem verae Theologiae et Christianae, nomine Actionis vel Operis non intelligimus seruile et mechanicum aliquod opus, sed actionem intellectus et voluntatis nostrae ad Deum tendentis nobilissimam. Nam ipsi illi Scholastici etiam constituunt int ell ectum practicum esse facultatem animi humani praestantissimam.  \pend\pstart Aiunt quarto, oculus est praestantior manu. Resp. Κατάτι. Neque etiam hic operis nomine intelligimus tantùm manuales actiones, sed cultum Dei, reformationem vitae quae sunt actiones nobilissimae  \pend
\section*{CAPVT. I. }
\marginpar{[ p.15 ]}\pstart Charitas) Et ad Deum, et ad homines referri potest haec charitas, quia Dei verbum vtrunque docet et Deum et homines diligere. Nos tamen de charitate ea, quae proximum spectat, hoc loco intelligemus.  \pend\pstart Huic igitur charitati addit, Cor purum, Conscientiam bonam, Fidem veram et synceram. Per quae tria non tantùm vera charitas descripta est, sed etiam fundamentum illius et vsus demonstratur, ne hanc charitatem, quam finem praecepti dixit, putemus esle quiddam otiosum, et contemplatiuum, vt quidam fingunt. Est.n. affectus charitas qui in animis nostris est, non vt lateat: sed vt sese proferat ipso opere.  \pend\pstart Quia igitur charitas est animi affectus huiusmodi, Duo habeat oportet nempe, Veritatem et Opus, vt docet Ioann: in 1. Epist.3. vers.18. Cor purum) habet charitatis veritatem. Itaque opponitur cordi ( Impuro, quod sordida et foeditates amat et diligit ( Ficto seu duplici, quod non est integrum et purum a simulatione et hypocrisi tum in Deum, tum in proximos.  \pend\pstart Conscientia bona) vsum, et opus illius habet. Nam et finem actionum nostrarum regit, et circunstantias prudenter tractat et perspicit in agendo. Tunc enim est bona conscientia in nobis cùm cor nostrum nos non condemnat, sed coram Deo licet attestari nos et rectum finem habuisse, et diligenter omnia, quae circunspicienda erant, perspexisse, quantum in nobis fuit.  \pend\pstart Fides autem est moderatrix, aut potius regula et norma vtriusque, quae docet ex Dei verbo nos  \pend
\section*{AD I. PAVL. AD TIM. }
\marginpar{[ p.16 ]}\pstart quo spectare debeamus, quid odisse, quid diligere, et quomodo pro cuiusque vocatione nos gerere. Itaque coniungenda est cum verbo Dei. Non vult autem Paulus hanc fidem esse hypocriticam quae speciem tantùm pietatis et metus Dei prae se ferat: sit tamen ab vtroque remotissima.  \pend\pstart Haec fides igitur ad cognitionem ipsius doctrinae refertur, quae in nobis esse debet ἐν ἀληθείᾳ vt loquitur Coloss.I vers. 6.id est, et qualem esse ex Dei verbo oportet, non autem qualis ex nostris aut hominum commentis, et sensu haberi, formarique potest.  \pend
\phantomsection
\addcontentsline{toc}{subsection}{\textit{6 A quibus quum quidam vt a scopo aberrarint, deflexerunt ad vaniloquentiam }}
\subsection*{\textit{6 A quibus quum quidam vt a scopo aberrarint, deflexerunt ad vaniloquentiam }}\pstart Confirmatio superioris sententiae a contrario. Nam qui ab hoc scopo deflexerunt, inciderunt in miseram vaniloquentiam, et inutile studium, quantumuis magni aliorum Doctores esse vellent vel viderentur. Haec vero Pauli sententia pulcherrimam confir mationem recipit non tantùm ex haereticorum doctrina, qui vanitati illi plane dediti fuerunt, relicto verae Theologiae scopo vt Valentiniani Manichei et alii: sed imprimis ex somniis Scholasticorum, qui nunc soli Theologi videri, et dici (si Deo placet) volunt. Hinc apud eos portenta quaestionum, et myriades natae inutilium disputationum, in quibus tota aetas tamen teritur. ac de quaestionibus quidem haec sunt pauca exempla.  \pend
\section*{CAPVT I. }
\marginpar{[ p.17 ]}
\phantomsection
\addcontentsline{toc}{subsection}{\textit{7 Quum velint esse Legis doctores, nec intelligant quae loquuntur, neque de quibus asseuerant. }}
\subsection*{\textit{7 Quum velint esse Legis doctores, nec intelligant quae loquuntur, neque de quibus asseuerant. }}\pstart Descriptio est istius generis hominum, et superioris dicti amplificatio, vt doceat quicunque ab illo scopo aberrant, eos incidere in vaniloquium, quantunuis magni sint doctores. His autem tribuit Paulus Epitheta quatuor, nimirum quod sint I Arrogantes Volunt. 2 Timoth.3 2 Ignari Coloss2. vers.18.3 Docendi munus prae se ferant, et sibi vsurpent Rom. 2.vocantur Ra  \pend
\section*{AD I. PAVL. AD TIM. }
\marginpar{[ p.18 ]}\pstart bi Marth.23. vers.7. 4 Temerarii, assertores et habere se clauem scientiae iactant. Matth. 23. Ex quibus omnibus intelligitur, quam graphice et vere hic nobis a Paulo proposita sit Scholastica Theologia et Mag strorum nostrorum mores et ingenia descripta, ne non facile hodie agnosci possent, nisi plane vel in ipso Meridie caecutire velimus. Denique cùm totam fidei doctrinam corrumpant, nulli tamen et doctiores et sanctiores videri volunt, quam furfures isti et Midae portantes auriculas asini.  \pend
\phantomsection
\addcontentsline{toc}{subsection}{\textit{8 Scimus autem bonam esse Legem, siquis ea legitime vtatur. }}
\subsection*{\textit{8 Scimus autem bonam esse Legem, siquis ea legitime vtatur. }}\pstart Hypophora, per quam occurrit calumniae, qua et ipse, et doctrina Euangelii grauabatur, quasi Dei Legem, quae sacrosancta est, et Prophetarum testimoniis confirmata, sperneret et damnaret Paulus tanquam vanam doctrinam, quia Genealogias, et huiusmodi quaestiones superius reprehenderat. Nec vero vno tantùm loco Paulus cogitur se ab ea graui accusatione defendere, quasi Dei Legem damnaret: sed srpius, vti Rom. 7.v. 12. et Oalat.3.vers.24. Ex quo facile int elligitur, quam valide et astute oppugnata sit Euangelii doctrina ab istis nebulonibus et aërem ipsum subtilitate suarum quaestionum fatigantibus hominibus: quasi manifesto Dei verbo esset contraria Euangelii puri praedicatio. Quod idem nobis a Doctoribus Papistarum obiicitur. Sed ex hoc ipso loco apparet eos qui tunc temporis tam sanctis admonitionibus repugnarent maxime fuisse  \pend
\section*{CAPVT I }
\marginpar{[ p.19 ]}\pstart Iudaeos. Illi enim de lege gloriabantur quam acceperant, et qua etiam abutentes circa huiusmodi quaestiones languescebant. Gentes vero ad fidem conuersae, vix, ac ne vix quidem istas de Genealogiis quaestiones attingebant: vel Legis a Mose datae meminerant, de qua ante Euangelium non audiuerant: neque in ea, vti Iudaei, innurriti fuerant. Praeclara est haec confessio Pauli Apostoli de vsu legis et veteris Testamenti in Ecclesia Christiana et sub Euangelio valde notanda contra Marcionitas, et Manichaeos. Obseruandum tamen est hic Paulum non in vniuersum de toto vsu Legis agere, sed de eo tantùm vsu, qui erat maxime ad praesentem quaestionem accommodatus, et vt vanam illam Iudaizantium doctorum ambitionem et iactantiam contunderet: qui quum nihil, nisi garrulitatem et vaniloquentiam, monstrarent, prae se tamen ferebant magnam Legis obseruationem, praedicationem, ac reuerentiam. Ergo docet quis sit verus Legis, quam praedicant, vsus, per quem Lex ipsa nobis salutaris et vtilis esse possit. Est autem, si ex ea cor purum et sanam conscientiam habere didicerimus per fidem, quae est in Christo.  \pend\pstart Lex bona est, si quis lagitime ea vtatur) Legitimus Legis vsus est, quum ad eum finem ea vtimur, ad quem nobis a Deo data est et promulgata. Est enim is demum rectus illius vsus, qui quis sit supra docuit vers. 5. et 9. Paulus alludit etiam apte paronomasia quadam ad vocem νόμου in verbis vόμος et νομίμως. Nec enim Lex est doctrina vtilis non vtenti illa legitime, et vt oportet. Sed obiicitur, etsi quis non vtatur recte Lege, tamen  \pend
\section*{AD I. PAVL. AE TIM. }
\marginpar{[ p.20 ]}\pstart ipsam legem per se bonam semper esse. Qui enim eam in maiorem peccandi occasionem trahunt, ea legitime non vtuntur, non desinit tamen Lex propterea esse bona. Simile quiddam de Sacramentis afferri potest, quae semper sunt signa salutis per se, quanquam impii vtantur iis ad hypocrisin suam fucandam et tegendam. Respon. Hoc verum esse, sed a Paulo demonstratur hoc loco non tantùm quid res sit ipsa, id est, Lex sit in sese: sed quid in nobis, et qualis eius fructus esse et produci debeat, ne ea abutamur. Erit autem tum demum fructuosa, si ad quem finem data est, cogitemus, id est, vt fidem cum charitate ex ea discamus et exerceamus.  \pend
\phantomsection
\addcontentsline{toc}{subsection}{\textit{9 Nempe sciens Legem iusto positam non esse, sed Legis contemptoribus, et iis qui subiici nesciunt: impiis et peccatoribus, nefariis et profanis, patricidis et matricidis, homicidis. }}
\subsection*{\textit{9 Nempe sciens Legem iusto positam non esse, sed Legis contemptoribus, et iis qui subiici nesciunt: impiis et peccatoribus, nefariis et profanis, patricidis et matricidis, homicidis. }}
\phantomsection
\addcontentsline{toc}{subsection}{\textit{Io Scortatoribus, maseulorum concubitoribus, plagiariis, et mendacibus, periuris, et siquid aliud est quod sanae doctrinae sit oppositum: }}
\subsection*{\textit{Io Scortatoribus, maseulorum concubitoribus, plagiariis, et mendacibus, periuris, et siquid aliud est quod sanae doctrinae sit oppositum: }}\pstart E’ξήγησις. Explicat enim quis sit is legitimus Legis finis, de quo dixit. Est autem vt per Legem hominum mores et vitia corrigantur, et emendentut: non autem vt in vanis quaestionibus occupemur. Posita enim est, vt scelera hominum, veluti homicidia, adulteria, mendacia coerceantur: non autem ad bonos, et qui iam sancti sunt,  \pend
\section*{CAPVT  I. }
\marginpar{[ p.21 ]}\pstart reformandos et reprimendos. Mirum autem hic videri potest, cur cùm illae questiones, quas dannat Paulus, ad Legem caeromonialem et politicam pertineant, tamen de lege morali agat, non autem de caeremoniali aut politica. Respon. Quod illa parte Legis, quae est moralis, istos sibi aduersarios homines magis premi posse videt, quia se pietatis et sanctitatis nomine commendabant. Deinde cùm praecipua Legis diuinae pars sit Moralis, reliquae vero duae sint illius quaedam appendices tantùm, de legis Moralis vsu apte et conuenienter agit Paulus, ex quo et Caeremonialis et Politicae legis vsus quoque pendet et intelligitur. Itaque non extra rem respondet et arguit Paulus. Similis est hodie nostrorum aduersariorum et Monachorum fastus, et arrogantia. qui cùm toti criminibus scateant, nihil tamen, nisi Legem Dei, mandata, et sanctitatem vitae crepant, vt videantur sancti, et nos oppugnant, quasi eadem illa spernamus et damnemus. Quibus respondemus eos esse pares istis hypocritis a Paulo descriptis. Quod ait Paulus, Legem non esse positam iusto, ex eius mente facilem sensum habet, nempe quod qui sponte iam parent Legi, et virtutem sequuntur illi nullis loris et habenis vinciri debeant, nullis aliis praeceptis egeant, quibus ad vitae sanctitatem adducantur: sed illi tantùm iis egeant, qui sunt immorigeri, rebelles, et dissoluti. Haec sententia dignitatem et commendationem Legis non imminuit, quasi iustis et iis qui Doi spiritu aguntur sit vel contraria vel inutilis, sed eos tantùm perstringit, cum quibus egit Paulus, quos facile quidem docet esse tales, quales hic a\pend
\section*{AD I. PAVL. AD TIM. }
\marginpar{[ p.22 ]}\pstart lieno nomine pinguntur, id est, homicidas a dulteros, etc. lusto etiam est vtilis Lex. primum quia viam iustitiae inchoanti perfectionem docet et omnis virtutis est noima. Pelagius putauit eatenus Dei verbum et Legem vtilem esle, quatenus tantùm est liber continens doctrinam morum, contra quem varie disputauit Augustinus.  \pend\pstart Exlegibus autem hominibus esse positam Legem potest videri mirum, sed verum est, quia vt coerceret hominum rebellionem posita Lex est, et exlegibus posita est, ne mancant exleges.  \pend\pstart Hic autem describuntur gradus nequitiae: et primùm fontes ipsi duo, R. bellio, et Contumacia in Legem, ex quibus riuuli isti fluunt, duo scilicet, Impietas ad primam tabulam, et Peccatum, id est offensio proximi ad secundam tabulam pertinens. Hae res paruae primum, postea crescunt, et dissolutiores fiunt, adeo vt tandem euadant homines, Nefarii, id est, publice et nullo pudore peccantes, et Profani. Inde infinita genera vitiorum nascuntur, quae reuocat ad tria genera, Vim Libidi nem, Fraudes siue furando siue mentiendo fiant. Decimum versiculum explicuimus in Ethicis nostris.  \pend
\phantomsection
\addcontentsline{toc}{subsection}{\textit{11 Quae est, secundum gloriosum Euangelium beati Dei, quod concreditum est mihi }}
\subsection*{\textit{11 Quae est, secundum gloriosum Euangelium beati Dei, quod concreditum est mihi }}\pstart Denique easdem quaestiones inutiles docet contrarias esse Euangelio ipsi. Commendat autem Euangelium nomine, Gloriosi a differentia scilicet quae nonest de scriptura, sed de effectu intelligenda, Dei ab authore, quem vocat non  \pend
\section*{CAPVT I. }
\marginpar{[ p.23 ]}\pstart tantùm benedictum, sed beatum. Definit denique illud Euangelium esse eam doctrinam, quae ab ipso et collegis praedicatur Christi mandato. Itaque videri potest argumentum ab enumeratione partium scripturae sumptum. Habet enim Legem et Euangelium, cui vtrique hae quaestiones inutiles repugnant. Opponit igitur Euangelium Dei isti Mataeologiae, quam et a Lege et ab Euangelio damnari docet. Quo ipso demonstrat Paulus consentire inter se vtranque doctrinam contra Manichaeos. Sanam quoque doctrinam hanc vocat, quia aed’ficat ad pietatem. Denique Deum illius authorem μακάριον dicit et quod in se felix sit et beatus, et quod aliis felicitatem donet et afferat. Benedictum alibi vocat eadem scriprura Deum id est dignum omni laude et inuocatione. Paulus autem appellat hoc idem Euangelium 2 Timot 2. vers. 8.suum Euangelium, id est, quod ipsi creditum erat, non cuius ipse sit author.  \pend
\phantomsection
\addcontentsline{toc}{subsection}{\textit{12 Gratiam igitur ei habeo, qui me robustum effecit, id est, Christo Iesu Domino nostro, quod me fidelem duxerit, vt qui me in ministerio constituerit. }}
\subsection*{\textit{12 Gratiam igitur ei habeo, qui me robustum effecit, id est, Christo Iesu Domino nostro, quod me fidelem duxerit, vt qui me in ministerio constituerit. }}\pstart E’'κβασις siue digressio ex iniecta mentione sui ministerii et Apostolatus, quem tuetur et commendat hoc loco. Cur autem excurrat Paulus duplex ratio est. Prima, vt haec sit hypophora contra calumnias, quas inuidi homines illi obiiciebant propter superiorem vitam. Quomodo enim videri potest Euangelii praeco, qui tam atrociter illud antea fuit insectatus. Captatas autem fuisse  \pend
\section*{AD I. PAVL. AD TIM. }
\marginpar{[ p.24 ]}\pstart huiuscemodi occasiones ab aduersariis Euangelii apparet tum ex I. ad Corinth.15.tum quod ad Epist. ad Galat. et hic suum ministerium tuetur. Secunda ratio cur suum Apostolatum Paulus commendet, est, vt Timotheo, quem exhortaturus erat ad difficile certamen suscipiendum, seipsum in exemplum proponeret, in quo summa Dei potentia, vis, misericordia, et virtus illuxit, ne fatisceret, aut expauesceret Timotheus. Ita vero totum suum ministerium asserit Paulus, vt illius gloriam et laudem vni Deo vindicet, non sibi, non bonis suis operibus, aut praeparationibus. Hoc significat Gartiam habeo, Gratiae enim aguntur bene ficiorum nomine acceptorum: non autem meritorum a nobis aut operum collatorum. Christo) Deum esse ex eo intelligimus Christum, cui et vocatio, et virtus conferens dona falutaria et nobis necessaria tribuitur 2 Corinth. 3.vers.5. omnis nostra idoneitas ex Deo est,  \pend\pstart Potentem reddidit) Non dicit Paulus, qui me tantùm iuuit, sustinuit, erexit: sed qui reddidit potentem, vt valerem: alibi ait ἱκάνωσε, et robustum vt persisterem. Ergo prorsus extenuat et deiicit vires suas Paul. quas se ab vno Deo fatetur habere Neque hic de prima sua ad Euangelium vocatione tantùm, sed de toto ministerio suo et Apostolatu agit. Vires enim a Deo nobis suggeri oportet ad vocationes nostras obeundas. Vires etiam ad perseuerandum. Vnde ait idem Paulus collocans me in Ministerio. Hoc idem de nobis sentire debemus, et ex animo quidem, non autem simulate et fict e:nec illud Leonis x. Pontif Rom. blasphemum vsurpare, nos a Deo vocari meritis  \pend
\section*{CAPVT. I. }
\marginpar{[ p.25 ]}\pstart licet imparibus.  \pend\pstart Fidelem duxerit) Illud ducere est efficere. Non inuenit enim nos Deus tales, sed efficit. Nec enim idcirco Paulum vocauit Deus quod fidelem inuenerit I. quia erat blasphemus 2. quia frustra superius dixisset Robustum me fecit. Erat enim dicendum, Robustum inuenit 3. Deinde et quod illa fide litas est idoneitas, quae non est a nobis, sed a Deo. Et seipsum superioremque sententiam plane euerteret Paulus, si quid in eo praecessisse boni oPeris fingamus, propter quod ipse a Deo sit vocatus.  \pend
\phantomsection
\addcontentsline{toc}{subsection}{\textit{13. Qui prius eram blasphemus et persequutor et iniuriosus: sed mei misertus est. nam ignorans id faciebam fidei expers. }}
\subsection*{\textit{13. Qui prius eram blasphemus et persequutor et iniuriosus: sed mei misertus est. nam ignorans id faciebam fidei expers. }}\pstart Comparatio est, ex qua Dei gratia maxime amplificatur, et maxime quod aduersarii in calumniae causam arripiebant, transfert in suum conmodum, et sui ministerii maiorem commendationem. Accusabant igitur quod fuisset Paulus antea ( Blasphemus Insectator ( Violentus Quae omnia concedit ille, quia et vera erant, et sic Dei gloria maior, et sua vocatio ad Apostolatum illustrior apparet. Cùm enim talis fucrit, mirum est vocatum tamen a Deo esse Actor.9. et 13. Non loquitur de se Paulus vt hypocritae, sua peccata minuens, sed ea exaggerat. est enim haec gradatio: Fui blasphemus. i. maledicere solitus Deo. Deinde persequutor. Denique rabiosus et contu\pend
\section*{AD I. PAVL. AD TIM. }
\marginpar{[ p.26 ]}\pstart melias omnes coniungens. Mirum Dei opus in Paulo vocando, vt de multis persecutoribus verbi Dei non sit hodie desperandum. Opponit autem Dei mifericordiam, per quam non tantùm veniam superiorum vitiorum consequurus est, sed etiam pulcherrimam illam ad Apostolatum vocationem adeptus, omnibus suis sceleribus ne futuris operibus vel quae Deus post Pauli vocationem edenda ab eo praeuiderat, electionis huius causam tribuamus.  \pend\pstart Addit tamen rationem, quia feci, et ignorans et ἐν ἀπιςτίᾳ. id est, carens fide et cognitione. Non tantùm igitur docet se ignarum fuisse prius: sed se de fide et verbo Euangelii quicquam antea audiisse vel fuisse edoctum negat. Hic igitur ἀπιςτία est species ignorantiae, quae est maxime excusabilis, quum quis non ignorat ea, de quibus oblata illi est cognitio: sed de quibus nunquam audiuit. Ergo ἀπιςτία non est hoc loco infidelitas et incredulitas cum quadam malitia et animi obstinatione coniuncta: sed simplex et praecedenti verbi Domini praedicatione carens. Vnde quaesitum est vtrum ignorantia excuset peccatores a peccato. Quidam distingunt inter ignorantiam Simplicem et Malitiosam. Alii inter ea quae ignorantur. Aut enim sunt vt illi aiunt, praecepta pertinentia ad Primam Legis tabulam, aut ad Secundam. Illa ignorare non conceditur, vt volunt. Haec autem licet, quae sunt Secundae tabulae. Res. Ignorantia excusat non a toto: sed a tanto Rom. 2.vers. 12. Nam si caecus caecum duxerit ambo in foueam cadent: et qui ducit, et ipse qui ducitur. vt est Mat th. i5, et apparet, Genes. 2o. vers.5 et7.  \pend
\section*{CAPVT I. }
\marginpar{[ p.27 ]}\pstart Hinc autem quidam colligunt definitionem peccati in Spiritum sanctum: sed non satis apte.  \pend
\phantomsection
\addcontentsline{toc}{subsection}{\textit{I4 Superabundauit autem gratia Domini nostri cum fide et dilectione in Christo Iesu }}
\subsection*{\textit{I4 Superabundauit autem gratia Domini nostri cum fide et dilectione in Christo Iesu }}\pstart Αὔξησις Dei gratiae, quae vice confirmationis addita est, et cx verbis ipsis Pauli apparet, vt omnis superioris vitae labes deleta esie doceatur. Probat autem hoc ipsom etiam ex effectis contrariis superiori illi vitae, nempe ex fide, quae opponitur blasphemiae priori: et ex Charitate in Christo, quae opponitur Persecutionibus, et Contumeliis prioribus in Dei Ecclesiam factis.  \pend
\phantomsection
\addcontentsline{toc}{subsection}{\textit{15 Fidus est hic sermo, et modis omnibus dignus quem amplectamur, Christum Iesum venisse in mundum vt peccatores seruaret, quorum primus sum ego. }}
\subsection*{\textit{15 Fidus est hic sermo, et modis omnibus dignus quem amplectamur, Christum Iesum venisse in mundum vt peccatores seruaret, quorum primus sum ego. }}
\phantomsection
\addcontentsline{toc}{subsection}{\textit{16 Verum ideo misertus est mei, vt in me primo ostenderet Iesus Christus omnem clementiam: vt essem exemplo credituris ipsi in vitam aeternam. }}
\subsection*{\textit{16 Verum ideo misertus est mei, vt in me primo ostenderet Iesus Christus omnem clementiam: vt essem exemplo credituris ipsi in vitam aeternam. }}\pstart Αʹνάκλασις eaque summo illustrandae Dei gloriae studio inducta, vt quod contra ministerium Pauli maxime videbatur antea obiici posse, hoc totum ad pios consolandos confirmandósque iam valeat, quum in Paulo, tanquam pulcherrimo quodam speculo, insigne et singulare Dei misericordiae exemplum etiam maximi quoque pec\pend
\section*{AD I. PAVL. AD TIM. }
\marginpar{[ p.28. ]}\pstart catores habeant. Ergo primùm se describit: deinde ex aduerso Dei misericordiam confert. Demum caeteros suo exemplo consolatur et exhortatur ad idem sperandum. Magnum autem huius spei et certum fundamentum primum omnium ponit, nempe, Christum factum esse hominem, vt peccatores nimirum redimeret, etiam sceleratissimos. Quod sundamentum, quia magni est momenti, conmendat solita praefatione, fidus est sermo et dignus infr. cap.3v.1.2. Timoth.2.V.19. Quanquam enim vniuersa scriptura sacra est per se, et in omnibus suis partibus amplectenda, et sermo fidus, tamen imprimis illa, quae summa fidei nostrae saluiis et spei capita complectitur, singulari quodam modo amplectenda est. qualis haec et inf. cap.3.V.16. Caeterum Christumvenisse, vt saluos faceret homines, est certissimum. idcirco.n. dicitur Iesus Math.1. et est hocipsum αἴτημα θεολόγικoν quod nulli Chri stiano esse ignotum potest nec debet, ideóque  diligem tius et retinendum et commemorandum. Vnde fit, vt mirum videri non debeat, si qui talis antea fuit Paulus, misericordiam iam est consequutus. Ea enim gratia factus est homo et Iesus Dei Filius Christus. At fuit maximus peccator Paulus: fatetur id quidem ipse, non tantùm modesta quadam sui deiectione vtens: sed vere et ex serio sensu animi sui loquens, et quali a nobis infidelitas omnis maximeque obstinata debet apprehen di cum pauore ac horrore conscientiae, Et hoc suum tam graue, et grande peccatum esse illustrissimum summae Dei misericordiae erga omnes peccatores speculum decet, ne desperent vlli, sed conuertentur ad Deum, et in eo spem vitae aeterng ha\pend
\section*{CAPVT I. }
\marginpar{[ p.29 ]}\pstart beant. Ex hoc loco apparet quam sit enorme scelus persequi Christi doctrinam et seruos, quanquam homines zeli velo conantur excusare tam impium facinus. Persequuntur enim Christum ipsum, quo scelere nullum esse maius potest Act. 9. quanquam ex iis tamen ipsis alii aliis grauiùs peccant:Ii grauius, qui scientes veram esse eam doctrinam nihilominus insectantur. Leuiùs autem qui eam esse Dei doctrinam ignorant, inter quos fuit Paulus. In se vero primo dicit patefactam Dei clementiam illam ingentem, quod non est accipiendum, ratione temporis. Omnes enim sumus natura filii irae, et neminem ante Paulum vocauit ad se Christus, aut hodie vocat, qui non sit natura illius hostis et persequutor, si votum spectes: Sed se primum vocat aliorum comparatio ne, et sceleris sui detestatione, id est, insignem peccatorem et caeteris capitaliorem Ecclesiae Dei Definit autem fidem et a subiecto, circa quod versatur. Est autem Christus: et a fine, quem spectat vel potius sperat. Est autem vita aeterna.  \pend
\phantomsection
\addcontentsline{toc}{subsection}{\textit{17 Regi autem aeterno, immortali, inuisibili, soli sapienti Deo honor, sit et gloria in secula seculorum. Amen. }}
\subsection*{\textit{17 Regi autem aeterno, immortali, inuisibili, soli sapienti Deo honor, sit et gloria in secula seculorum. Amen. }}\pstart Epiphonema, id est, sententia summi ardoris et vehementiae plena, in lquam ex diuinae erga se misericordiae meditatione erumpit Paulus, quoniam zeli sui feruorem, et suae gratiarum actionis silentium comprimere iam non possit amplius. Itaque effundit hanc praeclarissimam vocem de stupenda Dei bonitate, primùm quidem erga se  \pend
\section*{AD I. PAVL. AD TIM. }
\marginpar{[ p.30 ]}\pstart patefacta: deinde quae cernitur in caeteris rebus omnibus. Quae nanque de Deo hoc loco dicuntur, illi semper competunt, quanquam ad Pauli conuersionem optime quadrant, in qua Deus suae bonitatis, sapientiae, gloriae, et immutabilitatis specimen eximium ostendit. Ergo si quid mirum et incredibile inest ac occurrit in tanta PauIi conuersione, ne mirentur. Nam est Deus qui conuertit, Deus, inquam, ille rex, et aeternus, qui potest omnia, et ab aeterno fuit. Si quid alienum, et nouum apparet in eadem factum conuersione est Deus ille sapiens. Si quid praeter humanam rationem extitisse cernitur, idem ille Deus est inuisibilis, vt qui alia, quam carnali ratione operatur et exequitur sua consilia, sua beneficia confert, et Euangelium suum prouehit. Itaque haec omnia Epitheta sunt rot ἀιτιολογίαι, ne fictitium hoc esse videatur, et impossibile foctu quod Paulus de se narrat, et de Dei misericordia in conuertendis sceleratissimis.  \pend\pstart Similia Epiphonemata passim apud Paulum occurrunt, vti Romanor. 9. vers.5. et infra 6 v. 16  \pend\pstart Quatuor autem Epitheta Deo tribuit, Vocat enim Regem seculorum, immortalem, inuisibilem, sapientem et quidem solum. Alia quoque alibi. Deo nomina tribuuntur, sed hic selegit Paulus, quae magis circunstantiis sui facti et vocationis conueniebant, Romanor. 1.V.2o. et Exod.34.v. 67. et passim in Psalmis et haec, et alia leguntur.  \pend\pstart Rex seculorum) dicitur quia condidit tempora non tantùm haec quae dicuntur Αʹιώνιά  Tit.1.V.2 et a rerum conditarum siue creatarum initio ceperuntised illa quae fluxerunt ante mundi crea\pend
\section*{CAPVT I. }
\marginpar{[ p.31 ]}\pstart tionem, si modo nobis liceat bona Peripateticorum venia illam durationem tempus appellare, vt docet Augustinus in liber  Confe.11. Est enim Deus ante omnem temporis non tantùm durationem: sed etiam cogitationem. lstud tamen Epitheton quidam de Christo proprie scribi putant, quia Pater seculi nominatur Isai.9. vers. 6.sed alius est illius loci sensus. Idcirco autem Deus. i. Trinitas ipsa appellatur Rex seculorum. vt immutabilis et aeternus Deus esse intelligatur, qui nullum temporis principium habeat: Cui subsequens est, vt sit immortalis. Hoc enim significat ἄφθαρτος. Nec enim ad mores incorruptos pertinet haec vox, vt nonnunquam accipitur, sed finem temporis excludit, que madmodum aeternitas principium ὑπάρξεως.  \pend\pstart Inuisibilem) vocat iuxta illud Christi, Deum nemo vidit vnquam Ioan. 1.vers.18.id est, humano visu Deus conspici non potest, vt explicat Augustinus liber  I.de Trinitate cap.6. Beati quidem Deum vident et videbunt, sed alia ratione et visione, quam haec nostra est. Hoc de solo Patre quidam dici putant, quia Filius per carnem assumptam factus est nobis conspicuus Deus, quemadmodum etiam loquitur Paulus infra 3. vers.16. At nec in Christo ipsa deitas fuit conspecta, imo velo carnis fuit obducta et obtecta. Praeterea dicit Christus ipse, qui me videt, videt et Patrem Ioa. 1a v.9 vt ad Filium quoque hic locus referri debeat, quatenus Deus est.  \pend\pstart Sapientem) vocat non eo modo, quo nos homines saepe sapientes dicimur, sed incomprehensibili prorsus sapientia, qua solus sapiens est et  \pend
\section*{AD I. PAVL. AD TIM. }
\marginpar{[ p.32 ]}\pstart quae quemadmodum. ait id I Corinth 2. sapientia repugnat et aduersatur mundi et hominum sapientiae, vt si quid in ipsius operibus non comprehendimus, obmutescamus, et Dei consilium immensum adoremus, non autem in temperie quadam animi agitati obstrepamus Deo, quod multi faciunt.  \pend\pstart Soli) Quaesitum est vtrum haec vox ad Deum, an ad vocabulum sapienti sit referenda: quasi hic Deus vel solus Sapiens, vel solus Deus appelletur qui est aeternus, immortalis, inuisibilis. Variae sunt interpretationes. Ego vero potius existimauerim illud coniungendum esse cum proxima voce, Sapienti. Hic enim miram Dei sapientiam praedicat Paulus, quae nulli homini satis nota esse potest, per quam e miserrimis et desperatissimis peccatoribus, et Euangelii persecutoribus praecones, pastoresque Ecclesiae suae facit mutans eos mira sua potentia et virtute, vt ex lupis saeuissimis fiant oues placidissimae Isai.11.  \pend\pstart Honorem autem et gloriam tribuit Deo. Alibi alia quoque tribuuntur, vti in Epist. Iudae in fine et 2. Petri 3.vers.18. quae duo videntur ita separari et distingui posse, vt τιμὴ sit ipsius maiestatis Dei excellentia, splendor, admiratio, et virtus: δοξα vero illius tantae virtutis praedicatio, praeconium, et celebratio, quae fit ore hominum, vel Angelorum. Prius enim Dei maiestas agnoscenda est, post praedicanda.  \pend\pstart Quaeritur autem. Vtrum haec de solo Patre, an etiam de Filio dicantur, et de Spiritu sancto, quia Arriani ad solum Patrem haec per tinere volunt, quem vnum proprie quoque contendunt  \pend
\section*{CAPVT I. }
\marginpar{[ p.33 ]}\pstart esse Deum Chrysostomus autem ad Filium hic refert. Respond. Si verba spectentur, ad Patrem referri, cui supra quae sunt propria diuinitatis tribuerat Paulus: si tamen res ipsa definiatur etiam ad Filium, et Spiritum sanctum referri, quia paris sunt Filius et Spiritus sanctus cum Patre dignitatis: et confirmatur 2 Petri 3. vers. 18. nostra sententia. Quia vero superio. versiculo Christum nobis vti mediatorem, et hominem Paulus proposuit, ex more suo et scripturae, in Patris persona complectitur et restringit, quae sunt propria diuinae naturae. Est vero periculosa Erasmi annotatio ad hunc locum, qui id rarum esse in Apostolicis literis ait, vt Dei vocabulum Christo et Spiritui sancto tribuatur. Ac ne in symbolo quidem Apostolico, Filium et Spiritum sanctum Dei nomine concedit proponi nobis. In quo toto coelo errat. Nam qua ratione Patrem confitemur esse Deum, cùm dicimus Credo in Patrem: eadem etiam Filium agnoscimus Deum, quia sic enuntiamus, et in Filium: idem de Spiritu sancto et in Spiritum sanctum, praeterquam quod ita pronuntiatur prima illa pars Symboli Credo in Deum, vt illic dispunctio adhibeatur Credo in Deum, et susistat vox: deinde explicatur quis sit ille Deus, nempe Pater, Filius, et S sanctus distincte.  \pend
\phantomsection
\addcontentsline{toc}{subsection}{\textit{18 Hoc mandatum commendo tibi, fili Timothee, nempe vt secundum praegressas de te prophetias, milites per eas bonam illam militiam. }}
\subsection*{\textit{18 Hoc mandatum commendo tibi, fili Timothee, nempe vt secundum praegressas de te prophetias, milites per eas bonam illam militiam. }}\pstart Regressio est ad propositum, et huius Epistolae  \pend
\section*{AD I. PAVL. AD TIM }
\marginpar{[ p.34 ]}\pstart proprium argumentum, quod digressione illa superiori de vocatione miraculosa Pauli fuerat interruptum. Ergo postquam a quibus vitiis Pastori Ecclesiae sit cauendum docuit, adiicit quae sit summa totius muneris ipsorum, nempe et vt Fidem sanam habeant, et Conscientiam bonam et integram. Fidei nomine doctrinam vniuersam intelligit, quae ex puro Dei verbo docetur in Ecclesia:imprimis autem Euangelium significat, quae pars Scripturae fidem nobis proprie proponit, Bona conscientia ad ipsius pastoralis muneris executionem pertinet. Facit enim illa, vt syncere recte et integre, quemadmodum oportet, in toto et vitae et vocationis nostrae cursu versemur, et nauigemus, donec ad portum peruenerimus. Est autem is portus mors. De vtraque sup. copiose vers.5.  \pend\pstart Ex hoc autem loco facile intelligi potest, ecquod pastoris munus esse velit Paulus. Id quod ipse Tit.1.vers.9. amplius docet et Petrus I Pet.5. vers.2. Nempe vt pascant gregem Domini verbo Dei, quod proprium est animorum alimentum. Ad quod et doctrina pura, siue sana cognitione illius verbi opus est: et prudenti synceraque eiusdem dispensatione, et ὀρθοτομίᾳ vt loquitur Paulus ipse inf. 2 Timoth. 2 id est, ne sit personarum acceptator Minister verbi Dei: sed exactor a quoque officii, et cultus Dei. Quod si idem tunc de officio Episcopi receptum et constitutum fuisset, quod postea et in Papatu constitutum est, certe neque  doctrinam neque bonam conscientiam commen daret et requireret in Episcopo Paulus. Quanquam enim Episcopum explorandum et exanimandum illi  \pend
\section*{CAPVT I. }
\marginpar{[ p.35 ]}\pstart seribunt cano. Qui Episcopus ordinandus dist.25. tamen cùm munus Episcopi describunt et definiunt eo modo illi loquuntur, vt nihil plane cum hoc Pauli loco conueniat in cano. Perlectis dist.25 Ad Episcopum pertinet Basilicarum consecrati o Vnctio altaris, consecratio Chrismatis: ipse praedicta officia et ordines Ecclesiasticos distribuit: ipse sacras virgines benedicit, etc. Quae mera sunt tam sanctae dignitatis ludibria, nimirum tam procul a vero Episcopatus sine in Papatu receslum est.  \pend\pstart Commendat autem Paulus hanc suam, quanquam breuem pastoralis muneris descriptionem duplici argumento, Generali, et Speciali, quo diligentiùs et a nobis omnibus, et a Timotheo excipiatur et obseruetur diligentiùs. Generalis quidem ratio est quae a genere sumpta est, dum Praeceptum appellat, ergo non negligendum impune, et Depositum. Specialis ratio commendationis spectat Timotheum, et tum a suo in eum amore, tum ab ipsius officio ducitur, qui Prophetias de se ha beat plures Praecedentes, vt certus sit de Dei auxilio, et sua vocatione, Quod ad Prophetias, quidam, vt Chrisostomus, eas accipit de ordinaria via et electione, quam obseruatam esse vult in Timothei electione: alii de extraordinaria accipiunt, qualis illa, quae clara Dei voce facta est Actor. 13. Tale de Timotheo extitisse Dei testimonium volunt. Actor. 16 vers. 2. Ego extraordinariam electionem accipio hoc loco.  \pend\pstart Vt milites per eas) Explicat Paulus qui modus sit Euangelici ministerii fideliter exequendi, ninimirum, durum saepe certamen subeundum esse  \pend
\section*{AD I. PAVL. AD TIM. }
\marginpar{[ p.36 ]}\pstart sed in eo debere Timotheum sustentari bona spe quae Dei promissionibus generalibus, et praecedentibus illis Prophetiis incumbat. Ministerii vero Euangelici cursum comparat militari certamini, nec hic modo, sed in 2 Timoth. 2.vers.3 4.v. 7. vti et Deus ipse se comparat duci. Exod.15.V.5. vt intelligatur qualis sit Ministrorum conditio, nimirum dura et certaminum aduersus haereses et vitia hominum plena. Neque vero hoc facit, vt Episcopi militaribus se negotiis immisceant. Iure enim damnantur, qui arma tractant, Episcopi. Nam nemo militans Deo, se implicare debet illis negotiis 2 Timoth.2. vers.4 Hoc enim genus militiae est spirituale, non carnale, et arma Episcoporum sunt spiritualia:nempe, preces, lachrymae, verbum Dei:non autem carnalia, qualis est gladius bombarda, etc.  \pend
\phantomsection
\addcontentsline{toc}{subsection}{\textit{19. Retinens fidem et bonam conscientiam: qua repulsa, nonnulli fidei naufragium fecerunt. }}
\subsection*{\textit{19. Retinens fidem et bonam conscientiam: qua repulsa, nonnulli fidei naufragium fecerunt. }}\pstart Continuatio est: si militudine enim docet quantum sit erroris huius, qui in ministerio Euangelico committitur, imprimis autem in bonae conscientiae desertione periculum. Est enim fidei naufragium. In quoet magnitudo periculi, quia fit iactura fidei: et terror ostenditur atque metus, quia naufragium appellatur. Est igitur bona conscientia velut arca fidei et sanae doctrinae, ac intelligentiae thesaurus: tanquam illius cynosura, a qua qui deflectunt, se et nauem, id est, intelligentiam illam sanam perdunt, iusto Dei iudicio in eos grassante  \pend
\section*{CAPVT I. }
\marginpar{[ p.37 ]}\pstart perciti. Ex ea igitur scilicet bona conscientia torum vitae nostrae cursum regere debemus. Quod autem bonam conscientiam homines Ethnici commendant, quale est illud, Nil conscire sibi nulla pallescere culpa. Hic murus aheneus esto. et, Conscientia mille testes. item, Bono viro nihil praeter culpam praestandum, et caetera huiusmodi: vim quidem testimonii conscientiae demonstrant, quanta sit: praecipuum tamen periculum, quod ex ea spreta metuendum est, non ostendunt, quod est hoc, amissio et euersio verae fidei doctrinae, et cognitionis de nostra salute et Dei erga nos misericordia.  \pend
\phantomsection
\addcontentsline{toc}{subsection}{\textit{20 Ex quibus Hymenaeus et Alexander: quos tradidi Satanae, vt discant non blasphemare. }}
\subsection*{\textit{20 Ex quibus Hymenaeus et Alexander: quos tradidi Satanae, vt discant non blasphemare. }}\pstart Hoc ipsum duobus exemplis confirmat, Hymenaei et Alexandri, de quorum altero agit hic ipse in 2 Tim. 2.V. 17. cuius errorem qualis fuerit describit, in quem Alexandrum quoque  incidisse verisimile est. Fuit autem hic Alexander Ephesius Act. 19.V. 13. Vtrunque autem ad terrorem caeterorum pseudopastorum punit Paulus excommunicatione quam solenniter antea in Ecclesia factam fuisse verum est. Hic autem eam repetit, vt intelligant omnes, quo iudicio digni sint omnes haeretici obstinati. Verba autem haec tradidi Satana quid significent, D. Beza in I Corinth.5. vers.5. explicuit, ad quem lectorem remitto. Nam illic melius explicatur, quam in canon. Audi denique. De forma vero Papisticae excom municationis  \pend
\section*{AD I. PAVL. AD TIM. }
\marginpar{[ p.38 ]}\pstart vbi duodecim Sacerdotes circunstant Episcopum, et lucernas ardentes habent, quas in terram proiiciunt. vide canon. Debent duodecim Sacerdotes 1I. quaest.3.  \pend
\endnumbering\beginnumbering\section{CAP. II.}
\phantomsection
\addcontentsline{toc}{subsection}{\textit{\huge\textbf{A}\normalsize DHORTOR igitur ante omnia vt fiant deprecationes, preces, postulationes, gratiarum actiones pro quibus uis hominibus. }}
\subsection*{\textit{\huge\textbf{A}\normalsize DHORTOR igitur ante omnia vt fiant deprecationes, preces, postulationes, gratiarum actiones pro quibus uis hominibus. }}\pstart Transitio est ad alteram Euangelici muneris partem, quae est publicarum precum pro quolibet hominum genere conceptio et profusio coram Deo. Quae disputatio commode subiicitur tum ad susceptum argumentum, quia pastorem Ecclesiae de omnibus sui muneris partibus Paulus informat:tum etiam ad superiorem conclusionem, quia hoc ipsum est optimum pietatis Christianae exercitium, et verae fidei conseruandae modus tutissimus. Precibus enim impetramus a Deo, vt in sana ipsius doctrina et metu pergamus confirmemur, et crescamus. Hoc autem loco et capite Paulus breuiter complexus est, quaecunque de oratione siue precibus Christianorum quaeri possunt. I Quae sit, et quot eius genera. 2 A quibus sint preces fundendae. 3 Pro quibus 4 Vbi siue quo in loco. 5 Quando. 6 Quomodo et qua cum animi reuerentia, et totius coetus decoro: quae singula suis locis a nobis (Domino  \pend
\section*{CAPVT II. }
\marginpar{[ p.39 ]}\pstart dante) explicabuntur.  \pend\pstart Loquitur autem imprimis Paulus de publicis Ecclesiae precibus, non de singulorum Christianorum priuatis precationibus, quae fiunt a quoque domi et in conclaui. Nec enim has in domibus singulorum fidelium pastores facere debent: vel sufficiunt.  \pend\pstart Sed primo loco quaesitum est. Quot sint pastoralis curae et muneris partes, vt intelligatur quam recte et ordine omnia persequatur Paulus Respondent Canonistae esse tres. Astaria siue basilicas totas consecrare: virginibus benedicere: Ecclesiasticos ordines distribuere, vt diximus supra ex canon. Perlectis S Ad Episcopum Distinct. 25. Haec Euangelici ministerii definitio prorsus inepta est, et vana, cùm nihil huic simile doceat scriptura. Quare illi ipsi se reuocant, et alibi corrigunt, et melius sapiunt tota distinct. 85. Nam docent Episcoporum munus esse vt doceant corripiant, liberales sint. Priora duo reipsa pars sunt pastoralis muneris: tertium et postremum non item, sed illi cum omnibus Christianis commune. in Epist. vero ad Hebraeos cap.5.vers.1, 8. vers.3. 9. vers. 24. dicitur fuisse officium Sacerdotis offerre, et precari pro populo. Paulus in Epist. ad Titum 1.vers.9. inf. cap.4. vers. 11. 2 Timoth.4. vers. 2. ait, ministerium Euangelicum positum esse in eo, vt pastores doceant sanam doctrinam, refutent falsam, denuntient et corripiant siue increpent aberrantes in Ecclesia. Hic autem addit etiam ad eosdem pastores pertinere, vt ipsi precentur publice pro tota Ecclesia, et praecipuis illius membris. Ratio est, I. quia, vt ait Chrysostom.  \pend
\section*{AD I. PAVL. AD TIM. }
\marginpar{[ p.40 ]}\pstart Pastor est velut communis totius coetus pater, qui omnium curam gerit.2. Quia est publicum totius Ecclesiae os, vt nos in doctrina, et iis quae huic connexa sunt, erudiat, qualis est oratio. Ergo et precari debet, quia eo precante tota Ecclesia orat. 3 Hoc fuit officium olim sacrificatoris, in quo nulla caeremonia continetur. Hebr. 9. vers. 24. Numeror. 6.vers. 23. 4 Denique mos veteris Ecclesiae, quae Dei spiritu regebatur, in qua pro toto coetu, et ipsius nomine Pastores fundebant preces Deo, et populo benedicebant, vt ex variis historiis constat, et annotat Augustinus Epist.9o.  \pend\pstart Ergo cùm sint haec tria praecipua pastoris munera et officia, Docere, Corripere, et Precari, complexus est Paulus, Docendi verbo etiam correptionem. Precationem autem hic seorsim persequitur, quia est ministerii Euangelici pars praecipua, vti diximus.  \pend\pstart Magnam autem vim esse huius publicae precationis, et magnam illius habendam esse curam a pastore, docent haec verba Pauli et commendatio, ante omnia πρῶτον πάντων. Est enim comparatio, et commendatio huius exercitii, tanquam omnino necessarii in Ecclesia Christi, et vtilissimi, denique retinendi. Neque tamen haec verba confirmant, vel iuuant Audaeorum siue Messalaniorum haereticorum errorem, qui puppim et proram salutis in precationibus, id est, certis statisque horis factis orationibus et de murmurationibus (quod hodie faciunt Monachi Papistici) ponebant: sed tantùm docet Paulus, quam debeamus sedulo et diligenter precari Deum: et quam solicitum de eo debeat esse non modo pa\pend
\section*{CAPVT II. }
\marginpar{[ p.4I ]}storis, sed singulorum etiam fidelium studiym. Preces enim publicas quidam impiissime, etiam hodie contemnunt. \pstart Ait autem Paulus, Adhortor, Quae vox negligentiam nostram perstringit, et arguit. Quanquam enim et Dei ipsius maiestas, et promissio iussióque: item fructus ipse, quem ex precibus nostris vberrimum sentimus, satis nos ad precandum excitare deberet, sumus tamen natura nostra mire ad orandum frigidi, et torpentes. Itaque  exhortandi et euigilandi sumus. Christus ipse nos ad petendum et precandum adhortatur. Petite et inuenietis, Matth.7. vers. 7. Vae igitur somnio pigritiaeque nostrae etc. Luc. 17  \pend\pstart Ut fiant deprecationes) Primum praeceptum ponit quod est, faciendas esse preces in Dei coetu cui secundum adiungit, nempe pro omnibus: ex quibus tertium iam facile colligi potest, A quo fieri debeant: nempe, a pastore. Nam agit (vti diximus) de publicis precibus. Sed ex hoc ipso loco definiri potest, Quid sit oratio, et quot illius sint genera siue species. Precandus est igitur Deus quod extra dubium est, et concedunt non modo Christiani homines: sed etiam Ethnici et profani, vt et Tertullian. in libro de Testim. conscientiae, et Iustinus Martyr in liber  de Monarchia Dei copiose docent: item, Iuuenalis, Persius, et Horatius et alii poëtae id docent. Homerus, et Virgilius, qui suos inducunt preces fundentes ad Deum, et eum inuocantes. Nec gens vlla est, quae si testimonium et vocem propriae conscientiae audiat, hoc axioma neget esse verum, vel refutet.  \pend\pstart Vult autem Paulus fieri deprecationes, preces  \pend
\section*{AD I. PAVL. AD TIM. }
\marginpar{[ p.42 ]}\pstart postulationes, et gratiarum actiones. In quo primum est obseruandum non temere tot voces esse simul congestas, sed partim vt precandi assiduitatem et studium inflammet, incitet et iniiciat nobis his tot verbis, nósque excitet, quemadmodum dixi: partim vero vt diuersa esse precationum genera ostendat, quemadmodum diuersi sunt hominum affectus, diuersa hominum conditio, diuersa nostri status ratio, et ita pastor scite accommodet quae sunt cuiusque orationis generis propria. Denique, vt ostenderet Paulus, regum, et principum, caeterorumque hominum commemorationem in nullo precum genere omittendam esse, tot species commemorauit, itaque omnes negligentiae tollit excusationes.  \pend\pstart Cap.1. Quaeritur autem primum quomodo hae voces differunt δέησις, προς ἀχὴ  , ἔντάξις, et Εὐχαριστία Quod vt commode fiat primùm explicetur et definiatur, quid sit oratio in genere  \pend\pstart Est oratio (ait Damascenus) πρεπόντων αἴτησις, id est, Decentium postulatio petitióque a Deo. Igitur complectitur res, quae sunt a Deo postulandae. Illas enim solas fas est, et decet nos petere a Deo. Dicitur autem a Latinis Oratio ab ore, quia ore fit. Nec enim orationem vocabant quae intus et corde tantùm conciperetur, non etiam ore proferretur. Vnde orationis vim in dicendo maximam et in gestu esse respondent. Varro liber  5. de lingua Latina. Hebraei vocant תפלה a iudicando, quia deiici nos coram Deo necesse est, si vere precari velimus, id quod ex animo sensuque  nostrae inopiae facile fiet. Definitur etiam oratio Dei veneratio animi nostri vota et affectus illi  \pend
\section*{CAPVT II. }
\marginpar{[ p.43 ]}\pstart pandens, vt auxilium ab eo impetremus. Id colligitur ex Tertulliano, et Psal. 142.vers.2  \pend\pstart QQuotuplex autem sit oratio quaeritur, Resp. Est duplex, Implorans eaque  pro nobis, vt Bona det Deus, Mala auertat et auerruncet: vel pro aliis. Et Agens gratias, Est Εὐχαριστία Psal.22.  \pend\pstart Quae inuocat et petit αἴτησις appelletur nomine generali, et oratio siue ἀχὴ  Quae vero pro nobis ipsis postulat a Deo bona προς ἀχὴ  dicitur, vt Psal 61. Quae petit vt a nobis auerruncentur et abigantur mala, δένσις vti Psal. 14o.  \pend\pstart Quae pro aliis interpellat et intercedit apud Deum est ἔντάξις, vel μεσιτεία. Haec tot sunt genera precationum publicarum. Vide tamen August. Epist. 59.  \pend\pstart Cap. 2. Quaeritur autem a quo fieri preces debeant et possint. Resp. Ab omnibus in vniuersum: sed priuatae a priuatis. Publicae vero a solis iis qui Ecclesiae praesunt, quales pastores et olim Diaconi. Ac quoad priuatos scribit Chrysostom. iam olim fuisse solitos Christianos bis in die, scilicet, mane, et vesperi domi precari, adeo vt milites ipsi in castris id consuessent, et data est illis a Constantino breuis formula et precatio, quam mane et vesperi dicerent. Eusebius liber 4. de vita Constantini. Sed de publicis precibus hic agimus.  \pend\pstart Primùm igitur Pastores debent precari publice. Ad eos enim haec res pertinet. Sed etiam ex veteri Ecclesiae disciplina Diaconi possunt et Presbyterii etiam ii, qui doctrinam non tractant. Eorum enim muneri preces adiunctas fuisse scriptum est, in canon Perlectis 5 ad Diaconum dist. 25  \pend
\section*{AD I. PAVL. AD TIM. }
\marginpar{[ p.44 ]}\pstart et apud Socratem Scholast. lib, 2. Histor. cap.11. etiam praesente Episcopo, quanquam eodem praesente sacramenta administrare Diaconus non potuit, vt est in cano. Peruenit distinct 93.nisi ab Episcopo iussus esset.  \pend\pstart Quod autem obiici potest de Solomone, qui, adstante summo sacrificator, etamem preces publi cas pro populo fudit, vt est 1. Regum 8. vers.14. non obstat. Hoc enim semel tantùm factum est, et quidem extraordinarie. Nam Azarias rex et successor Solomonis, qui in Sacerdotii munus irruere voluit a Deo lepra percussus est, et segrex factus 2. Reg. 26. vers.18.  \pend\pstart Cap.3. Quaeritur pro quibus fieri debeant preces Resp. pro omnibus, id est, quouis hominum genere, religione, sexu, aetate, conditione. Sic enim vox πάντες hoc loco sumitur, vti et postea quemadmodum doctissimus Beza obseruauit. Ratio est, Quod homines quosuis pro Dei creaturis, et pro proximis nostris agnoscero debeamus. Ergo omnium vita, salus, et conditio nobis curae debet esse, nec quisquam negligendus, quantunuis pauper, humilis et abiectus.  \pend\pstart Excipiunt tamen quosdam Scholastici, imprimis inimicos nostros, pro quibus non putant specialiter orandum esse. Communia enim humanitatis officia tantùm illis tribui oportere: non autem haec tam egregia, quae solis amicis debemus, qualis est oratio. Durand. in liber 3. Sentent. distin. 3o. quaest. 1.et Thomas in 2. 2ae quaest. 25. Quae sententia quam sit falsa docet Christus ipse Matt. 5,vers. 44. Precamini pro iis, ait, qui vos infestant. item docet hic Paulus qui iubet nos precari pro.  \pend
\section*{CAPVY II. }
\marginpar{[ p.45 ]}\pstart regibus, qui tunc proculdubio erant infesti hostes Christiani nominis, et persecutores.  \pend\pstart Quaeritur vero vtrum pro excommunicatis et haereticis sit in publicis precibus precandus Deus. Nam qui orat cum excommunicato, est et ipse excommunicatus canon. Qui communicauerit 1I.quaest.3. qui canon est ex consil. Chartag. 4. cap 73. Deinde dicit loan.ne salutandos quidem eos esse, aut Aue illis dicendum 2. Ioan. vers.1o.11. quare nec pro iis videtur esse precandum. Resp. Cùm aberrantes sint in viam reuocandi. vt monet Iacobus cap.5. vers. 19. nec excommunicario ad subuersionem, sed ad aedificationem sit instituta, vt conuertatur ad Deum qui peccauit.2. Cor. 7.vers. 1o.1o. vers.8.etiam pro haereticis et excommunicatis est precandum, eóque vehementiùs et ardentiùs, quo magis sunt a Deo illi alieni et exitio suo proximi. Atque haec breuis responsio est.  \pend\pstart Quaeritur etlam Vtrum pro viuis solùm sit in Ecclesia precandum, an etiam pro mortuis: haec quaestio videtur ardua et valde difficilis propter morem veteris Ecclesiae, et propter ea, quae scribit Epiph. contra Ærianos haeres.75. qui hoc fieri debere negarunt. Omitto enim quod Papistae ex2. Machabus  12. vers.43, et a5. afferunt. Primùm enim non sic legitur, vt legunt, Bonum est orare pro mortuis, et totus ille locus pertinet ad spem de resurrectione carnis, non ad preces pro mortuis. Deinde quod est vers.4o. Res est igitur salutaris et sancta, plane glossema esse apparet. Denique nec fidem meretur ille totus liber cùm sit Apocryphus, et excusationem authoris longe a Canoni\pend
\section*{AD I. PAVL. AD TIM. }
\marginpar{[ p.46 ]}\pstart cis libris dissimilem contineat cap.15v.39. Afferuntur igitur testimonia veterum, imprimis autem Augustini liber  de Cura pro mortuis cap. 18. quod canonizatur (vt vocant) cano. Non aestimemus 13 quaest.2. item aliud eiusdem Augustini dictum cap. 11o. in Enchirid. Sed etiam Arnobius vetus scriptor Ecclesiasticus idem confirmat liber  4. Aduers. gentes sub finem libri nempe pro viuis et mortuis in Dei Ecclesia precatos esse Christianos. Resp. Nullo verbi diuini fundamento niti August. vel Arnobium, sed sola quadam erga mortuos beneuolentia et humanitate, solisque coniecturis. Id quod nec Augustinus ipse dissimulat. Praeterea possunt adduci alii eiusdem August. loci, ex quibus quam varius in ea re fuerit intelligetur, quemadmodum in Enchir. pridem annotauimus. Denique opponuntur Augustino scriptores, patresque Ecclesiastici, et ii quidem pii, et Orthodoxi, quemadmodum Ambrosius in libus  qui inscribitur Abraham cap.9. vbi id tantùm mortuis a viuis praestandum docet, vt eos sepeliant: non vt pro iis orent, id quod etiam videtur confirmari posse responso Christi quod est Matt. 8. vers. 22. Obiicitur etiam Hieronymus, cuius haec est apertissima sententia. In praesenti seculo siue orationibus siue consiliis inuicem posse nos iuuari. Cùm autem ante tribunal Christi venerimus, nec Iob, nec Noe, nec Daniel, rogare posse pro quoquam. Quam sententiam Papistae ipsi in suos canones retulerunt cano. in praesenti 13 dist.2.  \pend\pstart Origo autem huius mali ex paruis initiis cepit. Primùm ex eo quod cepit piorum, qui pro nomine Christi passi erant, et mortui, in Ecclesia et  \pend
\section*{CAPVT II. }
\marginpar{[ p.47 ]}\pstart imprimis in Synaxi, publica mentio publice fieri. Vnde Martyrologia nata sunt, qui sunt libri in quibus nomina Martyrum descripta erant, quae ex scripto in communicatione Coenae recitabantur. Inde ad eorum sepulchra ceperunt homines precari, et vigilias agere, et liba offerre, vt docet Augustinus in liber  de Moribus Ecclesiast. cap.14. et Epist.119. quem tamen morem idem Epist.64 sublatum et abolitum esse cupit.  \pend\pstart Addunt Papistae quasdam exceptiones alias, quae sunt nullius momenti. Ac primum interdicunt ne omnino pro iis oretur, qui sibi manus violentas attulerunt: qui si mortui sunt, vera est eorum sententia, sin minus, falsa. Possunt enim resipiscere. Est tamen hoc Papistarum dictum in cano. Placuit vt hi 23. quaest.5. et est ex concil. Braccarens. Act. 34.  \pend\pstart Deinde iubent, ne precemur pro clericis qui in bello pro Gentilibus occubuerunt, dum eorum partes et castra sequuntur. Est ex Tiburiensi concilio in cano. Quicunque clericus 23. quaest.8. Sed haec exceptio vana est, quia non est pro mortuis orandum.  \pend\pstart Excipiuntur igitur ab hoc Pauli dicto ii soli, qui in Spiritum sanctum peccant, pro quibus minime est precandum, quemadmodum docet loan. 1. Epist5.vers.16. et exemplo Saulis confirmatur 1. Samuel 16. vers. 1.  \pend
\phantomsection
\addcontentsline{toc}{subsection}{\textit{2 Pro regibus et quibusuis in eminentia constitutis: vt tranquillam ac quietam vitam degamus cum omni pietate et honestate. }}
\subsection*{\textit{2 Pro regibus et quibusuis in eminentia constitutis: vt tranquillam ac quietam vitam degamus cum omni pietate et honestate. }}
\section*{AD I. PAVL. AD TIM. }
\marginpar{[ p.48 ]}\pstart Πρόθεσις quae non modo superius praeceptum illustrat: sedetiam explicat. Addit enim illud ipsum, de quo magis aliqua et anceps quaestio erat pro ratione temporis, quia tum omnes pene magistratus: imprimis autem summus, qui erat Romanor. Imperator, erant et infideles, et persecutores Ecclesiae. lubet tamen vt pro iis Magistratibus precemur, non tantùm Summis, quales regum nomine significantur: sed etiam lnferioribus, qui describuntur ex eo, quod supra reliquum popu lum eminent Dubitari vero de iis maxime potuit, quod Euangelicae doctrinae essent hostes, qualis Nero Imperat. dissimilis locus 1. Tessalon.2.v. 16. vbi qui praedicationem Euangelii impediunt, in manifestorum reproborum numero et albo recensentur. Deinde vetus Ecclesia preces concepit aduersus Iulianum Apostatam imperatorem. Sed solutio ex eo est, quod ii, de quibus agit Paulus in Thessal. cap. secundo non ignorantia, non sola infidelitate animi : sed obstinata prorsus malitia peccabant, et tanquam in Spiritum sanctum, vti peccauit Iulia. Apostata. Distinguendi vero sunt qui hoc modo peccant ab istis, dequibus agit Paulus hoc loco.  \pend\pstart Similis locus est, Hier. 29. vers. 7. orate pro pace Babylonis. item orate pro iis qui vos persequuntur. Denique vetus mos Ecclesiae idem probat, in qua pro Imperatoribus et praesidibus prouinciarum, quanquam eam affligerent, orabatur, quemadmodum scribit in Apologet. Tertullia. et Iustin. Martyr.  \pend\pstart Donatistae tamen contra disputant, quod Imperatoris edictoipoena in eos, si in errore persta\pend
\section*{CAPVT II. }
\marginpar{[ p.49 ]}\pstart rent, indicta esset. Nullos enim' aut paucos omnino reges pios fuisse contendunt, vel Ecclesiae fauentes. Sed eorum argumentis ineptissimis et falsissimis respondet Augustinus in liber  Contra secundam Epist. Gaudentii.  \pend\pstart Addit autem rationem Paulus, quo et superiorem exhortationem, confirmet, et nos ardentiùs inflanmet, atque  stimulos addat et excitet ad officium. Est autem ducta haec ratio ab immensa quadam vtilitate, quae triplex hic enumeratur, nimirum quod Magistratuum ope et ministerio Pax, Pietas, et Honestas inter homines stabilitur et conseruatur.  \pend\pstart Ac pax quidem siue tranquillitas tum publica tum priuata, ad quam constituendam ordinatus est a Deo Magistratus: atque etiam Dei praecepto gladium gerit, vt docet idem Paulus Rom.13. Vtraque vero pax est Dei donum, et summum omninóque necessarium humanae societatis retinendae vinculum.  \pend\pstart Pietas vero, quae Dei cultum continet etiam ad Magistratus politici, non tantùm pastoris curam et officium pertinet, quia vtriusque tabulae constitutus est custos Magistratus. Conuenit autem haec sententia cum Psalmo 101. et cum exemplis Iosiae, Ezechiae, Theodosii, Constantini Magni et aliorum piorum regum, qui cultum Dei deprauatum ex ipsius verbo reformarunt et restituerunt. Ex hoc autem loco concludi certissime potest etiam inquisitionem de haeresi, et punitionem eorum, qui merito erroris et haeresis damnati sunt ad Magistratum pertinere, quemadmodum etiam disputauit Augustinus, quanquam Castallionistae hoc nostro seculo negant, et farraginem omnium\pend
\section*{AD I. PAVL. AD TIM. }
\marginpar{[ p.50 ]}\pstart que errorum licentiam concedendam scribunt et deffendunt, ne quis, aiunt propter suam opinionem puniatur, quasi vera religio sit opinio quaedam et hominum commentum.  \pend\pstart Honestas decorum proprie est, vt non tantùm nostra cum proximo commercia regantur rectis et aequis legibus, sed etiam omnis honesta modestaque conuersatio inter nos locum habeat, etiam in rebus mediis, et quas indifferentes vocant, veluti in communi vestitu, victu, officio, et quatenus pro quoque hominum genere et vocatione distinctio quaedam officiorum est inter homines adhibenda. Haec ἀταξία ordo et politia honesta magnam aedificationem habet, et verae pietatis quodammodo custos est.  \pend\pstart Confutat autem hic locus pulcherrime Anabaptistas, qui Magistratum ex Ecclesia Dei tollunt, vt pestiferam ἀναρχίαν inducant. Quam sit autem illius vsus vtilis et Ecclesiae necessarius vel vna haec Pauli sententia perspicue demonstrat.  \pend\pstart Sed quaesitum est, Num pro iis tantùm Magistratibus sit precandum, qui suo recte defunguntur officio, et a quibus pax, pietas, honestasque constituitur aut conseruatur. Resp. Ipsum Magistratus finem, qualis a Deo praescribitur, spectandum, non autem vitia personarum quae eos gerunt. Hic enim in vniuersum harum vocationum finis est, quem proponit Paulus, propter quem nobis commendati esse debent Magistratus, id est, qui munus publicum gerunt, quanquam male officio suo fangantur. Sed pro bonis orandum est, vt eos nobis Dominus conseruet. Pro malis autem vt eos conuer tat et ad officium faciundum  \pend
\section*{CAPVT. II. }
\marginpar{[ p.51 ]}\pstart excitet spiritu suo. Itaque  semper pro iis orandum est.  \pend
\phantomsection
\addcontentsline{toc}{subsection}{\textit{3 Nam hoc bonum est, et acceptum coram seruatore nostro Deo. }}
\subsection*{\textit{3 Nam hoc bonum est, et acceptum coram seruatore nostro Deo. }}\pstart Αιτιολογία est, eaque duplex. Nam altera ducitur a Nostro officio vel rei ipsius natura καλόν ἐστι altera a Consequenti vel connexis, sunt preces tribuendae Magistratibus, quia et eos Dominus ad gratiae suae participationem et Ecclesiae communionem vocat. Itaque hoc Ecclesiae subsidio priuandi non sunt. Ex hoc autem loco colligitur ecquod sit verum nostrarum precum fundamentum nimirum Dei voluntas et promissio. Haec enim vna est optima recte legitimeque orandi regula quemadmodum etiam tradit Ioan. in 1. Epistola cap. 5.vers. 14. Denique haec eadem valet in omni cultus Dei parte.  \pend
\phantomsection
\addcontentsline{toc}{subsection}{\textit{4 Qui quosuis homines vult seruari, et ad agnitionem veritatis venire }}
\subsection*{\textit{4 Qui quosuis homines vult seruari, et ad agnitionem veritatis venire }}\pstart Altera ratio quae a connexis sumpta est. Non sunt enim priuandi et excludendi a publicis Ecclesiae precibus ii, ex quibus Deus ipse colligit Ecclesiam, et qui ad eam spe promissioneque a Deo accepta pertinent. At Magistratus, etiam qui nunc sunt a Dei cognitione alienissimi, spe tamen ad eam pertinent. Quamobrem non sunt eo fructu, dono et ea spe defraudandi. Ac propositio quidem huius syllogismi verissima est, quae non tantùm hac ratione quae naturalis est confirmatur, quod quae sunt inter se connexa, non sunt  \pend
\section*{AD I. PAVL. AD TIM. }
\marginpar{[ p.52 ]}\pstart diuellenda, sed etiam authoritate Scripturae Act. 11.vers. 17. Non est nostrum( inquit Petrus) Deum prohibere, et iis gratiae testimonia aut aditum denegare et praecludere, quibus eam Dominus ipse largitur et concedit. Voto enim Dei subseruire debemus.  \pend\pstart Assumptio vero confirmatur a Paulo, et hic disertissime est expressa, adhibeturque argumentum a genere ad speciem, sic, vult Deus omnes homines saluos fieri, et ad fidem et Ecclesiam vocari, Ergo et Magistratus.  \pend\pstart Quaesitum vero est, quae sit huius tam generalis Pauli sententiae ratio, vult Deus omnes homines, etc. Respon. Explicari hac sententia antiqua Prophetarum vaticinia, quae de vocatione Gentium loquuntur, in quibus, Dei gratia omn ibus hominibus, sublato nationis, sexus, aetatis, et ordinis discrimine, promiscue promittitur, in pri misque  illa differentia quae olim inter Gentes et Iudaeos constituta erat, hodie cefsat. Quale vaticinium est Ps.2. Pete a me et dabo tibi Gentes in haereditatem tuam Malac. 1.V.11. Nomen meum a solis ortu ad occasum magnum est etiam inter Gentes lsai.11.v.10. Et illo tempore, erit, requirent Gentes radicem Iesai, etc. Imprimis autem de vocatione magistratuum ad fidem et Ecclesiam iidem Prophetae Dei concionati sunt, veluti Isai.32.V.1. et2. Ecce in iustitia regnabit rex, et principes in iudicio praeerunt. Et erit ille vir velut latibulum a vento, receptus ab imbre, riui aquarum in terra arida: vmbra magnae rupis in terra laboriosa. etc. item cap. 6o. vers.16. Et suges lac Gentium, mamillam regum suges. Quae omnia Deiverba non dubi\pend
\section*{CAPVT II. }
\marginpar{[ p.53 ]}\pstart tat Paulus impletum iri. Itaque vere etiam Magistratus ad Dei Ecclesiam pertinere pronuntiat.  \pend\pstart Deinde a natura et definitione Euangelii idem probari potest. Est autem Euangelium potentia Dei ad salutem omni credenti siue Iudaeo, siue Gentili, vt docet idem Paulus Romanor. 1. vers. 16. Nec enim, vt Lex, sic Euangelium vni tantùm hominum generi et nationi destinatum aut proponendum erat. Ex hoc autem loco intelligimus etiam nos pro Ethnicis et Gentilibus veluti pro Turcis, ludaeis, et iis qui adhuc in orbe terrarum idololatrae manent (quales in India et insula America innumerabiles pene sunt populi) Deum orandum esse. Id quod etiam hoc loco Chrisostom. annotauit, quanquam de eo variae fuerunt Augustini tempore quaestiones, quemadmodum ex ipsius Epistolis apparet.  \pend\pstart Modum etiam per ἐξήγησιν addit, quo ad salutem homines perducuntur, nimirum veritatis agnitionem. Hic autem veritas non est accipienda cuiusuis cognitionis doctoris et disciplinae certum effatum, certaque sententia, et cum reipsa consentiens, sed Euangelium tantùm quod κατ' ὀνομα- σίαν et insigniter veritas appellatur, tum quod illa sit certissima doctrina et a Deo ipso immediate profecta: tum etiam quod sola sit veritas aestimanda, a nobis et consectanda, et persequenda: reliquae vero artes quatenus huic subseruiunt et ad vitae huius commoditatem pertinent, exercendae. Ex quo etiam colligitur non alios velle Deum saluos fieri, quam qui ad fidem Euangelii perueni unt, illique credunt. Non sunt enim ista duo distrahenda, quae hic Paulus coniungit, ne quis se  \pend
\section*{AD I. PAVL. AD TIM. }
\marginpar{[ p.54 ]}\pstart putet contempto spretóque Euangelio salutem aeternam consequi posse, ad quam sola fides est via. Id quod tamen hodie multi factitant, qua quemque in sua, quam vocant, religione saluum fieri sentiunt. Quae fuit Rethorianorum haeresis, nunc autem est Turcarum impiissima sententia.  \pend\pstart Varie vero de horum verborum. Vult Deus omnes saluos fieri sensu quaesitum est, quemadmodum ex Augustino apparet in Enchirid. cap. 1o3. Primùm enim videntur Libertinistarum, et Origenistarum errorem confirmare, qui negant vllum hominem esse a Deo damnandum, et aeterna morte puniendum. Deinde etiam fauere eorum errori, qui reprobationem Dei prorsus tollunt, quasi omnes homines sint a Deo electi, nulli autem roprobati. Denique hoc ipso loco perperam explicato se tuentur et Pelagiani, qui liberum ad vtranque bene et male agendi electionem arbitrium in nobis statuunt: et Semipelagiani quoque, qui Dei gratiam de congruo quam appellant, cum libero nostro arbitrio tanquam duo simul et aequaliter cooperantia in bene agendo coniungunt. Hi vero errores omnes, ex eo nascuntur, quod vox ipsa omnes, non recte sumitur hoc loco. Ac Scholastici quidem nonnunquam ita sentiunt ideo omnes dici a Deo saluos fieri quia dedit Deus omnibus hominibus naturam per nos ordinabilem ad felicitatem vt loquuntur. Contra vero nostra felicitas non ex naturae nostrae conditione, sed ex mera Dei gratia pendet, et eatenus ordinabilis est ad eam nostra natura, quatenus Deus ipse nos ita fingit, destinat, et ordinat: non autem quatenus reliquis homi\pend
\section*{CAPVT II. }
\marginpar{[ p.55 ]}\pstart nibus pares et similes natura sumus. Alii sic explicant, vt vocem omnes coniungendam esse doceant cum eo quod sequitur ad veritatis agnitionem peruenire, quasi non vniuersaliter neque tam late sit accipienda quam sonat: sed ex sequenti illa sententia restringenda, vt ii tantùm intelligantur comprehendi, qui credunt, vel credituri sunt Euangelio. Consir matur ex Matth.14 vers.35.36.  \pend\pstart Tertia interpretatio est eorum qui volunt omnes saluos fieri a Deo qui salui fiunt: quasi hic non definiatur, qui salui futuri sint: sed a quo salui fiant, qui saluantur. Sic nonnunquam explicat Augustinus.  \pend\pstart Verior autem sententia et iustior, meo quidem iudicio, interpretatio est haec, vt vox omnes tollat discrimen ordinum, nationum, sexuum et huiusmodi rerum, quae inter homines percipiuntur. Ex omni enim hominum genere, sexu, aetate, Deus aliquos ad se per Euangelium vocat. Itaque  non pro singulis generum accipitur, sed pro generibus singulorum vt loquuntur in Scholis, id est, non pro personis, sed pro hominum generibus.  \pend\pstart Duplex est enim huius vocis significatio. Saepe enim ita dicimus omnes, vt singulos complectamur: saepe vero, vt quosuis, non autem singulos. Sic dicitur Christus sanasse πάντα νόσον, id est, quemuis morbum. Matt. 9.v.35. Sic Paulus in 2. Thess. 1.V.3. coniungit ἕκαστον cum voce πἀντων. Et hanc esse duplicem huius vocis significationem primus obseruauit Aristoreles liber  2. πολιτικ oprimus Graecae linguae author et interpres.  \pend
\section*{AD T. PAVL. AE TIM. }
\marginpar{[ p.56 ]}
\phantomsection
\addcontentsline{toc}{subsection}{\textit{5 Vnus enim Deus, vnus etiam mediator Dei et hominum, homo Christus Iesus. }}
\subsection*{\textit{5 Vnus enim Deus, vnus etiam mediator Dei et hominum, homo Christus Iesus. }}\pstart Αἰτιολογία superioris sententiae, ab effectu. Non esset vnus Deus, et vnus mediator omnium, nifi omnes, id est, quosuis homines saluos faceret. Si enim vnius tantùm hominum generis vel ordinis salutem procurat et perficit Deus, necesse est plures deos constitui:itemque plures mediatores Quorum vtrunque blasphemum est.  \pend\pstart Similis locus est Romanor.3.vers. 29. Psal.1o5. vers.7 et 33.vers.5.  \pend\pstart Obstat autem Psal.76. vbi in sola Iudaea notus Deus esse dicitur. Responsio est, Sublatum nunc esse inter Gentes et Iudaeos discrimen, quod olim fuit, quia vtrique, diruta per Christum maceria, in vnum populum coaluerunt.  \pend\pstart Deus autem quorumuis hominum vnus neque  esse, neque dici potest, nisi suae bonitatis, clementiae, misericordiae et electionis testimonia et effecta proferat in quosuis. Quamobrem quosuis ad salutem vocare et efficaciter quidem debet.  \pend\pstart Iam vero ex hoc Pauli responso Manichaeorum, Marcionitarum et huiusmodi aliorum haereticorum error fanaticus refellitur, qui duos Deos, duóque principia constituunt, et alium Iudaeorum Deum, alium autem nostrum somniant. Vnus autem est omnium Deus, non plures.  \pend\pstart Sed etiam mediator vnus est, non plures. Is autem Mt Christus homo. Quod variis rationibus confirmari potest. Prima, Quod vnicum est semen illud Abrahae, in quo promittuntur benedicendae omnes mundi nationes Genes.15. et 17. Galat.  \pend
\section*{CAPVT II. }
\marginpar{[ p.57 ]}\pstart 3.Vers. 16. Secunda, Quod ad vnum et eundem nos reuocat tota veteris et noui testamenti scriptura. Veteris enim Legis caeremoniae et sacrificia nos ad eundem Christum deducunt, ad quem etiam Euangelium. Itaque Christus is, qui per Euangelium praedicatur, dicitur esse finis legis. Tertia, Quemadmodum Deus non habet duos filios natorales: ita neque duos mediatores constituit, sed de vno tantùm illa Dei vox est, Ille est Filius meus dilectus, in quo mihi complacui Matth.3. Consentit cum hoc dogmate, et quidem verissimo Augustinus qui liber 2. contra Epistolam Parmenia. cap. 8. negat plures esse mediatores hominum: sed vnum tantùm Iesum Christum.  \pend\pstart Similis est locus in Epistola ad Hebraeos cap. 7.vers. 26. Non enim potuit quilibet Pontifex apud deum munere mediatoris pro nobis et officio fungi: sed is solus, qui pius, innocens, segregatus a peccatoribus, et sublimior coelis factus est. Is autem est Christus solus.  \pend\pstart Obstare vero multa videntur quae afferri possunt. Ac primum.  \pend\pstart Prima ratio, Quod alii pro aliis orare iubemur. Iac.5. Responsio est, quod nostrae illae preces non propter nos gratae sunt Deo: sed vnius Christi merito et intercessione. Eatenus enim De o placent, quatenus et fidei et mutuae inter nos charitatis sunt effecta atque fructus, quae vtraque  innititur Christo.  \pend\pstart 2 Moses et Sacerdotes Leuitici dicuntur fuisse intercessores et mediatores inter Deum et populum Exod.32.Resp. Puerilem esse obiectionem.  \pend
\section*{AD I. PAVL. AD TIM. }
\marginpar{[ p.58 ]}\pstart Illic enim intercedere nihil aliud significat quam internuntium et medium esse: non autem ipsam Dei gratiam nobis promereri.  \pend\pstart Tertia, Obiicitur, locus Galat.3.v,2o.Respo. Est sophistica obiectio. Negatur quidem internuntius et mediator esse vnius tantùm, non autem negatur esse vnus.  \pend\pstart Quarta, Angeli et Sancti vita functi, sunt nostri coram Deo intercessores. Nam sine mediatore non possumus accedere ad Deum, vti nec regem sine purpurato, qui nos intromittat, adimus. Resp. Praeter verbum Dei et Angeli et sancti nostri mediatores inducuntur. Nam de Angelis nominatim vetat Paulus ne colantur. Coloss.2. De sanctis vita iam defunctis ratio prohibet, quod et ipsi mediatore egent: neque fuerunt sancti et impolluti. noster autem aduocatus apud Deum debeat esse δίκαιος vt docet 1. Ioan.2. vers.2.  \pend\pstart Quod enim affertur de mediatore alio intercessionis et salutis, conuellitur vel ex hoc ipso Pauli loco, vbi agitur de precibus et intercessione. Item Hebr.7.vers.15, et 9.vers. 24.  \pend\pstart Addit Paulus Christum hominem esse. Primùm quod ea ratione Christus mediator noster est, quatenus homo pro nobis factus est. Deus enim manens pati et implere illa, quae sunt nobis ad salutem necessaria, non potuit. Deinde ne immensus ille diuinae maiestatis fulgor et splendor, ad quem nobis, qui sumus terrae vermes et homines, accedendum est nos perterreat. Habemus enim qui nobis facilem ad Deum aditum praebeat, Christum nempe, qui nostris infirmitatibus compatitur, et qui, quia nostram naturam assum\pend
\section*{CAPVT II. }
\marginpar{[ p.58 ]}\pstart psit, nobis iam formidabilis esse non potest. Heb 4.vers.15. Bernard. sermo.73. Canti. in tanta trepidatione electis fiduciam praestat naturae similitudo. Denique vt facilius quia naturam omnibus hominibus communem Christus assumpsit, intelligamus quorumuis hominum eum mediatorem esse, non vnius tantùm hominum generis, humanitatis vel hominis Christi mentionem fecit Paulus.  \pend\pstart Quaesitum est autem, Vtrum Christus, qua tantùm homo est, noster mediator sit, quod et Stancarus, Menno quidam, et hodie noui Arriani sentiunt, an etiam qua Deus est. Resp. Cùm mediatoris oeconomia et officium ad totam Christi petsonam pertineat, neque aliter saluator noster esse potuerit Christus, nisi esset Emmanuel id est, nobiscum Deus, idcirco et qua Deus, et qua homo est Christum mediatorem nostrum esse fateamur necesse est.  \pend\pstart Hoc autem docet et confirmat August. cùm in Enchirid. tum vero pulcherrime in libro de Ouibus  et liber  Confessionum. De qua re cum copiosissime doctiss.nostri temporis Theologi D. Caluinus in Epistola 312. et Theodo. Beza Epist. 28. disseruerint, plane superuacaneum puto latiùs hoc argumentum persequi.  \pend
\phantomsection
\addcontentsline{toc}{subsection}{\textit{6 Qui semetipsum dedit redemptionis pretium pro quibusuis, Christus inquam, teStimonium illud suis temporibus destin atum. }}
\subsection*{\textit{6 Qui semetipsum dedit redemptionis pretium pro quibusuis, Christus inquam, teStimonium illud suis temporibus destin atum. }}\pstart Αὔξησις est, per quam non modo superiorem rationem confirmat a connexis: sed etiam me\pend
\section*{AD I. PAVL. AD TIM. }
\marginpar{[ p.oo ]}\pstart diationis et intercessionis Christi pro nobis fundamentum explicat. Est autem ipsius sacrificium pro nobis, quod hic πεεεγρασν quadam describitur. Sunt enim res inter se connexae, sacrificium et intercessio apud Deum: atque etiam ita inter se comparatae, vt vna alterius causa sit. Id quod his Epistolae ad Hebraeos locis confirmatur cap. 4.verf.14. 5vers.1. 7. vers. 15. 8.vers.3. 9. vers.24. atque etiam hac ratione. Quod cùm hic sit intercessionis finis, vt Deum nobis propitium et beneuolentem reddat, is demum apud Deum pro nobis intercedere potest, qui eum nobis placare potest. Placatur autem Deus, non alia ratione, quam plena poenae peccatis nostris debitae persolutione. Id autem solum Christi sacrificium et mors potuit. Ergo illa mors Christi intercessionis pro nobis est fulcimentum et fundamentum. Vnde perperam Scholastici, qui salutis et intercessionis mediat ores diuersos faciunt. Quod esse non posse satis ex superiore argumento apparet.  \pend\pstart Definitur autem et ornatur hoc Christi sacrificium magna laude et encomio, quod ductum est ab ipsius effecto. Id autem est Redemptio nostra, cuius illud sacrificium fuit pretium integrum, plena merces, et iusta satisfactio siue persolutio. Dicitur autem Christus non tantùm Seipsum dedisse: sed etiam Ipse se dedisse. Nam neque  aliud pro peccatis noitris, quam seipsum dedit: neque cùm dedit, inuitus aut ignorans id fecit, sed volens. In quo ipso com mendatur impense Christi erga nos bencficium et charitas. Similis locus. Tit 2.V.14.  \pend\pstart Nec caret ἐμρέσι, quod sanguis Christi illo sacrificio effusus appellatur ἀν τιλvreςν, quia sanguis  \pend
\section*{CAPVT II. }
\marginpar{[ p.61 ]}\pstart ille Christi agni immolati longe est maioris pretii, et effectus, quam vitulorum et hircorum omnium sanguis. Est enim ille agni immolati sanguis, omni auro lapideque pretioso aestimabilior et potior, vt est Hebr. 9.vers. 12.1.Pet.1.vers.19. et vere plena et aequalis nostro peccato et debito satisfactio et pretium 1.Corinth.1.vers.3o. Vnde non tantùm λύτρον: sed ἀντίλυτρον dicitur, quod Latini dicunt contra auro venire, id est, iustum, et aequale rei ipsi pretium esse. Ex quo diluitur illa curiosorum hominum disputatio, Vtrum fuerit poena, quam pro nobis Christus pertulit, ae qualis peccatis nostris, vt peccata ex merito Christi deleantur. Fuit enim, vt hic docet Paulus, aequiualens, id est, aequata ipsi peccatorum nostrorum foetori apud Deum, et plane satisfactoria poena Christi perpessio et obedientia.  \pend\pstart Secundo vero loco notandum est, Nullam aliam satisfaciendi diuinae iustitiae pro peccatis nostris rationem esse a Deo constitutam, illique gratam et acceptam praeter sanguinem vnius Christi. ltaque nec opera quae vocant bona: nec caeremoniae, nec vllius alicuius rei, quantunuis pretiosae, donatio, aut ratio potest pars esse satisfactionis nostrae coram Deo, que in solidum in vna Christi morte quaerenda nobis est. Sed nec ex parte tantùm Christus satisfecit, veluti vt deleat ea tantùm peccata, quae dicuntur venialia: aut quae ante Baptismum commisimus, nobisque nocebant: sed etiam expungit, et redemit:  \pend\pstart Addit Paulus τὸ μαρτύριον καιροῖς ἰδίοις. Haec abrupta videtur esse oratio, itaque obscurum habet sensum, et diuersas interpretationes. Alii e\pend
\section*{AD I. PAVL. AD TIM. }
\marginpar{[ p.62 ]}\pstart nim ad sequentem versiculum referunt: alii huic coniungunt: alii mutandam censent vocem μαρτύριον in vocem μυςτήριον. Mihi vero videtur planus et facilis horum verborum sensus, si ab initio quinti versiculi ad hunc vsque locum parenthesin produci intelligamus. Itaque continua orationis serie haec, quae iam sequuntur, coniungantur, cum illa superiore Pauli sententia. Qui quosuis homines vult saluos fieri, et ad agnitionem veritatis peruenire. (Vnus enim Deus et c. iuxta testimonium propriis temporibus patefactum, ad quod ipsum implendum et exequendum, ego constitutus sum, etc. At supplendam esse post haec verba τὸ μαρτύριον καιροῖς ἰδίοις vocem φανερωθὲν apparet ex Tit. 1. vers.2. et Coloss.1.vers.26. Vocat autem hic testimonium Paulus, antiqua Prophetarum vaticinia de vocatione quorumuis hominum et gentium, ne id frustra et temere Ecclesiae polliceri ipse videatur Paulus nullóque in eo niti sacrae scripturae testimonio et fundamento, qualia tamen multa supra cap. q.annotauimus et ipse Paulus obseruauit et affert in Epist. ad Romanos cap.15. vers 9.1o.11.12. Ergo vocem μαρτύριον non refero ad Christum: sed ad totum complexum superius Vult Deus quosuis homines saluos fieri. Hoc testimonium et illa vaticinia antiqua ignota aut etiam obscura iis ipsis, quibus antea annuntiabantur, suis, id est, Euangelii temporibus patefacta sunt, et illustrata atque impleta. Nouit enim solus Dominus temporum articulos, et rerum opportunas maturitates, quas ipse decernit et constituit. Itaque huius rei perficiende gratia Paulus caeterique  Apostoli a Christo et vocati et missi sunt. Id quod ipse se\pend
\section*{CAPVT II. }\pstart quenti versiculo subiicit, et explicat.  \pend
\phantomsection
\addcontentsline{toc}{subsection}{\textit{7 Cuius constitutus sum ego praeco et Apostolus (veritatem dico per Christum, non mentior) doctor, inquam, Gentium cum fide ac veritate. }}
\subsection*{\textit{7 Cuius constitutus sum ego praeco et Apostolus (veritatem dico per Christum, non mentior) doctor, inquam, Gentium cum fide ac veritate. }}\pstart Α’πτιολoyια est et consirmatio proxime superioris sententiae ab effectu vel a consequente, Deum n.id decreuisse dubitari non potest, cùm executus sit, et propterea Paulum miserit, qui per Euangelii praedicationem quosuis vocaret ad agnitionem veritatis. Itaque suum apostolatum, qui a Deo erat, pro iusta huius testimonii et voluntatis Dein de vocandis Gentibus probatione affert Paulus. Vocat autem se Κήρυκα, id est, Praeconem et Aʹπόστολον, quae duo ita inter se videntur differre, quod illud est generalius: hoc specialius. Plures enim sunt verbi Dei κήρυκες et praecones, quam Apostoli: tum deinde, quod κήρυξ dicitur Paulus ratione executionis ipsius ministerii, et praedicationis verbi Dei: Α’πόςτολος autem ratione vocationis et gradus, in quem a Deo collocatus et assumptus erat. Hoc autem ipsum postea definit cùm ait διδάσκαλος ἐθνῶν. Quae et superiorem nostram sententiam et interpretationem confir mant, et finem ministerii Apostolici, imprimis autem Paulini, ostendunt.  \pend\pstart Obstat quod Romanor.1. se Iudaeorum quoque debitorem appellat. Resp. Gentium praesertim causa designatus erat Apostolus Actor.13, et Galat.1.vers.16.2.vers.8. quanquam Iudaeis Euangelium quoque pro re nata annuntiauit.  \pend
\section*{AD I. PAVL. AD TIM }
\marginpar{[ p.64 ]}\pstart Iureiurando etiam suam vocationem confirmat, quod vocatio Gentium tunc res noua prorsus et aliena a Dei consilio censeretur, qui tanto tempore solos Iudaeos pro suo populo agnouerat et res ista esset magni momenti. Quamobrem non videtur temere et de re nihili sumptum a Paulo Dei nomen. Similis autem iuramenti formula est etiam apud eundem Paulum Romano. 9.vers.1.  \pend\pstart Εν χρίστῳ dupliciter vel in Christo, id est, Christo linguam meam et mentem dirigente, et commouente: vel per Christum, id est, teste Christo ipso, quem huius rei et meae vocationis et muneris testem produco. In qua postrema sententia conueniunt docti interpretes. Quaeritur autem Num vero tunc per creaturas iuretur, quum per Christum iuratur. Sed responsio facilis est. Non iurari, quia Christus non tantùm est homo et creatura: sed etiam Deus. Deinde hic videtur Christus potius testis, quam iudex produci. Testes vero possunt a nobis, vt a Prophetis, appellari etiam mutae creaturae, nedum Christus.  \pend\pstart Addit denique in Fide, et Veritate. Quibus verbis et eam doctrinam confirmat, quam Gentes docebat cùm fidem appellat: et animi sui synceritatem, cùm huic fidei veritatem coniungit. His autem verbis breuiter comprehensa est veri pastoris definitio. Est enim is, qui in Ecclesia Dei legitime vocatus docet fidem, id est, sanam doctrinam in veritate.i.sana coscientia et recto fine. Quam definitionem confirmat Petrus 1. Pet.5.vers.2.3.4.  \pend
\section*{CAPVT II. }
\marginpar{[ p.65 ]}
\phantomsection
\addcontentsline{toc}{subsection}{\textit{8 Velim igitur viros precari in quouis loco, puras manus attollentes absque ira et disceptatione. }}
\subsection*{\textit{8 Velim igitur viros precari in quouis loco, puras manus attollentes absque ira et disceptatione. }}\pstart Hic versiculus varia capita complectitur, qua ad orationem pertinent, quaeque breuiter quidem sed commode Paulus hic tradit. Tria enim comprehensa sunt. 1 Quo animi affectu seu praeparatione orandum sit. 2 In quo loco. 3 Quo gestu.  \pend\pstart Sunt autem haec prae cepta non Pauli, fed Dei. Itaque vox βούλομαι non priuatum quoddam humanae mentis commentum significat, vti nec 1. Corinth. 1o. vers.1. sed praxin horum praeceptorum a se serio requiri docet his verbis Paulus, sine qua non potest probari, et Deo grata esse nostra oratio.  \pend\pstart Ac quod ad praeparationem (nec enim illotis manibus et impraemeditati ad Deum orandum debemus accedere) duo imprimis requirit a nobis, Sanctitatem siue puritarem vitae, et Fidem.  \pend\pstart Sanctitatem autem illam vitae designat et describit ab enumeratione partium, nempe a Manuum sanctitarte et puritate, et a cordis charitate. Vult enim et manus nostras (id est, opera externa) esse sanctas, vt est Isai.1.vers.13.et corda ab omni in proximum ira et irritatione vacua, vt est Matth.6. vers.15.5. vers.24. Vera enim sanctitas tum externis, tum etiam internis operibus constat, et definienda est: et impium est, si quis vel animo vel opere ipso sceleratus et profanus ad Deum orandum accedat, sine animi poenitentia et reIipiscentia Psal.5.vers.6. Id quod etiam homines  \pend
\section*{AD I. PAVL. AD TIM. }
\marginpar{[ p.66 ]}\pstart nrofani senserunt vti Hesiod. liber 1. ἔργ. καὶ ἠμερ. et Plato, et ex Platone M. Tull.libus 2. de Legibus.  \pend\pstart Videtur etiam haec sanctitas, quae vera est, hic esse opposita a Paulo omnibus illis externis ritibus, lotionibus, et carnalibus purgationibus, quae veteri Dei lege hominibus templum ingressuris praecipiebantur, de quibus agit Apostol. ad Hebus  9. vers 1o. et Moses Leuitici cap. 6. 13.15. 16.19. quibus ingens etiam cumulus a Pharisaeis additus erat postea per δευτερώσεις, vt apparet Mar.7.V.4. Quae lotìones antiquae a Deo praȩceptae verae quidem animi vitaeque sanctitatis signa erant et figurae: earum tamen implementum et corpus in Christo habetur et confertur. Quanquam vero de publicis precibus hic agit Paulus: in priuatis tamen eundem animum sanctum et purum requiri certissimum est, et ostendit Christus. Matt. 5.vers.24. Est in canon. Nihil. et canon. Non mediocriter, quod est dictum Hieronym. De consecrat. distinct.5  \pend\pstart Certa vero animi fiducia, quae in gratuitis Dei promissionibus acquiescat, orandum esse, docet et Christus ipse Ioan.4.v.23, et Iac.1.v.6. et Paulus Romanor. 1o.vers.14. quia vera precatio et Dei inuocatio est fidei effectus: quae cum animi haesitatione prorsus pugnat.  \pend\pstart Quae postea de certo quodam ieiunio, de lotione manuum addita sunt a Patribus, de quibus hic Chrysost. Hom. 7 vbi quosdam sollicitos fuisse notat de huiusmodi rebus: non modo legales, sed etiam Ethnicorum caeremonias, aut potius superstitiones reponunt nobis, et redolent, etsi animus sobrius longe ardentiùs Deum precatur,  \pend
\section*{CAPVT II. }
\marginpar{[ p.67 ]}\pstart quam satur et cibi plenus.  \pend\pstart Videndum iam est de loco, vbi orandum sit. Resp. vero vbique posse. Conuenit enim Paulo cum Christo Ioan.4. vers. 21. Quemadmodum enim Deus non est acceptor personarum, ita nec locorum. Nam Domini est terra, et plenitudo eius. Psal. 24. Videturque absurdum, vt Deus, cuius latissime per vniuersum orbem diffunditur maiestas et potentia, vbique agnosci inuocari et coli non possit, vt est Malach. 1.vers.11.  \pend\pstart Obstat tamen quod est scriptum Deuter.12.v. 5, et 2. Chronic.7.vers.12. templum Dei a Solomone constructum eum fuisse locum, quem Dominus specialiter elegerat, vt ibi inuocaretur. Itaque versus illud templum conuersus etiam in media Babylone precabatur Daniel, vt ipse scribit cap.6.vers.1o. Respond. vero pro tempore ita constitutum a Deo fuisse, vt quanquam vbique inuocari posset nomen ipsius, sacrificari tamen solùm in templo Hierosolymitano tunc temporis fas esset, quia et caeremoniae tunc locum habebant, et gratia Dei nondum diffusa erat in omnes Gentes. Denique hac paedagogia Dominus tunc volebat suos ad vnionem fidei et doctrinae exhortari, prouocare, et adducere, et futuram libertatem sub Christo maioremn ostendere, qualis etiam est.  \pend\pstart Ergo sublatum est illud locorumh vetus discrimen, vti nec iam cultus Dei in sacrificiis hostiarum consistit. Quo fit vt vbique pie colatur et adoretur Deus, et vt illa sacrificia, quae hodie exigit, quae sunt gratiarum actiones, vbique terrarum illi offerri possint. Atque eo praecipue per\pend
\section*{AD I. PAVL. AD TIM. }
\marginpar{[ p.68 ]}\pstart tinet hic locus Pauli, vt Iudaicum illud locorum discrimen sublatum esse per Euangelii praedicationem intelligamus, et quae quantaque sit iam nostra per Christum libertas, per quam et locorum et rituum seruitus nobis per Christum adempta est, cognoscamus.  \pend\pstart Neque tamen confusionem inducit vel inuehit in Ecclesiam Paulus, in qua iubet ipse vt omnia τάκτως et ordine fiant, quasi iam nolit vllum esse communi totius Ecclesiae consilio et delectu constitutum locum, in quem Christiani ad Deum precandum certo tempore conueniant, seque adunent: sed vt sparsi, et, prout quemque feret animi impetus, orent separati et disiecti. Hoc enim plane furiosum esset, et μανιῶδες: sed religionem propter loca vllam animis nostris inhaerere vetat quasi locus ipse sanctiorem gratioremque Deo nostram orationem efficiat. Quod superstitiosi homines, id est Papistae etiam non hodie putant.  \pend\pstart Vnde ex hoc loco votiuae illae ad terram sanctam, et alia quaedam Martyrum loca peregrinationes et precationes merito damnantur, quae ea de causa suscipiuntur a superstitiosis, quod nescio quid maioris sanctitatis in illis locis, et ipsi quam Dominus pedibus suis calcauit, terrae inesse iudicant. Quod falsissimum est. Si qua enim hodie est terra maledicta, est Iudaea: et huiusmodi cogitatio plena est idololatriae et blasphemiae. Nec iuuantur huiusmodi superstitiosi homines exemplo Naamani Syri, qui, vt est 2.Reg.5.terrae ipsius sanctae glebas quasdam et onera secum apportauit. Illo enim tempore caeremoniae locum adhuc habebant, et extra terram a se delectam sibi sa\pend
\section*{CAPVT II. }
\marginpar{[ p.69 ]}\pstart crificari Dominus nolebat. Quanquam aliquid in eo Naamani infir mitati concessum esse manifesto apparet, quod et sacrificat, et extra templum Dei. Mos autem iste visitationum monumentorum Martyrum, et terrae sanctae, votiuarumque peregrinationum ex superstitiosa et nimia Martyrum veneratione primùm ortus latius postea serpsit, et ex errore errorem produxit. Quam parum honorifice de Hierusalem loquatur Chri stus ipse apparet Matth.23.vers.37. dum eam homicidam Prophetarum appellat: et post eum Paulus, dum eam seruae et ancillae confert. Galat. 4.et sequens aetas, quae Apostolorum temporibus vicina fuit, idem sentiebat, quae ne monumenta quidem martyrum saepe norat. Sepultura enim eorum contenti, vt de Stephano docemur.Act. 7, et 8. venerationem istam, quae cum idololatria coniuncta est, omittebant: imo vero detestabantur et damnabant. Primùm Helena Constantini Magni mater mulier superstitioso et foemineo quodam zelo commota, terram, quam vocant, sanctam inuisit, non ipsius quidem terrae gratia, sed Christi causa et fidei suae vt verisimile est, magis confirmandae, quia fidem et soam et aliorum ipso aspectu rerum et monumentis passionis Christi (quae a Iudaeis Ethnicis, et quibusdam haereticis negabatur) et quae se dulo illa conquisiuit, confirmare voluit. Quanquam iam in eo peccatum est tamen, quod plus fidei de Domino cruci, clauuis, et sepulchro, quam ipsi Spiritui sancto, Apostoli scriptis et Euangelistis tribuisse videatur. Sed post orbem Christianum factum supplicarunt superstitiosi ad Martyrum monumenta, ad  \pend
\section*{AD I. PAVL. AD TIM. }
\marginpar{[ p.70 ]}\pstart quae concurrebant, tanquam augustiora quaedam loca, quod illic Martyrum, quorum fides Deo accepta fuerat, ossa iacerent. Sic Monica mater Augustini ex Affrica Mediolanum quotannis pergebat ad Geruasii et Protasii Martyrum sepulchra, quae tempore Ambrosii Episcopi eruta fuerant. Nec tamen ossibus et cadaueri ipsorum Martyrum hic honos tribuebatur primùm, sed ipsi defunctorum tantùm fidei: peccatum est in eo tamen, quia extra verbum Dei id fiebat. Quod ipsum etiam iam quibusdam in Ecclesia bonis viris displicuit, et Ærius extitit Constantini Magni seculo, qui omnia illa et merito quidem damnauit, quanquam ipse, quod iam communi errore haec recepta erant, tanquam haereticus propterea habitus est. Facit enim communis et receptus iam inter omnes error ius et confirmationem vt scribunt Iurisconsulti in L Barbarius D. de off. Praetor. Item Vigilantius tempore Hieronymi extitit, qui hoc totum superstitiosi cultus genus in reliquiis sanctorum colendis aperte idololatricum esse probauit. Error tamen obtinuit tum negligentia pastorum, tum quia homines suis commentis potius, quam ex verbo coelesti Deum adorare cupiunt. Inde furiosa aedificatio templorum consequuta, quae in diuorum honorem facta sunt, post annum praesertim 6oo. a Christo passo sub lustiniano imperatore Constantinopol. homine iis superstitionibus insane addicto, vt apparet ex Procopii liber  de aedificiis Iustinia. Et post eorum aedificationem, eamque etiam sumptuosam et magnificam religio et veneratio iis addita est, post deinde ipse cultus Dei iis conclusus, aut alligatus, vt nolla  \pend
\section*{CAPVT I. }
\marginpar{[ p.71 ]}\pstart sacra legitima et Deo grata oratio extra ea fieri posse ab hominibus etiam Christianis censeretur.  \pend\pstart Nec illa superstitio deffendi potest exemplo eorum, qui in sepulchrum Elizaei cadauer quoddam festinantes et coacti proiecerant. vt est 2. Reg.13.vers.21. Nam illi cadauer Elizaei non coluerunt: deinde id coacti fecerunt. Demum nulla inde superstitio nata dicitur ad sepulchrum Elizaei: sed tantùm proptere fuit doctrina Prophetae, quae in animis hominum adhuc recens inhaerebat, confirmatior facta, et Deus ipse maiori honore ab illius seculi hominibus cultus.  \pend\pstart Templa igitur in Ecclesiis Dei esse vtile quidem est: sed tamen nec sumptuosa, nec superstitiosa, sed quae ad capiendum populum sint satis commode extructa, et quibus fanctitas nulla religióque ascribatur. Hoc enim est plane superstitiosum et Iudaicum. Primùm enim sine templis propriis fuit Ecclesia Christiana tempore Apostolorum, et post eos etiam temporibus Iustini Martyris: falsa sant enim quae in Higini et Sexti Roman. Episcoporum Epistolis de consecratione templorum in decretis et alibi extant. Sub Diocletiano primum apparet Christianos habuisse oratoria quaedam, et προς ευκτηρίους  ὄικους, vt vocat Eusebus liber 9. Histor.cap.1o. et Ruffi in Histor. Ecclesiast. Templa igitur Christianorum imperatorum aetate et imperio demum aedificari ceperunt, et oratoria dicebantur, quae si ampliora erant, Basilicae: nondum autem ναοί aut ἱερα. Hae Basilicae etiam saepe a Caesarum nominibus, a quibus fuerant extructae, vocabantur, et denominabantur, vti Basuica Constantini Euag.libus 2.  \pend
\section*{AD I. PAVL. AD TIM. }
\marginpar{[ p.72 ]}\pstart cap.8, et3 cap.2. postea Apostolorum et Martyrum vocabulis nuncupari ceperant, in quorum erant honorem et memoriam constructa Sozom. libus  6.cap 8.quanquam ea Martyribus, non vt diis, fed vt hominibus, quorum memoriam colebant, consecrabant, vti ait Augustinus liber 22. de Ciuit. Dei cap.1o. et saepe Martyrium ipsum fuit totius Basileae tantùm pars quaedam cap.8. Demum eadem templa consecrari, dedicari, quibusdamque ritibus sancta credi et fieri ceperunt. Quod ante Constantinum Magnum factum et vsurpatum fuisse non videtur. Sed cùm primus ipse in loco Caluariae templum Martyrum exaedificasset, postea ad maiorem, vt illi volebant, venerationem, sed potius ad adulationem et ad morem Paganorum, profanatum illud est potius: quam dedicatum, et consecratum a quibusdam insulsis Episcopis, Sozomenus liber 2.cap.26. et Eusebus liber 4. de vita Constantini. Quo scelere etiam pollutum est templum illud augustum, quod Saluatori nostro Hierosolymis idem Constantinus extruxit. Hunc tamen morem et exemplum secuti sunt postea alii Episcopi auide, non spectaro, si ex Dei verbo id fieret, neque quid ex eo mali consequeretur: sed placuit aliis nouitas et ille ritus et inauguratio in Constantini templo iam vsurpota et probara a quibusdam, vt ab Eusebio Caesariensi Rethoricoteros laudata. Itaque Basilius ipse Episcopus vir alioqui doctus, ad Episcopi et ad templi Basilicae consecrationem solenniorem et venerabiliorem alios secum Episcopos conuocauit, vt tradit Sozomenus liber 4.cap.13. Quod idem factitatum fuisse in occidentalibus Ecclesiis, sed postea,  \pend
\section*{CAPVT I. }
\marginpar{[ p.73 ]}\pstart apparet ex Epistolis Ambrosii. Demum certae caeremoniae sunt institutae, adhibitae, et verba, et precationes, quibus sanctitas ipsis lapidibus inhaereret. Ita repetita et reposita sunt in Dei Ecclesia, quae profani homines in dedicandis idoliis suis obseruare consueuerant, de quibus agitur de Consecr. distinct.1. Caeterum illa iam dedicata templa appellarunt augusto nomine ναοὺς, ἐκκλησιας ἰερὰ vt apparet ex Epistolis Sidonii Apollinaris. Quorum etiam templorum inter Christianos, vt inter Iudaeos, tres partes constituerunt. Nempe Sanctum sanctosum, vbi est magnum altare Aʹ'γιον: et Sanctum siue χόρον, vbi est chorus Sacerdotum canentium: Ναὸν, vbi est plebs.  \pend\pstart Sed quaesitum est, versus quam mundi partem et plagam sit orandum. Respond. Siquidem verbum Dei, a quo solo pendere debemus, spectemus perinde est, neque refert in quam mundi regionem conuersi precemur. Sin autem morem Papistarum, versus Orientem orant, et eôdem etiam templorum suorum capita conuertunt. Vetus quidem Ecclesia etiam sub Constantino liberior fuit, vt docet Socrates liber  5. cap.22. Rationem Papistae afferunt, quod ad Orientem fuerit situs paradisus. Sed ex Ezechiel 8.vers.16.respondemus damnatos esse a Deo, qui versus Orientem quadam religione ducti orarent, nec quia ibi primum paradisum collocauerat Deus, voluit templum soum eo vergere. Vrbs enim Hierusalem versus meridiem. Templum autem, quod in monte Sion erat, ad Septentrionem situm fuit, vti apparet ex Psal.48. et Ezechiel.4o. vers.2. Depique videntur Papistae veteres Persas idolola\pend
\section*{AD I. PAVL. AD TIM. }
\marginpar{[ p.74 ]}\pstart tras imitari, qui praecise ad Orientem orandum esse docent, tantum abest, vt sit haec traditio Apostolica. Quod item idem Basil.sermo.2. de Batis. cap.8. in loco non sacrato mysteria non putat posse celebrari, fallitur et non videt vbi verbum Dei praedicatur, ibi locum esse sacrum 1. Timoth.4.vers.4.  \pend\pstart Tertio loco ex ipso versu quaeritur. Quo gestu sit orandum. De quo quia nihil est hoc loco diserte praescriptum, sed nec in toto Dei verbo, intelligitur habere ea res liberas obseruationes, modo ne quod offendiculum aliis nimia affectatione vel dissensione praebeamus. Christus ipse et stans et flexis genibus orauit, publicanus stans in templo precatus est, Luc.18.vers. 13.22 vers.41. In Actis videmus positis genibus Christianos veteres saepe orasse Actor.7.9. et 2o.Stantes orant, qui spe promissionum Dei erigantur. Flexis autem genibus id faciunt qui, propter peccatorum suorum sensum deiecti coram Deo humiliantur et prosternuntur. Quidam etiam prostrati orarunt, vt Elias 1.Reg.18.Nec in eo synodorum decreta magni esse momenti ad stabiliendam pietatem existimemus, sed morem potius regionis, in qua sumus, sequamur, modo ne sit in eo more et ritu aliqua idololatria apparens. De manibus etiam quaesitum est, quo earum situ, et gestu sit orandum. Resp. Alii iunctis, alii supinis, alii sublatis orarunt, vt docet Clemens Alexandr.libus 7. Stroma. Itaque totum hoc genus obseruationum et caeremoniarum liberum est. Et quod ait hoc loco Paulus, Attollentes manus, metonymico dictum est. Nam signum pro re signata positum  \pend
\section*{CAPVT II. }
\marginpar{[ p.75 ]}\pstart est, vti Isai.1.vers.15. nimirum pro ipso cordis affectu, qui ad Deum erigi debet. Hic enim gestus fauorem animi designat, quo preces ad Deum nostras concipi debere significat Apostolus: non quod homines Christianos huic caeremoniae in precandon velit esse alligatos et astrictos.  \pend\pstart Denique quaeri etiam potest de tempore, quo preces concipi et fieri a Christianis debent. Res. Quod quidem ad priuatas, non tantùm quotidie semel: sed etiam et bis et saepius in die fieri oportere, nimirum mane et vesperi, cùm surgimus aut cubitum discedimus, vt omnes actiones nostrae a Deo incipiant et finiant. Id quod exemplo Dauidis ita faciendum esse monemur Psal.55 vers. 18 et suo seculo in omnibus familiis factitatum scribit diserte hoc loco Chrysostomus. Quod autem ad precationes publicas nulla fuit ante Christianos imperatores de eo lex in Ecclesia constituta, cùm non auderent Christiani homines propter persecutiones libere conuenire. Post fancitam vero Ecclesiae libertatem toties orarunt, quoties ad audiendum Dei verbum, vel ad Sacramentorum participationem colligebatur Ecclesia, quia sine precibus nunquam dimittebatur coetus piorum, vti docent, et lustinus Martyr et Tertullian.in Apologet. et Socrates liber 5.cap. I6. Certis tamen horis non erant precationes primum institutae, alligatae, et vinctae: sed creuit postea superstitio, et deuincta est certis horis orandi consuetudo et necessitas, etiam sine publico Ecclesiae coetu, et sine predicatione verbi Dei. Id quod docet Hieronymus ad Eustochium, adeo quidem, vt etiam sioe conuocatione Ecclesiae o\pend
\section*{AD I. PAVL. AD TIM. }
\marginpar{[ p.76 ]}\pstart mnibus pene diei horis preces publicae fieri decernerentur, diluculo, deinde tertia diei, id est, a luce solis exorta supra horizontem hora, post sexta, praeterea nona. Denique sub vespertinum diei tempus, Vnde illae horae inter Papistas dictae sunt Canonicae nempe 1.3.6.9. de quibus etiam meminit Clemens Stromat.libus 7. Sedulius, Cassianus, Beda et alii. Atque haec tota res magnam secum fuperstitionem primùm traxit, et verum finem orationis, atque vim extinxit: atque in opinionem meriti, et satisfactionis conuertit. Denique Euchitarum errorem in Ecclesiam induxit, qui nihil, nisi precari, id est, demurmurare certa verba solebant, et in eo proram et puppim salutis collocabant, vti diximus. Haec nos omnia breuiter hic complecti voluimus, vt tanquam locus communis haberetur. Sequitur iam vt de decoro, quod non tantùm in precibus, sed in congressionibus publicis et Ecclesiasticis seruandum sit videamus, quod sequentibus versiculis explicatur.  \pend
\phantomsection
\addcontentsline{toc}{subsection}{\textit{9 Itidem et mulieres amictu honesto, cum verecundia et modestia ornare sese, non, cincinnis, vel auro, vel margaritis, vel pretioso vestitu. }}
\subsection*{\textit{9 Itidem et mulieres amictu honesto, cum verecundia et modestia ornare sese, non, cincinnis, vel auro, vel margaritis, vel pretioso vestitu. }}\pstart Μεριςμὸς est, supra enim de viris egit, nunc autem de foeminis Christianis, quales eas in Ecclesia Dei esse et conuenire oporteat, vt ad rite orandum sint comparatae. Quo fit, vt hic de Christianarum mulierum officio et decoro duplici ratione et respectu Paulus agat. Nimirum quatenus sunt spectandae tum In sese, tum in totius Ecclesiae coetu.  \pend
\section*{CAPVT II. }
\marginpar{[ p.77 ]}\pstart Ac in sese debent esse Castae, et Modestae. Coetus autem Ecclesiastici ratione, Silentes, et Discentes.  \pend\pstart Primùm igitur ait. Itidem, Nam ne se a consortio et com municatione precum Ecclesiae excludi propter sexus infirmitatem, et maiora virorum priuilegia putent mulieres Christianae, docet eundem illis aditum ad Deum patere, et eandem orandi fiduciam dari, quam supra viris esse concessam ostendit. Deus enim suas promissiones vtrique sexui communes proposuit: et tam dicitur Sara parens et mater fidelium mulierum quam Abraham fidelium virorum 1.Pet.3.  \pend\pstart Communis vtrisque Baptismus, Coena Domini, Praedicatio Euangelii, fides, et gloria aeterna, quia vtrisque communis est imago Dei. Ergo mulieres,vti et viros, orare, et Deum precari vult: sed tacitas, non autem in coetu eas preces fundere et concipere. Vult autem Paulus eas esse in coetu Domini amictas et vestitas Honeste, Verecunde et Modeste. Atque huic muliebii cultui et ornatui opponit, Cincinnos et calamistros, Aurum et argentum, margaritas in vestibus, Pretiosum siue sumptuosum vestitum. Denique opera bona ab illis requirit et efflagitat, quae opponuntur omni impudicitiae, vanitati, ambitioni, fastui muliebri. Decent aurem maxime opera bona foeminas Christianas.  \pend\pstart Ac primùm similis est locus.1.Pet.vers.3.Tit.2 vers.2.Isai.3.  \pend\pstart Dissimile videri potest exemplum et factum Estherae: Esth.cap.5 et ludithae, quae se comit et ornat aliquid actura egregium, Respond. Non agi illic de precibus, nec de coetu Ecelesiae adeun\pend
\section*{AD I. PAVL. AD TIM. }
\marginpar{[ p.78 ]}\pstart do:sed de marito rege, eóque infideli demulcen do: item de castris infidelium ingrediendis. Nec hoc ipsum exemplum satis est tutum imitari. Obstare etiam videtur cap. 24.Genes.vers.47. vbi annuli et ar millae dantur Rebeccae sponsae Isaaci Resp. Non ad lasciuum ornatum, sed vt esset pignus futuri coniugii, ista Rebeccae fuisse et missa ab Abrahamo et donata.  \pend\pstart Certum autem est, etsi de mulieribus Christianis agatur hoc loco: haec tamen praecepta etiam pertinere ad viros Christianos, quibus eo turpiùs est comi lasciue, sumptuose et dissolute ornari quod sunt viri et mares. Nam etiam Ethnici homines sic docent.  \pend\pstart Sint procul a nobis iuuenes vt foemina compti Fine coli modico forma virilis amat. et M. T.libus 1. De off.sic. Cùm autem pulchritudinis duo genera sint, quorum in altero venustas sit: in altero dignitas: venustatem muliebrem dicere debemus: dignitatem virilem. Ergo et a for ma remoueatur omnis viro non dignus ornatus: et huic simile vitium in gestu motuque caueatur.  \pend\pstart Et paulo post. Adhibenda est praeterea munditia non odiosa, neque exquisita nimis: tantùm quae fogiat agrestem et inhumanam negligentiam. Haec ille. Patres etiam scriptoresque Ecclesiastici veluti Tertullian. in liber  de Habitu mulierum et Cyprian. in liber  de Cultu virginum nobiscum faciunt, qui haec praecepta ad viros perspicue referunt. Idem in concilio Gangrensi cap.21. et est in canon Parsimoniam distinct. 41. ne res noua aut nullius momenti esse censeatur. Idem etiam statutum est aliis synodis veluti Laodicena.  \pend
\section*{CAPVT. II. }
\marginpar{[ p.79 ]}\pstart Praeterea haec ipsa praecepta non ad solas mulieres, quae virginitatem vouent inter Papistas, et Nonnae dicuntur, pertinent: sed ad omnes et coniugatas et celibes, et viduas, quemadmodum praeclare hoc loco docet Chrysosto. et Gregor. Nazianzenus in liber  περὶ γυυαικῶν κεκαλωπισμένων. Nam tempore Pauli nondum erant monasteria, nec istae virgines Vestales et Nonnae visae et natae. Hae enim cultus diuini corruptelae sub Vitaliano Pontifice Romano maxime receptae et probatae videntur, sub quo tria haec tanquam venena Christianae pietatis praesentissima nata sunt, et constituta, Monachorum coenobia erecta, Reges in Monachos detonsi, Scortatio quotidiana in statum sanctum canonizata.  \pend\pstart Quaesitum vero est de primariis mulieribus veluti reginis, principissis, duchissis et huiusmodi aliis primariis foeminis, vtrum ad eas quoque haec praecepta sint extendenda. Respond. Etsi discrimen personarum etiam in Dei Ecclesia habendum est, omnia tamen etiam in illis, si modu sunt Christianae, ad modestiam, et pudicitiam composita esse debent, quanquam illis multa concedi et possunt et debent pro ratione dignitatis, quae aliis priuatis et priuatorum mulieribus non sunt permittenda. Nam et Solomonis vxoris cultus sumptuosus ille quidem et magnificus describitur laudaturque Psal. 45.  \pend\pstart Denique haec eadem praecepta locum habent in mulieribus, etiam cùm domi sunt, non tantùm cùm prodeunt in publicum, et conueniunt in coetu Ecclesiae. Perstringit autem Paulus breuiter omnem eum ornatum ascititium, qui est vel In i\pend
\section*{AD I. PAVL. AD TIM. }
\marginpar{[ p.80 ]}\pstart pso corpore nostro velut cincinni, vel Circa corpus, qualis est vestitus: cultus noster nimius quem damnat, siue ille sit nimius ratione Materie, veluti quia sumptuosior, vel Additae exquisitaeque lautit iae, venustatis, aut operae manus. Ad primum pertinet ἰματισμὸς πολυτελής. Palliola enim illa fere esle solebant sumptuosissima, magnique pretii, vt ex disputatione Hieronym. aduers. Iouinia. intelligitur, purpurea scilicet. Ad secundum membrum, et ad exquisitas lautitias vestitus pertinent ornamenta margaritarum, auri, et phrygionicae operae quae adduntur. Irem torques, armillae, annuli, etc. huiusmodi de quibus Tertullianus liber  de Habitu mulierum: et Cyprianus quoque. Et pulchra est veraque illa Plauti sententia in Mostell. scen.3.Act.1.  \pend\pstart Postea nequaquam exornata est bene, si morata est male.  \pend\pstart Pulchrum ornatum turpes mores peius coeno collinunt.  \pend\pstart Sed tamen, nequa hic superstitio nascatur, auri vsum prorsus non damnat Paulus, quemadmodum nec margaritarum: vti nec Petrus Palliorum. sed tantum nimiam in iis lasciuiam, arrogantiam, sumptum, fastum, immodostiaeque plenam cultus muliebris rationem et Christianis indecoram. Quanquam enim longe potius animi in muliere dissoluti aut superbi ratio habenda est, quam externi ornatus: tamen nec externus iste cultus, qui impudicitiam aut proteruiam redolet, maleque modestae mulieri et probae conuenit, tolerandus est. Ac de ea ipsa re sunt leges sumptuariae a Magistratibus Christianis ferendae. Quae luxu etiam  \pend
\section*{CAPVT II. }
\marginpar{[ p.81 ]}\pstart vetus Ecclesia in Synodis repressit. Oportet enim eo modo mulieres Christianas vestiri, quo interna earum pudicitia apparere possit, etiam infidelibus hominibus. Nam vt, et recte quidem, suo seculo est conquestus Chrysostom. quaedam mulieres ita exornatae ad templum et coetum Domini incedunt, et ad sacram Domini Coenam sumendam accedunt, vt illic potius saltaturae, quam Deum precaturae videantur. Nec vero iis prodest haec exceptio, Sic culta viro et marito meo placeo. Nec enim illae, vt lenoni marito placeant, debent studere: et quae sic exornata, aut fuco pigmentata incedit, certe naturam formamque, quam a Deo accepit, odit. Sed neque haec altera earum ratio est audienda, Sunt haec omnia ex τῶν ἀδιαφορων genere. Neque  enim haec cùm offendiculum cuiquam praebent, sunt adiaphora: sed damnata sunt, iisque  est abstinendum. In summa cultum mulierum non damnat Paulus, modo sit Honestus, et Modestiae verecundiaeque Christianarum mulierum conueniens et consentaneus. Et recte etiam docet Plutarchus liber 6.Sympos.quaest.7. abstersionem sordium et spurcitiei a cultu inhonesto et sumptuoso differre.  \pend
\phantomsection
\addcontentsline{toc}{subsection}{\textit{10 Sed (quod decet mulieres pietatem spondentes) operibus bonis. }}
\subsection*{\textit{10 Sed (quod decet mulieres pietatem spondentes) operibus bonis. }}\pstart A'ντίθεσις est. Fastui enim et superbiae, atque lasciuae illi mulierum pompae superiori versiculo descriptae verum earum ornatum opponit, nimirum bona opera, id est, pro externo ornatu, vt loquitur Petrus, internum requirit qui in eo posi\pend
\section*{AD I. PAVL. AD TIM. }
\marginpar{[ p.82 ]}\pstart tus est, vt sit earum animus placidus et quietus, nullisque vitiis expugnari possit. Rationem autem addit ab officio ductam. Nempe, quod honesta et sancta opera vere Christianas mulieres et religionem profitentes deceant: non autem ille externus vestitus, quo pudor omnis ab illis proiectus et conculcatus videtur exulare.  \pend
\phantomsection
\addcontentsline{toc}{subsection}{\textit{11 Mulier cum silentio discito cum omni subiectione. }}
\subsection*{\textit{11 Mulier cum silentio discito cum omni subiectione. }}\pstart Alterum et secundum de mulieribus Christianis praeceptum, quod ad Ecclesiae publicique coetus politiam et decorum pertinet, Ne in eo loquantur, sed sileant, et discant. Similis huic locus est. 1.Corinth.14.vers.34. Id quod etiam consilio Chartaginensi cap.95. constitutum est, referturque  canon. mulier distinct. 23. quantumuis docta et pia sit ista mulier.  \pend\pstart Dissimilia tamen exempla videntur afferri posse. De Debora Iudic. cap. 4. et Olda mulieribus prophetissis 2.Reg. 22. versu 14. quae in Ecclesia docuerunt. Sed responsio parata est, illa nimirum exempla plane esse extraordinaria, itaque haec Pauli praecepta minime immutare. Affertur etiam illud quod Eusebus liber  1.devita Constantini scripsit, mulieres eligere e suo sexu solitas fuisse aliquas, quae eas docerent ne in publicis congressionibus conuenirent cum viris, quod Licinii imperatoris Rom. edicto, et lege in Syria simul cum viris congregari et conuenire vetabantur. Sed respondemus, Illud praeter morem Ecclesiae, et praeter Dei praecepta iussum esse a Licinio homine profano,  \pend
\section*{CAPVT II. }
\marginpar{[ p.83 ]}vt mulieres ipsae in coetu mulierum docerent, non autem cum viris ad preces et conciones conuenirent. \pstart Denique quod de Pepuzianorum haereticorum secta, et sententia dici posse videtur habet facilem responsionem. Voluerunt illi quidem, sed haeretici, mulieres esse posse? et presbyteridas in Ecclesia, et docere, et Sacramenta quoque administrare (id quod hodie Papistae ex Marcionitarum dogmate obstetricibus in casu necessitatis concedunt, vt nimirum paruulos baptizarent) Verùm propterea ab omnibus orthodoxis sunt Pepuziani damnati, atque etiam Synodo Chartaginensi a. vt diximus, et meritissimo.  \pend\pstart Hoc autem praeceptum Pauli locum habet, dum sunt in coetu publico mulieres: non autem cùm sunt domi suae, quia familiam suam instituere posfunt priuatim, maxime si desit paterfamilias, et sint viduae. Nam tunc precari Deum, et docere suos liberos Dei metum, et possunt et debent.  \pend\pstart Requirit autem Paulus et silentium, ne vices loquendi sibi dandas esse putent mulieres, quanquam fortasse, modestiae causa primae, in Ecclesia loqui nollent. Deinde ab iisdem subiectionem postulat, ne putent aequum, vt propter maritorum suorum dignitatem et opes:vel propter etiam familiam suam ipsae saltem secundum Pastores et Ministros verbi Dei aliquod munus Ecclesiasticum gerant. Hoc enim illis prorsus interdictum est a Spiritu sancto. Huius autem prohibitionis variae sunt rationes, quas affert Paulus sequentibus  versiculis.  \pend
\phantomsection
\addcontentsline{toc}{subsection}{\textit{12 Mulieri enim docere non permitto, }}
\subsection*{\textit{12 Mulieri enim docere non permitto, }}
\section*{AD I. PAVL. AD TIM. }
\marginpar{[ p.84 ]}
\phantomsection
\addcontentsline{toc}{subsection}{\textit{neque authoritatem vsurpare in virum, sed esse in silentio. }}
\subsection*{\textit{neque authoritatem vsurpare in virum, sed esse in silentio. }}\pstart Prima ratio quae explicatur hoc loco, a differentia, quae a Deo ipso constituta est inter vtrunque sexum virilem nempe et muliebrem, ducta est. Secunda vero a genere. Deus neque mulieres praeesse vult, neque docere, neque  in viros authoritatem vsurpare: At qui loquitur in Eccesia siue praeest ea et docet, in viros sumit quoddam imperium et authoritatem. Ergo mulieres loqui in Ecclesia neque possunt, neque debent. Nec obstat quod eas in Diaconatum alleget Paulus inf. cap. 5.vers.9. Nudum enim illae manuum ministerium praebent, Diaconis autem ipsis parent. Sed ex hoc loco de eo quoque inter quosdam disputatum est, Vtrum honestum sit mulieres regnare, id est, viris imperare, et summum imperium et ius in viros et mares obtinere. Id quod in Hispania, Anglia, Scotia variisque aliis regionibus locum habet. Cuius etiam rei extant exempla in Semiramide Assyriorum regina, Candace Ethiopum Actor.8. vers.27. Cleopatra Ægyptiorum sub Augusto, et Zenobia fortissima muliere sub Adriano imperatore (cuius Zenobiae imperio etiam multae Ecclesiae Christianae paruerunt. In populo autem Dei nihil quicquam tale habemus, siue Iudaicum siue Israëliticum regnum spectetur. Nam quod de Athalia potest afferri 2. Reg. 11. facile diluitur: fuit enim vsurpatio et inuasio iniusta regni ea Athaliae dominatio: non autem regnum et potestas legitima cui populus lubens assentiretur vt postea apparuit. Quare  \pend
\section*{CAPVT II. }
\marginpar{[ p.85. ]}\pstart uste quoque illa a summo sacrificatore loiada de regni solio deturbata et deiecta, est, atque etiam propter vsurpatum regnum necata. Quod de Amazonum regno commemoratur, primum vt fabulosum non fuerit (quemadmodum omnino fabulosum non putant Arrianus, et Q Curtius quique res gestas Alexandri Magni scripserunt) tamen praeter rerum naturam constitutum est illud regnum, et in virilis sexus apertam contumeliam. Quid enim magis portentosum et monstrosum et a muliebris ingenii mansuetudine alienum, quam videre armatarum mulierum exercitum contra viros concurrentem: ac latissimam quandam terrae plagam cernere, quae omnino viris et maribus careret. Itaque videntur illae contra ipsam foeminei sexus vel potius humanam naturam pugnasse, et auxilium, quod a Deo optimum et praestantissimum paratum est mulieribus, nempe maritos reiecisse. Vnde prorsus breuique  tempore tota illa Amazonum gens deleta est. Sed ad propositam quaestionem reuertamur. Quanquam vero Isai.3. vers.II. pro magno Dei maledictionis signo positum legimus, quod pueri et mulieres imperium in aliqua gente obtineant: tamen illud ipsum non. est perpetuum. Saepe enim et pueri reges, velut Solomon et Iosias, sanctissime et felicissime regnarunt, fuitque eorum imperium omni bonorum genere a Deo cumulatum. Idem de mulieribus quibusdam, et earum imperio dici potest, quibus Dominus mirum in modum benedixit, vt ex variis historiis apparet. Certe de illustrissima regina Angliae Elizabeta, quae nunc felicissime regnat, dici potest, nihil terrarum orbem vidisse.  \pend
\section*{AD I. PAVL. AD TIM }
\marginpar{[ p.86 ]}\pstart vnquam illius regno felicius et optabilius. Sapientertamen sibi consuluisse videntur ii populi, qui legibus suis et publico iure cauerunt, ne inter se, et in sese foeminae dominarentur, summumque ius et imperium haberent, si genus ipsum muliebre cum virili conferamus, quia ad multa munia, quae regni administratio requirit, sunt illae minus propter sexus sui naturam et imbecillitatem a ptae et inhabiles quale est praeesse exercitui, ius dicere in publico sedentes. Quae certe res dedecent prorsus muliebrem verecundiam. Vnde iure ciuili Roman. L. mulieres D. de Regul. Iuris mulieres a virilibus muneribus et officiis administrandis recte arcentur: qualia sunt etiam quae in publico exerceri debent. Et illud est August. libus 1. de Nupt. cap.9. Nec dubitari potest viros potius foeminis, quam foeminas viris principari. Vnde non potest mulier esse Roman. imperatrix et regina, et cùm Athalia in ludaea et Irene mater Constantini tertii imperium Constantinopoli. gerere voluit, vtraque omnia subuertit, Idolorum cultum in Dei Ecclesiam inuexit, et haec Saracenis Roma. imperium lacerandum obiecit. Vnde et Carolus Magnus fuit tunc temporis in Romanum imperium asciscendus in Occidente: et Nicephorus in Oriente. Vbi autem inferiores iurisdictiones quales Ducum, Comitum, Baronum, Castellanorum, sunt patrimoniales, vt in Gallia, illae a mulieribus haberi et possideri possunt, meo quidem iudicio, quia non sunt summa imperia istae dignitates, et officia: sed tamen ab iisdem mulieribus nec possunt, nec debent hae iurisdictiones exerceri: sed a viris per eas delegatis. Quan\pend
\section*{CAPVT II. }
\marginpar{[ p.87 ]}\pstart quam hoc omnino perperam et pessime receptum est vsquam gentium, vt sit iurisdictio vlla pars patrimonii, et reditus nostri ac dominii, sed tam late patet et grassatur in omnia auaritia, vt etiam res sacratissimas, qualis est Magistratus, fecerit patrimoniales et in censu numeret : non autem virtuti et doctrinae eorum, qui sunt capaces eorum munerum, tribuat. Vide quae scribuntur can. Mulierem 33. quaest.5. de hoc argumento.  \pend
\phantomsection
\addcontentsline{toc}{subsection}{\textit{13 Adam enim prior formatus est, deinde Eua. }}
\subsection*{\textit{13 Adam enim prior formatus est, deinde Eua. }}\pstart A’ιτιολoyίa est saperioris sententiae ab ordine naturali ducta, quem Dominus Deus inter vtrumque sexum constituit rerum initio, et statim ab ipsa vtriusque creatione. Is autem fuit, vt pareret mulier. Vir autem caput et superior esset. Id quod ex ipso ordine creationis apparet. Adam enim, qui est mas et vir, prior creatus est. Eua autem, quae est foemina, posterior. Sed tamen eodem die vt docetur tum Genes. cap.I.vers.27. tum 2.vers. 18. Non tam autem ex natiuitatis ordine Paulus argumentatur et probat virum superiorem esse: quam ex fine creationis mulieris. Nam creata est foe mina, vt esset viro in adiutorium, et ad eum, tanquam suum propriumque finem respiceret. Quae autem ad finem aliquem destinantur illique  subseruiunt, sunt eo fine minora, inferiora, ac illi subiecta. Id quod verum esse ratio omnis et naturalis et Philosophica docet. Nam alioquin a solo natiuitatis et productionis tempore, et ordine ratio ducta nihil efficit et concludit. Pisces enim  \pend
\section*{AD I. PAVL. AD TIM. }
\marginpar{[ p.88 ]}\pstart herbae, sydera et reliquae omnes creaturae, quae homini subiiciuntur ex Dei praecepto Psal. 8.praeferendae tamen homini essent contra mentem et Christi, et Pauli. Omnia enim propter hominem facta sunt, etiam Sabbatum ipsum: non autem homo propter Sabbatum.  \pend
\phantomsection
\addcontentsline{toc}{subsection}{\textit{14 Et Adam non fuit seductus: sed mulier seducta, causa transgressionis fuit. }}
\subsection*{\textit{14 Et Adam non fuit seductus: sed mulier seducta, causa transgressionis fuit. }}\pstart Alterum argumentum a causa tum efficiente, tum incurrente sumptum. Efficiens igitur causa est Dei praeceptum et voluntas, quae eam mulieri poenam inflixit. Incurrens autem, transgressio mulieris, propter quam talem poenam merita est et huic seruituti et subiectioni addicta. Æquum enim fuit, vt quae seipsam regere non potuit, alterius consilio et imperio regenda subiiceretur a Deo. Est autem subiecta viri imperio et potestati, vt est Genes. 3. Quanquam autem illud Mosis proprie pertinet ad virum et mulierem coniuges inter se: non autem vt quilibet viri in quaslibet mulieres imperium sibi arripiant:tamem ex eo intelligitur, in quem gradum viri supra mulieres euecti sint a Deo. Propter transgressionem autem mulieris facta iam est durior et seruilior mulieris conditio, quae ad virum tantùm suum tanquam ad suum finem tunc spectabat. Sed nunc aerumna et iugum additum est huic seruituti, quae prius erat omnino et liberalis et voluntaria. Nunc enim dura est, etiam cùm est ad virum et maritum proprium. Ex quo ipso apparet Deum et confirmasse priorem subiectionem, et eam propter  \pend
\section*{CAPVT II. }
\marginpar{[ p.89 ]}\pstart peccatum effecisse duriorem et austeriorem. Valde autem notandum est, quod ex verbo sedacta colligit Bernard. sermo. de Duplici Baptis. his verbis, Serpens ô Eua decepit te, non impulit, aut coegit, mulier tibi, ô Adam de ligno dedit: sed offerendo, non vtique violentiam inferendo. Neque enim potestate illius, sed tua voluntate factum est, vt eius voci plus obedieris, quam diuinae.  \pend
\phantomsection
\addcontentsline{toc}{subsection}{\textit{15 Seruabitur tamen liberos gignendo, si manserit in fide, ac charitate, et sanctificatione cum modestia. }}
\subsection*{\textit{15 Seruabitur tamen liberos gignendo, si manserit in fide, ac charitate, et sanctificatione cum modestia. }}\pstart Hypophora est quae magnam consolationem tamen continet, ne se prorsus spe salutis orbatas esse mulieres existiment. Id quod illis potuit in animum incidere, vel nobis etiamipsis, dum eas esse nobis tanti causas mali audimus. Per eas enim factum est, vt vniuersum genus humanum siue virile, siue etiam muliebre periret. Ergo prorsus salutis exortes videbantur. Sed contra Paulus. Hanc ipsam poenam transgressionis, et hanc in coniugio subiectionem mulierum docet in ipsarum commodum et salutem cedere mira Dei sapientia nimirum, qui e tenebris lucem eduxit, vt est 2. Chorinth. 4. vers. 6. Sic conuertit hoc suum in eas iudicium Dominus in earum ipsarum consolationem et salutis viam, dum hanc subiectionem erga maritos vocationem earum propriam esse statuit. In quo illud primum est obseruandum vocem διὰ non denotare hoc loco causam efficientem: sed mediam tantùm, per quam, tanquam per iter a Deo demonstratum, est illis  \pend
\section*{AD I. PAVL. AD TIM. }
\marginpar{[ p.90 ]}\pstart pergendum et progrediendum. In quo peccant Papistae et Pelagiani, omnesque Operistae, et quidem grauiter, qui salutis nostrae causas in nostris operibus quaerunt et constituunt. Quid enim de viduis, quae nullos liberos vnquam susceperunt: de virginibus, et iis etiam quae cùm sint coniuges, nullos liberos tamen pepererunt, futurum esset (vt recte hoc loco Chrysostom. animaduertit) si salus mulierum ἑκ τεκνογονίας tota pendet? Certe saluae non fierent. Sed et Thomas Aquinas hic delirat, dum laudes virginitatis et celibatus cum hoc Pauli dicto conciliare, et placare nititur, quod cum illis aperte et perpiscue pugnat. Ergo hic vocatio et munus mulierum describitur etiam in eo ipso, in quo poenam sui peccati sentiunt. Est autem, vt liberos non tantùm suscipiant et gignant: sed vt educent, curent, et onera illa primae liberorum nutricationi, scerte quidem molestissima, lubentes volentesque deuorent et suscipiant. Tεκνογονίας enim hoc loco non tantùm parturitionem ipsam mulierum significat: sed quicquid illi adiunctum est, et imminet iustae matrum curae, vt liberi editi educantur. Vnde merito matres puerperae damnantur, quae filios, si possunt, suos non alunt, de quo Plutarchus in liber  de Liberorum Educatio. et dictum Gregorii in canon. Ad eius distinct.5. Mille quidem sunt taedia, quae contra Iouinia. enumerat Hieronym.tam minutim, vt ipse puerpera fuisse aliquando videatur, qui tamen perpetuo vixit celebs. Mille etiam dolores et molestiae non tantùm leuiter degustandae: sed diutissime deuorandae, tum in partu, tum post partum. Quae tamen omnia si libenti et grato animo de\pend
\section*{CAPVT I. }
\marginpar{[ p.91 ]}\pstart uorat, atque suscipit mater, Deo grata est, et in sua vocatione incedit: ad trimatum autem vsque lex Ciuilis iubet vt mater ipsa alat liberos suos, eósque praecipue curet.  \pend\pstart Sed addit Paulus alias quoque virtutes, quas esse necessarias ostendit, vt mulier Deo placeat, quarum aliae communes illis sunt etiam cum viris, nempe Fides, et Charitas: aliae illarum propriae, Sanctificatio, et Modestia. Quae enim virtus nos Deo commendat est fides. Illa enim nos et deducit ad Christum, et cum eo coniungit. Illa charitatem quoque gignit, extra quam nihil a matribus ipsis fieri Deo gratum potest, quia si inuitae, repugnantes, et coactae potius officium faciunt, quam ex charitate, displicent Deo.  \pend\pstart Sanctificatio, quam requirit a mulieribus, non est tantùm generalis omnium vitae actionum ad Dei voluntatem conformatio 1.Thess. 4. vers. 4. sed maxime sanctimonia quaedam corporis, qua se vxores et mulieres ab omni turpitudine, immunditia, et lasciuia puras conseruant, quae impudicitie opponitur: et alio nomine dicitur castitas. Modestia, luxui et sumptuoso illi cultui, de quo superius dixit, opponitur.  \pend
\endnumbering\beginnumbering\section{CAP. III.}
\phantomsection
\addcontentsline{toc}{subsection}{\textit{\huge\textbf{F}\normalsize IDVS est hic sermo, Si quis Episcopatum apperit, praeclarum opus desiderat. }}
\subsection*{\textit{\huge\textbf{F}\normalsize IDVS est hic sermo, Si quis Episcopatum apperit, praeclarum opus desiderat. }}\pstart Transitio est. Post quam enim de pastorum of\pend
\section*{AD I. PAVL. AD TIM. }
\marginpar{[ p.92 ]}\pstart ficio disseruit, de eorundem electione, qualis esse debeat, agit, vt in ea parte regimen Ecclesiae et ratio illius instituendae et conseruandae esset nobis certa, Deique verbo praescripta. Commode tamen et molli transitione ad hunc locum et argumentum descendit Paulus, quasi per hypophoram, quae ex superiori disputatione, per quam mulieres a docendi munere et Episcopatu semotae sunt, oriri videbatur. Cùm enim sit officium et munus virile Episcopatus, videbantur omnes viri huius tantae dignitatis et muneris capaces: quod tamem negat esse verum Paulus, cùm certas quasdam conditiones in iis requirat, qui sunt ad hoc munus vocandi, vt sint idonei:propter quas caeteri ab eo munere arcentur. Ergo vt nullius mulieris proprium munus est, Episcopatus: ita non cuiuslibet viri opus est. Est autem vtilissima haec disputatio, tum vt interdictae mulieribus huius administrationis ratio intelligatur, quae sumitur propter ipsius difficilem et arduam curam, cuius non sunt illae capaces:tum vt hoc ipsum munus in sese quantum sit, quamque excellens et eximium demonstretur, ne ad illud quisquam temere cooptetur: aut ne iam cooptati et vocati perfunctorie et leuiter eo defungantur, sed intelligant, quae ab ipsis requirantur. Denique vt esset certa quaedam regula ab ipso Dei spiritu, non ab hominibus praescripta, ad quam legitimae vocationes in Ecclesia renocarentur, et examinarentur. Itaque non licet hominibus priuatis vel Magistratibus: sed ne synodis quidem ipsis aliquid de electione personarum Ecclesiasticarum ab his praeceptis dissentiens vel contrarium statuere. Nam hos canonas  \pend
\section*{CAPVT  III. }
\marginpar{[ p.85 ]}\pstart In electione autem personarum Ecclesiasticarum tria quaeri et possunt, et solent. Nempe 1 Qui eligendi sint. 2 A quibus eligendi. 3 Quomodo eligendi. Quod vltimum caput de forma electionis agit. Hic de personis, quae eligi possunt alibi in hac ipsa, et de iis, ad quos ius eligendi pertinet et de legitima electionis forma: dicetur. Primùm enim est natura, vt sint quaedam personae, quae eligi possint, de quibus idcirco primo loco fuit dicendum.  \pend\pstart Agit autem Paulus hoc toto capite de Episcopis, et Diaconis. Attamen aliae sunt Ecclesiasticae et necessariae dignitates praeter Diaconos et Episcopos, quales Presbyteri, de quibus nihil penitus hoc loco traditum esse a Paulo quidam existimant. Volunt enim Presbyterorum elelectionem explicari in Epistola ad Titum cap.1. Hic vero Episcoporum. Sed falsa est eorum sententia. Eosdem enim vocat Paulus Presbyteros et Episcopo, svti apparet ex cap.1. vers.5, et 7. Epistolae ad Titum. Atque etiam ex 1. Pet.5 vers.1 et annotat cum Chrysostomo Hieronym. ex quibus Isidorus didicit, cuius dictum est in canone Clero. Et clericos 5 Presbyter distinct.21. Itaque priusquant du ba, juu metradit d praecipit Paulus, accedamus, illud explicandum est, vtrum idem sit Episcoporum, et Presbyterorum, de quibus in sacra scriptura et praesertim in Apostolicis scriptis fit mentio, nomen, dignitas, vocatio, munus, et praestatio. Nam Diaconi ab Episcopis hic distinguntur et separantur aperte. Presbyterivox. Graecis est πρεβύτερος. Itaque est Graeca  \pend
\section*{AD I. PAVL. AD TIM. }
\marginpar{[ p.94 ]}\pstart differunt autem hx voces παλαιὸς et πρτσβύτερος vti vetus et senex apuu Latinos, quod παλαιὸς est aetatis, et πρεβύς etiam dignitatis, vt ex Plutarcho in Nicea colligi potest. Horum igitur Presbyterorum (qui a dignitate non autem ab aetate hic nominantur) in scriptura duplex est genus. Vnum eorum qui tantùm moribus inuigilant. Alterum eorum, qui et moribus et doctrinae attendunt, et in vtroque laborant. Id quod ex hac ipsa Epistola cap.5.vers.17. facile colligi potest. Quod ad primum igitur eorum genus, vti a secundo distinguitur: ita ab Episcopis et doctoribus est secernendum. Sunt enim illi Presbyteri tantùm morum et disciplinae custodes et inspectores: non etiam praecones doctrinae. Cur autem hoc genus fuerit in Ecclesia necessarium vel inductum, ratio mihi nulla videtur alia esse, quam quod etiam illo tempore, quanuis multis donis abundaret Ecclesia Dei, tamen non satis multi idonei viri ad vtrunque hoc munus obeundum reperirentur, vt etiam Epiphan.testatur. Ergo ad formam veteris et Iudaici Synedrii habuit Ecclesia Christiana hos Presbyteros disciplinae φύλακας, qui non erant omnes Leuitae, et concionatores verbi Dei. Et haec quidem de primo Presbyterorum genere, qui morum tantùm curam habent in Ecclesia Dei. Ergo de secundo Presbyteroiam genere agamus, eorum nempe, qui etiam in doctrina laborant, et qui morum, atque etiam regiminis Ecclesiae curam habent. Hos nihil ab Episcopis differre dicimus et, Dei verbo affirmamus siue ex eorum electionem, siue munus, siue personae conditionem spectes: sed iis omnino pares esse debere: et cùm recte  \pend
\section*{CAPVT . III. }
\marginpar{[ p.95 ]}\pstart administrata est Ecclesia, semper fuisse his rationibus demonstramus.  \pend\pstart Prima ratio, quod quos Paulus iubet Presbyteros oppidatim creari, eosdem vocat idem Episcopos Tit.1.vers.5, et 7.  \pend\pstart Secunda. Quod cùm Presbyteros Ephesiorum alloqueretur Paulus quos accersiuerat, et euocauerat Miletum, eos omnes Episcopos nominat, et iis eadem praecipit quae Episcopalis muneris pro pria sunt, et postea ad solos Episcopos pertinere existimata et definita sunt Actor. 20.v.17.18.  \pend\pstart Tertia, Quod Petrus tribuit Presbyteris hoe ipsum munus, vt pascant gregem Domini, et illius curam habeant 1. Pet.5.vers.2.3, Praelatos enim omnes Presbyteri vox complectitur ( vt explicat Bernard. sermo. in Psal. Audiam) quorum est prodesse magis, quam praeesse.  \pend\pstart Quarta, Quod ad Presbyteros negotia Ecclesiae deferuntur, ab iisdem illa geruntur, et adminhtrantur. Atqui hoc ipsum alibi tribuitur Episcopis Philip.1.vers.1. Actor.11.vers.3o, 16. vers.4. Itaque omnes aequali nomine ηγόμεναι a scriptura appellantur Hebus 13.vers.7. Vnde Cyprian. in liber  de Simplicit. Praelatorum, Vnus est Episcopatus, cuius in solidum singuli participes sumus.  \pend\pstart Quinta, Quod plures Episcopos in singulis ciuitatibus fuisse, praeterquam Alexandriae, quae est vrbs Ægypti, scribit Epiphanius haeres.68. At si Episcopi nomem esset προεςασίας vox, non possent eiusdem ciuitatis plures esse Episcopi, quia in eodem coetu plures esse primi, et capita non possunt.  \pend\pstart His autem obstant quaedam, quae in contrarium afferri videntur posse. Primùm quod haeretici ha\pend
\section*{AD I. PAVL. AD TIM. }
\marginpar{[ p.96 ]}\pstart biti sunt, qui Presbyteros Episcopis adaequarunt, quales Herianrre docet Epiphan. et Augustinus in liber  deHaeresibus  ad Quoduultd. vbi nos diximus  \pend\pstart Secundum, Quod praeter Presbyteros habent Episcopi apicem Pontificatus et praeeminentiae. Itaque multa licent Episcopis in Presbyteros hodie, quae non item Presbyteris in Episcopos. Solus enim Episcopus creat Presbyteros: nulli autem Presbyteri soli creant Episcopum, vt est in cano. Cleros et clericos 5 Presbyter distinct. 21. et tota distinct. 67. Nan in Episcopo omnium ordinationum potestas inest, ait Ambrosius.  \pend\pstart Tertium Quod vt libro 1.cap.6.refert Ruffin. Synodo Nycena cap. Io. vetitum est, ne in vna ciuitate duo essent Episcopi. Itaque ex vetere consuetudine Eccleliae vnus e toto Presbyterorum collegio caeteris praeficiebatur, ad quem velut caeterorum praesidem et praepositum Ecclesiae negotla reserrentur; quei ee lνειdτα et ωροε͂δρωυτα Iustinus Martyr in 2. Apolog.nominat.  \pend\pstart Quartum, denique quod aequissimum videatur vti in politia ciuili et humana: ita in politia Ecclesiastica vnum aliquem praeesse caeteris, qui velut totius chori dux sit et moderator perpetuus.  \pend\pstart Respond. Et quidem ad prima tria argumentas idque , ex Hieronymo in Epistola Pauli ad Titum vbi sic ait, Sicut Presbyteri sciunt se ex Ecclesiae consuetudine ei, qui sibi praepositus fuerit, esse subiectos: ita Episcopi nouerint se magis consuetudine, quam dispositionis dominicae veritate Presbyteris esse maiores, et in commune debere Ecclesiam regere. Nec enim hoc προεδρίας et praeeminentiae perpetuae ius, quod vni tributum est,  \pend
\section*{CAPVT  III. }\pstart ex Dei verbo sumptum est: sed humanum plane commentum est dum nimium Ecclesiasticae politiae formam et regimen ad ciuilis disciplinae normam homines exigunt contra praeceptum et Christi Luc 22. vers.25. et Petri 1. Pet. 5. vers.3. Idem docet Augustinus Epist. 19. haec honorum diuersa vocabula et distincta ex hominum tantùm consuetudine, non ex Dei verbo introducta et recepta in Ecclesia esse. Nam neque tempore Apostolorum haec distinctio in Ecclesia fuit: neque Epaphroditum, quem remittit Philippens.neq: Archippum, quem Collossis agentem exhortatur Paulus, aut Episcopum, aut Presbyteris praepositum appellat. Imo Irenae. Higinum Episcopum Roma. semper Presbyterum vocat, vt docet Eusebus liber 5 cap.26. nunquam autem Episcopum, aut Presbyteris praepositum. Et cùm nihil sit ad Dei verbum addendum dicimus omnino temere, id introductum fuisse et receptum in Eclesiis Dei: et haec illius tyrannidis, quae postea miserrima irrepsit in Ecclesiam, verissima semina et fundamenta dico fuisse: vti apparet ex Epist.3. Ioannis vbi Diotrephes φιλόπρωτέυων damnatur v.9. Atque  varia esse huius proestasiae incommoda docet Caluinus cùm Epist. 375. tum in suis commentariis in Epist. ad Philip. cap.1. quanquam is fuit vetus mos, et iampridem receptus, vt apparet ex Nicena Synodo prima vti ex Ruffino docuimus. Vnde igitur et cur primum est introductus? Res. Primùm quidem, vt ex Epiph. et Hieronym. apparet, hoc Alexandriae institutum fuisse, praeter morem caeterarum Ecclesiarum. Cur autem, ratio est, vt ait Hieronym. nimirum vt scismatum  \pend
\section*{AD I. PAVL. AD TIM. }
\marginpar{[ p.98 ]}\pstart occasio tolleretur, et dissentionum plantaria euellerentur. Ita ille in Epist. ad Euagrium, cuius dictum refertur in canon. Leginus in Isaia dist. 93. et ipse repetit in Epist. Pauli ad Titum cap.1. Contra vero res ipsa docuit, istud fuisse verissimum vniuersalis a fide Christi apostasiae semen et fundametum: et laeumimnae fius tyranidis, quae in Papatu etiamnum hodie viget, fulsimentum maximum.  \pend\pstart Quod ad quartum argumentum attinet, prorsus elumbe est. Neque enim regimen Ecclesiae est reuocandum ad ciuilis politiae rationem et normam vt docet Christus Luc.22.vers.25.1. Pet. 5 vers.3. quia Dominus ipse, qui suae Ecclesiae est rex, eius quoque regendae regulam nobis verbo suo praescripsit prorsus ab Monarchica alienam. Denique tantum inde malum serpsit, id est, ex hac Episcoporum praelatura, vt quemadmodum scribit contra Luciferia. Hieronym.non satis recte censerentur baptizari, qui a Presbyteris erant baptizati, nisi postea manus illis ab ipso Episcopo et praeposito isto imponerentur: a quo solo dari spiritum sanctum crediderunt. Atque inde putidum illud confirmationis suae Sacramentum Papistae fabricati sunt.  \pend\pstart Veniamus ad verba Pauli. Commendat futuram tractationem, tum a Modo ipso dicendi, tum a Re ipsa. Modus dicendi est, praefatio quam attexuit, quo futuram doctrinam veram esse, et retinendam intelligeremus, velut Dei ipsius vocem:non autem existimemus haec praecepta posse a nobis sperni, et quae hic traduntur, inter ἀδιάφνρα recensenda reponendaque esse. Est enim solius Dei in  \pend
\section*{CAPVT  III. }
\marginpar{[ p.99 ]}\pstart sua domo vocationes constituere.  \pend\pstart Res autem ipsa, est dignitas Episcopatus. Episcopi vero vox ex vsu communi sumpta est ab Apostolis. Dicti sunt enim Episcopi a Graecis qui rei alicui cum aliqua potestate praeerant. Suidas ab Atheniens.scribit esse dictos, qui ad res sociorum inspiciendas mittebantur: Plutarch. tum in Numa, tum in Pericle hac voce vtitur, quam esse rem pulcherrimam et praestantissimam docet hoc loco Paulus. Itaque non esse quoslibet illius capaces, χαλεπὰ γκ̀ δύσκολα τὰ καλὰ, vt ex veteri prouerbio dixit Plato: sed multa in iis requiri, qui ad tantum muhus obeundum le acemgunes aut a nobis eligi, vt sint idonei, debebunt. Quidam enim sunt nimium audaces, qui nec se, nec munus hoc, quod affectant, norunt. Itaque sunt illi tanti muneris difficultate deterrendi.  \pend\pstart Ne propterea tamen alii hoc munus tam arduum repudient, adiicit consolationem, eos, qui hoc munus expetunt, recte facere, et opus Deo gratum exoptare. Arque ita Paulus, Audaces reprimit, et Timidos excitat. Christi enim hominis imaginem gerit Episcopus, vti Magistratus ipsius Dei, vt Augustinus in liber  Quaestio veter. et noui Testamenti quaest.35.scribit. Loquitur autem hic Paulus non de ambitioso desyderio. Hoc enim maxime in iis damnatur, qui Ecclesiae praefici volunt, non item de auaro: sed de sancto illo animi affectu, per quem Dei gloriae prodesse, eanque promouere. i. quantum Dominus ipse legitime dederit, auent, qui Episcopatum desydetant: nec praeoccupantes locum et sese intrudenres,nec fugientes, cùm precibus Ecclesiae flagi\pend
\section*{AD I. PAVL. AD TIM. }
\marginpar{[ p.100 ]}\pstart tantur. Et vera est illa Augustini sententia liber 19 de Ciuitate Dei cap. 19. Otium sanctum quaerit charitas veritatis: negotium iuitum suscipit necessitas charitatis.  \pend\pstart Episcopatum autem definit Paulus ἔργον non τιμὴν, quia, vt ibidem dicit, et recte quidem Augustinus, episcopatus est nomen operis et oneris: non autem honoris. Vae igitur desidiosis et canibus mutis, qui se tamen vocant Episcopos, vt est apud Isai. 56. Videtur autem haec sumpta definitio ex Exod. cap. 39. sed imprimis vers.41. vbi quae Aaron et Leuitae habere singularia et peculiaria iubentur, ad ministrandum Deo, non ornamenti causa illis data esse commemorantur: sed potius ministerii et operis illis a Deo iniuncti gratia.  \pend
\phantomsection
\addcontentsline{toc}{subsection}{\textit{2 Oportet igitur Episcopum irreprehensum esse, vnius vxoris virum, vigilantem, sobrium, compositum, hospitalem, aptum ad docendum. }}
\subsection*{\textit{2 Oportet igitur Episcopum irreprehensum esse, vnius vxoris virum, vigilantem, sobrium, compositum, hospitalem, aptum ad docendum. }}\pstart Varias dotes postea recenset Paulus, quibus ornatum essefpiscopum oportet, vt sit maioris inter suos authoritatis, et tanquam exemplar totius gregis, vt docet Petrus 1. Epist. cap.5. Manant enim in vulgus exempla regentum, et pastorem oportet non tantùm doctrina: sed et vitae probi tate ac sanctitate caeteris praelucere. Matth.5.vt sit sal, quia quodammodo in ipsius Dei doctrinam, quam praedicant, redundant pastorum vitia et probra. Vnde fit, vt etiam in lege veteri symbolis quibusdam doceret Dominus quantas dotes in Sacerdot ibus  \pend
\section*{CAPVT  III. }
\marginpar{[ p.101 ]}\pstart Sacerdotibus inesse oporteret. Omnis enim ille cultus et summi Sacerdotis ornatus, qui Exod.28. vers.3o.31. et sequent. describitur, quanta pietas, sanctitasque a pastore Ecclesiae requiratur, ostendit: quemadmodum etiam explicat Isychius liber  6.in Leuitici cap.21, et 22. Externa enim illa Leuitici ministerii species et splendor fuit interni ornatus et vitae symbolum. Ex iis autem tot virtutibus, quas describit Paulus hoc loco, vt veram perfectamque boni Episcopi ideam habeamus, quaedam ostendunt, qualem Episcopum ipsum In sese esse oporteat, et Ratione suae familiae, Item, Extraneorum et infidelium, Quae omnia distribuo in tria haec capita, vt congestas hic a Paulo dotes distinguere aliqua ratione possimus. Berna. autem Sermo.77. in Cantica breuiter has omnes dotes videtur his verbis complexus. Ad solam ciuitatis, id est Ecclesiae custodiam, vt quantum satis est, procuretur, opus est viro forti, spirituali, fideli: forti ad propulsandas iniurias: spirituali ad deprehendendas insidias: fideli, qui non, quae sua sunt, quaerat. Sed etsi doctus fuerit, non sit autem bonus: verendum, ne non tam nutriat doctrina vberi, quam sterili vita noceat. Idiota autem pastor quomodo in pascua diuinorum eloquiorum deducet greges dominicos?  \pend\pstart Eo vero diligentius ista Pauli praecepta sunt obseruanda. Primum, Quod ex quibus causis Episcopi eligi prohibentur: ex iisdem electi deponuntur. Eas autem hic expressas a Spiritu sancto habemus. Secundum. Quod praeter has conditiones et qualitates, aliae sunt synodorum decretis postea additae, quarum quae dam prorsus  \pend
\section*{AD I. PAVL. AD TIM. }
\marginpar{[ p.102 ]}\pstart sunt reiiciendae, nec nobis Euangelii lucem adeptis iam retinendae, vel in Episcoporum electione requirendae, et obseruandae. Nec fere locus est, quem scupulosius canonistae tractarint, quam hic qui est de clericorum et Episcoporum electione vti apparet tota prima parte decretorum et a distinct.23.ad 1o2. Nos quantum videbitur, et ad vsum nostrum, et ad Ecclesiae Papisticae perturbationem et συνχύπν patefaciendam, ex illis etiam hauriemus. Primùm igitur ait, Irreprehensibilem, vel irreprehensum. Paulus in Epistola ad Titum 1. vers.7. ἀνέγκλητον, id est, non a quoquam publici criminis accusatum, sed inculpatum. Ratio est, quod multum de dignitate et authoritate Episcopi et pastoris detrahit publica non tantùm sceleris damnatio:sed etiam accusatio, quia in ea perpetrati criminis suspicio quaedam inest, quanquam non est conuictus. Et qui caeteros reprehensurus est, debet ne sordescat eius admonitio, omnibus vitiis carere. Debet enim omni vitio carere, qui in alterum dicere paratus est, vt ait M. Tull. Istud vero Pauli dictum, vt sit irreprehensus Episcopus, de publica infamia et de eo scelere, quod accusatione publica obnoxium, aut dignum est, intelligitur, quo carere et illius accusatione, debet futurus pastor: non autem de peccato vel infirmitate qualibet accipitur, cui etiam optimi quique sunt obnoxii, vt ea in futuro Episcopo nulla inesse omnino possit. Ἀνέγκλητοι quidem esse possumus, Αυναμαότητοι vero esse non possumus Peccatorum igitur, vt docet etiam saepe et apposite quidem ad huius loci explicationem Augugust quaedam sunt Crimina, et Peccata, vel infir\pend
\section*{CAPVT  III }
\marginpar{[ p.103 ]}\pstart mitates tantùm humanae naturae, et vitiositatis humanae appendices et communes omnibus hominibus propter inhaerentem nobis peccati originalis radicem. Peccata in futuro Episcopo inesse possunt, crimina in eo nulla inesse vel tolerari debent. Itaque si quis sit homicida, adulter, falsus testis, latro, pirata, vel seditiosus, Episcopus ne eligitor:electus per ignorantiam deponitor, vt est in Agathensi concilio cap. 50. et in canon. Si Episcopus distinct. 50. Idemque Ambrouus docet Eplitola ou.  \pend\pstart Hic autem variae quaestiones incidunt et oriumtur, quarum quasdam tantùm, quae sunt vtiliores, breuiter tractabimus.  \pend\pstart Prima est, vtrum hac eadem voce irreprehensus comprehendantur peccata quae scelera quidem sunt, publica poena tamen minime puniuntur, qualis sub Paganismo idololatria: et in Papatu, sacrificium missae, vt qui haec admiserit, id est, qui idololatra fuerit, qui Missam blasphemam cantauerit eligi non possit in Episcopum. Nec enim ista peccata publice in Papatu accusari possunt, nedum damnari. Respond. Distinguendum esse. Aut enim de eo quaerimus, qui ignorans talis fuit: aut qui fidem Christi semel susceptam abnegauit. Ignorans, et postea ad fidem conuersus potest exemplo Pauli, fieri Episcopus. Non potest autem is qui sciens et abnegata Christi fide talis fuit. Nec enim tam spectat Paulus in criminibus et vera infamia definienda, quid leges hominum statuerint, quam quid Dei verbum damnet. Ergo Christianus qui defecit a fide, qui apostata fuit, qui Christum abnegauit publice, qui Christum  \pend
\section*{404 AD I. TAVL. AATTIM. }\pstart professus, idolis postea sacrificauerit, eligi certe non potest in Episcopum. August. epist.5o. et est in can. Presbyteros et canon. Hi qui post, dist.50. Quod enim Mar cel. Episcopus Roma. post fidei abnegationem retentus est in Episcopatu, propter magnam eius poenitentiam factum est, quam demonstrauit. Qui igitur nondum Christum agnoscentes et professi idolis seruierunt, et postea ad fidem çonuertuntur: non impediuntur propterea, si caetera habeant, in pastores Ecclesiae eligi.  \pend\pstart Secundo loco quaesitum est, Qui accusati tantùm quidem sunt criminis publici, non autem dannati, vtrum eligendi omnino postea nunquam sint. Inhaeret enim quaedam labes in istis, et opinio turpis apud bonos et graues viros. Respon. Hoc totum ex facto pendere. Nam si verisimilis nota et suspicio infamiae aut malae turpisque vitae iis inhaeret, qui accusati sunt et non conuicti, eligi non debent, quanquam prorsus non sint conuicti, quia Episcopum non tantum carere ipso vitio et crimine apud Deum oportet: sed etiam apud homines criminis suspicione, vt Ecclesia melius aedificetur. Sin prorsus innocens est pronuntiatus, maiórque illius integritas inde apparuerit, et eluxerit, eligi et admitti potest, etiam is, qui publice criminis accusatus est.  \pend\pstart Quaesitum denique est, vtrum Episcopi vel Presbyteri, vel Diaconi, qui aperte in huiusmodi scelera inciderint, fuerintque propterea depositi, iterum eligi ad eosdem vel maióres vel inferiores gradus Ecclesiasticos possint, si appareat  \pend
\section*{CAPVT  III. }
\marginpar{[ p.105 ]}\pstart eos vere poenitere, et vtiles Ecclesiae esse posse. Respond. in eo variam fuisse veteris Ecclesiae sententiam. Primùm enim fuit seuerior Ecclesia: et erat moribus receptum, vt lapsi siue in haeresin, siue in apostasiam, siue in aliud aliquod crimen, et vitium, et depositi propterea, non admitterentur vnquam postea in gradum vllum Ecclesiasticum, propter aedificationem Ecclesiae. Deinde aliud receptum et constitutum fuit, nempe vt etiam recipi possent, si poenituissent. Quam secundam sententiam cùm Luciferiani damnarent et improbarent, haeretici sunt habiti. Denique tandem res tota in concilio Agathensi composita et conseieutu eneyte sapiimmrumce restitui et retineri in gradu possent: ad maiores tamen non possent prouehi. Excepti tamen sunt ex eo numero apostatae. canon. Contumaces et cano. Domino sancto distinct. 50.  \pend\pstart Vnius vxoris virum. Secunda Episcopi futuri conditio. Primùm autem huius loci explicatio vera videnda est. Deinde ratio praecepti. Denique  quae huius loci postea corruptela siue deprauatio facta fuerit.  \pend\pstart Ac quidem primùm hic locus a patribus vere et pure et syncere explicatus est, veluti a Iustino et hic a Chrysostom. Atque etiam a Hieronymo, quanquam infestissimo nuptiarum hoste:item ab Ambrosio, epist. 82. et Theodorito (etsi postea Graeca Ecclesia cum Latina de huius loci explicatione disputauit et dissensit.) Ne is nimirum in Episcopum eligeretur, qui vno eodemque tempore duas apud se haberet vxores. Tales enim tunc multi etiam Iudaei erant, qui exemplo Daui\pend
\section*{AD I. PAVL. AD TIM. }
\marginpar{[ p.106 ]}\pstart dis Salomonis et aliorum Patriarcharum abusi duas pluresque simul coniuges ducebant. Nunquam autem illi acceperunt ita hunc locum, quasi qui coniux sit, Episcopus fieri non posset, aut quasi qui iam factus esset Episcopus, vxorem ducere, si velit, non deberet aut auderet. Nondum enim impia Cyricii Rom. Episcopi sententia edita fuerat, quae Episcopos celibes effecit. Ergo apparet esse posse Episcopos coniuges, sed tamen non posse esse digamos. Digamum enim et quidem vere interpretabantur, qui duas simul vxores haberet, et eodem tempore. Sed et duas quoque vxores habere censendi sunt, qui primam quidem vxorem repudiarunt: sed extra causas verbo Dei comprehensas, et aliam illi repudiatae superinduxerunt, qualia multa ex non iustis causis fiebant olim diuortia tum inter Iudaeos ipsos, tum etiam iter Gentesirvee enn duiusmodi coniugia abfunt ab impuris adulteriis, vt docet Christus Matt 19. vers.9. quia extra causas a Deo explicatas fiunt. Ergo et hi quoque sunt digami, vti et priores illi, qui duas vxores habent, quos Episcopos esse non vult Paulus.  \pend\pstart Ratio vero huius praecepti est multiplex. Primùm aedificatio Ecclesiae. Cùm enim multae coniuges eiusdem viri, sint profanatio matrimonii a Deo instituti: et magnae libidinis et incontinentiae signum, resque ea sit ómnino per se foeda, et turpis, perinde atque si plura corpora vnum tantùm caput haberent, ab ea foeditate maxime pastori Ecclesiae est fugiendum et abstinendum. Denique  cùm Deus ipse grauissime huiusmodi coniugiis, quae sunt adulterio paria, offendatur, illa prohi\pend
\section*{CAPVT  III. }
\marginpar{[ p.107 ]}\pstart beri debuerunt, maxime futuris pastoribus. Cur autem diserte istud commemorare oportuerit ratio est, tum ludaeorum receptus et vulgaris iam mos, vt docet Chrysostom. qui sibi plures simul vxores sumebant: tum communis omnium pene populorum, imprimis autem Graecorum et Orientalium opinio, barbara quidem, sed tamen iam confirmata: qui id minime turpe putabant, maxime si priori vxori repudii libellus etiam sine causa missus esset. Nam et Diogenes Laert. scribit in Socrate id praeceptum fuisse et edictum Athenis, vt duas vxores vir quisque haberet: et Socrat. liber 4. cap.32. tradit eandem legem etiam a Valentinia.3. imperatore Romano repetitam fuisse et sancitam, tam leue peccatum habebatur Polygamia.  \pend\pstart Epiphan.in liber  Haeres. de Cataphrygibus  agens putat ad Christi summi Sacerdotis imitationem, cuius imaginem gerunt Ecclesiae doctores, istud a Paulo praecipi. Nam quemadmodum habet Christus vnam tantùm Ecclesiam:sic Episcopus ait vnam tantùm vxorem habere debet. Augustinus liber  de Bono coniugali cap.18. vnitatem Gentium nomine vnius vxoris significari ait, Inepte certe, pace tantorum virorum dictum sit, vtrique. Ergo proculdubio Episcopos, Presbyteros, et Diaconos coniuges habuit vetus Ecclesia et proxima Apostolorum temporibus aetas, vt apparet ex Historia Ecclesiastica: imprimisque Eusebus liber  6. cap.42. Socrat. Scholast.lib5.cap. 22. Sed et ex ipsis quoque Papistarum canonibus et decreto Gratiani. Sed et August. cùm scribit ad Paulinum Nolae Episcopum Tarracasiae quoque ipsius vxoris me\pend
\section*{108 AD I. PAVL. AD TIM }\pstart minit, nec eum quod coniux esset, indignum iudicat Episcopatu.  \pend\pstart Secuta est autem alia aetas, quae propter coelibatus et virginitatis peruersam admirationem, hunc locum primum a vero sensu deflexit et detorsit. Deinde prorsus corrupit. Ac detorserunt misere hunc locum ii, qui ita sunt interpretati, eos esse eligendos Episcopos, qui fuerant quidem ante electionem viri vnius vsoris: non autem, qui coniuges tunc erant, cùm eligerentur. Sic Ambronus accipit eplit. 62. Sed ii fanuntur pessime. Nam et hic et ad Titum 1. vers. 6. diserte dicit Paulus, si quis est tempore praesenti:non autem, si quis fuit, tempore praeterito.  \pend\pstart Corrumpunt plane hunc locum alii et ridicu le explicant, qui vxorem hic Ecclesiam appellant et a Paulo significari voce vxoris putant. Vnius ergo vxoris virum esse, est vnius Ecclesiae et Sacerdotii Episcopum et administrum esse. At tunc id est, tempore Pauli nondum erant Sacerdotia, et quaestus Ecclesiastici. Vt mera sit cauillatio et peruersio Scripturae, sic mentem Pauli explicare. Ergo sic varie corruptus est hic locus. Quod autem ait Paulus, vxoris, ita fuit stricte acceptum a Patribus, vt praeter mentem Pauli quidam laqueum postea induxerint conscientiis hominum.  \pend\pstart Vxorem enim eam tantùm esse volunt, quae antea honestae vitae mulier fuerit. Itaque si quis minushonestam mulierem, si quis scenicam mulierem, quae spectaculis inseruierit. Item aliam quampiam ex turpi causa olim notam duxerit, excludi ab omni Ecclesiastico ministerio debere putant, et Synodorum decretis sanciuerunt. Est in cano\pend
\section*{CAPVT  III. }
\marginpar{[ p.109 ]}\pstart in canonibus Apostolorum, qui dicuntur; estque relatum in canon. Si quis viduam et eiectam dist. 34. et a Neocesariensi concilio (quod Nicenum primum praecessit,) repetitum cap. 8. Quod etsi non improbo ab iis ita esse constitutum, ex eo tamen laqueum conscientiis vllum induci nolim, et ita stricte in omnibus obseruari. Nam hoc ita constitutum solùm est a Patribus illis, quod Ecclesiae aedificatio in tanta meretricum copia, et inter gentes profanas versantes exigere huiusmodi quiddam videretur. Oseas tamen duxit meretricem in vxorem, si quibusdam credimus. Sed et siquis eam vxorem habet, quam ex legum praescripto ducere non potest, veluti si intra eum gradum, qui legibus alicuius prouinciae politicis prohibetur, ad munus Ecclesiasticum vocari et admitti non debebit ex Synodo 6. generali, quae Ioit Gontaneinoponta in Piuno, et est in cano. Presbyterum non legalibus distinct. 28. vbi Gratianus se imperitum prorsus prodit: vel si per ignorantiam admissus est ad Episcopatum, interdici muneris sui exercitio debebit, ex iisdem Synodi decretis  \pend\pstart Additum quoque est a concilio Chalcedonensi cap. 14. et sancitum, ne qui vxorem haereticam, aut Paganam siue Gentilem et idololatram duxisset fieret Episcopus et clericus, ne fortasse illa virum a fide Christi abducat, aut pietatem Christianam profanet, et calumnietur, quia hostis est et inimica. Sic est in cano. Quoniam in quibusdam distinct.32. Sed et qui duas successiue et legitime ductas vxores habuerunt, Digami sunt appellati ab istis: sed perperam. Itaque ab officio pastorum  \pend
\section*{AD I. PAVL. AD TIM. }
\marginpar{[ p.110 ]}\pstart prohibiti sunt ex mala interpretatione huius loci. Nam nec secundae nuptiae, nec tertiae peccatum habent, si caste seruent Augustinus in liber  De fide ad Petr. quanquam Bernar sermo. 58 Cant. damnat exemplo turturis numerositatem nuperaruiloudumomumeu laperstitione, vti eam quoque Tertullianus damnauit ex Montanistarum errore.  \pend\pstart Denique vxorem vnam dixit Paulus. Quid si quis vxori vni quidem legitime ductae soperinducit concubinam, Vtrum is pastor et Episcopus eligi potest? Est enim haec concubina et pellex, non vxor. Respond. Tergiuersationem esse. Nam plana est Pauli sententia, per quam omnem in eo qui Episcopus est futurus, turpis et obscoenae libidinis notam damnat, ne propterea male audiat Dei verbum, aut ipse sit offendiculo cuiquam: Itaque qui pellicem cum vxore habet in consortio, est adulter, et prohibetur eligi et hic et can: Apostolorum 17. Aurelia. concilio cap. 4. Agathensi cap.55. et est in canon. I. distinct.33. et tota distinct.34.  \pend\pstart Si excipiatur separatam esse Euangelii doctrinam a vitiis Ministrorum, qui eam annuntiant, nec eam, quae per se pura est, illorum sordibus pollui posse. Respond. Per se quidem illud verum esse: opinione tamen hominum non item. Putant enim homines Ministrorum vitiis tanquam attrectatione quadam doctrinam illam coelestem pollui. Itaque illa minus et lubenter et reuerenter iam auditurs cùm huiusmodi quaedam personae foedae et turpes in diuinum ministerium asciscuntur, et alleguntur.  \pend
\section*{CAPVT  III. }
\marginpar{[ p.111 ]}\pstart Venio iam ad eas postremae aetatis constitutiones, per quas hoc Pauli dictum et praeceptum non tantum inflexum est a vero sensu: sed omnino immutatum, et sublatum ex nimia et praepostera celibatus admiratione.  \pend\pstart Ac primum quaesitum est anxie, Vtrum celibatus esset Ministris Ecclesiasticis, siue iam vocatis et ordinatis, siue vocandis et creandis praecise indicendus et praefiniendus. Respond. Cùm neque donum continentiae sit in nostra potestate, neque  semper, et vel in perpetuum a Deo detur, si cui detur: sed saepe tantùm ad tempus, perperam, et inique eos agere, qui illum praescribunt Episcopis et vllis Ministris Ecclesiasticis, vel iam electis, vel eligendis. Onerant enim hominum conscientias praeter Dei verbum, quod facere non licet, iisque leges imponunt, quod solius est Dei. Deinde Cataphrygum et Encratitarum haeresin, qui nuptias damnarunt, et reuocant et confirmant. Etsi enim nuptias non in totum tollere videntur: eatenus tamen illis haereticis aslentiuntur, quatenus sanctam illam Dei institutionem (quae est coniugium) putant polluere quosdam pios: aut saltem obstare, quo minus illi ad Dei obsequium sint apti et idonei, cùm tamen iis nihil eorum, quae Dei verbum in pastoribus requirit, pro captu humanae infirmitatis, desit.  \pend\pstart Quare cùm Synodo Nicena prima de ea lege clericis imponenda ageretur et disceptaretur, obstitit vir sanctus Paphnutius Episcopus et Christi martyr et ipse ἀγαμος tamen, qui ex Dei verbo omni hominum generi honestum esse coniugium docuit ex Hebus  13. vers 4. Atque haec Hi\pend
\section*{AD I. PAVL. AD TIM. }
\marginpar{[ p.112 ]}\pstart storia extat apud Sozom. liber  1. cap.23 quicquid Torrensis contra debacchetur. Nam Theodorus Beza in liber  de Polygamia ipsa Synodi Nicenae verba Graeca recitauit ex veteri codice Graeco ex cerpta, et illius ipsius rei fit mentio in ipso Dedereto cano Nicena distinct.31. vt mirum sit fuisse tam impudentes quosdam vt id negent. Et nostro seculo ad fratrum Boemorum interrogationes de ea re (obseruant enim illi in suis Ministris coelibatum) quod esset suum iudicium scripsit responditque D. Cal. epist. 339.  \pend\pstart Cogi quidem vxores ducere non debent ex hoc Pauli loco Episcopi, sed neque prohiberi ponunt. Pibe emm cotuin eit plane liberum et ἐδίάφορον. Vnde necessitas aut lex in neutram partem statui debet et induci. Nam etiam Eunuchus Episcopus eligi potest: Eusebius liber  7. cap32.qui certe vxorem ducere non potest.  \pend\pstart Quibus vero gradibus ad hanc tandem Coniugii prohibitionem, vel potius profanationem in Ecclesia ventum sit, iam explicemus.  \pend\pstart Primùm igitur coniugatos Ministros aeque et promiscue, atque coelibes habuit vetus Dei Ecclesia, et bene instituta. Id quod durauit diutissime, quando etiam Gregorius Nyssenus, vt ait Nicephorus, fuit vxoratus.  \pend\pstart Postea vero aliud, tum in Ordinandis in ministerium, tum in Ordinatis receptum est. Ac Ordinati iam Episcopi, qui erant coniuges, vxores suas retinebant, iisque vti poterant, neque id etiam obstabat eorum ministerio exercendo et muneri Socrat. liber 5. cap. 22. cano. de Syracusanae 5 Siue ergo distinct.28. cano. Apostolorum 5 et ita Gan\pend
\section*{CAPVT . III. }
\marginpar{[ p.113 ]}\pstart grenn etiannconemos quou luit lub Syluestro. Roma. Episcopo constitutum fuit.  \pend\pstart Ordinandi vero primùm coeperunt in ipsa electione interpellari, vtrum continentes esse vellent, et in perpetuum abstinere coniugio, an non. Qui si promisissent se continentes futuros: vel etiam tacuissent interrogati, post electionem ad ministerium Ecclesiasticum vxorem ducere, iam non poterant. cano. Diaconus distinct.27. Hic enim vel silentium ipsum pro voto perpetuae continentiae habebatur. Itaque  si postea matrimonium inibant, ordine suo mouebantur. Neque tamen vt scortatores; vel adulteri puniebantur, neque Sacramentorum vsu, et participatione interdicebantur, vt ex Neocaesariensi consilio apparet. Quod tamen decretum, et quam tam seueram Ecclesiae disciplinam neque August. neque etiam Hieronym. probarunt cano. Quidam distinct.27. Sin autem Ordinandi de promittenda perpetua ἀγαμίς interpellati vel omnino recusassent vouere, et polliceri continentiam, vel ad tempus tantùm promisissent, neque eligi prohibebantur, neque post electionem ad Ecclesiasticum munus vetabantur vxores ducere, ex Ancyranae Synodi cap.1o. Differebatur tamen saepe electio, et spectabant illo tempore, num quem allum inuenire possent, idoneum tamen, qui continentem et abstinentem a coniugio vitam polliceretur, vt eum potius eligerent: est in cano. De Syracusanae distinct. 28. Quod si neminem continentiam vouentem idoneum pastotali muneri reperiebant, hunc, qui continentiam nollebat promittere, eligebant, idque fas erat.  \pend
\section*{AD I. PAVL. AD TIM. }
\marginpar{[ p.114 ]}\pstart Et quidem in iis qui ad maiora munera et officia Ecclesiae vocabantur, obseruabatur ista stipulatio, et interrogatio, veluti in iis, qui ad Episcopatum, Presbyteratum, et Diaconatum adoptandi erant. In subdiaconis vero acolutis, lectoribus et aliis inferioribus, et minoribus ordinibus idem neque fiebat, neque habebat locum, quia poterant etiam el ecti vxorem ducere, neque de vouenda ἀγαμία tales eligendi interpellabantur cano. Si quis vero et can. Seriatim distinct.32 cano. Lex contmentiae diitiniet.51.Quanquam postea id immutatum est in subdiaconis, quibus fuit interdictum ne vxores ducerent: vt scribit BernaAbbas Clareua. epistola 203.  \pend\pstart Ratio vero, cur istae stipulationes et antestationes in ordinandis fierent, duplex afferri solet. Prima quidem paupertas Ecclesiarum, quae tot familiis alendis, Episcoporum nimirum, Presbyterorum, et Diaconorum satis esse non poterant, in tam tenuibus adhuc opibus. Nam Ecclesiae tunc vel nulla erant aeraria, vel exigua admodum, quia quae bona habebant Ecclesiae pauperibus erogabantur, vt plurimum. Etsi vero bona Ecclesiae semper distinguebantur a bonis Episcopi, tamem illius aetatis homines metuebant, ne coniugati Ministri bona Ecclesiae ad vxores, aut affines, aut filios suos clam deriuarent, et transportarent. Est in cano. De Syracusanae. distinct. 28. Sed non debuit diffidentia de Dei erga nos beneuolentia, et erga suos prouidentia inducere hos metus, et aspergere istam labem Christianae Ecclesiae. Si enim Dominus prospicit auium necessitati, et dat illis cibum et escam, multo magis suis prospiciet,  \pend
\section*{CAPVT  III. }
\marginpar{[ p.115 ]}\pstart et imprimis verbi sui fidelibus Ministris Matth.G vers.25. etsi Ecclesia iis non subuenerit. Deinde nonneiis praecipue, vti Christianis pauperibus, exaerario pauperum erogari debuit?  \pend\pstart Secunda vero istius postulationis ab ordinandis ratio fuit, quod continentes Ministros sanctiores, quam coniugatos esse existimarent et mundiores, quia Dominus iussisset suum populum aliquando vxoribus abstinere, vt se sanctificaret Exod. 19. Denique propter dictum Pauli 1. Cor. 7. vers.5. liberiores et paratiores esse putabant: vt orationi et suo muneri vacarent. Quae tamen opinio ex peruersa coelibatus admiratione et laude profecta est, et ex Encratitarum haereticorum coeno deducta, in vituperium et contumeliam coniugii can. Ante triennium distinct.31.  \pend\pstart Atque hie mos in Occidente potiùs et strictius siue religiosius, quam in Oriente primum obseruabatur. Imprimis autem Romae obtinuerat, quem tamen aperte improbat Gregor. Magnus liber 1. epist. 42. et apparet ex cano. Placuit Episcopos distinct. 32. cano. Ante triennium distinct. 31. quam hoc aegre fuerit receptum. Nam Graeci Episcopi et Orientales vel raro pollicitationem istam de continentia exigebant: vel eam etiam non obseruabant, vt docet Socrat.libii: caprril Ioc tantùm constituerunt, vt ne liceat ordinato Episcopo accipere aliam, priori vxore defuncta et mortua cano. De Syracusanae in fi. distinct. 28.cano. Quoniam distinct. 31.  \pend\pstart Denique tandem post varias dispurationes et laudes virginitatis et coelibatus eo res rediit, vt mariti et coniuges viri ad aliquam dignitatem  \pend
\section*{AD I. PAVL. AD TIM. }
\marginpar{[ p.116 ]}\pstart Ecclesiasticam non possent assumi Arelat. concilio 4. cap.2. can. Assumi distinct. 28. cui subscripsit Callixtus Papa Rom. At vero Huldericus Augustanus Episcopus, qui vixit ante annos 700. acerrime ei decreto restitit et Nicolao Papae. Quod nouum etiam decretum et inauditum appellat can. De Syracusanae et prorsus improbatur Gangrensis Synodi cap.4. Synodo Constantinopolitana, quae fuit sexta vniuersalis Synodus, quae vtraque Synodus multis seculis et Arelatense illud concilium 4. et Callixtum Pontificem praecessit. Sed et ante vtranque Nicenla quoque prima cap.13. et ante hanc Nicenam cano. Apostolico, qui est N.6. qui etiam est in cano. Si quis distinct. 28. adeo vt Leo imperator cognomento Philosophus Nouella sua 2. iubeat hoc decretum ex Dei Ecclesia facessere, et prorsus in personarum Ecclesiasticarum electionibus silere. Denique Synesius Episcopus ad Eutropium negat suam se vxorem dimissurum, quam ei et Deus ipse, et lex habere concedunt. Cùm autem ita varie de coniugio Sacerdotum disceptaretur, primus Syricius Papa, qui vixit circa annum Domini 400 et Ambrosii Mediolanensis Episcopi aetate, magna temeritate et nullo Sacrae scripturae nixus lufidamente definiuit et sanxit, Sacerdotes manerent coelibes. cui videtur Ambros. epist. 8o. et 81. assensus. Hispania vero vniuersa, et Germania contradixit, adeo vt haec etiam post longum postea tempus Episcopos et Presbyteros coniuges haberet. Syricii dictum est in distin. 8.Anno Domini 1074. vt ait Nauclerus contra Hildebrandum id est, Gregorium 7. Papam Germani de eo cer\pend
\section*{CAPVT  III. }
\marginpar{[ p.117 ]}\pstart tarunt grauissime. Et haec quidem in Ordinandis ad munus Ecclesiasticum obtinebant. In iam Ordinatis vero istud obseruatum primum est, vt vxores suas, quas tunc habebant, retinerent, iisque vterentur: nisi eas contra spontaneam promissionem iam factam duxissent, tunc.n. ab officio deponebantur, quia vxores duxerant. Tamem rata ista coniugia esse et haberi oportere iudicabatur: ne tales a fuis vxoribus diuellerentur. Sed inter digamos tamen poenitebant. Postea vero iussi sunt, etiam qui coniuges et mariti electi fuerant, vxorum suarum vsu abstinere, quod fuit Abelonitarum haereticorum dogma vt docet in liber  de Haeres. August. etsi retinere et eas, tamen sine vsu, domi habere poterant et concedebantur. Itaque fere in monasteria, quod ista vxoris sine vsu retentio mali esset exempli, or dinatorum ministrorum vxores detrudebantur Tandem vero crescente in Ecclesia superstitione et instante Syricio Papa impurissimo, et eius successoribus Papis, coacti sunt, vel inuiti, Episcopi et Presbyteri eas deserere, et huiusmodi coniugia, velut incestae nuptiae, propter violatum continentiae votum sunt, omnino et aperte dannata, et dissoluta tanquam nulla essent can. Presbyteris et can.seq. dist.27. et omne consortium cum muli eribus clericis interdictum demum est, nisi cum cognatis et sibi sanguine coniunctis foeminis, cum quibus etiam nimium frequens consuetudo propter suspicionem fuit prohibita cano. Interdixit dist. 32, et cano. Quorundam distinct.34. Sed dum homines temerarii honestum matrimonium ministris Ecclesiasticis prohibuerunt, et interdixerunt,  \pend
\section*{AD I. PAVL. AD TIM. }
\marginpar{[ p.118 ]}\pstart grauissima pellexerunt et inuitarunt vt ait etiam Berna. ipse sermo. de conuersis cap.29. quibus nihil est inter huiusmodi coelibes clericos frequentius. Saluianus Massil. Episcopus qui vetustissimus est, sic improbat istam coniugii continentiam et nondum tamen prohibitionem liber  5. de Prouidentia Dei. Licita non faciunt, illicita committunt, temperant a concubstu, non temperant a rapina. Quid agis stulta perfuasio? Peccata interdixit Deus, non coniugia. Post huiusmodi igitur constitutiones receptas tantum extitit coniugii inter clericos odium, vt pro haereticis et Nicolaitis vulgo damnarentur, qui coniugium iis licere docerent. Itaque Arnoldus Brixianus, et ipse Mronaenusjae tate Dernat ur damatus est, ob eam causam, sed iniustissime. Longior historia de hac tota re est apud Balaeum et Centuriatores. Quod fuisse inuentum Heliodori tradit Socrat. libus 5.cap.22.  \pend\pstart Addit Paulus vηφάλαιον quem alii sobrium, alii vigilantem vertunt: sed vox sequens σώφρονα sobrium significat. Nos hoc loco vigilantem interpretabimur. Hoc ipsum est quod ait idem Paul. in Epistola ad Titum 1.vers. 8 ἐγκρατῆ, cui virtuti ἐσωτία contraria est, quam etiam reprehendit in Episcopo praesertim et pastore, totaque ipsius familia Paulus Tit.1.vers.6. Hoc quidem praecipit Christianis omnibus Christus, vt sint vigilantes Matth.24. vers.42, et 25. vers.13.tamen a pastoribus et iis qui aliis praesunt vigilantia multo aequius et seuerius requiritur propter eorum munus. Est n.haec vigilantia pars illius fidelitatis et curae, quam ab illis exigit scriptura 1. Corinth.4.  \pend
\section*{CAPVT  lll }
\marginpar{[ p.119 ]}\pstart vers.2. et 1.Pet. 5vers.2. Et quia huius vigilantiae comes est sobrietas, vti negligentiae comes est ebrietas, vt docet Christus Luc. 21. v.34. idcirco ab eodem Episcopo postea requirit Paulus, ne sit vinosus, vt postea dicetur. Huic autem virtuti duo quaedam vitia opponuntur, Negligentia, et Absentia. Ac negligens sui muneris pastor esses non debet, nec illud per contemptum spernere: Paulus 1.Corinth.9. vers.16. Vae mihi nisi Euangelizauero, ad hoc enim missus sum. Sed neque absens esse debet diutius a suo grege. Quomodo enim vigilare potest, qui absens est? Vnde praesentiam et residentiam pastorum dicit esse prorsus necessariam Paulus, nedum vt ab ea Ecclesia cui pastores dati sunt, abesse perpetuo possint, quod faciunt praelati Papistici Philip.2.vers.25. Actor.2o. vers.28. vt dispensatio certe de non residendo: vel de pluralitate beneficiorum, quam vocant, omnino contra ipsam Pastoralis muneris naturam et deffinitionem instituta apparet, verissimaque subuersio est officii et muneris Episcopalis et pestis. Quae nullo modo danda pastoribus est, nec toleranda in Ecclesia Dei. Conquesti sunt Patres in Synodo Ariminensi collecti, quod diutius quam vtilitas Ecclesiarum postularet, a suis sedibus praetextu Synodi abessent, vt tradit Socrat. liber 2. cap. 37. Bernardus sermo.77. in Cant. Propterea attendite vobis quicunque opus ministerii huius sortiti estis. Ciuitas est Ecclesia vobis commissa, vigilate ad custodiam concordianque, sponsa est, studete ornatui. Ques sunt, intendite pastui. Porro custodia ciuitatis, vt sit sufficiens rifaria erit, a vi tyrannorum, a fraude haeretico\pend
\section*{AD I. PAVL. AD TIM. }
\marginpar{[ p.Ro ]}\pstart rum, a tentationibus Daemonum. Sponsae vero ornatus in bonis operibus, et moribus, et ordinibus. At pastus omnium communiter quidem in pascuis seripturarum tanquam in haereditate Domini etc.  \pend\pstart Huius autem negligentiae et absentiae castigatio non tantùm ad Synodum, aut consistorium Ecclesiasticum pertinet, sed etiam ad Magistratum pium et fidelem, vt docemur exemplo Constantini Magni, et refert Sozomenus liber  2. cap.4.  \pend\pstart Ex huius autem praecepti occasione, et recte quidem, factum est, vt pastores et Episcopi Venatores et Aleatores, non tantùm damnarentur: sed eriam deponerentur, quia iis rebus et studiis occupati, et rebus Ecclesiae attendere et moribus hominum inuigilare non possunt. Denique nimiae dissolutionis vitae huiusmodi studia et occupationes sunt certissima testimonia, quanquam Germanis et Gallis Episcopis semper nimium istae res assiduae et familiares fuerunt vt est in can. Quorundam et seq. distinct. 34. cano. Episeopus distinct.35.  \pend\pstart Subiicit Paulus σώφρονα, Alii prudentem vertunt, alii sobrium (Male prudentem, est enimoρóνιμος prudens, non σ́φρων) Huic autem temperan tiae coniungit Paulus iustitiam. Vult enim esse temperantem, et iustum Tit.1.vers. 8. Temperantem in seipso, ne abripiatur vllo vitioso affectu ad turpia perpetranda. Iustum vero ad alios, ne quemquam offendat, aut re aut verbo, aut manu, aut lingua vel rapiendo, vel fraudando. Quae hic disputant Canonistae de prudentia Episcopi, qualis esse debeat, nos in explicatione vocis Nεσφύτα dicemus.  \pend
\section*{CAPVT  III. }
\marginpar{[ p.121 ]}\pstart Κέσμιον, id est, Compositum, alii, vt Canonistae Ornatum, mens Pauli est, vt omnis vitae ratio ab Episcopo recte sit, sancteque instituta et conposita, vt neque leuitas in eo appareat: neque mo rositas quaedam. Itaque grauem Episcopum esse oportet, non quidem, vt non se affabilem et cuilibet accessibilem praebeat (est enim ista toruitas potius et inhumanitas, quam grauitas) sed ne se contemptibilem cuiquam ostendat et praebeat. Vnde recte ii explicant, qui concludunt ex hoc loco. Neque scurram, ne que adulatorem, neque leuem esse debere Episcopum Ephes.4. vers.5. Bernard.sermo 46. Cantic. cant. Cypressus item boni aeque odoris et imputribile lignum incorruptae vitae, et fidei, etiam quemuis de clero debere esse demonstrat.  \pend\pstart Papistae tamen et Canonistae quaedam argumento huius vocis afferunt, quae non videntur omnino praetermittenda, quaedam etiam quae sunt refutanda et praetermittenda. Ornatum igitur hunc, quem ab Episcopo requiri putant, faciunt duplicem, nimirum alium Externum, et alium Internum Externum ornatum collocant, in Cultu, et Incessu, quem vtrunque requirunt in Episcopo, honestum, et modestum, et recte quidem. Itaque quod ad cultum et vestitum priuatum Episcoporum, constitutum est, vt neque sordidus sit, neque  splendidus et fulgidus, sed sit illorum personae et ordini congruens, et decens et accommodatus, concilio Agathensi cap.2o. et Gangrensi cap.21. et est in can. Clerici qui dist.23.itemque 21. quaest.4. per totum. Quod ad incessum etiam est sancitum, ne sit mollis siue fractus: neue militaris can. Dar\pend
\section*{AD I. PAVL. AD TIM. }
\marginpar{[ p.422 ]}\pstart simoniam cum veste distinct. 41.  \pend\pstart Nominatim autem praescriptum est concilio Martini Papae, vt clerici vestes talares gestarent, et iis induerentur. Deinde ne barbam, et comam nutrirent, in quo manifesta concilii Chartaginensis deprauatio et falsificatio a Gratiano facta est, vbi vetantur clerici omnes barbam radere, pro quo iste supposuit, barbam nutrire. Ergo comam alere non debent clerici: sed attonso capite esse auribus patentibus, et desuper habere caput in modum Sphaerae abrasum, vt est in concilio Chartagin. 4. Additum est extra liber 3. Tit. de vita et honestate, clericor. ne discissas caligas, et indusia operose facta induant et gestent, quae fuerunt iam monachales superstitiones, et pars humanarum traditionum perieulola. Modus quidem in iis omnibus seruandus est, quem ratio locorum et personarum exigit et honestas ipsa postulat. Sed istae leges in nugis sanctitatem pastorum postea collocarunt:  \pend\pstart Quaesitum est etiam, Vtrum militari habitu et more vestiri possit et debeat Episcopus:aut habitum quem habet, vtrum in militarem vestem mutare possit, et Episcopus manere? Respondeo ex concilio Gangrensi cap. 21. Episcopo eo vestitus genere, nempe militati, non licere vti assidue et ordinarie. Ex causa tamen nonnunquam sumere posse, pro ratione locorum et temporum, vti docet Theodoret.libus  4.Histor. cap.13 exemplo Eusebii pii Episcopi militarem vestem assumentis. Coniungunt cum hoc loco disputationem, de Integritate membrorum, et de Nobilitate seu Diuitiis ordinandorum, quia haec etiam duo viden\pend
\section*{CAPVT  FII }
\marginpar{[ p.123 ]}\pstart tur quodammodo pars esse' externi ornatus, et species quaedam commendationis apud homines. Imprimis autem in persona, quae quotidie in publicum prodire debet, deformitas videtur esse odiosa et inconcinna. Vnde vulgo tritum prouerbium, forma digna imperio, et in Philopoemene (vt in eius vita scribit Plutarchus) deformitas vultus contemptum gignebat. Quamobrem quaesitum est, vtrum qui non est integro corpore et toto in Episcopum eligi possit et debeat. In quo distinguenda est, Deformitas Αʹκαλλία, a Mutilatione corporis, quam vocant, id est, a πειρώσει membrorum. Ac quidem de deformibus haec regula comrereuta cies vt detormes an Eecleliasticis dignitatibus non arceantur adipiscendis: nec adepti etiam deponantur. Deformitas autem eleganti formae opponitur. Mutilatio autem integritati corporis. Itaque ab ea differt.  \pend\pstart Ergo de mutilatis et corpore vitiatis et truncatis hominibus tum Ilerdensi concilio, tum canonibus Apostolor. tum responsis Pontificum, veluti Gregorii Maximi liber 2. epist. 25. statutum eity ve qur sespros serentes pruuemerque, et per impatientiam aliquo membro mutilarunt, ne dignitatem Ecclesiasticam adipiscantur: et adepti deponantur cano. Praecipimus distinct.34. et cano. Si quis distinct.55. Imprimis autem qui seipfos castrarunt, quia tales sunt homicidae sui cano.22. Apostolorum.  \pend\pstart Qui autem vi et renitentes mutilati corpore sunt aut castrati: vel ex iusta honestaque causa id passi sunt imse fieri, veluti valetudinis tuendae causa, et ex concilio medicorum, vt sanarentur a  \pend
\section*{AD I. PAVL. AD TIM. }
\marginpar{[ p.124 ]}\pstart morbo, vti saepe integra membra sint rescindenda, ne totum corpus putrescat.item propter nomen Christi, vt in perfecutionibus eruti sunt multis oculi, praecisae aures, praecisa membra, ii neque dignitatem Ecclesiasticam nancisci prohibeantur, neque adepti deponantur. Sic Paphnutius altero oculo vi afflictionum amisso mansit Episcopus, eóque laudabilior. Sic cano.21. Apostolo. et in concilio Martini siue Braccarensi cap 21.quo d refertur cano. Eunuchus et cap.seq. et cano. Lator distinct. 55. constitutum fuit.  \pend\pstart Alias, et his casibus et causis exceptis, sancitum est, ne mutilati vitiatique corpore in Episcopos vel clericos ordinentur, quales caeci, surdi, claudi, Eunuchi, exsecti, lippi etiam, et gibbosi cano. Hinc etenim distinct. 49. quia summus noster Sacerdos Dominus noster Iesus Christus, cuius, vt isti aiunt, imaginem Episcopi gerunt, fuit integer: et summus olim Sacerdos lege Dei vetabatur esse vel vitiatus corpore vel mutilatus Leuit. 21.Hebus 7.v 26. Respondeo vero iam totum hoc plane esse superstitiofum, nos ad illas veteres caeremonias reuocare: et dignitatem ministerii ex istis rebus metiri, non ex doctrina potius et morum sanctitate. Itaque vetus Ecclesia aliud obferuauit et probauit, et Eunuchos pastores habuit, vt et de Origene scriptum est et de Dorotheo Eusebus liber  7. Hist. cap.32. Ergo modo ne sit in eligendo vel electo ea infirmitas, quae muneri ipsius obeundo et praestando obstet, veluti caecitas et ἀλαλία, potest quis, et eligi et ordinari: nec electus deponendus est propter corporis mutilationem aut vitia membrorum.  \pend
\section*{CAPVT  III. }
\marginpar{[ p.125 ]}\pstart De nobilitate generis etiam quaesitum est, cùm in ea sit apud homines dignitas quaedam et commendatio magna, Vtrum nobiles in electione sint aliis non nobilibus praeferendi. Respondent August. epist.148. et Ambrosius liber 4. de paradiso, virtute vnusquisque gratiam sibi comparet: non loci, nomgeneris nobilitate, quanquam in Papatu aliter obtinuit, idque  pessimo exemplo. Nam apud istos generosi et nobiles aliis praeferuntur, etiam in Ecclesiastis muneribus adipiscendis, vbi minime omnium gratiae et fauori locus esse debebat.  \pend\pstart Pauperi vero diuitem praeferre in administranda Ecclesia, et munere Ecclesiastico deferendo noi taritam mutn en, leu cuum contra Pauli et Dei praeceptum, non multi nobiles estis, ait, non multi diuites 1.Corinth.1.vers.27. Nam quae sunt stulta et infirma in mundo, elegit Deus, vt mundi superbiam contundat.  \pend\pstart Sextam conditionem addit Paulus sit hospitalis, οιλόξινος, plus est, quam quod Latini dicunt esse hospitalem. Exigit enim promptum ad hanc ipsam liberalitatem exercendam et hilarem animum. Haec autem liberalitatis species omnibus Christianis praecipitur Hebus 13.vers.2. imprimis autem verbi Dei ministris, vt sint caeteris exemplo. Coniungit autem Paulus hanc cum superiore conditionem, vt sit ρηάγ ἀθος, id est, bonorum virorum amans Episcopus, quia similes fere iis sumus, quibuscum ver samur. Nam illam Chrysost. sententiam, quae est Homil 43. in Matth. 23. maxime probamus, vt sit omnis Sacerdos sanctus. Nec eo pertinet quod vtranque vocem coniunxit Paulus ριλόξνον, et φηάγεθον, vt excipiendorum  \pend
\section*{AD I. PAVL. AD TIM. }
\marginpar{[ p.126 ]}\pstart Αιδακπκὸν Septima conditio. Hoc autem ipsum quid sit Paulus explicat Tit.1.vers.9. nempe esse huiusmodi, qui possit, Sanam doctrinam apposite tradere, et contradicentes conuincere. Sunt autem tria contradicentium genera, Haeretici, Iudaei, et Pagani, imprimisque ipsi Philosophi. Plane haec virtus et conditio ad Episcopum pertinet, quoniam is debet esse doctor in Dei Ecclesia, et quidem vocalis: non autem canis mutus et idolum, quales Episcopi Papistici Isai.56.vers.11. non item solum praeceptor mutus et ἀφωνος, quales sunt qui scriptis tantùm erudire nos possunt, non autem viua voce, quod aut muti sint omnino natura vel lingua excisa: aut prorsus inepti ad docendum. Itaque mittit Christus discipulos vt praedicent, id est, quemadmodum explicat Paulus, vt denuntient et doceant, inf. cap.4.vers.11. Id quod etiam apparet ex dicto Petri in 1. cap.5.v. 2. Pascite, et ex Ezech.22.v.26, 830.cap.33. Quod  \pend\pstart hospitum delectum diligentissime faciat Episcopus, et solùm veros, cordatos, pios, bonósque Christianos excipiat: alios autem negligat, minimeque  excipiat, quia non tantùm omnibus est benefaciendum, etiam iis, qui sunt hypocritae: sed etiam mali et segnes homines saepe sunt beneficiis et huiusmodi liberalitate alliciendi, vel ad meliorem frugem reuocandi. Quaedam tamen conuiuia, quae fiebant olim in basilicis propter veteres agapas, quae fuerunt in vsu, fieri hodie ad hospites istos caeipleiidos, et benigne tractandos non probat August. epist. 1o9. Quarta pars redituum Ecclesiae olim excipiendis hospitibus pauperibus destinata et Episcopo commissa erat, vide distinct. 42.  \pend
\section*{CAPVT  III. }
\marginpar{[ p.17 ]}\pstart non nisi a loquente et vocali pastore fieri potest. Clamosos tamen concionatores, et Stentoreos, quosdam indoctos homines et inconditos, qui nihil nisi boant in templis et vaste vocifetantur, non probamus. Obtundunt enim tantùm auditorum aures, non aedificant neque docent.  \pend\pstart Iam vero ex hoc loco, et vere et apposite colligit in epist ad Titum Hieronym. Inertiam, desidiam, somnum et otium excuti Ministris verbi Dei et pastoribus Ecclesiae hac voce. Qui enim se his vitiis dedunt et segnitiei praesertim, non possunt ipsi esse διδακτποί. Oportet igitur sedulo et diligenter Episcopum et pastorem Ecclesiae euoluere scripturas, et earum lectioni intendere, quod etiam ipsi Timotheo, licet tanto viro, praescripsit Paulus inf. cap.4. vers.13. Ergo et sacrarum scripturarum libros versare, et legere eos oportet assidue, vt discant: atque etiam haereticorum scripta nosse, vt eos refotent, quemadmodum etiam diserte concil. Chartag. 4. cap.16. praeceptum est, et est in cano. Episcopus distinct. 37.  \pend\pstart Quaesitum est autem vtrum in artium, quas seculares vocant (qualis Poëtica, Historia, Dialectica, Physica, linguarum variarum cognitio et interpretatio) cognitione versatum et peritum esse Episcopum, aut pastorem Ecclesiae oporteat. Respondeo, Quidam ne has quidem Episcopis, aut eorum filiis (excepta et linguarum trium nenpe Hebraicae, Graecae Latinae cognitione) attingendas putant, quia corrumpunt bonos mores colloquia praua, et per Philosophiam nos a veritate abduci et saepe decipi docet Paulus Coloss2.vers. 8.Quae sententia in artes quidem ipsas, que  \pend
\section*{AD I. PAVL. AD TIM. }
\marginpar{[ p.128 ]}\pstart sunt Dei dona, certe est iniuriosa Quamobrem sequenda potius est illa, et Basilii, et Hieronymi ad Damas. De filio prodigo scribentis sententia, qui omnibus politioribus actibus et humanioribus instructum, si potest fieri, Episcopum vult esse. Ac ne seculares quidem literas, qualis est Poëtica, ignorare. Multa enim illae artes, ad docendum veram Theologiam, adiumenta conferunt imprimis autem Dialectica, et Physica, et sunt animi nostri ornamenta. Implendus sibi quisque bonis est artibus, illae sunt animi fruges, vt ait ille in Ætna. Paulum ipsum summum fuisse Dialecticum, ex illius scriptis apparet: et Augustinum in omni bonae literaturae genere versatum fuisse: sed maxime in Dialecticis intelligitur, ex illius opusculis variis, et libris contra Donatistas, et alios haereticos editis. Sed tamen in earum artium studiis, vel pertractandis, vel persequendis, modus adhibeatur. Placet enim hic illa Neoptolemi apud Ennium sententia, Philosophandum quidem esse, sed paucis: nam omnino haud placet. Nota est vero historia de Hieronymo ab angelo in somnis flagellato, quod Ciceronianorum librorum lectioni nimium esset addictus, vt ipse ad Eustoch.scribit de Custod. virginitat. Et hic vsurpemus illud, Bibi et fugi. Bernard. sermo.36. in Cantica cant. Videar fortasse in sugillatione scientiae nimius, et reprehendere doctos, ac prohibere studia literarum, absit. Non ignoro quantum Ecclesiae profuerint, et prosint literati siue ad refellendos eos, qui ex aduerso sunt: siue ad simplices instruendos. Ergo, vt recte ex variis Patrum seriptis deffinitum est, et est distinct.37. degosten\pend
\section*{CAPVT  III. }
\marginpar{[ p.129 ]}\pstart tur et legantur quidem seculares literae:itemque  aliae artes ediscantur, sed ne verbo Dei, eiusque lectioni praeferantur ab vllis Christianis, nedum ab Episcopis. Ex quo fit, vt recte reprehensus sit a Gregorio Magno Episcopus quidam, qui in coetu Ecclesiae pro eo, quod debebat ipsam Euangelii doctrinam explicare, et populum exhortari, ex ea consolari, et docere, Grammatices rationem in suis concionibus exponebat et persequebatur, vti qui hodie in Actis Apostolorum aut veteri testamento interpretando, in sola Chronologiae vel Topographiae ratione, et descriptione explicanda toti sunt et immorantur. Vnde mira et damnanda fuit Heliodori Thessalonicensis clerici responsio, qui cùm interpellaretur vt fabulosam illam Æthiopicam Historiam a se conscriptam, quae tota amoris insani imaginem pingit, deleret: respondit, se male Episcompatu priuari, quam vt huiusmodi scriptum suum periret. Hoc igitur vt non decet Eplscopum: ita nec illud tam esse communis literaturae ignarum, vt si Latine loqui cogatur, barbarismos aut soloecismos ridicule committat. Quanquam propterea tamen non erit is a munere suo deponendus, quemadmodum scribit Augustinus de Catechizand. Rudibus   \pend\pstart Hac autem voce Αιδακταὸν continetur et illud, vt Episcopus sit prudens et sciat οὐρθοτομεῖν, verbum Dei(vt loquitur Paulus 2.Timot.2.vers.15.) id est, verbum Dei pro tempore, personis, locis cuique ad aedificationem applicare. Docebant autem oe lim Episcopi stantes, et κολοβίωνα induti vt dicit Epiphan. haeres. 69. de Arrianis. Est autem κoλοeιoν breuis tunica sine manicis et pectore tenus  \pend
\section*{AD I. PAVL. AD TIM. }
\marginpar{[ p.130 ]}\pstart demissa Gall.vn Roquet postea coleron. Sed docebant in Ambone stantes qui erectus erat, in coetu, vt facilius audirentur docentes Episcopi Socr. lib 6.cap.5. Est autem Ambo sedes eminens in coetu aliquo.  \pend
\phantomsection
\addcontentsline{toc}{subsection}{\textit{ Non vinosum, non percussorem, non turpiter lucri cupidum: sed moderatum, ajienum a pugnis, alienum ab auaritia. }}
\subsection*{\textit{ Non vinosum, non percussorem, non turpiter lucri cupidum: sed moderatum, ajienum a pugnis, alienum ab auaritia. }}\pstart Supra virtutes explicauit, quae decent Episcopos: nunc exponit vitia, quae dedecent eos. Itaque est hic versiculus ἀντίθεως. Non omnia tamen enumerat, sed tantùm quaedam, ex quibus maius Ecclesiae offendiculum praebetur, ad quorum exemplum de aliis, quae praetermissa sunt, statuere possumus quid sit de toto hoc genere definiendum. Multa vero, a quibus non tantùm Episcopi, sed in vniuersum clerici abstinere debent sunt in cap.clerici officia Extr. de vita et honesta. clericor. his verbis. Clerici officia vel commercia secularia non exerceant, maxime inhonesta. Minus ioculatoribus, et histrionibus non intendant. Et tabernas prorsus euitent, nisi forte causa necessitatis in itinere constituti. Ad aleas et taxillos non leuant, nec huiusmodi ludis intersint. Coronam et tonsuram habeant congruentem. Et se in officiis Ecclesiasticis, et aliis bonis studiis exerceant diligenter. Clausa deferant desuper indumenta mmia breuitate vel longitudine non notandaPannis rubeis aut viridibus, necnon manicis aut sotularibus consuticiis, fraenis, sellis, pectoralibus, calcaribus deauratis, aut aliam superfluitatem ge\pend
\section*{CAPVT  ILI. }
\marginpar{[ p.131 ]}\pstart rentibus, non vtantur. Cappas manicatas ad diuinum officium intra Ecclesiam non gerant, sed nec alibi, qui sunt in Sacerdotio et personatibus constituti: nisi iusta causa rimoris exegerit habitum transformari. Fibulas omnino non ferant, neque corrigias auri vel argenti ornatum habentes, sed nec annulos: nisi quibus competit ex officio dignitatis. Pontifices autem in publico, et in Ecclesia superindumentis lineis omnes vtantur: nisi Monachi fuerint, quos oportet ferre habitum monachalem. Palliis diffibulatis non vtantur in publico, sed vel post collum, vel ante pectus hinc inde connexis. Videtur etiam in his ipsis paucis eniuinoranuio, a uun nanuille rucronem gentis, in qua versabatur Timotheus. Erant enim Graeci et Asiani, quanquam hic est doctrina generalis, quae ad omnes pertinet. Ergo vetat Episcopum esse  \pend\pstart Vinosum) Non damnat quidem vsum vini aut aliorum ciborum Paulus, id quod et Manichaei, et alii haeretici faciunt: sed modum in iis adhibendum esse ostendit praesertim a pastore, a quo Tenperantiam supra requirebat. Quaesitum autem qui sint πάροινοι appellandi, vtrum ebrii vel ebriosi homines tantùm: an etiam qui nimia vini voraplate titillantur, illique sunt nimium addicti Gall. frians de vin. Chrysostom. πάροινον accipit pro ferocis ingenii homine, quales esse solent qui nimium se vino ingurgirant. Canonistae ca.1. distinct.35. exponunt definiunt que eum, qui se non semel, sed saepius vino repleuit, et ingurgitauit et ebrius factus est. Vtraque interpretatio non satis quadrat. Ratio enim huius praecepti est ex ef\pend
\section*{AD I. PAVL. AD TIM. }
\marginpar{[ p.132 ]}\pstart fectis vini ducta. Ea vero nimirum sunt luxus, omnis intemperies et libido. Vera.n. est illa Hiero. sententia, ad Titum 1. Vbicunque saturitas atque ebrietas fuerint, ibi libido dominatur. Nunquam ego ebrium castum putabo. Fx quo, fit vt non tantùm ebrii et ebriosi: sed etiam qui ardentius vinum appetunt, hac voce comprehendantur, damnentur, atque prohibeantur, ne sint Episcopi aut pastores Ecclesiae: aut si in huiusmodi vitia inciderint, vt deponantur. Quare Vinolentia, Ebrietas, et Luxus, damnantur in Eplicopo non ab homine: sed a Dei spiritu, vt hoc notent pastores vinosi et ebriosi. Imo vero saepissime ebrii. De quo ipse Calu.episto.370. videndus.  \pend\pstart Hoc eodem praecepto Pauli vetantur pastores, imo vero Christiani omnes conuiuia et inuitationes illas ad potandum ex compacto facere. Id quod etiam Ethnicis quibusdam displicuit, vt Assuero regi. Estner 1. Et contilii Laodiceni capite 55. concilio Agathensi cap. 4. vetitum est seuerissime. Augustinus etiam epist. 65. ad Aurelium August. grauissime comessationes istas et ebrietates damnat et execratur in ipso grege, nedum in pastoribus Christianis.  \pend\pstart Denique  cauponas et popinas ipsas ingredi vetaitur clerici, nisi peregrinentur in aliena regione.Id quod diutissime et sanctissime fuit, tum in Oriente, tum in tota Graecia obseruatum: adeo vt miram hanc Christianorum clericorum sobrietatem, et tanquam castitatem et temperantiam laudauerit etiam Iulianus apostata, nominis Christiani ifestissimus hostis, volueritque  Sacerdotes ipsos Paganorum eodem modo cauponis et po\pend
\section*{CAPVT  III. }
\marginpar{[ p.133 ]}\pstart pinis abstinere, vt est apud Sozom. liber 5. cap.36. et extat de ea re decretum concilii Laodiceni cap. 24. Denique cano. Non oportet, et cano. seq. distinct.44. de iis rebus, et ministris Ecclesiae qui huiusmodi loca adeunt et frequentant seuerissime edictum est.  \pend\pstart Non percussorem) Decet enim famulum Dei mansuetudo animi atque in omnes humanitas. Vti enim Christus fuit animo miti, sic pastores, qui illius imaginem hic in Dei populo gerunt, mites esse oportet. Vnde nec inferiores clericos etiam peccantes pugnis caedere licuit vnquam Episcopo, cui tamen multa in clericos suos iure potestatis Episcopalis concedebantur et licebant. Sed nec quenquam e plebe et grege Domini caedere potest, sed ne Iudaeos quidem ipsos aut idololatras, qui sunt ab Ecclesia alieni, vt a blasphemiis in Deum eos cogat abstinere. Toleta. concilio 4. cap. 56.cano. Neminem quisquam distin. 35. Verborum enim correctione, non verberum metu timeri debet Episcopus, vt Grego. Magnus respondit. Nec pastores percussores excusat quod Dominus noster ementes et vendentes de templo flagello eiecit. Nam illud iure supre mae potestatis, et magistratus vtraque  quae illi soli in Ecclesiam competit, fecit Christus. Nam in Ecclesia ipsa ve re est et fuit et Pastor et Magistratus summus. Ius tamen patriae potestatis in liberos suos non amittit Episcopus propter hoc dictum Paulimneque ius suum in seruos, quin eos moderate cùm res postulat, castigare etiam verberibus possit. Pertinet enim haec Pauli sententia ad eos, qui feroci et militari animo et imperioso pugnos et ver\pend
\section*{AD I. PAVL. AD TIM. }
\marginpar{[ p.334 ]}\pstart bera adhibent, vbi verba et sanctas exhortationes aut admonitiones afferre et in promptu habere debebant: non autem ad eos, qui in suos domestica disciplina vtuntur.  \pend\pstart Denique quaesitum est de Episcopo, vtrum ad bellum proficisci possit, in acie armatus stare, duellum sibi delatum committere, et acceptare. Denique torneamenta, quae vocant, vel ipse obire, vel spectare? Ac quidem quod ad militiam pertinet, scripsit Origenes Homilia 6.ad cap.8. Leuit. Est et aliud opus, ait, quod fecit Moses, ad bella non vadit, non pugnat contra inimicos. Sed quid facit? orat. Id quod est relatum in canon. Si quis vult dist. 36. Adeo vt vetentur clerici manu authoritateque  propria ipsi in hostes arma sumere, et cum aliis fidelibus (nam de infidelium exercitu dubitandum non putabant, quin omnino ab eo Episcopi abesse deberent) ad arma concurrere, in hostem pugnare, seque  illi obiicere. Id quod cùm fieret a Gallicanis Episcopis et Siculis repre hensum grauiter est, vt est in can. Reprehensibile et can. Si quis Episcopus S repreheduntur 23. quaest. 8. Nam prorsus abesse a consiliis pugnae, secularisque  militiae negotiis Episcopos oportere sancitum est synodorum decretis propter illud Pauli 2.Timo. 2.vers.4. Nemo militans Deo, implicat se secularibus negotiis. Sed etiam ab omni caedis humanae participatione oportere eosdem immunes esse decretum est, quod ea res pertinere videretur tum ad tuendam ministerii Euangelici innocentiam et sanctitatem: tum etiam ad authoritatem illius commendandam et augendam apud omnes Braccarens. Synod.cap.26. et est in canon. Si quis  \pend
\section*{CAPVT  IIE }
\marginpar{[ p.135 ]}\pstart viduam distinct. 5o. Vnde postea synodo Ilerdensi fuit cap.I.constitutum, vt clerici etiam qui in obsidionis necessitate positi hominem, id est, hostem obsidentem occidissent, tribus tamen annis poenit entiam agerent, et ab officii sui executione suspenderentur. Nimirum tam est visum Episcopali mansuetudine et munere indignum eum arma tractare. Ex quo etiam factum est et constitutum, vt clerici pugnantes in duello etiam sibi oblato deponantur. Nec modo ipsi pugnantes castigentur, sed etiam si huiusmodi duella et homicidia ab aliis facta spectent cum caeteris, de gradu deiiciantur: et in vniuersum omnia athletarum spectacula et hominum gladiatorum, et in arena pugnantium ludi et certamina sunt Episcopis interdicta, et eorum etiam aspectus, et delectatio Concilii Laodic. cap.42. et est in canon. Non oportet de consecrat. distinct.5. Ex quo iam facile colligi potest eos non debere bellis interesse, in acie armatos stare, in exercitu versari, et certaminibus praesentes se praebere, ex istarum Synodorum decretis.  \pend\pstart Sed obstant huic sententiae multa et veterum dicta et exempla, et ipsa adeo sacrae Scripturae authoritas. Ac quidem dicta, veluti Ambros.libus  1. de Officiis cap.27. fortitudo, quae bello tuetur a Barbaris patriam, vel domi deffendit infirmos, vel a latronibus socios plena iustitia est. Hieron. in Sophoniam cap.I. idem sentit. August. liber  22. cap.74. contra Faustum et multa huiusmodi collecta sunt in tota 23.quaest 1.2. et 3. Exempla etiam extant tum praetorianorum, qui dicebantur, Episcoporum, qui imperatoris comitatum sequebani. iiii.  \pend
\section*{AD I. PAVL. AD TIM. }
\marginpar{[ p.139 ]}\pstart tur, siue in bello siue in pace: tum etiam Leonis 4. Roma. Episcopi quod refertur in cano. Igitur cùm saepe 23 quaest.8. Ratio quoque idem suadet. Operaepretium est enim interesse in exercitu homines pios et Episcopos, qui militum insolentiam reprimant, fundant quotidianas preces ad Deum pro toto exercitu, eum consolentur et animent, vt dicitur Ambrosius ipse pius Episcopus versatus esse in castris Theodosii aduersus Eugenium proficiscentis, et expeditionem parantis. Denique authoritas verbi Dei istarum synodorum omnibus decretis obiicitur et opponitur. Nam et Moses et Aaron bellis iis interfuerunt, quae suo tempore a Dei populo gerebantur: et Numer Io.vers.1.2.3. iubentur Sacerdotes sacras tubas deferre in bellum, illique presentes adesse. Apparet ex epist.225. Bernardum Monachum pro pace inter regem Galliae Ludouicum et Comitem Burgundiae ineunda imploratum. Sed ex his locis factum est, vt in cano. Si in morte Longobardorum s sed notandum est dist.23. quaest.8. distinguatur insulse inter eos, qui munus Episcopale habent. Alii enim ex iis, Solis decimis contenti sunt easque habent pro victu, qui sunt nudi et simplices Episcopi, quales in Italia: sed alii Praedia, villas, castella, et ciuitates possident, quales Episcopi qui vulgo Barones, Comites, Duces dicuntur et Pares, vt in Germania, et Gallia multi. Ipsis enim Episcopatibus annexae sunt dignitates Baronatus, Comitatus et Ducatus, electoratus et pariatus. Solis et nudis siue simplicibus Episcopis arma tractare, bellisque interesse concedi oportere insulsi homines Pontificii et Canoniste  \pend
\section*{CAPVT  III. }
\marginpar{[ p.187 ]}\pstart negant. Aliis autem, qui dignitates politicas suo Episcopatui adiunctas habent, licere putant, propter eas dignitates. Ego vero respondeo. Primùm vitiose omnino huiusmodi munera politica annecti et coniungi cum Ecclesiastico munere, siue Episcopatu, siue Archiepiscopatu, siue alio quolibet, cùm diuersa sit Ecclesiastica disciplina a politica siue ciuili: et ius soli, vt vocant, a iure poli sit toto genere differens diuersum et distinguendum. Itaque pessimo fundamento et magno prorsus abusui ista distinctio innititur. Deinde falsum est simplices Episcopos non posse interesse iustis et legitimis praeliis, et bellis. Tam enim gregem suum sequi in aduersis rebus, quam prosperis, et tam in bello, quam in pace debent Episcopi omnes, qui vere suo officio fungi volent, si modo per valetudinem id facere ipsi possunt. Praeterea cùm iusta bella sint Dei ipsius bella, non hominum, vt est 2. Chronic. 19. v.7. vti praeter Domini os, et verbum suscipi non debent:ita nec sine eo administranda sunt. Nisi autem adsint pastores et ministri verbi et oris Dei interpretes, non possunt bella vel suscipi, vel geri confulto Domini ore. Ergo eos interesse iustis bellis et expedit et oportet. Sed et in acie ipsa et praelio versari possunt, et dimicare, cùm sancta et honesta causa belli suscipiendi fuit, non modo ad reliquum exercitum confirmandum, sed etiam vt Deus ipse inuocetur et proteruia insolentiaque militum reprimatur. Quare sunt omnino inepti qui D. Zuinglium optimum Dei seruum et verbi ipsius praeconem egregium damnant, quod in aciem  \pend
\section*{AD I. PAVL. AD TIM. }
\marginpar{[ p.138. ]}\pstart fuerit. Id enim et bona conscientia in honesta belli causa potuit, et idem nunc quoque licet caeteris. Ministri vero stratioticis moribus et plane militaribus sunt omnino damnandi vt Cal. epist.354. docet, propter ferociam ingenii, et proteruiam quae mansuetudini pastorali contraria est.  \pend\pstart Non turpis lucri cupidum) Hoc tertium est vitium quod in Episcopo futuro atque etiam electo damnat Paulus, tanquam omnis mali et offendiculi maximi fontem, nimirum auaritia. Est enim mater et radix omnium vitiorum. Similis locus. Tit .1.vers.7.1.Pet.5.vers.3. Sed vt nec lucrum quaerere nec captare debet Episcopus: ita neque gloriam suam spectare, nec ventrem conmoditatemque suam aliquam pro fine sui ministerii sibi proponere. Galat. 6. vers. 13. Philip.I. Est quidem dignus operarius mercede sua, vt ait Christus Matth.1o. vers.1o. sed aliud est alimentum et victus necessarius: aliud autem lucrum et imprimis turpe, quod hic spectari a nobis vetat Paulus. Deinde nec hic ipse debet esse doctorum Ecclesiae finis, dum hoc munus tam sanctum, vel suscipiunt vel exercent, vt inde viuant et alantur, sed vt Dei gloriam promoueant. Alimentum enim et victus quaedam accessio potius est, quam se opus te finis huius muneris, et recte facientem officium sequi debet alimentum, vti vmbra corpus sequitur. Sed quaesitum est, Quod sit lucrum illud turpe, quod consectari vetantur Episcopi, et pastores Ecclesiae? Respond. Omne illud, quod legibus municipalibus, vel honestis moribus accipi prohibetur, pugnant que cum verbo Dei, aequitate na\pend
\section*{CAPVT  III. }
\marginpar{[ p.139 ]}\pstart turae, et vera Christianaque charitate. Nam ueberene enempium gregla partor, e totus ex Dei prouidentia pendere.  \pend\pstart De accessione igitur quaesitum est, etiam ea, quam ex constitutionibus Principum vel ciuitatis cuiusque nobis accipere et stipulari licet, vtrum eam exigere ex pecunia a se mutuo data possit pastor Ecclesiae et Episcopus? Et quidem canone Apostolorum 44. Niceni concilii cap.17. et Laodiceni cap. 5. omnino vetitum et interdictum prorsus est Episcopis, vt neque suo, neque etiam alieno nomine (veluti si tutores alicuius pupilli essent, vel alieni negotii procuratores) vsuras etiam aliis concessas ipsi exigant, et stipulentur: quod id praeter eorum dignitatem esse videatur, sitque haec summa auaritiae (quae ab Episcopis abesse debet) species. Idem tota distinct. 47. est constitutum et definitum. et Elibertino çoncil. cap.2o. Sed tamen duriora illa, vt verum dicam, decreta videntur, quod omnino sine exceptione perceptio qualiscunque fructus ex pecunia mutuo data pastori interdicitur, quae tamen saepe bona conscientia fieri potest. Itaque potius amplector D. Caluini consilium in epistola 346. vt censeam omnino quidem et in genere illud in Ministro non esse damnandum, si fructum aliquem ex sua pecunia sentiat, moderationem tamen non tantùm ex legibus eius ciuitatis, in qua versatur, sed potius ex verae et Christianae charitatis sensa et praeceptis esse adhibendam, et quidem eam, vt omne offendiculum iustum Pastores caueant, eamque legem adhibeant, quam ipse Caluinus ibidem praescribit, et necessariam esse docet  \pend
\section*{Ap l. PAVL. AD' TIM. }
\marginpar{[ p.140 ]}\pstart Quae tria sequuntur sunt trium superiorum vielum ex ἀντιSεσιι explicationes. Itaque longiori iam intepretatione non egent.  \pend
\phantomsection
\addcontentsline{toc}{subsection}{\textit{4 Qui suae domui bene praesit, qui libeTos conbineat in juviectiont tum omni honeJtate. }}
\subsection*{\textit{4 Qui suae domui bene praesit, qui libeTos conbineat in juviectiont tum omni honeJtate. }}\pstart Τάξις cum transitione coniuncta. Postquam enim docuit quae in Episcopo spectari debent ex ipsius persona, iam addit quae aliorum ratione in eodem requirantur. Ac primùm quidem Liberorum ratione et totius familiae. 2 Vocationis ipsius ad Euangelium. 3 Extraneorum. Nam quaedam in eo ex tribus his causis cnosyderari oportet, quis hic ordine tradit et tractat Paulus.  \pend\pstart Ergo familiae et imprimis Liberorum institutio spectari debet, qualis ab eo facta sit, qui in Episcopatum petitur. Qui enim domui suae recte praeesse non potest, quomodo idem poterit Ecclesiam, quae maioris est ponderis, gubernare?  \pend\pstart Simile praeceptum est Tit. 1.v. 6. Ratio huius praeceptionis est tum ea, quae a Paulo postea subneitur: tum quod quodammodo ingenium paternum referunt filiorum mores, adeo vt ex illis iudicari possit, quam sit fedulus et diligens pater in munere suo obeundo: et quam erga Dei gloriam vehementer sit affectus, et serio, maxime. cùm facile imperium habere possit pater in liberos, si semel recteeos instituerit. Nam qui domi suae gloriam Dei sperni videt, et patitur, ille facile, idem patietur fieri alibi, vbi tam assiduus non est, et vbi tam aequum et facile imperium non habebit.  \pend
\section*{CAPVT  III. }
\marginpar{[ p.14t ]}\pstart Haec autem cura familiae, quae in Episcopo laudatur, non est in eo, vt attentus sit ad rem, ditet et locupletet suos liberos: et vt optimum se oeconomum rerum, quae ad alimenta familiae pertinent (quanquam et haec curanda sunt) ostendat, spreta morum in sua familia cura. Denique non eo pertinent, vt astutum, vel callidum, vel addictum bonis opibusque terrenis familiae suae acquirendis Episcopum esse debere doceat Paulus. Auaritiae enim ista sunt praecepta potiùs et effecta, quam pietatis. Itaque vt ait hoc loco Chrysostom. non eo pertinet mens huius praecepti Pauli:sed ad curam pietatis et morum, quos sanctos et honestos a patre Episcopo futuro et tradi et traditos esse liberis vult Paulus. Id quond etiam ipsa eius verba satis probant. Vult enim liberos Episcopi institutos esse, in omni Reuerentia et subiectione praestanda patri ipsi, et in omni Honestate et sanctitate. Addit Paulus in Epistola ad Iitum. In fde ettam Christiana eos a patre edoctos, et informatos esse oportere. Nec enim is eligendus est Episcopus, cuius familia idololatra est: vel a quo liberi ipsi non sunt in fide et doctrina verae pietatis et Christianae imbuti et instructi.  \pend\pstart Ex hoc autem loco confirmatur nostra sententia, qui negamus coelibatum Episcopis futuris imponendum, cuius ipsam familiae et liberorum gubernationem certissimum argumentum esse indolis episcopalis, ingeniique ad maius munus strenue capessendum idonei demonstrat Paulus hoc loco.  \pend\pstart Item nominatim meminit, liberorum futurs  \pend
\section*{AD I. PAVL. AD TIM. }
\marginpar{[ p.142 ]}\pstart Episcopi. Ergo eum vxorem vel antea habuisse, vel adhuc habere satis colligitur. Nec enim illa Hieronym. sententia et allegoria probanda est, qui nomine filiorum et filiarum Episcopi cogitutionte te epr Aintelligi vult, vt scribit in Epist. ad Titum 1. Verba enim haec Paulisine allegoria debent accipi.  \pend\pstart Praeterea ex hoc ipso loco docetur, quam longe sint idoneiores ad Episcopatum coniuges, quam coelibes viri, et vt specimen quoddam maioris humanitatis, et industriae, prudentisque administrationis in iis appareat, qui familiam habent, id est, vxorem et liberos: quam qui perpetuo coelibes et solitarii vixerunt. Certe vt non debet homo Christianus esse ἀπολίς, et tanquam fera quaedam in solitudine vaga viuere:ita etiam non debet ab iis officiis humanitatis alienus esse et expers, quae quodammodo naturam nostram excolunt, redduntque  politiorem. Ex quo fit, vt totum genus illud hominum, qui in vita Monachica consenuerunt, sit ad munus Episcopale ineprissimum, non tantum propter ignorantiam disciplinae et communis hominum societatis: sed etiam quod in solitudine enutriti pene solent esse ἀυθαδας, et magna humanioris illius conuersationis suauitate curtt, per quam hominum societatem inter se ipsos tuentur. Et tamen aliquando hic fuit receptus mos, vt Episcopi ex hoc hominum Monachorum genere et coeno maxime eligerentur, vt apparet ex cano. Monachus distinct. 77. Id quod suo seculo probatum fuisse scribit Ambrosius epist.82. Sed hi homines secum in Ecclesiam maximas et pessimas superstitiones quoque, quas semel iam  \pend
\section*{CAPVT  III. }
\marginpar{[ p.13 ]}\pstart imbiberant, inuexerunt magno certe Ecclesiae damno. Hodie quidem Monachus sine Abbatis sui consensu, nec clericus, nec Episcopus fieri potest, non propter superiorem rationem: sed quia Monachi mancipiorum loco censentur. Sunt enim Abbatis serui, vt apparet distinct 54. et 58.  \pend\pstart Quaerit Hieronym. an sit aequum propter immodestos liberorum mores patres arceri abEpiscopatu, cùm eadem ratione videantur liberi ab Episcopatu prohibendi, nempe, propter sceleratam parentum vitam. Quod tamen prorsus esset iniquum. Deinde quam absurdum est, ait ille, censere, Isaacum propter Esau filii turpem vitam, si viueret, ab Episcopatu arcendum esse? Sed respondeo, primùm non esse parem rationem, cur propter vitia parentum filii prohibeantur hunc gradum assequi:atque cur parentes non admittantur propter sceleratos liberorum suorum mores: Nam filii parentes suos regere non debent, quemadmodum patres debent instituere liberos suos. Itaque animi morumque filil iudicium non fiet ex paterno luxu, quod sit neçessarium et validum argumentum. At ex dissolutis familiae moribus intelligitur quale sit ingenium ipsius patrisfamilias, Deinde respondeo. Hic Paulum agere de tota familia, non de vna tantùm eius parte. Nam si e pluribus liberis in eadem familia institutis, alii quidem sancte instituti cernantur: alii vero luxu diffluentes appareant, hic luxus non tam ex paterna disciplina et moribus nasci videbitur, quam ex prauo ipsorum liberorum ingenio, quod est procliue  \pend
\section*{AD I. PAVL. AD TIM. }
\marginpar{[ p.144 ]}
\phantomsection
\addcontentsline{toc}{subsection}{\textit{(Nam si quis propriae domui praeesse nescit, quomodo Ecclesiam Dei curabit?) }}
\subsection*{\textit{(Nam si quis propriae domui praeesse nescit, quomodo Ecclesiam Dei curabit?) }}\pstart Argumentum est a minori ad maius, quo vtitur Christus Matth. 24. de talentis agens. cuius ratio et firmamentum est, non modo certissimum: sed etiam notissimum, et ab omnibus concessum, hoc nempe, eum qui minorem nauim regere non potest, nec maiorem posse. Maior autem vis huius argumenti apparet, in his verbis Propriae domui: Ecclesiam Dei Est enim Ecclesia et ipsa do mus vt inf. vers.15. docebit Paulus, et est Hebr.3. vers.3. 4.5. sed domus non priuati cuiusdam hominis, sed regis: non etiam familia et domus alicuius humani regis mortalis: sed Dei ipsius summi et celestis omnium imperatoris. Ex hoc autem intelligitur quanta sit pastorum Ecclesiae dignitas, qui in domo Dei praeficiuntur. Ex hoc quoque colligimus Ecclesiae dominium ad quemquam nomiiiu non pertinere: nec cuiquam licere in ea sibi ius authoritatemque , nisi Deo ipso, id est, Domino deferente vsurpare, vt est in epist. ad Hebr.  \pend
\phantomsection
\addcontentsline{toc}{subsection}{\textit{6 Non nouitium: ne inflatus, in criminationem incidat diaboli. }}
\subsection*{\textit{6 Non nouitium: ne inflatus, in criminationem incidat diaboli. }}\pstart Non nouitium, vel Neophytum, Metaphora traducta a plantis ad homines fideles, qui etiam in Dei domo, et Ecclesia plante nominantur. Psal 92, vers. 14 et Christus ipse ita appellat, omnis arbor, vel plantatio, quam Pater meus non plantauerit, etc. Matth. 15. Ac primùm qui Neophyti dicti sunt explicandum est, propterea quod abusi sunt Monachi hac voce ad votorum suorum distinctionem  \pend
\section*{CAPVT  III. }
\marginpar{[ p.145 ]}\pstart stinctionem. Dicuntur enim Neophyti ab istis ii: qui propositum alicuius religionis recenter assumpserunt, suntque adhuc intra annum suae, quam vocant, probationis. Sic definitur cano. Quoniam 5 Neophytus distinct.48. Apud veteres autem alii Neophyti dicebantur, ii nimirum, qui nouiter baptizati fuerant, et excesserant ex Catechumenis. Catechumeni enim erant qui fidei Christianae nomen dederant quidem, nondum tamen erant baptizati, quod nondum satis in fide edocti et instructi viderentur. Sic in concilio Chartag 4.cap.3o. appellantur.  \pend\pstart Cur autem de hoc admonuerit Ecclesiam Paulus, ratio est, quod tanta Dei dona in Ecclesiam mire effusa erant, vt etiam saepe adhuc nouitii, et tunc primum recepti in fidem homines excellerent magnis Dei donis, praestanti doctrina, et ingenio ornati essent a Deo. Hos tamen vetat admitti indiscriminatim Paulus, quia nondum satis probati sunt. Nam qui tam cito promouentur non modo turgidi sunt, inflati et ambitiosi: sed etiam saepe regendae Ecclesiae sunt ignari. Vnde facile in Diaboli laqueos incidunt, et ab officio modestia que abripiuntur, quae in pastore Ecclesiae requiritur.  \pend\pstart Haec praeceptio Pauli sumpta est ex eo Mosis praecepto, vt docet et vult Bernard. sermo. 64. Cantic. vbi Deuter.15. vetat lex, ne in primogenito bouis quis aret. Quae legis interpretatio sane allegorica nimis est. Potius ex ipsa rerum experientia ducta est. Neophytos enim in omni actionum genere, si statim promoueantur, superbos esse certissimum est. Itaq' haec Pauli praeceptio etiam  \pend
\section*{AD I. PAVL. AD TIM. }
\marginpar{[ p.146 ]}\pstart Synodo Nicaena prima repetita est, vt apparet cap.2. et canon. Quoniam multa, distinct.48.A163 haec fuit vetus Ecclesiae disciplina etiam post Apostolos, tamen postea contempta est. Nam cùm Ambrosius adhuc Neophytus extruordinarie Mediolanensis Episcopus creatus fuisset, voluerunt homines statim post baptismum etiam quosuis eligi in clericos vel Episcopos posse. Id quod suo seculo iam factitatum fuisse scribit Hieronymus, qui fuit Ambrosio aequalis.  \pend\pstart Veruntamen haec Pauli regula non vsque adeo praecisa est, et stricta, vt omnino ab ea, ex iusta tamen causa, discedi non possit. Sed iustam causam esse oportet, quia huiusmodi personarum electiones sunt certe periculosae, et tanquam extraordinariae. Si tamen vel Ecclesiae necessitas, vel eximiae Neophyti dotes, vel summa eius modestia iam probata in aliis rebus, eum eligendum suadeant, eligendus est, neque a Pauli praecepto discedemus, quia rationem subiicit Paulus, propter quam vetat eos eligi, non quod Neophytos omnino neget posse ex aliqua iusta causa vocari ad Episcopatui ver doctoratunt: sed ne tales inflentur. Quod periculum ab iis abest, qui iam suae modestiae magnum et certum specimen ediderant.  \pend\pstart Iam vero huius praecepti occasione veteres synodi quaedam sanciuerunt iquae non sunt omnino vel praetermittenda, vel vt inutilia negligenda.  \pend\pstart Et primùm voluerunt, vt qui ad maiores gradus Ecclesiasticos promouerentur, illi ex minoribus ordinibus et inferioribus gradibus assumetentur. Id quod audio adhuc in Boëmicis Eccle\pend
\section*{CAPVT  III. }
\marginpar{[ p.147 ]}\pstart siis stricte et religiose obseruari. Hoc decretum extat tot. Tit. et distinct.52. Vt nemo per saltum ordinetur Dicebantur autem per saltum ordinari et promoueri qui praetermissis omnino aliis gradibus, vel mediis et inferioribus non acceptis et degustatis, ad maiores dignitates prouchebantur. Erant enim in ecclesiasticis muneribus certi ordines constituti, per quos gradatim erat ad Episcopatum conscendendum, ne quisquam non satis probatus, aut expertus in pastorem Ecclesiae, et tam arduum munus eligeretur. Erant igitur hi ordines Ostiarius, Lector, Exorcista, Cantor, Cerearius, Acolytus, Subdiaconus, Diaconus, Presbyter, Episcopus, can. illud dist. 77. Sed hos omnes tamen rigide obseruare in singulis eligendis ac velut ex necessitate quadam, est certe superstitio, cùm etiam ex istis quidam ordines non sint. Satis est, in puuigiima ratione io, qui m doctorem vel pastorem Ecclesiae asciscendus est, probatus sit in doctrina et moribus: siue quod in aliquibus inferioribus gradibus iam substiteritisiue quod alibi alióque  modo sit probatus. Nam ex istis tot ordinibus, quidam sunt, et inanes et superstiriosi.  \pend\pstart Denique hoc ipsum postea videmus in abusum, et ridiculum quiddam in Papatu fuisse conuersum. Omnes enim isti ordines sacrificulandis ab Episcopo conferuntur, sed frustra, et sunt tituli et dignitates vanae et sine effectu. Secundo Ioco, Certa quaedam aetas cuique ordini Ecclefiastico adipiscendo est Synodis constituta, infra quam haberi obtinerique illi gradus non possent, veluti Synodo sexta vniuersali, concilio Chartag. 3. cap. 4. Agathensi cap.16. Toletano 4. cap.17.  \pend
\section*{AD I. PAVL. AD TIM. }
\marginpar{[ p.148 ]}\pstart Itaque decreuerunt vt Diaconi quidem ante vicesimum quilitun ae tatis amum nulli eligerentur.  \pend\pstart Presbyteri vero et Episcopi nulli ante tricesimum aetatis annum, quae est iusta hominis aetas, crearentur cano. Placuit distinct. 77. et tota dist. 78. Quod ex Numerorum cap. 4. et 8. sumptum aut imitatum videri potest. Et certe aequum est iustam aliquam aetatem, et certam definiri, aut spectari in Ecclesiasticis personis eligendis, ad quam tamen dona Dei minime sunt restringenda. Nam et Timotheus quamuis ipse admodum adoIescens esset, tamen iam erat Euangelista: et Dominus, cur vuit, et quando ruit udna sua largitur: vt potius hoc totum ex dotibus ingenii, et Dei donis erga singulos sit aestimandum et diiudicandum, quinam sint idonei et eligendi, quam ex certo annorum numero a nobis constituto.  \pend\pstart Tertio loco quaesitum est, vtrum Episcopum futurum oporteat, secularium, quae vocant, negotiorum peritum esse, eaque tractare et admininistrare posser in quo videtur fuisse diuersa tum Patrum, tum Ecclesiae veteris disciplina, et consuetudo. Ac primùm in hac quaestione, disciplina Ecclesiastica a negotiis secularibus distingui debett Dheiplinae. n. Ecclesiasticae imperitus et ignarus non debet esse Episcopus, ne magna in Dei Ecclesiam perturbatio confusióque ab eo inducatur. Rudes enim disciplinae et ordinis Ecclesiastici homines saepe turpissime grauissimeque peccant. At que huiusmodi disciplinae ignorantes homines etiam electos, repudiatos tamen esse apparet ex distin.59. cano. Qui Ecclesiasticis.  \pend
\section*{CAPVT  III. }
\marginpar{[ p.149 ]}\pstart Negotia vero secularia quid sint, est videndum. Non sunt autem omnia, quae hic in mundo (qui seculum vocatur) a nobis aguntur, omnis enim rerum etiam priuatarum, et domesticarum administratio et cura Episcopis interdiceretur, cùm dicat Paulus 2.Timoth.2.vers.6. Nemo militans Deo implicet se secularibus negotiis. Sed ea tantùm negotia dicuntur secularia, quae ad huius vitae corporeae vsum, et commoditatem spectant omnino, et per se, nósque a cura rerum diuinarum auocant, siue nostro, siue alieno nomine a nobis illa gerantur, qualis est Negotiatio, siue Mercatura mercium, Aduocatio in Praetoriis et foro iudiciali, Ratiociniorum alienorum susceptio, Procuratio domorum, bonorumue alienorum, Conductio publicorum redituum: canon. Perlatum est ad nos. et canon. Consequens distinct. 88.canon. Placuit 21. quaest.3 canon. Te quidem oportet 1I. quaest.i, quia in his omnibus lucrum quaestusque per se spectatur, et solius vitae huius corporeae causa geruntur, et instituuntur haec negotia.  \pend\pstart Pauperum quidem, Viduarum, Pupillorum, Orphanorum, et Extraneorum, denique omnium miserabilium personarum, quae dicuntur, negotia Episcopo curae esse debent, quia est Dei seruus Non tamen vt illis sic implicetur, vt a lectione scripturae, concionibus de verbo Dei, reliquisque  sui muneris par tibus propterea abstrahatur. Hoc enim fieri vetat Paul.2. Timoth.2.V.6. et confirmatur exemplo Apostolorum Actor.6.vers.2. Adeo vt iam tempore Cypriani Chartaginensis Episcopi ex vetusta quadam Ecclesiae disciplina excusarentur Pa\pend
\section*{AD I. PAVL. AD TIM. }
\marginpar{[ p.1 Q ]}\pstart stores Ecclesiae a munere tutelae etiam testamentariae, quemadmodum est in ipsius Epistolis liber  1. Epistola 9. et positum extat in can. Cyprianus 21.quaest.3. Itaque saepius repetitum est in synodis, ne Episcopi, imo vero ne clerici vlli secularium negotiorum curam administrationemque susciperent aut gererent, veluti Chartaginensi Synodo, quae conuenit in Hispania cap. 6. et est relatum in canon. Peruenit et seq. 21. quaest.3. Idem antea fuerat constitutum concilio Chartaginensi 4.cap.15. Adeo vt ne vel tanquam arbiter vel iudex de huiusmodi negotiis cognoscere permitteretur, vt ex veteri quadam Clementis epistola (si modo eius est) colligunt Canonistae in cano. Te quidem II. quaest.1. Artificio tamen aliquo sibi victum quaerere non prohibetur Pastor, si suo quoque  muneri superesse potest, et propterea non abducatur. Otiosos enim eos esse vetant canones concil. Chartag.4. cap. 49. et si vide distin. 91. per totum. Quanquam propterea neque Anabaptistarum errorem, neque aliorum phanaticorum delirium probo, qui Apostolos quicquam proprium habuisse aut possedisse negant. Itaque stipendia a pastoribus Ecclesiae peti sana conscientia posse negant, illisque res soli etiam patrimoniales adimunt. Sed hic error fuit quod Ecclesiae temporalia bona, ingentia illa quidem et plena litium immensasque possessiones curare, et administrare posse crediderunt Episcopos et clericos, quia ad Ecclesiae quaestum et commodum ea cura pertinebat, id est ad paucorum hominum, qui se clericos dicebant, alimenta, et quorundam templorum conmodum, cùm tamen iis ipsis agendis magis a le\pend
\section*{CAPVT  III. }
\marginpar{[ p.151 ]}\pstart ctione scripturae, et reliquis partibus fui muneris auocentur, quam si regias opes administrarent. Sed quia ista Ecclesiastica negotia vocantur, eam curam et distractionem non damnabant: et ad eam dictum Pauli pertinere nolebant, ne lucris, et commodis ipsorum quicquam decederet et detraheretur. Atqui et haec ipsa cura, est cura rerum secularium, et implicat nos ita vt officio satisfacere et superesse non possimus, adeóque  haec ipsa cura a Paulo damnata est 2.Timoth.2.vers.6. Vnde rectius quidam Episcopi, tempore Iustiniani Imperatoris, qui alios praeter se, praeterque clericos ipsos, id est, praeter Diaconos et Presbyteros, habebant, qui ea bona Ecclesiae administrarent, quos homines Oeconomos Ecclesiae vocabant, vt apparet in constiturionibus Iustin. libus 1.Codic. Tit de Episcopis. Pro quibus postea seculo corruptissimo et crescentibus Ecclesiae bonis aduocatos nobiles habuerunt. Oeconomos autem bonorum Ecclesiae viros bonos habere praestaret etiam in iis locis, vbi post susceptum verum Euangelium, bona tamen Ecclesiastica abEpneopis retmemtur, ver au iplis Irissstris possidentur, ne in iis procurandis, et defendendis, et administrandis occuparentur, ii, qui in rebus melioribus versari, et toti esse deberent. Sed bonus est odor lucri. Peius autem illa bona ad homines aulicos transferuntur, tanquam praeda ex Papatu capta.  \pend\pstart Ergo quaeritur, num secularium negotiorum scientes esse Episcopos futuros oporteat. In quo responsum est varie et constitutum diuerse. Nam distinct 39. aperte sancitum, oportere tales esse,  \pend
\section*{AD I. PAVL. AD TIM. }
\marginpar{[ p.152 ]}\pstart Episcopos, adeo vt qui eorum negotiorum rudes sunt et imperiti, deponendi dicantur. At distinct. 88. et 21. quaest3 aliud omnino responsum est et variis patrum dictis confirmatum. Sed vtriusque responsionis tam diuersae conciliatio haec affertur a Papisticis. Nempe requiri quidem, vt Episcopi futuri eorum negotiorum rudes omnino et imperiti non sint, quoad scientiam: administratores tamen et gestores non sint, quo ad ipsam curam et solicitudinem, ne a munere suo auocentur et distrahantur. Sed quorsum in illis scientia huiusmodi rerum requiritur, si earum cura, et administratio iisdem denegatur? Et quod est in illa distinct. 39. secularium negotiorum peritiam futuris Episcopis, et iis qui iam sunt Episcopatum adepti, esse necessariam plane ex abusu ministerii Episcopalis profectum esse nemo non perspicue videt. Exempla enim eorum temporum afferuntur, in quibus Episcopi iam non tam legendo Dei verbo, et explicando vacabant, quam opibus et diuitiis nomine Ecclesiae congerendis et cumulandis, ad quas certe secularium commerciorum et negotiorum cognitio, et scientia omnino necessaria fuit. Desierant igitur iam Episcopi esse veri pastores et aurei Episcopi, et quales hic describit Paulus: sed erant iam astuti homines, auari mercatores, et ambitiosi, qui pietatis, et religionis praetextu, ampla latifundia, dignitates, amplasque  sibi possessiones comparabant, in iisque tuendis et augendis positum esse Episcopi munus censebant, quo ditesceret his mundanis opibus Ecclesia, et veras sperneret caeleslesque  diuitias. Id quod tandem euenit. Itaque bona eorum coeperunt esse  \pend
\section*{CAPVT  III. }
\marginpar{[ p.152 ]}\pstart aurea: ipsi vero facti sunt plumbei. Bernard.epistola 78. detestatur quendam sui seculi clericum, qui et Diaconus esset, et Decanus, et regius dapifer, quod intolerandum esse docet in Dei Ecclesia.  \pend\pstart Caeterum qui id Episcopis licere putant, secularibus se negotiis nimirum immiscere et implicare, aliquot rationes afferunt, quas operaepretium est examinasse, vt intra suos fines et a Deo praescriptos, tandem verus Episcopatus, et verbi Dei ministerium reuocetur et coerceatur.  \pend\pstart Prima est ex Deuteron.17. vers.9. vbi Sacerdos Magnus et Leuitae constituuntur iudices causarum, et litium. Ergo secularibus negotiis sese immiscebant. Respon. Cùm Dominus politiam ipsam, quam in suo populo vigere volebat, suo verbo comprehendisset, cuius interpretationem commiserat Sacerdoti summo, et Leuitis Malac. 3 aequum tunc erat, vt in dubia interpretatione legis politicae, quae Verbi diuini pars erat, ad eos recurreretur. Sed nunc leges politicae a verbo Dei separantur, nec eo comprehenduntur, vt olim. Praeterea ad sacrificatores tantùm recurrebatur, si de interpretatione alicuius loci, et legis scriptae dubitatio nascebatur, non in quauis controuersia, nec cùm erat quaestio facti. Tunc enim ad alios iudices remittebatur disceptatio.  \pend\pstart Secundo obiicitur cano. Ecce plenarie distin. 39. exemplum Christi, qui turbis sequentibus non tantùm spiritualia, sed etiam carnalia, veluti cibum, dedit. Debere igitur Episcopos carnalia suo gregi administrare. Itaque debere implicari negotiis secularibus. Resp. Extraordinaria Christi  \pend
\section*{AD I. PAVL. AD TIM. }
\marginpar{[ p.154 ]}\pstart facta non esse nobis in exemplum trahenda. Nec enim semper pauit Christus omnes turbas, quae ipsum sequebantur: nec quas pauit, ordinarie aluit: sed extraordinarie, non etiam se fatigans in quaerendo illis cibo.  \pend\pstart Tertio afferunt exempla Patrum veluti Danielis ipsius, qui cùm esset Propheta, tamen administrabat negotia regis Babyloniorum, et Persarum. Respond. Hoc esse extraordinarium. Voluit enim Deus ad subleuationem sui populi captiui, Danielem praeesse rebus et negotiis illius regis, a quo captiuus detinebatur, vt esset clemens seruitus sui populi. Ambrosius item Mediolanensis Episcopus laudatur, qui pro Valentiniano legationes saepius obiit ad Maximum tyrannum, qui erat Coloniae et Treuiris, vti ex ipsius epistolis apparet. Respond. Illud factum esse raro quidem, nempe cùm summa Ecclesiae et Reipubus  vtilitas postulauit, sed statim expedita legatione et explicata ad munus suum rediit Ambrosius. Itaque potest Ecclesia quidem suo pastore nonnunquam vti ad res suas et seculares agendas: si enim videat ad ea negotia aptum, et idoncum: sed tamen id rarissime faciendum est, et cauendum ne a suo munere, quod in verbi praedicatione et Ugm Ectienae cura consistit, abducatur vel tantillum. Longe igitur praestaret Episcopos, qui tam sunt ad secularia negotia tractanda idonei, vt maiorem per eos vtilitatem Ecclesia sentiat, quam per alios quosuis: praestaret (inquam )eos ex honesta causa a munere illo suo Episcopali honestam missionem habere, et liberari prorsus et dimitti, quam eos in Episcopatu retineri: et  \pend
\section*{CAPVT  III. }
\marginpar{[ p.153 ]}\pstart interim secularibus negotiis perpetuo eos distineri, distrahi, et occupari. Et haec mea est sententia.  \pend
\phantomsection
\addcontentsline{toc}{subsection}{\textit{7 Oportet autem eum etiam bonum hapere testimonium ab extraneis: ne in probrum incidat et laqueum Diaboli. }}
\subsection*{\textit{7 Oportet autem eum etiam bonum hapere testimonium ab extraneis: ne in probrum incidat et laqueum Diaboli. }}\pstart Tertium est quod requirit Paulus, sed aliorum ratione, nimirum ipsorum Extraneorum, id est, infidelium. Vult igitur futuros Episcopos tam esse probatae vitae, vt etiam a fide Christiana extranei homines nihil habeant, quod in ipsorum vita reprehendant. Hoc idem ab omnibus Christianis requirit Christus Matth.5. et Paulus Colo. 4. vers.5. quia vita nostra debet esse lux mundi: multo igitur magis haec lux in pastore desideratur. Nec obstat quod idem Paulus ait 2. Corinth. 6.vers.8. per gloriam et ignominiam nobis esse incedendum: item qui sunt foris, nihil ad nos 1. Corinth.5.vers. 12. Hoc enim verum est quod ad doctrinam ipsam. Nam doctrinam nostram nunquam probabunt infideles: quare in ea ab illorum iudiciis vllo modo non est dependendum, non magis, quam a caecis: sed in vita et moribus secus est. Nam debet esse vita illa nostra tam sancta, vt etiam ignorantes, et infirmos Christo lucrifaciamus. Confitmat suam sententiam Paulus a consequenti. Qui enim vitiis secure indulgent, sunt Probrosi, et In Diaboli laqueum incidunt, quia nullo pudore postea retinentur, cùm se vident etiam post sceleratam et turpem vitam ad tantam dignitatem tamen admitti, promoueri, et posci. Vnde fraena et habenas impune laxant omnibus  \pend
\section*{AD I. PAVL. AD TIM. }
\marginpar{[ p.20.  ]}\pstart vitiis, et a Sathana possidentur. Quod quaerit Chrysostom. Quid de iis statuendum sit, qui vel falso laudantur ab extraneis et infidelibus, vel immerito vituperantur. Responsio facilis est, Ecclesiam nimirum debere iudicium suum adhibener nec diimo ea iiidenam iudragio pendere, et statuere, quid de eligendo fit sentiendum.  \pend
\phantomsection
\addcontentsline{toc}{subsection}{\textit{8 Diaconos itidem honestos, non bilingues, non multo vino deditos, non turpiter lucri cupidos. }}
\subsection*{\textit{8 Diaconos itidem honestos, non bilingues, non multo vino deditos, non turpiter lucri cupidos. }}\pstart Transitio est, ad Diaconos enim pertinet sequeo umputatio d praeceptio, quos describit Paulus, et ordine quidem. Nam de Episcopis prius egerat: sequens igitur erat, vt ageret de Diaconis. Hos enim duos ordines Ecclesia Dei maxime habet.  \pend\pstart In quo duo quaeruntur. I Quid vox Αιάκονος hic significet. Ac ipsa per se et ex verbi etymo est latissima huius vocis significatio. Omnem enim Ministrum significat siue priuatum, siue publicum, id est, omnem eum, qui suam operam suumque ministerium alii exhibet. Vnde etiam Apostoli ipsi Diaconi nonnunquam in Scriptura vocantur. Ex quo etiam factum est, vt quidam existimarent, hic Diaconi vocabulum priuatum Episcopi famulum denotare, quem, qualis esse debeat, describeret Paulus, sed falluntur. Lata vero haec significatio postea restricta est ad eos, qui in Ecclesia Dei aerarii pauperum curam, administrationemque  habent, et circa pauperum ministerium et subleuationem occupantur, quales ii de quibus Acto.  \pend
\section*{CAPVT . III. }
\marginpar{[ p.157 ]}\pstart 6.et Philip.1.vers.1. et Rom.12.vers.7, et 8.quo fere significato accipitur cùm ab Episcopis in Ecclesia Diaconi distinguntur. Atque etiam hoc loco ista postrema et strictiori notione sumitur. Insulse vero Papistae a Diaconis Αακόνιον appellarunt locum quendam aedis sacrae penitiorem et secretiorem, in quo sacra Ecclesiae vasa reponuntur, et conduntur: quod Secretarium alio nomine nuncupant: et Diaconos Secretarios, quod huius loci claues gestent, curamque habeant can. Non oportet distinct. 23, etsi ex concilio Laodiceno videtur haẽc significatio sumpta. Sed iam erat corruptissima Ecclesiae disciplina, desierant que, qui Diaconi tunc vocabantur, esse veri Dianconi, quaeles vetus Ecclesia habuit: et hic describit Paulus.  \pend\pstart 2 Quaesitum est, cur praetermissa sit disputatio de Presbyteris, vt quales esse debeant, et eligi. intelligeremus. Respondet Chrysost. quod Episcoporum nomine Presbyteri continentur. Quod verum est vt apparet ex Epistola Pauli ad Tit. cap.1. Actor.2o. et aliis multis Scripturae locis. Differunt enim, ait, Chrysost. Presbyter et Episcopus solum ordinatione, non potestate, non gradu, non ministerio, quia aliae quaedam tantum et plures caeremoniae in ordinandis Episcopis obseruabantur, quam in constituendis Presbyteris.  \pend\pstart Ex huius autem loci et ordinis (quem hic Paulus obseruat et in Epistola ad Philip.1.vers.1.) ratione quaesitum est, Vtri gradu sint priores et superiores in Dei Ecclesia censendi Presbyterine, an Diaconi. Etsi enim Diaconi hic subiiciuntur Episcopis: videntur tamen illis ex sententia Hieron. in epist. ad Rusticum praeponi. Haee enim illius  \pend
\section*{AD I. PAVL. AD TIM. }
\marginpar{[ p.158 ]}\pstart verba sunt, vt si distinctio locorum creditur in mumdo etiam in altari Dei videant sibi Episcopi, qui superbi sunt, Diaconos anteponi. Deinde videntur ex veteri disciplina Ecclesiae etiam ad honestio ra quaedam officia et mimsteria quam Presbyteri ipsi admitti. Ipsi enim (vt ibidem tradit Hieron. quanuis conquerens) Sacerdotibus de mensa Domini calicem tradunt: non autem Diaconis Presbyteri. Id etiam transcriptum est in cano. Diaconi sunt quos distinct.93. Sed etiam videntur aliquando praedicasse, absente Episcopo, et administrasse Sacramenta Diaconi, etiam praesentibus Presbyteris in cano. In sancta distinct.92. qui canones a nobis non vt doctrinae Christianae firmamenta citantur: sed vt rei gestae historiae tantùm. Denique videntur Diaconi saepius, quam ipsi Presbyteri locum et vicem Metropolitanorum suorum repraesentasse, et obtinuisse in veteribus Synodis, ex eo quod Rom. Episcopi suas partes Diaconis saepiùs, quam Presbyteris demandarint: quod his facile carere non possent: illis autem Diaconis facilius possent Et hunc morem postea secuti sunt alii Episcopi, imprimisque Metropolitani, vt se ad Roman. Ecclesiae exemplar conformarent, adeo vt ex istis causis Presbyteris Diaconi censeri coeperint praeferendi, et iis esse gradu superiores maioresque. Quid igitur de hac re statuendum sit, quaeritur. Repo. Primùm omnes istas de proëdria et praeeminentia disputationes (quae sunt verissima dissidiorum semina, et damnabilis ambitionis effecta) de Ecclesia Dei, imprimis autem de eoetu praepositorum Ecclesiae (inter quos sunt et Presbyteri et Diaconi) tollen\pend
\section*{C APVT III. }
\marginpar{[ p.153 ]}\pstart das esse et facessere oportere, tanquam certissima et praesentissima Ecclesiae venena, vt docet Christus ipse Luc.22 vers. 25. tamen si quid de ea re dicendum est, tum ex veteri disciplina Ecclesiae: tum ratione muneris vtrorumque, certe Presbyteri sunt Diaconis priores. Itaque in veteribus Synodis veluti Nicaena prima cap. 14. Laodicaena cap 2o. Chartag. 4. cap.37.est statutum, vt ne Diaconi quidem inter Presbyteros sedeant, nimiram ne iis, aequales esse videantur. Id quod etiam iu tota distinct. 93. fuit positum et narratum. Ita nimirum ius istud canonicum perturbatum est, et sibi non constat. Nam quod ad praedicationem verbi Dei, quae illis aliquando concessa est, de eo quam recte factum sit, aut potius inepte, postea dicetur, sed primùm de personis ipsis agendum est.  \pend\pstart Eadem autem requirit Paulus in Diaconis, quae in Episcopis supra, siue ipsorum personae spectencur.iide raiiia. nue luuie ia ce teitimonia aliorum et infidelium. Addidit autem, ne sint Bilingues. Haec autem vox non tantùm impudentem memdacem, et impostorem significat, quo sensu Virgil. dixit.  \pend\pstart –Tyriósque Bilingues.  \pend\pstart sed etiam dissimulatorem simulatoremque. Denique qui id quod se praestare non posse intelligit, pollicetur: et qui cùm de promisso interpellatur, ita variat, mutat, et excusant promissionem a se factam, vt iam aliud videarur dicere, quam prius. Huic vitio simile est διδυχιας vitium, quod damnat. Iac. 4. vers. 8. cùm cor duplex gerimus. Hoc autem vitium inter reliqua maxime Diaconis fugiendum est, qui cùm saepe magna pauperum  \pend
\section*{AD I. PAVL. AD TIM. }
\marginpar{[ p.160 ]}\pstart importunitate vexentur et premantur, honeste quidem eos dimittere coguntur, et promissis implere, quae tamen postea praestare non possunt. Hoc autem offendiculum gignit, et a veritate Christiana, quae in omnibus nostris tum dictis tum leeus inesse debet, est alienissimum. Hoc idem vitium omnibus Christianis fugiendum est, et maxime Pastoribus.  \pend
\phantomsection
\addcontentsline{toc}{subsection}{\textit{ Tenentes mysterium fidei cum pura tonscientia. }}
\subsection*{\textit{ Tenentes mysterium fidei cum pura tonscientia. }}\pstart Α'κολόθησις est, postquam de moribus Diaconorum dixit, de eorundem doctrina, qualis esse debeat, disserit, nimirum eos fideles, probeque in fidei arcanis institutos esse oportere. In quo magnam et non tantùm vulgarem vel leuem in fide Christiana cognitionem eos habere vult Paulus, dum eos iuber, Non tantùm tenere fidem, id est, constanter et apprime eam didicisse et tueri sed etiam arcanum ipsum siue mysterium fidei ab iisdem perceptum esse postulat. Quo verbo altior quaedam cognitio significari videtur: quanquam salutis fidei, Christianaeque doctrinae quaelibet pars mysterium propter argumenti maiestatem. et rei ipsius difficultatem, nisi Dei spiritu simus muimdati, fere nominatur in scriptura. Refert Chrysost. hanc Pauli sententiam ad vers. 6. quasi haec ratio sit, cur Episcopum non oporteat esse nouitium: sed hic agit Paulus de Diaconis, non de Episcopis.  \pend\pstart Ergo quaeritur cur hoc requirat in Diaconis Paulus. Resp. Ratio est multiplex. Ac prima quidem quod cùm Diaconi sint non tantum pars quae\pend
\section*{CAPVT   III. }
\marginpar{[ p.161 ]}\pstart dam Ecclesiae: sed praecipua quaedam pars et nobilissima (sunt enim ex ηγκμένων et praepositorum numero) eos quoque oportet recte et pulchre in fidei Christianae doctrina, et arcano esse institutos et edoctos. Absurdum enim esset eos praefici Ecclesiae, qui verae et propriae ipsius Ecclesiae doctrinae ignari, et rudes essent. Altera ratio est, quod cùm pauperum et afflictorum qui sunt in Ecclesia curam habeant Diaconi, ad eorum quoque officium pertinet eos consolari, erigere, sustentare, et confirmare in fide. Quod non possunt, nisi doctrinam fidei apprime sciant et teneant. Quidam tertiam huius dicti rationem afferunt, sed falsam, nimirum quod ad Diaconorum officium pertineat concionari de Dei verbo, et illud praedicare, quemadmodum Philippus, et Stephanus fecerunt. Id quod etiam ex Iuone Carnutensi tradit Gratianus in cano. Perlectis 5Ad Diaconum distinct. 25. Haec autem sententia falsa est. Nam quod de stephano et Phmippo affertur eos praedicasse, verum est: sed non quatenus illi tantùm erant Diaconi, imo vero potius quatenus etiam Euangelista vterque fuit, vti tunc temporis varia Dei dona in vno et eodem fuerunt a Deo cumulata:  \pend\pstart Ergo admonet me hic locus, vt de Diaconorum officio aliquid dicam tum ex Dei verbo, tum ex ea, quae in veteribus Synodis extat, disciplina. Diaconorum in verbo Dei duplex genus esse videtur, vt colligitur ex Paulo Ramanor.12. vers.7. 8. Alii enim versabantur in Ipsius aerarii collectione, cura, procuratione tota. Hi sunt velut Oecouomi Ecclesiae. Tales fuerunt etiam Procuratores  \pend
\section*{162 AD I. PAVL. AD TIM. }\pstart dicti Hermi. liber  8.cap. 12. cuius generis exempla sunt Nehe. cap.13.vers.13. et hodie ii, qui hospitalium et πλοχεδοχέίων Prouisores appellantur. et in singularum partium, et elemosynarum, quae singulis pauperibus et egenis fiunt, distributione, et erogatione occupantur. Tales sunt qui dietim vel hebdomadatim eleemosynas cuique  pauperi fratri distribuunt, qui saepe habent tanquam subsidiarias operas et ministras adiunctas Diaconissas, de quibus postea. Hi autem septem primum sunt creati, et hic numerus retentus est Romae vt ait Sozo.libus  7.cap.19. postea ad imitationem Romanae Ecclesiae Neocaesariensi Synodo cap.14. sancitum est, vt etiam si non esset magna ciuitas, in qua esset Ecclesia, septem tamen Diaconi in ea crearentur. Est decretum in can. Diaconi septem distinct.93. Sed hoc decretum postea Chartag. Synodo prima fuit emendatum, cùm pro amplitudine Ecclesiae saepe plufes, saepe pauciores Dianconi requitantur. Itaque hoc totum iudicio cuiusque Ecclesiae est permittendum, non autem praecise definiendum. Nam et Romae postea, vt docet liber 2. Chronic. Martinus Pollonus auctus est Diaconorum numerus, et geminatus, vt pro septem hodie sint sexdecim Diaconi, ex quibus alii ipsius Papae ministerio et cubiculon sunt addicti, vt ait illeErgo Diaconorum verum munus est, tum Ærarium Ecclesiae colligere, curare et administrare, tum ex eo aerario pauperibus fratribus subuenire, tum Sanis, tum Ægrotis, tum etiam Ciuibus et Indigenis, tum Peregrinis, qui ad nos es aliis regionibus adueniunt. Ad eos autem per se et rasione muneris ipsius neque verbi Dei praedicatio  \pend
\section*{CAPVT  III. }
\marginpar{[ p.163 ]}\pstart pertinet: neque Sacramentorum administratio. Haec fuit vetustissima Diaconatus descriptio:quae tamen postea magna hominum temeritate immutata est, et cum Episcoporum munere pene confusa, quos tamen hic et Philip.I.aperte distinguit Paulus. 1  \pend\pstart Ac primùm Diaconorum muneris perturbatio facta est, cùm velut oculi Episcopi dici coeperunt, quoniam omnia totius Ecclesiae acta oberrantes perlustrabant, et ad Episcopum offendicula quae in Ecclesia nascebantur, deferebant, vt est in cano. Diaconi Ecclesiae distinct.93. Vnde apparet eos in locum Presbyterorum, ad quorum curam ea diligentia de inspiciendis hominum piorum moribus pertinet, temere fuisse substitutos. Ac quidam e Scholasticis magna ignorantia nullos praeter Episcopos et Diaconos in Ecclesia ordines constituunt, quod, excepto Episcopali officio reriqua inuma a Llaconis luicepta et obita fuisse dicunt. Quod est falsum, cùm etiam ex 1.Corint. 16. appareat saepe quosdam etiam extra ordinem ab Ecclesia delegatos, et electos ad curam et administrationem Ecclesiasticae pecuniae, quia Diaconi satis esse tanton oneri non poterant.  \pend\pstart Hoc Diacononrum munus iam tam late patens etiam amplificari postea coepnt jet pene cum Episcopi officio confundi, certe in multis adaequari. Nam in cano. Perlectis 5 Ad Diaconum pertinet distinct. 25. haec omnia Diaconis tanquam muneris eorum partes tribuuntur.  \pend\pstart 1 Assistere Sacerdotibus, id est, Episcopis.  \pend\pstart 2 Ministrare in omnibus, quae aguntur in Saeramentis Christi, veluti et in Baptismo, cheisman\pend
\section*{AD I. PAVL. AD TIM. }
\marginpar{[ p.164 ]}\pstart 6 Praedicare ad populum Euangelium, et Apostolicas Epistolas.  \pend\pstart Officium precandi.  \pend\pstart 8 Exhortatur et iubet populum clamare ad Deum, et ad eum cor et aures erigere et habere.  \pend\pstart Vnde hic ordo inter sacros praecipue refertur, tum propter materiam (aiunt Cañonistae) quam tractat (est autem Eucharistia) tum propter ipsam institutionis formam, quae in Papatu magnas caeremonias habet.  \pend\pstart Cùm vero tanta complecteretur hoc vnum Diaconorum munus, vt pene reliquas omnes dignitates Ecclesiasticas absorberet, contentiones variae natae sunt, et cogitatum est, de hoc officio intra certos quosdam cancellos et limites coercendo, Itaque Synodis constitutum est, vt Diaconi in posterum neque calicem benedicerent, neque corpus Christi ofrerrent ;iten neque praecinerent. populo, id est, cantum publicum moderarentur, nec ea tractarent et vsurparent, quae primis superiorioui jue ordinibus sunt decreta. Id quod ex conciliis Laodicensi et Chartaginensi apparet, et est relatum in cano. Non oportet et seq. dist. 93. cano. In sancta dist 92.  \pend\pstart Denique Durandus anxie versat hanc quaestionem. in liber 4. Sententiarum P. Lombardi dist. 13. quaest. 4. Vtrum ad Diaconum pertineat Eucharistiam populo ministrare, tam multa absur\pend
\section*{CAPVT  III. }
\marginpar{[ p.165 ]}\pstart da consequuntur, vbi semel a verbo Dei discessum est.  \pend\pstart Longe autem pessime et verbo Dei et munere Diaconi abutuntur, qui Diaconos tantùm titularios hodie habent, nulli rei, nisi pompae, et superstitioni inseruientes. Id quod fit in Papatu.  \pend
\phantomsection
\addcontentsline{toc}{subsection}{\textit{10 Atque hi etiam probentur priùs, deinde ministrent, si sint inculpati. }}
\subsection*{\textit{10 Atque hi etiam probentur priùs, deinde ministrent, si sint inculpati. }}\pstart Τάξις Ordine persequitur Paulus singula, quae ad legitimam Diaconorum electionem pertinent: et ordinem ipsum, qui seruandus est, praescribit. Ergo priùs eos examinari vult: postea vero eligi et ad officium admitti. Vnde qui sine examine: vel lusorio quodam examinis genere eliguntur, vti in Papatu, sunt omnino abdicandi, et reiiciendi. Ad quos vero hoc examen pertineat ad vniuersumne populum, an ad senatum, an ad vtrosque infra dicetur.  \pend
\phantomsection
\addcontentsline{toc}{subsection}{\textit{11 Vxores itidem honestas, non calumniosas, sobrias, fidas in omnibus. }}
\subsection*{\textit{11 Vxores itidem honestas, non calumniosas, sobrias, fidas in omnibus. }}
\phantomsection
\addcontentsline{toc}{subsection}{\textit{12 Diaconi sint vnius vxoris mariti, qui liberis recte prasint, et propriis domiDux. }}
\subsection*{\textit{12 Diaconi sint vnius vxoris mariti, qui liberis recte prasint, et propriis domiDux. }}\pstart Ista sunt superius explicata partim hoc ipso capite: partim cap.2. vbi de mulieribus Christianis disputauimus.  \pend
\phantomsection
\addcontentsline{toc}{subsection}{\textit{13 Nam qui bene ministrauerint, gradum sibi bonum acquirunt, et multam libertatem per fidem quae est in Christo Iesu. }}
\subsection*{\textit{13 Nam qui bene ministrauerint, gradum sibi bonum acquirunt, et multam libertatem per fidem quae est in Christo Iesu. }}
\section*{AD I. PAVL. AD TIM. }
\marginpar{[ p.166 ]}\pstart Α'ιτιλογία est ab effectu. Nec enim tales Diaconos frustra a se requiri docet Paulus: sed quis sit harum virtutum in illis fructus nunc ostendit, nimirum Gradus honorisque in Ecclesia maioris commendatio, et Bonae conscientiae, magnaeque apud fideles libertatis et authoritatis comparatio et testificatio. Hoc primum opponere videtur Paulus contumeliis, calumniis, conuitiis et probris, quae vulgo patiuntur ii, qui sancte et syncere Deo seruiunt, et ipsius Ecclesiae ministrant, qui prae caeteris apud profanos homines probris conspui solent. Sed et hic idem effectus Diaconatum tanquam praecipuum quoddam in Dei Ecclesia munus commendat. Denique Diaconos iplos consolatur, et docet gregem ipsum, et singulos pios et fideles, quantum huiusmodi personis, qui syncere hoc munere vel funguntur, vel functi sunt, debeant.  \pend
\phantomsection
\addcontentsline{toc}{subsection}{\textit{14 Haec tibi scribo, sperans forevt celeriùs ad te veniam. }}
\subsection*{\textit{14 Haec tibi scribo, sperans forevt celeriùs ad te veniam. }}\pstart Παράκλνσις est, Consolatur enim Timotheum, et admonet cur tam minutim singula sit persecutus. Denique spem facit proximi aduentus sui. Primùm ipsius Timothei causa, ne cùm ad absentes literae scribantur, existimet velle Paulum diutius abesse, quam Ecclesiae vtilitas postulet. Ad tempus enim reliquerat Timotheum Ephesi Paulus, vt ipse docuit cap.1. donec in Macedoniam ad constituendas Ecclesias proficisceretur. Itaque  ne sui promissi oblitus censeretur, aut hoc falso praetextu Timotheum reliquisse, hoc addidit. Si\pend
\section*{CAPVT  III. }
\marginpar{[ p.167 ]}\pstart mile quiddam apparet 2. Corinh.1. Deinde hoc ipsum Timothei excitandi causa videri potest additum, quanquam ipse per se diligens esset, sed vt alacriore animo in illud opus incumberet, in quo suae diligentiae testem oculatum Paulum esset posthac habiturus.  \pend\pstart Sed et aliorum quoque causa videtur hoc adiectum, vt si qui essent peruicaces et contumaces, metu aduentus Pauli comprimerentur. Hinc auceni taun sonestuud de moeroujs apparet, qui ne minimum quidem tempus differre potuit, quin ea quae videbantur egere correctione quadam, emendaret.  \pend\pstart Dicit autem, sperans, Neque enim nostrae vitae sumus nos ipsi arbitri: sed Dominus, et, vt monet Facobus 4 vers.15. ipse de nobis statuit, vt ex eius solo nutu pendere toti debeamus. Rediit autem Ephesum Paulus, vt apparet ex Act. 18. post hanc Epistolam scriptam.  \pend
\phantomsection
\addcontentsline{toc}{subsection}{\textit{15 Quod si tardiùs venero, vt noris quomodo oporteat in domo Dei versari, quae est Ecclesia Dei viui, columna et stabilimentum veritatis. }}
\subsection*{\textit{15 Quod si tardiùs venero, vt noris quomodo oporteat in domo Dei versari, quae est Ecclesia Dei viui, columna et stabilimentum veritatis. }}\pstart Occupatio est, quae eum commendatione ministerii Euangelici est coniuncta. Monet enim Paulus si tardiùs venerit, ecquod muneris sui recte gerendi remedium habeat Timotheus, nenpe haec ipsa praecepta, quae continentur hac Epistola, tum in superioribus capitibus: tum etiam in sequentibus. Itaque parasceue etiam est ad sequentes praeceptiones hie locus.  \pend
\section*{AD I. PAVL. AD TIM. }
\marginpar{[ p.168 ]}\pstart Ex hoc autem locoapertissime demonstratur, quis sit eorum praeceptorum, quae hic de administrandae Ecclesiae ratione traduntur, vsus. Quidam enim existimant (quos ipse aliquando audiui) quae hic de disciplina Ecclesiastica praescribuntur a Paulo, veluti de Episcoporum, et Diaconorum electione, et munere: de censura morum quae fit a pastoribus, denique quae regulae politiae Ecclesiasticae cap.5. et tota hac Epistola traduntur, esse tantùm singulares quasdam praeceptiones et particulares, quae tunc temporis vsum suum haberent, sed quae generalem et perpetuam politiae Ecclesiae normam et regulam minime continent, minimeque praescribant. Itaque non debere hodie ad hanc formam a Paulo descriptam politiam et rationem Ecclesiae Christianae administrandae vel reuocari, vel examinari. At contra hoc loco docet, affirmatque diserte Paulus, se Timotheo ea praescribere, quae in domo Dei obtinere, et institui debeant, vt pulchre illa et regatur, et administretur. Neque vero tempus certum, aut etiam locum aliquem designat Paulus, cuius ratione vtiles sint hae praeceptiones, et necessariae: sed generalis est affirmatio Pauli, vt scias, ait, quomodo te regere debeas in domo Dei, id est, Ecclesia, non dixit Ephesi tamtùm, aut in Asia: non dixit, hoc tempore nempe, quo sunt Magistratus infideles et idololatrae: sed doceo (ait Paulus) quomodo te, id est, omnes veros Episcopos (est enim in persona vnius Timothei omnium bonorum Episcoporum officium descriptum) se gerere oporteat in Dei Ecclesia, qualis illa Dei domus est, et ex Deivoluntate recte administrar i debet, in  \pend
\section*{CAPVT  III. }
\marginpar{[ p.169 ]}\pstart qua omnia ἀτάκτως fieri et tempore et loco necesse est. At iidem excipiunt disciplinam Ecclesiasticam esse mutabilem. Resp. Distinguendum esse. Quedam enim esse in ea Essentialia, aio, quaedam Accidentalia. Essentialia sunt omnia, quae et de electione et de munere personarum Ecclesiasticarum hic traduntur, quia cùm haec sint publica Ecclesiae munera, electio in iis omnino necessaria est. Et cùm haec sint munera, officia et munia quaedam habent, quibus definiuntur, vti pastoris Ecclesiastici est assidue verbo Dei pascere gregem suum: super eo vigilare, et morum vitaeque illius habere curam. Presbyterorum et Seniorum est attendere gregi Domini, ne haereses vel vitia obrepant in Ecclesia, siue aperte, siue sensim.  \pend\pstart Accidentalia sunt, quae ad particularem earum rerum modum seruandum, formamque aliquam constituendam, et cuiusque populi commoditatis rationem habendam pertinent: veluti, quoties in hebdomade sit Episcopo concionandum: quibus diebus, qua hora, quot locis, et caetera quae sunt huiusmodi. Sed quae praescribit hic Paulus sunt essentialia. Quare haec a nulla Ecclesia, sine piaculo praetermitti possunt: nisi domum Dei vere domum confusionis facere velimus. Nec vlla vllius regni vel reipublicae commoditas, vtilitasque aliud suadere, aut inducere nos debet, vt ab his Pauli, id est, Dei Spiritus praeceptis recedamus: aliaque vel Episcoporum, vel Diaconorum munera faciamus. Sed nec Synodi quidem ipsae sunt audiendae, si ab his Pauli regulis discedunt, quia a Dei verbo in Ecclesia Dei regenda, illiusque politia constituenda discedunt. Quod est ne\pend
\section*{AD I. PAVL. AD TIM. }
\marginpar{[ p.114 ]}\pstart Dome Dei) Commendat hic non tantùm Eeclesiam Paulus: sed ipsum potius ministerium Euangelicum a re, et obiecto, circa quod versatur, Nimirum ab Ecclesia, cuius aedificandae, administrandae, et conseruandae gratia constitutum est. Magnam enim esse eam dignitatem oportet, quae domo Dei praeest, eique gubernandae praeficitur. Illa vero dignitas et cura est ministerium verbi Dei. Excitat vero hoc ipsum etiam pastores pios vt de se suaque diligentia cogitent, et cui muneri sint a Deo praepositi serio meditentur, et saepe apud se proponant. Similis locus Actor.20.V.28. Magna vero commendatio Ecclesiae (cuius pascendae gratia Episcopi et Presbyteri eliguntur) continetur paucis his Pauli verbis, columna est, ait et stabilimentum veritatis Ecclesia. Hic enim doctr inam Ecclesiae praefert omnibus aliis, quia illae non sunt vel veritas, vel illa salutaris veritas, quam vnam pios et Christianos homines sectari oportet. Itaque longe vniuersae Philosophiae praeferenda est haec doctrina, quae in Dei Ecclesia docetuiiee prauicutur, quaa haec est sola veritatis nomine digna. Etsi vero quaedam vera a Philosophia traduntur, tamen illa veritas non est (vt idem Paulus ad Titum 1. vers.1. ait) κατ' αἀσέβειαν, sed magna ex parte hominibus inutilis, vel certe cum hac nullo modo conferenda, quia haec ad animorum vitam aeternam pertinet. Deinde, vt scribit Chrysostomus, haec vox ve\pend\pstart ritas discrimen notat inter doctrinam eam, quae olim Iudaeis sub lege annuntiabatur: et eam, quae nunc Christianis annuntiatur. Illi enim vmbras tantùm, id est, figuras et caeremonias habuere:nos  \pend
\section*{CAPVT  III }
\marginpar{[ p.171 ]}\pstart autem ipsam veritatem et corpus in Christo, vt docet Paulus Coloss.2.V.17. Ex quo fit, vt Ioann. I.verl.lj. dicatur per ivioien data Lex-per Christum vero veritas et gratia.  \pend\pstart Huius autem tam salutaris, et sacrae veritatis custos, conseruatrix, et publicus velut tabellio apud quem deposita est, est Ecclesia illa Dei viui, nempe per ordinariam et puram verbi Dei praedicationem, quae ab veris Episcopis et pastoribus facta est. Nam, vt ait Paulus Romanor.3. vers.2. Credita sunt (πςεύθησαν) Dei Ecclesiae, non autem Sathanae Synagogae, non Antichristiano coetui: sed Christianae Ecclesiae diuina oracula et eloquia, quia in hoc tantùm vno hominum coetu illa explicantur, docentur et conseruantur, per ordinarium verbi Dei ministerium, quod in Ecclesia viget. Vnde idem Paul. Rom.9. vers.4 Legem, et foedera, et eorum foederum tabulas fuisse apud Israëlitas, ait, quia et cum iis pactum erat foedus Dei, et eius foederis tabulae apud eos solos asseruabantur, et custodiebantur, velut apud fidelissimum et publicum tabellionem manent acta et instrumenta contractuum. Hoc sensu vocat Paulus veram et Euangelicam Dei Ecclesiam, Columnam, et Stabilimentum veritatis, non quidem tanquam istius veritatis Ecclesia sit causa, sed tantùm custos, Conseruatrix, et Propagatrix. Veritatis enim, quae in Ecclesia Dei praedicatur, causa est Deus ipse, cuius est haec veritas vox et verbum. Ergo ratione perpetui et ordinarii ministerii et praedicationis puri verbi Dei, quae in vera Ecclesia locum habere debet, dicitur Ecclesia Columna veritatis: non ratione reipsius, id est,  \pend
\section*{AD I. PAVL. AD TIM. }
\marginpar{[ p.172 ]}\pstart Ecclesiae ipsius, quasi certitudinem suam ab Ecclesia, id est, hominum iudicio et approbatione habeat sacrosancta Dei veritas. Neque enim ab hominibus pendet diuinae vocis commendatio, aut certitudo: sed in diuina voce hominum ipsoruin demntio d rententia de veritate aut falsitatate alicuius doctrinae fundata esse debet Ioan.5, vers.39.  \pend\pstart Ergo duobus modis Paulus hic nobis describit Ecclesiam I Cùm domum Dei appellat. 2 Cùm columnam veritatis. Ac domus quidem Dei vocatur Ecclesia, quia est piorum et fidelium hominum coetus Matth.21. vers.33. Certe si quo in alio mundi loco, maxime in eo coetu illustria specialiaque praesentiae Dei signa et testimonia extant. Ergo dicebatur Deus habit are in suo tabernacule, et in Sione, propter suae praesentiae et maiestatis signa peculiaria, quae in eo loco posita a Deo ipso erant Psal.76. quia illic et verbum Dei purum audiebatur, et cultus ab eo praeceptus exercebatur. Non quod Deus corporaliter illic versaretur: non quod extra eum coetum nusquani arior enet Deus, qui semper vbique est: non quod etiam Deus eodem, quo ipsa Ecclesia, loco concludatur: sed quod cùm se maxime homimo uo communicet, per verbum suum, per Sacramenta, et per cultum: hos omnes modos sola Dei Ecclesia ex omnious hominum coetibus habet, quae retinet  \pend\pstart Deinde quod Ecclesia dicitur Columna veritatis, ostendit hanc esse veram et vnicam Ecclesiae verae notam, nempe veritatem et verbum ipsum Dei. Huic enim vni fundamento innititur. Nam  \pend
\section*{CAPVT  III. }
\marginpar{[ p.172 ]}\pstart oues meae vocem meam audiunt, ait Christus Ioa.10. adeo vt sublata et abolita illa veritate, vel mutata in mendacium, desinat quoque esse illic Ecclesia, et vbi illa veritas annuntiatur, illic quoque nascitur Ecclesia. Ergo non est Ecclesia vel locis, vel cercio perionis aiipatajce aitrieea, sed doctrinae verbi Dei, vt etiam desiniuit Synodus Gangrensis de sanctitate agens.  \pend\pstart Abusi sunt igitur perspicue hoc loco Papistae dum hinc colligunt, verbi Dei certitudinem ex Ecelesiae approbatione et hominum suffragio pendere. Nam falsum est. Subiiceretur enim Deus hominibus, id est, creator creaturae, et omnipotens vermibus terrae: et possent homines suam Deo veritatem adimere si vellent, vel illi conseruare, quae tamen penes eum vnum residet, neque ex hominum iudicio et authoritate pendet, vel immutari potest Iacobus 1.V.17. Praeclare enim Iren. libus 3. cap.1. et 11. Euangelium est columen fidei nostrae, non hominum authoritas. Et quod contra fundamenta. Manichaeo. epistol. cap 5.scribit August. Ego vero Euangelium non crederem, nisi me Catholicae Ecclesiae commoueret authoritas, dictum est aduersus eos qui pro animi libito, ex sola suorum praeceptorum et quide recentium approbatione, hos sacrae scripturae libros recipiebant vt authenticos: illos autem, vt falsatos, reiiciebant, quod faciebant Manichaei, quibus opponit Augustinus maiorem et fideliorem hominum approbationem, nempe diuturnum totius Ecclesiae consensum.  \pend\pstart Denique quod iidem Papistae ex hoc loco colligunt, Ecclesiam errare non posse, falsum est, si  \pend
\section*{AD I. PAVL. AD TIM. }
\marginpar{[ p.174 ]}\pstart ex eorum sensu accipiatur. Nam Sacerdotum, qui Christum damnarunt, collegium ex eorum definitione erit Ecclesia. Atque hi tamen grauissime errarunt, et veram atque vnicam illam Dei veritatem, id est, Christum ipsum reiecerunt. Deinde cùm ipsi illi Papistae non sint Ecclesia Dei, frustra ad se transferunt hanc laudem. Fatemur quidem Ecclesiam, quatenus vnicum et purum illud Dei verbum pro fidei regula amplectitur, retinet, et sequitur, falli non posse: sed quia hominibus constat Ecclesia, in quibus saepe maior pars meliorem vincit, quique non semper ad solum Dei verbum collimant, aut attendunt: vel qui purum Dei verbum etiam propositum non retinent, idcirco recte diestur etiai fuiumour coetus, qui Dei Ecclesia dicuntur in his terris, errare posse, et saepe etiam grauissime errasse, vt apparet in secunda Synodo Nicaena, vbi idololatria constituta est. 16 Et sine controuersia magnum est pittati mysitrium: Heus tonspicuus factus est in carne, iustificatus est in spiritu, conspectus est ab Angelis, praedicatus est Gentibus, fides illi habita est in mundo, sursum receptus est in gloriam. Altera ministerii Euangelici commendatio a subiecto et materia, circa quam versatur. Hanc autem, vocat mysterium magnum, quod explicat quale, quantumque sit, quo magis ad suum rnunus attendant praepositi Ecclesiae, illudque re3 2 uerenter colant, ñon autem perfanctorie. Similis  \pend
\phantomsection
\addcontentsline{toc}{subsection}{\textit{}}
\subsection*{\textit{}}
\section*{CAPVT  III. }
\marginpar{[ p.175 ]}\pstart etiam est admonitio Pauli Actor.20. vers.28. Er go falsa est eorum sententia, qui putant hic a Paulo tantùm conferri legis caeremonias cum doctrina Euangelii, tum quod maior sit huius, quam illarum vis: tum quod non eandem iam requirat in futuro pastore Ecclesiae Euangelicae, quae in lege sua Dominus statuit, et praecepit obseruari a summo Sacerdote: Quae sententia est Chrysostom. et Erasmi. Locus hic est sane singularis, vbi auguste, et breuiter tamen praecipua doctrinae Christianae capita, et summa Dei in nos beneficia complexus est Paulus. Atque huiusmodi locos annotare operaepretium est, vt summam quandam fidei nostrae memoria conseruemus. Similis etiam est locus alius apud eundem Paul.1. Corint.2.v.2. apud Lucam Luc.24.v.47, et Hebr.13.v.8. et seq. Vt autem ordine et methodo quadam hic locus explicetur, est in suas partes distribuendus. Sunt autem duae. Primum enim hanc Dei veritatem explicat in genere ipso, cuius conseruandae et retinendae inter homines gratia ministerium Euangelicum constitutum est. Vocat igitur eam 1 Mysterium, id est, arcanum. 2 Magnum, idqi 3 Siue controuersia tale esse affirmat. Mysterii vox hic non significat Sacramentum, quo sensu Coenam Domini, et Baptismus Sacramenta dicuntur: sed doctrinam arcanam, augustam, humani ingenii viribus et industria incomprehensibilem et ex sola Dei gratia, et reuelatione notam, qualis profecto est doctrina fidei, et salutis nostrae, vti docet Paulus Ephes.3. vers.4,5. Quanquam hoc ipso sensu, quo vsurpat Paulus mysterium hoc loco, veteres Patres vocarunt doctrinam E\pend
\section*{AD I. PAVL. AD TIM. }
\marginpar{[ p.170 ]}\pstart uangelii, Sacramentum fidei, et ipsam fidem Sacramentum: delectati nimirum voce ipsa Sacramenti, ranquam augusta, d doctrinae huius maiestati consentanea.  \pend\pstart Auget vero nanc appellationem a comparatione Magnum enim vocat. Nam inter arcana ipla Deranud anlo maius eit, et augustius. Hoc autem est summum, et ipsa mundi creatione longe etiam maius, ac praestantius.  \pend\pstart Addit sine controuersia, id est, prout est omnibus in confesso. Amplificatio item est a consensu omnium piorum, et vere de pietate sentientium. Sed obstat quod hoc mysterium multi negant vnquam extitisse: alii huic non assentiuntur. Aliis vero est offendiculum vti Iudaeis, aliis stultitia vt Gentibus. Respond. Istud a Paulo dici non ex hominum infidelium iudicio, sed ex fidelium tantùm et piorum sensu. Praeclara igitur est illa Bernardi sententia sermo.45. Cantic. ab hoc loco non aliena. Quam mihi decorus es Domine mi, in ipsa huius tui depositione decoris. Etenim vbi te exinanisti, ibi pietas magis enituit: vbi charitas plus effulsit, ibi amplius gratia radiauit. Quam spectabilis et stupendis etiam virtutibus supernis in conceptu de Spiritu, in ortu de virgine, in vitae innocentia, in doctrinae fluente, in coruscationibus miraculorum: in reuelationibus Sacramentorum: Quam denique rutilans sol iustitiae post occasum, de corde terrae resurgens, quam formosus in stola tua: demum rex gloriae in alta coelorum te recipis. Sequitur.  \pend\pstart Altera pars huius versiculi, quae rei ipsius, id est, veritatis et doctrinae Euangelieae summam  \pend
\section*{CAPVT  III. }
\marginpar{[ p.2 ]}\pstart continet velut et differentiam. Ex quo ipso magis admirabilis, inscrutabilis, et augusta apparet haec veritas, eóque  ipso excellentius est Euangelicum ministerium, per quod illa dispensatur et annuntiatur. Habet vero haec explicatio quatuor partes. Proponit enim 1 Rem ipsam. i. Incarnationem Christi, in qua inest etiam Approbatio ipsius deitatis, 2 Testimonium huius mysterii. 3 Praedicationem et receptionem admirabilem eiusdem. 4 Conuicroneul denqyneucuiq marertatis Christi, in qua nunc est. Quae singula summam consolationem nobis afferunt, et miram Dei tum potentiam, tum etiam sapientiam ostendunt.  \pend\pstart Primùm igitur Deus dicitur conspicuus factus esse in carne (nec enim quae sequuntur, referenda sunt ad mysterium, sed ad Christum: etsi et vetus interpres, et Erasmus, et alicubi Fulgentius ita accipiunt, et corrupte in quibusdam exemplaribus olim legebatur o, proὸs, quod ôs etiam ex θεὸς factum est, vt fidei nostrae de deitate Christi doctrina certissima, quae hic continetur, conuelleretur.) Hic autem Deus, qui est conspicuus nobis factus, est Filius Dei et λόγος, non autem ipse Pater Hebr.1. Ioann.1.qui Filius vti Deus est Prouerbus 8. ita Deus, et magnus quidem, appellatur Rom.9. vers.5. Tit.2. vers. 13. quantumuis id negent Arriani, qui refutantur Ioann.1o.v.3o. Hic igitur Deus Filius conspicuus nobis factus est, non tantùm aliquo signo praesentiae suae hominibus dato, vt cùm in rubo ardenti apparuit: non etiam assumpta tantùm ad tempus humana forma, vt cùm apparuit Abrahamo sedenti ad quercum Mamre. Genes.8. et Exod.3. Nec quale illud et  \pend
\section*{AD I. PAVL. AD TIM. }
\marginpar{[ p.176 ]}\pstart Deus humana lustrn sub imagine terras. Sed per vnionem hypostaticam per quam vniuit sibi, id est diuinae suae naturae humanam naturam, vt fieret vnum et idem ύφιςάμενον, qui esset verus Deus, et verus homo. Diuina quidem natura non est humana dicenda: sed quia illa vtraque natura diuina, nempe et humana, vnica tantùm persona est, idcirco idem dicitur et Deus et homo, et homo ille qui apparuit nobis dicitur esse Deus. Quare cùm fuerit Deus, Deus in carne nobis est manifestatus.  \pend\pstart Modum autem addidit Paulus per quem id factum est. in carne, id est, per humanam naturam, quam verbum assumpsit Ioan.1.vers.14. Hebr. 2. vers.16. Quare idem ipse Christus air, Qui me videt, videt et Patrem Ioann.14.v.9. Hoc ipsum etiam explicat Ioan.1. Epist.1.vers.1.2.3. Quod enim erat per se aeternum, inuisibile et a principio assumpsit eam naturam, quae videri, tangi, palpariue potuit, per quam factus est Deus nobis quodanmodo propter vnionem hypostaticam conspicuus et visibilis, non conuersione diuinae naturae: sed assumptione humanae, et vnione cum diuina, quae vtraque in vnam personam Christi coaluit. Ergo carnis vox hic non vitiositatem illam (quam natiuam habemus) significat: sed humanam naturam et essentiam totam et veram, quatenus corpore et animo illa constat: non autem quae tantùm corpore constat, ne in errorem Apollinaristarum incidamus. 11  \pend\pstart Assumpta vero est haec caro siue humana natura in vnionem verbi hypostaticam, vt in perpetuum illi maneret in posterum coniuncta. Quare  \pend
\section*{CAPVT  III. }\pstart neque per se suum esse habet: neque extra Dei verbum, neque seorsim et separatim est illa humanitas: sed verbi Dei humanitas, et in Dei tantùm verbo habens suum existere, non per se et seorsim, vt humanitas Petri et Pauli et mea. Quare hic reuocamur ad mysterium illud augustum incarnationis Christi, per quam σωματικῶς, id est, υὐ̔ποςατικῶς et ad ἐνὸς ὔφιςαμέντν ἐπόρασιν vtraque illa natura coniuncta est, inconfusas tamen proprietates suas retinens et conseruans. Quare duae naturae sunt hic nobis propositae Deus et caro, id est, Deus verus et verus homo, et vtriusque vnio et coniunctio tamen, cùm Deus manifestatus per carnem dicitur. Si enim maneret separatum verbum ab humana natura et carne, vel si caro seorsim a verbo existeret, tanquam duae quaedam diuersae personae (quod sensit impius Nestorius) mentiretur, et falso diceret Paulus, Deum in carne nobis esse manifestatum. Si esset confusa et absorpta caro, id est, humana natura assumpta a diuina assumente, vt voluit blasphemus ille Eutyches, falso quoque distingueretur, ac commemoraretur hic vtraque natura. Si vnica tantùm in Christo natura, falso etiam diceretur idem Christus Deus et caro. Hoc igitur stupendum est quod cùm toto genere Deus et caro dissideant, vtraque tamen natura in vnicam Christi personam coaluerit, et vnica sit per incarnationem. Sed et illud quoque est hoc loco obseruandum, Paulum vti ea ratione dicendi, per quam de naturis loquitur in concreto, vt Dialecticorum scholae vocant, Deus, inquiens, est manifestatus: non  \pend
\section*{AD I. PAVL. AD TIM. }
\marginpar{[ p.180 ]}\pstart Nec enim vere in abstracto diceretur, etsi vere pronuntiaretur in concreto, quia concretionis vocabula totam Christi personam designant vtraque natura constantem. Abstracta autem vocabula naturas ipsas seorsim et per se consideratas.  \pend\pstart Justificatus Spiritu) Est adiunctio et ἔνδεξις et velut superioris mysterii amplificatio, quod in illa etiam vili natura, quam pro nobis assumpsit Christus (vt docet Paulus Philip.2.) eluxit, refulsit, et apparuit manifestissime etiam diuina Christi natura, quae hic Spiritus appellatur. Itaque conuenit hic locus cum Rom. cap.1: vers. 4. vbi vna quaedam species istius diuinae virtutis recensetur, nempe Resurrectio. Sic etiam 1. Pet.3.vers.18. vbi haec vis spiritus appellatur. Ergo iustificare est non iustum declarare et probum virum: sed verissime ostendere et inconcussis argumentis, quaeque refelli non possunt, demonstrare aliquid, probare, et deprehendere, Sic Matth. 11.vers.19. Quae enim exacte perfecta sunt et exquisite siue ad amussim probata, dicuntur iustificata esse etiam apud nos Gallos. Vnde vulgo dicimus, Iustifier vn arrest de conte. Iustitia enim appellatur omnis exacta ratio, quae obseruata est in re aliqua demonstranda. Hoc Pauli dictum ostendit potentiani enffernom pompis et apparatu quodam externo apparuisse: sed in spiritu, id est, in virtute diuina, miraculis, operibus, doctrina coelesti, et incomprehensibili ipsius erga homines efficacitate. Deinde tollit Paulus hac voce omnem illam haesitationem animi, quae potest inhaerere nobis, propter humilem Christi hominis statum, qui fuit seruilis, abiectus, et contemptus.  \pend
\section*{CAPVT  III }
\marginpar{[ p.181 ]}\pstart V'isus siue conspectus ab Angelis.) Haec est secunda pars, quae continet testimonium Angelorum de praecedenti mysterio. Auget autem haec circunstantia hoc miraculum, et Dei potentiam, quae in hoc fidei nostrae mysterio eluxit: quod non modo homines ad tantae rei dignitatem, id est, Christum hominem factum esse obstupuerunt: sed Angeli etiam ipsi illud idem videre gestierunt, et videntes delectati sunt, Deumque laudarunt, suntque admirati. Haec enim omnia complectitur haec Pauli sententia de Angelis, vt apparet ex Luc.2 vers.14. Viderunt igitur hoc Angeli, quasi rem sibi nouam et admirabilem. Etsi enim illi saluandum esse aliquando mundum sciebant, et Christum hominem futurum, et occidendum praedixerat etiam ipse Gabriel Angelus Danie.9, singula tamen mysteria, et Dei miracula, quae in hac incarnatione, passione, et resurrectione Christi facta sunt, et futura erant, ipsumque adeo modum, priurquam illa nerent, ignorabant. vnde Angelis ipsis admirabile fuit hoc tantum Dei opus, quod Euangelico tamen ministerio praedicatur, et dilucide explicatur. Hos autem Angelos testes hic produci huius tanti mysterii certum est, quo certior sit fides nostra, et augustior.  \pend\pstart Praedicatus est Gentibus) Tertia pars huius explicationis, et mysterii, in qua non minus summum Dei arcanum, quam in superioribus, continetur. Est autem manifestatio tanti mysterii facta hominibus, quae et ipsa plane est miraculosa. Nam facta est et excepta a Gentibus, quae erant antea prorsus infideles et a Deo omnino alienae.  \pend\pstart Quarta vero pars nimirum de gloria Christi,  \pend
\section*{AD I. PAVL. AD TIM. }
\marginpar{[ p.180 ]}\pstart explicatur fere in articulo, Sedet ad dexteram Dei Patris quae cùm sit ex Catechismis notissima, et ab aliis copiosissime explicata, maxime propter infoelices istas de Coena Domini controuersias, et de vbiquitate corporis Christi, idcirco nos eam hic omittimus. Studemus enim breuitati.  \pend
\endnumbering\beginnumbering\section{CAP. IIII.}
\phantomsection
\addcontentsline{toc}{subsection}{\textit{\huge\textbf{S}\normalsize PIRITVS autem diserte dicit fore ut postertorions temporibus desciscant quidam a fide, spiritibus deceptoribus attendentes, ac doctrinis daemoniorum. }}
\subsection*{\textit{\huge\textbf{S}\normalsize PIRITVS autem diserte dicit fore ut postertorions temporibus desciscant quidam a fide, spiritibus deceptoribus attendentes, ac doctrinis daemoniorum. }}\pstart Transitio est, per quam a contrario superiorem doctrinam non tantùm illustrat: sed etiam confirmat et commendat, ne verbi Dei praecones segnes sint in ea tuenda et conseruanda, cùm aliquando futurum esse audiant, vt misere corrumpatur et labefactetur, pessimaeque superstitiones pro ea in Ecclesiam inducantur. Similis locus et admonitio fit a Paulo Actor.2o. vers29. a Ioanne 1. Epist.2.vers.19. a Petro 2. Epist.3.vers3. et a Iuda Thaddaeo Apostolo, vt apparet ex ipsius Epist. vers.4. Neque eo magis terreri debemus: quod fit a pluribus ista admonitio: sed potiùs eo diligentiùs cauere nobis, et prospicere ne decipiamur: minusque offendamur oportet, cùm huiusmodi haereses obrepunt, quas fore iampridem spiritus Dei praedixerat.  \pend\pstart Duo vero quaedam praecipue sunt hic obser\pend
\section*{CAPVT . IIII. }
\marginpar{[ p.181 ]}\pstart uanda, 1 Quanti momenti sit haec praedictio. 2 Quid contineat. Magni enim esse momenti ex eo apparet, quod ρντως, id est diserte et praecise a Dei spiritu praedicta fuit, et reuelata. Etsi enim omnium eorum, quae Paulus annuntiauit, Spiritus sanctus est author: hanc tamen tanquam specialem quandam prophetiam et admonitionem vtilem futuram Ecclesiae manifestauit, et obseruari voluit Spiritus sanctus. Quaeri vero potest, vbi hoc praedixerit Spiritus. Respond. Alium librum, et alium locum, quam hunc ipsum, non esse quaerendum. Satis enim est hanc fuisse prophetiam, quae Paulo et Ecclesiae reuelata fuerit. Nec enim fuit Paulus mendax: sed Dei propheta: et de his rebus, vti de aliis multis, Dei voluntatem compererat quemadmodum apparet ex 2. Tess.3.  \pend\pstart Non posset fortasse prima fronte haec doctrina, quae hic a Paulo damnatur, vti Daemoniaca, et hypocrisis plena, tanti esse momenti videri, cùm de cibis tantùm agat, et matrimoniis. Respondeo tamen cùm hac ratione verae sanctitatis definitio, et cultus Dei natura prorsus peruersa fuerit, et in meram non modo superstitionem, sed impietatem vel hypocrisin conuersa, merito detestandam et periculi plenam dici et damnari. Cùm enim huiusmodi sententiae inter Christianos obtinent, iam rebus externis alligatur vera sanctitas. nimirum ciborum et coniugii continentiae. Deinde hac ratione censetur Deus melius coli. Demum hac fenestra semel patefacta certatim homines ex ingenii sui vanitate in ἐθελοSoησκέας ruunt, et Deum ex sui animi commentis colunt: non autem amplius ex sacrosancto ipsius ver\pend
\section*{AD I. PAVL. AD TIM. }
\marginpar{[ p.1o4 ]}\pstart bo: quae res omnis impietatis proculdubio fons est ac origo Coloss.2.  \pend\pstart Magnopere igitur huiusmodi doctrinae venenum nobis cauendum esse satis docent haec Pauli verba, per quae eam vocat, Defectionem a fide, DaeIomorumjuocerinan, vtretiain facobus 3.V.15. docet eam esse plenam hypocriseos, et Conscientiam cauterio infectam afferre.  \pend\pstart Alterum caput huius disputationis rem ipsam explicat, et Causam huius doctrinae, et Tempus et Doctrinam ipsam siue illius dogmata. Causa vero huius doctrinae est multiplex, quaedam Superior, Spiritus deceptores, et quaedam inferior Defectio a fide, et Hypocrisis, malaque  hominum, tum Docentium, tum Discentium conscientia. Vocat autem hic Paulus spiritus non homines quidem horum dogmatum et errorum praecones et ministros: sed Sathanam ipsum, et illius Angelos, quos postea Daemones nominat, quia illi sunt et flabella et authores verae doctrinae corrumpendae, ne homines ad salutem perueniant. Praeterea qui haec venena spargunt, etiam spiritum iactant, sed ille est spiritus vertiginis et erroris, cui Dominus ad puniendum nominis sui contemptum dat efficaciam. Hi igitur spiritus dicuntur πάνοι, id est, non tantùm deceptores, sed etiam seductores, quia nos a recta via, cui iam insistebamus, et in qua iam incedebamus, deflectunt, id est, a verbo Dei abducunt, cui vni tanquam soli veritatis columnae est acquiescendum: quod plus est, quam nos adhuc veritatis ignaros decipere, et in errore confirmare.  \pend\pstart Ex parte autem auditorum causa quoque no\pend
\section*{CAPVT  IIII. }
\marginpar{[ p.185 ]}\pstart tata est, nempe defectio a fide, et a sana doctrina. Doctores igitur appellat δαδολόγες: Eorum enim doctrina est falsa, et contraria velae et sanae doctrine Discipulos autem vocat ἀποςάτας et homines leues, qui veritatem iam perceptam facile deserunt, qua sunt ingenii vanitate. Omnis ergo mali istius et ruinae fons et origo ex eo est, quod homines a verbo Dei desciscunt, et deficiunt ad vana animi sui figmenta scilicet, et ad impias soggestiones Sathanae, quibus attendunt. Dixit enim attendunt. Est autem attendere diligentem animi curam et operam ponere, et non perfunctorie vel leuiter aliquid degustare: quod statim deseras: sed magno studio et diligenti nosse, perquirere, et ediscere. Sunt autem homines nouarum rerum et blasphemarum auidissimi natura: sed imprimis ii, qui huiusmodi doctrinae attendunt. Tempus quoque designauit Paulus, non quidem tempora omnium nouissima: sed tantùm posteriora. Quod notandum est, vt eo magis intelligamus, vtrum istum Pauli dictum ad Papistas pertinere possit, de quo seq. versiculo agetur.  \pend\pstart 2 Per hypocrisin falsiloquorum, quorum conscientia cauterio resecta est.  \pend\pstart Εʹπεξεσγασία, Operose enim et diligenter Paulus singula notat, quae in ipsis haereseon authoribus sunt obseruanda, vt eae sint nobis magis detestabiles, et earum çausae magis turpes et foedae appa reant. Itaque obseruanda sunt encomia, quae istius doctrinae authoribus hic a Paulo tribumntur. Tri buit igiturhis, Hypocrisin, et Conscientiam cau  \pend
\section*{AD I. PAVL. AD TIM. }
\marginpar{[ p.104 ]}\pstart teriatam istis inquam, authoribus haereseon, quos δαδολόγος quoque appellat.  \pend\pstart Hypocrisin hic quidam referunt ad eorum doctrinam, qui variis superstitionibus et caeremoniis verum Dei cultum corrumpunt. Hunc enim praetextum accipiunt haeretici et impostores homines, quo melius mendacia sua tegant, vt est Coloss.2. vers.23. nimirum quod augustiorem, sanctioremque Dei cultum inducere se iactent et cogitent. Itaque in eorum ἐθελοθρησκέία et superstitioso cultu apparet et humilitas, et abstinentia quaedam macerarióque corporis maior quam in reliquis hominibus. Et quanquam pessima conscientia praediti sunt, neque sint ipsi sibi conscii se quod agunt, propter Dei gloriam agere: hunc tamen fucum et praetextum purioris et sanctioris Dei cultus, clamosissime obtendunt. Sic quidam vocem hanc explicant.  \pend\pstart Alii hypocrisin ad impudentem potius audaciam, qua vtuntur, referunt, cùm eo hominum genere nihil sit magis effrons et impudens. Id quod etiam docet Paul.2. Timoth.3. Mihi vero videtur Paulus hic tria complexus esse, nempe 1 Doctrinam ipsam. 2 Mores. 3 Internum affectum animi. Sanae, doctrinae opponitur tum quod supra dixit eam doctrinam Diabolicam esse: tum quod hoc loco vocat istos εώδολόγος, vel δαδολογίας: vt alii legere malunt hoc loco. Moribus sanctis opponit Paulus istorum hypocrisin, vti fere in scriptura opponi solet.  \pend\pstart Interno animo affectui recto et pio, qualis in veris doctoribus inest, opponit Paulus conscientiam eam, quam cauterio notatam et infectam  \pend
\section*{CAPVT  IIIR }
\marginpar{[ p.185 ]}\pstart vocat, id est, non sanam, et rectam. Haec autem docent haereticorum hoc loco descriptorum neque esse sanam doctrinam, neque sanctos mores neque rectum animi affectum in suis illis commentis tradendis, quicquid illi tamen contra praetendant et praetexant, vel persuadere conentur.  \pend\pstart Cauteriatam vero conscientiam quam appellet Paulus varie etiam quaesitum est. Alii enim interpretantur putrem, et morbidam: alii nullam, vt pote quae iam perdita prorsus et excisa sit, vt sunt ea, que cauterio ambusta sunt: alii vexatam et perpetuo Dei iudicio tanquam cauterio quodam s'irbatam fluctuantem, et afflictam, atque ipse hanc postremam interpretationem sequor. Nam nulla non inest in iis conscientia, cùm eam cauteriatam esse dicat Paulus, non autem deperditam. Putris quidem est, id est, male sana, sed quae sunt putria non sunt huiusmodi, vt nihil plane sentiant et omnino emortua sint. Hic conscientiam explicat Paulus, quae assidue iudicio Dei voxetur et torqueatur propter malum testimonium et hypocrisin, quam simulat, et propter finem pessimum, propter quem istos errores propagat, cuius ipsa sibi testis ac conscia est, vt quae sit ἀυτοκατακριτὸς vt idem Paul. loquitur Tit.3.vers.1o.  \pend\pstart * Prohibentium contrahere matrimonium, iubentium abstinere a cibis, qous Deus creauit ad participandum cum gratiarum actione, fidelibus, et iis qui cognouerunt veritatem,  \pend
\section*{AD I. PAVL. AD TIM. }
\marginpar{[ p.188 ]}\pstart Εξήγησις, Rem enim ipsam et doctrinam illam falsam iam explicat. Duas autem illius species enumerat, nempe Ciborum prohibitionem, et Coningii, siue quod illae tuncipraecpue subnascerentur: siue quod primae. Certe ex iis duabus rebus superstitiose interdictis infinita postea aliarum superstitionum scaturigo et congeries vel genera sunt nata, et ’inducta. Magno vero nostro bono explicatio harum haerescon et diabolicae doctrinae a Paulo particularis addita est, ne vlli, ac ne haeretici quidem ipsi, generalem illam falsae doctrinae descriptionem a Paulo propositam a se detorquere, et in alios torquere poslent, essetque  incertum, quinam hic designarentur a viritu laneto. Piuius autem promoitionis meminit Epiphan. in Collyridianis, vbi etiam tertium caput addit, nempe Mortuorum adorationem siue eultum. Meminit etiam Socrat.libus 2. cap. 43. Eusebus liber 5. Histor. Eccles.  \pend\pstart Hic autem notandum est a Paulo non certas quasdam personas, sed doctrinam potius ipsam notari, quae vel, In vniuersum, vel In parte Cibos et coniugia prohibet, atque in earum rerum abstinentia cultum Dei et partem pietatis atque sanctitatis vitae nostrae constituit. Id quod prorsus cum verbo Dei pugnat.  \pend\pstart Ac plures sunt sectae, quae his Pauli verbis damnantur. Primùm, quae in vniuersum damnarunt ea quae hic commemorantur, et quae nobis concessa esse Deus voluit. Deinde, quae non in totum, sed ex parte tantùm haec ipsa prohibuerunt, de quibus et hic doctissime Cal. et nos in liber  Augu. de Haeres. ad Quodvultdeum diximus. Manichaei  \pend
\section*{CAPVT  IIII. }
\marginpar{[ p.189 ]}\pstart igitur comprehenduntur, Encratitae, Cataphryges, Montanistae, qui omnes in hac parte Pythagoreorum delirio fascinati sunt. Sed et Priscillianistae quoque, vt ait Bernard.sermo. 66. in Cant. Deinde Papistae ipsi, qui eadem ista tantùm ex parte, non autem in totum, vt illi priores haeretici, damnarunt, quos hic notat Paulus: quanquam cibos omnes nemo vnquam prorsus et in totum prohibuit, propterea quod sine his vita haec corporis sustentari non potest. Sed ex cibis, alii negaut concedi Christianis hominibus Carnes, vt Manichaei, alii vinum, vt Encratitae. Atque  hae prohibitiones habent speciem magnae humilitatis et corporis macerationis: sed tamen cùm sint ἐθελοθρησκείαι et Dei cultus in illis statuatur, sunt diabolicae et in Deum ipsum blasphemae superstitiones, quae sunt omnino ex Dei Ecclesia profligandae, vt docet Paulus Coloss2. et Tit.1. vbi de iisdem rebus agit. Ac de quadam ciborum differentia et distinctione, non ea quidem, quae postea ab haereticis et Papistis inuecta est, et de qua Paulus hic vaticinatur: sed quae ex Mosaica praeceptione defendebatur, vt apparet Leuitic. 20.V.25. et Deut. I4.iam tunc contentio et certamen erat in Dei Ecclesia, vt docet Paulus Rom.14. vers.3 1.Cor. lorrerliej, de eprre de futuro super his rebus maiori errore, peste ac pernicie periculosiore fuerit operaepretium Ecclesiam moneri. Neque vero quod inter detestandos haereticos hic censentur, qui ciborum vsum prohibent, facit, vt ingluuies, crapula, luxus et intemperies probetur. Nam nec eo pertinet haec prohibitio ciborum vt sobrie non viuatur:sed vt ne in his rebus  \pend
\section*{AD I. PAVL. AD TIM. }
\marginpar{[ p.186 ]}\pstart pars cultus Dei censeatur posita: nec quod addit Paulus, eos nobis a Deo concessos esse, aliter intelligendum est, quam vt iis, conseruata moderatione et temperantia Euangelica et Christiana, vtamur. Cauendum enim est, ne crapula et cibis grauentur corda nostra, vt docet Christus ipse Luc.21.vers.34. Sic hune locum explicat quoque August. liber  contra Adimantum Manich. cap. 14. Et aliud est prohibere cibos, aliud autem intenperantiorem eorum vsum damnare, quod etiam nos facimus ex Dei verbo.  \pend\pstart Sed neque ieiunia quoque tolluntur et abrogantur, quorum vsus frequens fuit in Ecclesia Dei, atque vtinam inter nos esset frequentior et religiosior eorum obseruatio, cùm ab Ecclesia ex iustis causis publice indicuntur. Aliud enim est cibos ex superstitiosa opinione damnare, et in eorum delectu et vsu regnum Dei collocare, quod facere vetat Paulus et hic et Rom. 14. vers.17. aliud autem iis ad tempus abstinere tantùm ad conpuibendam carnis lasciuiam, quae nimio pastu, tanquam ferox equus, superbit, et furit. Itaque ieiunia non damnantur hoc loco a Paulo. Sic August aduersus Manichaeos liber 6. contra Faustum cap.8, et 3o, cap.5. Sic Bernar. sermo. 66. in Cant. qui plane cum hoc loco conuenit, et huius loci Pauli interpres est optimus.  \pend\pstart Hanc autem esse haeresin, et diabolicam doctrinam variis argumentis et rationibus probat Paulus: sed imprimis a fine, propter quem cibi a Deo creati sunt. Duplex quidem erroris genus hic a Paulo propositum est, vnum de cibis: alterum de coniugio:tamen hoc primum solum persequitur,  \pend
\section*{C AAVTIIII. }
\marginpar{[ p.187 ]}\pstart et refutat : alterum praetermittit, quod hoc malum nondum vicinum esse praeuideret Paulus. Sed apertissime in Epist. ad Hebr. cap.13. vers.4 ref atatur, et iugulatur iste error de prohibitione coniugii. Etsi vero praecipuus confutationis locus a fine ducitur, alia tamen argumenta admiscentur, quae singulis pene Pauli verbis comprehensa sunt. Haec autem optima erroris cuiusuis tollendi via, et ratio est, si ad Dei ordinationem et verbum nostras sententias et definitiones reuocemus. Quod ipium m'unputatione de donafacit alibi Paulus 1. Corinth.11. Sed singula verba examinemus. Nam in iis varia latent argumenta, quemadmodum diximus.  \pend\pstart Deus creauit) Locus refutationis a temeritate istorum hominum et a sacrilegii crimine. Nemo inrem alienam, et praesertim Dei, ius habet. At cibi sunt Dei creaturae, non istorum. Ergo Dei est, non horum, leges de cibis praescribere. Et qui aliter faciunt, ingerunt se in aliena quae ad eos non pertinent, de quibus nullum ius habent definiendi et statuendi, vt ait Paulus Coloss.2. vers.18. Est enim quisque rei tantùm suae, non autem alienae legitimus moderator. Hinc sane refutantur varii haeretici, quorum alii vinum, alii carnes, alii aliud esse Diaboli, non autem Dei opus et creaturam pessime atque falsissime contendebant.  \pend\pstart Ad participandum) Ait Paulus ἐις μετδ λητη. vox ἐις finem declarat. Μετάληδις vsum, more dicendi in quauis lingua vsitato, vt etiam in Odysseam Homeri annotat Didymus interpres antiquissimus, vt ex antecedenti consequens, et contra concludatur et intelligatur. Nec enim tantùm a  \pend
\section*{AD I. PAVL. AD TIM. }
\marginpar{[ p.192 ]}\pstart nobis cibos assumi posse vult, et concedit Deus, vt tractentur, spectentur, conseruentur: sed vt edantur, consumantur, et iis famem nostram subleuemus. Sunt enim cibi ex istarum rerum genere) quae vsu ipso consumuntur. Itaque in genere suo functionem, vt aiunt Iurisconsulti, recipiunt. Est autem hic locus ex praecipuis refutationis superioris erroris argumentis. Ducitur a fine, propter quem a Deo ipso illae res sunt conditae et creataes Ad eum enim sunt omnia referenda, nisi et Dei ordinationi resistere, et eius placitum euertere impie, at que impudenter velimus. Ergo cùm Dominus hominem creasset, tamen noluit eum sibi Iuo aronttio en ercutario ipllus rapere ad victum quas vellet: sed ipse tanquam optimus paterfamilias quas et quantum vidit necessarias esse, illi concessit benignissime. Concessit autem eas, quas cibi nomine propter vsum appellamus.  \pend\pstart Confirmatur autem hoc Pauli dictum Genes.1. vers.29.9. vers3. Matth.15.vers.11. et Actor.1o.V.13 Explicatur etiam a Fulgentio in liber  de Fide ad Petrum cap. 42:  \pend\pstart Dissimiles tamen loci quidam, qui solebant a Manichaeis huius erroris summis patronis afferri et a Papistis etiam regeruntur, sunt sane a nobis diluendi. Ac primùm Papistae excipiunt, cibos istos non ipsos per se suaeque substantiae ratione polluere edentes, sed propter Ecclesiae prohibitionem et authoritatem, quae violatur a sumentibus. Respondeo vero ipsam Ecclesiam id prohibere non posse. Neque enim in res alienas ius habet, neque conscientiis nostris legem et iugum imponere potest, quod Dominus ipse non impo\pend
\section*{CAPVT  IIII. }\pstart soit. Itaque petunt illi principium, vt loquuntur, cùm volunt illam Ecclesiae prohibitionem de ciborum vsu et differentia firmam et legitimam esse. Quam quaestionem breuiùs tracto, quia ab infinitis nostro seculo, imprimis autem a D. D. Luthero, et Caluino explicata est.  \pend\pstart Manichaei vero et Tatiani alios aliquot scripturae locos obiiciunt aduersus Paulum. Primùm Genes.1.vers.29. vbi nulla fit carnis mentio inter eos cibos, qui homini conceduntur a Deo: sed tantùm herbarum, olerum, et fructuum arboris. Ergo si ad primam Dei institutionem reuocemur, carne nos abstinere oportere aiunt. Resp. i uuramipoouno lonmiuminquo concessionem, quae homini primùm a Deo facta est: sed eam, quae postea, tanquam illius beneficii ampliatio quaedam, accessit, et extat Genes.9.vers.3 vbi carnium esus diserte conceditur, et exprimitur.  \pend\pstart Secundo locum Pauli afferunt, ex 1. Corinthe 9. vers.7. vbi lactis tantùm, et eius esus meminit Paulus, non etiam carnis. Resp. Non excludere illam Pauli sententiam esum carnis, cuius ipse meminit, vti rei licitae etiam in eadem Epistola 8. vers.13. Nam illic alia Pauli mens est et scopus, quam vt explicet vtrum homini Christiano carnem edere liceat, nec ne.  \pend\pstart Tertio dicunt inter res ad victum humanum necessarias non recenseri carnes in Ecclesiastic, cap.39, vers.31, et 32. cùm tamen multa ciborum genera enumerentur. Respond. Quod illic praetermissum est, alibi additum esse. Neque ex eo libro regulam fidei vitaeque instituendae petendam esse, cùm sit liber Apocryphus. Deniqs illic oepen\pend
\section*{AD I. PAVL. AD TIM. }
\marginpar{[ p.194 ]}\pstart Cque. ex specie vna reliqua ciborum genera intelligi oportere.  \pend\pstart Cum gratiarum actione) Tertium argumentum  \pend\pstart Quarto, docent Deum ipsum hanc ciborum differentiam instituisse maioris cuiusdam sanctitatis gratia: et olim carnium esu interdixisse populo suo. Deut. 22. Quae res non est visa leuis momenti, cùm laudentur qui etiam mortem obire non recusarunt, ne vetitas a Deo carnes esitarent, vt apparet ex historia Macchabaeorum. Respond. Cùm adhuc sub paedagogia versaretur Dei populus, nec eius libertatis, quam per Christum adepti sumus, plenum vsum haberet (quo inter Christi nondum manifestati et manifestati tempora perspicuum discrimem appareret) Dominum voluisse quibusdam vmbris vitae spiritualis sanctitatem, quam exigebat, illis repraesentare, quae illis quidem partim variis lotionibus: partim ista ciborum distinctione proposita fuit, tanquam rudioribus adhuc puerulis. Quae causa in nobis cessat, qui per Christi aduentum et gratiam ad legitimam Ecclesiae aetatem peruenimus, et pueri esse desiimus. Itaque haec causa fuit instituti istius delectus ciborum a Deo, nempe vt illis adhuc rudioribus esset vitae sanctioris delineatio et significatio:non quod ciorprouibitripsi per se impuri essent, aut quod per fele ipn ciorconcessi puriores, cùm tota ea res carnalis fuerit (vt ait Scriptura) quae proprie animos non attingit, vt docetur in Epist. ad Hebraeos cap.9. vers. 1o. 7. vers. 16. Atque haec ipsa hodie firmissima est Papistarum ratio, quam a Manichaeis sumpserunt.  \pend\pstart ab haereticorum blasphemia ductum. Est ea do\pend
\section*{CAPVT  IIII. }
\marginpar{[ p.195 ]}\pstart ctrina blasphema, quae Deum spoliat suo honore, et ea gratiarum actione, quae propter summa ipsius in genus humanum beneficia illi debetur. At haec Deum isto honore priuat. Ergo, etc. Addit autem quibus maxime hic creaturarum vsus sit concessus a Deo, nimirum piis, qui eum nouerunt, et in eum credunt. Hic enim, nosse, non simplicem et historicam agnitionem significat: sed Corurinnurencent de cum ami naucia coniunctam, per quam in Deum homines acquiescunt vera animi fiducia et πληρςφορία. Similis locus Romanor.4.vers.13.  \pend\pstart Dissimilis locus Matth.5. vers.45. et Psal.17.y. 14. vbi dicuntur etiam infideles a Deo impleri ipsius thesauris. Respond. Dantur quidem a Deo terrena bona etiam infidelibus hominibus: sed in eorum exitium, et perniciem, quo magis reddantur inexcusabiles. Itaque non tam iis vtuntur, quam abutuntur: neque tam vsus abiis fit, quam depraedatio quaedam, quoniam Deum dantem spernunt, in cuius contumeliam et contemptum nta dona lumunt et coniumunt, quem non agnoscunt eorum authorem.  \pend
\phantomsection
\addcontentsline{toc}{subsection}{\textit{4 Nam quicquid creauit Deus, bonum est, nec quicquam reiiciendum est, si cum gratiarum actione sumatur. }}
\subsection*{\textit{4 Nam quicquid creauit Deus, bonum est, nec quicquam reiiciendum est, si cum gratiarum actione sumatur. }}\pstart Est hic locus canon, et quidem generalis, ac maxime necessarius. Iis omnibus nimirum vti nobis licere, quae in nostrum commodum creauit Dominus, modo cum gratiarum actione illa sumamus. Ex quo colligitur, 1 ad eum finem  \pend
\section*{AD I. PAVL. AD TIM. }
\marginpar{[ p.196 ]}\pstart rebus creatis vtendum ad quem a Deo conditae sunt. 2 Non modo Deo gratias agendas, postquam iis rebus vsi sumus: sed etiam Deum inuocandum, vt vti liceat, cùm ex eius manu sint eae res sumendae: cuius proprie sunt, vt apparet ex oratione Dominica.καλὸν vero non refertur hoc loco ad mores honestos: sed ad vsum, qui a nobis Dei nomen inuocantibus recte fieri et salua conscientia potest. Alias enim perinde, atque fures dannaremur etiam ex minutulae cuiusuis rei vsu: et merito quidem, quasi rem alienam contrectantes: nisi Dominus ipse nobis eam in vsum concederet, eamque commoditatem ex rebus a se creatis permitteret. Bonitas item ciborum hic non pertinet ad valetudinem et sanitatem corporis, nec ad salubritatem ipsorum ciborum. Neque enim hic medicum agit Paulus, sed Theologum et Apostolum: ne quis forte obiiciat sibi aliquos cibos non esse salubres, dum iis vtitur. Atque propterea hunc Pauli locum calumnietur, quia generaliter dicit Paulus ὀδὲν ἀπόβλντον, id est, nihil esse reiiciendum ex iis, quae Dominus in vsum et cibum hominum creauit. Quid enim liceat conscientiae nostrae ratione spectat Paulus, non autem quid conferat stomacho nostro, vel valetudini corporis.  \pend\pstart Hic tamen et inferiori versiculo duplicem ciborum bonitatem proponere videtur, Vnam, quae ex Ipsorum substantia: et Alteram, quae ex vsu nostro consideranda et aestimanda nobis est. Quod ad primam: illa nobis testata est Genes. cap.1. Quaecunque enim condidit Deus, illa substantiae, in qua sunt creata, ratione bona sunt, vt Dominus ipse,  \pend
\section*{CAPVT  IIII. }
\marginpar{[ p.1 ]}\pstart pronuntiauit. Nec obstat quod quaedam ex iis sunt nobis venena. Respondeo enim non idcirco mala esse. Malum enim, quod bonitati rerum creatarum opponitur, non est vsus noxius (nam arbore scientiae boni et mali impune vti non licuit Adamo) sed est malum hoc loco qualitas interna quae peccatum in nobis producit: et Dei ordinationi contraria est, nobisque propterea noxia, et detestanda. Deinde quae creaturae sunt nobis venena, illae non sunt nobis in cibum a Deo destinatae. Denique hoc ipsum vitium, quod rebus creatis nunc inest et superuenit, non ex ipsa Dei creatione, sed ex peccato hominis accessit. Distinguendum autem est vitium et accidens a rei natura et substantia, vt docet Augustinus in articulis contra Felicem Manichaeum cap.8.9. et in liber  de Fide ad Petr. cap.42.  \pend\pstart Si cum gratiarum) Altera bonitas ciborum ex vsu ipso, et nostri ratione bonitas appellanda. Vsus autem hic certis legibus a Paulo circunscriptus est, vt a profanorum mensis piorum hominum conuiuia et victus separetur. Cur autem has leges addat Paulus, ratio est, quod pro concesso sumit, Nihil nobis, tanquam a nobis ipsis, in Dei creaturis iuris esse vel licere, sed id tantùm, quod Deus, qui est earum Dominus, quantumque concessit, et largitur. Ergo si, vel alio Modo, vel fine, illis vtamur, sane peccamus. Quoad modum autem, Dominus non vult nos immoderate suis rebus et donis vti, quia hic plane abusus est, non vsus: atque  ipsius dona non sunt spernenda, aut contemptim habenda, quod ii faciunt, qui per crapulam et gulam sese ingurgitant, sed magni sunt ae\pend
\section*{AD I. PAVL. AD TIM. }
\marginpar{[ p.Io ]}\pstart stimanda. Quoad finem vero, ideo sumi vult, tanquani a luu iplruo manunoo cidos sumamus, ve ipse eorum author agnoscatur et laudetur.  \pend\pstart Similis locus Psal. 16. et Psal. 5o. Huius autem totius reirationem reddit lequeti versiculo Paufus.  \pend
\phantomsection
\addcontentsline{toc}{subsection}{\textit{5 Sanctificatur enim per verbum Dei et preces }}
\subsection*{\textit{5 Sanctificatur enim per verbum Dei et preces }}\pstart Α’ιτηλογία est ducta a causa instrumentali (quae refertur ad genus causarum efficientium) per quam ciborum vsus fit nobis purus, sanctus, et legitimus. conscientiae nostrae ratione. Ergo sanctificatio hoc loco opponitur non tantùm profanationi ciborum:sed etiam malae conscientiae, et ei vsui, qui nobis in exitium et condemnationem apud Deum cedit. Sanctificatur cibus cùm eum salua conscientia sumimus, et eo vtimur, quia nobis Dominus ipse eum largitur ad finem huiusmodi. Neque vero sic sanctificarl cibus dicitur, quemadmodum homo. Nam qui a spiritu Dei sanctificatur homo, is immutatur, et nouam tanquam naturam induit. At qui cibus sanctificatur in nostrum vsum, totam integramque naturam suam retinet, fit tamen Deo gratus ex eo, quod Deus nostra inuocatione illius author esse a nobis agnoscitur et praedicatur. Ex hoc autem loco perperam explicato natae sunt magicae illae ciborum, velut ouorum, vini, pernarum, aquae execrationes et incantationes, quas Pontificii faciunt, quósque D. Martinus Chemnitius prolixe in Examinis Trident. concilii postrema parte persequitur et refutat doctissime. Duplex autem huius cibi san\pend
\section*{CAPVT  IIII }
\marginpar{[ p.19 ]}\pstart ctificandi ratio proponitur, verbum Dei, et Precatio. Verbum Dei, quia tum demum sana conscientia sumimus hos cibos, cùm per Dei verbum intelligimus eos nobis in esum concessos et datos esse. Itaque hoc Dei verbum fide a nobis apprehendi prius debet, quam hic cibus sanctificari possit. Cur autem ad Dei sermonem sit nobis recurrendum, ratio est, quod per Adamum omne ius illud, quod prius in Dei creaturas, et mundum hunc habuimus, perdidimus. Quod nobis restitutum esse nomnisi ex Dei verbo fide apprehen so possumus agnoscere. Hanc fidem sequitur precatio nostra, quae a Deo petit, qu'od est ipse mera luu belignitate pomeltus, et agnoscit eum tantorum bonorum authorem esse. Atque tunc demum nostrum esse et a Deo nobis donatum, quod apponitur, intelligimus, vt sine animi scrupulo vtamur.  \pend\pstart Haec autem sanctificatio communis mensae differt ab ea, quae fit in mystica Coena Domini, vti praeclare docet hoc loco Caluinus.  \pend\pstart Quaerit Chrysostom. de deliciis ciborum, vtrum iis ver niecat homini Christiano. Respond. vero cùm nobis profusa liberalitate sua dona largitus sit Dominus, nec vnius tantùm generis cibos in vsum nostrum creauerit, posse nos cum gratiarum actione, etiam cibis lautioribus vti. neque tantùm vulgaribus, modo omnis luxus, ingluuies, intemperantia, et crapula absit.  \pend
\phantomsection
\addcontentsline{toc}{subsection}{\textit{}}
\subsection*{\textit{}}
\section*{AD I. PAVL. AD TIM. }
\marginpar{[ p.200 ]}
\phantomsection
\addcontentsline{toc}{subsection}{\textit{bus fidei, bonaeque doctrinae, quam assectatus es. }}
\subsection*{\textit{bus fidei, bonaeque doctrinae, quam assectatus es. }}\pstart Παράκλησις et exhortatio, per quam quae de pastoris officio et vera doctrina contra haereses tuenda in vniuersum dixit, applicat ad Timoth. Sic enim nobis agendum est, vt quae funt nostri muneris diligenter perpendamus, et ex iis intelligamus, vtrum fungamur officio. Duo vero praecipit Paulus I.vt superiorem doctrinam Timot. suggerat, 2 vt fratribus ipsis. Suggerere (inquit Chrysostomus) est frequenter proponere. Vti enim cibis quotidie ad pastum corporis egemus: sic Dei verbo et sana doctrina ad pabulum animi, vt semel sanam doctrinam proposuisse non fit fatis:neque  odiola censeri debeat huius doctrinae tam necessariae et salutaris repetitio, vt est 2.Pet. I.vers.12. et 3.vers.1. cùm Sathanae impetus in oppugnanda Dei doctrina nunquam fatiscat et conquiescat.  \pend\pstart Dixit autem Paulus ὐποτιθέμενος i.subiiciens, vel suggerens non autem imperans, vt obseruat hoc loco Chrys. quod distincta sit muneris et ministerii Ecclesiastici ratio a ciuili imperio. Magistratus imperat: denuntiat autem tantùm. Minister verbi Dei. Quae tamen denuntiatio habet coniuncta Dei iudicia et minas, si non illi pareatur, longe omni poena corporali formidabiliores et tetriores. Paulus tamen 2. Thess.3. vers.12. vtrunque  coniunxit et imperium et obsecrationem, quia Domini ipsius nomine agunt Ministri Euangelici:non autem suo et priuato.  \pend\pstart Ronus erus minister) Duplex exhortationis locus  \pend
\section*{CAPVT  IIII. }
\marginpar{[ p.201 ]}\pstart nimirum tum ab officii ratione, tum a prima Timothei vitae institutione. Ergo vtrunque nobis sedulo cogitandum est, et quid munus, in quo versemur, requirat: et quid susceptae vitae ratio et institutio. Vtrunque enim est nostrae Vocationis a Deo factae certissimum testimoninm,  \pend\pstart Quatenus igitur Minister munus sulim considerat tria haec spectare debet, vt sit Minister vt Iesu Christi, vt Bonus siue fidelis. Quae his tribus opponuntur, nempe Ministris Otiosis, Humanae doctrinae, et Infidelibus. Hi vero sunt, qui vel Seipsos, id est, suam gloriam, vel quaestum suum quaerunt. Ergo sana doctrina a fidis, Ministris verbi Dei proponenda est: non autem inanes, non futiles, et subtiles quidem, sed futiles quaedam quaestiones, quae neque nostras conscientias aedificant:neque Dei gloriam, salutisque doctrinam vllo modo explicant vel illustrant.  \pend\pstart Quod autem addit, snnutritus sermonibus fidei. etc, magnam vim et pondus exhortationi addit. Est enim prima educatio tanquam tacitum futurae vocationis indicium a Deo demonstratum.  \pend\pstart Denique turpe est illa talenta, quae Deus in nos contulit, vel omnino perdere, vel infructuosa habere, vt docet Christus Matth. Vtitur autem his verbis Paulus, ex quibus intelligitur, seriam prius operam Timotheum nauasse in sacris literis, et doctrina fidei Euangelica. Id quod apparet etiam ex 2.Timoth.1. vers.5. sup.cap.1. vers.18. Ex hoc autem loco quaeri duo possunt 1. Vtrum eos solos in ministerium Euangelicum eligi oporteat, qui a primis et teneris annis sacris sçriturie sunt innutriti et eas consectati:an etiam alii  \pend
\section*{AD I. PAVL. AD TIM. }
\marginpar{[ p.202 ]}\pstart possint. Respond. Deum quidem in quem vult dona sua conferre. Itaque eum nobis eligendum esse, cui maiores dotes et χαρίσματα De? ipse dederit. Tamen longe tutiùs esse, vt, nisi magna aliqua causa subsit, ii potissimum eligantur, qui diutius nii iueru beripture lectione versati sunt, et a teneris annis sanam Domini doctrinam tum didicerunt, tum etiam professi sunt.  \pend\pstart Secundo quaeritur, An a suscepta semel professione Euangelici ministerii abduci et reuocari homines pastores possint. Vult enim Paulus Timotheum in ea cognitione et arte perseuerare, quam semel assectatus est. Deinde damnantur et Demas et Marcus, qui Paulum deseruerant, et fusceptam semel ministerii vicariam delegationem reliquerant, mutato vitae genere ac instituto. Res. Deum nullum huiusmodi tam arctum et strictum vinculum inter pastorem et gregem instituisse, vt nunquam solui possit. Saepe enim ex multis, et iis quidem iustis, causis liberantur pastores, et habent a suo susceptóque munere legitimam missionem. Id tamen, nisi ex iusta causa fiat, proculdubio in temeritatis et eius quidem turpissimae non tantùm suspicionem, sed etiam criminationem merito incidit, qui consilium suscepti semel ministerii mutauit et deserit. Et quod Demas, et Marcus damnantur, ex eo est, quod seculum Deo praetulerant: non tamen propterea excommunicantur, aut eiiciuntur ex Ecclesia a Paulo, etsi nulla fuit deserendi ministerii iusta, quae ab iis praetenderetur causa. Et quod ait Bernard.epistola 87.Aut ergo oportuit te gregem dominicum minime seruandum susciperetaut susceptum nequaquam relin\pend
\section*{CAPVT  IIII. }
\marginpar{[ p.203 ]}\pstart quere iuxta illud Alligatus es vxori, noli quaerere solutionem. Primum falsa similitudine nititur coniugii, et vinculi pastoralis: deinde nec eum ipsum Ogerium, qui gregem suum deseruerat, dannat, tanquam qui sit ab Ecclesia expellendus, sed qui sui facti et temeritatis poenitere debeat, non de eo gloriari. Et sane in eo falluntur homines, qui non putant magis licere pastori Ecclesiastico ministerii professionem semel susceptam mutare, quam maritato coniugem non adulteram dimittere. Etsi sine ratione ista vitae mutatio fieri non debet.  \pend
\phantomsection
\addcontentsline{toc}{subsection}{\textit{Caterum profanas et aniles fabulas reiice: sed exerce teipsum ad pietatem. }}
\subsection*{\textit{Caterum profanas et aniles fabulas reiice: sed exerce teipsum ad pietatem. }}\pstart Ἀντίθεσις. Sanae enim et superiori doctrinae opponit Paulus fabulas, quas Prophanas, et Aniles vocat. Fabularum nomine intelligit tum superius damnatam de prohibitione ciborum doctrinam: tum etiam reliquas, quae in rebus externis cultum Dei constituunt Prophanas voca, quod sint blasphemae. Aniles autem, quod sint nihili. Est autem hoc postremum dictum ex forma quadam vulgaris prouerbii sumptum, quia solent anus mulieres fabulas et fictitia quaedam commemorare lubentiùs, quod propter aetatem multa sibi tribuunt: atque sunt natura verbosae. Has tamen fabulas Chrysost. refert non ad superiorem doctrinam, sed ad Iudaicas traditiones, et caeremonias: sed perperam, vt quilibet facile animaduertere potest.  \pend\pstart Sed exerce teipsum) Transitio est, per quam ad alteram ministerii et muneris Euangelici parte acce  \pend
\section*{AD I. PAVL. AD TIM. }
\marginpar{[ p.204 ]}\pstart dit Paul. Est autem haec altera pars Pia, et religiosa vita. Neque enim pastores tantùm sanam doctrinam tradere debent, sed etiam vitae suae honesto exemplo eam confirmare, et caeteris praelucere. Quemadmodum et Christus ipse Matt.5.V.13. et Petrus 1. Pet. 3.V.2.3. docent. Ratio igitur huius additionis duplex est. Prima, quod Ministri verbi Dei caeterorum comparatione esse debent lux mundi, et lucidum exemplar sanctae vitae non tantùm fidelibus, sed etiam ipsis infidelibus, vt docet inf.v. Paul.12. et Tit.2.v.7. Secunda vero, quod doctrinae suae fidem illi derogant, qui sunt turpis, et inhonestae vitae pastores. Si quis enim docet non esse furandum et ipse furetur, eleuat suae praeceptioni vim et authoritatem, quemadmodum ait Paul. Rom. 2. Contra vero ipsa suae doctrinae praxi et praestatione constanti caeteros in ea doctrina confirmut Caueliuom igitur nobis est, ne praedicantes aliis ipsi reprobemur, nobisque vere obiiciatur illud, quod Christus Pharisaeis, Dicunt, sed non faciunt 1.Corinth.9.v.27. Matt.23.vers.3. Pulchra igitur est harmonia recte docentis et sancte viuentis Ministri verbi Dei.  \pend\pstart Similis huic locus est 2. Timoth. 2.vers.15.Tit. 2.vers.7.  \pend\pstart Dissimiles vero loci afferri videntur posse. Primum dictum Christi Math. 23. vers.1. Etsi enim erant Pharisaei moribus impuris, eorum tamen sanam doctrinam sequendam et retinendam esse docet Christus. Respond. Gregis quidem ratione illa Christi sententia vera est, eorum nempe pastorum doctrinam esse amplectendam: qui a veritate non discedunt. Sed tamen illi ipsi si male  \pend
\section*{CAPVT  IIII. }
\marginpar{[ p.205 ]}\pstart vixerint, eo grauius sunt damnandi, quod eos, quos verbis aedificare student, vitae turpissimae exemplo perdunt. Postea vero aduersus hunc locum obiicitur illa Christi sententia, Martha sollicita est circa plurima Luc. 1o. vers. 41. Ex qua vitam contemplatiuam Actiuae, quam vocant, et quam hac voce, γύμναZε, significasse videtur Paulus, prtferri colligunt. Respond. Nihil illum locum cum hoc Pauli habere commune. Alia enim plane Christi mens est illic, quam vt piae et actiuae, quam vocant, vitae studium restinguere in nobis conetur Christus: quod actiuae vitae genus hic procul dubio etiam Paulus commendat: at que vanae et Monachicae hominum speculationi. i. vitae contenplatiuae postponit, quicquid de loco illo Lucae Monachi imperiti sentiant.  \pend\pstart Ait γύμναζε Exerce) Qua voce duo demonstrat Paulus, 1 Continuum vitae nostrae studium et exercitium in eo esse debere, vt sancte, et pie, et religiose viuamus. Is enim est finis, propter quem salutaris illa Dei gratia per Christum apparuit. 2 Vt intelligamus quam natura inepti et indociles ad huiusmodi opus simus, ad quod nos idcirco formari longa et perpetua exercitatione, cura, meditatione, et vigilantia necesse est. Ataque ad illud non nascimur apti. Nam ad male agendum non est nobis vlla exercitatione opus, cùm sponte in illam partem propendeamus et procliues simus. Denique ex hoc loco verum esse intelligitur, quod ait Bernardus epist. 201. Pastoris officium esse, vt pascat gregem suum verbo: pascat exemplo, pascat et sanctarum fructu orationum. Nam vocis virtus est opus, et operi et  \pend
\section*{AD I. PAVL. AD TIM. }
\marginpar{[ p.206 ]}\pstart Pietatem) Non tantùm doctrinam ipsam de side complectitur hoc loco vox ἀσεθέας: sed ipsius doctrinae praestationem et praxin, quae in vita, conuersatione, charitate, cultu Dei, caetePirque friliibus ais reuus ponta est, vt Paulus ipse inf.vers.12. interpretatur. Id quod ex sequentis versiculi ἀντεθέσει manifestissime apparet.  \pend\pstart Ex hoc loco facile colligi potest, Christianae Theologiae cognitionisque veri Dei verum finem non in mera contemplatione et θαρεία constitui a Paulo: sed in vita et actione ipsa, per quam tales efficimur, quales nos esse oportere doctrina ipsa iubet, atque  ipso Christianorum nomine profitemur. Vnde fit, vt omnino Monachismus ille Papisticus subruatur, qui Christianae religionis professionem in nuda rerum contemplatione, et doctrinae meditatione collocat et deffinit, si modo veram Christi doctrinam haberent tamen.  \pend
\phantomsection
\addcontentsline{toc}{subsection}{\textit{8 Nam corporalis exercitatio paululum habet vtilitatis: at pietas ad omnia vtilis est, vt quae promissionem habeat vitae praesentis ac futurae. }}
\subsection*{\textit{8 Nam corporalis exercitatio paululum habet vtilitatis: at pietas ad omnia vtilis est, vt quae promissionem habeat vitae praesentis ac futurae. }}\pstart Τ. ποφορὰ est. Occurrit enim tacitae obiectioni, quae ex corrupto iam hominum iudicio fieri potuit: vel quam in posterum fieri debere sanctus Apostolus et spiritu prophetico plenus iam tum praeuidit. Cùm ergo mirum in modum semper homines ad externas illas actiones quae speciem pietatis prae se ferre videntur, stupeant, mirum etat Paulum vnius tantùm pietatis meminisse, in  \pend
\section*{CAPVT  IIII. }
\marginpar{[ p.209 ]}\pstart qua nos exerceri, et operam ponere vellet: non etiam externarum vllarum caeremoniarum, quas tantopere homines tamen laudant et admirantur. Respond. Paulus, illam vnam pietatem vtilem esse tum in sese, tum aliarum exercitationum et operum, quae homines consectantur, ratione et comparatione.  \pend\pstart Apparet vero tum ex hoc loco, tum etiam ex epistola ad Coloss2.iam aetate Pauli plus satis externae illi vitae rationi et austeritati homines etiam, qui Christum profitebantur, tribuisse, illique addictos fuisse. Quod ex eo profectum videri potest, primùm quod fermentum illud Pharisaeorum, quod tamen toties Christus reprehenderat, iam inualuerat apud plerosque, quemadmodum ex quorundam animis falsa illa de rituum Papisticorum, et habitus monachalis sanctitate persuasio euelli hodie non potestMMatth.23. Deinde haec haereseun et superstitionum semina, quae postea creuerunt, et Papismum et Monachismum penitus, vti foetus infelicissimos, pepererunt, a Sathana iaciebantur, et serebantur in animis Christianorum hominum, id est, mysterium iniquitatis iam agebatur vt docet Paulus 2. Thess.2. Denique quod humanum ingenium vti suapte natura vanum est ita vanarum rerum amans est et auidum, quales sunt omnes illae externae cultus Dei vel vitae sanctitatis species. Itaque mature tantae huic verae fidei et religionis pesti occurrere voluit Dei spiritus: sed propterea tamen non est excisa tanti, tamque  noxii mali radix: sed statim in Ecclesia pullulauit, primùm per Nicolaitas, et Ebionaeos: deinde per Gnosticos, denique per Mona\pend
\section*{AD I. PAVL. AD TIM. }
\marginpar{[ p.208 ]}\pstart chos, per quos tandem vis et authoritas huius tam lacrolanetr t'aui praecepti, non tantùm elusa, sed abrogata et prorsus antiquata est.  \pend\pstart Quaesitum autem est, Quid corporalis exercitationis nomine complectatur Paulus. Putat Chrysostom. corporis agitarionem, quae sanitati conferat, significari, qualis est venario, ludus, fossio, aut quae exercitamenta iubentur a Medicis ad superfluos corporis humores consumendos fieri. Rationem addit ille, quod corporis exercitationem cum ea, quae animi est, et spiritualis, confert, vt hanc vtilem: illam nullius momenti homini Christiano esse doceat. falsa est haec in terpretatio. Nam non agit medicum Paulus, sed Theologum: Et vt supra cùm de cibis agebat, de iis eatenus disputauit, quatenus ad conscientiae tranquillitatem pertinere censebantur, non autem ad corporis valetudinem: sic hoc loco de exercitatione corporis disserens eam intelligit, in qua cultus Dei et nostrae sanctitatis pars quaedam constituitur, quoniam in rebus externis eam homines collocabant.  \pend\pstart Est igitur exercitatio ista, de qua nunc agit Paulus, alia, quam censet Chrysostomus, et est ea, per quam carnis mortificatio, vitae sanctitas, et externa quaedam pietatis cultusque Dei professio demonstrari videtur inter homines, qualis quae in ieiuniis et afflictione magna corporis versatur. Haec tamen dicitur per contemptum σωματικὴ, ne per eam animus vere reformari censeatur. aut quisquam reipsa propterea religiosior sapienriórque fieri iudicetur: cùm tota vis et effectus istarum rerum circa corporis tantùm maceratio\pend
\section*{CAPVT  IIII. }
\marginpar{[ p.209 ]}\pstart Sed et σωματικὴ quoque nominatur a subiecto. Est enim corpus, quod his modis atteritur, affligitur, et has externas actiones agit: quanquam ratione finis, quem spectant qui haec taedia perferunt, dici quoque potest haec exercitatio spiritualis, quod spiritualis vit ae gratia fit.  \pend\pstart Haec igitur omnia externa sibi imponunt homines, et quae vocantur nunc regulae ordinis suscepti et professi, eas appellarunt Patres A’σκnσwν. Vnde veteres Monachi et Μονάζοντες et Α’σαηταὴ dicti sunt: Paulus autem nominat γναναςίαν, quod Αʹσκνσις magis quiddam spirituale significare videtur: yuncu vero quandani tantum corporis vexationem ac agitationem, quae non tam commendat monasticum statum, quam illi conmendatum esse inter homines voluerunt. Sed praestabat vocem Spiritus sancti retinere.  \pend\pstart Iirnis autem extermis rebus praescriptis obseruandis versatur tota Monachismi essentia, quam vocant. Nam in tribus hisce potissimhm posita dicitur, in voto Obedientiae, Paupertatis, et Continentiae, tum matrimonii, tum certorum quorundam ciborum.  \pend\pstart Sed et illud quoque notandum est, Paulum hic non loqui de superstitiosa earum rerum obseruatione, qualis fit in Papatu, et praesertim a Monachis (eam enim prorsus damnasset Paulus, est enim superstitiosa, cùm pars salutis nostrae aliqua in iis rebus statuitur vti fit a Monachis) sed de rerum indifferentium discrimine constituto, et in iis obseruatione quadam carnis conterendae gratia, ad tempus indicta, qualis a piis viris pro  \pend
\section*{AD I. PAVL. AD TIM. }
\marginpar{[ p.210 ]}\pstart raria rerum, temporum, locorum, circunstantia fieri et suscipi potest.  \pend\pstart Hanc tamen totam vitae rationem et obseruationem, etiam ad tempus factam, dicit Paulus esse minimi momenti. At qui Monachi in iis ipsis rebus vitam Angelicam sitam esse volunt, et proram et puppim salutis et sanctitatis collocant.  \pend\pstart Neque tamen prorsus tales exercitationes damnat Paulus. Inter eas enim est ieiunium, et huiusmodi seruitus et maceratio corporis, qualem se etiam obseruasse ad carnem subigendam scribit ipse I, Corinth.9. vers.27. et Psal.35.vers.13. Similis locus est Matth.9. vers.14, et 15.vers.2. Romano.14 vers.17.1. Corinth.7.vers.5. Dissimilis autem, primùm, quia ait Paulus Ephes.5.v.29. fouendam esse carnem: at qui istas rigidas exercitationes obseruat, vexat sese, suumque corpus. Ergo videtur inhumanus et saeuus, vt hae non tantùm parum vtiles dici debeant, sed sint prorsus damnandae. Respond. Cum moderatione istas esse a nobis suscipiendas, et ad tempus tantùm: non autem cum superstitione, vel in perpetuum, vt docet Paulus 1. Corinth.7.vers.5. Nam omnes illae austerae, et seuerae Monachorum regulae, uo iis etiam cum periculo vitae et valetudinis dectimento obseruatae, fuerunt supelstitiones merae et Diaboli praestigiae, qualia quaterduana ieiunia, perpetuae humi cubationes, stationes continuae, vt sibi sedere non permitterent: et reliquae huiusmodi saeuae praeceptiones vel potius crudelitates in ipsam humanam naturam. Obstat etiam quod ait Paulus Rom.14. vers.17  \pend\pstart Regnum Dei non esse in esca et in potu. Ergo frustra  \pend
\section*{CAPVT  IIII. }
\marginpar{[ p.211 ]}\pstart stra suscipiuntur ista vel ad breuissimum tempus. Respond. Cùm ista sine causa et temere fiunt, frustra reuera suscipi a nobis. Cùm autem ad carnem lasciuientem domandam, vel vt paratiores promptioresque ad Deum precandum simus, non frustra: quanquam in iis rebus tamen regnum Dei nion est politum, sed in vera fide, regeneratione, spe, et obsequio eo, quod mandatis Dei praestatur.  \pend\pstart Pietas autem adomnia) Α'ντίθεβις est, ex qua superior sententia illustratur. Pietatem autem vocat verum Dei cultum, qui ei ex ipsius verbo exhibetur. Primum autem hic fidem continet, deinde charitatem: deinde cordis sanctitatem et puritatem, quae externis actionibus se patefacit. Haec vna res vtilissima est, et inf. capx̃. vers.6. magnus quaestus appellatur. Caeterum exemplo Christi, qui nihil praeter caeterorum hominum morem habuit: et Ioannis Baptistae, qui istas exercitationes retinuit, nec propterea tamen Christo melior fuit, aut sanctior:imo vero nec cum eo vlla in re conferendus, hoc Pauli dictum perspicue confirmatur. Fuit enim sanctior illa Christi vita conmunis, quam austera et rigida vita, quam tenuit Ioannes.  \pend\pstart Huius autem suae definitioms et laudis de pietate argumentum Paulus profert ab effectu, vel signo, quod pietas vera sine his exercitationibus habet promissiones, nimirum vitae, tum Praesentis, tum futurae. Magna certe laus verae religionis. Sed cur vitae praesentis promissiones habeat pietas, ratio est, quod Deus est suorum in vniuersum et per omne tempus Pater et ουτὴς, non tantùm  \pend
\section*{AD I. PAVL. AD TIM. }
\marginpar{[ p.212 ]}\pstart post hanc vitam: sed etiam in hac ipsa vita. Bern sermo.4. in Psal. Qui habitat. Credo in his scapulis geminam Dei promissionem intelligendam, vitae scilicet eius, quae nunc est: pariter et futurae. Si enim solum promitteret regnum, et in itinere deesset viaticum, omnino conquererentur homines et responderent: Magnum quidem est, quod promittitur, sed illuc perueniendi facultas nulla daturt vnde omnem excusationem Deus ademit. Deus igitur iam se eorum Deum ex hac vita esse ostendit, et iis benefacit Psal. 36. Incipit enim Dei promissio impleri in iis, etiam in hac vita.  \pend\pstart Obstat quod fides res spirituales, et eas etiam absentes spectat, non autem praesentes. Neque proprie corpus: sed potius beat et spectat animum. Bona enim spiritualia pollicetur Deus potius, quam temporalia Hebr.11. vers.1. Respond. Bonorum quidem iprastuanun promissionem potissimum et praecipue habent et spectant fides et pietas: sed tamen etiam de suis curam habet in hoc mundo Deus, et illis bona terrena, quantum latis esse peripiere jrargitut. Nam ea a se peti et vult et praecipit in oratione Domimca.  \pend\pstart Sunt autem quaecunque a Deo pii obtinemus ex mera ipsius bonitate, et ea gratuita, non ex melito: ilde bona terrena, siue coelestia spectentur, vt docet perspicue hic locus, qui ea omnia Dei promissiones appellat.  \pend\pstart Admonet hic locus, vt de Monachis aliquid dicamus, quia hi Α’σκηταὶ primum sunt appellati: veluti a Philone Iudaeo, si modo de Monachis Christianis potius, quam de Essaenis et Iudaeorum Monachis agit Philo:deinde a Basilio, et aliis, qui e\pend
\section*{CAPVT  IIII. }
\marginpar{[ p.213 ]}\pstart tiam Monasticas regulas Irc umpor homine inscripserunt.  \pend\pstart Μονάζοντες et Μοναχοι appellati sunt, quod vitae solitariae genus elegissent, qui hoc modo viuere instituissent. Μοναχὸν autem putant Canonistae ex can. Placuit S Agnoscat 16. quaest.1. dictum composita voce, ἀπὸ το͂ μόνον et ἄχος. Quae sane vera est hoc tempore Monachorum etymologia. Sunt enim illi re vera vnica pestis mundi, vt et ex veteri quoque Epigrammate iampridem intelligebamus. Sed falluntur et Papae et Papistae et Monachi in nominis sui etymo. Est enim syllaba xos ἐπενθέ Cς tantùm, ad vocem μόνος, a qua fit deriuatio vocis huius, Μοναχὸς. Vnde et Μοναχᾶ et Moναχο͂ et Μοναχῶς vocabula Graecis scriptoribus etiam ante monachorum exortum defunctis vsitatissima. Basil de Trinitate agens quanque  personam Μοαχῶς ἐκφωνεῖςς, id est seorsim pronuntiari dixit. Latine ipn se’Acengioros vocarunt, quan renquus populus Christianus sit irreligiosus. Item et Canonicos quoque et Regulares, et aliis nominibus ad sui commendationem. Sed et Mandritae etiam sunt appellati a cauernis, quas isti Troglodytae subibant, et vbi latebant.  \pend\pstart Monachus quid sit nondum perspicuede finiri potuit, etsi variae a variis descriptiones tentatae sunt. Neque enim ex vestitu et professionis vestimento definiri potest, quia, vt vulgo dicitur, vestitus non facit Monachum: et tale est monstrum Monachus, vt vna forma non constet. Tandem definitum est, esse eum Monachum, qui haec tria vota solenniter suscepisset seruanda, Paupertatem Continentiam coniugii, et Obedientiam vni cui  \pend
\section*{AD I. PAVL. AD TIM. }
\marginpar{[ p.214 ]}\pstart dam certo praeposito, et in eadem vitae regula degenti.  \pend\pstart Viuendi quidem instituta et regulas posse esse in quoque  ordine Monachorum diuersas, haec tamen tria vota esse ipsius vitae Monasticae essentialia, et ab ea semper indiuulsa. Quod Thomas Aquinas in 2 22 et in opere Contra impugnant. Religion. prolixe docet. Nam qui definiunt Monachum esse militem Christi: vel agricolam terrae:vel Angelicum hominem, non modo rem non explicant: sed quod experientia huius aetatis demonstrat, mentiuntur, cùm nullus hodie Monachus sit huiusmodi. Vtcunque maxime superstitiosum vitae genus est, et a Turcis etiam hodie in blasphema ipsorum doctrina semper retentum, quia fuco faciendo hoc genus hominum est aptissimum. A Iudaeorum Essonis, vt multi existimant, primum introductum fuit, et institutum: et inde postea ad Christianos translatum a superstitiosis hominibus primùm Ægyptiis, deinde Syriis: postea Asiaticis, demum Graecis et Europaeis, ne vlla pars orbis tanti mali expers esset, et immunis. Certe post Christum natum Anno 300. maxime in Ægypto commendari coepit. Fuit enim ea regio, vti Arrianismi, et Meletiasmi, maximarum nereseom:ita etia Monachismi satrix et ferax aut siue propter solitudines, quas hoc hominum genus ferum prorsus et agreste captat: siue propter innatam illi genti superstitionem praeter caeteras mumdi nationes. In occidente hoc hominum genus fuit posterius et notum et receptum, et demum sub Vitaliano Pontifice Roman. coeperunt Monachorum coenobia erigi. Sed tamen, vt externis  \pend
\section*{CAPVT  IIII. }
\marginpar{[ p.215 ]}\pstart rebus homines fere faciliusque capiuntur, quam vera sanctitate: sic hoc nouo hominum genere exorto et terris monstrato, omnes miro stupore affici, nouum hoc vitae genus in coelum tollere, et cùm ad scripturae regulam non examinarent, (quod tamen Basil. Magnus vitae Monasticae praeco fieri vult) solum beatum praedicare coeperunt. Hinc laudes horum hominum immensae extiterunt, et hyperbolice decantatae, etiam a priscis scriptoribus quibusdam, qui Basilium sunt sequutiNam Iustini Martyris, Tertulliani, Cypriani, Irenaei, Ignatii seculo hoc monstrum hominum nondum natum in orbe Christiano fuerat, quicquid Eusebius ex Philone narret: cùm Philo non de Christianis, sed de Essaenis Iudaeis loquatur, quae res non fuit, vt alsae mustae, ab Euseblo obseruata, quemadmodum praeclare Chemnitius in Examine Tridentini concilii probat, quicquid Dionysius ille consecrationis Monachorum in Hierarchia meminerit.  \pend\pstart Quidam hoc hominum genus ex persequutionibus, iis praesertim, quae saeuissimae contigerunt sub Maximino, et postea Diocletiano, exortum censent, quod cogerentur pii homines in vastas solitudines confugere, vt persequutorum immanitatem euitarent. Et certe haec causa erat probabilis nisi ἔξ ἀιρίβως potius animi, quam ex vlla necessitate appareret homines hoc sibi iugum imposuisse, et eo tempore, quo nullae erant persequutiones: sed summa pace fruebatur Ecclesia Deinde nisi ipsa historia Ecclesiastica eo tempore, quo Monachismus adhuc pene recens erat, scrip ta, aliud diceret, quae refert hoc institutum vitae  \pend
\section*{AD I. PAVL. AD TIM. }
\marginpar{[ p.216 ]}\pstart ad quendam Ammonium virum Ægyptium et coniugem. Eremitatum enim author fuit Salas, vt ait Nicephorus, quanquam Isidorus in liber  de Officiis non tantùm a Ioanne Baptista: sed etiam ab vsque Elia Propheta huius vitae primordia repetit, tam est bonus Monachorum patronus. Certatim igitur omnes ingenii sui aciem in hoc vitae genus extollendum contulerunt. Alii enim, vt Ambrosius epist. 82. Monasteria sçholas et officinas virtutum appellarunt, Chrysostom.dixit, Vt in coelo sunt astra:sic in terra esse ordines diuersos Monachorum. Alius dixit, vt e paradiso terrestri quatuor flumina fluebant, sic ex Ecclesia quatuor Monachorum ordines extitisse. Alius, velut Hieronym. et Bernard. susceptionem Monachismi vocat seçundum Baptismum, per quem remissio omnium peccatorum obtineatur. Alius Monachismum esse dixit Angelicam vitam, vt Thomistae omnes, quo sibi maiores in vitae suae errore delicias faciant. Nos et ex superstitione primum inuentum esse hoc vitae institutum defendimus: et mere esse Diabolicum, cùm pars quaedam cultus Dei, et vitae sanctitatis in hoc externo rerumgenere posita sit, et praeter Dei verbum falutaris, et beata vita definita.  \pend\pstart Etsi enim admittimus minus fuisse olim corruptos Monachorum mores, quam nunc sunt, tamen haec ipsa vitae institutio fuit a Christi vita prorsus aliena, et superstitionis plena, et aulla alia de causa, nisi ex ἐθελοθρηβεία et idololatrica rerum externarum admiratione inuenta, et fuscepta. Admitto igitur quod de Monachis suae aetatis scribit Chrysost. Homill. 14. in Timoth.5.  \pend
\section*{CAPVT  IIII. }
\marginpar{[ p.217 ]}\pstart his verbis, quae recito vt appareat, quantum nostri Monachi ab illis veteribus differant. Verae, ait ille, domus luctus sunt Monasteria, vbi cinis atque cilicium: vbi solitudo, vbi risus nullus, nullus negotiorum secularium strepitus, vbi ieiunia, vbi terrenorum duritia lectulorum, vbi munda sunt omnia. Nullus ibi nidoribus locus, nulli sanguinis riui, etc. quae quantum ab Monastica vita nostri seculi differunt vident omnes. Et Bernardus in Apolo. ad Guill. Abbatem verissime multa de luxu Monachorum et ἀσωτία dixit, quae hodie, vti alia vitia omnia, in iis maior est et intolerabilior, quam vnquam: imo vero prorsus insanabilis et inueterata. Sed et hoc vitae genus viris quibusdam bonis statim displicuit, et postea Synodis vti Gangrens. et Matisconensi, improbatum est et velut superstitiosum damnatum. Modum vero non habuit. cùm quotidie noui ritus, inauditaeque  superstitiones inducantur, vt nouus Monachorum ordo excudatur, veluti Stilidae olim:nunc Iesuitae, adeo vt non si mihi centum sint linguae, sint ora que centum, ferrea vox, possim eorum varias comprehendere formas. Sed et illud immane, quod inter eos ipsos huiusmodi genera reperiuntur, quae humanam exuant naturam prorsus et in belluinam mutent qualia sunt haec Armenta Bασκοὶ et ἀπαθεῖς vt est apud Nicephorum liber  14. cap.50. Tamen nunquam istae laudes tantùm obtinuerunt, vt propterea reliqua Christianorum genera, quae nihil cum ista ἀυθαδεία et vitae genere conmune habent, minus ad Christum pertinere censerentur. Dixit Paphnutius Episcopus in Nicaena Synodo, in omni ordine, et non in solo Mo\pend
\section*{AD I. PAVL. AD TIM. }
\marginpar{[ p.218 ]}\pstart nachatu, posse homines placere Deo. Bernard.in in sermo.in Psal. Audiam:tam coniugatos, quam Monachos esse faluandos aisi mat, et idem in sermone 66.in Canticum cantic.  \pend\pstart Eorum genera fuerunt primùm duo tantùm, Eremitarum, et Coenobitarum, vt ait Sozomen. libus 3.cap.16. Quae multiplicata postea, vt docet Euagrius liber 1 cap.21. in quatuor excreuerunt, vti etiam ex Cassiano apparet Collat. 18.cap.4. Alii enim fuerunt Coenobitae, alii Anachoretae, alii Sarabaitae, alii Rhomoboitae, vt etiam docet Hierony. Coenobitae qui sub vno capite in coetu viuunt: Anachoretae qui in coenobiis primum instituti, seorsim postea vel singuli, vel bini in solitudines secedunt, ibique habitant, non in coetu. Horum author fuisse dicitur ab aliis Paulus quidam, ab aliis Antonius. Quod vitae genus plane ferum est et agreste et ab ipsa hominis natura et definitione dissentaneum, qua dicitur homo esse ζῶον πολιταὸν, et praeclare M. Tullius docet ne si quidem alicui omnia virgulta diuina suppeterent, optaturum tamen eum vt viueret in solitudine et procul areliquis mortalibus more ferarum, quod isti Anachoretae faciunt. Sarabaitae in cellulis habitant, et Ana choretis sunt peiores. homines vagi ce aduri qurihirgypto praesertim fuerunt. Rhomoboitae qui terni habitant in cellulis suo arbitratu viuentes, effrenes, et nulli capiti parentes. Hi pessimi omnium vt in epist. ad Eustoch. testatur Hierony. qui eos vidit Stilidae siue Columnarii a Simeone instituti, vt ait Euagrius, qui in columnis habitant, vt nec in terra, nec in coelo viuere videantur. Socrat. liber 4.cap.23 4.haec gene\pend
\section*{CAPVT  IIII. }
\marginpar{[ p.219 ]}\pstart ra enumerat Gnosticos, Coenobitas, Practicos, Eremitas, vt varium fuisse hoc totum genus facile appareat. Hi etiam ordines et distinctiones posterioribus temporibus sunt mutatae, et aliis nominibus appellatae. His enim successerunt quatuor alii Monachorum ordines, nempe Benedictinorum, Augustinianorum, Carmelitarum, Praemonstrensium, qui postea se in infinitas propagines diffuderunt, Mendicantium (quos verius alii fratres Manducantes vocant) et non Mendicantium, et variae distinctiones inter eos inductae a colore vestis, a genere cibi, a forma calceorum, ab amplitudine manicarum, a forma pallii, et huiusmodi aliis nugis, quaineu luo mnt portu cunce lanitis hostrae causae. Antea tamen tres tantùm ordines Monachorum enumerabantur scilicet, Sancti Basilii, Sancti Augustini, et Sancti Benedicti, vt est in cano. Pernitiosam 18. quaest 2.  \pend\pstart Anno vero Domini 500. ista non apum, sed fucorum potius examina varie et sparsa sunt et multiplicata per totum terrarum orbem, imprimis autem per Palaestinam, Cappadociam, Asiam, Graeciam, et Ægyptum, tanquam Apocalypseos locustae, sese diffuderunt, et praecipue in Sceti, siue Staitide Ægypti Nomo siue praefectura, vt Sozomenus liber 6.cap,7. tradit.  \pend\pstart Monasteria vero in vrbibus etiam furiose erigi coeperunt Anno Domini 700. et postea Velut autem scholae quaedam aedificari primùm coeperunt, et velut officinae literarum in summa totius mundi barbarie: quae per Gothos, Vandalos, Hunnos, et Erulos inducta fuerat. Vnde Iuo scribit in Iustiniano Florum quendam Gallorum  \pend
\section*{AD I. PAVL. AD TIM. }
\marginpar{[ p.220 ]}\pstart Francorum procerem prope Aurelianensem ciuitatem coe nobium sancti Benedicti condidisse, vbi Maurus ipsius Benedicti discipulus, vt ait Sabell. Ennead. 8 lib 2. docuerit, et Floriacenum Monasterium adhuc hodie appellatur. Certe ab hoc vitae instituto prorsus nunc nostri Monachi degenerarunt, quorum coenobia, non sunt literarum officinae, sed harae porcorum, et pascua hominum cantillando viuentium, more cicadarum. Tamen postea vsurpatum est, vt poenae loco in Monalteria detruderentur nomies inertes, scelerati et immorigeri, et tonderentur, vt ex historiis praesertim Francogallorum apparet, et Gregorio Turonensi.  \pend\pstart Quis autem vsus fuerit istorum hominum varie quesitum est, quae item regulae.  \pend\pstart Ac regulae vt Nicephorus ait liber 9. cap.14. hae primae erant, vt omnia essent libera, votum autem perpetuum ista professio non haberet: adeo vt nec ciborum, nec ieiuniorum certa tempora et indicta essent, sed quantum vellent ederent, et biberent, et operarentur: item quando vellent, ieiunarent. Atque ita Pachoimus primas regulas instituerat. Sequuti sunt alii qui restrinxerunt illam libertatem. Ac primum in vestitu: post autem in cibis. Qui primus autem institota ista scrupulosius et minutius vel explicasse vel tradidisse censetur, est Basilius, siueiis est, qui Magnus dicitur: siue alius, eóque recentior. Ergo quoad habitum, cucullis sunt vsi sumpto more ab Ægyptiis Monachis, vt Cassia.scribit. Melotis item. i. ouilibus vestimentis, et vilibus pannis, Colobiis lineis ad imam vsque cubiti partem pertingentibus, vt expeditas ad  \pend
\section*{CAPVT  IIII. }
\marginpar{[ p.221 ]}\pstart operandum manus haberent, Cilicio, et Succinctoriis siue subligaculis, et pelliceis, quae lineo colobio operiebantur, vnde superpellicea dicta, Gal. des superplis. Colores in habitu varios habuere, et gestauere: calceamentorum nullus inter cos, vtpote in calidis regionibus primum habitantes, fuit vsus, vt Cassianus docet. Quoad ipsam autem corporis eorum formam, comas alebant et barbam resecabant, contra quam tamen Epiph.haeres. 8o. decorum viris Christianis esse censet. Barbam enim μορφην ἀνδεος vocat, et in concilio Carthaginensi clerici barbas radere veriti sunt: sed non sunt Monachi de numero clericorum. Accincti semper lumbis incedebant, et pecunias cumusare vetabamtur. Nupins stem et festiuitatibus interesse prohibebantur. Operari autem iubebantur quotidie, id est, in agro aliquid fodere, aut arare, aut ferre, vt victus vel ipsis vel toti sodalitio in quo viuebant, superesse posset.  \pend\pstart Praecepta quoque in ciborum delectu et genere contra Pauli regulam illis instituta sunt. Alii enim herbas tantùm, alii oua, alii pisces: alii omnia comedere permittuntur. Ieiunia tamen quaedam illis ad certam mensuram et dies certos praescripta fuerunt: et mirum silentium, quod non est hodie in nostris Monachis, cùm nusquam inuidia et murmur maius frequentiusque locum habeat, Denique nulli Monacho esse sine cella licuit, ne velut piscis siue aqua viuere existimaretur. Hae sunt nugae, sed potius hae sunt superstitiones et idololatricae caeremoniae, in quibus hoc vitae institutum, et Christiana (si Deo placet) sanctitas fuit tandem, abolitis verae sanctimoniae praeceptis, collocata.  \pend
\section*{AD I. PAVL. AD TIM. }
\marginpar{[ p.222 ]}\pstart Quasi illa in his externis etἀδιαφόρσις rebus insit: quae in mentis renouatione, vetefis hominis depositione, et noui susceptione posita est, in fide Christi, in charitate proximi, in omni internae et externae immunditiae expurgatiohe. Sed ita homines sibi praestigias istas fecerunt, vt aliud quid esse sanctum, quam istam viuendi rationem non arbitrarentur. Denique isti homines et seditionum in Ecclesia magnarum authores fuerunt, vt Euagrius liber 3.cap 32. totaque historia Ecclesiastica docet: et pessimorum errorum tum ipsi inuentores, tum propagatores, veluti Messalianorum, Antropomorphitarum, Pelagianorum, Iudaeorum, Abelonitarum, Samaritanorum. Quod non tantum teitatur I'neodorit.nb reap.1 Histor. sed ipsa rerum veritas et experientia: cùm nulli alii sint hodie Antichristi, totiusque Papisticae perfidiae maiores patroni, et fulcimenta, quam tota ista Monachorum factio et coitio.  \pend\pstart Sed quaesitum est, cuius in Ecclesia vsus fuerint isti homines. Respond. Certe nullius, cùm nec Presbyteri, nec Diaconi, nec Episcopi essent. Itaque  de se vere illi istud Horatianum dicere possunt. Nos numeri sumus. et fruges cousumere nati. Proci Penelopes. nebulones, Alcinoique.  \pend\pstart Mionaeni officia populo celebrare non possunt. Monachus officium plangentis habet, non docentis. Ergo nullum munus Ecclesiasticum gerunt, sed nec gerere, quatenus Monachi sunt, debent aut possunt, non concionari apud populum: non Sacramenta administrare. Claret enim et certum est, ait Bernardus, quod publice praedicare Monacho non conuenit, nec Sacramenta admini\pend
\section*{CAPVT  IIII. }
\marginpar{[ p.223 ]}\pstart strare, quia ne clericorum quidem loco et numero censentur. Nullam enim in Ecclesia vocationem legitimam habent Monachi, vt est in cano. Alia est causa et cano. Nemo potest 16. quaest.1. et apud eundem Bernardum epist.42. labor et latebrae, ait, et voluntaria paupertas, haec sunt Monachorum insignia, haec vitam solent nobilitare Monasticam. At, Audite nostri seculi Monachi, vestri oculi omne sublime vident, vestri pedes omne forum circumeunt: vestrae linguae in omnibus audiuntur conciliis: vestrae manus omne alienum diripiunt patrimonium. Quod si Monachorum sceleratissimam vitam pingere aut narrare vellem: nulla mihi Ilias satis esset. Sed quaero, cuius primum rei et vtilitatis gratia hoc vitae et hominum genus institutum, inuentum, et productum fuit? Respondeo ego quidem, Nullius. Nam in canon. Generaliter 5 Hoc idem 16. quaest. 1. Abbatem quemlibet et cuiuslibet Monachorum coetus praefectum simplici ostiario Ecclesiae postponendum esse et subiiciondum aperte responsum est, durauitque vsque ad Syricii Papae spurcissimi tempora, vt Monachi Monachi simpliciter manerent, clerici autem nullo modo fierent, non praedicarent, non celebrarent officia in Ecclesia. Postea vero paulatim gens ista, instar locustarum, e suis desertis, et e suo coeno velut ranae emergere coeperunt, et tandem effecerunt, vt speciali priuilegio (sed toti ordini a Bonifacio quarto, vt alii volunt: vt autem alii, ab Innocentio secundo concesso) quilibet Monachus Sacerdos et sacrificulus fieri possit et Doctor, et praedicator, et Sacramentorum administrator de consensu Abbatis, et Episco\pend
\section*{AD I. PAVL. AD TIM. }
\marginpar{[ p.224 ]}\pstart pi, et curionis, in cuius paroecia ista agit. Ex quo postea factum est, vt in maximas Ecclesiae Papisticae dignitates, veluti Episcopatus, Archiepiscopatus, Papatum denique ipsum, isti crabrones et locustae inuolarent, et iam ipsius Papismi solae columnae restarent. Quo hominum genere nihil habet vniuersus orbis pestilentiosius, superstitiosius, impudentius, audacius, crudelius, ambitiosius, aut denique magis auarum et rapax. Hoc autem omne vitae institutum quam procul absit a Pauli praecepto, et ἀσκύσει, cuius hic memtionem facit, vident omnes: imo vero vt nihil quicquam cum Dei verbo consentaneum habeat sed potius mortis Christi beneficium, et nostrae coram Deo iustificationis causam (quae in vno Christi sanguine est) ad cibos et colores vestium merasque superstitiones et blasphemias transferant, in hac Euangelii luce intelligunt omnes, vt mihi non duxerim necesse diutius ista refellere, quae vel indicasse satis est ad refutasse. Redimus igitur ad Paulum nostrum.  \pend
\phantomsection
\addcontentsline{toc}{subsection}{\textit{9 Fidus est hic sermo, et dignus qui omnibus modis recipiatur. }}
\subsection*{\textit{9 Fidus est hic sermo, et dignus qui omnibus modis recipiatur. }}\pstart Commendatio superioris doctrinae ab vtilitate et illius praestantia. Tum enim fidelis est haec doctrina: tum vero omnibus vtilis, qui illam audiunt, illique  credunt. Praecautio quoque  et praemunitio quaedam haec esse videtur aduersus varias animi dubitationes, quae ex superiori doctrina solent oriri. Tantum.n. abest, vt fideles felices iudicentur, vt et miseri et probrosi passim et vulgo esse censeantur. Quod iudicium falsum esse pronuntiat hoc loco  \pend
\section*{CAPVT  IIII. }
\marginpar{[ p.223 ]}\pstart Spiritus sanctus et seq. versiculo. Quae sit autem sententia huiue oracomo uoproapehv. 15. fusius ex posuimus.  \pend
\phantomsection
\addcontentsline{toc}{subsection}{\textit{10 Nam idcirco etiam fatigamur, et probris afficimur, quod speremus in Deo viuo, qui est conseruator omnium hominum, maxime vero fidelium. }}
\subsection*{\textit{10 Nam idcirco etiam fatigamur, et probris afficimur, quod speremus in Deo viuo, qui est conseruator omnium hominum, maxime vero fidelium. }}\pstart Occupatio est, quae per digressionem iniecta est, ne superior illa pietatis commendatio falsa esse censeatur, propter varios rerum euentus et duram piorum sortem, quae in hoc mundo prorsus misera esse apparet et afflicta, tum in corpore ipso, tum in honore. Respondet autem et negat eam sortem Paulus propterea esse miseram et infelicem existimandam, idque respond. propter Causam propter quam patiuntur, pii: et Exitum, qui cernitur, et inde speratur. Acfinis siue causa, propter quam pii affliguntur, praeclarissima est. Est enim pietas et ipsa doctrina coelestis, et spes in Deo viuo, non in idolis, qua nihil melius.  \pend\pstart Exitus etiam laetus et felix est. Id quod probat argumento sumpto a minori ad maius. Deus enim, qui omnium hominum est conseruator, multo maxime suorum et fidelium custos est et liberator. Haec est summa totius doctrinae Pauli. Disputat Chrysost. atque ambigit, vtrum vox συτὴρ hoc loco referri ad vitam, et salutem aeternam debeat. Sed hic non potest. Neque enim omnes homines vitam aeternam consequuntur, sed soli fideles, quanquam tamen omnium est custos, conseruator, et benefactor Deus, vt est Ps.36. quatenus nempe omnium hominum est creator et opifex.  \pend
\section*{AD I. PAVL. AD TIM. }
\marginpar{[ p.226 ]}\pstart Spem autem coniungit cum doctrina Paulus, cùm ait (quod speremus in Deo) quia haec cognitio, quae Per fidem et Euangelium a nobis habetur, non est euanida quaedam phantasia de Deo, sed est huiusmodi agnitio, quae spem, eamque  certam in Deo nobis afferat, et inserat. In quo differt a reliquo cognitionis genere. Nam aliae sunt tantùm in animo et intellectu versantes cogitationes, neque cor hominis mouent, afficiunt et mutant, haec vero etiam in corde ipso est certissima fiducia, quod afficit et inclinat ad Deum amandum, eiusque promissiones certissime sperandas et expectandas, vt est in cap. 5. Epistolae Pauli ad Romanos. Sed et vox naec speramus pro credimus posita videri etiam hoc loco potest.  \pend
\phantomsection
\addcontentsline{toc}{subsection}{\textit{11 Haec denuntia et doce. }}
\subsection*{\textit{11 Haec denuntia et doce. }}\pstart Commendatio superioris doctrinae ab vtilitate, quae ex illius praedicatione colligitur. Est enim haec vtilis doctrina, cuius praedicatio est sic a Paulo praecepta, vt et saepius inculcetur, et Pro imperio et potestate ministerii Euangelici proponatur, vti Dei ipsius vox. Ergo comparatione quadam superiorem doctrinam magni esse in Ecclesia momenti docet probatque. Nam etsi alia capita doctrinae doceri a pastore debeant: hoc tamen doctrinae caput, per quod falsus Dei cultus oppugnatur, inanes caeremoniae et superstitiones euelluntur, saepissime et frequentissime repeti, doceri, ac inculcari debet. Atque hoc faciendum esse affirmat Dei spiritus, id est, doctor ille, qui optime nouit, de quibus sedulo et saepius  \pend
\section*{CAPVT  IIII. }
\marginpar{[ p.227 ]}\pstart edocendi et commonefaciendi sumus. Quo minus nobis molestum aut fastidiosum esse oportet, si saepe saepius in Ecclesia hanc eandem doctrinam a pastoribus audiamus, quia maxime nobis est salutaris. Paulus Philip.3.vers.1. dicit tutum et vtile esse, vt eadem nobis saepius repetantur. Idem docet et affirmat Petrus in 2. Epist. cap.1. vers.13. Vnde vae delicatis hominum auribus, qui abusus doctrinae et superstitiones Papisticas nunquam a concionatoribus Euangelii taxari vellent.  \pend\pstart Denique ex hoc splo apparet, eam maxime a pastoribus doctrinam proponendam, quae necessaria est, et vtilis: non autem quae inutiles et vanas quaestiones habeat.  \pend\pstart Colligit vero ex hoc Pauli dicto Chryso. duplex esse pastoris munus in proponendo Dei verbo pro rerum, quas tractat, et personarum, apud quas agit, ratione, nimirum vt Doceat, et Ea quae sunt obscura, et Eos, qui sunt obsequentiores. Item Denuntiet, propotestate et imperio sibi a Deo dato, Ea, quae omnes concedunt esse vera, Eis, qui sunt contumaciores et magis rebelles. Quod certe verum est. Debet enim pastor nonnunquam increpare et virga.i. seuetitate vti, ve ait Paulus 1. Corinth 4.vers.21. et 2. Timoth.4.  \pend\pstart 12 Nemo tuam iuuentutem despiciat, sed esto exemplar fidelium in sermone, in conuersatione, in charitate, in spiritu, in fide, in puritate.  \pend\pstart Νοθέτησις. Admonet enim qua ratione et grauitate superior doctrina sit a Timotheo propo\pend
\section*{AD I. PAVL. AD TIM. }
\marginpar{[ p.228 ]}\pstart nenda ac illi ipsi sit in omni ministerii Euangelici quod obibat ratione conuersandum et agendum. Adeo vt etiam haec praecepta ad reliquos pastores pertineant, qui nobis in Timotheo, tanquam in perfecti ministri idea, a Paulo depicti proponuntur. Est vero hic locus etiam tacita regressio ad superiùs institutam disputationem de officio veri et Euangelici pastoris, quam et hoc, et sequenti capite persequitur. Haec est igitur ratio ordinis, ne putemus hoc scriptum Pauli esse scopas dissolutas. Quod Athei obiiciunt. Summa vero huius doctrinae eo pertinet, vt doceantur pastores superiorem doctrinam reuerenter et grauiter tractandam esse: et in reliquis ministerii et muneris suipartibus ita illis esse conuersandum, vt nemo eos, quantunuis aetate iuuenes, despiciat vel contemnat: sed sint totius gregis exemplar in Verbis, et factis, et in Doctrina, et in vita, Videri autem potest non ad omnes in vniuersum verbi Dei ministros pertinere haec Pauli exhortatio: sed ad eos tantùm, qui adhuc iuuenes sunt: neque dum sibi per aetatem eam authoritatem compararunt, per quam sint populo commendabiles et venerandi, quales sunt adolescentes adAuc pastores, et nondum prouectae aetatis. Sed cùm etia iiilius fionorentur, qui, quanuis aetate grandaeui, moribus tamen sunt iuuenes, certe ad senes etiam haec praecepta pertinent, ne hi ipsi in ea aetate propter morum dissolutionem contemnantur, si idem agant, quod homines moribus iuuenes et proterui, et lasciui et leues. Quae certe vitia adolescentiam sequi et comitari solent.  \pend\pstart Dicit Paulus reóreroe, id est, Iuuentutis. Differt  \pend
\section*{CAPVT  IIII }
\marginpar{[ p.229 ]}\pstart haec aetas, a Pueritia, Adolescentia, Virili maturaque aetate, et Senectute. Est enim iuuentus post adolescentiam, et ante virilem aetatem veluti a decimo septimo anno aetatis ad vicesimum quintum, intra quos annos existimo tunc fuisse Timotheum: quemadmodum adolescentia est post pueritiam, et ante iuuentutem a septimo anno aetatis ad decimum septimum. Sic Pythagorici distinxerunt vt apparet ex Laertio liber 8.in vita Pytha. gorae etsi plures annos cuique aetati tribuit. Solon Atheniensis in epigrammate distinguit hominis totam vitam in septem aetates. Sed Ioannes in prima Epistola cap.2.vers.13. et 14. tres tantùm constituit. Nos nostram distributionem sequimur, ex qua curus tunc atatis iuerit Timotneus intelligi potest, et quantis in ea aetate Dei donis, et iis quidem summis et egregiis, ornatus fuerit. Id quod etiam apparet ex secunda Epistola ad eundem et ad Philip.2.vers.18.  \pend\pstart Remedium autem per subiectionem affert Paulus, per quod efficiatur, ne vlla pastoris aetas contemni vel possit, vel debeat. Est autem si illi se exemplar caeterorum fidelium praebeant in vita et doctrina. Igitur ex hoc loco colligitur iuuenes eligi posse in ministros, si iis Dei donis excellant, quae ad tantum munus requiruntur. Supra cap.3 vers.6.diximus, quae aetas olim cuique muneri Ecclesiastico exercendo praescripta fuerit a Synodis vt hic repetere nihil sit necesse. Sed apologo lepidissimo docet Æsopus non semper canitiei et barbae inesse sapientiam.  \pend\pstart Describit autem fusius Paulus illa duo, quae diximus, id est, quibus in rebus pastores gregis ex  \pend
\section*{AD I. PAVL. AD TIM. }
\marginpar{[ p.230 ]}\pstart emplar, vt ipse cum Petro 1.epist. cap.5. loquitur esse debeant. Ac primum summatim quidem inquit, ἐν λόγο id est, Verbis siue dictis, ἐν ἀναςροφη id est, Conuersatione et factis, toti gregi praelucere pastorem oportere. Postea vero quasdam virtutes addit, ex quibus distinctius haec duo parantur et docentur: et tum vitae, tum ipsius linguae vera reformatio oritur, tanquam ex ipsis fontibus. Obiter dixerim hoc genus dicendi esse ἀσόνθετον, quod maiorem vim habet, et facit, vt singula diligentius perpendantur a nobis. Ergo inter cas virtutes praecipue commemorat Paulus, Charitatem, cui subiicit Spiritum siue zelum. Charitas autem hic tum ea, quae in Deum, tum in proximum est, significatur, cui subseruire debet spiritus ipse, id est, feruor animi: atque illa etiam dona, quae Dominus cuique fidelium donat et concedit, charitati ancillari oportet. Deinde addit, charitati suum fontem, nempe Fidem, cui subiicit Puritatem siue Castitatem. Ergo fides haec non tantùm est doctrina, vel cognitio: sed etiam ipsa in Dei promissiones per Christum acquiescens assensio et voluntas, quae parit ex se tanquam vberrimum fructum ἁγνείαν, id est, castitatem. Haec autem castitas non tantùm abstinet alieno thoro, et libidine lasciua: sed omni vitae impuritate, quae Christianae sanctitati contraria est. Ex his igitur quatuor fontibus pene omnis vitae fanctitas reformatio et exemplum nascitur. Ergo vanae illae Episcoporum commendationes ex insulis, lituo siue pedo chirothecis, annulis, cothurnis, sandaliis, et caeteris huiusmodi nugis facessant, quae in  \pend
\section*{CAPVT  IIII. }
\marginpar{[ p.231 ]}\pstart insignia, ornamenta, et testimonia dantur, et soIenniter gestantur.  \pend
\phantomsection
\addcontentsline{toc}{subsection}{\textit{13 Interim dum uenio, attende lectioni, exhortationi doctrinae. }}
\subsection*{\textit{13 Interim dum uenio, attende lectioni, exhortationi doctrinae. }}\pstart Α’κολύθησις est. Nam sequuntur variae exhortationes, quas muneris Ecclesiastici dignitati consentaneas esse vidit Paulus, et necessario addendas. Ex his autem quantum sit hoc onus, quanque sedulo capessendum et diligenter exercendum intelligere certe debemus, ne segnes socordesue ipso Episcopatus titulo, vel fastigio turgidi et insolentes, vel etiam reditu et bonis contenti torpeamus. Sunt enim hae cohortationes, tot Spiritus sancti aduersus pastores negligentes et muneris sui grauitatem non apprehendentes fulmina vel tonitrua. Neque enim tam Timothei gratia dicuntur ista, quam omnium pastorum, quorum verissimam et optimam ideam nobis in Timotheo format et pingit. Denique si tam ardens in Dei metu et tantus vir Timotheus iis stimulis eguit, quid nos? et quot quantisque exclamatiomious et puneturis smus cxcitanun Est autem hic exhortationis et admonitionis locus non vnus, sed quadruplex. Omnes quidem eo tendunt, vt sedulo et diligenter pastores suo muneri intendant et incumbant. quilibet tamen locus exhortationis habet quoddam peculiare et necessario considerandum.  \pend\pstart Prima igitur haec est exhortatio et admonitio, attendendum et assidue versandum Ministro Euangelico in Sacrae scripturae lectione. Vox enim, lectionis, generaliter hoc loco enuntiata ad Sacrae ....  \pend
\section*{AD I. PAVL. AD TIM. }
\marginpar{[ p.432 ]}\pstart scripturae tamen lectionem specialiter restringi debet, quia non simpliciter hoc requirit a pastoribus Paulus, vt aliquid legant quodcunque  tandem sit, sed vt facrum Dei verbum euoluant. Hanc esse Pauli mentem apparet ex fine, propter quem attendendum lectioni docet, nimirum vt pastor Doceat gregem suum, et Exhortetur, quae sunt duae muneris Ecclesiastici partes praecipuae. Doctrina, tum in Verorum dogmatum traditione, et Falsorum refutatione posita est. Quorum neutrum nemo facere, nisi apprime et abunde verbo Dei instructus et diligenter in eo versatus potest, cùm ex eo vno fonte et veritas Dei fit haurienda, et mendacium illis armis oppugnandum, illóque gladio iugulandum Actor. 18.vers.28.  \pend\pstart Exhortatio est in Bene iam currentibus confirmandis, Improborum acri reprehensione: denique in frigidis excitandis. Illi autem stimuli:non nisi ex Dei promissionibus et minis, addi possunt. Haec igitur omnia Dei verbo continentur, quod idcirco et versandum et meditandum est diligenter. Ergo non satis est Episcopos et pastores quaedam legere ( nam nec fabulas Poëtarum, nec vanas et inutiles Philosophorum scriptitationes terere debent, aut iis immorari) sed oportet eos Dei verbum scriptum assidue contrectare. Vt autem ad Pharisaeorum scripta et scholas non remittit suum Timotheum Paulus, neque nos etiam ad Scholasticorum et Sophistarum Sorbonistarum Academias reuocandi sumus, vt veri Theologi euadamus: sed ad sacrum Dei verbum et ad eos scholas, vbi illud pure docetur.  \pend
\section*{CAPVT  IIII. }
\marginpar{[ p.233 ]}\pstart membranas, in quibus aliquid descripserat priuatae memoriae causa, sibi Troade afferri vult Similis locus 2. Pet.1.V.19. Dissimilis Mat. 10.v19. Respo. Christum quidem nolle pastores esse in legendo Dei verbo negligentes: terreri tamen nos non vult vel conspectu vel doctrina vel fama nostrorum aduerfariorum, cum ipse nobis sit suo spiritu potentissime affuturus, et quae ad conuincendos veritatis Dei hostes necessaria sunt suggesturus.  \pend\pstart Addit Paulus, donec venio. Quae limitatio non excludit tamen reliquum tempus, quasi praesente Paulo Dei verbum contemnere Timotheus debuerit. Nec facit pastorum, et praeceptorum nostrorum magna et excellens doctrina, vt in Dei scriptura euoluenda fiamus desidiosiores: sed contra nos magis illorum egregiae dotes debent excitare, vt quae illi proponunt doctissime, ad exemplum Thessalonicensium, cum scripto Dei verbo conferamus, diligentiusque et lubentius a tantis viris ediscamus.  \pend\pstart Denique quod nullum tempus praescribit Paulus, ac ne longissimum quidem, post quod cessare debeat a studiis Timotheus, ostendit nos per totum vitae tempus proficere posse: tamen non facit, vt videatur inepta çertorum annorum certique spatii definitio, intra quod tum Theologiae: tum aliis artibus et facultatibus sit studendum, operaque ponenda, quanquam potius ex ipso cuiusque profectu id diiudicandum est, qui apti sint ad ministerium vel gradum adipiscendum, quique  plurimùm profecerint, quam ex ipso temporis curriculo et spatio, quo versantur in Academiis. Quod autem quidam putant, et ex hoc loco col\pend
\section*{AD I. PAVL. AD TIM. }
\marginpar{[ p.254 ]}\pstart ligunt fuisse scriptam hanc epistolam, cùm Paulus in Syriam postremum venturus esset, non video nauere lacio hihum argumentum et fundamentum.  \pend
\phantomsection
\addcontentsline{toc}{subsection}{\textit{14 Ne negligito donum quod in te est: quod datum est tibi ad prophetandum, cum impositio ne manuum presbyterii. }}
\subsection*{\textit{14 Ne negligito donum quod in te est: quod datum est tibi ad prophetandum, cum impositio ne manuum presbyterii. }}\pstart Secunda exhortatio ab officio vel fine donorum Dei, quam amplificat etiam a circunstantia, eaque  duplici nempe a Doni Dei excellentia, quo praeditus erat Timotheus, et Prophetia quae de Timotheo praecesserat, quod vtrunque delebit et deturpabit Timotheus, nisi suo muneri diligenter operam nauet et impendat. Eodem vero modo debet quisque animaduertere, quae a Deo dona peculiaria acceperit, ne apud ipsum otiosa et infructuofa lateant, aut inutilia maneant, et defossa: sed vt illa in Ecclesiae vsum communem, propter quem illa talenta data sunt a Domino, conferantur.  \pend\pstart Similis locus Matth.25.vers.14.1. Pet. 4.V.1o.  \pend\pstart Quo autem maiora sunt Dei dona erga nos, eo diligentiores in caeteris lucrandis esse debemus, ne illa apud nos otiosa iaceant, vel rubiginem contrahant. Nam cum foenore a Deo repetentur.  \pend\pstart De Prophetia autem, quae de Timotheo praecessit, quod dicit hoc loco Paulus, etsi quidam sic explicant, quasi finis illius Dei doni ostendatur: non autem vlla insignis praedictio' Dei et testimonium de futuris! in Timotheo donis antegressa sit:ta men malo accipere:vti supra 1.vers.18. et 2.  \pend
\section*{CAPVT  IIII. }
\marginpar{[ p.235 ]}\pstart Timoth.1. vers.6. Explicatum autem est a nobis supra eo ipso cap.1.vers.18,  \pend\pstart ΣάριQμα autem dixit Paulus, non χe Quα, ne quis hebes Papista putet chrisma, quo ipsi superstitiose vtuntur ad suos Sacerdotes et Episcopos vngendos et inaugurandos, hic esse a Paulo commemoratum et quaesitum.  \pend\pstart De impositione manuum infra, vti et de voce hac Presbyterii, quae coetum significat Presbyterorum, contra quam sentit Erasmus: putat enim ille Episcopatum significari hoc loco, et finem designari non instrumentum et modum, quo vocatus est Timotheus: infra dicetur cap.5.vers.22.  \pend\pstart Hic locus commendat authoritatem Presbyterii, et vocationes ordinarias, quae hominum ministerio fiunt, quam qui maxime, ne quis praetextu ingentium donorum et egregiae suae doctrinae in Ecclesiam intrudat ipse sese, cùm Timotheus tantus vir etiam de quo praecesserant spiritus Dei apertae prophetiae, tamen ordinario modo designatus sit minister et praeco verbi Dei.  \pend
\phantomsection
\addcontentsline{toc}{subsection}{\textit{15 Haec exerce, in his esto: vt tuus profectus manifestus sit inter omnes. }}
\subsection*{\textit{15 Haec exerce, in his esto: vt tuus profectus manifestus sit inter omnes. }}\pstart Tertia admonitio de ipsa operis difficult ate, vt ea perspecta diligentius pastores suo muneri incumbant. Habet autem summos stimulos ipsa styli, quo Paulus vtitur, ratio, meditare (ait) in his esto. Vrget enim pastores ne alibi diutius aut lubentius occupentur, quam in hoc munere suo, tanquam in sua Sparta exornanda. Ergo damnantur ex hoc Pauli dicto pastores omnes, qui in iis,  \pend
\section*{AD I. PAVL. AD TIM. }
\marginpar{[ p.430 ]}\pstart quae propria sunt sui muneris, perfunctorie et naτιργως versantur: in aliis autem rebus, quae proprie ad eos non pertinent, serio, assidue et diligentius insunt eaque satagunt. Faciunt enimη̃ργον πάἑσργον: et πάρεργον ἐργον, vt ait Prouerbium vetus, contra exemplum Apostolorum, qui cùm se viderent a vero sui muneris officio propter curam mensarum distrahi et auocari, quanquam ea res et sancta erat et pia, tamen voluerunt se ea cura leuari, vt liberius meliusque veris et potioribus exercitiis suae vocationis inseruirent ac incumberent. Nec valet haec quorundam pastorum, qui sunt πολυπράγμονες excusatio, quod Reipubl. commoditas, et ratio postulet, vt in istis ciuilibus et extraneis rebus ipsi versentur, et occupentur. Facit enim haec ratio, vt hodie quidam Episcopi potius sint militares Duces, vel Praefecti aerarii: vel forenses iudices, vel Cancellarii regni, vel Consili arii, regni, quam pastores, ministri et praecones verbi Dei coelestisque doctrinae, et pabuli gregi suo dispensatores. Denique quid illi Episcopatus est aliud, quam reditus quidam pinguis et opulentus vel ad fruendas voluptates, vel ad familiae onera sustinenda, quae ampla est, interea dum aliis negotiis vacant, quae tamen ab Episcoporum munere, sunt alienissima?  \pend\pstart Μελέτα Meditatio non est leuis quaedam cogitatio, sed assidua et frequens eiusdem cogitatio. nis repetitio, cura, et mentis volutatio, vt aptiores ad id, quod meditamur, fiamus. Sed non solum horum praeceptorum cogitationem requirit Paulus, imo vero etiam executionem et praxim, ne contemplatiuae vitae homines hinc sibi patro\pend
\section*{CAPVT  IIII. }
\marginpar{[ p.237 ]}\pstart cinium quaerant. Ex hoc etiam ipso ostenditur, quam sit difficile et arduum hoc ministerii opus, quocirca diligenter est in illud iincumbendum. Profectus tuus) Est finis quem in munere suo faciendo spectare debet pastor Euangelicus, nenpe propagatio Dei gloriae, et ipsiusmet Ecclesiae mnalor in dies prorectus et incrementum. Itaque vocem, Profectus, non refero ad personam Timothei:sed ad Ecclesiam potius,  \pend\pstart Manifestus sit) Quid si nullus profectus, nulIum incrementum Ecclesiae appareat, estne propterea pastori desperandum. Estne etiam ipse semper propterea, vt deses et negligens damnandus? Respond. Hoc quidem lugendum esse, sine vllo profectu praedicari Dei verbum:si tamen pastor quantum in se est, syncere pascat gregem suum, vigilet, in eo sit totus, et officium faciat, neque illi desperandum esse:neque eum, vt negligentem, damnandum l Pet.5.verl.2.  \pend
\phantomsection
\addcontentsline{toc}{subsection}{\textit{16 Attende tibiipsi et doctrinae: persiste in istis, id enim si feceris, et teipsum seruabis et eos qui te audierint. }}
\subsection*{\textit{16 Attende tibiipsi et doctrinae: persiste in istis, id enim si feceris, et teipsum seruabis et eos qui te audierint. }}\pstart Quarta exhortatio a fructu huiusmodi laboris, et eo quidem vberrimo. Est autem is fructus, salus animarum, tum in Pastore ipso, tum in Grege et Ecclesia. Quod vt fiat vult Paulus vt quilibet pastor Euangelicus, Attendat sibi, et Persistat in officio. Attendat autem, id est, diligenter videat, et non perfunctorie consideret, et Qualis ipse sit in sese et moribus, et Qualis sit ea doctrina, quam docet gregem Domini: vt si quae vitia  \pend
\section*{AD I. PAVL. AD TIM. }
\marginpar{[ p.230 ]}\pstart vel in sese:vel in sua doctrina inesse animaduertat corrigat, et emendet. Denique ea vitia potius a leiplo, quain ab anis difcat, ne videatur negligenter aut perfunctorie officium suum facere.  \pend\pstart Persistat vero, id est, continuet, neque fatigetur aut fathiscat, neue susceptum semel opus propter infinitas pene molestias deserat. Id enim est non modo leuis animi: sed etiam infidelitatis et diffidentiae de Dei misericordia signum verissimum. Ergo perseueret, quantunuis difficile opus hoc videatur.  \pend\pstart Finis tandem et fructus suauissimus, atque adeo salutaris sequitur, salus nimirum ipsius pastoris dum in munere suo obeundo fideliter Deo paret. Habent enim in eo ipso optimum misericordiae Dei erga se testimonium pastores, quia opera bona sunt nostrae adoptionis et aeternae salutis certa testimonia et indicia, praesertim quae in cuiusque vocationis opere implendo et obeumdo apparent 1. Timoth.2. vers.15.  \pend\pstart Ipsius quoque gregis salus consequitur, quia ea ratione et per ministerium Euangelii Dominus adducit suos ad salutem. Vnde Ministri Euangelici sunt salutis hominum ministri et organa. Deus autem ipse est illius author et causa. Eadem phrasi alibi quoque vtitur scriptura, vt Iacob5. vers.21. Prouerbus  23. vers.14. Haec autem magna laus est ministerii Euangelici, per quod seruari dicimur, vti quod operarii cum Deo iidem ministri vocantur 1. Corinth.3.  \pend
\section*{V. }
\section*{CAPVT  }
\marginpar{[ p.239 ]}
\endnumbering\beginnumbering\section{CAP. V.}
\phantomsection
\addcontentsline{toc}{subsection}{\textit{\huge\textbf{S}\normalsize ENIOREM ne increpato, sed hortare ut patrem, iuniores vt fratres. }}
\subsection*{\textit{\huge\textbf{S}\normalsize ENIOREM ne increpato, sed hortare ut patrem, iuniores vt fratres. }}\pstart Μετάβαζις. Transit enim, et ordine quidem Paulus, a doctrina ad mores, cuius vtriusque rei cura ad veros pastores et Ecclesiae praepositos pertinet. Sed prius de vitiis, quae doctrinae ratione obrepunt, et haereses dicuntur, fuit dicendum. Postea autem de morum vitiis, et eorum emendatione, quam hoc loco persequitur, vti et in Epistola posteriori ad hunc ipsum Timotheum, et ad Titum cap 2. Dicitur autem alio nomine haec correptio et emendatio, censura Ecclesiastica, quae potestatis clauium pars est. Sic autem et hoc loco et alibi loqui videtur Paulus, vt vni episcopo eam tribuat, non autem toti ipsi Ecclesiae, non item coetui et senatui Presbyterorum. Id quod perperam intellectum duplicem errorem peperit. Primum eorum, qui totam Ecclesiae ipsius iurisdictionem vnius tantùm personae, id est, Episcopi arbitrio vindicarunt et asseruerunt, ex quo ipso inducta est postea superba Episcoporum tyrannis. Alterum eorum, qui vt hunc errorem corrigerent, ipsi secundum induxerunt, et ad vniuersam, quam vocant, Ecclesiam, id est, ad singulos de Ecclesia hanc potestatem retraxerunt, et ita per malum malo mederi studuerunt, malo nodo vt est in Prouerbio, malum cuneum quaerentes. Denique inciderunt in Scyllam, dum cupiunt vi\pend
\section*{AD I. PAVL. AD TIM. }
\marginpar{[ p.240 ]}\pstart tare Charybdin. De quo toto argumento et ttilissimo, et futurae disputationis sequentiumque toto hoc capite praeceptorum fundamento nobis aliquid est dicendum. Est autem hic status causae et controuersiae (quae etiam non his nostris temporibus acerbe a quibusdam agitatur) Ad quos nimirum Ecclesiasticae censurae ius et potestas pertineat. Hoc enim pro concesso sumimus, quod tota docet scriptura, cùm vitae sancte et honeste instituendae gratia, et vt Dei gloriae inseruiamus, scriptura et doctrina coelestis nobis praedicetur, omnino necesse esse, vt aliqua sit Ecclesiae disciplina et turpium morum emendatio, per quam homines tum in officio contineantur, tum, si semel abducti fuerint, reuocentur. Nam qui Ecclesiasticam disciplinam prorsus tollunt, quid aliud Ecclesiae quam laruam nobis inducunt, et omnem peccandi effrenem licentiam probant? contra quos Pral.0. et Tiebr.12.veri.6. d dietum Tertulliani de Habitu virginum, item Cypriani liber  2.epist.7. valere debet. Disciplina enim Ecclesiae, inquiunt illi, et castigatio Ecclesiastica morum est fidei custos, quam qui repudiant, fidem ipsam nomine tenus habent. Secundo autem loco postulamus et sumimus  \pend\pstart noe nobis concedi, quod est verissimum. Hanc morum correctionem et vitiorum emendationem, quae Ecclesiastica iurisdictio et ius clauium dicitur, latissime differre et toto genere dissimilem esse a ciuili et politica, quam vocant, iurisdictione, quae Magistratibus competit.  \pend\pstart Tertio denique et hoc quoque tanquam ἀιτηua praesupponendum est, quod alibi verum esse  \pend
\section*{CAPVT  IIII }
\marginpar{[ p.241 ]}\pstart probabimus, Deo dante, et iam ex parte docuimus supra cap.3.ex veteri et bene constituta Christianae Ecclesiae institutione, non vnicum Episcopum id est, non vnam quandam duntaxat personam ad regendam totam Ecclesiam electam et ordinatam fuisse (id quod nunc fit in Papatu) sed plures simul delectos, quibus Ecclesiae cura demandaretur, tum in iis quae ad doctrinam pertinent et hi erant Pastores, et Doctores: tum etiam quae ad mores, et Pauperum aerarium, quales erant Pastores iidem, Presbyteri, Diaconi. Ex quo fit vt a Tertulliano in Apolog. cap, 39. Rectores Ecclesiae dicantur non duntaxat soli Episcopi, sed cum illis etiam Diaconi, et Presbyteri, tanquam ipsorum Episcoporum collegae. Quod idem in liber  de fuga in persequutione confirmat, et infinitis pene locis Cyprianus. Atque etiam Augustinus, veluti in Enchirid. cap.65.  \pend\pstart Ergo his ita praemunitis quaesitum est. Vtrum ad vnum tantùm Episcopum, vel delegatum ab eo aliquem, ius censurae et emendationis morum, (quae Ecclesiastica dicitur) eiusque cognitio tota Fertineatean porius au prures. etuou li pertinec ad plures, vtrum ad totum coetum Ecclesiae, an ad singulos fideles qui sunt de eo coetu, an ad eos tantùm qui Ecclesiae praefecti, Authores, ηγομένος, et Praepositi dicuntur, quales sunt Presbyteri et pastores simul, qui alio nomine dicuntur senatus Ecclesiasticus.  \pend\pstart Ac primùm ad vnum duntaxat totam eam potestatem et cognitionem pertinere non posse probat Christi et Pauli dictum. Christus vult peruicaciam rebellium denuntiari et dici Ecclesiae,  \pend
\section*{AD I. PAVL. AD TIM. }
\marginpar{[ p.242 ]}\pstart Ergo non vni tantùm, Matth. 18. vers.17. Paulus quae geruntur in Ecclesia a Presbyterio fieri docet, id est, a quodam coetu hominum. Ergo non a solo illius Ecclesiae Episcopo, supra cap. 4.V.14, Philip. 1.vers.1. Actor.11. vers.30. Deinde Ecclesia ipsa Dei transformaretur in regnum terrenum, essetque iam Monarchia, non Aristocratia, si vnus in ea dominaretur, et omnia solus ageret. Quod fieri et esse vetant hi scripturae loci Matth.2o.V.25 1.Pet.5.vers.3. Hieronym. ad Rustic. penes coetum Presbyterorum, quem vocat senatum morum, censuram esse aperte scribit.  \pend\pstart Denique mos ipse vetustissimus primae Ecclesiae huic vnius dominationi in Ecclesia repugnat. Nam omnia ex communi consensu totius coetus praepositorum Ecclesiae geri solita fuisse, non vnius autem hominis duntaxat arbitrio, probant hi scripturae loci, vbi congregata fuisse Ecclesia dicitur ad res et negotia Ecclesiastica gerenda, et decernenda et dispicienda Actof. 15. vers.4.22. vers.3o. Ioan.9. vers24. Actor.21.vers.18. 1. Cor. 5.vers.4. 16.vers.3.  \pend\pstart Nec in contrarium efficit quicquam eorum ratio, qui omnia, quae tum in coelo, tum in terra fiunt, ad vnum principium, non ad plura reuocari variis exemplis docent, quae fuit ineptissimi cuiusdam Itali et Veneti Cardinalatum ambientis recens pro Papistica et Romana tyrannide in Ecclesia tuenda defensio. Hoc enim verum esse concedimus, et ipsam Ecclesiam vnicum caput habere ipsi contendimus: sed Christum, non autem hominem vllum mortalem, a quo capite, et eo quidem vnico et summo, sunt caeteri, tanquam a sum\pend
\section*{CAPVT  V. }
\marginpar{[ p.243 ]}\pstart mo iudice, delegati, qui pastores in Ecclesia prae ficiuntur. Et vt inter ipsos delegatos nullus sine sacrilegio et laesae maiestatis crimine summi magistratus ius sibi vindicare potest: ita nec in Ecclesia vel vniuersa, vel particulari quisquam ius summi pastoris et ipsius capitis sine aperta in Christum blasphemia nomem sibi tribuere debet. Quod item obiiciunt, olim sub lege fuisse inter Leuitas vnum quendam summum sacrificatorem et Sacerdotem, habet facile responsum. Quia certis figuris sub lege, Christi dignitas, vti Sacrificium adumbrabatur, donec ipse apparuisset, fuisse vnum quendam inter multos fratres summum Sacerdotem, qui Christi figura erat. Itaque fuit ea res legalium caeremoniarum pars, quae omnes Christi aduentu abolitae sunt, et quae sine piaculo a Christianis in vsum reuocari non possunt.  \pend\pstart Nec efficit etiam quicquam haec quorundam obiectio, qui idcirco purant ab vno quodam (cui hae partes sunt totius Ecclesiae consensu demandatae) veluti Episcopo, aut ipsius officiali, censuram Ecclesiasticam fieri, quia quod agit, vt delegatus ab vniuersa Ecclesia agit, non vt priuatus. Delegatio enim ista nunquam vel ab Ecclesia fieri, vel a quoquam suscipi sana conscientia potest, cum alium ordinem ad seipsam regendam instituere Ecclesia ipsa vniuersa non possit, quam illum, quem Christus et ipsius Apostoli praescripserunt, et obseruari praeceperunt. Illi autem regiam istam vnius in Ecclesia regenda authotatem, vti docuimus, siue per vsurpationem, siue per delegationem institutam damnarunt. Neque enim vti homines in rebus politicis, quae ad hanc corporis vitam  \pend
\section*{AD I. PAVL. AD TIM. }
\marginpar{[ p.244 ]}\pstart pertinent, multa mutant, instituunt, corrigunt, figunt et refigunt, idque iure et prudenter: ita in ipsius disciplinae et censurae Ecclesiasticae essentia immutare et nouare quicquam possunt, quia huius disciplinae finis, et scopus animas spectat, et ad earum gubernationem refertur, vtque ad poenitentiam homines adducantur. Cuius rei cognitio, efficacia, modus, ratio, et iurisdictio solius Dei propria est, ei nota et relicta:nullius autem regis aut coetus mortalium imperio aut prudentiae concessa. Ergo immutare ipsius censurae et disciplinae Ecclesiasticae fundamentum (quale est, vt ex Aristocratica fiat monarchica administratio Ecclesiae) nemo homo, quamuis Rex et Imperator potest, ac ne vniuersus quidem ipsius Ecclesiae coetus. Accidentalia quidem moderari possunt, quae pertinent ad faciliorem istius disciplinae et censurae vsum vel praxin. Atque eo sensu concedimus formam disciplinae vel censurae Ecclesiasticae eandem in omnibus gentibus constitui non posse. Nempe cùm de accidentalibus, non autem essentialibus vel disciplinae, vel censurae Ecclesiasticae disputatur. Nam essentialia eadem vbique gentium in Ecclesiis Dei esse debent: accidentalia mutari et variari ex variis causis possunt. Ergo ad plures illud ius, ea cognitio, et potestas gerendi et decernendi et dispiciendi pertinet. Sed vtrum sic ad vniuersam Ecclesiam, nimirum, vt suffragium singuli de Ecclesiade qualibet re Ecclesiastica ferant, vbi aliquis admonendus vel reprehendendus erit. Res. Per solos Ecclesiae ηγεμένος et prepositos eam potestatem oportere exerceri, etsi totius Ecclesiae communis potestas est, et in ipsius coe\pend
\section*{CAPVT V. }
\marginpar{[ p.245 ]}\pstart tus aedificationem a Christo concessa et relicta: non solum autem praepositorum ipsorum gratia aut ratione est. Atque nostra haec sententia, cui multi obsistunt, plane est confirmanda, scilicet ad solos Ecclesiae praepositos eius iuris et pocortutin cacreitiuin pertieremon autem ad singulos de coetu et populo Ecclesiae. Probatur autem his fere rationibus.  \pend\pstart Primùm authoritate verbi Dei Actor.11.V.15. 21.et 22.Petrus et Paulus coram Apostolis et praepositis rationem reddunt sui facti, non coram vno quodam Episcopo: quanquam, si quibusdam credimus, iam erat Iocobus frater Domini Hierosolymorum Episcopus. Non etiam singuli de Ecclesia in suffragia mittuntur. Sed coram praepositis Ecclesiae, qui Ecclesia vocantur, ea tota disputatio disceptatur. Petrus ipse priore Epistola cap.5.V.5. negat se vllum in coetum Domini habere principatum. Iudaeorum mores tempore Christi multa quidem in doctrina corrupta habebant: tamen veteris disciplinae et a Deo institutae vestigia retinebantur. Itaque cùm Christus ipse suae vitae et doctrinae rationem redditurus sistitur, conuocantur Scribae, Pharisaei, principes Sacerdotum, id est, praepositi Ecclesiae, totumque synedrion: non autem solus summus Pontifex sibi de Christo sumit iudicium. Non item iudicant singuli de populo, sed praefecti Ecclesiae duntaxat. Denique Christus nos iubet ad plures: non ad vnum tantùm aliquem referre, si qui nos offenderunt nobis sponte non reconciliantur Matt, 18.vers.16.1. Timoth.5.vers.19.  \pend\pstart Secundùm: ratione. Est morum coercitio et  \pend
\section*{AD I. PAVL. AD TIM. }
\marginpar{[ p.446 ]}\pstart censura Ecclesiastica cum verbi Dei praedicatione indiuulse coniuncta, et illius tanquam appendix quaedam 2. Timoth.3.vers.16. Malach.2.vers.4. Itaque Paulus ipse praecipit, vt qui Ecclesiam pascere volunt, iidem eam increpent 2. Timoth.4. At Presbyteri sunt Pastores et ποιμένες Ecclesiae, non autem quilibet de populo. Ergo et ii ipsi, non autem singuli de populo, Ecclesiae correctores, censores, iudices esse et constitui debent. Deinde iidem’illi curam animarum nostrarum haberedicuntur, vt est Hebr. 13. vets.7. Iidem reprehemnduntur, nisi nos admonuerint sceleris nostri, Ezechiel. 34. Iidem verbo Dei iubentur Ecclesiae curam habere. Ergo ad eos, non ad populum cognitio et exercitium huius censurae et potestatis pertinet. Denique  vetustissimus Ecclesiae ipsius Christianae mos et vsus idem probat. Nam semper Presbyteri cum eo, qui Episcopus dictus est, iudi. carunt de scandalis et censuris quae cuique fieri debent. Ambros.epist.28. Cyprian. epist.39. et 40. non autem vnus ille Episcopus, quantunuis mapuis aiimi dotibus excellens, non item quilibet de populo. Ergo haec iurisdictio est totius quidem Ecclesiae, ratione potestatis, Praepositorum autem, ratione exercitii et administrationis.  \pend\pstart Qui vero contra sentiunt, ad totum Ecclesiae coetum et multitudinem hoc exercitium potestatis pertinere, afferunt.  \pend\pstart 1 Dictum Christi Matth.18. vers.17. Dic Ecclesiae. Resp. Ecclesiae nomen illic accipi pro iis, qui totius Ecclesiae authoritate fulti et vocati legitime Ecclesiae praesunt ac inuigilant. Sic quod Senatus Roman. decreuit, dicitur Ro\pend
\section*{CAPVT V. }
\marginpar{[ p.247 ]}\pstart 2 Obiicitur Paulus 1. Corin.5. qui conuocasse videtur totam Ecclesiam, cùm ait vers.4. vobis et meo spiritu conuocatis in nomine Domini, cùm de incestuoso iudicando ageretur. Ergo ad illius suffragia haec censura pertinet.i. ad singulos de populo. Resp. Paulum decernere quidem quid sit faciendum, sed a Presbyteris et a Senatu Ecclesiae:id enim suffragiis populi non permittit. Itaque non pertinet ad populum ea deliberatio: sed ad praepositos Ecclesiae, populo assentiente. Aiunt Quod ad omnes pertinet, ab omnibus fieri debet. Respond. ex Ambrosio in liber  de Dignitate Sacerdotal. cap.3 Agi quidem debere a quoque pro loco et munere, quod in Dei Ecclesia gerit. Aliud enim est, quod ab Episcopo requirit Deus: aliud quod a Laico. Sed cùm sint pastores et Presbyteri Ecclesiae ηγομενοι, ad eos censurae iudicium pertinet, ad quos eosdem pertinet immorigeros Deo diuini iudicii metu iniecto compescere. Ad totam vero Ecclesiam spectat discere alieno exemplo, quenque Deo rebellem poena esse dignum. ltaque publicae censurae in Ecclesia fieri debent ad ipsius aedificationem, et vt intelligant omnes scelera in coetu Christianorum non esse tolleranda in can. Certum est 5 sed adhuc 24. quaest.4. Aliae obiectiones afferuntur, sed omnino leues Executioni igitur publicae censurae interuenire debet notitia et consensus Ecclesiae audientis et admonitae de poena scelerati hominis et obstinati:ipsius tamen iudicium et executio ad solos praepositos Ecclesiae pertinet.  \pend
\section*{AD I. PAVL. AD TIM. }
\marginpar{[ p.248 ]}\pstart Quod autem quidam existimant, vbi est vere Christianus Magistratus, ibi nullam esse censuram Ecclesiasticam debere, illi certe Ecclesiasticam et Politicam iurisdictionem inter se pessime confundunt. Deinde non vident se refutari Exod.1o. vorde Hecienaitieu censura, et poena politica locum in vno et eodem homine et negotio habet, diuerso tamen respectu.  \pend\pstart Augustinus certe vixit sub imperatoribus Christianis, et tamen idem liber 11.de Genes.ad lite. cap. 4o. docet solitos esse Episcopos vti aduersus scandalosos homines excommunicatione et censura Ecclesiastica. Item confirmatur, omnibus Synodis, quae sub Christianis imperatoribus habitae sunt, vbi censurae Ecclesiasticae et confirmatae et retentae sunt. Idem denique, si quis hoc etiam genus probationis admittat tota 24. quaest. 4 crucentee pu poenam exigit a nobis scelerum nostrorum tum in corpore, tum in anima. Nec ea sane est duplex poena: sed vna tantùm. Tanta enim peccato debetur, et ea, per quam et animus et corpus peccatoris puniatur, quia peccatum hominis et totius suppositi actio est.  \pend\pstart Quibus expositis iam accedamus ad ipsa Pauli verba. Distinguit autem hoc loco homines, Paulus ab Ætate, et Sexu. Ab aetate, alii sunt Seniores, alii Iuniores. A sexu, alii sunt Mares, alii Foeminae. Hic enim πρεσβυτέρων vocem ad aetatem refero, non autem ad dignitatem Ecclesiasticam, quemadmodum tamen saepe accipitur, velut inf. vers.17.19. et 1. Pet.5 vers.1. Id quod ex ἀντιθίσει et voce νεωτέρους apparet. Cur autem haec praece. pta danda existimauerit Paulus, ratio est, quod  \pend
\section*{CAPVT  V. }
\marginpar{[ p.249 ]}\pstart oporteat ministrum verbi Dei et pastorem ὀρθοτομιῆν, id est, recte et pro cuiusque captu et vtilitate dispensare sanam doctrinam, et quae ex ea proficiscuntur admonitiones 2. Timoth. 2.vers.15.Id quod fecisse Ioannem animaduertimus 2. Epist. cap,2, vers. 12.13.14. Sed obstat quod hac ratione, quam inducit hoc loco Paulus, videtur haberi et praescribi acceptio personarum, quam in vero et fideli dispensatore verbi sui esse vetat Deus Malach.2.vers.9. Praeterea cùm in Christo neque sit mas, neque foemina, neque seruus, neque liber, neque senex, neque adolescens, cur ista vel sexus vel aetatis ratio et distinctio fit? Respond. Personarum acceptionem non induci, cùm omnes reprehendantur qui peccarunt. Sed aliud est prudens et fidelis verbi doctrinaeque coelestis dispensatio, quam Paulus ὀρθοτομίαν appellauit: aliud autem personarum gratiosa et damnata acceptatio, quae dicitur πρπωποληδία. Quod autem affertur in Christo neque marem esse, neque foeminam Galat.3, vers.28. alio prorsus pertinet, quam ad hoc argumentum.  \pend\pstart Quaeritur etiam vtrum haec praecepta ad priuatas tantum, an vero ad publicas etiam reprehensiones, quae fiunt a Ministris, pertineant. Publicae vero sunt admonitiones, quae et in consistorio fiunt, et nomine totius Ecclesiae a senatu Ecclesiastico, vel ab aliquo Presbytero, quanquam priuatim et domi fiant. Respond. Quidam ad priuatas tantùm admonitiones et reprehensiones haec praecepta referenda esse putant. Quidam autem ad publicas duntaxat. Ego vero ad vtrasque, quia semper ratio aetatis et sexus a nobis habenda est.  \pend
\section*{250 AD I. PAVL. AD TIM. }\pstart In quaque enim aetate, vt oportet tractanda magna quaedam non tantùm prudentiae, sed etiam charitatis pars posita est, quae vtraque et a publicis et a priuatis admonitionibus nunquam abesse debet, vt suum vsum et fructum habeant.  \pend\pstart Illud autem notandum est, cùm ita Paulus nos agere iubet, nolle tamen nos vel personarum acceptatores vel assentatores esse, vt alienis vitiis blandiamur et fau eamus: sed ne nimium acerba ce uerfugenurratione et oratione vtentes peccatores ipsos a Dei doctrina auertamus. Quae certe acerbitas prorsus est inutilis, quanquam tamen grauitas est a pastore retinenda Tit.1.vers.7.  \pend\pstart Haec igitur sunt de his rebus Pauli praecepta et rationes I Pastor Euangelicus, tanquam peritus animorum medicus, pro suae aetatis et sexus ratione tractato quenque, neminem imprudenti et acerba agendi ratione offendito. 2 Seniores tum mares, tum foeminas, peccantes tanquam patres et matres pastor Euangelicus leniter admoneto. 3 Iuniores, tum mares, tum foeminas, peccantes velut fratres et sorores idem commonefacito. 4 Foeminae imprimis autem iuniores ita a pastore Euangelico moneantur, vt summus castitatis pudor in ipsa agendi et loquendi cum illis ratione eluceat et appareat.  \pend\pstart Επιπλήξης. Quidam ad pugnos et verbera referunt, quae quidem infligi viris senioribus a Timotheo vetet Paulus: sed perperam. Pertinet enim hic locus ad verba tantùm acerbiora in pastore damnanda: non ad verbera. Nam supra Paulus docuit non debere pastorem esse percussorem et τλήκτην. Ratio vero horum omnium praece\pend
\section*{CAPVT  IIII. }
\marginpar{[ p.251 ]}\pstart ptorum est in vniuersum duplex. Prima, quod omnis reprehensio est ipsa per se iam acerba. Itaque quadam verborum lenitate est mitiganda, et tanquam condienda, vt sit vtilis peccatori, lubentiusque ab eo audiatur.  \pend\pstart Secunda, quod ista ratione pastores agentes, non impetu quodam animi ferri, sed summa charitate eos impelli apparebit, cum vitia reprehendunt. Maiorem autem vim et pondus habet nostra oratio, vbi sibi persuadent homines, se a nobis non animo iniurandi: sed summo beneuolentiae affectu commoneri, et reprehendi.  \pend\pstart Quod ad speciales cuiusque praecepti rationes, sunt duae. Prima, quod aetas illa iuuenilis, qualis fuit in Timotheo, omnino rationem hanc agendi necessariam ene suadet. Multo enim odiosiores sunt reprehensiones, quae nobis ab adolescentibus fiunt, quam quae a senibus et maturae iam aetatis viris. Senum enim prudentiam admiramur: atque illis propter aetatem facile cedimus.  \pend\pstart Secunda ratio, quod haec senum et iuniorum appellatio, quam praescribit hic Paulus, cum lege Dei maxime conuenit et consentit. Illa enim iubet, vt senes, et ipsam senectutem veneremur. Vnde Seniores pro patribus agnoscere suadet. Qui vero sunt nobis aetate aequales, fratrum nominum et sororum fere vulgo, amicitiae ergo, a nobis appellantur.  \pend
\phantomsection
\addcontentsline{toc}{subsection}{\textit{2 Mulieres natu grandiores, vt matres iuniores et sorores, cum omni puritate. }}
\subsection*{\textit{2 Mulieres natu grandiores, vt matres iuniores et sorores, cum omni puritate. }}\pstart Cum omni castitate) Etiam in ipso agendi cum illis iunioribus foeminis modo, ne praetextu mu\pend
\section*{AD I. PAVL. AD TIM. }
\marginpar{[ p.252 ]}\pstart neris et officii nostri, imo vero verbi Dei, videamur impurae libidinis semina quaedam aut irritamenta iacere. Illud enim certe est accipere nomen Dei in vanum.  \pend
\phantomsection
\addcontentsline{toc}{subsection}{\textit{3 Viduas honora quae vere viduae sunt. }}
\subsection*{\textit{3 Viduas honora quae vere viduae sunt. }}\pstart Πρόθεεις. Nam inter eas personas, quarum maxima et peculiaris quaedam a nobis ratio habenda est in Dei Ecclesia, merito censentur viduae, de quibus idcirco hoc loco quaedam praecepta Paulus tradit. Viduarum autem, quae in Ecclesia versabantur, aetate Pauli, duplex fuit genus, Vnum earum quae Simpliciter viduae erant: Alterum earum, quae Allectae et cooptatae erant in ministerium Diaconorum. De primo genere agit primo loco Paulus, de secundo postea hoc ipso capite a versiculo 9. et seq.  \pend\pstart Viduae nobis a Deo commendantur passim in scriptura veluti Exod. 22.vers.22. Zach.7.vers. 1o. Psal.146. vers.9. et Luc.7. vers.12. Haec vidua erat, et aspexit super eam Dominus. Itaque in iure ciuili Vidua, Pupillus, et Peregrinus dicuntur Miserabiles personae liber  2. codic. id est, tales, quae sint ipsae per se miseratione dignae:non autem Deo exosae, sed cui sunt plurimum curae, vt mirum videri non possit, si earum curam habuit Ecclesia Dei, praesertim fidelium et piarum. Itaque hos tres canonas de viduis subleuandis tradit Paulus.  \pend
\section*{CAPVT  V. }
\marginpar{[ p.83 ]}\pstart destitutae bonis, et aliis auxiliis ab ipsa Ecclesia aluntor, et exhibentor.  \pend\pstart Hos tres canones ex tota hac disputatione colligimus. Ac quidem duos ex hoc versiculos nempe 1. et 3. Secundum autem canonem ex sequenti versiculo, qui est 4. Tίμα Hebraeorum more dictum est, vti in praecepto quarto Honora patrem et matrem, qua voce non tantùm reuerenter et sollicite curandas esse viduas a pastore intelligimus, sed etiam alendas ab Ecclesia, et exhibendas, si aliunde non possint se sustentare. Hoc idem praeceptum transferri ad vniuersum pauperum fidelium genus et potest et debet, qui si a suis possunt, debent ali sine onere Ecclesiae: sin minus, ab Ecclesia. Debet enim inter nos ea vigere charitas et tanta, vt nemo inter nos egeat, quemadmodum Dominus ipse praecepit, neque sit memdicus Deut.15. vers. 4. non quidem vt pauperes qui sunt eiiciantur: sed vt alantur a nobis et subleuentur.  \pend\pstart Vnde sicolim etiam nondum plane corruptis Ecclesiae Romanae moribus a Gelasio primo constitutum fuit, vt ex omnibus Ecclesiae reditibus fierent quatuor partes (in quibus ipsis etiam fidelium oblationes continebantur) ex quibus vna tantùm pars erat Episcopi:secunda clericorum: tertia insolidum pauperum, quarta vero fabricae templi applicanda, quemadmodum relatum est in cano. Quatuor autem et can.seq.12. quaest.2. Adeo vt Bernard. epist. 348. velit etiam sacros Ecclesiae calices vendi et distrahi, vt pauperibus suecurratur. Quod antea fecerat quoque Exuperius Tholosanus Episcopus, vt Hierony. scribit,  \pend
\section*{AD I. PAVL. AD TIM. }
\marginpar{[ p.254 ]}\pstart Vere vidua) Quaenam autem illa sit vera vidua postea docet paulus ins.vers.5. est autem quae Pia est, et Destituta omni alia ope. Vnde illa Hieron. definitio, quae est in illius quadam epistola ad Fabiolam non satis huic loco conuenit, vbi de officio pastoris euangelici agitur. Vidua est, ait, cuius maritus mortuus est. Primum enim de piis tantùm viduis hic locus intelligitur, non de omnibus: deinde proprie quod ad subsidium dandum Ipeetat, de ns tuneumeomgaut, quae destituuntur omni alio auxilio et ope: non de omnibus egenis, etiam viduis.  \pend\pstart Πιαρακολήθησις, siue Τάξις. Prospicit enim viduis et docet a quibus illae sint alendae. Nempe a propriis liberis, id est, filiis, nepotibus, abnepotibus in infinitum, si quos tamen illa habet, qui id facere possint. Videtur autem in genere ipso dicendi Paulus alludere ad cap.21. Genes. vers.23. vbi vidua sic describi videtur.  \pend\pstart Vox igitur Μανθανέτωσαν ad liberos viduarum pertinet. Æqua sane ratio, vt etiam hoc loco nepotes, et in infinitum alii, liberorum nomine conprehendantur, quia ab auia tanquam a propria stirpe suam originem habent, et vt sint. Itaque eos, vti authores, post Deum generis vitaeque suae  \pend
\section*{CAPVT  V. }
\marginpar{[ p.255 ]}\pstart Ratio huius praecepti hic multiplex affertur a Paulo. I A consequenti vel consentaneis. Sibi ipsis hoc quodammodo liberi impendunt, quod in parentes suos conferunt: habent enim a parentibus, vt sint, et sunt cum illis vna caro et vnus sanguis 2 Ab officio, oportet liberos rependere parentibus quodammodo quod ab iis acceperunt. Id enim lex ipsa naturae iubet et dictat, quae nos obligat ad ἀντιδῶρα, vt Iurisconsulti docent, et ad ἀντιπελαργίαν, quam hoc loco vocat ἀμαλὶν Paulus. Hoc exemplo Ciconiarum facere docemur. 3 Ab honesto. Hoc ipsum Deo gratum est et acceptum. Quod ex praecepto Honora patrem, et matrem, satis intelligi potest. Illa enim voce, Honora: non tantùm reuerentia quaedam externa significatur, sed etiam subsidium vitae, quod est a liberis parentibus praestandum. Quod officium etiam pietas hoc loco nominatur, et a Platone iustitia. At haec omnia iura semel violant, qui religionis et Monasticae vitae, quam profitentur, praetextu, se ab ea de parentibus alendis cura nberatos ene contendunt, quod Monachi facti sint. Denique omne liberorum in parentes obsequium abiiciunt, et officium suum, propter votum Monachatus a se semel susceptuii, ccrlare et nullum iam vsum aut locum habere debere sentiunt. Quae sane est horrenda impietas et omnis humani sensus violatio. Quod autem de parentibus alendis a propriis  \pend\pstart liberis et nepotibus Paulus hoc loco praecipit et tradit, omnino latius patet. Omnes enim eos, qui parentum loco sunt, etiam complectitur:non tan\pend
\section*{AD I. PAVL. AD TIM. }
\marginpar{[ p.256 ]}\pstart tùm eos, qui nos e suo sanguine genuerunt, qualis est patruus, amita, matertera, frater, et soror, denique quales sunt omnes, qui nobis aliqua sanguinis necessitudine sunt coniuncti: quos egentes et inopes deserere, cùm ipsi abundemus opibus, est certe criminosum, nefas et impium. Est enim Charitatis regula amplianda, et ex pari ratione, quid Dominus nobis hic praescribat, satis intelligi potest. Itaque Leuitic. 18, illae personae, quas commemoraui, nobis parentes ipsos referre dicuntur et pietas Abrahami nora est, qui patrem Thare assumpsit. Dauidis, qui omnes e sua familia secum in deserto habuit, et aluit, itemque patrem et matrem tuto loco deposuit apud regem Moabitarum 1.Samuel.22.Idem in l.alimenta C. de negotiis Gest. ab hominibus profanis responsum est.  \pend
\phantomsection
\addcontentsline{toc}{subsection}{\textit{5 Porro quae vere vidua est ac sola, sperat in Deo, et permanet in supplicationibus et precibus nocte ac die. }}
\subsection*{\textit{5 Porro quae vere vidua est ac sola, sperat in Deo, et permanet in supplicationibus et precibus nocte ac die. }}\pstart Ε’ξηγησις est siue ὁρισμός. Exponit enim quaenam sint illae vere viduae, quas supra honorare praecepit. Ex hac autem ipsa descriptione, quales viduae piae esse debeant, intelligitur. Haec autem verba poti? malcandi, quam imperandi modo sunt accipienda, quemadmodum fecit Ambrosius, idque omnino recte, et ex propria Graecarum vocum significatione. Ait igitur Paulus vere vidua est ea, quae manet in precibus: non autem sic ait quasi praecipiens, vere vidua maneto, vel maneat in precibus. Primum igitur harum, quarum cura ab Ecclesia maxime suscipi debet, conditionem proponit, et quam sit  \pend
\section*{CAPVT V. }
\marginpar{[ p.257 ]}\pstart pastorum Euangelicorum sollicitudine digna, docet Paulus: deinde ab effectu siue proprio earum munere easdem depingit, nempe a precibus et pietate eximia. Est in cano. Vidua distinct. 34 aliqua de viduis disputatio ex Hieronym. sed quae nihil, nisi superstitiosum et mancum, habeat. Eae igitur viduae Ecclesiae eleemosynis subleuari debent, quae describuntur hoc loco. Nam exhis Pauli praeceptis etiam aliae regulae et canones de pauperibus ab Ecclesia alendis sunt decerpendi, vt et Pastores, et Diaconi Ecclesiae intelligant, quos in albo pauperum habere, quibusque ex aerario Ecclesiastico succurrere debeant, ne, si aliud hominum genus alant, fucos potius quam pauperes fuer lant ce paicarit. et temere tendes Ecclesiae opes profundant. Magna enim cura est adhibenda, vt prudenter, et quibus oportet, fidelium eleemosynae distribuantur, vt etiam 2.Thessalon.3 docet Paulus.  \pend\pstart Videtur autem hoc loco proponi duplex verae et pauperis viduae descriptio, vna ab Etymi ratione et explicatione sumpta, altera ab officio et munere earum. Nam cùm ait Paulus χῆρα ἐςτιν ὸ μεμονωμένη vocis ipsius χύρας significationem synonymo explicauit. χῆρς enim ἀπὸ το͂ χιρο͂ν, quod est desolare, destituere, orbare proculdubio deducitur. Hoc autem ipsum vox sequens μεμονομem plane designat, quae solam esse, desertam, et praesidio destitutam significat. Id quod ex Genes. 15. vers.3.videtur sumptum. Sed a Latinis ipsis viduertas dicta est calamitas, vt scribit Sext.Pompeius. Haec igitur μονώτης et solitudo non tantùm marito priuatam esse oportere hanc viduam do\pend
\section*{AD I. PAVL. AD TIM. }
\marginpar{[ p.258 ]}\pstart cet (neque enim aliter vidua dici posset) sed et liberis, et cognatis, et omni alio humano praesidio destitutam. Quae iam conditio plane est afflictissima et commiseratione, atque Ecclesiae et pastorum cura dignissima. Quare ipso viduae nomine, id est, desertae et destitutae mulieris appellatione, quae certe statum miserabilem denotat, illae nobis conmendatissimae esse debent. Vbi enim melior eleemosina, aut vbi iustior misericordia esse potest, quam vbi est maior miseria atque afflictio? Neque  tamen, si quae viduae eos liberos aut propinquos habeant, a quibus praeter officium deserantur, sunt ab Ecclesiae eleemosyna arcendae, et deserendae, quia perinde est tales cognatos habere, atque omnino nullos. Haec igitur est prima viduae definitio.  \pend\pstart Sequitur altera, ex qua maior etiam commendationis earum ratio apparet. Eas enim vere viduas appellat, quae Pietatem habent, et Eam exercent suo specialique  modo, qui eas deceat. Quod enim de viduis ait hoc loco Paulus, proprium quodammodo earum esse videtur, non etiam conmune caeterarum mulierum. Est vero vtraque definitio coniungenda, vt quae sit vera vidua, possit agnosci. Neque  enim omnis quae sola est et deleitutu alieno auxilio, est vera vidua, si piam se non praestet. Neque omnis mulier, quae pia est, vidua quoque est, aut ea vidua, de qua hic Paulus agit, quae ab Ecclesia debeat ali, veluti si non est μεμoνομένη. Neque enim omnes, etiam, quae solae sunt et desertae, piae sunt: sed quae vtrunque misera experitur, ea demum vere est vidua.  \pend\pstart Spes autem in Deum hoc loco continet πειγραφην  \pend
\section*{CAPVT  Y. }
\marginpar{[ p.29 ]}\pstart fidei, cuius spes est effectus. Itaque perinde est, atque  si diceret Paulus eam esse viduam, quae veram fidem habet. Apposite tamen spei potius (quae aduersus res duras et afflictas pugnar) quam fidei nunc meminit, quia de afflictis personis agit, quas spes corroborat, vt intelligamus qualem se praestare Christianam viduam oporteat in rebas accisis, et cùm se omnino destitutam videt. Nimirum igitur in Deum sperare debet, non autem animi diliiuentia detperare, aut propter huiusmodi res aduersas veram religionem abiurare, aut instar infidelium mulierum in Deum ipsum impie murmurare, et blasphemare  \pend\pstart Permaneat in supplicationibus,) Hoc iam proprium est Christianarum viduarum munus, et earum quae sunt maxime ab Ecclesia subleuandae, vt ipsis externis actionibus et operibus se totas Deo esse deditas atque consecratas ostendant. Itaque συνεκδοχικῶς Paulus a parte et specie totum genus piarum externarumque actionum expressit. Simile quiddam de Anna Phanuelis filia et prophetissa narrat Lucas 2.vers.37. et 1. Samuel.2.v. 22 facta est eius consuetudinis mentio. Ergo assiduitas ista precandi demonstrat huiusmodi viduam esse debere, quae, omnibus aliis rebus postpositis soli Deo, ipsiûsque negotiis et Ecclesiae inseruiat et vacet. Denique ea curet, quae sunt Domini, non autem mundi huius, vt loquitur idem Paulus 1.Corinth. 7. vers.32. Nam quae nupta est mulier, varia sollicitudine, variisque curis et cogitationibus saepe propter familiam distrahitur, ne tam libere tamque assidue, et frequenter precibus publicis interesse, et de Deo ipso  \pend
\section*{AD I. PAVL. AD TIM. }
\marginpar{[ p.260 ]}\pstart cogitare possit. Haec igitur sunt quodammodo propria viduarum opera, quia haec assiduitas precationum non potest pari modo a reliquis exigi quae non sunt, vti viduae, ab omnibus impedimentis solutae, sed occupantur in iusta rerum domesticarum cura, qualis erat Martha Luc 1o.vers.4o.  \pend\pstart Obstant vero duo huic Pauli praecepto: e quibus primum est, videri Paulum magnam hominum animis superstitionem iniicere, cùm Christianas mulieres vult diem et noctem in precibus permanere. Atque istae videntur fuisse superstitiosae Monacharum et Monialium instituendarum causae, nimirum vt illae nihil nisi precarentur: et tum pro se, tum pro aliis, preces Deo assidue funderent.  \pend\pstart Cuius etiam vitae exemplum Messalianorum haereticorum erroribus, et moribus inductum est, sed est prorsus a Dei praeceptis alienum. Itaque  Psallianorum haereticorum surculus vel potius propago miserrima fuit istud Monachorum institutum, vt ex Cassiano, qui prolixe omnes istas nugas descripsit, facile colligi potest. Respond. Non laudari hic a Paulo otium illud inutile et superstitiosam precationum institutionem, qualem hodie habent obseruantque Monachi, vel eciuiii asii Papistae superstitiosi, qui suas etiam preces venales habent, quemadmodum olim inter Ethnicos praeficae illae in funeribus adhibitae mercenariae: sed tantùm ostendit Paulus ecquod sit verum et legitimum earum, quae nullis negotiis domesticis distrahuntur, exercitium, nempe, vt in iis quae sunt Dei, prorsus se occupent, qualis est nimirum precatio, concionum assidua frequenta\pend
\section*{CAPVT  Ve }
\marginpar{[ p.261 ]}\pstart cioscereinequu huiuimour pletutis d endritatis verissima officia et exercitia.  \pend\pstart Neque tamen cùm haec officia a viduis requirit Paulus, ab iisdem reliquos homines et mulieres, veluti coniuges, excusat: sed ex comparatione, quae etiam versiculo proximo sequitur, docet, quod sit potissimum viduarum mulierum, quae solae sunt atque sine familia, iustum studium ac exercitium. Neque etiam hoc loco vlli verborum demur murationi, quantunuis assiduae, Paulus pietatem alligat, sed ex opere externo, quae sit interna fides et pietas, et vt deprehendi possit, solùm demonstrat. Nam potest in ea ipsa assiduitate precandi inesse damnabilis superstitio, quam non probaret certe quidem Paulus, quales quaedam mulierculae cerauntur maxime in Papatu, quae omnes aras atque omnia templa quotidie circumeunt, tanquam in iis Deus habitet, aut colatur sanctius: aut etiam ipso precum et vocum repetito numero delectetur, contra Domini nostri Iesu Christi sententiam Matth. 6. vers.7. Itaque res suas domesticas, curamque familiae et mariti praeter officium deserunt. Quid de iis dicam, qui coniuges cùm sint, etiam in Iesuitarum, id est, pessimorum idololatrarum ordinem, inuitis maritis, cooptari se volunt?  \pend\pstart Obstat item quod in Actis Apostolorum 9.v. 39. affertur, vbi alia viduae Christianae exercitia describuntur, nempe opera manuum, labores corporis, qualia sunt nere, suere, sarcire et caetera huiusmodi, quae magnam humanae societati commoditatem afferunt. Respond. vtrunque recte inter se conuenire et consentire. Nempe Deï  \pend
\section*{AD I. PAVL. AD TIM. }
\marginpar{[ p.262 ]}\pstart cultui interesse assidue, quod hic describitur: et pauperum curam habere, ac ea, quae prodesse illis ponunt, parare, quou laelebut Dorcas. Vnum cnim alterius comes est.  \pend
\phantomsection
\addcontentsline{toc}{subsection}{\textit{o At quain tuxu viuit, ca viuens mortua est. }}
\subsection*{\textit{o At quain tuxu viuit, ca viuens mortua est. }}\pstart Ἀντίθεσις est, eaque duplex, Vna quidem, quae superiorem sententiam illustrat ex contraria viduae descriptione. Altera, quae in ipsius versiculi sententia, quanquam breunisisia, satet, habetque magnam emphasin, pulchramque insolentis viduae picturam, quam inf. vers. 13. copiosiùs postea persequitur Paulus.  \pend\pstart Primum igitur lasciuientem et luxu diffluentem viduam opponit verae viduae Paulus. Neque enim huiusmodi ab Ecclesia ali vult, imo potius serio admoneri, vt exitium suum intelligat et apprehendat. Vocat autem eam παταλῶσαν, qua voce idem significare videntur Graeci quod vocabulo ἐντρυφαν, vt apparet ex Iac.5. vers.5. Eae vero sunt huiusmodi omnes, quae cùm viris sunt orbatae, delitiis, corporis voluptatibus, lasciuiae indulgent, item sumptuosis vestimnentis, magnis virotum coetibus inuisendis delectantur, seque stude corpore et vultu et cultu componunt, comunt, et comparant, vt viris placeant, arrideant, ac allubescant, quales pene sunt hoc nostro seculo infinita. Itaque illae etiam ipso Dei flagello miserrimae gaudent, et laetantur, quod viros suos, id est, optimum vitae suae subsidium amiserint. Nam hac ratione se ad omnem lasciuiam perpetrandam liberiores esse iam putant. Certum est  \pend
\section*{CAPVT  V. }
\marginpar{[ p.263 ]}\pstart autem semper apud probas mulieres, et bene moratas etiam ciuitates alium fuisse viduarum, quantunuis iuuenes adhuc essent relictae, et alium coniugum mulierum habitum, et cultum vti ex libro ludith. et historia Thamaris Genes. 38.v. 14. apparet et ex Tertull. non quidem, vt aliqua superstitio in cultu induceretur: sed vt modestia et honestas personarum et ordinum in Ecclesia conseruaretur. Quae certe distinctio honestati, atque etiam ipsius naturae sensui consentanea propter dissolutissimos nostri huius seculi mores iam sublata et antiquata est, adeo vt magis comptae et dissolutae viduae nostris his temporibus nunc incedant et cernantur, etiam quae se Christianas, si Deo placet, appellant, quam et coniuges et maritatae.  \pend\pstart Alterum autem est, quod etiam hoc loco ἀντίθεσν continet, nempe quod ait Paulus, ζῶσα τέθνηκe id est, viuens mortua est. Quae enim potest esse haec vita, quae tôta offendiculi plena est, et certissimum futurae mortis aeternae indicium ac testimonium? Haec autem sententia quanquam contradictionem implicare videatur, vt aiunt, et prorius αδυνατος esse, facile tamen explicari potest, quemadmodum docet hoc loco Chrysosto\pend\pstart Saepe autem vita corporis cum animi morte coniuncta est, vt hoc loco. Saepe autem disiuncta ab ea, vt cum pii corpore moriuntur, vt animo vi\pend
\section*{AD I. PAVL. AD TIM. }
\marginpar{[ p.264 ]}\pstart uant: mundo moriuntur, vt Deo viuant. Ergo viduae lasciuae viuunt quidem corpore, animo tamen mortuae sunt. Sic infideles dicuntur et a Christo et a Paulo mortui 1. Corinth.15. vers.29. Matth.8.vers.22. Sunt enim ad omnia quae ad vitamroternal perunent, immobiles, et inepti, et inutiles, vt ait Chrysostomus: vnde perinde, atque  mortuae merito inter pios censentur hae viduae.  \pend
\phantomsection
\addcontentsline{toc}{subsection}{\textit{Haec igitur denuntiato, vt sint irreprehensae. }}
\subsection*{\textit{Haec igitur denuntiato, vt sint irreprehensae. }}\pstart Haec repetitio, siue potius admonitio a Paulo facta superiorem doctrinam valde vtilem esse docet, neque semel tantùm pronuntiandam, sed saepius in Dei Ecclesia inculcandam, vt animis hominum inhaereat altiùs. Eiusdem autem maximus fructus hic ostenditur, cùm, si obseruetur, omne offendiculum sublatum iri dicat Paulus, foreque Ecclesiam ipsam irreprehensibilem, et boni apud omnes ouoria ridelibus autem et infidelibus bene olere debemus omnes.  \pend
\phantomsection
\addcontentsline{toc}{subsection}{\textit{S Quod si qua suis et maxime domesticis non prouidet, fid em abnegauit, et est infideli deterior. }}
\subsection*{\textit{S Quod si qua suis et maxime domesticis non prouidet, fid em abnegauit, et est infideli deterior. }}\pstart Α’ιτιολογ́α est superioris canonis de parentibus alendis a liberis, et contra, quae tum a turpi ducta dici potest, tum etiam a repugnantibus: est autem haec sententia generalis, quae tam ad mares, quam ad foeminas pertinet. Itaque relatiuo communis generis vsus est Paulus. Hoc autem pro certo et confesso sumit, quod pii fatentur, debe\pend
\section*{CAPVT  V. }
\marginpar{[ p.265 ]}\pstart re Christianum quemlibet, siue marem, siue foe minam longe sanctiori viuendi ratione omnibus innotes cere, quam infideles homines soleant: adeo vt quaecunque vel honestas ipsa, vel iustitiae et pietatis leges postulant, ea lubentissime omnia praestare studeat. Sic Christus Math.5.vers.46. Nam si dilexeritis, eos qui diligunt vos, quam mercedem habebitis, nonne publicani haec faciunt. Est etiam huius hypotheseos et sententiae ratio apertissima ex definitione Christiani viri, et Christianae doctrinae siue religionis. Est enim vir Christianus, qui caeteris hominibus vitae sanctioris facem in omnibus praefert. Math.5. Est etiam Christianismus vt Basil. definit Homil.1o. in Examero. nostra cum Deo similitudo, quantum humaha natura illius capax esse in hoc mundo potest. Est igitur vir Christianus non tantum, vt caeteri, ad Dei glo riam natus: sed ad Dei optimam similitudinem et perpetuum obsequium deo praestandum et for matus et intentus. Denique vt ait Paulus. 2.Timo.3.vers. 17. ὁεην ἀνθρωπες τιὸς πᾶν ἔργον ἀγαθὸν ἐξηρτισμένος. Id quod de infidelibus hominibus di ci minime potest. Itaque Paul Rom.33. et Petrus I.epist. cap.3.latum discrimen inter fideles et insi deles constituunt: veluti quod appellanturfideles filii lucis, homines Dei, et serui Ben Hre infideles appellantur filii tenebrarum, filii diaboli, et serui peccati. Ergo hic canon est, Christianus quilibet vt caeteris exemplum sit pietatis, suorum omnium, imprimis autem domesticorum curam habeat. Qui secus faxit, fidem abnegasse, et infideli deterior esse censetor. Voce προνοεῖν et Animi et corporis no strorum curam complectitur demandatque Pau  \pend
\section*{AD I. PAVL. AD TIM. }
\marginpar{[ p.266 ]}\pstart lus quanquam de alimentis maxime agit hoc loco. Neque enim qui corporis et vitae huius caducae curam haber, is animum, qui longe potior hominis pars est, debebit negligere: sic interpretatur Chrysosto: et confirmatur Deuter.11.vers.19. Ephes.6.vers.ꝗ in verbis ἐν παιδεία καὶ νυθεαία κυρία. j  \pend\pstart Suorum et Domesticorum) omnes intelligit Paulus, qui ad nos siue iure sanguinis, siue iure fa miliae et subiectionis pertinent. Est enim omnium illorum a nobis cura habenda. Quanquam enim om nes non sunt pari gradu coniuncti nobis, etiam qui sanguinis necessitudine sunt deuincti, nulli tamen sunt a nobis spernendi et deserendi. Quanquam etiam saepe in alienas familias transiisse videntur nostri liberi, vel per adoptionem ciuilem: vel per coniugia: propterea tamen ius illud et vinculum langoins non est abolitum, neque abruptum, vt eos negligere debeamus. Iura enim sanguinis nullo modo tolli possunt. Vnde perpetuo ad nostrum erga consanguineos et cognatos officium praestandum, ac curam de iis habendam iure naturae, di uinóque praecepto manemus obligati.  \pend\pstart Hos autem gradus statuit hoc loco Paulus, vt Domesticos primum nobis commendet, et praeferat iis, quos ἰδυς appellauit. Ratio est, quod hi domestienquos itenigit Paulus et sanguis sunt noster: atque etiam nostrae familiae pars, quae nobis praeci pue commendatur: et cui a Deo praeficimur specialiori atque arctiori quodam vinculo. Non vult autem seruos a nobis propinquis et cognatis prae ferri, etsi serui sunt domesticorum numero fidem abnegat qui secus facitait Paulus: duplicem notam huius modi hominibus inurit et Infidelitatis ferinae vel  \pend
\section*{CAPVT  V. }
\marginpar{[ p.267 ]}\pstart beluinae crudelitatis: vel etiam saeuioris quam ferina est. Quae vtraque ab homine debet certe quidem esse alienissima. Itaque hoc argumentum est a repugnantibus humanae et Christianae charitati mo ribus ductum.  \pend\pstart Fidem abnegare est non verbis ipsis disertis abiurare fidem in Christum, illius doctrinam, quam antea amplexus eras, iam exspuere: sed vim et fructum illius ac praecepta, quae fides Christiana tradit, spernere et conculcare. Qua ratione dicit idem Paulus Tit. 1. vers. 16. multos ore quidem profiteri Christum, qui factis eum abnegent, quoniam illi sunt Dei verbo immorigeri et rebelles. Docet autem fides Christiana non tantùm omnem paternitatem et cognationem hominum inter ipsos essea Deo. Itaque sperni sine Dei ipsius contemptu non posse, vt docet Paulus Ephes.3 vers.15. sed etiam eam, ex illa paterni tate, qua Pater Filium generat, quodammodo profluere. Docet etiam nos esse hac lege Deifilios, vt illi tanquam patri obsequamur, et quam cuique hoittui deuit paternie esementissimus curam, prouinciam, locum, et vocationem, eam lu bentissime tueamur: neque eam per incuriam aut negligentiam deseramus. Praeceptum enim Dei Honora patrem et matrem etiam parentes ad curam liberorum et Domesticorum suorum obstringit.  \pend\pstart Videtur autem dici posse in istos crudeles aliquid amplius, nimirum eos qui suos negligunt, esse non tantum infidelibus peiores: sed etiam bru tis ipsis animantibus deteriores et stupidiores, quae ex se nata suscipiunt, fouent, amplectuntur, et educant,  \pend
\section*{AD I. PAVL. AD TIM. }
\marginpar{[ p.268 ]}\pstart ipso naturae instinctu tantum impulsa, id quod cer te verum est. Obstare autem videtur huic comparationi quod Infidelitatis et contempti muneris notam fere facit ἀτορίαν Paulus, vti in epistola ad Roma 1.vers.31.2. Timoth.3. vers.3. At hic etiam σοργάς φυτικὰς Infidelibus hominibus relinquit. Resp. Docere Paulum, cùm ita Roma.1.vers.31. lo quitur: quo tandem delabantur infideles, qui Deum contemnunt: non propterea tamen eos omnes, qui caelestis veritatis luce destituuntur, ἄςόργυς esse omnino pronuntiare. In quibusdam. n. infidelibus manent istae storgae. At obiici potest, Multi infideles fuerunt ἀςτοργον, et suorum negligentes, quos etiam feris ac bestiis saepe obiecerunt ac exposuerunt. Respon. Paulum hic loqui de iis, qui naturam humanam non prorsus exuerunt quales inter eos sunt multi: qui naturales illos et insitos a Deo omnibus hominibus affectus retinent et sentiunt. Disputat hic insulse Tho mas, vtrum ex hoc Pauli dicto colligi debeat patres fideles suorum liberorum et domesticorum negligentes, esse ἀπλῶς et omnino peiores et deteriores infidelibus et ethuicis, vt nullaiam in re iis sint praeferendi: sed sint vtrique inter se aequa les et per omnia conferendi, adeo vt quemadmo duir imide rm paruuli baptizari non possunt, neque etiam istorum sidelium paruuli baptizari debeant. Respon Hoc Pauli dictum esse κατάτι accipiendum. Quoad ipsos sideles hypocritas, plane in exitium illis cedunt Dei dona vt mitior sit futura conditio et poena Sodomae, quam istorum contemptorum Dei. Quod autem ad ipsa Dei dona, qualis est baptismus suum habent effectum, et firma  \pend
\section*{CAPVT  V. }
\marginpar{[ p.269 ]}\pstart sunt, atque adeo valent ad eum finem, ad quem Deus dedit illa. Itaque vti ingratitudo Iudaeorum non aboleuit signum foederis nempe circuncisio nem: sic neque improbitas parentum efficit quo minus liberi ex foedere in populo Dei censeantur. Particularisquidem effectus donorum Dei in ipsis improbis nullus est, generalis autem effectus eorundem manet propter baptismum et foe dus Der, vt propterea non sint omnino cum infidelibus exaequandi, nec debeant eorum liberi a Dapillllo atceil  \pend
\phantomsection
\addcontentsline{toc}{subsection}{\textit{9 Vidua allegatur non minor annis sexaginta, quae fuerit vnius viri vxor. }}
\subsection*{\textit{9 Vidua allegatur non minor annis sexaginta, quae fuerit vnius viri vxor. }}\pstart Τάξις. Ingressus enim Paulus disputationem de viduis, quae ab ecclesia aluntur, canonem addit de iisdem valde necessarium, maxime pro ratione et illius temporis, et huius quod postea sub secutum est. Nam nostro hoc seculo in omnibus reformatis ecclesiis cessat huius canonis et praecepti vsus. Ex quo apparet, quae ad particularem ecclesiae politiam pertinent et pro ratione tantùm temporum aut locorum aut personarum constituta et recepta sunt, ea non esse perpetua:sed pro diuer sa temporum, locorum, personarum etc.cir constantia mutari posse: quemadmodum etiam ex iis apparet, quae a prima synodo Hierosolymitana, quae ab Apostolis habita est, decreta sunt Actor.15. Neque enim illa semper fuerunt obseruata, quod in eo cernitur, quod de suffocato constitutum est. Illud enim prorsus hodie et cessat et sublatum est.  \pend\pstart Hic igitur canon de Viduis a Paulo sancitur,  \pend
\section*{AD I. PAVL. AD TIM. }
\marginpar{[ p.279 ]}\pstart vt inter diaconissas ecclesiae, ea demum vidua al legetur, quae sexagenaria est. In quo vox haec, καταλογόθω non tantum additionem ad numerum iam constitutum et ascriptionem significat: sed etiam nouam eius ordinis in aliqua ecclesia, si quae iis egeret, nullas dum tamen haberet, electionem. Id quod fieri necesse erat in iis ecclesiis, vti dixi, in quibus nullae adhuc diaconissae erant, et quae tamen iis opus haberent. Itaque et deligi et allegi significat hoc loco de nouo creare, si oporteat, et in album siue numerum iam electarum adscribere.  \pend\pstart Loquitur autem Paulus de Vidua fideli, et quae pietatem veram iam professa sit: non autem de infideli, id quod exsuperiori disputatione satis apparet.  \pend\pstart Duo autem spectari in iis viduis iubet Paulus e quibus alterum Abesse debet, alterum Adesse. Abesse haec duo iubet, vt sit minor sexagenaria et Duos maritos habuerit. Adesse autem vult vt I Liberos suos educauerit. 2 Hospitalis fuerit erga extraneos. 3 Misericors erga afflictos. 4 De omni bono opere testimonium idoneum habeat. Itaque ait Chrylosto. eadem pene hic a Paulo in vidua requiri, quae supra in episcopis ipsis. Et certe cùm hae viduae munus laliquod in ecclesia habuerint, idque publicum, tales deligi oportuit, quae essent aedificationi omnibus siue fidelibus, siue insidelibus et extraneis futurae: ne Dei gloria, et ecclesiae progressio impediretur.  \pend\pstart Hoc autem vt intelligatur, totum hoc argumem tum altius arcessendum est, nempe ex eo loco, qui Actor.6.vers1. et 9. vers.39. Rom.16. vers. 1.2. Nam  \pend
\section*{CAPVT V. }
\marginpar{[ p.271 ]}\pstart legimus primam et Apostolicam ecclesiam ab ipso sui exordio habuisse mulieres Viduas, imprimis ad certa quaedam et quotidiana tum erga vniuersam ecclesiam ipsam, tum erga fideles pauperes et egenos fratres ministeria destinatas, quae postea Diaconissae ad imitationem marium, qui Diaconi dicebantur, appellatae sunt. Sic enim Phoe ben appellat Paulus Rom.16.vers.1. Cur autem id factum sit et inductum in ecclesiam quaeri potest. Nam non modo Paulus 1. Corint.14. et supr. 2.vetat mulieres loqui in ecclesia, et publicum aliquod munus gerere:sed etiam quaecunque leges politicae recte vtriusque sexus discrimen agnouisse videntur, quales sunt Romanae, idem faciunt. Illis enim mulier prohibetur publicum aliquod munus exercere l. Mulier.5.de Regul. Iuris. Cur igitur et sibi ipsi iam non consentire videtur Paulus: et recte constitutum illud inter mares et foeminas discrimen tollere, atque adeo naturae ipsi tanquam aduersari? Resp.neque sibi, neque naturae repugnare hoc Pauli praeceptum: vel antiquissimum illum ecclesiae Christianae morem, quae huiusmodifoeminas ad pauperum ministerium accersiuit, deditque illis in ecclesia munus atque aliquam dignitatem: publicam quidem ratione fructus et vtilitatis, quae ad oiines ex ta proueniebat. priuatam autem ratione iniunctae ipsis prouinciae, et descripti deffinitique  certis augustisque finibus carum muneris. Differunt enim hae Diaconissae, atque earum munus a Diaconis: quod Diaconi toti aerario ecclesiastico et curae vniuersorum ecclesiae pauperum praeficiuntur: hae vero certo cuidam ministerio tantùm, atque adeo sub  \pend
\section*{AD I. PAVL. AD TIM. }
\marginpar{[ p.272 ]}\pstart Diaconis sunt, quibus subiiciuntur et parent, redduntque rationem suae vocationis: Diaconi vero ipsi, toti ecclesiae. Vnde non tam fuit munus hoc publicum, quam subsidium aliquod a Diaconis quaesitum: partim, quod vniuersis ecclesiae pauperi bus subuenire propter numerum tunc non possent: partim, quod vt et Epiphanius et Chrysosto. ait, pleraque circa pauperes vel aegrotos, vel men sas ipsas ecclesiae parandas commodius a foeminis administrantur, quam a maribus: nec modo commo dius: sed etiam honestius fiunt et geruntur: vt si aegrotet Christiana mulier, vel aemorroide laboretjaut aliis huiusmodi morbis occultis, et quae ho neste a maribus tractari non possunt. Ad eam enim honestius alia mulier admittitur, quam mas aliquis Diaconus et medicus: ergo honestas hoc munus suasit admittendum.  \pend\pstart Deinde necessitas. Neque enim omnia munia et ministeria primum obire potuerunt soli Diaconi, cùm tota ecclesia, more Laconum, simul epularetur, et Α’γάπας et conuiuia communia celebraret in testimonium mutuae charitatis. Sic enim fiebat et vt alii aliis. i. ditiores pauperibus subuenirent, ac ea quae erant victui necessaria tanquam in conmuni suppeditarentur Actor.4. et 6. Itaque alique subsidiariae oper ae et ministeria, eaque commoda Diaconis quaerenda ad mensarum administrationem fuerunt, ad quae ascitae sunt hae foeminae. Quas vetustiss. et primaeua ecclesia habuit. Deinde vero caeterae ecclesiae, imprimis autem eae, in quibus diligens pauperum cura et ratio habita est, eum morem, easque Diaconissas retinuerunt etiam post Apostolorum tempora.  \pend
\section*{CAPVT  V. }\pstart 273  \pend\pstart Tertio videntur hae mulieres diaconissae ad Mar thae imitationem toleratae vel inductae in ecclesia, atque ad exemplum earum mulierum piarum, quae Christo ministrabant, de quibus Lucas 1o. vers. 41. 8.vers.3. Vtcunque vtilis visa est ecclesiae pro tempore illa accessio et subsidiaria mulierum opera, antequam essent in Christiana ecclesia Ptochodocheia, Nosodocheia et Hospitalia quae dicuntur, quod commodissimum esset et ma xime idoneum mulierum ministerium quibusdam in rebus, atque adeo aptius, quam Diatonorum.  \pend\pstart Vnde quae aetas secuta est tempora Apostolica, Diaconissas habuit, vt ex conciliis Nebcaesariensi, Laodiceno, et Chartaginensi apparet, et ex veterum scriptis, imprimis autem Epiphanio, qui vixit sub Theodosio maiore, et historia ecclesiastica: adeo vt magnae dignitatis postea huiusmodi Diaconissae in eccrenaruerme rua fienyrias Diaconissa interfuit synodo Chalcedonensi quemadmodum Euagrius scribit liber 2. cap.16. et Pulcheria etiam ipsius Theodosii Iunioris Constantino pol. Imperatoris soror in Diaconissam electa est vt ex Hermio Sozom. facile colligi potest.  \pend\pstart Praeter Diaconissas autem neque Presbyteridas, neque episcopissas mulieres habebat Christiana ecclesia, quanquam Pepuziani haeretici nullam Maribus dignitatem in ecclesia tribuerent, quam eandem etiam foeminis non concederent.At ex Pauli praecepto aliud semper orthodoxa eccle sia obseruauit. Nullum enim huiusmodi munus vnquam foeminis delegauit, quod in docendo versaretur. Itaque Epiphani. Haeres. Collyridianorum quae est 79. sic scribit. Hρεςθυλερίδες ἠἰερίσταιοὔκ ἐιοιν  \pend
\section*{AD I. PAVL. AD TIM. }
\marginpar{[ p.274 ]}\pstart Taii vero do naui remaihus, et ea quae requirit in istis Diaconissis.  \pend\pstart Ne minor sexaginta annis. Ratio est duplex, quod Iuniores oportet dare operam liberis et q lasciuiunt quae nondum eam aetatem attigerunt, quia nondum plane in iis extinctus est vigor generandi. Quaeritur vero cur sexagesimum annum statuat Paulus, cùm quinquagesimus annus foeminis censeatur prorsus sterilis et infoecundus, vt vulgo an notant interpretes in Luc.1.vet.36. et Iurisconsulti in l.viuiC. de Caduc. Toll. Resp. Anno qudem 50. fieri vt plurimum steriles foeminas: non taile lempp̧r Qae dam enim post 5o. annum peperisse dicuntur. Ac post óo nulla est, quae concepisse referatur, nisi miraculo id factum sit, et praeter naturae ordinem, id quod Sarae vxori Abrahami accidit Genes.18. vers.ii. 3.  \pend\pstart In sexagenariis autem foeminis videtur vigor il Ie naturae et prolificus prorsus periisse et defecisse.  \pend\pstart Sed cur tantum tribuit aetati senili Paulus, qui supra Timotheum adolescentem nobis com mendabat.c.4. vers. 12. Resp. Chrysosto, proprie non tribuere Paulum id honoris aetati, quasi sexagenariae lasciuae fieri non possint: sed quod minus in hac aetate offendicula illa nascantur, quae euitanda postea docet. Neque enim ipsa per se aetas facit sapientes vel prudentes homines: sed Deus potius et rerum praeteritarum experientia. Deinde cùm leges ferantur de iis, quae ὥς δ̓ι τοπο»o accidunt, non debet propter paucas aliquot sexagenarias insolentes toti etati suus honor  \pend
\section*{CAPVT  V. }\pstart 275 detrahi. Iam vero, quia haec capacitas non ex aetatis senio et beneficio proprie pendet, imminutum est annorum tempus et spatium a Paulo hoc loco praescriptum, voluerunt que homines viderispiritu ipso Dei sapientiores. Itaque concilio Chalcedonensi constitutum est, vt statim post 4o annum çtatis eligi mulieres ad istum Diaconatum possent est in cano. Diaconissas 27. quaest. 1.id quod est confirmatum concilio Carthaginensi 4. vt est in cano. sanctimoniales 2o. quest. I.in quo postea va rie peccatum est, vt docebimus.  \pend\pstart Addit Paulus. Vnius vir vxor. Hoc vero probari potest exemplo Annae, de qua Luc.1.ver.36. Respondet autem illi, quod de episcopis et Diaconis agens idem Paulus dixerat. Vnius vxoris vir. Itaque ex illo loco facile iam intelligi potest quae sit vis huius loci sententia atque mens. Primum enim non damnat nuptias Paulus siue primas, siue secundas, siue tertias siue vlteriores, Hanc enim esse diabolicam doctrinam supra asseruit cap.4. vers.1.2.3. Deinde non iniicit hunc laqueum viduarum conscientiis, quasi praecipiat, vt in viduitate permaneant, etiam eae, quae denuo vel iterum nubere vellent: neque hoc loco etiam viduitatem aut coelibatum tanquam per se sanctio rem praefert honesto conugio Pauius, vti neque I.Corinth.7. vers.8.9. 34. et 35. Denique cùm neque ante Pauli aetatem, neque eius seculo vsquam terrarum obtineret: vt vna mulier eodem tempo re duobus nupta esset viris (etsi inter Iudaeos vir vnus solebat ex vitiosa patrum imitatione duas vel plures vxores ducere) satis intelligitur eam dici plurium maritorum vxorem esse, quae quan\pend
\section*{AD I. PAVL. AD TIM. }
\marginpar{[ p.276 ]}\pstart quam est ex iniusta causa repudiata a viro, tamen eo viuo alteri non nubet, sed propter vinculum coniugii, et Dei praeceptum, quandiu prior ille maritus, a quo iniuste repudiata fuit, viuit, estque superstes, alterius coniugio abstinet. Quae enim ex iniusta, vel potius nulla causa repudiatur, manet vxor: et vel reconciliari debet vel innupta ma nere 1. Corinth.7.vers.1o. Neque tamen impeditur, quae ita immerit o repudiata est a viro, viro ipso ad alteras nuptias conuolante, quin et ipsa alii nubere possit. Sed hoc praesertim in ea, quae in Diaconatum allegenda est, requirit Paulus, quoniam maioris continentiae signum et exemplum in ista muliere cerni debeat. Erit autem hoc certissimum, si immerito repudiata, tamen alteri non nupsit, sed continuit. Sed quid si merito illa repudiata est, veluti propter adulterium, vel quod ipsa sit desertrix coniugii? Nemo dubitare tunc potest, quin ea ad hoc munus ecclesiasticum eligenda omnino non sit, sed ab hochono rene drcenda.  \pend\pstart Vult autem Paulus, vt quae est allegenda, fuerit Vxor. Ergo si fuit vel concubina, vel pellex, vel meretrix eligi non poterit. neque enim sine magno ecclesiae offendiculo tales cooptari in vllam ceciellaiticam dignitatem possunt, quamlibet poeniteant.  \pend\pstart Ait Paulus. Quae fuerit. Quid si adhuc illa sit vxor, etiam in illa aetate sexagenaria, num allegi et ascisci poterit, cùmillud, quod de episcopo supra dicitur, Sitepiscopus vnius vxoris vir, locum habeat: siue adhuc sit ille maritus et vir vnius vxoris: siue mortua vxore maritus fuerit et iamesse  \pend
\section*{CAPVT  V. }
\marginpar{[ p.277 ]}\pstart desierit. Resp. non esse ex mente Pauli eligendam sexagenariam, quae adhuc vxor manet, quia de viduis loquitur Paul': et variis illis muneribus et cu ris nimium illa distraheretur, si iam alligata sit viro suo: et tamen cura etiam pauperum illi sit gerenda. Vtrique enim oneri ac curae vna satis esse non potest. Nam quae nupta est, inquit Paulus, curat ea, quae sunt mundi, et quomodo placitura sit viro 1.Corinth.7.vers. 34. Ergo ea, quae funt Dei et pauperum, libere, quemadmodum Diaco nissam decet iam curare aut administrare non potest. Alia vero et diuersa est ratio viri et episcopis qui quanquam cùm est maritus, curat, quomodo sit vxori placiturus. 1.Corint.7.ver.31. tamen non perinde in ea cura subiicitur vxori, quemadmodum vxor viro. Itaque manens maritus, munus ecclesiasticum gerere ac obire potest.  \pend
\phantomsection
\addcontentsline{toc}{subsection}{\textit{to In operibus bonis idonec testimanió ornata, si liberos educauit, si fuit hospitalis, si sanctorum pedes lauit, si afflictis subuenit si omne bonum opus est assidüe sectata. }}
\subsection*{\textit{to In operibus bonis idonec testimanió ornata, si liberos educauit, si fuit hospitalis, si sanctorum pedes lauit, si afflictis subuenit si omne bonum opus est assidüe sectata. }}\pstart Αύξησις Amplificat enim ornamenta futurae et eligendae Diaconissae, quod eam publica voce totius ecclesiae et aliorum etiam extraneorum hominum, qui a fide Christi sunt alieni, consensu vult habere idoneum testimonium et probitatis, et charitatis, et beneficentiae, et curae et diligentiae Hoc idem supra in Episcopis requirebat cap. 3vers.7. ex quo apparet quanti sit momenti publica illa fama et opinio de nobis concepta, modo non temere, sed idoneis argumentis et testibuz fulta spargatur. Est enim velut quidam odor?  \pend
\section*{AD I. PAVL. AD TIM. }
\marginpar{[ p.270 ]}\pstart quem nobiscum circumferimus: et siquidem illa bona est, bonus est et suauis gratusque ille odor, quem emittimus: sin vero illa est mala, est etiam odor a nobis effluens malus, et male olens qui alios, cum quibus versamur, offendit. Itaque iubet Paulus, vt quaecunque sunt ἔυφημα.ι. bonum nomem bonis conciliant, et cogitemus et procuremus. Philip. 4. vers. 8. Id quod etiam Hesiod. lib1. ἐργ. καἰ ημές. praeclarissime versibus monet. Ἀύξνσιν superiorem sequitur ἐξηγησις. Aliquot enim exem plis, et generibus actionum explicat atque illustrat quod generaliter antea dixerat Paul. Affert autem quatuor eorum bonorum operum species, quae maxime futuro muneri et allegendae diaconis sae quadrant atque conueniunt. Cùm enim pauperum cura ilircommittendant, quae res et Misericordem animum, et diligentem atque etiam de missum requirit, ea bona opera commemorat Pau nus, ex quibus fuiunnour amius facile argui et deprehendi in nobis potest.  \pend\pstart Diligentiae signum erit, si liberos proprios studiose, et diligenter ipsa educauit et curauit.   \pend\pstart Minicriturs te hamanus mui deprehendetur animus, si hospitalis fuit, et afflictis subuenit.  \pend\pstart Demissus autem illius animus, si sanctorum pedes lauare non est dedignata.  \pend\pstart Denique pietas et probitas eius intelligetur, si quodlibet opus bonum assidue et studiose consectata est. Haec est breuitaer mens Pauli.  \pend\pstart Quod si haec tam diligenter in iis, qui ad minora ecclesiae munera vocantur obseruari vult Paulus, id est, Dei spiritus, et ab illis requirit has virtures, quid sentiendum statuendumque nobis  \pend
\section*{CAPVT  V. }
\marginpar{[ p.273 ]}\pstart est de iis, qui ad maiora munera ecclesiastica accersuntur et eliguntur? Haene virtutes ab iis abesse debebunt? Minime sane. Itaque in iis et diligentia et misericors animus, et demissus et pius spectetur.  \pend\pstart Quod ait. Si liberos educauit mouet dubitationt quasi Paulus steriles sexagenarias, quae vxores fuerunt nullósque liberos pepererunt, ab hoc munere arceat, illisque tanquam hanc notam inurat propter sterilitatem, vt eligi ad hoc munus non possint. At ee saepe sunt aptissimae, An non igitur quae fuerunt steriles allegi et cooptari poterunt? Sane poterunt. Neque enim vniuersum iudicium, quod de mulierum piarum diligentia haberi potest, Paulus restringit ad vnicam curam, et educa tionem susceptorum liberorum: sed hoc tanquam verissimum et optimum inter caetera earum diligentiae testimonium protulit. Ergo quae mulieres nunquam liberos sustulerunt in coniugio, aut etiam quae nunquam fuerunt coniuges, possunt, si erunt idoneae, in diaconissas eligi.  \pend\pstart  Quod ait. Si fuit hospitalis, si afflictis subuenit, mouet etiam quaestionem, quasi solas diuites eligi et cooptari ad hoc munus velit Paulus. Illae enim sunt, quae et extraneos excipere apud se et suis opibus afflictis subuenire ἱ̓πόρκεσιν enim dixit Paulus id est, quantum satis est, sumministrauit egenis) possunt. Quae autem sunt egenae et ipsae 41 pauperes idem praestare nequeunt. Ergo ab hoc munere summouebuntur, quas tamen saepe ma4. gis esse ad hoc ministerium idoneas, quam dites mulieres certum est, et experientia comprohatum Resp. Non excludit pauperes Paulus imô maxi\pend
\section*{AD I. PAVL. AD TIM. }
\marginpar{[ p.A0O ]}\pstart me de iis allegendis agit hoc loco. Sed neque etiam diuites ab eodem munere suscipiendo arcet: caeterum quae pauper est, potest tamen et misericors et benefica, et hospitalis esse, et afflictis subuenire, si quantum in se est, illa egenis praestat, si corporis ministerium exhibet, si operam suam illis impendit et praebet, quando pecuniam confer re non potest. Denique et vidua paupercula, quae minuta duo in gazophylacium misit, maiorem, quam alii Iudaei opulentissimi eleemosynam contulit: et Petrus ac Ioannes cùm pecuniam non haberent, summo tamen beneficio claudum affecerunt, vt est Luc.21.vers.2.Actor.2.  \pend\pstart Quod idem Paulus ait. Quae lauit pedes.) sumptum videri potest partim ex Gene.18.vers.4.19. vers.2.43.vers.24 vbi haec eadem ratio dicendi vsurpatur: partim ex recepto tunc more inter Orientales homines. Nam cùm qui fere pedibus iter faciunt, fole, nimióque labore incalescant, sunt eorum pedes abluendi vt refrigerentur, recreentur et refocillentur. Qui autem equo incedunt et iter conficiunt, non egent fere hoc auxilio et ministerio. Ergo ex consuetudine loci hoc annotauit Paulus. Hoc autem ipsum magnam animi demissionem significat, cùm vilissima et sordidissima hominis pars fere sint pedes: quos quaecunque ita contrectare, et curare non dedignatur, certe nec vilia quaelibet alia munera et ministeria, quae pauperibus aegrotantibus impendi necesse est, obire eadem vnquam recusabit. Christus ipse in signum humilitatis lauit pedes discipulorum Ioan 13. Idem fecit et illa mulier, quae lacrymis suis lauit pedes Christi. Luc.7 vers38.  \pend
\section*{CAPVT  V. }
\marginpar{[ p.281 ]}\pstart 11 Porro iuniores viduas recusas postquamenim lasciuire coeperint adnersus Christum, nubere volunt.  \pend\pstart Αʹντίθεσις est, ex qua superior sententia non tantùm illustratur: sed etiam probatur. Ex eo enim, quod subiicit Paulus, apparet, cur demum sexagenarias mulieres allegi voluerit et praescripserit. Varia enim offendicula, quae a iunioribus oriuntur, id suaserunt. Itaque non ea est causa cur sexagenariae demum eligantur, ne diutius alantur diaconissae ab Ecclesia: sed ne scandalum praebeant vitae mutatione et lasciuia.  \pend\pstart Vult autem vt recusentur iuniores, etiamsi se offerant, et paratae sint hoc munus acceptare: vel etiam si alias idoneae viderentur. Quod vero, ait Paulus παραιτο͂, non est eo modo accipiendum, quasi in solius Timothei arbitrium totam hanc electionem conferat, quam tum ex Presbyterii, tum ex totius. Ecclesiae consensu fieri solitam esse certum est: sed ita loquitur Paulus, quod pastor omnium Ecclesiae actionum et electionum moderator est. Notandum est autem, quod viduas iuniores tantùm recusari et electas a munere diaconissae deponi ac reiici vult, etiamsi postea nubant: non autem praecipit eas excommunicari et eiici de Ecclesia: vel vsu et participatione Sacramentorum arceri et prohiberi, quemadmodum de Nonnis et Virginibus postea Synodis ineptissimis sancitum est. Atque  vtinam haec Pauli moderatio semper in Ecclesia locum habuisset: sed postea longius progressi sunt Episcopi, maxime postquam nata est superstitiosa illa de coelibatus merito et  \pend
\section*{AD I. PAVL. AD TIM. }
\marginpar{[ p.202 ]}\pstart dignitate opinio. Primùm enim sanciuerunt Ecclesiastici canones, vt vidua iunior, atque etiam vt Monialis, quae post votum nubit, inter digamos tantùm censeretur, et reponeretur, can. Quotquot 27. quaest. I Post autem inualesc ente superstitione sancitum est, in eas aliquid acerbius. Voluerunt enim, vt si qua Monialis post votum nupsisset, vt non tantùm Sacramentorum vsu priuaretur: sed etiam a toto Ecclesiae coetu excommunicaretur, quia votum de coelibatu et continenitia vibiarat, quou Deo gratum esse falso et insulse crediderunt. Est haec poena constituta in can. Virgines 27. quaest.1.qui ex consilio Elibertino decerptus est.  \pend\pstart Sed iam quae sint illa offendicula, quae praebent viduae iuniores ad Diaconatum electae videamus. Sunt autem imprimis haec duo nimirum, quod I Institutum vitae genus in magnam contumeliam Christi commutant magnam animi leuitatem et carnis lasciuiam demonstrantes. 2 In hoc ipso vitae genere suscepto si perstent, magna tamen Ecclesiae praebent offendicula, dum sunt Otiosae, Discursatrices, Nugaces, Garrulae. Caeterum cùm improbat Paulus quod nubant hae viduae, quae iam electae sunt, non ideo facit, quasi secundas nuptias damnet, vel etiam tertias, et quartas, id quod Montanistarum haereticorum, aut Papistarum Montanizantium proprium est, sed quod tale quid fit, oritur ex animi magna, tum Leuitate, quae graues matronas et eas quae sunt hoc munus publicum adeptae, non decet: tum Lasciuia siue infreni ac indomita carnis petulantia, quam tamen abiecisse debuerant, et exuisse. Id quod etiam cum Chri\pend
\section*{CAPVT V. }
\marginpar{[ p.283 ]}\pstart stiani nominis et doctrinae contumelia coniunctum est. Irridetur enim Dei Ecclesia, quod tam leues homines habeat, vt qui munus a se susceptum statim abdicent, relinquant, et huic mundum et carnis delicias praeferant. Quasi Christus a auis parui nat de contemnatut aperte etiam iis, ipsis qui tamen maiores in fide progressus fecisse videbantur.  \pend\pstart Vox καταςρηνιάζην cum sequenti, id est, χειςο construitur. Est autem ςρηνιαζεην ἀπὸ τωὺ ςερεῖν vel σρηνὲς deductum, quod durum significat et pertinax: qualia sunt animantia nimium saginata. Itaque sunt petulantiae plena et indomita, vnde cibus est illis detrahendus, vt in ordinem cogantur et pareant, quemadmodum et Xenophon et ipla rerum experientia faciendum docet. Hac igitur similitudine vtitur Paulus, vt carnis nostrae lasciuiam eo turpiorem ac foediorem esse demonstret, quod brutis et petulantibus animalibus conparatur.  \pend\pstart Chrysostomus autem interpretatur hanc vocem fornicari, sed Apocalyp.18.vers.7. et 9. explicatur eo significato, quo nunc accipimus.  \pend
\phantomsection
\addcontentsline{toc}{subsection}{\textit{12 Ex eo damnandae quod primam fidem reiecerint. }}
\subsection*{\textit{12 Ex eo damnandae quod primam fidem reiecerint. }}\pstart Αʹὐξισις est. Amplificat enim harum iuniorum viduarum lasciuiam a consequenti, quod hac ratione, et vitae genere tandem sibi extremum exitium acceasant. Nam eo tandum prolabuntur, vt Munus a se susceptum deserant, et Fidem ipsam primam Christo datam abnegent. Voces autem has πρώτηνπιςιν varie interpretes accipiunt. Alii  \pend
\section*{284 AD I. PAVL. AD TIM. }
\marginpar{[ p.]}\pstart de voto non nubendi sumunt, quasi ab iis interponi huiusmodi votum solitum fuerit, cùm eligebantur: quod nullo certe scripturae loco probari potest. Alii de fide Christo in Baptismo data:alii de munere Diaconatus acceptato, quod diligenter se praestaturas esse, in eóque perseueraturas pollicebantur. Quod postea cùm nubendi pruritu et libidine deserunt, fidem a se datam violant, ac frangunt. Sed ea sententia magis placet arg. vers. 15.infra, quae fidem primam refert ad eam promissionem, quae primum fit a Christo nobis in Baptismo, quam reiicere dicuntur, quae a fide Christi et professione Euangelii postea deficiunt, et ad Gentilismum, vel priorem idololatriam reuertuntur. Sic Athanas.explicat liber 6. de Trinit. Hierony. in praefat. ad Tit. Vincentius Lyrĩnensis, et est similis locus Apocalyps.2.v.3. Veritae enim probrum inter Christianos propter eam animi sui leuitatem malunt prorsus a Christo desciscere, quam husuimodi contumelias audire et pati. Nicolao, de quo est in Actis Diacono id sane accidisse narratur, qui ne opinionem leuitatis et temeritatis inter Christianos subiret, maluit condere nouam haeresin, nouum sibi coetum colligere, et ab Ecclesia discedere.  \pend\pstart Quae autem hic narrat Paulus, rerum ipsarum experientia edoctus didicerat, quemadmodum ipse testatur.  \pend
\phantomsection
\addcontentsline{toc}{subsection}{\textit{13 Simul autem etiam otiosae discunt circumire domos: imo non solùm otiosae, sed etiam nugaces et curiosae, garrientes quae non oportet. }}
\subsection*{\textit{13 Simul autem etiam otiosae discunt circumire domos: imo non solùm otiosae, sed etiam nugaces et curiosae, garrientes quae non oportet. }}
\section*{CAPVT  V. }
\marginpar{[ p.285 ]}\pstart Altera superioris amplificationis pars, in qua iuniorum viduarum in munere Diaconatus permanentium offendicula et mala describuntur. Tria vero enumerat Paulus, non quod sola haec obseruari possint: sed quod maxime soleant. Deinde quod ex his paucis, quae sint reliqua incommoda, quae ex eodem fonte nascuntur, id est, ex huiusmodi iuniorum mulierum electione, facile iam per se potest quilibet intelligere. Ergo eas, etiam dum hoc munus administrant, docet esse, Otiosas (cui vitio discursationem adiungit) et Nugaces, et Garrulas. Garrulitatis autem vitium breuiter definit Paulus dum ait, çasdem loqui quae non ubuon: porto noed viela luuo parsunt magna inter Christianos offendicula, quod sint contra diuinae legis praeceptum 9. quod est, Non falsum testimonium aduersus proximum tuum dices, quódque ex his infinita mala nascuntur.  \pend\pstart Sunt otiosae) At obstat, quod eas ipsas postea vocat σεμ́ργύς, id est, plusquam oportet, occupatas et curiosas. Respond. Haec duo vitia inter se recte conuenire. Nam nihil agentes illae rerum, quae ad se pertinent, sunt in alienis perscrutandis curiosae. Itaque ἀργαι appellantur, quod nihil suum agunt. Iεριίρyοι vero quod aliena negotia plus quam sua curant, atque etiam plus quam oportet, quemadmodum ait ille. Excussis propriis aliena negotia curo. Imo vero etsi saepe habent domi, quae agant: magno et turpi animi veterno illa negligunt, vt nimirum de rebus alienis co gitent, et disputent.  \pend\pstart Circumeunt domos) Annexum pene semper otio et segnitiae vitium, nimirum discursatio. Hoc  \pend
\section*{AD I. PAVL. AD TIM. }
\marginpar{[ p.286 ]}\pstart autem praetextu maxime id facere videbantur istae mulieres, quod tanquam publicae quaedam personae, et pauperum curam habentes in alienas aedes possent, solerentque liberrime ingredi 2. Timoth.3.vers.6. Videtur autem etiam his verbis detractrices easdem appellare Paulus, more dicendi Hebraeis vsitatissimo, qui et detrectare et discurrere eadem voce לx7 Ragal dixerunt, quod detractores iugiter eant, et referant verba hinc inde ex aliis audita. Peccant igitur hi in nonum legis diuinae praeceptum.  \pend\pstart Imo vero) Amplificatio est superioris mali per adiectionem et προθεσιν. Non solum enim sunt otiosae: sed etiam Nugaces, et Curiosae vel plus quam oportet negotiosae φλυάρυς appellat Pau. lus, quos Latini Nugaces, quae vox videtur deriuari a voce. φλύαξ, ακος, quae tumulentum hominem significat, quod haec sint duo maxime inter se connexa vitia, Temulentia nempe, et Nugacitas. Itaque vt Plutarch. liber 8. Symposiacon Problema. 1. nihil aliud esse videtur gλuzeία quam λήρησις πάροινος. Nugari quoque, ait idem, proprie nihil aliud est, quam vano et futili sermone vti, et eo, qui vel mendax est: vel qui quanquam verus est, ad propositum tamen nihil facit et attinet. Vnde postea rixae, iurgia, et contumeliae oriuntur, atque varia offendicula, Vocem autem ipsam nugae Latinam ab Hebraea nan hagah quod more auium garrire significat deduci putat Iosephus Scaliger in liber  Varronis de lingua Latina. Qanquam autem Plutarchus in liber  πεὶ αδολεσχίας videtur facere φλυκρίας siue ληρήσεως, duo  \pend\pstart genera, Temulentiam quamὁ̓ηωσιy vocat, et ἀδο\pend
\section*{CAPVT  V. }
\marginpar{[ p.287 ]}
\marginpar{[ p.O ]}\pstart λσ χίαν quam garrulitatem vertunt, quas ita distinguit, vt temulentia sit nimia inter pocula vel inter potandum orta ex vino sumpto loquacitas. Garrulitas vero sit stultiloquium nimiaque loquacitas vbique, et ex qualibet, vel etiam nulla oblata occasione. Itaque garrulus in foro, in theatro, in deambulatione, domi, interdiu, et noctu nugatur, et loquax est. Temulentus autem tantùm in mensa et inter pocula: ista tamen distinctio etiam hunc Pauli locum illustrat, et excoule. oeu videtur vtranque quidem, ied maxime secundam speciem hoc loco damnasse Paulus vti etiam Psal.14o. vers.1o. aperte inter vitia pessima recensetur ista nugacitas,  \pend\pstart Curiosas vocat tamen eas, quae plura quam deceat, scire affectant, tum etiam, quae diligentius et penitius ea, quae sciri possunt, curant et inuestigant, quam sit necesse. Itaque sunt illae nimium negotiosae, sed vbi non oportet. In quibus autem oportet otiosae. Alio nomine πολυπράγμονες eaedem appellantur. Quanquam Plutarch.in libro de Curiositate eos tantùm videtur definire πολυπράγμονας, qui in alienis erratis siue peccatis, vnde alteri nascitur aliqua infamia, curiosiùs indigandis occupant sese. At siue in peccatis, siue in aliis etiam rebus, quae ad nos non pertinent, inuestigandis nimium laboremus, hic πράργων nomine a Paulo significamur. Quanquam vero haec vitia faere sunt omnibus foeminis communia, maxime tamen eas dedecent, quae caeteris praelucere vitae exemplo debent. Inquirere autem in alienum factum, quatenus illud ad nos spectat: vel Dei gloria proximiquextilitas postulat, vel inde promouetur: denique  \pend
\section*{AD I. PAVL. AD TIM. }
\marginpar{[ p.286 ]}\pstart quantum ex muneris nostri ratione vel Christianae charitatis officio tenemur, non est curiositas: sed est officium, quod est sane laudandum. Eo denique pertinere videntur postrema huius versiculi verba loquentes quae non oportet. Nec enim de alio loqui prohibemur semper, sed videndum est. Quid, Quomodo, et Quatenus loqui de eo liceat.  \pend\pstart Priusquam autem totum hunc de Diaconissis locum finiamus, duo scitu necessaria sunt enodanda. Primùm quidem, vtrum eas hodie in reformatis Christianisque Ecclesiis retinere sit operc̨- pretium, cùm nos quam maxime fieri potest, proxime ad veterem et Apostolicam Ecclesiam accedere tum in doctrina: tum etiam in disciplina et politia Ecclesiastica oporteat. Quid igitur? Respondeo aliquod discrimen esse nobis statuendum inter doctrinam veteris Ecclesiae, et inter generalem qua illa vsa est, disciplinam. Nam cùm doctrina Apostolica sit Christi ipsius vox et sapientia, cui tanquam vero et vnico fidei et salutis fundamento innititur Ecclesia, certe ab ea doctrina vniuersa ne tantillum quidem a nobis recedi vllo tempore debet. Itaque nulla occasio afferri vel fingi vnquam potest, propter quam yel in vno apice, quanquam minimo, doctrina Apostolorum immutetur Matth.5. vers.18. At illus disciplinae vniuersae, quam illa vsurpauit, non eadem omnino ratio est. Cùm enim sit disciplina tanquam decorum quoddam vestimentum ad res, tempora, personas accommodatum, fit, vt quemadmodum haec variari possunt:ita possit quoque altquibus in rebus mutari etiam disciplina vetus. Tantum distinctionem hanc adhibeamus, vt in  \pend
\section*{CAPVT  V. }
\marginpar{[ p.289 ]}\pstart disciplina ipsa Ecclesiastica, distinguamus ea, quae sunt Fundamentalia politiae Ecclesiasticae, ab Accessoriis et leuioribus. Fundamentalia vero sunt haec in politia Ecclesiastica. Primùm vt sit aliqua politia Ecclesiastica, et quidam ordo in Del domo, qui tollat confusionem. Deinde, vt legitimae personarum vocationes, per quas Dei verbum legitime administratur, charitasque et sanctitas vitae exercetur in Ecclesia, retineantur, veluti Pastorcomreiojteri, Diacom  \pend\pstart Tertio, vt ii electi legitime et ex praescripto Dei verbo suo munere fungantur, vel deponantur, et remoueantur.  \pend\pstart Accessoria dico, quae, vt haec fiant et obseruentur, in Ecclesia quaque pro tempore, personis et locis statui possunt. Accessoria incolumi manent disciplina in Ecclesia ipsa variari possunt. Maneat enim semper disciplinae fundamentum et pars praecipua, sed tantùm ea mutentur, quae commodioris illius obseruationis gratia tunc visa sunt necessaria: veluti vt futurus et eligendus Episcopus a tribus Episcopis, non a duobus tantùm eligeretur. Hoc mutari potest incolumi Ecclesia et disciplina. Vt in vrbe Metropolit. electio fiat et examen futuri Episcopi, non alibi. Hoc etiam mutari potest. Itaque vetus Ecclesia in iis rebus saepe diuersam formam habuit et secuta est, incolumi nihilominus Dei verbo, ipsóque fundamento Ecclesiasticae disciplinae retento. Ex quo fit, vt cùm hae Diaconissae fuerint tantùm subsidiariae Diaconorum manus et operae (propter egenorum et aegrotorum ingentem numerum et multitudinem pene infinitam) quo tempore Nosodochiiss  \pend
\section*{AD I. PAVL. AD TIM. }
\marginpar{[ p.290 ]}\pstart Xenodochiis, et Ptochodochiis careret Ecclesia postea pauperum numero decresente et Ptochodochiis constitutis, non fuerunt eaedem Ecclesiis necessariae. Itaque neque sunt hodie a nobis retinendae, nisi eadem et par aliqua causa, qualis illo seculo fuit, occurreret, et eas retinere suaderet. Quae cùm nulla sit, cessante causa, cesset effectus necesse est. Eodem modo videmus veteres habuisse plures officiorum Ecclesiasticorum gradus (vti apparetex aliquot Synodis etiam quando adhuc erat tollerabilis Ecclesiae forma) quam nunc habeamus. Itaque quemadmodum illae dignitates tanquam nobis inutiles hodie propter diuersam locorum rerum, et personarum inter nos et illos rationem explosae sunt, et cessant, qualis Subdiaconatus, Acolythatus, Archidiaconatus, Archipresbyteratus: sic et hic Diaconissarum muIierum or do desiit, nullo prorsus detrimento vel in doctrina, vel in disciplina ipsa ab Ecclesia, accepto.  \pend\pstart Secunda vero quaestio, quae ad hunc quoque locum maxime pertinet, est huiusmodi. Vtrum recte a Papistis sanctimoniales sint in istarum Diaconissarum, quas nullius iam esse in Ecclesia vsus videbant, locum substitutae: at que vtrum eorum consilium et institutum sit probandum, et sequendum necne. Respond. Minime id quidem, quia in eo statu hominum et genere instituendo nihil nisi plane Satanicum, superstitiosum, et in Deum ipsum blasphemum a Papistis inductum est. Primùm, quia coelibatum illis suis sanctimonialibus tanquam legem ad vitam aeternam assequendam necessariam, et verissimae sanctitatis  \pend
\section*{CAPVT . V. }
\marginpar{[ p.291 ]}\pstart partem indicunt. Deinde, quod in eo ipso cultum Dei situm esse contendunt. Tertio quod in iis plis a Patrum consuetudine recesserunt, et a meliori more in Ecclesia iam recepto et inueterato.  \pend\pstart Quod vt melius intelligatur, rem totam fusius sic explico. Cùm κακοζηλία quadam homines omnia quae Apostolorum temporibus obseruata fuerant, sibi putarent imitanda, neque locorum, neque temporum, neque rerum dissimilium rationem haberent, etiam et ipsi suas Diaconissas retinere praecise voluerunt. Sed cùm posterioribus temporibus et mutatis circunstantiis earum nullum in Ecclesia vsum esse cernerent, magno Ecclesiae malo, magnóque conscientiarum detrimento pessima tamen eius ordinis instituendi ratio reperta est.  \pend\pstart Instigabant virgines Christianas quidam Episcopi, vt se Deo dicarent, atque, vt illi loquebantur, consecrarent. Deinde vt in signum istius consecrationis se velo, tanquam castitatis et verecundiae testimonio, semper tegerent, et obnubilarent. Quod ex Pauli consilio facere se et iubere iactitabant 1I.Corinth.11. Hinc variae Hieronymi, Augustini, Ambrosii, Fulgentii epistolae ad varias virgines scriptae, quia etiam ante aetatem horum omnium huiusmodi votunii ni magno apud Ecclesiam pretio iam haberetur. Tamen, vt ex iisdem scriptoribus apparet, et imprimis ex Cypriano et Tertulliano in liber  de cultu et habitu virginum, hoc totum voti genus etiam conceptum et nuncupatum adhuc liberum erat. In priuatis enim parentum vel cognatorum aedibus huiusmodi virgines, quae sanctae vel consecratae  \pend
\section*{AD I. PAVL. AD TIM. }
\marginpar{[ p.292 ]}\pstart Deo dicebantur, manebant non in vllo monasterio et simul, etiamsi in eadem ciuitate plures huiusmodi essent. Atque etiam saepe ipsae domos seorsim, quaeque suam, habebant, ex quibus vel solae, vel comitatae cùm volebant, et cùm libebat egrediebantur, nihilque a caeterarum virginum conditione differebant, nisi quod eas caeteris castitatis exemplum esse oporteret. Id apparet ex can. Nuptiarum 27. quaest. 1. Quin etiam postea eaedem, si volebant, poterant nubere. Et quanquam in eo opinione quadam leuitatis apud pios grauabantur propterea tamen neque sacra Coena iis interdicebatur, neque etiam ipsae Ecclesia excommunicabantur, vti postea factum est. Hoc igitur mali initium et exordium fuit, sed latiùs postea Satanae astu prouectum atque propagatum est. Primùm quidem his nominibus vocari coeperunt, vt Moniales, Sanctimoniales, Virgines et Nonnae dicerentur. Postea vero Monachae et ad prophanarum imitationem, virgines Vestales: ad quarum instituta reuocatae vetitae sunt, ne priuatas aedes haberent, id est habitacula vel cellulas feorsim incolerent: sed iussae sunt vt in communi habitarent et viuerent. Hoc primum fuit, deinde fancirum est, ne aliae Sanctimoniales censerentur, quam que uiedi Monasterio certo, certaeque  regulae sese in perpetuumaddixissent can. Perniciosam 18 quest. 2.concil.tamem Agath. vetabantur ista Monacharum Monasteria et cellae (quibus istae concludebantur) iuxta Monachorum et virorum caulas et lustra propter periculum et offendiculum aedificari, vti est in cano. Monasteria puellarum quaest. 2Coepta sunt autem aedificari earum coenobia ma\pend
\section*{CAPVT . V. }
\marginpar{[ p.293 ]}\pstart xime post annum Domini 700. Antea enim etsi fuerant quaedam, omnino tamen pauca et rara erant foeminarum et virginum Monasteria extructa. Atque tres tantùm earum ordines concessi sunt primùm, et probati eodem illo seculo, nenpe post annum Domini 700. quod tunc demum ex Synodo Hispalensi cap.11. ad instituta Monachorum coeperunt reuocari.i. ex magna superstitione regi et institui Moniales: id quod antea non fiebat. Tres autem illi Monacharum ordines primùm hi fuerunt, ordo Basilii, Benedicti, Augustini, extra quos nulli foeminae esse licuit Moniali, quodcunque  votum vouisset et suscepisset can. Perniciosam 18. quaest.2. Hodie vero infinita prorlus lune carui genera, et adeo diuersa, vt nihil inter se simile pene habeant. Sunt enim hoc nostro seculo Monachae, Benedictinae, Bernardinae, Franciscanae, Dominicanae, Carmelitanae, Cleranae, Augustinianae, et aliae plures species.  \pend\pstart Ita vero primùm fuerat in suscipienda professione constitutum, vt nulla ante 4o. aetatis annum votum vllius ordinis susciperet, aut si suscepisset ne obliga retur, can. Diaconissam 27.q. 1.can. Sanctimoniales 20.q.1. Deinde hoc tempus, tanquam longius esset, restrictum est, atque constitutum vt statim post vicesimum quintum aetatis annum posset virgo quaelibet votum Monasticum suscipere, et 20.quaest. 1. Sed et hoc ipsum quoque annorum spatium postea est imminutum, atque sancitum, vt statim post pubertatis annos possent mulieres virginitatem perpetuam vouere: quanquam consecratio earum differebatur vsque  ad annum aeta\pend
\section*{AD I. PAVL. AD TIM. }
\marginpar{[ p.294 ]}\pstart tis vicesimum quintum: sed postea receptum institutumque est, vt statim post annos pubertatis, id est, post duodecimum aetatis annum posset, quae votum ordinis et vitae Monasticae suscepisset etiam a Praelato consecrari cano. Puella 2o. quest. 2. quem ex concilio Triburiensi sumptum esse volunt. Ante autem duodecimum aetatis annum, ne ex parentum quidem iussu et confirmatione, votum ab iis factum ratum esset constitutum est. Denique ne qua inuita in Monasterium detruderetur, sed si qua ergastulum illud iubentibus parentibus esset ingressa, decimo quinto anno aetatis interrogetur a Praelato de voto, cui si ipsa non assentitur, ne votum ante illam aetatem factum omnino valeat. Sic can. Firma can.Sicut et seq.20. quaestione I.sancitum est, a quibus tamen hodie plane discessum est. Coguntur enim etiam antequam pene vsum rationis habeant miserae puellulae, vt se in Monasterium aliquot detrudi sinant, quod vel non satis formosae sint, vel caeteris liberis pars hereditatis maior accrescat his in pistrinum istud et Monasterium relegatis.  \pend\pstart Quae virgo vel vidua votum ipsa sibi imposuerat, et habitum Monialis assumpserat domi suae vel in aedibus patris:non autem ab Episcopo vel praelato illud acceperat, primum poterat illud, si ronet portea deponere, neque cogebatur in eo vitae statu inuita permanere, nisi fuisset postea ab Episcopo consecrata. Post autem aliud decretum est. Nam neque deponere illud velum potuit, ac ne ea quae illud priuatim domi suae assumpserat: sed cogetur etiam ante consecrationem factam in eo statu et habitu permanere perpetuo vitae tempore, et Monasterium ingredi can. Si quis sa\pend
\section*{CAPVT V. }
\marginpar{[ p.295 ]}\pstart cra 27. quaest.1.et can. Mulieres cano. De viduis et Puellis eod. Ita nimirum paulatim auctae sunt luperititiones, et dursor conscientiis laqueus iniectus.  \pend\pstart Cùm sanctimonialis a praelato consecratur, stola alba induitur, et chrismate vngitur, vt ait Alcuinus. Primus autem eam consecrari instituit Basilius magnus, vt plaerique volunt. Post autem eadem foemina vel virgo tondetur. Quanquam ne tonderetur Christiana mulier, et concilio Gangreno interdictum fuerat, et expresso Pauli praecepto, quod ipsius naturae voci et sensui consentaneum est 1.Corinth.11. Socrat.libus 2. cap.43. Sed boni scilicet illi patres omnia iura diuina et nuinanla, vt iuas luperititiones Itabilirent, peruerterunt: et susque deque habuerunt. Nam et ex cano. Quaecunque distinct.3o.ita in Ecclesia fuisse olim obseruatum quemadmodum dicimus, constat.  \pend\pstart Consecratae Moniales habent Monachos qui illis praesint, nimirum oues lupis committuntur, (quanquam tamen vt ab earum accessu isti arceantur ex concilio Hispal. et can. in decima 18.quaest2.sancitum est, sed frustra) Post istam consecrationem nubere istae mulieres ex iure Pontificio iam non possunt. Qanquam olim infamia tantùm quadam in Ecclesia notabantur, et tanquam Digamae tantùm censebantur, si quae nupsissetst, atque adeo ex concilio Ancyrano et can. Quotquot 27. quaest. ita erat constitutum. Postea vero ipsa superstitione nullum iam modum seruante prorsus ab Ecclesia sunt excommunicatae, quae post istam consecrationem nupsissent can. virgines 27.  \pend
\section*{AD I. PAVL. AD TIM. }
\marginpar{[ p.296 ]}\pstart quaest.1.Atque ita miserrimus conscientiis hominum laqueus semper durior, et maior inuectus est: et προς το ἀδύνατον tandem contra rerum naturam miserae illae mulieres sunt obligatae, nempe vt perpetuo coelibes manerent, quantunuis intus aestuarent et magno igni vrerentur.  \pend\pstart In Monasteriis autem illis et coenobiis, tanquam carceribus, conclusae horas canonicas, quas Papistae appellant, et aliquot Psalmos tantùm cantillant. Nam neque  Missam cancre ipsae possunt, neque vasa sacra aut vestimenta, quibus in Missa sacrificuli vtuntur, contrectare, ac ne quidem sine piaculo attingere audent cano. Sanctimonialis et can seq. distinct.23.Atqui etiam seculo Hieronym. Sacris literis quae virginitatem vouebant diligenter operam dabant, illudque exercitium, non autem cantillandi in templis habebant, quod magis erat tolerabile. Atque has superstitiosas pro iis Diaconis, de quib Paui. nic agit, r’apistae sustituunt: quae quales sint. et vt a verbo Dei alienae, atque etiam a veteribus canonibus et more Ecclesiae dissentaneae, omnes iam vident qui modo aliquid videre possunt. Sed pergamus ad Paulum.  \pend\pstart 14 Velim igitur iuniores nubere, liberos gignere, domum administrare, nullam occasionem dare aduersario ad maledicendum.  \pend\pstart Αʹντίθε Qς est vel potiùs Ε’πανοέρθοβις. Rationem enim subiicit, per quam superiori malo Eeclesia medori potest. In quo ipso canon vtilissimus, et quidem generalis, qui ad disciplinam Ecclesiasticam pertinet, continetur. est autem hic canon Iu\pend
\section*{CAPVT  V. }
\marginpar{[ p.297 ]}\pstart niores viduae Christianae nubunto, et liberos sus cipiunto, familiasque suas habento, quas curent, et administrent, vt omnia superiora mala euitent, et omnem nominis Dei blasphemandi occasionem cuilibet calumniatori et aduersario adimant. Singula vero huius canonis verba momentum et pondus aliquod habent. Volo, inquit Paulus. Primum videtur alienum ab Apostolica dignitate et doctrina praeceptum, mulieribus iam plus nimio nu pturientibus et nuptiales laetitias concupiscentibus praecipere vt nubant, Illud enim est oleum addere camino, non autem frigidam soffundere, vel lasciuas illas carnis libidines extinguere Resp. Duplici nomine hoc praeceptum subiiei a Paulo, 1.vt ecquod sit optimum vitae genus, quod viduae iuniores eligere debeant, vt sint aedificationi, doceat. Nec em ad nuptiales laetitias eas inuitat Pau lus, sed docet quae sit ratio omnium superiorum of fendiculorum tollendorum. 2. vt non minus Deo gratas esse, quae nubunt, quam quae viduae manent contra superiores errores, de quibus. c.4. sup. dixit demonstret. Nam ex hoc ipso apparet secundas nu ptias in mulieribus a Paulo minime damnari, imo vero et iuberi et probari, nedum vt viris sint interdictae. Q as tamen Papistae ex Encratitarum deliriis et haeresi damnant.  \pend\pstart Quaesitum est autem ad quas viduas hoc praeceptum pertineat, vtrum ad eas tantùm, quae sunt pauperes, an ad omnes: Vtrum etiam ad eas tantum, de quibus ad Diaconatum eligendis actum est in ecclesia:an ad quaslibet. Resp. Ad om nes. Id quod ex ratione, quae subiecta est, apparet manifeste. Non minus enim diuites, quam pau\pend
\section*{AD I. PAVL. AD TIM. }
\marginpar{[ p.Ayo ]}\pstart peres esse possunt otiosae garrulae, et nugaces, itemeque propter perieurum earmis, atque nuptiarum offendiculum infidelibus hominibus praebere.  \pend\pstart Volo) Sic et sup.cap.2.vers.8. et 1. Corinth. 7. vers7. sed diuerso sensu. Haec enim vox modo vim praecepti habet, vti in cap.2.vers.8. Modo autem consilii tantùm cuiusdam et exhortationis, vti hoc loco. Itaque pro subiecta materia huius vocis, volo, significatio accipienda est. Nec tamen illam Papistarum distinctionem probamus inter consilia et praecepta, quasi consilia quiddam et pius et sanctius contineant: cum hoc ipsum consilium Pauli etiam ex eorum sententia minime contineat eam vitae rationem, quae ab illis sanctior cen setur, sed eam potius quae minus ex eorum sententia, sit et seuera et sancta. Suadet enim Nuptias, non coelibatum. Praeclare autem hoc loco Chrysosto. negat hic tradi praeceptum a Paulo absolute cui quisquam obstringatur, nisi vt superiora mala euitet. Alias enim nubere aut non nubere viduis liberum est, vti apparet ex 1.Corinth. neque  cuiusquam conscientia irretitur et laqueo constringitur hoc loco, si nullum est offendiculum. Atque etiam ex hoc ipso loco apparet secundas nu ptias, dequibus proculdubio Paulus agit cum de Vi duarum nuptiis hic disputet, minime esse vel dannatas, vel Deo exosas, quicquid vel Montanistae, vel post eos Papistae sentiant: e quibus illi secundas nuptias omnino improbant: hi vero eas tanquam infames impedimentum esse censent ad ecclesiastica munera subeunda. Nihil certe moror dictum vel Tertullia. vel Bernard.serm.59 cant. quorum ille fuit Montanista, iste monachali superstitione  \pend
\section*{CAPVT  V. }
\marginpar{[ p.299 ]}\pstart iam fascinatus. Quod si secundae nuptiae nondam nantur, nec tertiae Aug. liber  de fide ad Petr. cap.3. in fine. Bernard.ser. 66. cantic. eam quoque sententiam refutat quae vult honestiora esse inter vir gines connubia, quam quae iam nupserunt.  \pend\pstart Uiduas. Ex hoc loco apparet, quam sit lasciuiae dedita natura, et ingenium omnium hominum: sed imprimis mulierum, de quibus hic agit Paulus: quanquam tamen ab eo malo non sunt mares liberi et immunes. Nam et quae.1. Corint.7. tradit Paulus et hic, aeque ad vtrunque sexum pertinent, quod si offendiculum oritur, et nascitur.  \pend\pstart Iuniores sexaginta annis supra definiuit Paulus: non quod tamen omnes, quae ea aetate minores funt, nlubere vent. leu cas tantum, quae liberos suscipere, et familias administrare possunt. Nam si quae 5o annos nata est, vel ea sit corporis consti tutione, vt haec mala fugiat, illi nullam legem dat Paul.hoc loco.  \pend\pstart Nubere Tria praescribit, in quibus omne muliebre officium contineri videtur. Nimirum Nu bere, liberos procreare familiam administrare. Nu bere. Hoc pertinet, vti diximus, ad secundas nuptias: et si iam secundo nupserat, et vidua restet, ad tertias et quartas et quintas nuptias, qualis illa, de qua Hierony. ait, septies eam nupsisse. Erat Vercellis apud Insubres. Liberos procreare τεκνογονεῆν. Quid significet vide supra, cap:2. ver.15. Oἰκοδεσ ωοτεῖν. Est in verbis Graecis ομοιτελευτὸν, et orationis ornatus. Hoc autem ὀικοδεσποτεῖν est familiam administret, non vt marito suo praesit, et sit superior vxor, sed pro eoiure, quo debet tanquam mater familias caeteris quidem prae\pend
\section*{AD I. PAVL. AD TIM. }
\marginpar{[ p.300. ]}\pstart esse eam vult Paulus interim tamen agnoscere virum caput suum, illique subiici et parere. Nullan occasionem. Subiectio est, quae totius vitae Christianae scopum demonstrat. Nos enim aedificationi omnibus esse oportet Fidelibus et Infide Iibus. Insirmis et Validis. Itaque  hac ratione vsus est saepissime Paulus vti sup.3.ver.6.Tit.2.vers.8. sed et Petrus quoque 1.epist. cap.2. vers.12.15. et Christus ipse Math.5.in si.idem argumentum afferunt. Ἀν τικειμένον autem generaliter accipio pro calumniatore et aduersario vti.2.Thessa.2.ver.4. Isec emim Sata iple maiedieit per se de doctrina, sed homines profani, Infideles et calumniatores detrailunt d de nobis et de doctrina Christi, quam profitemur.  \pend
\phantomsection
\addcontentsline{toc}{subsection}{\textit{15 Iam enim nonnulla deflexerunt, sequubae Satanam. }}
\subsection*{\textit{15 Iam enim nonnulla deflexerunt, sequubae Satanam. }}\pstart Confirmatio est superioris sententiae ab exemplo. Debent enim aliena pericula reddere nos cautiores ac prudentiores. Nam felix, quem faciunt, aliena pericula cautum. Hoc ipsum pro exemplo et documento Papistis esse debuit, vt tyranni cas suas de coelibatu leges tollerent, et abrogarent, cùm inde tot mala nasci et emergere quotidie cernerent: sed mordicus doctrinam suam non modo retinent, verum etiam ipsi Dei legi impiissime praeferunt. Neque tamen sic semper con cludendum est. Illud omne nimirum esse mutandum, vnde aliquod periculum, et aliqua discessio ab ecclesia sit exorta. Sic enim omnis et correptio et disciplina ecclesiastica tolleretur, propter quam multi a fide defecerunt, et ab ec\pend
\section*{CAPVT  V. }
\marginpar{[ p.301 ]}\pstart clesia Christi recesserunt, vti Nouatiani et Dona tistae. Resp.igitur caussam ipsam secessionis spectandam esse. Nam si quis idcirco deflectit a fide, quod iugum Christi ferre nolit, discedat. Sin autem quod tyrannidis praeter Christi legem iugum sibi impositum ferre non potest et recusat, certe illud iugum tollendum est, et demendum. Neque enim quisquam iugum Christi durum sentiet, nisi qui omnino sceleratus et impius esse volet. Ergo quaedam nimio isto rigore coelibatus sam offensae discesserant atque deflexerant a fide, et ad Satanam reuersae erant, abiurata nimirum Euangelii doctrina.  \pend
\phantomsection
\addcontentsline{toc}{subsection}{\textit{10 Quod si quis fidelis aut si qua fidelis habet viduas, eis subuenito, et non onerator Ecclesia, vt iis qua vere viduae sant subueniat. }}
\subsection*{\textit{10 Quod si quis fidelis aut si qua fidelis habet viduas, eis subuenito, et non onerator Ecclesia, vt iis qua vere viduae sant subueniat. }}\pstart Conclusio est totius superioris disputationis de viduis pauperibus, quae quidem sunt fideles. Continee aucein mne loeuo eclai cauonem illi primo si millimum, qui sup. vers.4. explicatus est. Est autem hic canon sextus numero, Quisque fidelis suas viduas alito: Illis, quae a cognatis ali et exhiberi possunt, ecclesia ne grauator, vt aliis omni auxilio destitutis subuenire et suppeditare possit quod satis est.  \pend\pstart Primum quoad verba, Erasmus censet haec a σςὴ a studioso aliquo qui mentem Pauli plenius explicare glossemate voluit, fuisse addita Vult igitur hunc locum sic restituendum esse ὐτις πσὸς ἰχει χόρας. Sed perinde est. Nam etiam  \pend
\section*{AD I. PAVL. AD TIM. }
\marginpar{[ p.302 ]}\pstart voxista πιςις non addatur, tamen ex mente Pauli intelligitur. Nihil enim hic de viris praeceptum est, qui propinquas et cognatas viduas alere possunt, quod idem ad foeminas non pertineat. Nam etiam illae si sunt diuites, pauperes consanguineos suis opibus sustentare et subleuare certe debent, exemplo illius Tabithae de qua est facta mentio Actor.9. et Electae, de qua Ioannes 2.epistola, quae vtraque suis opibus alienos fouit, nedum cognatos fideles egentes, si quos habuit, Item et Ioannae mulieris, et aliarum de quibus est Luc.8. vers. 3. Quaesitum est autem. Quid si illae viduae quae se a piis propinquis ali postulant, ipsae non sunt piae. Resp. Etsi omnibus est benefaciendum, Vt dotet iuem l'auius Galat. o. imprimis tamen de domesticis fidei hoc Pauli praeceptum intelligendum esse, et de Viduis piis, quae nomen dederunt Christo, quarum maxima et potissima ratio sem per haberi debet a cognatis piis et Christianis.  \pend\pstart Dicit autem Paulus Επαρκεῖτο id est, abunde et copiose sumministret. Αρκεῖν enim est copiose, et quantum satis est atque necessarium, suppeditare. In quo tamen praeclare Chrysosto. annotauit, non praecipi a Paulo, vt deliciae illis Viduis vel a cognatis vel ab ecclesia, a qua aluntur, praebeantur: sed quod ad victum modestum, sobrium ac earum conditioni aptum sit satis et conueniat. Neque tamen sic, vt Viduae illae alantur otiosae et inertes: sed vt illae ipsae victum sibi, quantum arte vel manibus suis vel labore possunt, comparent etiam, veluri lanificio, serificio, suendo, nendo, aliisque rationibus honestis. Quod si istae viduae et pauperes laborent, iam cognati vel ec\pend
\section*{CAPVT  V. }
\marginpar{[ p.303 ]}\pstart clesia addat, si quid illis ad necessarium victum deest. Idque pro ipsorum praebentium facultatibus. Neque enim vel cognati, vel ecclesia cogitur omnia quae viduis et aliis pauperibus desunt, suppeditare, nisi id fieri ipsorum facultates et opes patiantur. Vera est enim etiam inter Christianos illa Ennii sententia a M. Tul. Cicerone recte explicata, ita aliis largiendum esse, vt nostrae familiae etiam subuenire possimus. Pro eo enim quod quisque habet dare debet, et in eo Deo acceptus est: non autem pro eo, quod non habet. Exhaurire enim seipsum prorsus nemo tenetur, ne sit huiusmodi subuentio ipsi subuenienti oppressio, et ne postea ipse mendicet, vt docet Paulus.2. Corinth.8.vers.12.13. Etsi parce tamen grauate et tenuiter non est subueniendum proximo sed luben ter et vti dixi, quantum possumus.  \pend\pstart Vtiis. Æquissima ratio quae ducta est a fine aerarii ecclesiastici. Ecclesiam enim ad suos pauperes alendos pecunias a piis collatas habuisse apparet ex Sozom. liber  7. cap.26. et Euang. libus  5. capite. 5. ne praeter aequum se onerari existiment propter hoc praeceptum istarum viduarum cognati. Etsi enim ecclesiae aera ium est commune omnium pauperum aerarium tamen videndum est et prospiciendum, ne statim exhauriatur: omnibus emm sufficere diu non pote rit Reseruetur igitur illud maiori inopiae, et eorum praesertim egestati, qui non habent, vnde se aliunde exhibeant, siue sint Viduae: siue Mares pau peres. Qui habent vnde aliunde adiuuentur, eccle siam ne grauent. Haec sane de pauperibus viduis cura perstitit diu in ecclesia Christiana puriori\pend
\section*{AD I. PAVL. AD TIM. }
\marginpar{[ p.304 ]}\pstart bus seculis maximeque episcopi iubebantur illi prouidere et prospicere vt apparet ex variis Synodotum decretis et imprimis Synodo Chartag.4. postea pontificii παρεπισκοποι, qui nihil omnino, ni si suum lucrum ex ecclesia captant, totam hanc pauperum curam neglexerunt, minimeque ad se pertinere existimarunt: sed ipsi dona Viduis, pupillis, et aliis ecclesiae Dei pauperibus destinata sacrilege inuaserunt. Quorum certe bonorum sunt Deo rationem reddituri: at que ii omnes etiam, qui in reformatis ecclesiis, bona, quae ad pauperes pertinent, hodie alio transferunt, et sibi propriisque commodis applicant. Quod faciunt. nimium multi, imprimis aut Nobiles et Principes Christia. Quid autem de bonis ecclesiasticis sit inter nos constituendum, vide Cal. epist. 372. 394.384.  \pend
\phantomsection
\addcontentsline{toc}{subsection}{\textit{17 Qui bene praesunt Presbyteri, duplici honore digni habentor: maxime qui laborant in sermone et doctrina. }}
\subsection*{\textit{17 Qui bene praesunt Presbyteri, duplici honore digni habentor: maxime qui laborant in sermone et doctrina. }}\pstart Πόεσβασς Transitio est. Tractat enim iam aliud caput disciplinae et politiae ecclesiasticae Paulus, quod ad Ministerii ecclesiastici munus eque pertinet, atque superior cura et disputatio. Est aucenrnottput de Presbyterorum, et eorum, qui praesunt ecclesiae munere id est, tum aliorum erga ipsos: tum eorum ipsorum erga seipsos: tum eorundem etiam erga alios officio et cura in vniuersum.  \pend\pstart Ac quidem de aliorum erga ipsos, et ipsorum cura erga seipsos, quid nimirum praestare inter se debeant, ostendit Paulus vsque ad vers.24. variós  \pend
\section*{CAPVT Ve, }\pstart que canones constituit omnino et semper ecclesiae necenarios;i inouo cam recte administrare ipsi velint. 101100  \pend\pstart Cur autem haec praescribat Paulus facilis ratio est, Nam si vel minimi cuiusque ratio in Dei ecclesia haberi debet, multo magis eorum, qui praesunt ecclesiae. Si singuli, qui sunt in ecclesia ordines, regendi sunt certis legibus, vt omnia recte fiant in Dei ecclesia, multo magisii, a quibus maius periculum iminuumgo mlla quaedam hic propeules.  \pend\pstart I Quid praestandum Presbyteris.  \pend\pstart 2 Quibus Presbyteris.  \pend\pstart 3 Cur praestari illis aliquid debeat.  \pend\pstart Presbyterorum autem nomen hoc loco, et in toto noc argumento et capite, dignitatis nomen est, non aetatis, quicquid Chrysosto. et alii contra sentiant vti saepe alibi et I.Petr.5.vers.1.accipitur. Id autem esse ita accipiendum apparet ex his ver bis, Qui laborant in sermone.  \pend\pstart Ergo illis, debetur honor, isque duplex, ait Paulus Honoris voce subsidium omne piumque officium, more Hebraeorum, significatur, vti supra eodem significato vers.3. Viduas honora et 1. Petr.3.vers.7. Vir vxori honorem tribuat. Sumpta est haec dicendi ratioex lege Dei et praecepto 5. Honora patrem. Sunt enim nomine parentum, sed spiritualium comprehensi, qui nos verbo Dei pas cunt: et praesunt nobis, nósque in Christo vel genuerunt: vel iam genitos coelesti cibo alunt, quod faciunt hi Presbyteri. I taque hic canon Pauli ex quinto legis diuinae praecepto ductus est. Consentit autem hic locus cum 1. Corinth.9. etHebr.13.  \pend
\section*{AD I. PAVL. AD TIM. }
\marginpar{[ p.306 ]}\pstart vers.7.1.Thess.5.vers.12.13.  \pend\pstart Cùm autem hoc praecipit tam diserte Paulus, oblique taxat illius aetatis ingratitudinem in suos Pastores, nec tantum cauet in futuram posteritatem. Hoc tamen malum regnat adhuc hodie, et hac ipsa ratione verbo coelesti priuare ecclesias tentat Satan: multos certe a suscipiendo coelesti ministerio deterret, quod eos efficiat famelicos, qui Deo ipsi in tam honesta necessariaque causa seruiant: quorum operae et labori nullus alius, vt ait Chrysostomus, eit tontereuus. Et recte. Ergo quod aliis collatum diceretur eleemosyna, Presbyteris praestitum dicitur honorarium, stipendium, Debitum, a Christo ipsoμιθὸς.  \pend\pstart Ait Paulus Duplicis, vt caeterorum comparatione, quibus aliquid etiam specialiter debemus, eos nobis magis commendet. Sic tamen quidam interpretantur, Duplici honore perinde et tantum esse, atque si dixisset, Magno. Alii vero comparate hoc dici putant, alii absolute. Absolute quod et in hac vita digni sint verbi Dei ministri stipendio et subsidio victus: et in futura vita, felicitate aeterna. Quicerte errant, si nos Ministrósve ipsos aIiquidapud Deum mereri posse putant. Chrysosto. rectius, qui Duplicem honorem explicat Veneracionem ipiotum personae id est reuerentiam: et Subsidium vitae ac alimoniae. Certe vtrumque Pa storibus, vti Parentibus ipsis, deberi verissimum est. Qui comparate vero hunc locum accipiunt, referunt ad superiora praecepta de viduis alendis et honorandis. Istis enim viduis longe sunt magis honorandi Presbyteri et pastores. Quae sententia videtur huic loco consentanea et com\pend
\section*{CAPVT  V. }
\marginpar{[ p.307 ]}\pstart moda. Non vult autem magnas opes et honores congeri in Ministros Paulus, etsi duplici honore dicit dignos. Recte Chrysosto. Ego ecclesiae prae sules audenter dixerim, ait, nihil praeter victum et vertieu habcre oporteresverim.cap.o.ver. 8. expli cabitur. Praeclare quoque Erasmus. Sacerdotes et praedicatores quidam nihil sunt, nisi decimatores.  \pend\pstart Obstat autem quod propriis manibus victum sibi quaesiuit ipse inseruiens ecclesiis: ergo ipsi θιωτικa sibi comparare debent. Actor.2o. ver.34 Respond. ex Chrysosto ne pastores inopia coacti a studiis et melioribus rebus distrahantur, iis victum esse ab ecclesia suppeditandum. Nam et victum sibi manibus artificióque quaere re, et suo muneri superesse simul non possunt Mi nistri, nisi dona illa prophetiae et cognitionis scri pturarum egregia haberent, extraordinaria qua lia certe habuit Paulus. Plusquam belluina ingrati tudo est denegare carnalia ei, qui tibi coelestia praebet vt loquitur Paul. Ro.15.v.27. Neque negat Paulus se ab vllis ecclesiis vnquam accepisse stipendium.2. Corinth.11.vers. 8.sed eo tm abstinuisse, cùm ecclesiae aedificatio id fieri postulabat. Eodem modo est respondendum, si quis obiiciat Ne he.5.ver.15 Res. Nehemiam id fecisse propter popu li paupertatem: non autem qu victum exigere ab eo non licuisset. Itaque  non est lex generalis illud Pauli vel Nehemiae factum, sed singulare quoddam pro temporis ratione assumptum. Et certe nimis sunt impudentes Catabaptistae, qui illius facti et exempli praetextu omnes verbi Dei Ministros ad manualia opera cogendos esse contendunt, vt victum comparent: aduersus quos praeclare Bulling liber 2.  \pend
\section*{AD I. PAVL. AD TIM. 308 }
\marginpar{[ p.]}\pstart Tract.7.contra Catabaptistas. Nec fures nec latrones sunt verbi Dei Ministri, qui iosta a grege suo stipendia, quantum quidem ad victum et vestitum satis est, exigunt, vt isti calumniantur. Nec sunt fuci cùm strenue verbo Dei nauantes operam, ita viuunt siue manuali labore.  \pend\pstart Secundum caput, est quibus Presbyteris debeatur? Resp. Iis presertim qui in verbo Dei et Doctrina laborant. Ergo est duplex Presbyterorum genus, Vnum eorum, qui Doctrinae coelestis sont Ministri: qui nihil ab episcopis et Pasto ribus differunt: Alterum eorum, qui moribus fidelium tantum inuigilant, et eos inspiciunt, ne quid offendiculi in ecclesia Dei ex quoquam oriatuc. Hi sunt, quos in Gallia vulgo Supervigilantes dicimus qui vna cum Ministris verbiDei con sistorium siue Senatum ecclesiasticum constituunt. Sed quaeritur, vtrum hoc postremum genus Presbyterorum sit negligendum ab ecclesia. Respon. Minime vero, siquidem isti quoque egeant: sed primum genus magis commendatum nobis esse de bet Ex quo fit, vt Paulus addiderit. Ma xime. simi lis locus Galat. 6.vers.6. Ait enim communicet qui instituitur in sermone ei, qui se irstituit, omnia bona. Obstat quod ait Christus Matth.1o. vers. 8.Gra vis ac cepistis, gratis date. Resp. Eum gratis praeesse et pascere gregem, qui nec lucrandi nec ditescen di animo id facit: quanquam victum sibi necessarium a grege, idest ab ecclesia accipit, quemadmodum docet et Christus eodem loco, et Paulus 1.Corinth.9.2. Corinth.11.vers.8.  \pend\pstart Vult autem istos Presbyteros laborare, ergo non otiosos esse vult, non etiam perfunctorie  \pend
\section*{CAPVT  V. }
\marginpar{[ p.309 ]}\pstart tantùm fungi suo munere: non denique pro sua tantùm commoditate inseruire ecclesiae: sed magno animi zelo, totisque et animi et corporis viri bus incumbere, vt optimam muneris sui rationem Deo reddere possint. adeo vt saepe de se testetur Paulus se non tin laborasse, sed etiam lassitudinem, quasi exhaustis viribus, passum esse, tantae aerumnae corporis in eo faciendo nimirum saepe sunt (vt 2. Thes.3.ver.8. docet ipse) perferendae, ne sint muti canes, vt ait Isaias. Sed et laborare in sermone et Doctrina, non qualibet, sed coelesti et Euangelica iidem iubentur, per quam homines adveramn agnitionem Dei deducant. Quidam sermonem pro exhortationibus accipiunt. Doctrinam pro ipsa traditione methodica coelestis mysterii Quod vtrunque fit concionando. Optime Chrysosto. qui docet exigi a pastore his verbis, vt Predicet, doceat, et concionetur: non tantum vt studeat et vr operam det priuatim Dei verbo, illudque n ditetur et speculetur ipse, vel vt vite sanct? „ plo gregem suum instituat. Docere igitur det pastor, et eatenus officium facit, quatenus viua vo ce illud exponit. Ite, ait Christus, docete. Ergo qui absunt a suo grege: qui in grege ipso muti sunt: qui etsi per alios docent, ipsi nunquam id faciunt, omnino sunt indigni nomine Pastorum et Presbyterorum: et duplici illo honore, de quo agit hoc loco Paulus, priuandi  \pend
\phantomsection
\addcontentsline{toc}{subsection}{\textit{18 Dicit enim Scriptura, Boui trituranti non obligabis os. Et, Dignus est operarius mercede sua. }}
\subsection*{\textit{18 Dicit enim Scriptura, Boui trituranti non obligabis os. Et, Dignus est operarius mercede sua. }}\pstart Αιτιολovα est, in qua tertium caput eontinetue  \pend
\section*{AD I. PAVL. AD TIM. }
\marginpar{[ p.310 ]}\pstart nempe. Cur hic honor Presbyteris debeatur. Affertur autem a Paulo ratio duplex, quarum Prima ducitur ab exemplo, quod etiam diuina lege comprobatur. Secunda ab officio et communi hominum sententia. Ergo videtur hic vti duplici autoritate Paulus, nimirum Diuina et Humana: Dlustia, cum ast. Diit entin seriprura, Boui trituran ti etc. qui locus sumptus est ex Deut.25.ver.4.et1. Corint.9. ver. 9. explicatur fusius ab eodem Paulo.  \pend\pstart Humana vero autoritas est, cùm ait, Dignus operarius mercede sua: quae in prouerbium abiit, ipsa naturae voce inter omnes mortales confirmata: adeo quidem vt cùm Christus ipse hoc idem dixit Math.1o.vers.1o. Luc.1o.vers.7. ex communi prouerbio sumpsisse videatur. Nihil enim hoc di cto magis tritum est, et vulgare. Nisi quis fortasse ductum malit ex eo, quod saepe in scriptura legitur, esse dignum mercenarium mercede sua Deut. 24.vers.15. Genes.31. vers.12. Ad quod ipsum quoque respexit Iacobus cap.5.vers.4.  \pend\pstart Neque tamen, etsi dignus est Pastor suo stipendio et suo honorario, tanquam merce de, propterea pascere gregem debet eo animo, vt inde stipendia accipiat, et mereatur. Nam huiusmodi animus est Auaritiae (quae est omnibus fugienda, imprimis autem Pastoribus Luc.12.vers.15.1. Petr.1. vers.2.)surculus, quanuis non expetat inde ditescere, vel opes cumulare. Praecipuus enim Pastorum scopus esse debec Dei ipsius gloria. Victus autem ipsorum est accessio quaedam tantùm, quem idcirco petunt et spectant, quia ad faciendum et melius, et alacrius, et liberiùs officium, ille ne ceslarius est.  \pend
\section*{CAPVT V. }
\marginpar{[ p.31 ]}\pstart Mirum autem est praetermissam hic esse confir mationem huius sententiae, quae fieri potuit exemplo Leuitarum, qui a reliquo populo Israelitico ale bantur Dei iussu ex decimis, et aliis oblationibus: sed ab eodem Paulo in 1.Corinth.9. vers 13. haec ratio non est neglecta.  \pend\pstart Vnde errant, qui stipendia Ministris tribui putant tanquam meras et gratuitas tantum eleemo synas putantque  liberum esse ea persoluere vel non: debentur enim vti iustissima merces, etsi nulla est tanto eorum labori aequalis, quemadmodum ait Chrysosto. Scholastici et Canonistae in cano. Qui pro pecunia 1.q.1.cano. Altare.1. q.3. distinguunt inter illud quod accipitur tanquam Pretium ope rae et Sustentatio vitae, Illud quidem accipi a Pastore hejdiit ponemue cutem concedun nec est sane inutilis vel incommoda.ista distinctio. Vide Durand.in liber 4. Sent. dist.25.quest.3. et 4.  \pend
\phantomsection
\addcontentsline{toc}{subsection}{\textit{19 Aduersus presbyterum accusationem ne admittito, nisi sub duobus aut tribus testibus: }}
\subsection*{\textit{19 Aduersus presbyterum accusationem ne admittito, nisi sub duobus aut tribus testibus: }}\pstart Transito est: iam enim agit desecunda huius ar gumenti parte nimirum ecquod sit ipsorum Pres byterorum aduersus seipsos et collegas suos offr cium, ne quid ex ea parte offendiculi nascatur. Imprimis enim cauendum est, ne illi se miutuo excusent et sibi parcant, dum tamen alios castigant seuere: atque ne in illis locum habeat prouerbiumvetus, Mu tuo muli scabunt. Potius illis in mentem veniat illa Christi sententia, vos estis sal terrae si sal euanueritin quo salietur? vos estis lux mundi. Math.5.verf. 13.14. Item illa, sint lumbi vestri praecincti et lucernae ardentes in manibus vestris. Luc.12.ver.35. Sed et illa quoqs  \pend
\section*{A D I. PAVL. AD TIM. }
\marginpar{[ p.312 ]}\pstart Petri debere pastores toti gregi esse exemplar vi tae sanitate doctrinaeque puritate, caeterisque  praelu cere.1.Petr.5. vers.3. Ergo quemadmodum vna ex parte diligenter illis cauendum, est ne propriis collegarumve peccatis conniueant: sic ex altera parte prospiciendum, ne ad facilem criminationem et accusationem aditum aliquem proteruis hominibus et iugum disciplinae ecelesiasticae renuentibus et excutientibus praebeant. Falsa est qui dem illa Papistarum ratio, qui postquam vniuersum Domini gregem distribuerunt in Laicos et Clericos, fingunt laicos esse infestos acperpetuos Clericorum hostes, eósque innato quodam ex clericalis dignitatis fulgore et priuilegiis odio prosequi. Itaque cano. Nullus laicus, et cano. Laicos. 2.quaest.7.omnes laicos, quamuis probatae vitae vi ros pontificii ab accusatione et testimonio in clericos dicendo summouent. Sed peruerse. Primum enim contra Christia. charitatis praecepta illud odium in omnibus, qui subiiciuntur, inesse praesumunt, cùm non sit charitas suspiciosa. Deinde res ipsa docet, qui Christiani aliis subsunt, eos propter Dei verbum ex animo pastores, praepositósque suos diligere, illisque subiici lubenter velle, quia non sunt pastores euangelici animorum tyranni. Quid igitur? Possunt fateor in ecclesia Dei esse efftenes homines: possunt esse simplices: possunt esse hypocritae et astuti, qui cùm doctrinae ipsius firmitatem concutere non possunt: neque aliter mysteria arrodere: in personas ministrorum debacchantur impudenter, vt eorum autoritatem eleuent et imminuant quantum possunt.  \pend\pstart Atque ea semper arte Satan tacite tanquam  \pend
\section*{CAPVT  V. }
\marginpar{[ p.333 ]}\pstart per cuniculos Ecclesiam Dei conatus est euertere. Itaque huic et Satanae technae, et improbo rum audaciae fuit occurrendum, praecludendusque  huiusmodi criminationibus aditus. Imprimis quod multorum accusationi et calumniis obnoxius esse solet, qui in omnium oculis versatur, qui in vitas omnium inquirit, qui omnium vitia reprehendit, tam sunt nimirum homines correctionis impatientes, vt recriminatione etiam falsissima sese vlcisci studeant aduersus eos, a quibus tamem merito arguuntur. Est autem hoc pastorum officium, vt est Ezechiel. 18. tanquam hominis in specula collocati qui despicere omnia, ac turpia quaeque  redarguere debet. Ex quo fit, vt sint multis odiosi, et malignorum obtrectationibus pateant, sintque obnoxii. Id quod in causa. Bonifacii. docet Augustinus. epist. 137. Vnde huic quoque periculo cauendum et prospiciendum fuit Vtrunque  autem hoc malum, de quo diximus, excludit et summouer canon, quem sot loco sancit Paulus, aut potius Dei ipsius spiritus. Est autem huiusmodi, Presbyteri ne aduersus Presbyterum facile vel accusationem ipsam admittunto, nisi duorum ver trium testimonio crissien cùm defertur, sit statim confirmatum.  \pend\pstart Qui canon vt commode declaretur, sunt haec tria ordine nobis perquirenda.  \pend\pstart Primum caput facile est, et ex superioribus versiculis satis apparet. Agit enim de iis, qui dignitatis, non autem aetatis ratione Presbyteri sunt.  \pend
\section*{AD I. PAVL. AD TIM. }
\marginpar{[ p.34. ]}\pstart quorum officium et munus in fine huius capitis persequitur et exponit sigillatim. Quanquam aliud sentit Chrysostom. et ad aetatem refert, cui, tanquam minus ad peccandum procliui, vult parci, minusque  facile aduersus eam vel solam ipsam acculationen adinitti oed naee ratio Chrysost. saepe quidem apparet falsissima cùm senes quosdam videamus ad omne scelus propensissimos.  \pend\pstart Secundum igitur caput est. Quid de Presbyteris hoc loco Paulus praecipiat. Resp. Ne contra eos vel sola duntaxat accusatio, citra duorum vel trium testimonium, excipiatur vel a toto collegio Presbyterorum, vel ab vno aliquo Presbytero, nedum vt condemnatio de iis temere a quoquam fiat. Quaerunt autem, quid hoc nouo praecepto fuerit opus, et in Presbyteris praesertim, de quibus hic aliquid speciali' sancire Paülus videtur, cùm in omni accusatione hoc ipsum locum habere videatur. Nam Deut. 17.vers.6. sic ait lex Ex ore duorum vel trium stabit et ratum erit verbum et accusatio, quae fiet de aliquo. Itaque hoc ipsum tam leuis momenti quibusdam hoc loco visum est, vt omnino delendum censerent, tanquam inutile et superuacaneum, cùm tamen singulare priuilegium, et a Deo ipso pastoribus et Presbyteris concessum contineat. Respon. igitur aliud esse Accusationem ipsam, aliud Iudicium, quod de accusato fit, continetque vel absolutionem vel condemnationem rei a iudice factam. Iudicatio igitur siue sententia sine testibus et legitimis probationibus in vniuersum a nemine fieri debet. Ac eo pertinet locus Deuter. 17. Accusatio autem ipsa certe admitti sine testibus a Iudice  \pend
\section*{CAPVT V. }
\marginpar{[ p.315 ]}\pstart de quouis alio potest, praeterquam in hoc casu verbo Dei expresso: si modo, qui accusant non sunt omnino infames personae: vt est praeclare responsum ab Eusebio Vercellensi Episcopo, et extat in canone De accusatoribus3 quaest.5. Presbyteri igitur admitti a collegis sine testibus, iisque duobus vel tribus, ne accusatio quidem ipsa potest. In quo est Presbyterorum praerogatiua praeter caeteros homines hic a Paulo expressa, quia variis hominum calumniis, multorumque inuidiae sunt obnoxii.  \pend\pstart Testes autem eos hic requirit Paulus, qui in caeteros reos audiri et admitti possunt, non alios. Ergo huius canonis et priuilegii ratio (in quo tertia hui' versiculi pars posita erat) apparet esse, non quod presbyterorum peccatis, et vitiis parcere et conniuere suadeat Paulus (quos grauius propter dignitatem peccare certum est) non quod priuilegia fori et immunitatis personis Ecclesiasticis tribui velit, qui Romanor. 13. omnes atque etiam pastores ipsos magistratibus subditos esse iubet, vti et Petrus I.Pet.2.vers.17. Non quod semina odii et suspicionum inter gregem et pastorem serere, aut alere cupiat: sed quod rerum iplarum experientia doceat quosdam odisse principes et eos omnes qui sibi praesunt: atque etiam tantam esse humanae mentis malignitatem, vt eos qui vitia nostra reprehendunt, maxime in odium ac inuidiam aliorum adducere conemur. Cui malo et effraeni licentiae occursum esse sapienti isto remedio voluit Paulus.  \pend\pstart Illud autem priuilegium, vt Episcopi nonnisi a septuaginta testibus, iisque non Laicis homini.  \pend
\section*{AD I. PAVL. AD TIM. }
\marginpar{[ p.310 ]}\pstart bus conuicti damnari possint prorsus impunitatem peccandi dat Episcopis. Quod tamen canon: Placuit 2.quaest.4. continetur. Item illa clericorum praeuilegia, per quae vetantur in vniuersum Laici clericos accusare, et inter clericos ipsos ii, qui inferioris sunt gradus, prohibentur superiores deferre, aut in eos testimonium dicere, sunt omnino tollenda et abroganda, tanquam nefandissima omnium scelerum in Episcopis vela. Atque illa omnia late 2.quaest.7. per totum explicantur.  \pend\pstart Fenestram enim ad nequitiam patefaciunt, et nullum Episcoporum corrigendorum locum relinquunt, planeque pugnant cum mente Pauli.  \pend
\phantomsection
\addcontentsline{toc}{subsection}{\textit{20 Eos qui peccant, coram omnibus argue, ut et reliqui timeant. }}
\subsection*{\textit{20 Eos qui peccant, coram omnibus argue, ut et reliqui timeant. }}\pstart Alter canon de poena Presbyterorum sic conuictorum. Est autem hic, vt Presbyteri, qui publice et notorie peccarunt, publice, id est, in coetu totius Ecclesiae arguantur.  \pend\pstart Ratio autem est, qu huiusm odi crimina et facta, quae pluribus testibus probantur, etiam antequam delata sunt, pluribus sunt nota, praesertim vero si a Presuyteris coma fuerit. Vere enim ille Nam lux altissima fati, Occultum nihil esse sinit. Cùm igitur crimina illa pluribus iam innotuerunt, et in ore hominum et cognitione versantur, sunt iam publica. Ex quo fit, vt etiam sint publica cenfura notanda, et publice arguenda, quemadmodum hoc loco praecipit Paulus.  \pend\pstart Deinde cùm Presbyteri caeteris exemplo esse debeant vitae sanctitate, ita et poenitentia, si deliquerint: et censura atque castigationc, si non resi\pend
\section*{CAPVT  V. }
\marginpar{[ p.317 ]}\pstart puerint, vt sint de populo reliqui ad iniuriam, et peccandum tardiores. Sed etiam olim reliquis Ecclesiis significabantur ii, qui ob scelus erant depositi, vtex concilio Carthaginensi 1. et 2. apparet, et ex Cypria. Atque has literas formatas appellabant. Certe Presbyter Flauianus Episcopum suum publice reprehendit, quod in Christiana religione veram fidem simularet, vt ait Theod.libus 5.cap.3. et idem liber  5.cap.28. Laudat Chrysostom. quod suos Presbyteros libere, tum moneret, tum publicis censuris castigaret. Vt autem male agentes Presbyteri sunt semper reprehendendi: ita bene agentes sunt et cohortandi et laudandi, vti Episcopi Ægyptii Athanasium exulem propter fidem aliis Ecclesiis commendarunt, vt apparet ex Tripar. histor. liber  4.cap.19.  \pend\pstart Plato annotat tria esse supplicii, per quod crimen punitur, genera κόλασην, τμιν, οράδεγμα. κόλασις est in ipsius delinquentis corpore et sensu inflicta poena, vti ictus fustium, abscissio auriculae. Tiuὴ poena ea, quae ob publicam tranquillitatem violatam imposita est, qualis est infamia, exilium, multa pecuniaria. Neutra castigatio pertinet ad Ecclesiam: sed ad magistratum quanquam ab Eugenio primo constitutum est, vt liceret Episcopio luos Presbyteros criminis conuictos carcere multare, qui certe fuit manifestus potestatis abusus. Παράδ'ειγμα est poena, quae caeterorum erudiendorum vel terrendorum causa publice constituta est, atque haec ab Ecclesia imponi potest. Verbum inde est οραδειγεόζεθς, quo vtitur Mat.1  \pend\pstart 1 Quod est Matth.18. vers15. Mitius repre\pend
\section*{AD I. PAVL. AD TIM. }
\marginpar{[ p.JI0 ]}\pstart hendendos esse qui peccant, quam hic Paulus iubeat, nimirum inter nos et illos solum. Res.Illud Chri sti dictum de priuatis offensis intelligendum esse, nimirum cùm nos priuatim laesi sumus, necdum offensio publice innotuit. Hoc autem delictum, quod publice docet arguendum Paulus, est publicum. Nec enim nisi propter publica aut publice nota crimina vel peccata reprehendi in coetu Ecclesiae quisquam debet, siue is sit pastor et Presbyter, siue priuatus. Atque ita hunc locum explieat Aoguitmus, nempe de peccatis Presbyterorum, quae non latent, sed publica sunt. De Corrept. et Grat. cap. 16. cuius sententiae plane assentior.  \pend\pstart 2 Obiicitur, peccata quae occulta esse oportuit, nota fieri hac ratione. Id quod et ab officio hominis Christiani alienum est, et pugnat cum aedificatione Ecclesiae, ac illa veneratione, quae deferenda est praefectis nostris. Respond. Ea tantùm argui publice ex hoc Pauli praecepto posse, quae publice nota sunt: non autem quae occulta sunt, quaeque celari, et ad caeterorum aedificationem tegi oportet, vt in exemplo Petri a se publice in Ecclesia reprehensi apparet Galat. 2. vers.14. In summa omnino disciplina aliqua Ecclesiastica opus est in Ecclesia Dei, si modo hominum effraenem licentiam coerceri, et morum sanctorum neruum aliquem, atque vinculum inter nos retineri volumus. Id quod, etsi verissimum est, Doctoris tamen eximii, et primarii D. Ro. Gualteri verbis potius, quam meis intelligi malo, qui in praefatione doctissimorum illorum suorum in Galatas commentariorum ita scribit. Quare, ait,  \pend
\section*{CAPVT  Vs }
\marginpar{[ p.319 ]}\pstart Ecclesiis disciplina opus, est sine qua non magis consiStere possunt, quam res domestica, quae patrefamil. et moderatore idoneo destituitur: non magis item, quam Respubus  quae nullis legibus gubernatur, aut nullos habet legum vindices. Etenim non minus vere quam prudenter tum puiuamuracimnt, matu esse sub principe viuere, sub quo nihil liceat: sed multo peiorem esse eorum sortem, qui principem habent, sub quo licent omnia. At qui huic loco commodùm adiicere libet quod idem Doctor praestantissimus et pientissimus in iisdem commentariis ad cap.2. Homil. 14 his verbis disertissime, sanctissimeque scribit. Admonemur hoc exemplo (de Petro a Paulo publice reprehenso agens) ait nullam vel in Ecclesia, vel in Republ. dignitatem esse tantam, quae censurae et reprehensioni eos eximat, qui aliquid praeter officium faciunt. Etenim cum omnium Dominus sit Deus, eius quoque verbo omnes homines cuiuscunque  sint loci, aut crcunis,uυiei, et voeuire vonnent, et ad omnes veroi Detpracones pertinet, quod olim Hieremie dicebat Deus, Ecce ego te facio hodie ciuitatem munitam. columnam ferream, et murum aeneum contra omnem terram, contra reges Iuda, contra principes eius, contra sacerdotes, contra vniuersum populum terrae.  \pend\pstart Haec Gualterus mihi obseruandissimus.  \pend\pstart Ex hoc autem loco colligitur, iam tempore Pauli in Dei Ecclesia duplex censurae Ecclesiasticae genus fuisse vsitatum, vnum Priuatum, alterum Publicum. Priuatum est, cùm qui offendit, agnoscit culpam eamque fatetur ei, quem offendit, nullo confessionis teste adhibito, sed inter se et eum, quem offendit, tantùm et priuatim id facit. Hanc descripsit Christus Matth. 18. descripsit  \pend
\section*{AD I. PAVL. AD TIM. }
\marginpar{[ p.320 ]}\pstart quoque Iacobus cap.5.vers.16.quae saepe a Pastore fieri praescribitur ad mutuam reconciliationem: et non praecepta, tamen ex officio a quoque nostrum fieri debet, vt bonum conscientiae testimonium habeamus.  \pend\pstart Publicum censurae genus est, cùm agnitio offensionis non fit fratri soli offenso: sed adhibitis aliis velut illius confessionis testibus. Haec etiam a Christo commemorata est Matth.18.vers.16. Hie autem Bernard. sermo.44. consilium locum habet, qui oleo suauium admonitionum mordacia medicamenta, et vinum compunctionis saepe docet esse adhibendum, duris videlicet et praefractis. Denique si duritia tanta est, etiam censurae Ecclesiasticae baculo percellendum contenptorem docet, qui priuatim reconciliari cum fratre nolit. Huius autem postremae censurae multae sunt species et differentiae, quae ex loci et multitudinis, in qua fit, ratione sumuntur. Eo enim magis publica est haec confessio et censura, quo maior est coetus, et frequentior celebriórque locus, in quo fit, atque quo plures illius testes adhibentur. Ergo sic a nobis videtur diuidi posse vt Alia dicatur Consistorialis: alia Ecclessastica censura. Consistorialis est, quae coram solo senatu Ecclehastico fit ab eo, qui peccauit. Coetum enim Presbyterorum senatum Ecclesiae vocat Hieronym. Ecclesiastica, quae fit praesente et aduocata vniuersa illius loci et vrbis Ecclesia, et non tantùm senatu Ecclesiastico praesente, id est, eo tempore et loco, quo coetus Ecclesiae conuenit, et frequens est.  \pend\pstart Publica vero nunquam, nisi pro publico deli\pend
\section*{CAPVT . V. }
\marginpar{[ p.524 ]}\pstart cto, indici debet. Publicum autem delictum duobus modis dicitur, vel quod Pluribus, vel quod Omnibus notum est: idque, vel ex ipsius facti perpetrandi ratione vt quod in publico loco, vel coram magno coetu hominum committitur: ex et sparso postea illius facti rumore, et pluribus communicato, vti multa peccata priuatim et clam primum admissa postea ex fama innotescunt, et diuulgam, tur. Ergo quae delicta omnibus innotuerunt, ea coram omnibus agnosciep huprouum est, et de iis publicam in totius Ecclesiae coetu exomologesin fieri oportet, si verum remedium morbo adhibere velimus.  \pend\pstart Quae pluribus tantùm et non toti Ecclesiae nec maiori parti innotuerunt peccata, illa in senatu tantùm Ecclesiastico:non autem coram vniuersa Ecclesia confiteri oportet, ne latiùs illud patefiat, quod occultum esse debuit.  \pend\pstart Illa vero delicti nostri, quo offendebatur Ecclesia, confessio publice praesenteque tota Ecclesia facta, exomologesis speciali et antiquo nomine dicta est, quemadmodum ex Latino Cypriano, et Graecis scriptoribus apparet.  \pend\pstart Haec in vsu fuit etiam primis Ecclesiae vetustissimisque  temporibus, vti ex hoc loco satis probari quoque  potest: et ex Ecclesiastica historia pene tota Philippus imperator qui post Decium Romae regnauit, sic poenituit, atque  confessus est, vt testatur Eusebus liber 6.Histor. cap. 34. Ecebolius quidam, etiam ex consuetudine iam recepta, sic peccata sua agnouit Socrat. liber 3.cap.13. Theodosius imperator, qui Maior dictus est, sic suum delictum, quod erat publicum in Ecclesia Mediolanensi a\pend
\section*{322 AD I. PAVL. AD TIM. }\pstart gueuit, Theodorit.libus 5.Histor. cap.18.et Sozomlibus 7, cap.24. Atque hic locus est vtilissimus.  \pend\pstart Aliquot igitur quaestiones in hoc argumento sunt nobis explicandae, quae magnum in exercenda disciplina Ecclesiastica vlum habere solent. Ac primùm quaeritur.  \pend\pstart In quo peccati genere locum haec publica exomologesis habere debeat. Resp. Cùm differat iurisdictio ciuilis ab Ecclesiastica, et haec tantùm voluntaria sit, illa vero contentiosa: haec solius poenitentiae animique salutis testificandae gratia comparata sit: illa publicae tranquillitatis tuendae causa sit constituta: in eo tantùm locum habereposse siue publicam siue priuatam delicti confessionem, qui agnoscit se deliquisse, idque sponte fateri paratus est. Nam in eo qui negat a se peccatum esse (quia vel factum ipsum a se commissum pernegat, vel hoc suum factum negat esse peccatum, et damnandum) nulla certe vel publica vel priuata exomologesis proprie locum habere potest, etsi hic testibus ipsis fuerit in senatu Ecclesiae conuictus fecisse Id, quod tamen pernegat. Itaque huiusmodi pertinax est ad magistratum remittendus, vt de facto ipso amplius inquirat. Sed tamen obstinatum hominem et peruicacei Eecielia honhunquam excommunicatione aut suspensione Coenae perstringere potest, id est si factum a se commissum esse agnoscens, tantùm qualitatem facti, id est, ἀνρμίαν et turpitudinem illius non vult talem videre, et idcirco non poenitere. Negantem autem id a se, de quo accusatur factum, nunquam Ecclesia damnabit, neque ad confessionem delicti, in cuius perpetrati dene\pend
\section*{CAPVT  V. }
\marginpar{[ p.323 ]}\pstart gatione persistit, coget, aut adiget: sed confitentem tantùm et vere poenitentem, inquam, Ecclesia ad confessionem adducet. Probo igitur id quod in can. Presbyter xv. quaest.I.in eandem sententiam scriptum est. Et haec quidem in vniuersum de iis, qui censura Ecclesiastica ad delicti confessionem compelli debent et possunt. Quod si de iis, qui publice confiteri debent, quaeritur, Respondeo eos tantùm ad id esse compellendos, iisque folis indicendum esse, quorum delictum et factum iam publicum et notorium est Ecclesiae.  \pend\pstart Quaesitum est, Quomodo fieri haec publica confessio debeat. Atque hoc partim ad formam externam spectat: partim etiam ad ipsius peccatoris publice confitentis officium. Ac quidem quod ad ipsum confitentem attinet, nemo per nuntium, neque per scriptum facere illahi confessionem olim potuit, sed corporali praesentia, et ipso ore suo, vt est praeclare responsum ab Augustino in libus  de Poenitent. et est relatum in cano. Quem poenitet distinct. I. de Poenitentia, et hoc ipsum retinendum censeo.  \pend\pstart Quod autem ad formam et ritus externos. Primùm haec publica confessio simplicissime in Ecclesia fieri solebat. Deinde vero huic simplicitaostentationem et pompam. Ergo concilio Agathensi constitutum est, vt publice confitens initio Quadragesi mae staret nudis pedibus pro foribus templi, sacco indutus, cilicio caput tegens, et supra caput cinerem spargens, vultu in terram demisso. Deinde exomologesi facta, Ecclesiae con\pend
\section*{AD I. PAVL. AD TIM. }
\marginpar{[ p.324 ]}\pstart ciliabatur per manuum impositionem a pastore siue Sacerdote. Quae res fuerunt pompae potius externae, quam vera castigatio delicti: sed quaesita ratio est nugarum in Ecclesiam pro vera poenitentia substituendarum, et affectata inepta quaedam seueritas. Ista referuntur in can. Poenitentes et cano. In capite quadragesimae distinct.50. quae prorsus a nostris Ecclesiis exulare debent. Quid enim signa haec poenitentiae praescribi necente cies cui qui vere poenitet, satis illa exhibeat:imo vero potius attollendus sit, et erigendus, quam deprimendus, ne incidat in desperationem?  \pend\pstart Quaeritur qui publicam exomologesin facere possint et permittantur. Respo. Omnes debere, et non tantùm posse, qui eo peccati genere deliquerunt, cuius poena debet esse consessio publica. Itaque nullius dignitatis, sexus et conditionis ratio habenda est, ne sit in Ecclesia Dei damnabilis illa personarum acceptio. Quod fieri vetat Dominus Malach.2.vers.9. Apparet etiam id ipsum ex hoc loco, vbi ne Presbyteris quidem ipsis parcit Paulus:item exemplo Theodosii imperatoris qui ab Ambrosio Mediolanensi Episcopo ad eam compulsus est. Quare perniciosissima est Canonici Iuris et Romanici decisio, atque sententia, quae negat clericos vllos quantumuis grauiter errantes ad hanc confessionem et exomologesin publicam damnandos ac compellendos, vel etiam si velint, admittendos, ne toti ordini clericorum fiat iniuria. Itaque  elericos detrudunt in Monasteria vbi poenitentiam agant solitarii, et nulli noti. Ita prae ceptum est in canone Alienum, et can. Con\pend
\section*{CAPVT  V. }\pstart firmandi distinct.5o. qui plane et apertissime cum hoc Pauli praecepto pugnant.  \pend\pstart Quaesitum est quoque, quoties in vita fieri publica haec exomologe sis possit. Respond. Augustinus in epistola ad Macedoniam semel tantùm in vita, ne vilescat Ecclesiae authoritas: aut ne fiant audaciores homines ad peccandum, si saepius in vita cuiquam liceat publice poenitere canou. Quanuis caute distinct. 50. Sed cùm nulli peccatori poenitenti sit spes veniae, et remissionis peccatorum deneganda: toties in vita fieri haec publica confessio potest, quoties quis eo modo deliquit, quo iuste publica confessio indici atque imponi debet. Et iste fuit certe nimius Patrum rigor, et pene Nouatianorum error, nimiaque austeritas, quae homines ad desperationem tandem adducebat, quasi denegata peccatorum remissionis impetrandae spe. Quoties igitur vocat ad se peccatorem Dominus, toties ille est ab Ecclesia admittendus!, etiam in die septuagesies septies, vti docet Christus ipse Matth. 18.nisi (vti diximus) velimus esse Nouatiani.  \pend\pstart Praeclare vero Bernard. quatuor esse docet, quae nos a publica hac exomologesi deterrent, quae sunt omnino nobis et cauenda et debellanda, si peccauerimus Sunt autem illa quatuor, Pudor, Timor, Spes, Desperatio. Itaque recte idem Ad milit. Templi cap.12. Et quidem verbum in corde pectoris operatur salutiferam contritionem. Verbum vero in ore noxiam tollit confusionem, ne impediat necessariam confessionem. Ait enim scriptura. Est pudor adducens peccatum, et est pudor adducens g loriam. Bonus pudor est, quo peccasse, aut certe pec\pend
\section*{AD I. PAVL. AD TIM. }
\marginpar{[ p.326 ]}\pstart care, confunderis. Huiusmodi proculdubio pudor fugat opprobrium, parat gloriam, dum aut peccatum omnino non admittit: aut certe admissum et poenitendo punit, et confitendo expellit: si tamen gloria etiam nostra haec est, testimonium bonae conscientiae. Hic igitur pudor est secundum Deum, qui nos ad eum reuocat, facitque vt peccati nos pudeat: non autem resipiscentiae, et melioris frugis. Praeclare enim Augustinusin Ioan. tract. 12. Qui confitetur et accusat peccata sua, iam cum Deo facit. Cum autem coeperit tibi displicere quod fecisti, inde incipiunt bona opera tua, quia accusas mala opera tua. Denique haec censura est pars potestatis clauium, quae a Domino Deo data est Ecclesiae. Quam qui impugnat, optimum Ecclesiae doctrinaeque Christianae retinendae neruum et vinculum tollit, nimirum authoritatem illi a Deo ipso concessam, et ius tuendae conseruandaeque Deigiorlae at retorsanidae nomhum vitae in melius.  \pend\pstart Obiiciunt tamen quidam atque afferunt aliquot rationes, per quas istam disciplinae Ecclesiasticae partem abrogent, vel saltem in odium adducant si possunt.  \pend\pstart Primum autem eorum argumentum est. Quid prodest haec confessio, si iam peccator voce dominica et interna Spiritus sancti vi resipuit et resuscitatus est. Cui argumento respondet August. in sermone 8.de Verbus  Domini secundùm Matthae. Quid prodest Ecclesia et authoritas ipsius confitenti, cui Dominus dixit, Quae soluetis in terra, etc. Ipsum Lazarum attende, cum vinculis prodit, iam viuebat: sed nondum liber ambulabat  \pend
\section*{CAPVT . V. }
\marginpar{[ p.227 ]}\pstart vinculis irretitus. Quid ergo fecit Ecclesia et confessio, nisi quod ait Dominus continuo ad discipulos, Soluite illum, et sinite abire. Ergo non est satis ei qui publice peccauit, interni animi compunctione affici, nisi datum a se offendiculum tollat, et resarciat aperta culpae suae confessione, vt si qui prius illius exemplo peccare vellent vel didicissent, iam contrario eiusdem exemplo resipiscant, et auertantur. Fit enim haec confessio non ad peccatorum remissionem impetrandam a Deo: sed ad Ecclesiam aedificandam, quae scelere illius destrui videbatur.  \pend\pstart Obiiciunt irem, Haec confefsio publica affert infamiam, cùm iam peccatum negare non possit is qui publice agnouit. Respond. Minime vero, Imo et Ecclesiae, et ipsis adeo Angelis illa gaudium affert, cùm sit certissimum resipiscentiae animi testimonium, et aberrantis ouis receptio, atque inuentio. Malach.18.vers.1o.Luc.15.vers.1o. Neque  certe laudem dare Deo (quod fit hac publica confessione apertissime, quemadmodum docet Augustinus in Psal.141.) cuiquam Christiano vnquam irrogauit infamiam. Et longe distat haec Ecclesiastica censura a ciuili confessione. quae Amenda honoraria vocatur: quod haec hominum ratione fit, illa Dei: haec pro poena statuta est, illa pro aedificatione Ecclesiae, et testimonio aeternae salutis.  \pend\pstart Caeterae obiectiones sunt leues et nullius momenti. Itaque a nobis praetermittentur hoc loco.  \pend
\phantomsection
\addcontentsline{toc}{subsection}{\textit{21 Obtestor in conspectu Deï, et Domini Iesu Christi, et electorum Angelorum, ut haec serues absque eo, vt vnum alteri prae}}
\subsection*{\textit{21 Obtestor in conspectu Deï, et Domini Iesu Christi, et electorum Angelorum, ut haec serues absque eo, vt vnum alteri prae}}
\section*{AD I. PAVL. AD TIM. }
\marginpar{[ p.828 ]}
\phantomsection
\addcontentsline{toc}{subsection}{\textit{feras, nihil faciens in alteram partem declinando. }}
\subsection*{\textit{feras, nihil faciens in alteram partem declinando. }}\pstart Alius canon, qui tamen ex superiori pendet. Itaque hic versiculus est ἀκολέθησις. Cùm enim de iudicio aduersus Presbyteros instituendo supra praeceperit, docet, quale hoc iudicium esse et fieri debeat, ne in iudicando peccetur.  \pend\pstart Est autem huius canonis summa, summam in Ecclesiasticis iudiciis synceritatem adhibendam esse, ne vel in personarum iudicandarum acceptatione, vel ipsius iudicantis temeritate peccetur Hoc autem praeceptum quanquam a Paulo in iis iudiciis, quae aduersus Presbyteros et Pastores Ecclesiae exercentur traditum est, tamen in vniuersum ad omnia iudicia, siue politica, siue Ecclesiastica pertinet, de quocunque siue priuato, siue ἠγυμένω illa fiant. Atque hunc canonem in omnibus Ecclesiasticis iudiciis obtinere debere praeclare traditum est a Theodorito lib Histor. 2.cap.16. et a Libero Rom. Episcopo responsum aduersus Constantii imperatoris minas Athanasium ex praeiudicio quodam deponi praecipientis. Hoc idem etiam praeceptum in politicis iudiciis locum habere apparet ex 2. Chronic.19. vers.6.adeo vt maximi sit vsus in omnibus rebus, sitque ipsius iudicum conscientiae dirigendae tanquam norma, vt recte officium faciant. Nos tamen de Ecclesiasticis tantùm iudiciis hic agemus Pauli vestigiis insistentes. Non nude autem, neque simpliciter hoc praecipit Paulus, sed cum obtestatione, eaque grauissima, et terribili, quemadmodum loquitur Chrysost. Ratio autem huius obtestatio\pend
\section*{CAPVT  V. }
\marginpar{[ p.329 ]}\pstart nis tam seriae est duplex. Primùm quod de re ma ximi momenti, et late fusa et patenti acturus excitare eos voluit quos affatur, et ad quos hoc praeceptum pertinet, quos hoc fulgure et stimulo per stringit. Agitur enim de Iudiciorum synceritate, ex qua vna re tum vitae et societatis humanae conseruatio:tum ecclesiae ipsius vinculum, pax et disciplinae omnis ratio pendet. Itaque hac synceritate neglecta, aut violata multam rerum confusionem et perturbationem tum in politicis, tum in ecclesiasticis rebus sequi necesse est. Altera ratio ad dendae huius obtestationis est rei ipsius, quae prae cipitur, difficultas. Nihil enim tam arduum et dif ficile etiam optimo cuique etcordatissimo, quam illam synceram, aequam, et mediam viam in iudicando tenere, vt in nullam partem vel personarum acceptatione et respectu: vel ipse proprio animi motu a veritate vel officio deflectas. Admonet emnime Paulus Pimotneum, non gregarium quempiam et nouitium pastorem: sed egregium et insignem. Itaque proponit Paulus nlobis alios nostrae sententiae iudices futuros, quorum reuerentia et metu affici debemus, ne quid vel hominibus vel nostris affectibus tribuamus.  \pend\pstart Singula vero huius obtestationis verba suntdi ligenter obseruanda et perpendenda. Ac primum quod vel ludices vel testes quos producit presentes esse testatur. Quod enim dicit perinde est, atque  si spectatores, ac αυτόωτας eos esse affirmaret deinde testes adhibet grauissimos nempe, Deum ipsum, Iesum Christum, et Angelos sanctos. Dei nomine hic proprie patrem intelligimus, quod Iesu Christi Filii seorsum postea facta mentio  \pend
\section*{AD I. PAVL. AD TIM. }
\marginpar{[ p.N ]}\pstart est. Patris enim et Dei voce complectitur Paulus quicquid generaliter ad diuinam naturam pertinet, et est trium personarum commune: quo eodem dicendi genere passim vtitur scriptura. Ergo et Patrem et filium ipsum, qua Deus est, et Spiritum sanctum testem esse nostrarum actionum; et earum imprimis, quae in ecclelia et in iudiciis fiunt, quibus duabus nobis ipsa Dei maiestas maxime representatur spectatorem pronuntiat Paulus. Nam Deus etiam nudas cordis nostri cogitationes nouit, et intuetur non tantùm ipsas actiones nostras perspicitque quae ab hominibus cerni non possunt. Inde illud est etiam apud profanos homines vulgare ὁμμα θεο͂ πάντα καθετα. Dei oculus omnia perspicit. Itaque recte testis et vindex ad uocatur Deus noltrarum actionum. qui vbique est, omnia videt, et omnium futurus est iudex. Similes autem huic scripturae loci sunt pene infiniti veluti, 1. Samuel.2o.ver.42.2.Chron.24. vers.22. Atque hic primus est testis.  \pend\pstart Iesum Christum secundo loco testem citat et producit Paulus, quem Dominum vocat. Dominus est Christus quodsit constitutus a patre rex et Dominus omnium, dataque illi potestas in colum et in terram. Item quod omne iudicium pater concessit filio, neque iudicat quenquam vt est Ioan5. vers.22. Sed et quatenus ecclesiae caput est Christus ecclesiasticis iudiciis praeest, neque tantum vt testis et spectator iis interest: sed vt iudex ipse futurus, ac rationem eorum a pastoribus, tam quam a Ministris a se delegatis, exacturus. Etsi ve ro Christus mediator et quatenus Iesus Christus est, homol est et Deus: tamen qua homo est,  \pend
\section*{CAPVT  V. }
\marginpar{[ p.3 ]}\pstart non est iam in his terris, neque in medio nostri versatur siue visibilis, siue inuisibilis praesentie, vt nostras actiones oculis cernat, modo: sed tantùm qua Deus est et omnia haec inferiora videt. Recte tamem Iesus Christus dicitur testis omnium actio num nostrarum, quia ille Deus et homo, est vnica tantùm persona, quae hypostatica vnione vtraque illa natura constat. Neque hic curiose disputandum est, Vtrum humana Christi natura, quae in coelos euecta est, omnia, quae hic in terris geruntur, ex sese sciat. Quanquam enim supra omnes angelos dignitate! est euecta, tamen illa non est in deitatem mutata, vel deificata. De solo Deo dicitur illum, quae hic fiunt, omnia nosse et videre: quanquam Christum, qui homo est ad Dei dexteram exaltatus, non tantùm ea scire, quae angeli, fatemur: sed omnia, quae ad ecclesiae suae administrationem, totiusque  huius mundi curam quam a patre acceptam habet, pertinent et spectant, habere nota et perspecta.  \pend\pstart Addit, electis Angelis. Primum electos vocat, vt eos distinguat a malis Angelis qui vulgo dicuntur, et qui a Christo Angeli Diaboli, id est, maledicti nominantur in Math.25. vers. 41. deinde docet, qu hi ipsi angeli propter custodiam fidelium sibi a Deo demandatam intersint ecclesiae coetibus, et actionibus ecclesiasticis, tanquam testes et spectatores pastorum negligentiae vel temeritatis: vel perfidiae, si quam admittant. Merito igitur ter tio loco citantur. Sunt enim spiritus Ministratorii, qui ad subsidium fidelium et piorum a Deo mittuntur. Itaque actiones nostras norunt. Incidunt autem in hunc locum variae quaestiones, quas  \pend
\section*{AD I. PAVL. AD TIM. }
\marginpar{[ p.B ]}\pstart breuiter tantum perstringemus, et explicabimus Ac primum igitur quaeritur, velutia Basilio libro de spiritu sancto ca.13. vtrum licuerit Paulo, Ange los Christo et Deo adiungere. Num vero ista sit blasphemia parem facere creaturam creatori, eamque  illi honore adaequare. Id quod plane idololatricum est. Chrysostomus responder per modestiam id sieri a Paulo, vt pote qui non solum Deum testem, qui nobis est inaccessibilis longeque dignissimus adhibere voluerit: sed etiam creaturas ipsas, quae sunt nobis notiores ac familiariores : imprimis igi tur Angelos bonos commemorasse, qui sunt nobis conciues Ephes.2.ver.19. Haec Chrysostomus. Respondeo vero distinguendum esse finem, proter quem et Deus et Angeli in testimonium pro ducuntur. Etsi enim vtrique appellantur, vt sit maioris momenti haec obtestatio, et percellantur ve hementius lectorum, et eorum ad quos haec admonitio spectat, animi, tamen et Deus et Christus non tantùm producuntur vt testes a Paulo, sed etiam vt vindices futuri, et ludices pastorum culpae, et negligentiae: Angeli vero proponuntur et appellantur tantùm, vt testes. Ex quo fit, vt non adaequentur Deo; Angeli: neque habeat hic locus quicquam, quod idololatriam sapiat, aut in Deum blasphemiam.  \pend\pstart Testes vero harum rerum, quae in ecclesia geruntur, citari merito posse Angelos certum est, et eos imprimis, quibus a Deo ecclesiae cura demandatur vt ex Daniele. 1o. vers. 13. apparet. Et huic nostrae sententiae consentit Bernar. ser. 7. cantic.cantic. Doleo perinde aliquos vestrum graui in facris vigiliis deprimi somno, nec coeli ciues  \pend
\section*{CAPVT . V. }
\marginpar{[ p.333 ]}\pstart reuereri: sed in praesentia principum tanquam mor tuos apparere, cùm vestra ipsi alacritate permoti vestris interesse solenniis delectentur. Angelos igitur inuisibiles cum piis versari scribit. Itaque Christus ipse venturus dicitur cum Angelis illis beatis Luc.9. vers.26. Iudae. vers.14 Et ipsi gaudent si quis peccator ad Deum conuertatur. Luc.15.ver. 1o. quia id sciunt, cùm cura piorum quibusdam sit commissa. Quod idem de sanctis et in coelum receptis hominibus dici non potest. Nam idem munus non est illis a Deo impositum et demandatum, vt nos in viis nostris conseruent. Quod tamen est a Deo Bonis angelis praeceptum. pl. 91. Denique Paulus in ti Gormtmnin rerl.lo.et Epuel.3.ver.1o. praesentes et spectatores eos interesse coetibus ec clesiasticis docet et asserit. Propter quos etiam ait honestum esse et decorum, vt foe minae Christianae velentur in Dei ecclesia et conuentu.  \pend\pstart Secundo loco quaeritur, Cur hos electos appellet: Num etiam alii, qui mali dicuntur Repro bati a Deo fuerint, et dici possint. Resp. Etsi illa electio et reprobatio, de qua tam saepe in scriptu ra fit mentio, proprie ad homines pertinet, non ad Angelos: tamen cùm nihil praeter Dei voluntatem acciderit etiam in lapsu angelorum, imo vero cum ex Angelis alii ceciderint, alii perstite rint, Deo ipso ita ordinante et praedestinante. Deni que cùm ex ipso Dei decreto hoc totum ita euenerit, dubitari non posse, eos Angelos, qui perstiterunt in sua origine, a Deo electos quidem me ra ipsis gratiua fuisse: Eos autem, qui ceciderunt, reprobatos ab eodem Deo: idque iusta de causa, quanquam nobis ignota. Voluit enim Deus hac ra\pend
\section*{AD I. PAVL. AD TIM. }
\marginpar{[ p.334 ]}\pstart tione, et in angelis, et in hominibus, suae tum misericordiae, tum etiam iustitiae diuitias ad gloriam nominis sui, patefacere et demonstrare. Nam licet opifici facere de opere suo, quod ipsi libucrit Ergo boni Angeli electi fuerunt, etiam ante quam a Deo crearentur. Non autem ex vllis eorum meritis: sed ex sola Dei ipsius dignatione et gratia, vt loquitur Bernardus: quicquid Scholastici ex Augustino male intellecto dicant et philosophentur: Mali vero reprobati fuerunt, vt Boni electi.  \pend\pstart Tertio loco quaesitum est, ecquae sit istius electionis Angelorum causa: Num Christus quatenus me diator ecclesiae: cuius etiam illi sunt mem bra: praesertim cùm Paulus Coloss.1.vers.2o. scribat Deum sibi reconciliasse omnia per Christum. Ex quo colligunt quidam Christum esse etiam electorum Angelorum Mediatorem, alio tamen modo, quam hominum piorum et electorum Resp. vero ipse, ceterorum rationibus breuitatis causa, omissis, Christum quatenus creator eorum est, non autem quatenus passus in carne, videri horum bonorum Angelorum etiam mediatorem esse. Cùm enim Deus numquam iis fuerit offensus, nulla reconciliatione certe opus habuerunt. Dein de haec eorum electio a Iuda desinitur conseruatio et confirmatio ipsius eorum primi status et ori ginis, adeo vt nihil prorsus illis restitutum fuisse per huiusmodi electionem dicatur, quoniam nihil prius fuerat ademptum, neque quicquam ab iis amissum. Ergo ipsi quales a Deo primùm conditi fuerunt, Deo autori suo certe placuerunt Ne que vera tamen huius eorum electionis causa est,  \pend
\section*{CAPVT  V. }
\marginpar{[ p.335 ]}\pstart qualitas illa et status in quo tantùm creati sunt. ( Omnes enim angeli, etiam qui ceciderunt, dicerentur et fuissent electi, quia in eodem statu felici omnes conditi sunt) sed praecessit haec Dei e-- lectio etiam illorum creationem et ortum, et ab aeterno in Deo fuit. Hunc statum tam excellentem omnes quidem a Deo mera ipsius gratia erant con secuti, sed Dei illa gratia per quam sunt alii dicti electi, et prior est, et maior gratia creationis. Hoc vero qui a Deo electi fuerunt, plus iam sunt adepti post malorum lapsum, qu de feliciss. illius conditionis suae statu perpetuo futuro iam a Deo certissimi effectit et in eo stabiliti sunt, et in aeternum confirmati: Id quod, Deo illis reuelante, iam cognoscunt et intelligunt. Bernandus in Sermo. 5de festo omnium sanctor. Inter Angelos et homines assignari potest ista diuersitas sanctitatis. Neque enim tanquam triumphantes honorari pos se videntur, qui nunquam pugnasse noscuntur. Aliter tamen honorandi sunt etiam ipsi, tanquam amici tui Deus, cuius nimirum voluntati semper adhaeserunt, tanta vtique felicitate quanta facilitate. Aliud igitur in Angelis sanctitatis genus. Et ser.5. de visio. Non est innata eis sed a Deo collata iustitia, quae inferior est diuina iustitia. Idem Bernardus, serm.22. cantic. cantico. Si angeli, ait, nuquam redempti sunt, alii vtique non egentes, alii non promerentes: illi quidem quia nec lapsi sunt, hi autem quia irreuocabiles sunt, quo pacto tu dicis Dominum Iesum Christum eis fuisse redemptio nem? Audi breuiter. Qui erexit hominem lapsum, dedit stanti angelo, ne laberetur: sic illum de captiuitate eruens, sicut hunc a captiuitate defendens.  \pend
\section*{AD I. PAVL. AD TIM. }
\marginpar{[ p.530 ]}\pstart Et hac ratione fuit aeque vtrique redemptio soluens illum hominem, et seruans istum angelum. Liquet ergo sanctis angelis Christum Dominum fuisse redemptionem, sicut iustitiam, sicut sapien tiam, sicut sanctificationem. Haec Bernard. Augusti in Enchirid. ad Laurentium cap. 28. idem sentit.  \pend\pstart Ac de obtestatione quidem haec dicta sunt. Vi deamus iam, a quibus vitiis pastori iudicium ecclesiasticum exercenti sit cauendum, atque fugiendum.  \pend\pstart Duo autem sunt nimirum. 1. Praelatio personarum propter ipsarum dignitatem 2 Temeritas, propter iudicantis affectum. Neminem enim propter dignitatem excusari debere censet Paulus. Neminem etiam propter affectum iudicis praecipitanter, et inconsiderate damnandum esse vult. Quae duplex ratio magnos in Iudicando errores inducit, ex quibus tandem euertitur ecclesia. Caussae igitur propter quas peruertuntur iudicia sunt in nobis: vel in illis ipsis, de quibus iudicamus, a nobis quaeruntur. Vocem Προκρίμα- τος ad praeiudicia illa referunt, per quae non persona: sed causa potius Presbyteri, qui iudicandus est, vel grauatur, vel leuatur, id est, antequam cognitum sit a nobis de eo ( de cuius causa agitur) iam animi quadam praecipiti sententia domi concepta definitum est. πρόσκλισιν autem eam esse volunt in iudicando rationem, quae leges ordinarias et statutas praetermittit, vt faueat illi, de quo agitur. Sed Graecae vocis significatio non patitur hanc interpretationem, potiusque sequenda nobis est D. Bezae explicatio. Vtcunque Paulus  \pend
\section*{CAPVT  V. }
\marginpar{[ p.337 ]}\pstart non omnia iudieiorum peruersorum vitia nunc enumerat: sed ex iis aliqua, ex quibus paucis caete ra colligere licet. Gregorius Magnus.4.esse summas iudiciorum corrumpendorum causas scribit scilices Timorem, Cupiditatem, Odium, Amorem, cuius sententia descripta tie mi Ouh. Quatuor modis.Il quaeit.3. Quae sane aliqua ex parte exponit, quae breuius hic Paulus recenset.  \pend
\phantomsection
\addcontentsline{toc}{subsection}{\textit{22 Manus cito ne cuiimponito, neque communicato peccatis alienis: temetipsum purum serua. }}
\subsection*{\textit{22 Manus cito ne cuiimponito, neque communicato peccatis alienis: temetipsum purum serua. }}\pstart Canon praestantissimus, qui etiam ad pastorum erga pastores officium pertinet. Nam superiores canones ad pastores iam creatos spectabant: hic autem de Pastoribus adhuc creandis agit. Hoc vno breuiter complexus est Paulus optimam et totam pene eligendorum pastorum rationem, quae nobis hodie obseruanda quoque est, vt legitime vocatos pastores habeamus. Quae res tanti est in ecclesia momenti, vt propter eam superiorem obtestationem a Paulo praemissam esse censeat Chrysosto. Certe cùm de ea re monitum Timotheum esse oportuerit: et magni periculi plenam esse, et grauis momenti in ecclesia ordinanda, iam dubitari non debet. Haec autem praecepti summa est. Neminem in ecclesia Dei cito et temere promouendum esse ad aliquod munus eccle siasticum, quantumuis exiguum: de eo enim muneris genere, idest ecclesiastico, Paulum agere, non de politicis Magistratibus apparet ex phrasi ipsa qua vtitur. Ait enim Ne manus imponas. Non quod  \pend
\section*{AD I. PAVL. AD TIM. }
\marginpar{[ p.338 ]}\pstart tamem et ipsi ciuiles magistratus θιῦ διάκονο et sacrretiamnomdney seu me cecienastica tantum persequitur Paulus.  \pend\pstart Quaeritur autem, quae sit huius praecepti occasio, et cur de eo fuerint etiam prudentissimi cuiusque ecclesiae pastores, qualis Timotheus, qualesque nobis sub Timothei persona proponuntur, monendi. Resp. fuisse multiplicem necessitatem ac rationem cur hoc praeceptum spiritus sanctus adderet: sed imprimis fuit triplex, nimirum sci licet propter Impudentiam, Importunitatem tum Populorum tum Ordinandorum, et Nouitatis studium erat necessarium hoc praeceptum. Nam fere homines nouis Ministris gaudent: Ergo ne qui semel arrisit, statim probetur et eligatur, cauendum est. Nam permulti sunt, in quibus egregiae quidem animi dotes elucent, qualis est vel ze lus in Dei gloriam, vel etiam summa doctrina, quique propterea sunt commendabiles: qui tamen ad munus ecclesiasticum sunt prorsus propter alios defectus inepti, illiusque incapaces. Et certe qui ad huiusmodi muner a alios commendant, hoc notent diligenter, quod ait hoc loco Paulus, quódque Horatius Flaccus, quanquam poëta prophanus, recte monuit. Qualem commendes etiam atque etiam aspice, ne mox Incutiant aliena tibi peccata pudorem.  \pend\pstart Affert autem remedia aduersus omnes superiores rationes, quibus impulsi eligentes peccare solent. Primum igitur remedium primaque excusatio iustissima, quae contra afferri debet est haec, Neminem Dei praeceptis praeponendum, et idcir  \pend
\section*{CAPVT . V. }
\marginpar{[ p.339 ]}\pstart co vtcunque precibus instent et vrgeant nos alii, non esse alienis tamen peccatis communicandum. Nam qui indignum eligit, vel eligi consentit, reus est omnium errorum et peccatorum, quae postea electus ille committit. Deinde consentit cum male agentibus, quatenus ipse temerariis eorum suffragiis consensum suum addit atque adiungit. Verum est enim illud, Agentes et Consentientes pari poena puniuntur: Iacob patriarcha Genes.49. vers. 6.negat se malo consilio Simeonis et Leui de necandis Sichemitis assensum esse. Ps.5o. vers.18. Qui cum furibus currunt, pro furibus habentur apud Deum. Denique ait Paul.1. Thess.5.ab omni malo malique specie esse nobis abstinendum. Neque hoc cuiusquam gratiae concedendum a nobis est, vt in illius fauorem peccemus, officium nostrum negligamus, et Deum ipsum deseramus atque offendamus, cui nos nostri suffragii et vocis rationem reddere aliquando oportebit.  \pend\pstart Quod si iterum instent et premant nos alii, qu plura sint aliorum et contraria nobis suffragia, nos tamen puros seruemus, quid nobis videatur, libere declaremus, et nostra consilia ab huiusmodi temerario suffragio indignis hominibus dando separemus. Atque hoc est secundum remedium, siue secunda legitima excusatio et praemunitio boni pastoris aduersus instantes quorundam preces et suffragia iam lata itemque aduersus calumnias multorum, qui nos putant, dum nostra suffragia quibusdam iam ab aliis electis et probatis denegamus, vel tam cito in illis non assentimur, duci aliquo in eos odio vel inuidia, vel etiam nos esse  \pend
\section*{AD I. PAVL. AD TIM. }
\marginpar{[ p.340 ]}\pstart crudeles et nimium seueros. Sed valeat potius apud nos haec tutissima consolatio, quod puram conscientiam seruare volumus ac studemus. Neque tamen, si a maiori suffragiorum parte vincamur, solique dissentimus, seditionem in senatu, aut coetu aut ecclesia Dei mouere debemus: aut quicquam perturbate agere, vt vincat nostra vnius vox: sedliberauerimus animas nostras atque officio functi erimus abunde, si primùm nostram sententiam illiusque rationes exposuerimus: deinde, si res sit tanti momenti, eo ordine vtamur ad huiusmodi electionem impediendam, vel retardandam, qui et legitimus est, et in ecelesia Dei constitutus et receptus, qualis est nimirum prouocatio a lenatu ecelenlartieo au maiorem Ministrorum vicinorum coetum, quem Classem vocant: a Classe ad Synodum ipsam prouincialem: a Synodo prouinciali ad synodum Generalem et Nationalem quam vocant.  \pend\pstart Ait Paulus, Cito. Fit autem Cito ex Pauli mente, quando vel non legitimum: vel non satis exactum et certum examen de eligendo siue promouendo pastore praecessit: sed vel nullum omnino, vel leue tantum et perfunctorium, quale fit in Papatu, et in multis etiam mediis ipsis Dei viui eccleliontque de omnibus iis rebus, quae ad tale ministerium obeundum rite desiderantur. Ergo quem admodum sunt diuersae vocationes in ecclesia Dei: et aliae aliis sunt digniores et grauiores: ita serius et maius in his, quam illis examen esse et fieri debet, veluti in Pastore ipso, quam in solo Presbytero vel Diacono plura inquirantur necesse est, quia maius est pastoris quam Diaconio\pend
\section*{CAPVT  V. }
\marginpar{[ p.241 ]}\pstart nus.Ac quidem examen, quod in pastore requiritur, est fere triplex:fieri enim debet de Doctrina ipsius Vita et moribus, et Aptitudine ad docendum gregem et populum Dei. Id quod satis ex Tit. 1.et eadem epistola c.3. apparet.  \pend\pstart Et quia luculentus est hic locus, continetque maximum totius Politiae ecclesiasticae caput et pulcherrimum (quod est de electione et vocatio ne dignitatum ecclesiasticarum) de eo nobis bre uiter aliquid arbitramur esse dicendum. Habebit aute lic iocus 4 capita ichicer l'Ouid sit electi 2 Qui eligant 3 Quando 4 Quomodo fieri debeat electio. In primis autem hoc pro certissimo fundamento constituatur, nempe quod  \pend\pstart Electio vocationem ordinariam praecedere debet, quemadmodum vocatio ipsa anteit muneris ipsius susceptionem et functionem, si modo bona conscientia volumus munus aliquod in ecclesia Dei gerere ac administrare. Vocationem autem, de qua nunc agimus, κλησίαν aut χειροτονίαν potius quam ὀκλογην appellandam esse dicimus, quod ὀκλογὴ ad causas salutis nostrae pertineat, vti docet Paul. Rom.11.vers.5.quae disputatio est reiliota ab noc argumento. Eit enim vocatio haec sepositio designatióque tantùm alicuius personae ad munus aliquod ecclesiasticum gerendum.  \pend\pstart Ac primùm neminem sine vocatione legitimain ecclesia Dei recte sanaque conscientia munus vllum abtinere posse docet tum autoritas sacrae scripturae: tum etiam ratio.  \pend\pstart Ac scriptura quidem Hier.23. vers.21. currebant, et ego non mittebam eos. Hebus 5.vers.5. Ne mo sibi sumit honorem, sed qui vocatur a Deo.  \pend
\section*{AD I. PAVL. AD TIM. }
\marginpar{[ p.42 ]}\pstart Roma. 1o. vers. 15 Quomodo praedicabunt, nisi missi fuerint.  \pend\pstart 1.Timoth.4. ver.14. Ne negligitodonum quod mte est, quond eit tibi datum ad prophetandum per impositionem manuum Presbyterii.  \pend\pstart Tit. I. vers.5. Idcirco reliqui te, vt oppidatim constituas Presbyteros.  \pend\pstart Denique infiniti pene sunt loci his similes, quos illi studiose collegerunt, qui de ea re communes locos scripserunt: nos autem summa tantum rerum fastigia sequimur, et carpimus, vt simus oieulores.  \pend\pstart Ratio vero, cur sine legitima vocatione praecedente nemo bona conscientia vllo munere ecclesiastico fungi possit, est multiplex.  \pend\pstart Prima, quod in ecciena Der omnia fieri ordine certaque ratione, et decenter debeant 1. Corint. 14. Ergo nemo pro animi sui libito hanc vel illam dignitatem sibi in ecclesia sumere debebit, quia haec est summa ἀταξία et rerum confusio.  \pend\pstart Secunda. Quod in veteri ecclesia, quae Christi aduentum praecessit, nemo nisi ex Deo vocante et praecepto munus habuit: In ea etiam ecclella, qua Emfitus per euaiigeii praedicationem instaurauit, Apostoli ipsi a Christo vocati sunt. Ergo praeter et legalis, et Euangelicae eeclesiae for mam, praecepta, et instituta idest praeter ipsum Dei verbum in munere aliquo versantur, qui illud citra vocationem vllam legitimam vsurpant, et in ecclesia exercent.  \pend\pstart Tertia. Neque Moses, et Aaron, neque  Apostoliac ne Christus quidem ipse quanquam Dei filius, nisi prius legitime vocatus, munus sibi assum\pend
\section*{CAPVT V. }
\marginpar{[ p.343 ]}\pstart psit in ecclesia Dei. Ergo nemo sibi honorem tri buere meccienla Derjeitra legitimam vocationem debet.  \pend\pstart 4. Dominus Deus est ecclesiae tanquam peculiaris suae domus et proprii peculii solus paterfamilias et legislator, et dispensator. Ergo qui anquiu mea citra iprius voruritatem tiarpat, fur est, et latro: non autem legitimus illius administrator et gestor.  \pend\pstart 5 Qui, praeter legitimam vocationem ecclesiasticam dignitatem adit, perinde facit, atque qui domum alienam alio’aditu et modo, quam per ostium ingreditur. Hoc autem furum est, et latronum, non autem legitime ingredientium. ergo qui hac ratione, id est, sine legitima vocatione ec clesiasticum munus gerunt, sunt fures, non Pastores ecclesiae.  \pend\pstart Quod autem nobis obiicitur de Prophetis, eos sine vocatione saepe fuisse, et munus illud suum in ecclesia Dei exercuisse, falsissimum est, vt suo exem plo docet Amos, cap.7.vers.15. Item Hier.1.ver. 5.Sed vocatio illa prophetarum vt plurimum fuit extraordinaria. Eos enim Dominus ipse vocabat et excitabat, ex qua tribu volebat, quo tempore volebat, sed vocabat tamen. Deinde eorum vocationem Deus vel ipso rerum ab iis praedictarum euentu: vel aliis modis aperte confirmabat, adeo, vt illi ab eo missi esse satis ab omnibus ad eorum prophetias attendentibus intelligerentur. Idem est respondendum de Christo cuius vocatio et a Deo ipso in ipso baptismo fuit patefacta, etiam olim a prophetis praedicta, etsi non fuit ordinaria.  \pend\pstart Aeriani haeretici damnati sunt: et nunc Anabs\pend
\section*{AD I. PAVL. AD TIM. }
\marginpar{[ p.STF. ]}\pstart ptistae, et Enthousiastae quidam nebulones merito reiiciuntur, qui vocationes in ecclesia Dei se cure spernunt, vt se ingerant, atque intrudant ad ecclelialtiea munera magno animi fastu et ambitione.  \pend\pstart Ergo propter tria praecedere legitima vocatio debet, antequam quis munus aliquod in ecclesia Dei gerat vel sumat. I Propter Dei praeceptum. 2 Propter ecclesiae ipsius ordinem et disciplinam. 3 Propter ipsius, qui dignitatem illam accipit, conscientiam, quae alias nunquam tuta et quieta esse potest:non magis quam furis quandiu re a se furto ablata vtitur, vel eam retinet. Magnam enim animis nostris consolationem affert legitima vocatio, quoties a peruersis, vel obstinatis nommous vexamur, et negotia nobis praebet Satan. Quid est? Resp.  \pend\pstart Vocatio igitur haec, est ad munus aliquod facta ab alio, quam a seipso, adoptio ac designatio. Haec igitur duplex est. Interna et Externa.  \pend\pstart Interna est a Deo, ac, ea quidem immediate, Externa vero, etsi a Deo est, quia ipse illius celebrandae modum verbo suo praescripsit: mediate tamen est ab hominibus illam Dei ordinationem sequentibus. Neutra sine altera tuta satis esse potest vel salutaris: Interna tamen potior est: sed externa nostri ;atione tutior.  \pend\pstart Interna vocatio est interna et latens animi nostri affectio atque inclinatio a Deo immissa ad aliquod legitimum munus gerendum, quae illius desiderium in animo nostro parit. Igitur affectus infunditur a Deo, neque ex ambitione, neque ex auaritia,nec ex alio aliquo turpi vel carnali motu  \pend
\section*{CAPVT V. }
\marginpar{[ p.345 ]}\pstart in nobis inesse debet. Bernard.serm.58.in Cant. sic eam definit, Est inuitatio et stimulatio quaedam charitatis pie nos sollicitantis aemulari fraternam salutem, aemulari decorem domus Dei, incrementa lucrorum eius, incrementa frugum iustitiae eius laudem et gloriam nominis illius Istiusmodi ergo erga Deum religiosis affecti bus quoties is qui animas regere, aut studio praedicationis ex officio intendere habet, hominem suum interiorem senserit permoueri, toties pro certo sponsum adesse intelligat: toties se ab illo ad vineas inuitari.  \pend\pstart Externa vocatio est, hominum qui iudicare ponune de nostra ad aliquod munus capacitate et idoneitate testimonium libero, non redempto non captato ipsorum suffragio, in Dei timore et via a Deo praescripta atque legitima patefactum et declaratum.  \pend\pstart Vt sit autem haec externa vocatio nostra legitima, duo haec concurrere debent:e quibus tamen vnum alterum praecedit, nimirum Electio, quae praecedit: Ordinatio seu approbatio electi, sequitur electionem.  \pend\pstart Electio est, eorum qui possunt et debent syncerum atque sanctum de nobis testimonium, per quod designamur ad certum munus gerendum.  \pend\pstart Ergo hic videndum 1 A quibus fieri debeat electio. 2 Quando, 3 Quomodo. Ac quidem primum illud statuendum est, quod Nemo a seipso eligi vel debet, vel potest. Verum enim illud est, Nemo iudex ferendus in propriae ratione causae:item et illud, Laus in proprio ore sordescit. Neque obstat quod supra de interna eligen\pend
\section*{AD I. PAVL. AD TIM. }
\marginpar{[ p.340 ]}\pstart dorum vocatione diximus. Etsi enim esse debet, interna vocatio et affectus in eo qui munus aliquod acceptat: ea tamen interna vocatio etiam externam, tanquam sensus illius interni testem, adhibet, et expectat, nisi quis extraordinarie a Deo ipso impelleretur. Interna enim vocatio, quae externam spernit, non tantùm falsa est, et mera arrogantia, sed etiam violat Dei praeceptum: et turbat disciplinam Ecclesiae: neque a Deo est, sed est animi impia phrenesis. Id quod etiam in quadam sua ad Magnum epistola Cyprian. respondit. Sed et ipsum electionis et vocationis nomen huic a seipso de seipso factae ordinationi plane repugnat, cùm is demum dicatur Vocari, qui ab alio accersitur.  \pend\pstart Sed neque a quolibet Electio facta legitima est, in muneribus praesertim Ecclesiasticis. Nam ab vno tantùm aliquo facta alicuius ad dignitatem Ecclesiasticam electio, rata et legitima non est, siue ab ipso principe illa fiat, siue ab eo, qui Ecclesiae patronus dicitur, siue ab eo solo qui in Dioe cesi Metropolitanus est, et Ecclesiae πρνεσῶς et πρόεδρος. Neque  enim Ecclesiae, quae in terris est, regimen et administratio nobis a Christo relicta, est Monarchia quaedam: ac multo minus tyrannis: sed est Aristocratia, vt docemur parabola Christi, quae est Matth.24. vers.49. et 25.vers.14. Itaque vetat Petrus etiam ipsos Episcopos Kυριέυαν in Ecclesia: et Christus ipse dissimilem plane regum in suos subditos: et Pastorum in Ecclesias suas potestatem esse dixit 1.Pet.5.Luc.22. Itaque qui omnes similes faciunt, omnia perturbant. Quare ne a rege quidem ipso solo, vel ab vno dun\pend
\section*{CAPVT  V. }
\marginpar{[ p.347 ]}\pstart taxat Episcopo vel Archiepiscopo, facta electio legitima est, quantunuis magna censeatur inter homines istorum authoritas, vt praeclare annotat Theodorus lector lib2. Collecta. Sed in Dei Ecclesia Christi vnius lex dominari debet: quae nulla ratione infeingi potest, vt recte administrata Ecclesia dici possit. Praeter enim vnum Christum nemo est in Dei Ecclesia ἀυτοκράτωρ: nemo illius domus Dominus: sed solus Christus est ille vnus, extra cuius praeceptum nemo in ea dignitatem legitimam obtinet. Non vult autem Christus ab vno quodam solo patrono, vel Rege, vel Pontifice et Episcopo aliquos eligi et asciri ad munera Ecclesiae suae, vt perspicue postea apparebit. Quod ipsum recte quoque vetitum est fieri ab Adriano Rom. Pontifice, vt est in canone Nullus. et a Nicolao in cano. Porro, et can. seq. Distin.63. Atque  inde orta sunt grauissima illa, et plus quam nefaria bella inter Henricos imperatores Rom. et Gregorios Pontif. Rom. cùm vtrique ad se solum ius illud electionum trahere conarentur. Nec obstat, quod Titus, Timotheus, Paulus soli videntur Presbyteros et Episcopos Ecclesiis praefecisse, vti ex hoc loco videtur colligi posse, et Act. 16. Tit.1.Respondeo enim illud genus dicendi Constituas, imponas, etc. ex mente Pauli esse sic intelligendum: Non vt vni cuidam soli electio tota committatur iusque illud tribuatur: sed vt synecdochicons complexus esse ea dicen di ratione intelligatur omnes Paulus, et omnia, quae fieri in legitima vocationis Ecclesiasticae ratione et forma debent et solent (vt ex his verbis Tit.1.V.5. Vt cibi mandaui) facile concludi potest. Sed cùm in  \pend
\section*{AD I. PAVL. AD TIM. }
\marginpar{[ p.oi5 ]}\pstart iis electionibus praecipuae quaedam partes essent, tum Titi, tum Timothei, tum ipsius Pauli, et Barnabae, vt pote qui Ecclesias ipsas quomodo id fieri oporteret, erudiebant, et qui huiusmodi electionibus, tanquam totius actionis moderatores intererant, atque praesidebant: idcirco tota actio illis quodammodo tribuitur, quae tamen et torius Preibyterii, et totius ipsius Ecclesiae communis est, et tunc erat, vti apparet ex capite 4.v. 14. supra et Actor. 14 vers.23.  \pend\pstart Neque etiam obstat, quod Episcopus προεδ ρός appellatur a Patribus, quasi vnus ille totam Ecclesiam pro arbitrio possit administrare, et in ea quos velit, praeficere. Nam vox illa πρόεδ pος ordinem tantùm declarat, quod inter collegas suos, eticompresbyteros sedere deberet ipse, qui Episcopus dicebatur: non autem regiam, aut summam praetoriamque potestatem in Ecclesiam illi triDuit: Quod diemas confi mat fimbrosius liber  de Dignit. Sacer dot. cap.6. Solio in Ecclesia editiore Episcopus residet, ait, vt cunctos aspiciat. Etictum quoque Christi quod supra ex Luc.22. conmemorauimus, idem probat. Est enim solus Christus Ecclesiae rex, non autem mortalis quispiam.  \pend\pstart Pertinet autem electio, et designatio ipsa ad totam Ecclesiam, in qua fit: sed diuersa ratione. Quemadmo dum enim totius Ecclesiae pastor est futurus: ita aly omnibus debet ordinari et approbari, ne quisquam gregi inuito pastor obtrudatur. Praeterea etiam illud verissimum est, quod omnes tangit, et spectat, ab omnibus fieri debere. Huiusmod enim negotium totius Ecclesiae, non vnius autem illius partis negotium est, cùm qui  \pend
\section*{CAPVT  V.. }
\marginpar{[ p.349 ]}\pstart pastor eligitur, Ecclesiae vniuersae detur.  \pend\pstart Ergo pessime faciunt, qui plebem quam vocant ab omni suffragio in Ecclesiasticarum dignitatum electionibus ferendo repellunt et semouent, tanquam non sit et ipsa plebs pars Ecclesiae Dei, eaque maxima. 4  \pend\pstart Pessime quoque hallucinantur qui non distingunt quae sint Presbyterii, quae autem plebis, in electionibus istis partes: sed quod est illius, isti tribuunt, inducentes in Ecclesiam Dei magnam omnino confusionem. Praeclare enim Ambrosius in liber  de Sacerdot. dignitate cap.3. Aliud est, ait, quod ab Episcopo requirit Dominus: aliud quod a plebe. Idem Chrysost. in libro3. de Sacerdotio infra medium. Idem quoque Bernard.sermo 49. Cant. cuius aetate iam haec perturbata et popularis in electione ordinandorum ratio, a quibusdam videtur fuisse inuecta. Ordinauit me in charitatem, ait. Factum autem est hoc cùm in Ecclesia quosdam quidem dedit Apostolos, quosdam Prophetas, alios Euangelistas, etc. Oportet autem vt hos vna omnes charitas liget et contemperet in vnitatem corporis Christi. Quod minime omnino facere potuerit, si ipsa non fuerit ordinata. Nam si suo quisque feratur impetu secundum spiritum, quem accepit: et ad quaeque volet, indifferenter, prout afficitur, et non rationis iudicio conuolarit, dum sibi assignato officio nemo contentus erit, sed omnes omnia indiscreta administratione attentabunt, uon plane vnitas erit, sed magis confusio. Haec Bernardus.  \pend\pstart Sic igitur tota haec res est distinguenda, vt in vocatione et electione duos quosdam actus di\pend
\section*{AD I. PAVL. AD TIM. }\pstart stinctos obseruemus et separemus, nempe Electionem ipsam siue praesentationem personae: et Acceptationen, quati alii Oruinationem, alii Confir mationem appellant.  \pend\pstart Haec igitur Electio siue praesentatio, quam dico est ipsius personae ad munus aliquod Ecclesiasticum vocandae prima agnitio, agnitae tentatio, degustatio, probatio, et examinatio, quae tum deipsius vita, tum de doctrina fieri debet, et ita examinatae atque consentientis denuntiatio et praesentatio, quae fit toti populo et Ecclesiae simul collectae. Hae vero sunt Presbyterii, et quidem totius, partes, non plebis aut populi.  \pend\pstart Approbatio autem eligendi est, de persona iam a Presbyterio examinata diligenter et proposita, Iiberum, sed tamen iusta aliqua ratione nitens totius populi et Ecclesiae suffragium, per quod post commodum aliquot dierum interuallum concessum, illam personam populus vniuersus, vel acceptat, vel reiicit atque repudiat. Haec igitur approbatio ad plebem et totum populum Ecclesiae sane pertinet.  \pend\pstart Hoc verum esse atque ita factitandum, vt sit legitima ordinariaque eligendorum et electorum vocatio, apparet tum exemplis veteris, et primitiue Ecclesiae: tum etiam exemplis illius, quae subsecuta est prima illa tempora. Denique firmissimis rationibus idem comprobatur.  \pend\pstart Ac primùm plebem non esse ab ordinationibus vocandorum ac praeficiendorummuneribus Ecclesiasticis excludendam, demonstrant exempla veteris Ecclesiae, in quaproculdubio χεροτονία toti' Ecclesiae interueniebat, vt ex Act.6.et 14 ostendi fa\pend
\section*{CAPVT V. }
\marginpar{[ p.35  ]}\pstart cile potest. Itaque suo iure perfide priuant Ecclesiam, qui pastorem populo inscio et non assentienti, obtrudant. Quod fit in Papatu, faciunt enim maximam Ecclesiae iniuriam, cùm eam iudicio et suo suffragio spolient. Qui propterea, vere sunt Sacrilegi nominandi. Nec enim est pastor legitimus qui inscio vel inuito, vel non assentienti gregi prae sidet et datur, vt ait Gregorius Magnus epistola 87. et 9o. Itaque olim etiam iam corruptis Ecclesiae moribus et inclinata disciplina, tamen irritae factae sunt omnes Ecclesiasticae vocationes, quae ἀνδι λαο͂ συνέσεως, id est, vt vertit Cyprianus, sine conscientla et assensu populi fierent. Quod confirmat Chrysostomus liber 3. de Sacerd. Cyprianus liber 1.epistola 4. liber 2.epist.3. et 5. Theodoret. in Histor. cap.2o. et 22.libus 4 Gregor. liber 1. epist 5.libus 2.epist.69. Atque de hac ipsa electionis forma seruanda extat constitutio Caroli et Ludouici Imper. in canon. Sacrorum canonum distinct.63.Leo quoque primus Pontifex Rom. negat vllum verum esse Episcopum, nisi qui a clero electus et a plebe expetitus fuerit. in canon. In nomine distinct.23. Quam ipsam etiam suo seculo adhuc fuisse obseruatam testatur variis in locis Bernardus epistola 2o2. et 150.164.13. et 27.id est, vt populus electioni assentiretur. Denlque hoc ipsum verum esse iure canonico etiam confirmatur, quemadmodum apparet ex canone, Plebis Diotrensis et canon. seq. et cano. Nolle distinct. 63 can. Factus, vbi dictum Cypriani aureum recitatur.7.quaest.1.Idem testatur Ambrosius liber 1o epistola 82. Vnde quam impii sint Trident. concilii canones huic contrarii manifeste apparet.  \pend
\section*{AD I. PAVL. AD TIM. }
\marginpar{[ p.552 ]}\pstart Neque tamen propter tumultum popularem, aut votum populi temerarium in vnum aliquem, qui erit indignus, debet ille a Presbyterio designari et eligi. Ostendit enim hoc loco Paulus nullo modo esse cum alienis peccatis communicandum. Quod etiam a Leone Episcopo Roma. responsum est, et refertur in cano. Miramur. distinctione 61. Sed aliud est denegare plebi iustum suumque suffragium: aliud autem illius temeritati obsequi, et obtemperare. Ac quidem haec sunt iura totius populi et Ecclesiae.  \pend\pstart Presbyterorum autem est, primùm personam ipsam, quae idonea ad munus vacuum videri potest, inquirere ex toto populo, eamque designare oculis, et de ea inter se agere atque conferre. Praeire enim debent toti Ecclesiae Presbyteri. Itaque non tantùm eum agnoscere de vultu vel facie debent, quem designant: sed etiam de qualitate. Vnde fit vt primum examen ipsius eligendi ad eosdem Presbyteros pertineat, ne indignum proponant populo. Nam parum est personam ipsam nosse, nisi ingenii vires noueris. It aque doctrinam ordinandi debent Presbyteri tentare, tum quaestionibus iln propontis, tun periculum facientes de illius dexteritate in tractanda propoheirdaque et interpretanda Sacra scriptura.  \pend\pstart Cur autem hoc primum suffragium ad Presbyteros pertineat, non autem ad vniuersam Ec. clesiam et populum, quemadmodum tamen quidam censent, et contentiose disputant, ratio est:  \pend\pstart Prima, Quod Presbyteri et Pastores dicuntur ποιμένες et ηγέμενοι Ecclesiae. Hebr.13.vers.7.id est, Praepositi (vt loquitur Tertull. in Apolog cap.  \pend
\section*{CAPVT . V. }
\marginpar{[ p.353 ]}\pstart 9.) Vnde sua voce suóque consilio praeire populo debent, et viam futurae deliberationis ostendeIe. Eigo non plebs praene debet iua voce Presbyris: sed Presbyteri populo et plebi.  \pend\pstart Secunda, Quod illi ipsi sunt, vt est in cano. Ecclesia habet 16. quaest.1. Ecclesiae senatus (vt ait Hieronymus) ad quem prima rerum Ecclesiasticarum deliberatio et cura pertinet, atque etiam referenda est:itemque, vt idem Hieronymus ait, morum censura. Sic Paulus Philip.1.vers.1.ad eos scribit, vt per eos ipsius Epistola innotescat reliquae Ecclesiae. Sic ad eos per manum Pauli et Barnarbae eleemosynae Ecclesiarum deferuntur. Act. 11.vers3o. Sic illi ipsi Presbyteri Hierosolymis conueniunt, et primum audiunt Paulum, tum eumdem Ecclesiae vniuersae sistunt Actor. 20. vers.28. 21.vers.18.11.vers.13.  \pend\pstart Tertia, Exemplum ipsius primitiuae Ecclesiae. Apostoli enim ipsi primi de Diaconis eligendis deliberant. i. de personis, quas ad eam dignitatem promouent. Ideo verba, quae sunt Actor ó. ita explico, vt semper penes Apostolos ipsos prima totius actionis cogitatio, deliberatio, et adminiitratioruerst, Secuta autem nt Ecciesia eorum consilium. Quanquam etiam in eo loco aliquid speciale fuisse videri potest, nimirum vt toti multitudini prima illa suffragiorum ferendorum potestas deferretur, propter murmur Graecorum, qui conquerebantur Hebraeos sibi in Ecclesiae muneribus praeferri.  \pend\pstart Quarta, Ex ipsa loquendi ratione, qua vtitur scriptura Actor. 15. vers.22 16.vers.4.1. Corinth. 16. vers.3. ὃς ἀν δοκεμάσητε Ex qua phrasi luce cla\pend
\section*{AD I. PAVL. AD TIM. }
\marginpar{[ p.513 ]}\pstart rius intelligitur a populo et vniuersa Ecclesia probari tantùm, vel improbari, quae Presbyteri et ipn preponti Eceienae cie Ecclesiae esse iudicarant atque proposuerant.  \pend\pstart Quinta, Exempla posterioris Ecclesiae, quae tum canonibus suis sanxit, tum etiam vsu ipso et praxi demonstrauit a plebe tantùm assensum requiri, vt quae a Presbyterio gesta sunt, iusta aliqua ratione plebs vel rata haberet, vel irrita faceret. Laodice. Synod. can.13. qui est relatus cap. Non eit perittenaum uistinet. 63. Sic Cyprianus electiones Ecclesiasticas celebrauit liber 1.epistola 3 et 4. Sic suo quoque seculo factum esse scripsit Socrates liber 2, Histor. cap.6.  \pend\pstart Sed si erit summus Magistratus fidelis, in cuius territorio aliquis ad dignitatem Ecclesiasticam est asciscendus, debet certe specialis istius Magistratus consensus praeter eum, qui a toto populo praestatur, expectari et accedere. Id quod ex Socrate liber 5.cap.2. Theod.libus 5.cap.9. can. Reatina dist.63, et Bernar. epist.282. Platina in Bonif. 3 facile colligi et confirmari potest. Nam ad eum Ipluiiii Raglitiutam pertmiet, il Tresbyterium sit nimium negligens in substituendo aliquo in locum vacantis et mortui vel depositi pastoris, illud cogere et edictis suis increpare, vt apparet ex Hermio liber 1.cap.13.Euagr.libus 4.cap.37. Neque aliud (si recte intelligantur) in can. Porro, et cano. seq. distinct.63. constitutum est.  \pend\pstart Accedere quoque debet ipsius electi consensus, antequam populo denuntietur et proponatur et hoc recte in canone Sicut alterius 7. quaest. I. bonstitutum est, Ne inuitus quisqua m'et omnino  \pend
\section*{CAPVT . Ve. }
\marginpar{[ p.355 ]}\pstart repugnans praeponatur, praeficiaturque Ecclesiae. Quou oino perleurolam ent neqde fieri vllo modo debet.  \pend\pstart Trinundinum autem siue spatium trium hebdomadarum populo siue plebi concedi solitum fuisse, intra quod de biecto u rrte prerio viro et illi oblato atque etiam audito populus iudicaret, apparet ex Socrate Scholast.libus  5.cap.5.  \pend\pstart Ipsius autem populi assensum χεροτονία, id est, manuum eleuatione praestari et fieri solitum fuisse apud Graecos ex more gentis illius certissimum est, idque constat tum ex profanis:tum etiam ex Ecclesiasticis scriptoribus infinitis. Itaque  et in Actis Apostolorum: et saepe apud patres electiones ipsae χερροτονίας voce nominantur. Videtur vero etiam haec consensus praebendi et testificandi ratio vsitata fuisse inter Hebraeos ex Nehem. 8.v.7 Apud nos Gallos in omnibus fere conuentibus pablieis comsemus noiter lolet modeltissime solo silentio significari, post propositam negotii, de quo agitur, consultationem: modestiae, inquam, causa. Id quod in electionibus Ecclesiasticis hodie quoque in reformatis Galliae Ecclesiis obseruamus.  \pend\pstart Negant autem quidam obstinatissime, liberum populo suffragium eligendi relinqui, nisi duo aut tres populo presententur, siue proponantur a Presbyterio, quorum sit libera populi optio. Nam si vnus, isque duntaxat offeratur et praesentetur, ecquae electionis libertas populo restat, aiunt, cùm sit electio de duobus vt minimum vnius optio? Respondeo vero vbi copia atque facultas plurium idoneorum populo offerendorum datur, et sup\pend
\section*{AD I. PAVL. AD TIM. }\pstart petit, non vnum tantùm, me authore, populo proponendum esse:sed duos aut plures. Vbi autem vnus tantùm idoneus esse potest, etiam in eo ipso liberam esse populi electionem deliberationem, et assensionem. dico adeo, vt si quem etiam magis aptum agnouerit populus, aut vnus aliquis e populo, de eo Presbyterium monere et possit et debeat. Atque  haec omnia quae diximus in ordinaria vocatione locum habent. Nam extraordinaria vocatio alia est, quae nulla habet praecepta.  \pend\pstart Secundum caput est. Quando electio sit celebranda. Responsio vero facilis est, nimirum vbi prior pastor vel defunctus est:vel legitime depositus. Nam vel viuente adhuc, vel manente et ofcio fungente priori, secundus eligi non debebit nisi fortasse in subsidium defessi, veluti si nimia sit iam superstitis infirmitas senectus morbus, caecitas, aut aliud quippiam simile, alius aduocari potest, qui illi primo tantum accedit, non autem succedit, vt docet Socra.libus 7. cap.4o. Herun liber  8. cap.26. Bernard.epistola 126. et tota 7. quaest.I. in Decretis responsum est. Quod si prior pastor defunctus erat, illud fere fuit obseruatum, vt non nisi demum post tertium sepulturae diem in locum defuncti de alio eligendo ageret Ecclesia, vt est in canone rranas i ontifice distinctione 79. Viri etiam graues, et Episcopi soarum Ecclesiarum amantissimi saepe sibi Successores ipsi designarunt, etiam viui quod futurum periculum prospicerent in futura post suum obitum electione. Hos autemita ab ipsis designatos, postea populus atque vniuersa Ecclesia comprobabat ant equam legitimi fierent pastores vti Augustinus et Athana\pend
\section*{CAPVT  Ve }
\marginpar{[ p.357 ]}\pstart sius, et alii sibi successores elegerunt. August. epi. 1e nncor lincanbapeupinoi Bureuuibus 7.cap.32. Atque  haec desecundo capite et de tempore electionis  \pend\pstart Tertium caput est, Quomodo sit electio celelebranda. Respond. Sic esse celebrandam, vt et Persona idonea reperiatur, et Ratio animusque eligentium sit syncerus.  \pend\pstart Persona vero cognoscetur idonea si fiat serium examen ac diligens ipsius tum in Doctrina ipsa, et illius explicandae methodo, tum Vita et motibus, sintne in eo tales, quales antea praescripsit Paulus cap.3.  \pend\pstart Imperitia vero psalmodiae et cantus, doctum et probum hominem a munere Episcopi excludere minime vel debet, vel potest, quicquid a Gregorio Magno contra responsum sit, et probatum in cano. Florentinum Archidiaconum distinct. 85. Sed iam Episcopi officium in cantilenas conuersum fuerat, relicto abdicatóque ab illis munere et onere concionandi. Fuit etiam olim diligenter et sancitum, et obseruatum, vt constitutionum Ecclesiasticarum nemo ignarus in Pastorem Ecclesiae ordinaretur, vt tradit Sozom. libus 4. cap. 24. Quae sane conditio non est in futuro Pastore negligenda, si quis talis reperiri potest. Sin minus, eligi tamen poterit, qui non erit omnino ordinis Ecclesiastici ignarus, modo reliquas animi dotes habeat, quae sunt ad pastorale munus necessariae. Neque enim cantus Ecclesiastici neque  istarum regularum Synodicarum ignoratio aditum piis et doctis viris ad pastoris officium excludere debent, quia illa duo vel vsu, vel exerciatione facile postea edisci possunt.  \pend
\section*{AD I. PAVL. AD TIM. }
\marginpar{[ p.339 ]}\pstart Animus autem eligentium erit syncerus, si neque Auaritia (Bracarensi Synodo 3. cap.6.) Fauore, Ambitione, vi et armis et carnali aliquo affectu ad electionem progrediantur, vel impellantur qui ferunt suffragia: sed solam Dei gloriam, et Ecclesiae vocantis aedificationem spectent, sibique  pro scopo vere animo proponant.  \pend\pstart Ieiunium etiam publice indici solitum fuisse, cùm fiebat aliqua Ecclesiastica electio, testatur Bernard.epistola 202. Et certe quatuor illa tempora, quae dicuntur in Papatu (quibus publicum ieiunium toti Ecclesiae indictum est) erant eae anni tempestate, squibus passim et publice in tota Ecclesia electiones Ecclesiasticorum munerum celebrari cum precibus et ieiunio solebant. Apparet ex can. Ordinationes dist. 75. atque  etiam ex eo quod clericis adhuc hodie in Papatu Episcopi ordines quos vocant conferunt illo tempore, quicquid aliter Leo, et Calixtus sentiant et scribant, in canone Ieiunium et cano. Huius obseruantiae distinct.76.  \pend\pstart Post Presbyterii electionem et populi approbationem, olim quidem (vt ex Cypria. liber 2. epistola 1I. apparet) qui electus fuerat publico scripto proposito publice significabatur, vt si quod esset impedimentum, quod obiici posset, intelligeretur. Hoc male praetermissum est hoc tempore. Post talem igitur electionem ordinatio siue confirmatio, siue manuum impositio fieri solebat electo. His enim tot nominibus vna res et eademappellata est a veteribus.  \pend\pstart Est autem ordinatio haec ipsius electae atque approbatae personae in munus illud, ad quod ido\pend
\section*{CAPVT  V. }
\marginpar{[ p.359 ]}\pstart ne a visa est, receptio et missio, quae fit praesente vniuerso populo, cum inuocatione Dei nominis. Itaque tunc in illius muneris professionem mittitur ab vniuersa Ecclesia electus. Haec etiam cum ieiunio et publicis precibus fieri certe debet, appellaturque fere in scriptura, Impositio manuum veluti hoc loco Actor.6.versó. supra 3.vers.14.2. Timoth.1.vers. 6. a caeremonia quae in Ecclesia obseruabatur. Ergo electus Presbyter, vel Pastor vel Episcopus in possessionem mittendus est, ordinandus vel confirmandus, vt loquuntur, quod vt fieret more vsitato tempore Apostolorum manus imponebantur Electo atque approbato can. Presbyter distinct.23. et hoc loco. Fuit enim haec caeremonia ex veteribus et legalibus ritibus ad tempus retenta, vel potius ad imitationem impositionis manuum, quae fiebat super hostiam sacrificandam, recepta in Christiana Ecclesia, vt ap paret ex Leuitic.1. et 8.Iacob benedicturus filios Iosephi, manus illis imponit. vt est Genes. 48.v. 14. Christus etiam ipse paruulos benedicturus manum illis imponit Matth.19. vers.15. Atque etiam illa in Baptismo locum habuit, vti apparet Hebr.6.vers.2. In aegrotis etiam et dono miraculorum Mar.16.vers.18.  \pend\pstart Cur autem ista manuum impositio fieret duplex ratio est.  \pend\pstart I Vt intelligeret is, cui manus imponebantur, se totum Deo asseri et destinari. Denique in se a Deo manus iniici, vt iam abrenuntians rebus omnibus suis et mundi negotiis totus muneri suo ad Dei gloriam incumberet at que inseruiret. Vnde preces adiungebantur, vt eum Dominus ipse  \pend
\section*{360 AD I. PAVL. AD TIM. }
\marginpar{[ p.]}\pstart 2 Ratio est, vt Spiritus sancti dona in electum effunderentur, eumque quatenus, quaque mensura Dominus ipse videbat conducere, impleret. Quod ipsum etiam tota Ecclesia pro electo summis votis a Deo petebat. Vnde ab August. manuum impositio nihil aliud esse definitur, quam oratio super hominem. de Baptismo liber 4.  \pend\pstart Haec tamen caeremonia, cùm sit tantùm ritus quidam, non vsque adeo certe necessaria est, modo preces adhibeantur in confirmando electo, et Ecclesia vniuersa vota sua eum eo coniungat, vt spiritu Dei corroboretur. Id verum esse apparet, quod Paulus ipse non semper hac voce et caeremonia vtitur, cùm de electione Episcoporum agit. vti Tit.1. vers.5.  \pend\pstart Deinde quod etiam ipsis Papistis et canonibus testibus, non requiritur haec impositio manuum in omnibus ordinibus et dignitatibus Ecclesiasticis veluti in acolytho, et Subdiaconis ordinandis, vt expresse refertur can. Subdiaconus et cano. seq. distinct. 24. quicquid ibi glossa glosset, et distinguat.  \pend\pstart Tertio, quod etiam ipsis Papistis testibus, vis tota huius impositionis manuum a precibus pendeat, non ab ipso ritu et iniectione manuum in ordinandos, adeo vt sit tantùm illa caeremonia signum quoddam externum commendationis cuiusdam, quae a tota Ecclesia fit Deo, illius personae, quae iam electa est, et quae in possessionem sui muneris est mittenda Vnde praeter missa precatione nihil omnino valet et efficit ipsa manuum impositio, etiam ab Episcopo facta vt est in cano  \pend
\section*{CAPVT  V. }
\marginpar{[ p.361 ]}\pstart Quorundam. Dist. 24. Ergo cùm preces retinentur, ipsis canonibus Papisticis autoribus ipsa manuum impositionis vis, finis, et substantia conseruatur. Quod est satis. Neque est manuum impositio essentialis pars et ritus legitimae vocationis  \pend\pstart Denique dum hanc caeremoniam praecise vrgent homines, videmus eos incidere in vanas quaestiones et ineptas, veluti De triplici manuum im positione Vna Ordinatoria, altera Confirmatoria tertia Curatoria ve eit in canone. Manus. I. quaest.I.  \pend\pstart Item vtrum, repeti in eodem possit manuum impositio.  \pend\pstart Item vtrum sic manus imponi debeant, vt ca put ipsum electi attingant, de quo anxie quaeritur in cano. Episcopus dist 23. Ergo fatui sunt, qui haerent in ea nuda caeremonia, et proram et puppimlegitimae vocatloms cecienaitieae constituunt quod faciunt Papistae, cùm illa sit rerum indifferem num numeromeque sit pars vlla essentialis legitimae vocationis: sed tantùm illius testificandae adminiculum externum et signum quoddam confirmatae electionis.  \pend\pstart Ergo et retineri et omitti potest pro more regionis, in qua electus ordinatur.  \pend\pstart Additum est denique, vt esset legitima ordina tio episcopi, nimirum vt tres episcopi vicini ordinationi siue confirmationi interessent siue ma nus vna cum vniuersa ecclesia, de cuius pastore eligendo agitur, imponerent. Id quod ex Apostolorum canonibus sumptum putant, et studiose retinent Papistae, quod tamen et onerosum est pau peribus ecclesiis, et ad ambitionem et pompam  \pend
\section*{302 AD I. PAVL. AD TIM. }\pstart sit ab illis videturque inuentum. Quanquam tamen operaepretium est ad conseruandam fidei concordiam, et vnionem tum doctrinae tum charitatis electum episcopum quaerere et habere manus associationem a vicinis episcopis et pastoribus: quam se ab aliis Apostolis accepisse testatur Paulus Galat.2.vers. 9. et concilio Mileuit cano. 15.ordinati in vna ecclesia aliis episcopis vicinis statim significabantur publico scripto misso ab ordinatoribus: sed praecise tres alios episcopos accersere et aduocare est plane ecclesiis graue et onerosum, tamen Graecanicae ecclesiae studiose, vti fastum omnem externum, et alia pompernalia retinuerunt, quemadmodum apparet ex Historia Euagrii liber 2.cap. 8.Socrat.libus 5.cap.7.libus 4.cap. 3.cano. episcopus dist.23.  \pend\pstart Item constitutum erat et obseruatum in ista ordinationevt episcopus non nisi in vrbe metropolitana confirmaretur, vt ex concilio Toletano apparet et est in cano. Qui in aliquo Dist.62.  \pend\pstart Denique nisi electus curaret sese intra quinque menses consecrandum ab Archiepiscopo suo amitteret ius sibi quaesitum cano. Quoniam Dist. 100. Sed haec omnia cuiusque ecclesiae iudicio funt relinquenda, quando quis sit ordinandus, et eligendus. Nam velle praecise vt ab Archiepiscopo quis confirmetur et ordinetur, vtque in vrbe metropolitana, id plane redolet veterem tyrannidem Antichristi, quae prorsus est in ecclesia Christiana abolenda: et quilibet pastor melius inter gregem sibi destinandunconfirmatur. Demum additae sunt quaedam vestes a communibus diuer sae, et ornamenta, vt augustior haberetur inaugu\pend
\section*{CAPVT  V. }
\marginpar{[ p.363 ]}\pstart ratio, Item oleum chirothecae et similes quaedam nugae plane superstitiosae et pueriles. Quae sunt omnia reiicienda. Nam dum formosam ecclesiam illiusque dignitates speciosas repraesentare nobis conantur, meretricios illam cultus inducunt: et ex casta et Christi sponsa verum scortum et foetidum effecerunt. Postremo circa annum 7ooa Christo passo in Italicis et occidentalibus ecclesiis reliqui episcopi ceperunt confirmationem suae digni tatis petere a Roma. episcopo. Vide quae Isych. scribit. liber 6.in cap.21. et 22.Leuitic. Ex his autem omnibus apparet quam nulla sit, vel non legitima eorum verbi Dei Ministrorum, vel ecclesiae pastorum Vocatio, qui solius regis, vel reginae: vel patroni, vel episcopi, vel Archiepiscopi autoritate, diplomate bullis, iussu, et iudicio fit vel eligitur. Id quod dolendum est adhuc fieri in iis ecclesiis, quae tamen purum Dei verbum habent et lequuntur, veluti in media Anglia. Nam Anglos homines alioqui sapientissimos, acutissimos, et pientissimos in istis tamen Papisticae idololatriae et tyrannidis reliquiis agnoscendis et tollen dis, scientes. prudentesque caecutire, mirum est. Itaque praeclare sentiunt, qui omnem illam chartula riam et episcopaticam curionum, et pastorum ec clesiae creandorum rationem item ex solo episcopi con sensu et diplomate ministrorum verbi caelestis vo çationem, approbationem, et inaugurationem damnant, tollendamque  ex reformata ad Dei verbum ecclesia censent: quod ordo Dei verbo praescriptus in ordinatione huiusmodi personarum sit praetermissus, ac violatus sicut i perspicue appareat. Denique et Senatui ecclesiastico et Populo  \pend
\section*{304 AD I PAVL. AD TIM. }\pstart Christiano ius omne suum atque suffragium misere sit hac ratione et in hoc genere vocationum ecclesiasticarum ademptum, et in vnum quendam episcopum magna tyrannide atque abusu translatum. Dominus Deus talibus corruptionibus, quae adhuc in ecclesiis ipsius supersunt et deffenduntur, mederi magna sua misericordia dignetur et velit, quae tandem certe magnam ecclesiae Dei ruinam secum trahent, et ipsum sacrosanctum verbi Dei ministerium reddent efficientque vel mer cenarium: vel omnino contemptibile et abiectum. Quod Dominus auertat.  \pend\pstart Vna modo quaestio superest iil hoc argumento, Nimirum, quid de eorum administratione sit sentiendum, qui non sunt legitime vocati: atque vtrum iis fana conscsentia adhaerere debeamus. In quo certe distinguendum est. Aut enim est Omnino illegitima vocatio eius, qui ecclesiam aliquam administrat: vel Non omnino illegitima, sed tantum ex parte.  \pend\pstart Illegitima omnino est ea, in qua omnia Dei verbo praescripta, sunt praetermissa vel violata, vt si quis seipsum intruserit priuata autoritate. Non omnino sed tantùm ex parte illegitima est ea, in qua tantùm quaedam ex illis non obseruata fueTunbieu praetermissa. Ergo cuius omnino illegiti ma vocatio est, is certe pro priuato habendus est, ac huiusmodi vocatio pro nulla. Itaque neque sacramenta conferre potest, neque reliqua negotia ecclesiae gerere: et quae gessit pro nullis habentur, nisi fortasse esset extraordinaria illius vocatio, id est, quae signis testimoniisque certis a Deo confirmata esset. Sola autem propa\pend
\section*{CAPVT .V. }
\marginpar{[ p.365 ]}\pstart gati euangelii, et fructus, qui inde multus emergit, con sideratio non confirmat huiusmodiomnino miegiemna vocationem, vtine que loetus enatus illegitimam copulam, quae intercessit inter adulteros.  \pend\pstart Vitium autem istud tolli potest, si postea ordi ne et legitime vocetur is, qui primum illegitime in ministerio versabatur.  \pend\pstart Qui autem ex parte tantùm est illegitime electus et vocatus, veluti: si per Simoniam quis munus adeptus sit, et non legitimo populi vel senatus ecclesiastici consensu, consensu tamen gradum obtinet, non est habenda pro omnino nulla huius modi vocatio, sed vitium illud est corrigendum Itaque isti sacramenta conferre possunt, quia vtcum que non omnino veram et satis legitimam vocatio nem habeant: habent tamen aliquam. Sic Scribae et Pharisaei in Cathedra Mosis sedebant. Sic Caiaphas in pontificatu summo erat, quem pretio redemerat. Itaque exemplo et Apostolorum, et Chri sti ipsius, qui eos monuit et reprehendit non autem secessionem a corpore ecclesiae fecit, ab iis nos se 9r egure manumedeui non denemus, nec a toto ecclesiae corpere. ( Id enim esset Donatistarum haeresim reuocare) sed pacem ecclesiae, quantum in nobis est, colere nos oportet, etsi istud vitium quod in eorum personis et electione inhaeret, damnare oportet: et quantum possumus tum emendare, tum etiam emendatum velle.  \pend\pstart Hac ratione fit, vt etiam a sacrificulis Papisticis collatus Baptismus non sit in ecclesia Dei repeten dus: a priuato autem collatus idem baptismus sit tamen repetendus, et pro nullo habendus. Etsi enim in ecclesia Dei illegitimam vocationem Pa\pend
\section*{AD I. PAVL. AD TIM. }
\marginpar{[ p.320 ]}\pstart pistici sacrificuli habent: tamen ex consensu populi habent aliquam. Est autem aliud, vt ait Augusti liber 2. contra epistolam Parmeni. cap. 13 aliquid prorsus non habere: aliud autem aliquid perniciose habere: aliud aliquid salubriter habere. Quod tertium solum eligendum nobis est. Sed tamen vti magistratus vitio creatus, Magistratus tamen est propter populi consensum, quemadmodum.lib 5.de lingua latina docet Varro: ita isti sacrifici etsi vitio creati sunt episcopi et presbyteri, Vt, παρελιςλομοι potlus, quain ἐπιςλοuοι sint merito appellandi, tamen quadamtenus presbyteri sunt et Episcopi, praesertim apud eum populum, qui in eo consensit.  \pend\pstart Caeterum obiter monendi sunt lectores, quosdam referre totum hunc Pauli locum, vbi de non imponendis cito manibus agitur, ad reconciliationes publicas lapsorum, quae fieri solitae sunt in ecclesia per impositionem manum quemadm. ex Cypriano liber 1.epist.14. facile colligi potest, et ex ve terum scriptis et canonibus Synodorum. Sed tamen aptius videtur, si referamus hunc locum ad electiones ecclesiasticarum functionum: imprimis autem Pastorum. Illae enim reconciliationes lapsorum veluti idolis sacrificantitum, et libelariorum, qui dicebantur, quae per manuum impositionem fiebant, videntur demum post obitum et tempora Apostolorum notae et vsitatae fuisse ea forma in ecclesia, quam refert Cyprianus: vti et multae aliae ceremoniae post Apostolos in ecclesiam sunt inductae, et obseruatae.  \pend
\phantomsection
\addcontentsline{toc}{subsection}{\textit{23 Ne amplius esto abstemius: sed vino }}
\subsection*{\textit{23 Ne amplius esto abstemius: sed vino }}
\section*{Chla le v. }
\marginpar{[ p.113 ]}\pstart paueulo vtere, propter stomachum tuum et trebras ihas inprmirares.  \pend\pstart Etsi hunc versiculum quidam subiiciunt sequenti ta men nossequemur ordinem quem codices commu nes habent. Est igitur hic locus quoque continuatio siue Α’κολέθηηις superioris argumentum, in quo de Presbyterorum officio agitur continet autem hic versiculus praeceptum, quod ad omnes in vniuersum Presbyteros et praepositos ecclesiae perti nets quucenuoii enemplas pregis don tantum in doctrina: sed maxime in vita esse debent, iuxta dictum Christi. Math.5. Vos estis sal terrae. et Petri in 1.epistola cap.5.  \pend\pstart Quanquam enim hic locus quiddam peculiare, et quod ad Timotheum sigillatum spectat, trade re videtur: vis tamen et ratio praeceptl est generalis, et ad omncs, qui aliquo fungantur munere in eccrenla Denjeatendieur. Psuc aute ordinem in huius versiculi explicatione sequemur, vt primum quae ad specialem huius loci explicationem faciunt, dicamus: deinde praeceptum Pauli ad omnes generaliter producamus ex ipsius mente et scopo.  \pend\pstart Iubet igitur Paulus, vt Timotheus vtatur pauculo vnio, idque propter ipsius valetudinem et stomachi debilitatem, vt suo muneri melius supereesse possit.  \pend\pstart In quo quibusdam videtur nimis abrupta et perturbata orationis huius et styli Pauli ratio, quod a praecepto generali, de electione presbyte rorum, statim transeat, aut potius transuolet ad praecepta medicinae: et ea, quae peculiariter ad  \pend
\section*{AD I. PAVL. AD TIM. }
\marginpar{[ p.368 ]}\pstart valetudinem et stomachum Timothei pertinent. Quae igitur seriei et continuationis huius loci afferri ratio potest? Respondent quidam apte hunc locum cum superiori connecti. Cùm enim superiori proxime versiculo voluerit Paulus Timotheum esse ἀγνὸν, id est, castum, vt interpretantur illi:merito quoque de vini potione praeceptum subiici, quae castitati repugnare videri poterat. Itaque ne castitatis retinendae praetextu nimia et immodica abstinentia ab eo suscipe retur, vtilem hanc admonitionem Paulum subus  iunxisse volunt. Haec Chrysosto:haec Anselmus: haec pene coeterorum omnium sententia est. Nos tamen aliud existimamus. Nam quemadmodum voκάγνὸ ad castitatem corporis siue coelibatum quemadmodum tamen illi fingunt, minime per tinet: sic neque hoc praeceptum additum est propter castitatem. Quid igitur? Cùm oporteat omses iiueles ceciella pauorto ce ministros verbi Dei vitae suae luce et sanctitate praeire reliquo gregi, docet hoc loco Paulus exemplo Timothei quatenus eos victus sobrietatem et rationem ha bere deceat, vt in cibo ipso caeteris quoque vitae exemplum pastores praebeant.  \pend\pstart Timotheum autem vult Paulus vti vino, sed mouieo: ex quo apparet eum fuisse abstemium, non tam quidem natura: quam vitae magna conti nentia. Non vetat enim nude, ne Timoth. aquam bibat ὀλίγω. Haec vox non tam ad qualitatem vini, quam ad quantitatem referenda est, vt ὀλιγον οινον intelligamus, non ὀλιγορόρον. i. quod aquae minimum ferat: sed paucum, et quo quis inebriari non possit. Diserte ver o vocem hanc, ὀλίγον, addidit, vt  \pend
\section*{CAPVT  V. }
\marginpar{[ p.369 ]}\pstart hoc suum praeceptum atque doctrinam Euangelii ab omni calumnia vindicaret. Nam videri poterat ex hac Pauli iussione euangelium incitare, ac impellere homines ad vinum hauriendum, seseque ingurgitandum, nisi hoc tempera mentum atque hanc velut exceptionem subtexuisset Paulus. Id quod sobrii homines inter Ethnicos vituperassent, et merito quidem. Nam et Os.7. vers.5. et Prouerbus 31. quanta, quamque grauia Dei iudicia in ὀινοπότας id est, vinolentos et insobrios homines denuntientur, apparet. Denique Luc.21. vers. 34. Christus et crapulam et ebrietatem damnat. Quare merito moderationem at quesobrietatem hac voce ὀλίγον praescripsit Paulus, etiam Timotheo sua iam sponte sobrio et temperanti, nedum vt praeceptum de sobrietate seruanda cuiquam graue aut intempestiuum videri possit.  \pend\pstart Cur autem praescripserit vsum vini Timotheo, ratio est, non quod medicum agat Paulus: sed quod in suae ipsius valetudinis detrimentum et virium corporis debilitationem nimiam abstemius esset, idque sponte, hac sibi lege imposita, ad domandos carnis impetus et libidinosos motus, qui in iuuenili illa aetate magni et feruentes esse potuerunt, nisi hac tanquam frigida suffusa, id est, vini detractione, temperarentur et extinguerentur. Itaque Timotheus habita tum suae Ætatis, tum Vocationis ratione, vino sponte abstinebat: non superstitione aliqua, quasi vini vsu Deus offenderetur: non haeresi, quasi hanc Dei crea turam damnarer, quod postea fecerunt Aquarii et Seueriani: Non voto aliquo suscepto, quale fuit veterum Nazaraeorum votum. Non denique  \pend
\section*{AD I. PAVL. AD TIM. }\pstart odio et auersatione ipsius gustus et odoris, quod faciunt ii, qui natura sunt abstemii, et vinum gustare non possunt: sed certo animi proposito et instituto, vt carnis libidinem in tam feruente aetate, tanquam aestate, reprimeret, atque compes ceret: vti qui equis, ne nimium ferociant, pabuli partem adimunt, ne nimium impinguentur.  \pend\pstart Ergo qui superstitione aliqua ducti, qui haeresi alicui addicti, qui voto constricti, vino abstinent hoc Timothei exemplo iuuari vel deffendi non possunt, cuius factum hic a Paulo commemoratum est ab illis dissimillimum.  \pend\pstart Denique ratio addita est a Paulo, et modus, ex quo Timotheum aestimare vult, quatenus et vino et caeteris cibis abstinere debeat. Nimirum quatenus ratio valetudinis et corporis sanitas patitur id fieri: non vt longaeui tantum fiamus: sed vt commodiùs et faciliùs Deo inseruiamus. At que fiae quidem partieularieer pertinent ad Timotheum: Videamus iam quae ex his colligi debeant, et ad omnes Presbyteros et Ecclesiae praepontos diffundantur. Duo vero canones ex hoc versiculo a nobis excerpentur, iique generales et vtiles.  \pend\pstart I Presbyteri eam victus rationem ac sobrietatem instituunto, quae et ipsorum muneri facile obeundo conducat, et caeteris ad exemplum tem perantiae esse et praelucere possit.  \pend\pstart 2 Presbyteri atque praepositi Ecclesiae in ea vi ctus sui sobrietate ac ratione instituenda, et bonae valetudinis suae, et virium stomachi atque totius corporis sui rationem habeuto.  \pend\pstart Atque hic vterque canon est exp licandus.  \pend
\section*{CAPVT  V. }
\marginpar{[ p.371 ]}\pstart Ac primùm quidem primus.  \pend\pstart Iam inde a primis Apostolorum temporibus propter et Muneris eminentiam, et functionum ipsarum Ecclesiasticarum finem receptum est, vt maior quaedam sanctitas a praepositis Ecclesiae requireretur, quam a reliquo populo Christiano. Qui enim in eo gradu collocantur, vitae exemplo caeteris praelucere debent, quemadmodum docet Petrus. 1.Epist.5.vers.3. Vnde a Christo appellantur et Lux et sal terrae, in quo caeteri saliuntur. Math.5. Neque modo in iis, quae ad cultum Dei, veramque vitae reformationem et sanctimoniam pertinent, quale est vt scortatione abstineant, fur to, caede, idololatria, haec seuerior vitae ratio in praepositis Ecclesiae requirebatur: sed etiam in re bus ἀδιαφόροις et mediis: et quae etiam licent reliquis Christianis, nempe propter muneris ipsius Ecclesiastici dignitatem et sanctitatem, maioremque commendationem. Vnde Paulus 1. Corinth. 9. vers. 27. ait se contundere corpus suum, et subiicere, ne cum aliis praedicauerit, ipse reprobus fiat, cedatque illiin exitium, quod ab aliis requirit. Et quae secuta sunt Apostolos tempora austeriorem, rigidiorem, seueriorem que semper praepositis Ecclesiae, quam reliquo populo, disciplinam indixerunt, quemadmo dum et ex canonibus Apostolorum, qui dicuntur et decretis synodorum apparet, maximeque ex eo, quod Thelesphorus Episcopus Roma. ieiunium Quadragesimale primùm solis clericis indixerat, a quibus seruabatur, non etiam a reliquo Ecclesiae populo, vt apparet etiam ex ipso decre\pend
\section*{AD I. PAVL. AD TIM. }
\marginpar{[ p.372 ]}\pstart Ratio huius seuerioris regulae illis imponendae fuit duplex.  \pend\pstart Prima, quod debent propter doctrinae (quam ardentius et purius profitentur quam reliqui) confirmationem et commendationem, quemadmodum diximus, in omnibus quae non modo ho nesta sunt, sed speciem honesti habent, praelucere reliquo gregi. Itaque non modo omnem crapulam, luxum, lasciuiam fugere debent praepositi Ecclesiae, sed vel solam illorum vitiorum vmbram et speciem, vt alios suo exemplo aedificent, caelestemque doctrinam commendent. Nam sic vitia aliorum liberius reprehendere possunt. Est enim haec et verissima et maxime notabilis Hieronymi sententia, An non confusio et ignominia est Iesum Christum crucifixum et pauperem, et esuFieniten turuie praedicare corporibus, et ieiuniorum doctrinam rubentes buccas, tumentiaqueora proferre? Si in Apostolorum loco sumus, non solum sermonem eorum imitemur, sed etiam conuersationem et abstinentiam. Haec Hieronymus in Mich.  \pend\pstart Secunda, quod fere homines, quale est humanum ingenium, iis rebus externis plurimùm tribuunt a quibus si quis abstinet, is et Dei timentior ; ce fanctior a vulgo iudicatur. Coloss. 2. in fin.  \pend\pstart Tertia quoque ratio addi potest. Quod etiam (quae est nostra vitiositas et corruptela) in vsu re rum indifferentium nimium libero, minûsque seuero caro fit lasc iuior, et intemperantior, quam tamen contundere necesse est. Id quod fieri nisi rerum earum, etiam concessarum, detractione sae  \pend
\section*{CAPVT  V. }
\marginpar{[ p.373 ]}\pstart pe non potest. Adeo, vt ei, qui vere temperanter, et sobrie viuere velit, carnisque obscoenos motus comprimere, omnino sit necesse multis rebus ad tempus, vel in perpetuum abstinere, quae tamen licent: sed earum vsu carebit ipse sponte, non quidem propter conscientiam, quasi vetita sint edi, sed quod hoc fraeno carnis libido et rebellio domatur, et rationi subiicitur. Partiar autem totam hanc disputationem in tria capita, nempe vt agam de fine huius abstinentiae, De Cibis: et De Modo illius.  \pend\pstart Finis autem, propter quem quis iis cibis et re bus abstinet, varius esse potest, vnus tamen, isque solus probandus est, quem secutus est Timotheus, vt ex hoc loco apparet, nempe vt Deo liberiùs inseruiamus.  \pend\pstart Ac primùm quidem Ethnici ipsi et gentiles quibusdam cibis abstinuerunt propter superstitiosas opiniones, et plane in Deum blasphemas. Ægyptii porris, caulibus, cepis et aliis oleribus id genus seuere abstinuerunt, quemadmodum M. Tullius, Diodorus Siculus, Iuuenalis, aliique plures authores tradiderunt. Pythagorici philoso phi carnibus abstinuerunt. Philistinis, qui et Palae stini, piscium genus nullû edere concedebatur, et quidem hi omnes, quod eos cibos quibus abstinebant, Deos esse crederent, aut ita diis consecra tos at que sanctos, vt in hominum mortalium vsu et commercio esse minime deberent. Hic finis prorsus et blasphemus est in Deum, qui e creatu ris creatorem facit: et proculdubio damnandus.  \pend\pstart Deinde Haeretici successerunt, qui quibusdam cibis abstinent, sed non ea ratione, propter quam  \pend
\section*{AD I. PAVL. AD TIM. }
\marginpar{[ p.374 ]}\pstart Timotheus, et pii debent. Nam Seueriani et Aquei vino abstinuerunt, quod opus diaboli esse contenderent. Manichaei carne: alii aliis cibis, qu eos esse nefandas et diabolicas creaturas crederent, vt in liber  de Haeresibus docuimus. Hic finis damnandus quoque est, et impugnat dictum Pau Il sup.cap.4.verl.4.omnis cibus creatura Dei est, et bonum quiddam. Denique etiam cano. Si quis Presbyter Dist.3o.improbatur.  \pend\pstart Tertii etiam et ipsi peccant in fine, propter quem cibis abstinendum est. Etsi enim omnes ci bos a Deo creatos, productósque fatentur:ita tamen eos distinguunt, vt ex iis alios aliis sanctiores et ad vitam aeternam comparandam aptiores existiment, abusi eo Dei consilio, per quod in veteri lege quosdam cibos interdixit suo populo, vt est Deut.14. Leuit.11. Tales fuerunt Tatiani, Encratitae, et eorum discipuli:tales hodie Monachi. Et hic quoque finis est improbus, et transfert ad terrena mundi elementa quod vnius Dei Spiritus est proprium, nimirum sanctificare nostras conscientias. Itaque damnandus est iuxta illud Pauli sup.4. vers.4.cibus quilibet a Deo conditus in vsum hominis minime reiiciendus est, et Roma 4.vers. 17. Non est regnum Dei esca et potus, sed iustitia et pax, et gaudium per spiritum sanctum et I.Corinth. 8.vers.8. esca nos non commem dat Deo. August. in liber  quaest. Euange. nec in abstinendo, ait, nec in manducando est iustitia. Et ne quis sibi nimium in ea ciborum abstinentia placeat, haeretici etiam plerique ea excelluerunt, quem admodum Socrates liber 7.cap.17. tradit.  \pend\pstart Alii voto, eóque perpetuo, hanc abstinentiam  \pend
\section*{CAPVT V. }
\marginpar{[ p.375 ]}\pstart ciborum sibi indicunt propter illud Pauli in eadem epistad Roma.14. vers.21. Bonum est non ves ci carnibus: neque bibere vinum: et illud etiam quod est 1.Corinth.8.vers.13 si esca facit, vt offendat frater meus, non vescar carnibus in aerernum. Vtróque  modo, hoc nimirum, et superiore in Christiana Ecclesia grauiter peccatum est, et pene primis ecclesiae temporibus, dum inconsiderato quodam sanctitatis zelo homines ducuntur. Atque hic finis proxime a nobis commemoratus ex eo satis refutatur, quod Pauli dictum conditio naliter et κατὰ τι prolatum tanquam simpliciter sit dictum et effatum isti accipiunt. Docet enim Paul. quibusdam cibis abstinendum tunc, cùm proximi aedificatio, et iusta ratio id exigit: non autem ex voto id fieri vult, neque simpliciter et omnino. Verus igitur et legitimus huiusmodi ciborum abstinentiae finis, isque probandus, est ille, quem scriptura praescribit. Est autem is duplex vel Carnis nostrae extraordinaria mortificatjo, vel Infirmorum fratrum ratio et aedificatio, qui ad tempus sunt toletandi a nobis.  \pend\pstart Primus finis commemoratur a Paulo 1.Corin. 9.vers.27. contundo corpus meum, et in seruitutem redigo. Item Coloss2.vers.23. Sic ieiunium institutum et vsum habuit in Dei Ecclesia vel a tota, vel ab vno aliquo priuatim susceptum. Atque  certe ibi recte obseruat Chrysostomus dici a Pau lo Castigo: et in seruitutem redigo:non dixit, peredo et Punio. Non enim inimica caro est, sed casti go, et redigo in seruitutem, quod Domini est, non hostis: magistri, non inimici: exercitatoris, non adu ersarii Bernad sermo. 66.in Cantico. docet.  \pend
\section*{AD I. PAVL. AD TIM. }
\marginpar{[ p.310. ]}\pstart Nos cibis posse abstinere triplici de causa vel ex medicorum praeceptis, quod non reprehenditur, si tamen nimia non sit cura corporis: vel ex regu lis abstinentiae, vt mortificetur caro, quod etiam si sine superstitione fiat laudandum est: vel ex insania haereticorum, quod est omnino damnandum. Et in epistola ad Coloss. idem Chrysosto. Deus, ait, corpori honorem suum contribuit, ipsi vero (de nimium abstinentibus loquens) non cum honore vtuntur corpore. Denique concilio Gangrensi, quod vicinum fuit primo Niceno, nimia ista abstinentia, et si quavis in ea ponatur est grauiter reprehensa. Quod vero scribit Gregor. Mag. pios a licitis rebus fercubitmere, vel vt sibi meculpas vitae anteactae deleant, sitque haec abstinentia pro poena, est falsissimum, etsi probatur cano. Sunt qui 27.quaest.2.  \pend\pstart Secundus finis expressus est saepe a Paulo in epistolis, praesertim in ea, quae est ad Romanos cap.14, et 15. 1. ad Corinthios 9. atque etiam in Actor. ca.15. a Synodo Hierosolymitana. Propter aedisicationem enim infirmorum, qui non poterant videri peccare malitia, sed nuda tantùm ignorantia, Paulus se omnia omnibus factum esse dicit: et iubent Apostoli collecti quibusdam cibis Christianos abstinere ob eandem causam.  \pend\pstart Alterum caput huius disputationis est de rebus, quibus abstinere solerent veteres, et nunc debeamus ipsi abstinere. In quo nullum est praeceptum diserte ab ipso Dei verbo positum: sed debet quisque iis cibis abstinere, quibus libidinem in sese maxime vel accendi, vel excitari sen\pend
\section*{CAPVT . V. }
\marginpar{[ p.377 ]}\pstart tit et exper itur. Itaque qui carnis esu se proritari vident, carne abstinere debent. Qui vero piscium esitatione se lasciuiores fieri sentiunt (est enim cibi genus delicatissimum, vt testatur in Sympossac. Plutarchus piscium esus) piscibus abstinebunt. Qui vino se libidinosiores effici comperient, vinum ne bibant, quia quo maxime caro nostra proritatur, illud est subtrahendum, quemadmodum ferocienti equo nimium pabulum (quenadmodum in liber  de Arte equestri tradit Xenophon) est subducendum, vt contundatur, et dometur facilius. Qui vero esu leguminum se ad peccandum magis incitari sentiunt, leguminibus potius quam carne debent abstinere: quia fontes et incentiuum peccati est tollendum, non illud, quod libidinem in nobis non accendit. Irem quorum ciborum esu maxime offenduntur infirmi adhuc fratres, quibuscum versamur, illis est abstinendum, non est autem par ratio in aliis cibis Fieret enim temere et sine causa. Neque vero hic nobis vel Nazaraeorum votum, de quo est Numeror. 6. quisquam in vsum et praeceptum reuocet. Fuit enim illud caeremoniale, et iam per Christum implementum atque finem habuit. Itaque desiit esse in vsu. Neque etiam quisquam regerat, reponatque vetus illud discrimen ciborum, quod in lege sua caeremoniali Dominus instituit Deuter. 14. Leuitic. 11. quia illud eriam fuit plane caeremoniale. Vnde Apostoli ipsi de vsu ciborum interrogati minime illud inter homines retinuerunt in concilio Hierosolymitano Actor. 15. Quin etiam illud ciborum discrimen saepe Paulus oppugnat, imprimis autem in Epi\pend
\section*{AD I. PAVL. AD TIM. }
\marginpar{[ p.370 ]}\pstart stola ad Galatas. Denique neque etiam aliquis nobis obtrudat illud, quod est Leuitic. 1o.vers.9. quia et not ipluin plane antiquatum est a Christo, vti et altare et Sacerdotium.  \pend\pstart Fuerunt tamen certa quaedam cibortim genera, quibus maxime ἐγκρατέυοντες et μονάζοντες Chri stiani abstinere solebant, partim ex nimia iam superstitione, partim ex inepta virorum piorum imitatione. Ac quidem imitatione, quod antiquissimi Christiani propter persecutiones in deserta ac vastas solitudines fugere coacti, saepe herbis tantùm, aqua pura, secundo pane, et siliquis, similibusque rebus illic vixerunt, cùm neque caro, neque vman, neque alii cior propter incultum locum suppeterent illis, neque illic reperirentur. Horum vitam imitari se putauerunt, qui tam vilibus cibis postea, etiam sine necessitate, et in mediis vrbibus victitarunt.  \pend\pstart Superstitione vero quidam impellebantur, quod iam plus satis externae isti vitae austeritati, et curiositati quadam verae sanctitatis ignoratione homines tribuebant. Ergo qui se reliquis puriores atque mundiores esse inter Christianos volebant, vt ex Eusebio liber 2.cap.17 liber  6. cap.3. Socrate liber  4. cap.23. aliisque veterum locis colligitur, primum fere astinuerunt his cibis (quanquam tamen, vt ait Paulus Tit. I.omnia sunt munda mundis) abstinuerunt, inquam, Carne, Oleo, Vino: Oleribus interim, et reliquis rebus indiscriminatim vescentes, puram etiam aquam potantes, non conditam, non mistam, non saccaro, non herbis vim illius augentibus, aut vsum gustumque  gratiorem facientibus temperatam. Post\pend
\section*{CAPVT  V. }
\marginpar{[ p.879 ]}\pstart ea creuit superstitio, et haec initia latius postea sunt propagata, adeo vt postea, et iam tempore Gangrensis Synodi, quae fuit sub Constantino Magno, Sagimen omne et pinguedo, et non tantùm caro ipsa Monachis et abstinentibus interdice retur. Demum accessit decretum Gregorii Magni, quod abstinere cupientibus prohibet et interdicit esu omnis rei, quae sementinam carnis originem habeat, adeo, vt iam non solum caro ipsa sit prohibita: sed et caseus, et ouum, et lac, et butyrum, quia sementinam carnis, id est, a carne originem habent. Hoc iam sequuntur omnes Papistae, et est in canon. Denique distinct.4. tanquam hoc sit perfectissimum abstinentiae genus. Bernardus tamen sermo 66. Cantic. videtur sentire hane sententiam esse ex Manichaeorum haeresi deriuatam. Caeterum, ait, quid sibi vult, quod ita generaliter omne, quod ex coitu generatur, vitatur? Nempe horrent lac, et quicquid ex eo conficitur, vt in eo Bernardo neque cum Gregorio Magno, neque cum superstitiosis nostri temporis Papistis conueniat, qui generaliter omnes cibos istos prohibent. Mirum tamen cùm et votum Nazaraeorum vini interdictionem contineret, et Timotheus vino potius, quam aliis reliquis cibis ad carnem domandam abstinuisse dicatur: et veteres Monachi vino potissimum non vterentur. Praeterea cùm tota distinct. 35. vini incommoda atque pericula varia enumerentur, maximeque Episcopis, et sanctioribus viris cauendum praedicetur. Denique cùm Prouerbus 31. serium aliquid acturis vinum interdicatur, mirum, inquam, est, nullos tamen Monachos etiam abstinentissimos et rigidissimae vitae regulam professos.  \pend
\section*{An I. PAVL. AD TIM. }
\marginpar{[ p.Joo ]}\pstart Vino tamen abstinere voluisse, nisi fortasse Bernaexcipias, qui se vino docet abstinuisse. Atqui tamen Paulus Ephes.5. vers.18. in vno hoc cibi genere lasciuiam et ἀσωτίαν inesse testatur potius, quam in reliquis omnibus cibis. Ex quo apparet Monachismum nostri temporis, addo etiam lesuismum, esse meram hypocrisin, in quo non vera abstinentia praescripta est, sed sola tantùm ostentatio et abstinentiae vmbra, cùm vino potius, quam carne esset illis abstinendum, si vere carnem contundere studebant. Atque haec de cibis.  \pend\pstart Sequitur tertium et vltimum huius disputationis caput, in quo altera regula a nobis in vitae ratione, quae a praepositis Ecclesiae seruanda est, constituta continetur et exponitur. Quaeritur enim De modo suscipiendae istius abstinentiae, et instituendae sobriae victus rationis, quis esse et statui debeat. Breuiter autem respondet Paulus, eum esse modum legitimum, atque laudandum, in quo valetudinis ratio habetur a nobis. Ex quo fit, vt tria colligamus. primum non esse nobis omnino omnibus cibis abstinendum. Neque enim tunc, nisi modo plane miraculoso, viuere possemus. Deinde cousequitur, non semper, aut in perpetuum extraordinaria illa ciborum abstinentia vtendum nobis esse, sed tantùm ad tempus. Nam deficerent vires, et paulatim ita debilitare tur natura, vt nulli negotio superesse possemus. Sobrietas quidem nobis est per totum vitae curriculum seruanda: sed haec quoque non senper in iisdem cibis fugiendis versatur. Praeterea cùm de extraordinaria quaeritur, quae ad domandum magnum carnis impetum adhibetur, ea cer\pend
\section*{CAPVT V. }
\marginpar{[ p.381 ]}\pstart te perpetuo obseruanda et tuenda non est. Vt igitur non sunt idem impetus senilis, atque iuuenilis aetatis:ita neque eadem vitae ratio ab vtrisque est suscipienda, vel vtrisque praescribenda. Mutantur enim haec pro tempore, aetate, locis, personis, vt etiam praeclare scripsit, sensitque Augustinus in quaestion. Euangelii. Non enim, ait, interest omnino quid alimentorum sumas, vt succurras necessitari corporis, dummnodo congruat in generibus alimentorum cum his, cum quibus tibi viuendum est: neque quantum sumas interest multum, cùm videamus aliorum stomachum citius saturari, etc. Deinde ita concludit idem Augustinus, Magis ergo interest, non quid vel quantum alimentorum pro congruentia hominum atque personae suae, et pro suae valetudinis neceslitate quis accipiat: sed quanta facilitate, et seueritate animi iis careat, cùm iis vel oportet, vel necesse est carere, vt illud in animo Christiani compleatur, quod ait Apostolus, Scio et minus habere: scio et abundare. Tertium denique quod ex hoc Pauli dicto colligitur, est huiusmodi. Nihil in iniuriam naturae, et nimiam corporis debilitationem vel instituendum, vel institutum et semel susceptum retinendum nobis esse. Quod Timothei exemplo docemur hoc loco, qui, quemadmodum verisimile est, post hanc Pauli commonefactionem vino vsus est, quo tamen abstinere decreuerat. Sed quia ea abstinentia debilitabantur vires corporis naturales, postea vinum in victum adhibuit. Est enim non modo vita haec donum Dei: itaque a nobis minime spernenda: sed virium naturalium et sanitas et integritas quo\pend
\section*{5on 41i fl A AVne Ab TIM. }\pstart que donum Dei est, quae quantum quidem in nobis est, fouenda, sustentanda, tuendaque est, vt nostro muneri satis esse possimus. Alias enim nostra culpa fit, quod inutiles aut minus idonei euadimus. Praeclarum est exemplum veteris Ecclesiae, quae Alcibiadem quendam nimis austere viuentem monuit, vt ea viuendi ratione abstineret, vinumque biberet. At que ille sapienti Ecclesiae consilio paruit libens, vt refert Eusebus liber 5. Histo. cap.3. Vnde falsa et plane respuenda est illa laus, quam Bernardus serm.6.in Psal. Qui habitat, quibusdam Monachis tribuit, nimirum quod variis corporis exercitiis et abstinentiis contererentur supra vires, supra naturam, supraque consuetudinem. Male quoque idem Bernard. sermo.3o. Cantic. Monachum te fore professus es, non medicum: nec de complexione tibi iudicandum, sed de professione. Verior illa est eiusdem Bern.sententia ser. 1o.in Ps. Qui habitat. Nec sane dixerim vt vel ipsam odio habeas carnem tuam. Dilige eam tanquam tibi datam in adiutorium, et ad aeterne beatitudinis consortium praeparatam: caeterum sic amet anima carnem, vt non ipsa in carnem transiisse putetur. Sed et illa quoque eiusdem sententia sermo. 4. in eundem Psalmum verissima est Dominum voluisse, vt in manna caelitus compluta omnis saporis, omnisque odoris delectatio inesset, vt liberum esse ciborum vsum atque a Deo sibi concessum homines intelligerent. Nam hic vini vsum praescribit Paulus, in quo suauitas voluptasque gustus vna cum vtilitate coniuncta est, vt est Psal.1o4.vers.16.nec ignorauit Bernar.tam seueram vitae rationem merito illo etiam seculo  \pend
\section*{CAPVT V. }
\marginpar{[ p.383 ]}\pstart dici crudelitatem, vt ex eius epistolis apparet, de qua purgare quidem Monasticam vitam nititur: sed non potest, vti neque ad hoc argumentum respondit Thomas in Apologia et defensione Monasticae vitae, quia est ineuitabile aduersus cam telum. Modus tamen is adhiberi debet in cibis, in quo nullum ebrietatis, ac intemperantiae specimen aut vestigium appareat.  \pend\pstart Atque eo pertinet, quod ait Paulus, modieo vino vtere: non enim ad ebrietatem et satietatem sed ad sanitatem sumi et vinum, et reliquos cibos vult Deus cum gratiarum actione.  \pend\pstart Denique  ex hoc ipso loco futilis et inepta prorsus deprehendi porest illa Bernardi sententia, qui epist.321.negat honestati et puritati vitae Monasticae, id est, sanctae conuenire, vti medicamentis et potionibus ad expellendum morbum idoneis. Nimium certe fuit ille seuerus, et praeter aequitatem ac rationem. Æquior Chrysosto. qui et hic ait Timotheum ad medicos remissum, non autem miraculo curatum a Paulo: et in Homil. 14. in hanc epistolam de vita Monastica agens, vluii inouicamrentora monaenimjeum est opus, non adimit, sed concedit: quanquam tamen rigidior eorum censor habitus est Chrys. Sed Satan semel via superstitionibus patefacta, eas in immensum postea adauxit, et in iugum intolerabile tandem aggregauit ad torquendas hominum conscientias.  \pend
\phantomsection
\addcontentsline{toc}{subsection}{\textit{24 Quorundam hominum peccata ante manifesta sunt, praeeuntia damnationi: quosdam vero subsequuntur: }}
\subsection*{\textit{24 Quorundam hominum peccata ante manifesta sunt, praeeuntia damnationi: quosdam vero subsequuntur: }}
\section*{AD I. PAVL. AD TIM. }
\marginpar{[ p.324 ]}
\phantomsection
\addcontentsline{toc}{subsection}{\textit{25 Itidem et bona opera ante manifeSta sunt: et ea quae secus habent, occultari non possunt. }}
\subsection*{\textit{25 Itidem et bona opera ante manifeSta sunt: et ea quae secus habent, occultari non possunt. }}\pstart Diuersa quoque est huius versiculi cum superiore nectendi ratio. Putant enim quidam cum vers.22. coniungendum, vers.autem 23. per parenthesin insertum et legendum quoque. Alii aliud. Nos autem hunc locum tum ωρακολόθησιν, tum etiam νοθέτησιν esse censemus, et idcirco cum superiore argumento recte conuenire atque coniungi posse, quod praecepta duo contineat ad Pastorum et Presbyterorum munus spectantia, ex quibus vniuersa pene totius censurae Ecclesiasticae ratio definiri potest. Id quod maxime ad Presbyterorum officium pertinet. Atque id iam hoc toto capite persequitur Paulus. Sic igitur statuo canonem tradi iam a Paulo, per quem quomodo praepositi Ecclesiae se in censuris Ecclesiasticis gerere aduersus gregem suum debeant, instruantur, id est, quale sit eorum munus erga vniuersum populum sibi commissum et demandatum a Deo. Quae sane tertia pars est muneris pastoralis, vti supra diximus, maxima. Prima enim docet quomodo erga seipsos gerere debeant. Secunda quomodo aduersus collegas suos. Tertia haec, quomodo erga vniuersum Ecclesiae suae coetum. In qua ipsa parte posita est ratio et disputatio de censura Ecclesiastica. Quanquam huius praeceptionis, quae nunc traditur, tantam vim esse tamque late patentem fateor, vt etiam ad superius praeceptum, quod ad electionem atque ordinationem pastorum Ecclesiae pertinet, referri possit  \pend
\section*{ChRVIV. }
\marginpar{[ p.385 ]}\pstart et in eo magnum vsum habeat. Quare hic totus locus hunc canonem vthimnimanee valde necessarium ex nostra interpretatione continebit.  \pend\pstart In censuris Ecclesiasticis praepositi Ecclesiae ne temere progrediuntor: sed Deum ipsum praeeuntem sequuntor, qui suo tempore bonos et malos detegit.  \pend\pstart Malos, antequam a Deo patefacti fuerint, praepouiti ne eiiciunto. Bonos autem, antequam ipsis operib' tales apparuerint, ne speciali honore praepositi decoranto: aut caeteris praeferunto.  \pend\pstart Horum autem omnium canonum ratio est illa quam affert Paulus, quod nimirum multorum peccata saepe latent idque diu: multorum etiam saepe bona opera et candor animi et synceritas diu occulta manet, ita Domino dispensante res humanas e asque administrante.   \pend\pstart Cur autem id accidat magna sapientique Der prouidentia fit. Nam si quicunque sunt improbi et hypocritae in Dei Ecclesia, statim apparerent, neque piorum fides explorari et probari tam certo posset: neque pii ipsi tam officiosi in hoc mundo viuerent, neque etiam tuti essent, sed ab impiis opprimerentur. Facit igitur haec commistio, vt melius, certiûsque qui sunt probi et veri Dei serui manifesti fiant, cùm, tam multis diffluentibus et eiectis, illi firmi tamen maneant. Sic haereses probant piorum fidem, quemadmodum Paulus ait I.Corinth.11.vers.19. Praeterea si pii omnes agnoscerentur, impii vero aperte quoque statim ostenderentur a Deo, nulla spes resipiscentiae relicta illis videretur, omnisque nostrae in delinquentibus conuertendis industriae et charitatis  \pend
\section*{AD I. PAVL. AD TIM. }
\marginpar{[ p.380 ]}\pstart effectus cessaret. Dominus igitur vult aliorum statum et animum nobis occultum esse, vt cum iis et lubentius versemur: et magis sollicite nobis ipsis caueamus. Sed etiam illis ipsis, qui peccarunt, haec dilatio datur a Deo, tanquam occasio maturius resipiscendi, atque ad Deum redeundi.  \pend\pstart Denique Dominus ipse, qui rerum omnium maturitates et tempora nouit, atque etiam constituit videt quando et quae manifestari tum ad ipsius gloriam, tum etiam ad nostram aedificationem oporteat. Neque vult nos omnia illa scire quae scit ipse.  \pend\pstart Huc pertinet quod Ecclesia Dei comparatur a Christo ipso sagenae siue reti, quo boni et mali pisces capiuntur Matth.13.vers.47. Areae, in qua bonum granum et paleae promiscue conduntur. Matth.3.vers.12. Item agro, in quo bonum semen et lolium quoque nascitur Matth.13.vers.25. Denique domui amplae, in qua sunt quaedam vasa honesta, quae dam fordida 2. Iimtoth. 2. vers. 20. Nam quos domi suae adhuc tolerat Dominus ipse, non est nostrum eos expellere vel arcere:neque Deo ipso vel sapientiores, vel iustiores, vel ardentiores sumus ad ipsius gloriam promouendam.  \pend\pstart Hoc etiam ipsum pertinet tum ad Pastorum ipsorum excusationem, tum ad Bonorum consolationem. Ac quidem excusantur pastores, qui quanquam diligentiam, quam potuerunt, summam adhibuerunt, vident et animaduertunt tamen in Ecclesia Dei adhuc latere quosdam improbos et hypocritas, qui paulatim deteguntur, et a Deo ex suis tenebris in lucem pertrahuntur, pa\pend
\section*{CAPVT  V. }
\marginpar{[ p.387 ]}\pstart tefactis eorum sceleribus. Ergo non sunt accusandi fideles et diligentes Ecclesiae praepositi, si ex iis, qui ab ipsis electi sunt et ordinati, aliqui postea ap paruerunt esse nebulones, et pestes Ecclesiae: aut si quosdam in Ecclesia Dei fouerint et retinuerint qui postea detecti sunt improbissimi, et pessimi. Nam ad tempus multorum peccata latent, et tempore patent, vt ait etiam ille Ethnicus. Neque  statim produntur improbi. Bonis etiam viris haec eadem commixtio et occultatio consolationem affert, si quando videntur sperni, et prae aliis negligi, quod cogitant vnius Dei esse hominum animos et synceritatem mentis prodere atque  patefacere, id est, in lucem proferre, cùm voluerit. Interim igitur sua sorte et obscuritate contenti viuunt, neque murmurant, vel in Deum ipsum, vel in homines.  \pend\pstart Quod autem ait Paulus vers.25. Quae secus sunt id est quae praua sunt, non posse occultari, et in perpetuum celari, ita accipiendum est, vt hoc fieti saepissime, et vt plurimùm, intelligamus: nou autem vt semper. Saepe enim multorum et bona opera, et turpissima scelera latent nos et occulta manent perpetuo, quia ita Deo visum est. Sed vt plurimum Dominus detegit improbos: et piorum syncerum animum prodit, vt est in Psal.37. vers. 6. et Psal.112.  \pend\pstart Ex hoc igitur loco colligimus iudicium de moribus: et censuram, per quam peccata hominum Ecclesiastica poena coercentur, pertinere ad Praepositos Ecclesiae, quos hoc loco alloquitur Paulus. Quam sententiam confirmat Christus Matt.16 Quaecunque ligaueritis in terra, erunt ligata et  \pend
\section*{AD I. PAVL. AD TIM. }
\marginpar{[ p.388 ]}\pstart in caelis. Item et illud Bernardi, pastorum est vigilare super gregem propter tria necessaria, nimirum ad disciplinam, ad custodiam, ad preces. In disciplina rigor iustitiae: in custodia spiritus consilii: in precibus affectus compassionis requiritur. Censura vero haec est multiplex, vti supra est demonstratum. Restat autem vt de excommunicatione dicamus, quia de reliquis generibus censurae Ecclesiasticae quid statuendum sit, ex iis, quae hoc ipso capite dicta sunt a nobis, facile potest intelligi. Ergo de excommunicatione dicamus, de qua septem haec faere quaeri solent.  \pend\pstart Ac quidem, vtrum excommunicatio locum in Ecclesia habere debeat, dubitari non potest. Primùm, propter Dei ipsius praeceptum ita fieri iucentis: deinde propter perpetuum ipsius Ecclesiae vsum. Ac praeceptum Dei tum in veteri, tum nouo testimonio extat variis in locis. Iam inde ab orbe condito, et ab ipsis mundi initiis exconmunicationem fuisse ex Dei praecepto constitutam apparet ex Genes.4.v.14. et 16.17.item Iosue 6.vers.18. Exod.1o. vers.8. et aliis locis veteris Testamenti. In nouo etiam extat praeceptum Chri\pend
\section*{CAPVT  V. }
\marginpar{[ p.389 ]}\pstart sti, quod Matt.18. et Ioan.2o.recitant. Ex quibus ipsis locis eam spiritualis censurae et castigationis speciem in vsu fuisse in Ecclesia docetur, vti ex Ioan. cap.9. et 1.Cor.5 et 2.7. et hoc ipso loco. Neque de eo dubitari potuit, nisi duo quaedam argumenta contra afferrentur. Primùm neminem inuitum cogendum esse, vt recte viuat aut sentiat. Nam, vt vulgo dicitur, Inuitum qui seruat idem facit occidenti. Id autem faciunt, qui inobsequentes Deo excommunicant. Huic argumento saepe respondit Augustinus in disputationibus contra Donatistas et negat veram esse hanc sententiam. Saepe enim qui primùm inuitus ab errore vitae vel doctrinae abductus est, postea volens et lubens in recto itinere permansit. Praeterea cùm quis corrigitur inuitus, si hoc ipsum illi non prodest, quia inuitus id facit: aliis tamen prodest, qui metu poenae coercentur, ne ad eadem scelera exemplo perenananeuriDemquoide cauemrucio facit ne in sceleratos homines vllos, quantumuis impios et derestabiles, poenae vllae etiam a Magistratu decernantur, quia ea ratione inuiti resipiscere et modeste viuere coguntur. Itaque haec prima ratio est absurdissima. O  \pend\pstart Alterum argumentum, quod contra vsum excommunicationis affertur, ducitur ex loco Pauli 2 Corinth. 1o. vers.8. Nempe datam esse pastoribus, et Ministris verbi Dei potestatem ad aedificationem, non ad subuersionem et destructionem. At excommunicatio destruit gregem Domini, quem minuit, et eum ipsum qui excommunicatur, perdit, quia exponit Diabole:ergo ea vti pastoribus non licet in Ecclesia. Respond. Aliquid  \pend
\section*{AD I. PAVL. AD TIM. }
\marginpar{[ p.90 ]}\pstart fieri in aedificationem Ecclesiae duobus modis, vel Per se, vel Per consequens. Per se, cùm aliquid in ea plantatur: per consequens, cùm quod illi noxium est, euellitur. Est enim similis agro Ecclesia Dei, vt est apud Matth.13. At in agro siue quid noxiarum herbarum (quae bonas et vtiles suffocant) extirpemus: siue quod vtile semen seramus, agro prosumus. Ergo eodem modo de Ecclesia statuendum, qui ex ea eiicit improbos multum iuuat et prodest Ecclesiae.  \pend\pstart Quod autem obiicitur saepe quosdam ea abusos esse ad priuatam animi sui vindictam: saepe bonos ex Ecclesia Dei depulsos ea ratione, non oppugnat nostram sententiam. Ea enim sunt vitia persouarum male sua potestate vtentium, non autem rei ipsius. Sed distinguunt sic quidam, nimirum, vt velint quibus in locis et Ecclesiis fidelis est Magistratus, in iis nullum excommunicationis vsum esse debere: in quibus autem non est Magistratus fidelis, in iis tantùm debere: quia fidelis Magistratus si fungitur officio, punit scelera, quae contra Dei verbum admittuntur. Nemo autem vnius criminis et peccati ratione duplicem poenam ferre et pati debet. Quod tamen accideret, si idem etiam ab Ecclesia excommunicaretur. Res. Peccatum omne quod contra Dei praecepta fit, esse huiusmodi, vt ex Dei ipsius definitione et corporis et animi poenam mereatur, quia vtriusque partis nostrae creator est Deus. Et cùm dupliciilla poena afficitur peccator, vna tantùm poena afficitur, id est, ea, quae debetur, et indicta est a Deo ipso: etsi in diuersis partibus a nobis vna illa poena sentitur. Itaque Esdrae cap. 1o.  \pend
\section*{CAPVT . V }\pstart vers.8.vtraque poena in rebelles indicitur: ciuilis et Ecclesiastica siue excommunicatio, quia vtriusque  partis nostrae Dominus author est. Est enim peccatum totius suppositi actio, itaque totum suppositum ad poenam ipse exigit.  \pend\pstart Ergo in Dei Ecclesia vsum habere debet excommunicatio.  \pend\pstart Quaeritur autem, quid sit excommunicatio. Ac dubitari non potest, eam esse censurae et coercitionis Ecclesiastieae speciem quandam. Ergo sic eam definiunt nonnulli, vt generaliter esse velint a qualibet societate fidelium hominum separationem: itemque ab aliquo, vel omnibus donis Dei in Ecclesiam effusis interdictionem et publicam hominis scelerati prohibitionem. Verum haecdefinitio generalior est. Vnde sic rectius definiri potest excommunicatio, vt sit vel ab ipso Ecclesiae coetu toto et societate: vel ab aliquo Dei dono, quod quidem omnium Christianorum commune sit, publica et legitima interdictio siue separatio hominis impuri et non poenitentis. Ex quo fit, vt a quibusdam, velutia Scholasticis duplex excommunicatio statuatur, vna nimirum Minor, quae est interdictio tantùm a Sacramentis, vel ab vno eorum, non autem a coetu Ecclesiae: altera Maior, per quam non tantùm a Sacramentis, sed ab ipso toto Ecclesiae coetu quis separatur, quam ipsam etiam diuidunt in Excommunicationem quae appellatur hoc nomine: et Anathema, cuius etiam fit mentio apud Paulum Galat.1.et 1.Corinth.16. et alibi in scriptis Patrum.  \pend\pstart Qui inter Hebraeos ista scrupulosius et diligentius scrutantur, sanctionis legis constituunt ge\pend
\section*{AD I. PAVL. AD TIM. }
\marginpar{[ p.392 ]}\pstart 1 Auersationem, quam appellant uu, quae est mcerdietio au igrenu tempir ad tempus proptercorporis immunditiem.  \pend\pstart 2 Deuotionem, quam vocant uon quae est extremo exitio addictio et respondet nostrae exconmunicationi, per quam quis a coetu Ecclesiae seponstur donec resipiscat. Hoc ipsum aliis verbis dicitur exterminari e populo Dei, item ἀποσυνάvωyoν fieri Ioan.12.vers.42.16.vers.2.  \pend\pstart Adiudicationem, quam dicunt nnou est siHilis ahathemati.  \pend\pstart De secundo genere tantùm loquitur hic Paulus, et nos quoque, quod etiam vtranque excommunicationem continet et maiorem, et minorem.  \pend\pstart Ethnici homines videntur simile excommunicationi quiddam habuisse ex praua imitatione Ecclesiae Dei: quam plane simiarum more faciebant. Quosdam sacris suis interdicebant, quos υαγεῖς, id est, detestabiles appellabant, et ipsam sententiam ἐναγός. Sic Alcibiades factus est, et excommunicatus Athenis, vt Plutarchus docet.  \pend\pstart Ergo quaesitum est secundo loco, a quo excommunicatio fieri debeata Vtrum ab vno tantùm, an a senatu Ecclesiastico, an vero a tota ipsa Ecclesia. Dixit Christus. Matth. 18.vers.17. Dic Ecclesiae. Dicit contra Paulus, Hunc notate, tanquam rem totam ad se deferri velit, sibique ipse authoritatem eam vendicet. Et refertur a Theodorito Cyri Episcopo historia.lib 5. cap.37. de quodam Heremita, qui solus Theodosium imperatorem Constantinopolitan. huius nominis secumdum excommunicauerit. Scholastici dicunt Epi\pend
\section*{CAPVT V. }
\marginpar{[ p.393 ]}\pstart scopos, et eo superiores omnes ordines posse excommunicare: at qui sunt Episcopo inferiores non posse, adeo vt et Presbytero et Diacono idem non concedant quasi ex vnius Episcopi nutu et sententia id fieri debeat. Sic igitur illi censent. Certe ab vno solo, quanquam Episcopo, fieri neque debet, neque potest excommunicatio. Dicit enim Christus, dic Ecclesiae: et Paulus de excommunicatione agens, vult Ecclesiam cum suo spiritu congregatam esse 1.Corinth.5. vers.4. Ratio est, quod differunt regna politica ab Ecclesia, et Ecclelsaitico regimine. Nam regna vnius regis nutu administrantur, at Ecclesia non item. Nec enim pastor est rex et dominus Ecclesiae (vt ait Petrus) 1. Pet.5.vers.3.  \pend\pstart Ait igitur Christus reges et principes dominantur inter suos subditos: vos autem non sic in gregem Domini et Ecclesi.am Quare vnus quidam quantumuis magna in Ecclesia praeditus potestate non potest solus excommunicare. Nec aliud sibi vult Paulus cùm ait, Notate per Epistolam. Vult quidemteneneo nuingnincari, ue puuore maiori afficiantur, quod tam probrose innotescant apud eum, quem reuereri debent. Quod si dicamus et legamus, Notate, referentes vocem (Epistola, ad verba sermonem nostrum, haec erit mens Pauli vt velit rebelles a tota Thessalonicensium Ecclesia excommunicari.  \pend\pstart Et quod dicitur pastor πρόεδρως non tribuit illi ius et imperium regium in Ecclesiam: sed in publicis coetibus eminentiorem tantùm locum, et primam in decernendo deliberandóque vocem. Sed vtrum a senatu Ecclesiastico, an vero a to\pend
\section*{AD I. PAVL. AD TIM. }
\marginpar{[ p.334 ]}\pstart ta ipsa multitudine ecclesiae et a quoque fideli soffragium sigillatim ferente fieri debeat, quaesitum est. Res. Senatu suffragium ferente in consistorio fieri debet excommunicatio, ad quem senatum morum et doctrinae cura, diiudicatio, et aestimatio pertinet, conscia tamen et approbante postea ecclesia tota, cui illa res, actio, sententia senatus et caussa nota fieri debet, vt eam approbet ip sa, vel reprobet vt supra docuimus. Quod enim omnes tangit, ab omnibus fieri debet, sed a quoque pro eo gradu (vt docet Ambrosius liber  de Dignitate sacerdotali. cap.31) et ordine, quem habet in Dei Ecclesia. Haec sententia editis libris confirmata est, et abunde antea a nobis probata. Obiter autem, quid in hoc dicto Christi, dic Ecclesiae, vox Ecclesiae valeat et significet explicatum est etiam in ca. Nullus in verbo Ecclesia dist.63.irem cano. Ecclesia de consecratio. dist.1.et cano. scire. 7.quaest.1. ne putent qui nobis contradicunt, nouam afiquam interpretationem a nobis excusam, et fabrefactam. Est igitur illic Ecclesia accepta pro ipso coetu praepositorum Ecclesiae, id est, pro pastoribus et Presbyteris. sed et Dauid Kimhi in Ol5 ait nomine Domus Israel saepe solum Synedrion contineri, et designari in sacra scriptura.  \pend\pstart Quaeritur autem tertio loco. Propter quas caussas fieri debeat excommunicatio. Scholastici distinguunt, et volunt Minorem quidem propter peccatum quodlibet fieri posse. Maiorem autem excommunicationem non nisi propter crimen, id est, tale peccatum, quod et graue sit, et poenam publicam a magistratu mereatur, quale est adul terium, homicidium voluntarium, falsi crimen et  \pend
\section*{CAPVT . V. }\pstart caet. Vtranque tamen non nisi propter rem et deli ctum iam notorium, fieri debere. Haec illi. Alii aliter distinguunt. Volunt.n ex iis tantùm caussis et criminibus excommunicationem iaci posse in aliquem, ex quibus quis damnatus lege Dei lapidaba tur et moriebatur. Ex quibus autem caussis damnatus neque lapidaretur, neque vltimo capitis supplicio afficeretur, non posse quenquam excommunicari. Vtraque distinctio non nititur authoritate scripturae. Ac quidem proxima ciuilem Iudaeorum politiam cum Ecclesiastica plane confundit: illa autem altera Scholasticorum distinctio causas excommunicationis veras non animaduertit, quas in 2. Thessaloni.3. vno verbo tradit Paulus, dum ait, ὠτις ὀυχ ὐπακόει. Ergo et in peceato, et in crimine contumacia animi et obstinatio in perseuerando vel non poenitendo est verissima excommunicationis causa: vti resipiscentia et dolor peccati est causa absolutionis aut reuocationis in Dei coetum et Ecclesiam. Confirmatur nostra sententia dicto Christi Matth. 18.vers.17. Si te non auditdic Ecclesiae, si non audit Ecclesiam, sit tibi Ethnicus. Ergo contumaciam, et pertinaciam animi damnat Christus, quia qui corrigi non vult etiam in re leui, per quam tamen offendiculum praebet Ecclesiae, sed in ea perseuerat, monitus saepe saepiùs, poterit tandem excommunicari si aedificatio Ecclesiae id requirat.  \pend\pstart Notorium autem esse delictum oportet, quia propter scandalum fit excommunicatio. Quod autem est ignotum, offem diculum praebere cuiquam non potest. Ergo ignotum peccatum, vel crimen quantumuis magnum, non est excommunicabil e\pend
\section*{AD I. PAVL. AD TIM. }
\marginpar{[ p.J8. ]}\pstart Quaeritur, Quid si ea res, propter quam quis reprehenditur, verbo Dei non est expresse dannata et definita, Vtrum perseuerans in ea excommunicari possit, velut propter vestitum nimis sumptuosum, nimias mensae lautitias. Respond. Cùm rebellio sit excommunicationis causa, rebellio autem esse non possit, si in nulla re Dei praeceptum violetur, certe quando illud quod facimus, et quomodo facimus, lege Dei non est vetitum, ex eo neque offendiculum, neque iusta reprehensio nascitur, neque etiam excommunicatio fieri vlla debet vel potest. Sed quoniam in tota vita nostra ex Dei verbo, Dei et proximi ratio est habenda, siue in re ipsa quam facimus, nue mmouo, quo quie germaspotest dupliciter anobis peccari, vt in eo peccato nostro obstinate et contumaciter perseuerantes possumus excommunicari, tamen temeraria Ecclesiae iudicia esse non debent.  \pend\pstart Quartùm autem quaeritur, quomodo excommunicatio fieri debeat. Primùm vtrum Voce ipsa, an potius ex Scripto. Deinde qua verborum forma. Ac quidem videntur ex scripto excommunicati ii, de quibus agit in 2. Thess.3. Paulus, non autem viua voce. Ait enim, notate per Epistolam. Et fortasse ad euitandam excommunicatorum inuidiam et odium hoc prodest. Sed contra vsus Ecclesiae obtinuit. Fieri enim vina voce solet: scripto tamem potest. Vtcunque ea postea debet aliis innotescere, quorum intererit scire. Et hoc sibi vult in 2. Thessalonis. Paulus, Sic autem, siue verbis siue scripto fiat, concipienda est, vt et persona ipsa excommunicanda: et factum, propter quod eiici\pend
\section*{CAPVT .V. }
\marginpar{[ p.397 ]}\pstart tur quis e Dei Ecclesia, populo per pastorem, vel Presbyteros significetur aperte, non autem occultetur, ne sit aliis inutilis ista censura. Qua autem forma verborum debeat fieri, disputatum est. Et certe varia et diuersa verborum formula vsa est semper Ecclesia pro locorum, rerum, et factorum, de quibus agebatur, varia ratione. Eius rei exempla extant apud Synesium Episcopum epistola 58.et 66. et apud Euagrium historicum liber 4 cap.38. Qui enim putant his Pauli verbis vtendum praecise, Tradimus Satanae ad interitum carnis, vti est 1.Corinth.5. vers.5.1. Timoth.1. vers.2o. sunt nimium in verbis religiosi. Deinde vim illorum verborum Tradimus Satanae ad interitum carnis, non recte capiunt et intelligunt. Ideo enim tradit Satanae Paulus, vt caro, id est, vitiosa animi libido in iis intereat prorsus: non (vt interpretatur Augustinus liber 3 cap.1o. contra Parmenia. et Greg. in Leuit. Homil.15)in carnis et corporis afflictionem. Nanque  hoc falsum saepe apparet. Neque enim omnes excommunicati in carne et corpore suo hic viuentes cruciantur et vexantur a Diabolo. Vixit enim Iulianus apostata sanus corpore, et validis membris etiam excommunicatus. Sed ostendit Paulus finem excommunicationis, Nimirum, vt praua animi libido euellatur in iis, per quam Dei verbo sunt immorigeri: et spiritus, id est, vis et afflatus spiritus Dei sit in illis ipsis potens, et viuus. Ergo haec verborum forma libera est, modo vis excommunicationis et istius censurae exprimatur.  \pend\pstart Quinto loco quaeritur, quando fieri debeat excommunicatio, an statim atque quis publice de  \pend
\section*{AD I. PAVL. AD TIM. }
\marginpar{[ p.390 ]}\pstart liquit:an demum quando emendari iam amplius non potest, vel quia non est viuus, vel quia est obstinatus et desperatus. Respond. Quanquam huiusmodi quaedam vitia saepe committuntur, vt statim videantur eüciendi, qui ea perpetrauêre, sintque tanquam Ecclesiae pestes, tamen si qui commisit, ex animo dolet, poenitet, et emendatur, non debet excommunicari. Ergo statim fieri non debet, sed postquam agnitus est is, qui deliquit, contumax et obstinatus perdite. Iuberi quidem potest etiam poenitens sceleris, vt si ita sit ex aedificatione Ecclesiae Coena abstineat, sed ad tempus tantùm, non in perpetuum. Praeterea illud sit explorandae huiusmodi hominis poenitentiae causa, vt quantum fieri potest, sciatur, verane sit, an falsa resipiscentia eius, qui se poenitere facti sui testatur. Hoc ipsum docet Paulus saepe suo exemplo. Nam quos excommunicari iubet, ii sunt qui post plures admonitiones vitam tamen in meliorem frugem commutare spreuerant. Ergo ne porcis margaritae contra Christi praeceptum dentur, examinari et explorari eorum poenitentia debet.  \pend\pstart Quod autem quaeritur de mortuis hominibus et iam vita functis, vtrum excommunicari possint. Responsio est, eos posse excommunicari. Exenpla sunt non illa vulgata Stephani 6. et For mosi primi et Sergii primi Pontificum Romanorum, qui suos praedecessores iam mortuos excommunicarunt et persecuti sunt, vti L Sylla C.Marium iam defunctum: sed illa potius, de quibus historiae extant apud Euagrium liber  4. cap.38. item apud Socratem Scholasticum liber 7.cap.45. Et ita etiam  \pend
\section*{CAPVT  V. }
\marginpar{[ p.399 ]}\pstart scribunt et sentiunt Scholastici: atque huius omnis excommunicationis vis in mortuo ad illius nomen et doctrinam, non autem ad personam, quae iam nobiscum non versatur euitandam, pertinet.  \pend\pstart Ter vero an quater prius debeat moneri excommunicandus, definiri certo non potest. Ex rerum enim, et locorum circunstantia hoc totum diiudicandum est. Quanquam quidam leges Roman imitati et sequentes existimant, vt minimùm bis monendos esse, qui excommunicari debebunt. Ter autem admoneri eos satis esse putant, quia trina denuntiatio facta ex iuris ciuilis regulis sufficit, et satis interpellat, vt quis contumax esse iudicetur. Afferunt Paulum Tit.3.V.1o. et Theophylact.in Matth.18.Sed tota haec quaestio est facti, et relinquenda prudentiae senatus Ecclesiastici.  \pend\pstart Sextum quaeritur, quis excommunicationis sit finis, quem spectare debeamus. Ac multi quidem spectari possunt, imprimis autem tres debent obseruari. Vnus est ipsius peccatoris sanitas, id est, excommunicatum, ipsiusque salutem respicit nempe vt hac ratione pudore afficiatur, atque  ab errore reuocetur. Hunc finem annotat Paulus, cùm ait, vt eum pudeat 2. Corinth 2.vers.7.7. vers.12. Pudore enim solent homines ad resipiscentiam adduci. Hoc est quod aliis verbis dicit Paulus vti 1. Corinth.5. vers. 5. ad interitum carnis et vitiosae naturae, quae in nobis inest Neque tamen quod homines vereri magnus sit profectus, si ipse per se spectetur: sed qu est optimus et maxime nostrae naturae consentaneus modus peccatorum corrigendorum, pudor et verecundia  \pend
\section*{AD I. PAVL. AD TIM. }
\marginpar{[ p.403 ]}\pstart peccati. Volunt enim omnes recte de se et honeste ab aliis sentiri et iudicari.  \pend\pstart Secundus finis, est ipsius Ecclesiae incolumitas id est, spectat ipsum Ecclesiae coetum, ne propter vnius priuati scelus male audiat tota Ecclesia, videaturque, si talem retineat, potius sentina malorum, et porcorum esse hara, quam Dei sacratissima ara et templum. Sic improbos et impios e Dei Ecclesia eiici vult et praecipit Paul. Ephes. 5.vers.3.4. Nihil enim immundum in Dei Ecclesia tolerari debet, quantum quidem percipi et deprehendi potest.  \pend\pstart Tertius excommunicationis finis est aliorum fidelium vtilitas, vt hoc exemplo territi ab huiusmour peocacie ce erlinimous uostineant. Corrumpunt enim facile bonos mores colloquia et exempla aliorum praua. Ergo hac seueritate perspecta etiam mali in officio continentur, et timent.  \pend\pstart Septimo autem et vltimo loco quaeritur, quae sit legitimae et iustae excommunicationis poena et effectum. Et certe, vti diuini cuiusdam et coelestis fulminis, magnus est fragor et vis magna. Duo autem effecta solent enumerari excomnicationis praecipua et potissima. Vnum, qui ratione Dei:alterum qui Hominum, spectari debet. Ac quidem Dei ratione excommunicatio legitime et iuste facta excommunicatum separat a Christo et vita aeterna. Dicit enim Christus quaecunque  ligaueritis in terra Matth.18. erunt ligata in caelis ne huius teli vis imbecilla esse iudicetur: aut Dei verbi authoritas vilescat apud homines. Certe excommunicatio conscientiam proprie spectat,  \pend
\section*{CAPVT  V. }
\marginpar{[ p.401 ]}\pstart et animarum salutem, quam adimit obstinato: non est autem poena politica, vt quidam putant. Vti enim distincta est iurisdictio ciuilis, et Ecclesiastica:ita sunt diuersi vtriusque fines, poenae diuersae, et effectus diuersi. Ciuilis enim iurisdictio in corpus aut bona haec terrena exercetur: Ecclehaitieu m aninamiue comerenciam, et in bona aeterna, quibus nos priuat. Itaque a pastoribus infligitur, non a ciuili magistratu.  \pend\pstart Hominum autem ratione priuat excommunicatio excommunicatum non omnium quidem hominum: sed Fidelium tantùm et Christianorum communione. In quo quaeritur, quanam conmunione Fidelium, vtrum omni prorsus communione cum Fidelibus: an certarum tantùm rerum et bonorum cum iis participatione et vsu communi. Respondent Scholastici maiorem excommunicationem priuare excommunieatum omni licita communione cum Fidelibus, siue in bonis spiritualibus, siue corporalibus, terrenis, et ciuilibus. Nos in vniuersum statuimus, excommunicationem legitimam arcere per se excommunicatum ab omnibus spiritualibus et communibus bonis, quae Dominus suae Ecclesiae, id est, omnibus Fidelibus largitus est. Cùm enim excommunicati extra Dei Ecclesiam eiecti sint, non sunt iis magis, quam canibus et porcis danda illa bona, quae Dominus suorum filiorum propria et peculiaria esse vult, illisque solis reseruata, iuxta dictum Christi. Non sunt margaritae obiiciendae porcis: itemque illud, Non est bonum dare panem filiorum canibus Match.15. vers.26.7. vers.6. Aliis autem donis, quae non sunt generaliter omnibus veris et  \pend
\section*{402 AD I. PAVL. AD TIM. }\pstart piis Christianis data non priuat excommunicatio excommunicatum, qualis est donum linguarum, sanationum, et alia huiusmodi. Vnde priuantur excommunicati vsu Sacramentorum, Coenae, et Baptismi, si non fuerint adhuc baptizati, quanquam nondum baptizatos excommunicari moris non fuit, quia nondumin Ecclesia esse huiusmodi censentur. i. qui baptiz ati non sunt, etsi iam Catechumeni fuerunt, et tempus catecheseos improuorantllaguiuauanovride Genesi ad liter. cap. 4o. Solent, ait, in hoc paradiso, id est, Ecclesia a Sacramentis altaris visibilibus homines excommunicati disciplina Ecclesiastica remoueri.  \pend\pstart De filio excommunicati quaeritur, vtrum baptizari possit, si pater maneat obstinatus, neque satagat et curet Ecclesiae reconciliari. Respon. Etsi dubia et agitata fuit haec quaestio, in qua in diuersas alii ab aliis sententias iuerunt, ex Augustino tamen liber 2. de Adulter. coniugiis et liber contra Parmenia. satis intelligitur etiam esse baptizandos excommunicatorum filios, quia filius non debet portare iniquitatem patris: et in ipso excommunicato baptismi character non est per excommunicationem deletus, etsi vis eius nulla est in excommunicato, quandiu quidem in suo scelere manet obstinatus. Vide de hac quaest. Cal. epi. 136. et Bezae epi.1o. Sed quaesitum est, vtrum excommunicatio arceat excommunicatum ipso ingressu et aditu Ecclesiae, vt is ne in coetum quidem admitti debeat, vt audiat Dei verbum. Putant Scholastici arcendum, adeo vt in regno interdicto iuste nolint sacral fieri et celebrari publice et assa voce. Irenaeus liber 2. cap.4. scribit Cerdonem haereticum, postquam excommuni\pend
\section*{CAPVT  V. }\pstart catus fuit, ipso Ecclesiae ingressu prohibitum esse D. Caluinus tamen liber  4. Institutio. cap. 11. sect. 1o. docet diuersam fuisse in eo Ecclesiae disciplinam et consuetudinem, neque semper prohibitos excommunicatos ingredi Ecclesiam. Quando enim Christianis palam conuenire non licuit per imperatorum edicta ad audiendum Dei verbum, tunc excommunicati, tanquam lupi, arcebantur ab ingressu Ecclesiae, ne eam proderent hostibus. Postquam autem coeperunt publice et sine metu coadunari Christiani homines. i. qui Christum profitebantur, neque sibi metuerent a Magistratibus, tam seuera disciplina in excommunicatos non est vsurpata, vt arcerentur ab ingressu Ecclesiae, et abipsa auditione verbi Dei, ne fores et spes omnis resipiscentiae illis praecludi videretur. Vis tamen ipsa verborum et excommunicationis hoc habet et notat, vt excommunicati aditu Ecclesiae arceantur, sed poenae haec pars illis ex quadam Ecclesiae indulgentia remittitur Itaque excommunicati dice bantur ἀποσυνάγωγον, id est, segreges, et sunt eorum loco, quos Auuerpro Eemiseis nabebant. Rtaque aditu templi arcebant, et is etiam vetus Ecclesiae mos fuit Tertul. in apologet. cap.39. canone Cum excommunicato 11. quaest.3. Sic Ambrosius Theodosium imperatorem a se excommunicatum ex regulis Ecclesiasticae disciplinae arcuit ingressu templi Mediolanensis.  \pend\pstart Quod ad terrena autem haec et temporalia, distinguendum est. Aut enim de Publica communione, aut de Priuata et familiari nostra conuersatione cum excommunicatis quaeritur. In iis o\pend
\section*{AD I. PAVL. AD TIM. }
\marginpar{[ p.404 ]}\pstart nibus licet nobis communicare cum excommunicato, quae publica pax et ciuitatis politia et ratio cum eo agere postulat, veluti habet suffragium in creando Magistratu vna cum reliquis ciuibus. Itaque licet cum eo, tanquam cum ciue, agere et communicare: non autem tanquam cum fratre, vel Christiano. Manet enim excommunicatus etiam ciuis, quanquam frater proprie non manet. Nec obstat quod ait 2. Thess.3. vers.15. Paulus monete vt fratrem. Id enim quem sensum habeat, postea explicabitur. Vnde perperam Pontificii excommunicatis omnem sui iuris in iudicio ciuili persequendi potestatem adimunt. Etsi enim potest encoianieatus conuemiri et appellari de debito, et si qua nobis in eum actio competit: ipse tamen ex Pontificiorum iure et sententia debitores suos conuenire non potest, neque vllum ius suum iudicio et in foro ciuili persequi.  \pend\pstart De priuata autem nostra, qui sumus fideles et Christiani, cum excommunicato communicatione etiam quaeritur, quae esse debeat. In quo etiam ipso est distinguendum. Nam cum eo communicare volumus vel tanquam cum homine, vel cum amico, fratre et socio. Omnis quidem arctior cum eo societas, familiaritas, priuataque amicitia nostra dissoluenda, at que dissuenda est. Quae sunt autem ad vitae huius societatem tuendam necessaria, et quae humanitatis, cum eo fieri et haberi possunt. Vnde et ab eo emere possumus, item cum eo contrahere, communem haereditatem habere, ex conmuni cum eo puteo et fonte aquam haurire, illi item possumus nostras merces vendere. Sed etiam si sit vxor, parens, cognatus, aut noster filius ex\pend
\section*{CAPVT  V. }
\marginpar{[ p.405 ]}\pstart communicatus, possumus et officia praestare quae illi debemus: et ab eo exigere, quae debet ipse. Denique quaecunque homo homini debet, ea non sunt deneganda. Quae autem amicus amico ἐκ προαιρέσεως pollicetur, aut offert, ea nobis non sunt vel ab excommunicato quaerenda, vel cum eo contrahenda et ineunda. Veteres tribus maxime Symbolis istam cum excommunicato familiaritatem tam arctam describebant, quibus iubebant pios abstinere, neque  cum eo communia habere. Ea autem Symbola erant, lectus, mensa, domus, quae cum excommunicato communicare nos prohibent, quia ea fere amicitiae, et arctioris familiaritatis signa sunt inter homines Synes.epist.66. et Plutarch. Symposia. liber  8.quaest.7. ea putant esse sacra. Nam ea iis tantùm a nobis communicanda, et offerenda esse, qui communibus nobiscum sacris vtuntur: non autem excommunicatis, non quod tamen vetent isti patres hominem excommunicatum egentem alere, errantem lecto excipere aegrotantem lecto accommodare: sed ita loquebantur, ne vllam amicitiam cum iis contrahamus non autem vt iis vel communia humanitatis vel priuata officia denegemus. Eodem sensu est accipiendum, quod est 2 Ioan. vers.1o. Ne dicamus iis Aue, non vt nos aduersus eos inciuiles et inhumanos ostendamus: sed ne iis blandiri videamur et ita eos in errore foueamus. Ex quo apparet, quae sit significatio vocis eυναναμίγνυθς, qua vti tur Paulus et 1. Corinth.5.vers.9. et 1I. Quanquam enim plus est non habere commercium, vel non misceri, quam se subducere: vtrumque tamen  \pend
\section*{AD I. PAVL. AD TIM. }
\marginpar{[ p.400 ]}\pstart dera, iura, et commercia tantùm pertineat. Id quod satis docet et interpretatur Paulus 1. Cori. 5.vers.11. voce συεθίεην. Ratio est, quia hac voce et prohibitione Paulus non damnat publica commercia, quae cum iis necessario fiunt, quale est ab iis emere, conducere, stipulari, et reliqua huiusmodi. Nec item vetat, si cuius mercis societas antea cum iis contracta fuerit, fidem iis datam seruare, et in eo contractu manere: sed vetat familiares nos iis esse, ne eos in peccatis obstinatiores reddamus.  \pend\pstart Vnde quaesitum est, vtrum quae ex Dei praeceto priuatim quisque alteri, pro vocationis vario genere debet, excommunicatis praestare teneamur: veluti vtrum subditi principi suo excommunicato parere, teneantur, vxor marito, filius patri, seruus domino. Papistae non putant eis haec officia praestanda esse, sed bona conscientia denegari et posse et debere obsequium huiusmodi excommunicatis volunt, quasi soluta sint vincula subiectionis et iurisiurandi propter contumaciam principis in suum ipsius principem, et eum quidem summum, id est, Deum, vt quemadmodum ille non paret Deo:ita neque subditi ipsi principi suo contumaciin Deum et excommunicato parere teneantur. Falsa est haec sententia, et omnino impia et blasphema. Primum enim per excommunicationem non tolluntur iura naturae et sanguinis, veluti cognationum et affinitatum : ergo filius manet filius patris etiam excommunicati, vxor manet vxor: cognati manent illius cognati.  \pend\pstart Deinde non tolluntur iura et officia humanitatis per eandem poenam.Itaque excommunica\pend
\section*{CAPVT VO }
\marginpar{[ p.40 ]}\pstart tis etiam omnia humanitatis officia sunt praestanda.Ad quod affero quod est Romanor. 12.V.20.  \pend\pstart Tertio non tollit excommunicatio iura politica et ciuilis societatis, quemadmodum nec ipsa infidelitas vel apostasia aperta a Deo tollit ea. Nam infidelis et apostata imperator, qualis Iulianus olim, manet Imperator, et Magistracus, et pro tali a Christianis est agnitus. Ergo falsa est Scholasticorum sententia, qui sentiunt non esse parendum principibus excommunicatis: et contra eos ipsos scripsisse videtur Thomas Aquinas ea, quae extant in 2.22 quaest. 1o. Sane Sigibertus author probatus et Monachus scribit in Odone fol.1o1. hanc doctrinam esse et nouam et haereticam, quae iubet ne hypocritis principibus, id est, excommunicatis pareamus, quam Hildebrandus qui et Gregor.7. Pontifex Roman. primus promulgasse videtur, cùm veteres Christiani Iuliano apostatae paruissent. Ambrosius Theodosio a se excommunicato: et ipsi apostoli etiam omnino infidelibus et persecutoribus Ecclesiae principibus, quales Nero et Domitianus, paruerunt.  \pend\pstart Denique quaesitum est, vtrum bonorum confiscatio in excommunicatum sanciri debeat et possit, vt etiam ea poena multetur. Huius rei exemplum extat Esdrae. 1o.ver.8. Respond. Minime. Nam est excommunicatio conscientiae poena:non autem corporis vel bonorum. Et quod illic de Esdra dicitur nihil mouet. Erat enim Magistratus Esdras non solùm Leuitarum, et coetus Ecclesiastici praeses. Ergo vtriusque iurisdictionis, id est, ciuilis et Ecclesiasticae authoritate illud edictum sancitum est, quo celerius Iudaei dispersi conuenirent.  \pend
\section*{AD I. PAVL. AD TIM. }
\marginpar{[ p.Feu. ]}\pstart Haec autem omnia locum demum habent postquam publice quis ex Ecclesia eiectus est: non autem prius quam publice sit excommunicatus. Et quod affertur ex August. liber  de Vnitate Ecclesiae cap. 20. etiam antequam quis visibiliter excommunicatus sit, praecisum esse ab Ecclesia, pertinet ad mentis et conscientiae terrorem incutiendum et piam commonefactionem haereticorum hominum:non autem vt eos pari loco habeamus, atque qui iam palam et totius Ecclesiae suffragio recte sunt excommunicati.  \pend\pstart Videtur tamen omnibus, quae hactenus dicta sunt, de effectis excomunicationis obstare quod ait Paulus, etiam excommunicatum, vt fratrem esse agnoscendum. Cui ipsi repugnat quod ait Christus Matth. 18. excommunicatum habendum esse pro Ethnico et infideli. Ergo non pro fratre. Hae sententiae in speciem pugnantes ita primùm sunt inter se conciliandae, vt quod ait Christus, non interpretemur ἀπῶς et κτ πάντα. Nec enim omnino par aut aequalis debet iud icari excommunicati et infidelis conditio, sed in quibusdam tantùm. Est enim dissimilis vtriusque status in eo quod excommunicatus Baptizatus est: itaque Christianus et Prope est. At infidelis et Ethnicus Baptizatus non est, et Procul est. Sed in eo est simillima vtriusque conditio, quod legitime et vere excommunicatus. 1 vti et Ethnicus in Ecclesia non est, nec de Ecclesia 2 Quod nulla priuata et familiaris nostra cum vtroque conuersatio esse debet. 3 Quod si vterque talis esse perseueret, regno Dei vitaque aeterna excluditur.  \pend
\section*{CAPVT VI. }
\marginpar{[ p.409 ]}\pstart Quod autem ad dictum Pauli, sane et ciuiliter intelligendum est. Ex charitate enim dictum est:non autem ex vera, exacta, et propria fratris, id est, Christiani hominis definitione. Sed quando illud verum est, Nemo desperet meliora lapsis et potest resipiscere qui hodie desipit. Praeterea quos excommunicat Ecclesia, non odio personarum ipsarum, sed eorum tantùm sceleris et contumaciae ratione eiicit, fit, vt qui excommunicati sunt, quatenus et homines sunt, et aliquando nobiscum in Dei Eccesia versati sunt, eorum desyderio, amore, zelo, affectuque salutis ipsorm, moueri debeamus, neque sint prorsus a nobis abiiciendi, sed vt fratres agnoscendi, admonendi atque etiam, vt in viam reducantur, vti oues perditae, sedulo a nobis conquirendi. Ergo etiam excommunicati χτ τι fratres nostri sunt, non tamen omnino et absolute, nec pares iis habendi, qui in Derinetu peritaerut. La quipus omiibus explicatus videtur vberrimus et vtilissimus ille de excommunicatione locus. Sequitur,  \pend
\endnumbering\beginnumbering\section{CAP. VI.}
\phantomsection
\addcontentsline{toc}{subsection}{\textit{\huge\textbf{Q}\normalsize VICVNQVE sub iugo sunt serui, suos dominos ouni honore dignos ducunto, ne nomen Dei et doctrina blasphemetur. }}
\subsection*{\textit{\huge\textbf{Q}\normalsize VICVNQVE sub iugo sunt serui, suos dominos ouni honore dignos ducunto, ne nomen Dei et doctrina blasphemetur. }}\pstart Videri potest abrupta huius capitis et loci series, et minime consentiens cum toto superio\pend
\phantomsection
\addcontentsline{toc}{subsection}{\textit{}}
\subsection*{\textit{}}
\section*{AD I. PAVL. AD TIM. }
\marginpar{[ p.410 ]}\pstart ri argumento proxime sequens tractatio quae est de seruorum officio erga dominos. Quid enim haec commune habet cum pastorum electione, cum censura Ecclesiastica, cum officio Presbyterorum, cum tota denique politia disciplinaque Ecclesiastica, quae omnia superius in tota hac Epistola persecutus est Paulus? Praesertim vero cùm serui eo tempore, quo haec scribebat Paulus ex iuris Romani constitutione pro mortuis inter homines haberentur, id est, perinde atque si omnino nonessent homines, sed pecudes censerentur quae certe pecudes ad coetum Ecclesiae minime pertinent, nullamque nobiscum iuris communionem habent, siue terreni et ciuilis, siue coelestis et aererni. Ecquae igitur huius loci cum superiore continuatio esse potest? Respond. Cùm propositum sit Paulo in hac Epistola depingere atque  informare eam boni pastoris ideam, in qua nihil absit, nullave pars ministerii necessaria desideretur, etiam de seruorum officio erga dominos fuisse dicendum. Cuius rei duplex est ratio.  \pend\pstart I Quod serui etiam ipsi pars sunt Ecclesiastici coetus. Itaque minime a pastoris cura et disciplina semoti. Ergo his, vti et aliis, sua praecepta danda erant, ne propter eos Dei nomen, quemadmodum docet hoc loco Paulus, blasphemetur. Quod vero nobis ex iure ciuili Romanorum obiicitur, in politicis quidem contractibus et conmerciis exercendis verum esse et locum habere concedimus, seruos pro mortuis haberi, nullamque  esse corum voluntatem. Sed in iis, quae ad animarum salutem et ad Ecclesiae atque conscientiae statum pertinent, omnino secus est. Sunt enim  \pend
\section*{CAPVT  yI. }
\marginpar{[ p.411 ]}\pstart etiam ipsi pars Ecclesiae, vti et Domini, quanquam ratione hominum, non tam honorabilis et digna sit eorum conditio, quam ipsorum dominorum. Sed in Christo ( inquit Paulus) non est seruus, neque liber Galat.3.vers.28. Ergo vtrique ad pastorum curam pertinent.  \pend\pstart Altera ratio est Quod si qua maxime quaestio rixose agitata fuit sub origine nascentis Ecclesiae haec vna fuisse videtur. Vtrum serui Christiani suis dominis parere deberent et subiici, vt antea Itaque variis in locis docet Paulus, sed et alii quoque Apostoli, quod sit seruorum erga dominos suos officium, veluti Ephes.6.Colos.4. Denique in epistola ad Tit.2.vers.9. vbi cùm summa tarituni Eplicopalis muneris capita complecti vellet Paulus, de seruorum tamen officio meminit. Meminit etiam eorum Petrus luculenter I. epistola cap 2. vers. 18. vt magni momenti fuisse haec praecepta, et tunc necessaria appareat.  \pend\pstart Cur autem agitaretur haec quaestio de seruorum Christianorum subiectione, aut potius liberatione multiplex fuisse ratio videtur.  \pend\pstart Prima quod propter susceptam Christi fidem et vocatsonem ad Euangenum, dommis suis iam vel adaequati, vel etiam superiores dignioresue censebantur. Adaequati quidem Christianis dominis, quia in Christo neque est seruus, neque liber quisquam. Digniores autem dominis infidelibus, quibus eos Christianos homines parêre videbatur iniquissimum.  \pend\pstart Secunda ratio est exemplum Hebraeorum, qui serui erant. At hi post sextum annum, nempe anno septimo, liberi dimittebantur quemadmo\pend
\section*{41a AD f. PAVT. AD TIM. }\pstart dum extat praescriptum Exod.21.vers.2.  \pend\pstart Tertia ratio est, ipsa saeuities dominorum, imprimis infidelium, quae cùm supra modum esset dura et intolerabilis, videbatur aequum esse, vt se ab eo onere, iniusto quidem et saeuo, liberarent serui Christiani. Praetermitto vero rationem quae ex vitiosa scripturae interpretatione manifeste ducitur: nimirum quod homines libertatis Christianae nomen facile ad licentiam carnis transferunt, et cùm se liberatos esse a Christo audiunt, omne iugum excutere conantur, quod faciunt Anabaptistae.  \pend\pstart Denique vt hae rationes omnes cessent, aequum est tamen pastores consolari miseras illas personas, et ad suum officium alacriter faciendum exhortari.  \pend\pstart Quod enim ad superiores illas tres rationes, plane sunt ineptae. Ac prima quidem, quoniam dissimillima confundit, nempe ius poli, vt loquuntur, cum iure soll. Nam dictum illud Pauli, quod ad nostram salutis rationem prorsus referri debet, transfert ad negotia ciuilia et politias humanas, cùm tamen vtriusque rei sit plane diuersa et dissimilis ratio, etiam ex ipso Dei verbo constituta.  \pend\pstart Secunda ratio continet paralogismum a dicto secundum quid ad dictum simpliciter. Quae enim de Hebraeis seruis sanxit Dominus: ea neque in perpetuum, neque apud omnes homines locum habere debent, sed in populo Iudaeorum. Fuit enim ea lex politica, non moralis. Denique tertia ratio, peccat in modo. Etsi enim remedium aduersus inhumanam dominorum saeuitiem que\pend
\section*{CAPVT  VI. }
\marginpar{[ p.41 ]}\pstart rendum est, non idcirco tamen priuata authoritate sibi ius dicere serui Christiani conceduntur, seque in libertatem asserere, iugum dominorum excutere: sed Magistratus authoritatem implorare debent, Deo prius inuocato. Est ergo, vt ad seriem huius loci veniamus, hic versiculus μετάβοσις, id est, Transitio. Constituuntur enim canones duo vtilissimi de seruis, vti spectari potest eorum conditio duplex ex persona dominorum, quibus seruiunt. Ivam ade imidenbus yuue Piuelibus dominis illi subiiciuntur. Ac primus canon ad eos pertinet, qui dominos infideles habent. Est autem huiusmodi.  \pend\pstart Serui Christiani dominos suos etiam infideles omni honore prosequuntor, eisque subiiciuntor. ne et nomen et doctrina Dei blasphemetur. Secundus autem canon in sequente versiculo continetur. Similis est huic locus apud Petrum 1.Pet.2.v.28.  \pend\pstart Sed singula Pauli verba perpendamus. Nam et Anabaptistas refutant, et Α’ναρχικὸς homines, qui omne discrimen personarum sublatum, omnesque humanas politias abolitas vellent.  \pend\pstart Ait Paulus οι υπὸ ζυγὸν, id est, Qui sub iugo sunt, Iugi voce notare voluit Paulus duram ipsam et miseram seruorum conditionem, et tanquam istius detrectationis seruitutis causam. Itaque studet ipsis seruorum querelis obuiam ire, quasi diceret, De hoc maxime conquerimini quod iugum vestrum durum est, et conditio misera, quia seruilis. Hoc concedo ait Paulus, sed quia dominus ipse vos iis subiecit, et iugum hoc vobis imposuit, aequo animo ferendum est. Nam cùm vos dominus ipse depresserit, alios autem vobis im\pend
\section*{AD I. PAVL. AD TIM. }
\marginpar{[ p.414 ]}\pstart perare voluerit, Deo paretis, cùm illis hominibus paretis, et vestrum est eos a Deo eminentiori loco collocatos vobis digniores censere, et omni honore colendos existimare. Ergo Epicurei homines ἀτακτοι, Ἀναρχικὸt, et Libertini, qui omne imperium tum priuatum, quale est dominorum in seruos, mariti in vxorem, parentum in liberos: tum etiam publicum, quale est Magistratus in ciues et subditos, sublatum vellent, prorlus pugnant cumnoor auirdiete, et distinctionem personarum, cuius Deus ipse author est, euertunt: atque ita euadunt in Deum ipsum iniurii, non tantùm in homines.  \pend\pstart Ait Paulus, omni honore dignos existimanto, Duplex est horum verborum vtilitas. Prima ne serui, neve subditi dominos, aut eos, quos Deus illis praefecit, dedignentur, quod se praestantiores atque longe excellentiores existiment. Nec enim ex industriae, vel corporis, vel ingenii viribus aestimandum est, vtrum nos illis praestemus, atque vtrum eis parere subiicique debeamus: sed ex eo potius, quod dominus ipse illos praefecit: nos autem subditos esse voluit. Dat enim Dominus regna aut potestatem cui vult, non illi semper, quem maioribus et excellentioribus ingenii dotibus et donis ipse ornauit, Quod cùm facit Dominus, iuste tamen facit. Ergo ne subditi, ne serui suas et personas et dotes cum praefectorum vel dominorum suorum dotibus et personis comparent, vt eorum imperium detrectent: sed agnoscant honorent que illam diuinae maiestatis scintillam, quae communicata est a Deo superioribus.  \pend\pstart Secunda vtilitas est, ne perfunctorie tantùm,  \pend
\section*{CAPVT  VI. }
\marginpar{[ p.415 ]}\pstart et quantum ab hominibus obseruari possunt, pareant, et sibi praepositos honorent serui, et subditn.sed ryneere, et ex animo id racsat. Nam ompbαAμοδολύα a Paulo vbique damnatur.  \pend\pstart Quod autem ait, omni honore dignos esse dominos, hanc exceptionem habet, etiam ipsol Paulo teste et interprete Ephes.6.v.6. Ne vltra ea, quae Deus tribui dominis vult, serui pareant. In Domino enim obsequendum est illis, id est, eatenus quatenus Deus ipse ius suum tum in seruos, tum ail fubuitos incegrum et totui retmet; et non offenditur: non autem omnino simpliciter et ἀπλῶς quicquid domini aut magistratus iusserint, statim est praestandum atque exequendum.  \pend\pstart Addita denique ratio est, quod nisi dominis suis serui paruerint etiam infidelibus, Dei nomen et doctrina blasphematur. Haec ratio ducta est a contrario eius finis, propter quem sanctum Dei nomen super nos inuocatur, et Christianam religionem profitemur. Est enim, vt per nos, et propter nos Dei Nomen, et Doctrina bene audiat, propagetur, et glorificetur, idque inter omnes homines. Nos enim in suum populum peculiariter adoptauit Deus propter suam gloriam non vt ipse blasphemetur siue a nobis, siue propter nos ab aliis. Quod cùm fit, grauissime conqueritur Isai.51.vers.5.Rom.2.vers.24.  \pend\pstart Nomen Dei est ipsius gloria, maiestas, sapientia, et tanquam diuina ipsa essentia. Doctrina Dei est illa cognitio, quam de se patefecit per verbum suum, quam nos profitemur. Ergo vtrunque in calumniam rapitur ab improbis, cùm non facimus officium, sed refractarii, improbi, et are  \pend
\section*{AD I. PAVL. AD TIM. }
\marginpar{[ p.416 ]}\pstart πότακτοι sumus. Item vbique docet Petrus in t. epistola cap.2.et3.  \pend\pstart Cur autem blasphemetur Deus ex ista rebellione seruorum et subditorum in dominos suos et praefectos, ratio est multiplex: sed hae duae imprimis  \pend\pstart Prima, quod censetur et Deus ipse, et ipsius doctrina istius confusionis atque ἀ ταίας causa. Atqui Deus neque seditionis, neque confusionis author est, vt docet Paulus 1. Corinth. 14. vers.33. Imo vero nemo potestatem habet nisi a Deo, quemadmodum ait Dominus noster Iesus Christus Ioan.19. vers.11. Ex quo fit, vt qui potestati siut domttiutae, siue publicae resistit, Deo ipsi resistat. Quae vero rerum confusio, quantaque perturbatio inducatur, si neque serui Christiani suis dominis, neque subditi suis magistratibus pareant nunquam satis cogitari potest, nedum dici.  \pend\pstart Secunda ratio, cur Dei nomen ob istam rebellionem blasphemetur, haec est. Quod Deus ipsiusque Euangelium, humana iura tollere, et quod cuiusque proprium est, illi adimere tunc censetur. Itaque furta turpissima praecipere. Nam seruus cuiusque nostrum tam noster est, et ius nobis in eum tam competit, quam quaelibet alia res nostra, nostra est. Quod ius si Christiani nominis professio a seruo facta adimit, furtum rei alienae facere praecipit Deus, et adimit Domino quod est suum ipso, etiam inuito. Neque obstat quod de Onesimo obiici potest, Paulum pro eo deprecatum esse. Hoc enim verum est. Sed neque  seruum inuito domino eripuit Paulus: neque idcirco liberum esse pronuntiat Onesimum, quod  \pend
\section*{CAPVT  VI. }
\marginpar{[ p.417 ]}\pstart Christianus esset factus: imo vero ad dominum suum remittit. Haec autem secunda ratio tanti momenti etiam inter Christianos visa est, vt saepe canonibus Synodorum fuerit sancitum, ne quis seruus, aut addictus alteri, veluti curialis, aut cliens, inuito domino, aut patrono ad clericatum admitteretur, ne ius cuiquam suum praetextu religionis et cultus Dei eriperetur. Est in canonibus Apostolorum. Est etiam in canone Praesenti et seq. vel potius in tota distinct.54. Illud tantùm suit obreruatum, vt qui Audaeis leruiebant, si Christiani fieri vellent, possent redimi ab inuitis dominis, quod in odium Iudaeorum maxime fuit constitutum.  \pend\pstart Ex hoc autem loco manifestissime et pulcher. rime refutatur Anabaptistarum et Αʹτάκλων omnium insania, qui Magistratus de Ecclesia Christiana tollunt, tanquam Magistratuum potestas sit et Euangelicae doctrinae, et Christianae libertati non modo repugnans: sed omnino contraria. Atqui aliud docet hoc loco Paulus, Doctrinam ipsam Euangelii nimirum blasphemari, et merito improbari, si per eam dicantur tolli et peruerti iura politica, confundi rerum dominia, et ius suum cuiquam eripi. Neque Enim euangelium tollit humanas politias, quin imo eas confirmat et sancit.  \pend\pstart Afferunt obstinati homines duos Scripturae locos, e quibus primus est is, qui 1. Corinth.7. V. 23 Nolite fieri serui hominum. Secundus autem qui est Ioan.8. vers.32.36. Veritas vos liberauit. Ad quorum vtrunque vna haec responsio satis est nimirum eos confundere diuina cum terrenis  \pend
\section*{418 AD I. PAVL. AD TIM. }
\marginpar{[ p.]}\pstart Iura conscientiae et seruitutem peccati, a quanos Christus liberauit, cum iure politiae et seruitute hac corporis, quae etiam in Ecclesia manet, et a Christo ipso sublata non est, quia idem iubet reddi Caesari, quae sunt Caesaris, quae distinguit ab iis quae sunt Dei. Denique Christus ipse tributum Publicanis et Caeiaris onciatious persoluit, ne hic latiùs hoc argumentum persequamur.  \pend
\phantomsection
\addcontentsline{toc}{subsection}{\textit{2 Qui vero fideles habent dominos, ne cos despiciant quod fratres sint: sed potius inseruiant, quod fideles sint ac dilecti, vt qui sint beneficii Dei participes. Haec doce, et exhortare. }}
\subsection*{\textit{2 Qui vero fideles habent dominos, ne cos despiciant quod fratres sint: sed potius inseruiant, quod fideles sint ac dilecti, vt qui sint beneficii Dei participes. Haec doce, et exhortare. }}\pstart Ἀνταπόδοσις vel subiectio. Alteram enim propositae dominorum diuisionis speciem persequitur, eorum nempe, qui sunt fideles, Pii, et Christiani domini, quibus quantùm serui Christiani parere teneantur, docet hoc loco Paulus. Est igitur hic alter canon. Serui fideles dominos Christianos quam maximo honore colunto, et quam possunt summisse illis obsequuntor. Vtitur autem hic pulcherrima ἀντιθέσει, qua non modo praeceptum hoc illustrat: sed tacite seruorum obiectionibus respondet, contunditque praefractorum dura ingenia breui responsione. Ait igitur Paulus, serui dominos fideles, Ne spernant, quod sint fratres, sed Potius tantóque magis eos hono rent, quod sunt Fideles Dilecti, Gratiae Dei participes. Æquius est enim vt amicis obsequamur, quam vel inimicis, vel hominibus non ita nos diligen tibus. Æquius est, vt mansuetis dominis  \pend
\section*{CAPVT  VI.. }
\marginpar{[ p.419 ]}\pstart quis seruiat, quam asperis, et immitibus  \pend\pstart Æquius est, vt iis, qui nobis sunt tanquam fratres, nos fummtainus, quai iis, qui nobis sunt velut alieni.  \pend\pstart Æquius denique est, vt iis subiiciamur, quos Dominus ipse Deus tanto honore affecit, vt sint fideles et Christiani, quam iis, quos tanta et pari gratia Deus ipse non est dignatus. Quae omnia argumenta breuiter continentur verbis Pauli. Horum autem argumentorum locus est duplex: Partim enim ducuntur a minori ad maius, partim a personarum ipsarum, quibus subiiciuntur serui, conditione et differentia.  \pend\pstart Ne spernant igitur serui dominos suos fideles id est, Christianos, quod ipsi sint fratres in Chriito Pideles Domi appellantur, non qui pacta seruis suis aut hominibus facta seruant: sed quia fidem Euangelii profitentur. Hic enim fideles appellantur a fide in Christum: non a fidelitate et promissionum praestarione.  \pend\pstart Seruos autem fideles vocat dominorum suorum fidelium fratres, non tam propter eiusdem naturae humanae similitudinem, quae est inter omnes homines siue seruos, siue dominos: omnes enim homines Dominus creauit ex vno sanguine Actor.17.vers.26. quam propter eiusdem adoptionis et fidei participationem. Qui enim per Christum a Deo sunt adoptati, illi sunt inter se fratres, eiusdem baptismatis, eiusdem paris, eiusdem haereditatis et earundem primissionum consortes et κοινωνοὶ vt docet Paulus Ephe.4 vers. 45.6. siue illi statu serui sint, siue liberi: siue Domini,siue non Galat.3.vers.28. 1-  \pend
\section*{AD I. PAVL. AD TIM. }
\marginpar{[ p.420 ]}\pstart Respond. Earum quidem rerum ratione, quarum fratres et κοινωνοι sunt, sublatum est discrimen inter seruos et dominos: reliquarum autem rerum respectu, non item, quales sunt politia humana, corporum sexus, aequales vires, dotes ingenii, externa bona et dignitas, quae omnia et reliqua iis fia laiua manet, et ditincta atque  etiam inter ipsissimos et germanos fratres diuersa et dissimilia, nedum inter seruos et dominos. Id quod rerum experientia docet.  \pend\pstart Haec autem ratio et distinctio latissime patet. Nam siue regimen ecclesiasticum, sine Ciuile considerem9, afferri et dici poterat in ecclesia Dei Christianos subditos et Magistratus: Gregem et Pastores esse fratres inter se, ergo inter eos omne personarum sublatum discrimen videri, neque iam alterum alteri parere oportere: quod homines ἀναρχικοι et ἀτακιοι vociferantur. Sed respondet Paulus ipse fraternitatis nomem differentias vel personarum vel dignitatum siue in politia Ecclesiastica: siue in humana et tetrena mininime tollere, confundere, aut imminuere, vt iam doliiant houis iaicuitere, nósque calumniari, qui Euangelium, quod annuntiamus, quódque est vere Christi Euangelium, criminantur monarchias et regna euertere, atque propterea regum animos ab eo audiendo quantum possunt auertunt, et deflectunt. Imo vero nulli magis parent suis Magistratibus et regibus, quam veri Christiani: quod Remigius etiam Clodoueo pri\pend
\section*{CAPVT  VI. }
\marginpar{[ p.421 ]}\pstart mo Francorum regi respondit, cùm ille falsa ista opmione mmoueus rengionem emiidanam profiteri formidaret.  \pend\pstart Cur autem lubentiùs seruire parereque fidelibus dominis debeant fideles serui, atque subditi, quam prophanis et Ethnicis, triplex ratio videtur afferri, nempe quod tales domini sunt Fideles, Dilecti Deo, Beneficii Dei consortes, vel participes. Itaque triplicem qualitatem designat Paulus, quae nobis tales dominos commendat: neque hos tantùm, sed etiam fideles magistratus. Nam quae de dominis fidelibus ait Paulus, ad fideles magistratus sunt transferenda omnia, atque  accommodanda. Ergo fideles vocat, vt ex communi illo sacro fidei vinculo intelligant et subditi, et serui sese magis illis esse deuinctos, atque obstrictos, quam infidelibus, quibus tamen parendum esse supra docuit.  \pend\pstart Dilectos vocat, quod tales dominos, cùm sint fideles, speciali amore et honore Deus ipse prosequatur, nedum vt eos subditi vel serui spernant. Dilectionem igitur hanc ad Deum refero, et passiue accipio, non autem actiue, nec eam dilectionem esse puto, qua Domini seruos suos complectuntur: sed qua Deus ipsos dominos. Participes item vocat gratiae Dei. Ergo quanto maiore a Deo ipso beneficio et gratia donantur, decorantur, atque  honorantur magistratus, et domini fideles, quam infideles: tanto sunt a subditis et seruis suismagis venerandi atque  colendi. Facit. n.hoc Dei gratia, vt longe quibusuis infidelibus Dominis praecellant quemadmodum docet lacobus 2.v.5 Alii vti Chrysost. Erasm. et quidam, aliter haec  \pend
\section*{AD I. PAVL. AD TIM. 422 }
\marginpar{[ p.]}\pstart verba τπεργισίας ἀντιλαμβανόμενοι accipiunt, et ad curam, quam de suis seruis habent Domini, referunt, quasi hoc argumento vtatur Paulus, eo magis Dominos fideles reuerendos vere esse, quod de suis seruis solliciti sint, illisque benigniores atque humaniores. Sed haec sententia facit seruorum obsequium mercenarium. Praeterea nihil, quod sit speciale aut proprium fidelium dominorum continet, cùm infideles etiam domini, de suis seruis curam gerant sintque illis saepe humanissimi et clementissimi. At Paulus ea notat et obseruat hoc loco, quae sunt fidelium domnotum prepria et peculiaria: non autem illis cum infidelibus dominis communia. Itaque vox ἀρργεσίας hoc loco passiue accipitur: non autem actiue: et ad Deum refertur: non autem ad dominorum liberalitatem et curam de seruis.  \pend\pstart Haec doce) Αʹυξησις. Commendat enim superiorem doctrinam ab vtilitate, cùm eam saepius et repetendam, et inculcandam esse velit. Sic supra 4vers.1I.  \pend\pstart Docere, est non semel tantùm, sed saepius ingerere. Exhortari est viuis aculeis, et stimulis homines ad eam doctrinam amplectendam excitare,ce commouere. La quonles vt non tantùm leui doctrina: sed etiam authoritate et increpatione nos egere propter naturae nostrae peruicaciam intelligamus. Neque enim haec fieri praeciperet Spiritus sanctus, nisi iis omnino haberemus opus, tanquam necessariis animi nostri medicinis.  \pend
\phantomsection
\addcontentsline{toc}{subsection}{\textit{? Siquis diuersam docet doctrinam, ne}}
\subsection*{\textit{? Siquis diuersam docet doctrinam, ne}}
\section*{CAPVT . VI. }
\marginpar{[ p.423 ]}
\phantomsection
\addcontentsline{toc}{subsection}{\textit{que accedit sanis sermonibus Domini nostri esu Christi, et ei quae secundum pietatem est doctrinae. }}
\subsection*{\textit{que accedit sanis sermonibus Domini nostri esu Christi, et ei quae secundum pietatem est doctrinae. }}
\phantomsection
\addcontentsline{toc}{subsection}{\textit{4 Is turget, nihil sciens, sed insaniens circa quaestiones ac verborum pugnas: ex quibus nascitur inuidia, lis, maledicentiae, suspiciones malae, 5 Peruersae exercitationes hominum mente corruptorum, et qui prinati sunt veritate, et quaestui habent pietatem, secede ab iis qui eiusmodi sunt. }}
\subsection*{\textit{4 Is turget, nihil sciens, sed insaniens circa quaestiones ac verborum pugnas: ex quibus nascitur inuidia, lis, maledicentiae, suspiciones malae, 5 Peruersae exercitationes hominum mente corruptorum, et qui prinati sunt veritate, et quaestui habent pietatem, secede ab iis qui eiusmodi sunt. }}\pstart Confirmatio est superioris doctrinae ex ἀντιSέσηι et collatione doctrinae illi contrariae. Atque hac ratione non tantùm illustrat superiora praecepta Paulus: sed etiam commendat, et verissima esse ostendit, cùm quae illis aduersantur, quaeque doctrina illis contraria est, huiusmodi deprehendatur, qualem hic describit Paulus, vel potius Spiritus sanctus, nimirum sit impia.  \pend\pstart Primùm autem quaeritur vtrum hic colophon, qui a Paulo nunc epistolae imponitur, ad vniuersam superiorem doctrinam pertineat, an tantùm ad proximas praeceptiones de seruorum officio erga dominos suos? quidam existimant ad superiores tantùm versiculos duos, qui proxime praecesserunt, referri, quia magni momenti sit superior praeceptio de seruorum et subditorum officio atque obsequio erga dominos et magistratus. Adducor tamen vt sentiam ad vniuersam doctrinam, quae totis quinque superioribus capitibus tradita est, pertinere hanc con\pend
\phantomsection
\addcontentsline{toc}{subsection}{\textit{}}
\subsection*{\textit{}}
\section*{AD I. PAVL. AD TIM. }
\marginpar{[ p.414 ]}\pstart firmationem et ἀπεξεργασιαν, tum quod in calce epistolae addita est, et post quam nulla amplius praecepta, sed particulares tantùm exhortationes subiicit Paulus:tum etiam quod idem solet in aliisepistolis eodem stylo vti, vt suam doctrinam ex ἀντιSέβμ confirmet, ac corroboret, veluti Rom 16.vers.17. Tit.3.vers.9. Coloss.2.vers.18.  \pend\pstart Sunt autem in tota hac ἀντιθέσει et comparatione, quae durat vsque ad versum sextum, obseruanda haec tria nempe, 1 Quid damnet Pausus. 2 Cur dainet. 3 Qua poena damnet. Damnat autem tum Diuersam a superius tradita doctrinam, tum Eos, qui eam vel Docent, vel Sequuntur.  \pend\pstart Ait igitur, ἐτις ἐτεροδιιδασκαλῶ. Vox est composita ἐτεροδιδασκαλέίν ex iis vocibus, quarum vna rem designat, altera docentis studium. Ergo diuersa doctrina per sese damnatur. Similis locus sup. 1.v.3. et Rom, 16 et 17. Tit.3.v.9. Est autem doctrina diuersa, siue Substantia doctrinae ipsius et praeceptorum sit diuersa a superiore: siue Forma et typosis docendi, vt loquitur Paul.2.Tim. 1.V.13 sit vel ab analogia fidei: vel a stylo Spiritus sancti discrepans et dissentanea. Vtrunque  enim in recte docendo debet obseruari, medulla ipsius doctrinae, vt sit sana: et ratio illius tradendae, quae debet esse solida, vtilis, grauis, et ad aedificationem conscientiarum comparata. Et haec de ipsa doctrina.  \pend\pstart Cùm autem ait Paulus, Docet, imprimis doctores ipsos prauae doctrinae complectitur et ferit: postea vero etiam discipulos. Nam et qui docent, et qui attendunt doctrinae falsae, damnantur sup.4 vers.1. quia peccat, qui illicita discit, vti et  \pend
\section*{CAPVT  VI. }
\marginpar{[ p.425 ]}\pstart qui ea docet. Hoc primum est ἀντίθέσιως caput. Secundum autem, rationes varias complectitur, propter quas huiusmodi homines, et merito quidem, damnandi sunt: et tam grauiter, vt quidem excommunicentur, quemadmodum postea apparebit. Quas autem rationes affert Paulus, videntur omnes eae ex tribus locis sumptae nimirum ex. I Verae doctrinae et falsae epithetis et differentia. 2 Causarum, a quibus falsa doctrina emanat, descriptiont, voriplae uocentium personae depinguntur a Paulo. 3 Effectorum falsae ipsius doctrinae (quae quidem sunt teterrima et pessima) enumeratione.  \pend\pstart Ergo ad primi loci argumenta pertinet, quod Pauius verau doctimam vocat Euangelicam, sanos sermones, et doctrinam pietatis: falsam autem doctrinam contr a nominat insaniam, ignorantiam, verborum pugnas, et peruersas exercicationes nue peitheras ainii attiitlones.  \pend\pstart Ad secundum locum referuntur eae voces, quae verissimas falsae doctrinae cudendae causas continent, quales sunt hae inflatus est. Sunt homines mente corrupti, veritate priuati, auari et pietatem in quaestum conuertentes.  \pend\pstart Ad tertii loci argumenta reduco omnia effecta, quae ex huiusmodi praua doctrina consequi docet Paulus, qualis est inuidia, maledicentia, lis, superstitio mala, etc. quae variis in locis enumerantur. Quae omnia eo magis annotasse operaepretium est, quod Satan mataeologiam iis nunc titulis ornauit, atque  tam splendidis nominibus decorauit, vt nihil simile iis omnibus, quae Paulus recenset, habere videatur: sed ita perstringit hominum  \pend
\section*{AD I. PAVL. AD TIM. }
\marginpar{[ p.42o ]}\pstart simplicium oculos, vt sola censeatur esse et vera et luculenta Theologia. Quam tamen hoc loco prorsus de gradu deiicit, et prosternit Spirit sanctus, vt impiam vt falsam atque omnino profanam, et eos etiam ipsos, qui illam docent, vti pessimos atque pernitiosissimos Christi hostes execratur.  \pend\pstart Sed singula iam videamus. Sanos sermones vocat puram Euangelii doctrinam, tum quod illi in sese sunt sani et purissimi, tum quod animas nostras sanas et valentes efficiant, vti dicitur Ps. 19 Praeceptum et verbum Domini lucidum, illuminans oculos. Sanitas animae est vera Dei agnitio, atque cum eo reconciliatio Psal.41.vers.4.Ps.32. quod vtrunque beneficium affert syncera praedicatio Euangelii in nobis efficax. Hoc epitheton saepe tribuit Paulus Euangelio, vt illius vtilitatem commendet prae caeteris artibus, vt haeresibus eam opponat. Sic 2. Timoth. 1. vers.13. sup.1.vers.11. Tit.2.V.1. Huic autem laudi Euangelii et suae doctrinae opponit diuersae doctrinae νόζον quam alii languorem, alii insaniam vertunt. Certe ad animum haec vox pertinet, non autem ad corpus. Est autem non tantùm haeresis: sed omnino mataeologia animi languor. qualis Sorbonica Theologia magna ex parte, mentis aegrotatio et morbus, pestisque certissima animorum, ne sibi in vanis suis philosophamentis placeant insani isti, et memte capti homines neue alio pertinere et spectare censeatur vera Theologia, quam ad animorum salutem et conscientiarum consolationem.  \pend\pstart Appellatur autem secundo loco idem Euangelium Doctrina pietatis, vti et Tit.1. vers.1. et  \pend
\section*{CAPVT  VI. }
\marginpar{[ p.427 ]}\pstart Colloss.1.vers.5.Ac quidem doctrina, quod per eam aliquid et veri et solidi discant atque percipiant homines. Cui laudi opponit, haec duo, nimirum, quod qui diuersam doctrinam sectatur, nihil sciat: et si quid didicit, tamen Nihil Veri, vel Solidi nouit: sed tantùm consectetur λογομα- κίας, et ver borum inutiles pugnas, quemadmodum idem sup.1.v.6.Tit.3.v.9. Vtitur hic Paulus voce Ε'πιςάμενος, quasi fidei nostrae doctrinam vocet ὀιςήμην, qua voce tamen raro vtitur scriptura, sed tamen extat apud lacobus 3.vers.13.  \pend\pstart Αογομαχίας autem nomine non tantùm de vocabulis ipsis et vocibus disputationes inutiles damnat Paulus, sed etiam de rebus, quae nullius sunt momenti, neque ad conscientiarum nostrarum sustentationem, neque ad aedificationem perprortoeutii quo graphice nobis tota Scholasticorum et Sorbonistarum ratio, atque Theologia describitur. Imprimis autem damnamus Quodlibetarias quaestiones, quae appellantur, quae Academicorum potius, quam Christianorum hominum sunt propriae, et mera ludibria pietaris.  \pend\pstart Quaestiones autem neque omnino, neque omnes damnat Paulus, quae de Dei verbo fieri possunt: sed eas tantùm, quae sunt inutiles, otiosae, et sine aedificatione vel fidei nostrae, vel conscientiae, vel charitatis. Nam quae discendi studio vtiliter a nobis fiunt, vti multa de rebus obscuris, quae in vera Theologia versantur, quaeri saepe necesse est, eae minime quidem a Paulo damnantur, quales quaestiones multas tractauit Paulus in epistolis, et Apollos apud Ephesios, vt est in Actis et ov\pend
\section*{AD I. PAVL. AD TIM. }
\marginpar{[ p.420 ]}\pstart Cντητο͂ meminit Paulus 1. Corinth.1.vers.2o.  \pend\pstart Pia vero eadem doctrina Euangelii ab effectu nominatur, quia pietatem nos docet. Est autem pietas, Verus Dei cultus, qui veram proximi dilectionem parit ex sese, tanquam primarium fructum.  \pend\pstart Huic pietati opponit Paulus vitia contrariae doctrinae et Mataeologiae sorbonicae, haec nimirum Mendacium, quia veritate priuantur: Auaritiam, Pietatem in quaestum conuertunt: et Vitia pessima quae profert, qualis est Inuidia, Maledicentia, Suspicio mala, Vaniloquentiam, Contentiones, Pugnae. Denique peruersa exercitatio Tit 3.vers 9. sup.1.vers.6. Caeterae voces sunt faciles et perspicue. Vna haec, πθαδιατριβαὶ, obscura. Attingit enim Paulus eorum speciosos quidens sed tamen falsos praetextus, cùm huiusmodi quaestiones instituunt. Hoc enim se et conari et velle ad Philosophorum imitationem (qui de sue sectae dogmate, studiose inquirunt et disputant, scholasque institutas habent) vt aliquae de pulcherrima omnium scientia, nempe de vera Dei Theologia, et modestae disputationes instituantur, et vigeant inter Christiamnos, quo sint aduersus et haereticorum et Ethnicorum argumenta paratiores atque  munitiores. Philosophi suas illas Scholas et disputationes vocabant διατειβας. Hi quoque  suas de rebus diuinis dissertationes, vt ex hoc loco apparet, διαταβὰς nominabant. Ergo Paulus alludit ad hanc ipsam vocem, quam in eos ipsos lepide retorquet. Docet enim huiusmodi disputationes quae a sana doctrina, recedunt, non esse δια- τθας, id est, honestas scholas et disputationes  \pend
\section*{CAPVT . VI. }
\marginpar{[ p.429 ]}\pstart vtiliter institutas, quales Philosophorum pleraeque fuere: sed ωρα διατριβὰς potius esse. i. pestiferas, peruersas, inutiles et damnabiles λογομαχίας, quia in iis verbum Dei pessime peruertunt. Chrysostom. et Theophylact. aliter interpretantur, confricationes nempe, quod illi sese, tanquam morbidae oues, confricantes corrumpunt. Et certe quicunque ancea sani in Sorbonistarum collegium cooptantur, tanquam in reprobum sensum traduntur, et pess imi Christi ipsius hostes tandem euadunt.  \pend\pstart Quod autem pertinet ad causas peruersae istius doctrinae, illae ex doctorum moribus turpissimis facile intelliguntur, et a Paulo breuiter hic describuntur, vti et 2.Timoth.3. et Colossen.2. et Tit. I. Sunt enim fere hae, nempe, Superbia, Mentis iudiciique corruptio, Mendacium, et ignorantia, siue Profana et impia auaritia. Sunt enim iitius pertis doctores iati nuc luperoi: mente corrupti, priuati veritate, et ignarissimi, auarissimi, pietatem in quaestum turpissimum conuertentes. Quid aptius dici potuerit, quod totam Sorbonicam Scholasticorumque  mores referret, non video, modo nos eorum mores obseruemus?  \pend\pstart Secede ab iis) Tertium jest caput superioris arηθέσεως: in quo, qua poena sint superiores et dannati illi doctores coercendi explicat vno verbo Paulus. Etsi vero in quibusdam codicibus Latinis vetustioribus haec verba non reperiuntur: tamen in aliis extant: et apud Ambrosium leguntur: apud Chrysostomum et Chrysostomi abbreuiatorem Theophylactum, praeterquam quod in omnibus exemplaribus Graecis magno consensu reperium\pend
\section*{AD I. PAVL. AD TIM. }
\marginpar{[ p.430 ]}\pstart tur. Ac primum conuenit hic locus cum alioeiusdem Pauli loco nimirum cum Tit.3. vers.1o. Quod enim hic ait, Secede, illic ait, Reiice. Nam quemadmodum docet Augustinus liber 3.de Anima cap.2.ad pastorem pertinet reiicere, et dannare haereses. Nec vero haec tantùm in huiusmodi doctores et Mataeologos poena statuitur a Paulo, per quam illi priuentur familiari pastorum consuetudine et colloquio: sed per quam e totius Ecclesiae coetu pellantur et eiiciantur. Est enim hic nobis persona Timothei consideranda, quem alloquitur Paulus non vt priuatum aliquem de populo Dei, sed vt pastorem gregis et publicam totius Ecclesiae curam gerentem, adeo vt haec sententia, secede ab iis, idem sonet, atque haec reiice eos de coetu Ecclesia. Quem enim propter haeresim, aut vitia non licet pastori familiarem naberemec aliis fidelibus licet, qui multo minus sibi cauere possunt, ne ab istis corrumpantur. Praeterea non ita istos euitandos pastoribus docet Paulus, vt caeteris sint ignoti: sed quod a Timotheo fieri vult, idem ab omnibus semel fieri oportet.  \pend\pstart Itaque vox haec, Abscede, totum dominicum gregem et Ecclesiam complectitur, qui ab istis fugere debebit. Fateor ego quidem primùm ab ab ipsis pastoribus Ecclesiae tales vitandos esse, ne si iis familiares appareant, reliqui libentius vel audacius cum illis congrediantur, et communicent. Sic de Ioanne Euangelista scribit Irenae. libus 3.cap.3. eum etiam loco ipso abstinuisse, in quo erat Cerinthus haereticus: sed solis pastoribus eos notos aut interdictos esse oportere per\pend
\section*{CIEVIVI. }
\marginpar{[ p.431 ]}\pstart nego: quando quod inde malum nasci porest, ad omnes peruenire potest, et diffundi. Quare haec poena statui videtur a Paulo in huiusmodi doctoles et mataeologos, vt excommunicentur, quae aperte Tit.3. vers.1o. quoque praescribitur.  \pend\pstart Neque vero hoc nouum est, vt qui doctrinam Christi subuertunt, aut ea abutuntur manifeste, quales isti, omnino excommunicentur, atque eiiciantur de Ecclesia Dei. Si quis venit ad vos, inquit loannes, et hanc doctrinam non affert, ne recipite eum domum, ne Aue ei dicite 2.Ioan. vers.1o. Nam si qui Ecclesiae admonitionem et vocem de reconciliatione ineunda cum fratre, spernit, habendus est, ex Christi praecepto, pro Ethnico Matth. 16. quanto magis qui ipsius Christi vocem et Dei veritatem pene totam non tantum contemnit, sed etiam aperte corrumpit? Tantùm abest, vt is pro legitimo Ecclesiae doctore sit agnoscendus, quicquid Sorbonistae (qui hic graphsee depiguntur) contra sentiant atque decernant, vt se suamque Metaeologiam tueantur. Sunt enim isti non tantùm mercenarii pastores, sed lupi, quia animas subuertunt, ideóque a toto coetu Ecclesiae arcendi, quemadmodum praeclare etiam sentit Bernard.epist.237.  \pend\pstart Caeterum ex hoc loco duplex excommunicatio constitui non potest, quemadmodum quidam volunt. Vna nimirum Priuata, altera Publica.  \pend\pstart Priuatam vocant, per quam quis proprio animi iudicio sese ab eo segregat, quem impium vel improbum putat, nullo Ecclesiae iudicio interueniente, aut accedente. Itaque haec excommunicatio est priuatum cuiusque de altero et fratre  \pend
\section*{AD I. PAVL. AD TIM. }
\marginpar{[ p.432 ]}\pstart iudicium, per quod eum vti ab Ecclesia hominem alienum censet. Haec excommunicatio tum schismatum ansam praebet: tum etiam confundit ordinem Ecclesiae. Quanquam enim ab eorum consortio fugere debemus, quos scimus esse Dei irrisores, et homines improbos: tamen eos pro excommunicatis habere propterea non debemus. Aliud enim est abstinere familiari alicuius consuetudine: aliud autem est eum pro excommunicato habere, id est, pro membro reciso, et putri, et cum quo nobis in Sacramentorum administratione communicare salua conscientia non liceat. Haec enim excommunicationis vis est. Haec igitur priuata excommunicatio locum in Ecclesia Dei habere non potest, nisi forte quisque sibi priuatos caetus colligere velit, et nisi nos Donatistarum haeresim comprobemus. Etsi vero multi Deo noti sunt scelerati et improbi, qui tamen nobis pulchri videntur, alia tamen est Dei ratio, alia hominum, qui de occultis non iudicant. Praeterea etsi tibi soli hic sceleratus notus est, propterea tamen non est excommunicandus.  \pend\pstart I uuncan encommameudonem appellant eam, que praecedente Ecclesiastico iudicio, eóque legitimo facta est, vt quis tanquam alienus a Dei Ecclesia censeatur, idque iusta de causa. Hanc solam nos agnoscimus, atque admittimus, de qua etiam loquitur hoc loco Paulus, non autem de illa priore et priuata.  \pend\pstart Duo quaedam autem hic nobis obiiciuntur, Primùm, quod ait Iacobus 5.V.19. tales potius reuocandos esse, quam dimittendos, aut ita semel segreges faciendos. Cum quo consentit quod ait Pau\pend
\section*{CAPVT  VI. }
\marginpar{[ p.433 ]}\pstart lus Galat.6.vers.I.  \pend\pstart Deinde obiicitur nobis quod non dixit Paulus occide illum, trade magistratui puniendum:sed tantùm remoue eum abs te:vel, vt maxima poena statuatur, excommunica eum. Itaque ex hoc loco negant haereticos a magistratibus puniendos esse. Respondemus vero ad vtrunque superiùs obiectum. Ac quidem ad primum, loqui Paulum de iis hominibus, pro quibus conuertendis, et ab errore reuocandis omnia sint prius tentata, et suscepta: sed qui obstinati nihilominus manent, quod pietatem quaestui habeant, sint que auarissimi mortalium. Itaque haec sententia et censura Pauli praesupponit nos omnibus aliis, et mitioribus quidem remediis, vsos esse priusquam ad hoc extremum, id est, ad vrendum et secandum accedamus.  \pend\pstart Quod autem spectat ad secundum obiectum, quod Paulus istos ad magistratum non remittit coercendos, ratio est, quia hoc loto de officio magistratus non agit, neque quod sit eorum munus praescribit, cùm satis sup. cap.2.vers2. et Rom 13. docuerit curam verae pietatis, atque syncerae doctrinae tuendae ad magistratus quoque pertinere. Sed hic tantùm agit de officio Timothei, id est, pastorum. Neque quod de eo nihil addit iam Paulus, probat eos plecti a magistratu ciuili poena vel non posse, vel non debere, non plus quam quod ebrios, adulteros, furaces, etc. in 1. Cor.6.ad magistratum non remittens Paulus non demonstrat huiusmodi improbos homines castigari a ciuili magistratu in Dei Ecclesia vel non posse, vel non debere. Sed etiam haec altera responsio  \pend
\section*{AD I. PAVL. AD TIM. }
\marginpar{[ p.150 ]}\pstart afferri ad superiorem obiectionem potest, quod tunc temporis, cùm haec scribebat Paulus, nullus erat fidelis magistratus, qui piae et verae doctrinae euersores et corruptores puniret. Itaque  de mataeologis aut hypocritis istis ad magistratum, qui tunc nullus erat, remittendis nihil, et merito quidem, praescribit Paulus. De quo argumento prolixa est disputatio August. aduersus Donatistas.  \pend
\phantomsection
\addcontentsline{toc}{subsection}{\textit{6 Est autem quaestus magnus pietas cum animo sua sorte contento. }}
\subsection*{\textit{6 Est autem quaestus magnus pietas cum animo sua sorte contento. }}\pstart Mihi videtur totus hic locus esse διάβασς siue digressio in auaros, et ditescere cupientes homines, accepta occasione ex superiori argumento, nimirum, quod qui sanae doctrinae non acquiescunt, id lucri et ditescendi cupiditate vt plurimum faciant. Sic Ephes.5. vers. 24 inducta comparatione carnalis et spiritualis coniugii, de Christi erga suam Ecclesiam beneficiis late disserit Paulus. Artificiose autem Paulus et lepide etiam hoc loco et versiculo, vti et superiore, vtitur ἀνακλάση retorquens in istos homines auaros vocem ipsam quaestus, quem tam studiose et auideconsectantur. Ac primum grauiter perstringit pungitque eos qui pietatem lucrum esse existimant. Est sane, inquit Paulus, pietas et lucrum et quaestus, isque magnus, qui animum sua sorte contentum semper adiunctum affert vere piis, adeo vt de maioribus opibus comparandis minime sint solliciti et anxii. Secundo vero loco docet in quibus rebus sitae sint verae, stabilesque diuitiae. Demum quid causae sit, cur angi homines de maioribus copiis  \pend
\section*{CAPVT  VI. }
\marginpar{[ p.431 ]}\pstart parandis non debeant: sed paruo debeant esse concentu rido laimnu dlu couus amputationis, quam aliquot versiculis persequitur Paulus.  \pend\pstart Caeterum Pietatem hoc loco Paulus appellat efficacem inhaerentemque cordibus nostris Dei cognitionem, quae ipsius et cultum et metum, tanquam fructum, parit, et profert. Denique eam doctrinam, per quam scimus nos in Christo Dei filios esse, eumque tanquam et Deum et Patrem nostrum colimus, veneramur, et certa animi fiducia agnoscimus.  \pend\pstart Similis locus est Isa 33 vers.6 Pietas enim, ait Propheta, siue metus Dei, est piorum thesaurus atque summae diuitiae. Ergo appellat hoc loco Paulus eam non tantùm πορυς, id est, reditum quemdam et vectigal quod semel rantùm in anno: vel in tota vita percipiatur: sed πορισμὸν esse, ait. i. fundum et perennem siue inexhaustum fontem, quisenper summos vberrimósque prouentus fundat nobis et producat. Ihειζαος enim plus est, quam τορος vel πόρςι. Confirmatur haes eadem Pauli sententia Psalm.4. vers.7. Psalm 35. vers.8. et Psal. 37.112. Denique in eandem sententiam August. libus 1. de Ciuitate Dei cap.19. et Bernard. sermo. 5.in Psalm. Qui habitat copiose disputarunt.  \pend\pstart Dissimilis autem locus 2. Corinth. 6. versao. esse videtur vbi Paulus se et pauperem esse, et nihil possidere dicit. Item dissimilis huic locus fortasse censeatur Philip.4. vers.12. Noui pauperiem pari, ait, quae omnia pugnare videri possunt cum ea laude sofficientiae et opum copiae quam hic tribuit et vendicat Paulus, Pietati et viris piis. Res. Ista minime inter se pugnare. Neque enim abun\pend
\section*{AD I. PAVL. AD TIM. }
\marginpar{[ p.430 ]}\pstart dantia bonorum terrenorum hic promittitur nobis ea, vt nulla re egeamus, aut nunquam premamur: sed potius copia bonorum spiritualium hic notatut, et significatur. Aliter enim pietas dicitur quaestus et lucrum, quam auari opinentur et alias opes affert quam terrenas, nimirum coelestes et spirituales. Praeterea quae inopia et penuria a piis sentitur nonnunquam, ea plane illis ipsis salutaris est. Cedit enim illis bono, nimirum ad fidei probationem, caelestis vitae meditationem, seriam poenitentiam, et ad mundi huius odium et abrenuntiationem meditandam. Ergo pietas alius est quaestus, quam terrenus: alia quoque est abundantia et ἀυτάρκεια, quam vt nihil hic omnino, aut nunquam nobis desit, quoad corporis huius commoditatem ac victum. Nonnunquam enim deest, sed quod deest, tamen veras diuitias nobis non adimit, quia sumus eo ipso, quod adest, contenti, Bernard.sermo. 4. in Psalm. Qui habitat. Ersi non dat ad sa ietatem Deus, dat ad sustentationem. Magnam suis dat penuriam, ne multitudine praegrauentur. Idem sermo.5. in eundem Psal. eodem modo scribit. Sic interpretor quod est Psalm.37. vers.25. Puer fui, etiam senui, nec vidi iustum derelictum, nec semen eius quaerens panem. Quid enim de Lazaro diceremus, de quo est apud Luc. 16.quem abiectissimum mendicum, et Iro pauperiorem fuisse scriptura sacra testatur? Quid de Dauide ipso fugiente et famelico 1. Samuel.21? Sentiunt igitur etiam pii et fideles inopiam terrenorum bonorum, sed eam, quae longe illis salutarior sit, quam copia vel abundantia, et quae (vti dixi) meliores verioresque illis opes minime adi\pend
\section*{CAPVT VI }
\marginpar{[ p.437 ]}\pstart mit, quas opes semper habet secum pietas, et quibus nunquam priuantur pii. Quae sint autem verae Christianorum opes vide August. liber 1. de Ciuitate Dei cap.19. Quare genus ipsum diuitiarum, quas affert pietas, et de quibus intelligit hoc Paulus, totam superiorem quaestionem alias inexplicabilem et implacabilem soluit et dirimit. Est enim hoc genus diuersum a terrenis corporeisque opibus. Quanquam et eas ipsas habent pii, quantum Dominus ipsis expedire nouit. Habet enim fides promissiones vitae tum futurae, tum etiam praesentis, vt supra dictum est cap.4.V.8.  \pend\pstart Cur autem quaestus magnusque prouentus sit pietas, ratio est, quod Dei haeredes, Christi autem cohaeredes nos facit. Rom. 8.v. 17. Est enim nostrae adoptionis in numerum filiorum Dei certissimus fructus et effectus. Nulli enim vete Deum timent, nulli vere colunt, nisi qui vere ipsius sunt filii. Ij autem sunt pii. Ostendit igitur pietas nos esse et Christi ipsius et donorum Christi participes. Nam in Christo insunt omnes caelestes et spirituales salutaresque thesauri, vt docet Paul. Coloss.2. vers.3. Quare ditissimos fideles efficit ac esse arguit pietas. Ex quo fit, vt fideles a Iacobo cap. 2.vers.5. Πλόσινι, id est, diuites appellentur. Et praeclare Ambrosius epistola 1o. Omnes Christianos aduersus profanorum sannas et sapientes et diuites esse contendit. Id quod etiam Psalm. 112. vers.3.5. confirmatur, quanuis fortasse id ipsum multis de terrenis opibus cogitantibus videbitur οάδοξον et ἀτοπον. Sed tamen est verissimum, si quo sensu dicitur, intelligatur, id est, de caelestibus diuitiis. Adiungit autem huic  \pend
\section*{AD I. PAVL. AD TIM. }
\marginpar{[ p.434 ]}\pstart quaestui Paul. Α’υτάρκειαν. Quae vox quid significet quaesitum est. Alii enim eam ad animum nostrum referunt: alii ad res, quibus vitae huius corporeae infirmitas et conditio eget. Certe non est haec Pauli mens, vt dicat, si omnia nobis affatim primùm suppetant, posteaque accesserit ad immensam opum vim Dei metus, nos tunc contentos esse oportere, et magnum prouentum habere.  \pend\pstart Pietatem enim opibus terrenis, tanquam accessionem quandam, adderet et subiiceret. Neque certe haec esset vlla pietatis laus, si cùm omnia necessaria adsunt, contenta esse diceretur. Ex quo fit, vt ἀυτόρκειαν hic minime accipiam, pro rebus omnibus, quae sunt nobis ad victum necessariae, quas cumulatas nobis esse oporteat, priusquam in Christianis Pietas magnus prouentus dici possit. Sed sumo ἀυτέρκοιαν primùm pro ipso hominis pii animo, qui sua sorte contentus et laetus est, quantumuis exiguum possidet, multaque illi etiam de rebus ad victum necessariis deesse videantur. Supplet enim illa omnia pietas ipsa animi tranquillitate, et spe in Deum, cuius voluntati acquiescit qui vere est pius. Deinde sumo ἀυτάρκειαν pro ea opum mensura, quam Dominus dat cuique nostrum, quae semper satis nobis esse et debet et potest, quia si quid praeter eam opus erit, Dominus suggeret. Nam et contenti esse debemus iis, quae in diem ad victum et vestitum suppeditat Dominus ipse, non autem insudare opibus augendis, de crastino solliciti esse, studio inexplebili acquirendi torqueri,, opibusque cumulandis inhiare, et amore, vt ait ille, habendi senescere.  \pend
\section*{CAPVT  VI. }
\marginpar{[ p.435 ]}\pstart Inde quaeritur, vtrum liceat homini pio et Christiano opes suas honestis iustisque rationibus augere sed iis ipsis quas habet contentus esse debeat, neque in id incumbere moderato quodam studio, non auaro, non etiam insatiabili habendi ardore nec ea animi attentione, per quam vitae aeternae curam et meditationem deserat. Respon. Licere, quia modo auaritiae affectus, et animus absit, quod honesto labore et moderata cura comparatur acquiriturque, Dei donusn est, beneficium, atque ipsius innos amoris et benedictionis signum, vt est Psalm.127. Sed hanc curam sollicitudinemque animi moderatissimam esse oportet, ne a melioribus et salutaribus studiis auocemur et abducamur, dum toti in re facienda sumus Matth. 7.vers.34. Ergo si neque Dei cultum, neque Chrimlanl nomis endritatem et officium negligimus, rei familiari augendae recte quidem prospicere et licet, et honestum est, quoniam cura familiae nostrique ipsius est habenda supra 5.vers.8.  \pend
\phantomsection
\addcontentsline{toc}{subsection}{\textit{Nihil enim intulimus in mundum, videlicet nec efferre quicquam possumus. }}
\subsection*{\textit{Nihil enim intulimus in mundum, videlicet nec efferre quicquam possumus. }}
\phantomsection
\addcontentsline{toc}{subsection}{\textit{8 Sed habentes alimenta, et quibus tegamur, his contenti erimus. }}
\subsection*{\textit{8 Sed habentes alimenta, et quibus tegamur, his contenti erimus. }}\pstart Νυθέτοης, quae modum praescribit, qui nobis in terrenorum bonorum vsu et comparatione seruandus est, quem modum qui excedit, is plane et stultus est, et frustra sese cruciat. Fuit autem haec praescriptio valde necessaria, et cùm in hunc locum semel ingressus esset Paulus, aducienda  \pend
\section*{AD I. PAVL. AD TIM. }
\marginpar{[ p.310. ]}\pstart propterea quod si qua in re iustos limites praetergrediuntur homines, in vna hac cumulandarum diuitiarum cura nullam moderationem seruant. Haec autem est admonitionis summa, vt quia haec terrena et externa bona neque nascentes inserimus: neque efferimus asportamusque nobiscum morientes: sed habemus tantùm in hoc mundo ad victus vestitusque subsidium, ipso eorum vsu contenti simus, neque plus in iis quaerendis fatigemur, quam victus vestitusque nostri moderata ratio postulat et requirit.  \pend\pstart Vtitur autem Paulus argumento duplici:quorum alterum ducitur ab absurdo, alterum a Fine istorum bonorum. Ab absurdo. Quod absurdum est in iis comparandis fatigari, quae neque retenturi, neque nobiscum exportaturi sumus, cùm ex hoc mundo et vita decedemus.  \pend\pstart A fine, quod omninum bonorum, quae terrena sunt, largiendorum vna haec causa fuit, vt victus vestitusque nobis ex iis suppeditaretur atque  paratus esset. Itaque praeter victum et vestitum, id est, ipsum vsum, habere et aufferre ex iis nihil possumus.  \pend\pstart Primum autem suum argumentum confirmat Paulus a comparatis, Sic morimur, quemadmodum nascimur. At bona terrena in hunc mundum non inferimusnascentes: ergo nec efferimus morientes. Similis n.est ratio et conditio nascendi atque moriendi. Nudi nascimur et ex vtero prodimus: ita nudi illisque bonis omnibus terrenis vacui morimur, et in communis matris, id est, terrae aluum et gremium redimus et condimur. Similis locus Iobus 1.va1.Eccles. cap.5.v.15. vbi vis huius ar  \pend
\section*{CAPVT  VI }
\marginpar{[ p.444 ]}\pstart argumenti explicatur. Quemadmodum egressus homo est ex vtero matris suae nudus, sic nudus reuertitur, ut discedat, quemadmodum venit, neque quicquam reportat de labore suo, quod venire fecit in manum suam. Illud etiam praeclare Ethnicus quidam Infantem nudum cùm te natura crearit, Paupertatis onus patienter ferre memento.  \pend\pstart Sed obiici posse videtur, si haec non auferimus nobiscum, sint autem ad huius vitae cursum conficiendum necessaria, eo maiori studio videri conquirenda et paranda esse quod iis in hac mortali vita carere non possumus. Respond. Deum ista nobis omnia affatim atque abunde parasse, et ab eo tantùm illa esse petenda, qui gratis et benigne omnia nobis postulantibus largitur et suppeditat: adeo vt propter vexationem illam nostram circa Βιωτια atque anxiam de iis curam, nihil nobis copiosius aut vberius ex illis sit affuturum Mat. 6.vers.26.sed a Deo datur, quicquid benignus ille pater videt ene nobis commouum atque necessarium, et quantum videt expedire.  \pend\pstart Secundum autem argumentum est vberrimum, lateque patet, in quo nobis sunt duo spectanda Primùm, quod ait, habentes victum et vestirum. Secudùmquod ait nos iis contentos esse dobere.  \pend\pstart Neque vero dubitari potest hominem in hunc mundum iam conditum, tanquam in domicilium sibi a Deo paratum, rebusque omnibus abundantissimum et instructissimum inductum fuisse, vt apparet Genes.1. Post peccatum autem et pluribus et aliis rebus egere coepit homo, quam ante peccatum. Cùm enim primum victu tantùm et cibo egeret quia animalis, non autem spiritualis con\pend
\section*{AD I. PAVL. AD TIM. }
\marginpar{[ p.130 ]}\pstart ditus erat postea propter suae nuditatis horrorem egere et vestitu. Itaque Dominus ipse ex iis rebus, quas in hoc mundo condiderat, aptauit confecitque vestimentum Adamo, vt est Genes.3. Adeo vt mundi huius et rerum omnium, quae hominis causa a Deo creatae sunt, vsus siue fructus omnis in his duobus sit positus nimirum vt et Alamur, et Vestiamur in hac vita vt apparet Prouer. 27.vers.26. Eccles.39. vers.31.32. In futura vero vita neque cibo, neque potu, neque vestimento opus habebimus: sed in eo erimus similes Angelis beatis. Vtroque autem hic egemus tum propter naturae nostrae conditionem, in qua primum a Deo facti sumus (nam etiam ederet Adamus, etsi purus et qualis creatus erat mansisset) tum etiam propter peccatl labem et miseriam, in qua iam versamur. Inde est medicamentum, vestiment um varia propter stomachi infirmitatem rerum genera Sic Genes.28. ver s. 2o. duo tantùm haec petit a Deo Iacobus Patriarcha.  \pend\pstart Victus autem nomine etiam potum: Vestitus autem etiam habitationem complectimur, quod etiam Iurisconsulti docent, ne hic angustissimam esse harum vocum significationem putemus. Neque tamen sic stricte Victum at que vestitum nobis a Deo conc edi debemus interpretari, vt lautitiarum omnium vsum denegatum esse existimemus. Nam et odoribus et vnguentis vsos esse patres apparet Psalm. 23. et nobis quoque nonnunquam liberius epulari, suauioribus condimentis vti licere, modo luxus, atque abusus omnis, ebrietasque et crapula a nostris conuiuiis absit, verum est, ne dum auaritiam vitare studemus, in  \pend
\section*{CAPVT VI. }
\marginpar{[ p.439 ]}\pstart miseram superstitionem incidamus. Nam non ad, metitur parce nobis sua dona Dominus, etiam tetrend, ieunuerame protign aeque profundit, quibuscum gratiarum actione vti licet.  \pend\pstart Ex hoc etiam loco duo colliguntur. Primùm cùm nos victu et vestitu contentos esse iubeat Paulus, eóque et tenui et mediocri, non exquisito, non diti, verum esse quod dicitur paucis contentam esse naturam. Id quod etiam Ethnici ipsi senserunt, et Aulus Gellius ex Euripide docet. Praeclare quoq, Horatius. Pracc. Ild.I. cpistol. Pauper enim non est, cui rerum suppetit vsus: Si ventri bene, si lateri est, pedibusque tuis, nil Diuitiae poterunt regales addere maius.  \pend\pstart Ex quo falsa apparet eiusdem Horatii sententia illo eodem liber 1.epistol.  \pend\pstart Exilis domus est, vbi non et multa supersunt. Euripid in Medea.  \pend\pstart Secundum autem est. Non impediri tamen nos, quo minus de bonis et facultatibus nostris testemur, easque aliis relinquamus, etsi huius mundi et earum rerum, quae sunt in eo, proprie domini proprietariique non sumus: sed tantùm vsufructuarii. Nam etsi haec, quae habemus, non asportamus nobiscum, relinquimus tamen, quia  \pend
\section*{47 AD T. TAVL. AD TIM. }\pstart ista dominus nobis ea lege dedit, vt aliis ius, idem vsusfructus relinquete possimus non tantùm vt iis vtamur viuentes. Ergo et morientes legem certam imponere testamento, contractu, legato, aliaque legitima ratione possumus haeredi nostro.  \pend\pstart Denique insulsum est, quod Monachi hoc Pauli dictum ad se solos pertinere putant. Pertinet enim ad omnes Christianos, qui procul a se omnem auaritiae animum abiicere atque remouere debent. Longe verior est concilii Antiocheni cap.35. sententia, quod iubet ne de Ecclesiae bonis Episcopi vel clerici plus sibi, quam victum et vestitum accipiant, si tamen se suis sumptibus exhibere non possunt. In proprios aut domesticos vsus alios, aut in agrorum suorum, vel cognatorum commoda ne ea bona conuertant. Id quod etiam in can. Episcopus 12.quaest 1.relatum est. Denique  de bonis Ecclesiae Sacerdos viuere, non autem luxuriari debet, vt est in dist. 44. Et quod iam ad luxum concessae sunt Episcopis amplissimae opes de bonis Ecclesiae, id factum est cum pauperum magno detrimento et iniuria:atque etiam ipsius Ecclesiae magno malo. Nam quod Canonistae respondent in glossa necessitatem canon. I.21. quaest.1.quia hodie creuit decus Ecclesiae, idcirco praeter necessaria praeterque victum et vestitum debuisse quoque opes maiores Episcopis concedi, quas in aliis rebus, vti in equis, pretiosis vestibus etc.insumerent, vt splendorem Ecclesiae conseruarent, illud est pernitiosissimum et falsissimum, refutaturque canon. Vasa de consecrat. distinct.1 vbi omnem eum splendorem externum Episcoporum et Ecclesiasticae supellectilis, ve\pend
\section*{CAPVT  VIs }
\marginpar{[ p.445 ]}\pstart stitus et victus luxum esse et semper fuisse verissimam Ecclesiae corruptelam respondet Bonifacius Martyr. Id quod verum esse rerum experientia manifestissime docet hodie in Papatu, quo magie nor ao nta magnarum opum copia cauere debent Episcopi Euangelici et pastores.  \pend
\phantomsection
\addcontentsline{toc}{subsection}{\textit{9 Qui autem volunt ditescere, incidunt in tentationem et laqueum, et cupiditates multas amentes et damnosas, quae demergunt bomines in exitium et interitum. }}
\subsection*{\textit{9 Qui autem volunt ditescere, incidunt in tentationem et laqueum, et cupiditates multas amentes et damnosas, quae demergunt bomines in exitium et interitum. }}\pstart Εʹπεξεργασία Persequitur enim susceptam exhortationem de fugiendo opum accumulandarum studio Deterret autem ex effectis et periculo, quod huiusmodi studium comitatur. Ita vero comparat et enumerat Paulus haec effecta diuersa et varia, vt alia aliorum respectu et ratione sint antecedentia et priora ac velut causae: alia vero postersora lequentia, et tanquai teterrimi e superioribus fructus promanantes. Sunt igitur nut checta Cupiditates, quas, et Amentes, et Damnosas vocat. Naufragium siue defectio a fide. Illud effectum prius est: hoc autem posterius est et secundum.  \pend\pstart Est autem totus hic locus ὀθικὸς et παραμνεπαὸν id est, prauorum morum studiique hominum praeposteri et insani increpatio  \pend\pstart Est vero eo magis necessaria haec exhortatio, quod commune, vile, et pene nullum vitium habetur inter homines auaritia: imo vero laudatur nomine diligentiae (vt ait ille tanquam frugi laudatur auarus) nomine industriae, et prudentiae.  \pend
\section*{AD I. PAVL. AD TIM. }
\marginpar{[ p.44* ]}\pstart Vnde praeclare Iuuenalis Saty.14. Sponte tamen iuuenes imitantur caetera, solam Inuiti quoque auaritiam exercere iubentur. Praeterea hominum ingeniis inhaeret tum diffidentia pessima de Dei promissionibus et prouidentia: tum etiam fiducia quaedam in rebus ipsis: quibus cùm caremus, prorsum periisse nobis videmur. Atque vtrunque  hoc vitium facit, vt opibus, diuitiisque comparandis ardenter admodum incumbamus, tanquam solis vitae huius tuendae mediis causis, omnino et anaus neceisariis. Itaque euellendum est ex animis hominum tam tetrum malum, tamque fallax opinio, quae vt venenum serpit. Quod hoc loco conatur Paulus.  \pend\pstart Quaesitum est autem ad quos haec exhortatio pertineat, vtrumne ad solos doctores et verbi Dei ministros, de quibus superioribus versiculis egerat: an in vniuersum ad omnes Christianos. Respond. Ad omnes. Neque enim hic quicquam obseruari potest, quamobrem ad solos, vel doctores, vel pastores Ecclesiae haec doctrina spectet: quanquam hoc lubenter concesserim, si quos oporteat hoc tam foedo vitio, qualis est auaritia, et φιλοπλεσία carere, vnos maxime doctores pastoresque debere, qui caeteris vitae exemplo praelucere debent, et in quorum animos obrepens hoc venenum longe sit futurum nocentifsimum exemplo. Praeterea quicquid affert hic paulus de pessimis φιλεπλούσ́ας effectis, ex aliis Sacrae scripturae locis sumptum est, vbi vniuersi pii ab hoc tam tetro desiderio et vitio deterrentur.  \pend\pstart Duo vero sunt hoc loco notanda. Primùm  \pend
\section*{CAPVT  vk. }
\marginpar{[ p.443 ]}\pstart quod ait Paulus, qui volunt ditescere. Secundum quod subdit, incidunt in cupiditates, etc.  \pend\pstart Similes vero huic in sacra Scriptura loci sunt pene infiniti, sed hi maxime cum hoc conferendi Prouerbus 1.vers.17 18.19. Eccles.5. vers. 1o.11.12.13. Prouerbus 11.vers.4.20.21. Psalm. 119. vers.36. Psalm. 49. Quod si hoc loco profanorum de eadem re sententias colligere liberet, reperirentur innumerabiles. Imprimis autem locus est ille Iuuenal. Satyr. 14quem citauimus. item Q. Flacci Horatii epist. I. et ad Iccium, et Numicium liber 1. epist. et Satyr. 1.libus 1. Satyr. quae profanosum hominum sententiae certe minime sunt nobis negligendae.  \pend\pstart Dissimilis autem huic locus est Prouerbus 12.v. 15. Diuitiae diuitis sunt vt vrbs fortitudinis Eccles. 10.vers.19. argentum ἱCδυναμε͂ rebus omnibus. Itaque vino aequiparantur diuitiae, quod laetificat.  \pend\pstart Horum locorum solutio facilis est, si animaduertatur et recte notetur dictum Pauli. Non enim dicit diuitias per se esse malas, vel diuites ipsos miseros esse et per se Deo exosos: sed desyderium, ardorem, et cupiditatem ipsam auri argentique congerendi esse damnabilem, amentem et exitiosam. Itaque diuitiae hac Pauli sententia minime damnantur, vti neque illo Christi dicto quod est Matth.19.vers.24. tanquam Dei maledictio quaedam in iis inesset. Sunt enim potius benedict ionis Dei erga nos argumentum, si modo iis recte vtamur: sed damnatur studium et voluntas ditescendi, quam φιλεπλεβιαν appellemus licet ardens desyderium et solicitudinem opum accumulandarum vel retinendarum. Fuit enim Abrahamus diues, fuerunt Isaac, et Iacobus  Fuit  \pend
\section*{AD I. PAVL. AD TIM. }
\marginpar{[ p.440 ]}\pstart Dauid, fuit Salomo. Inter fideles, quorum facta est mentio in nouo Testamento fuit diues Manahem educatus cum Herode tetrarcha Actor. 13.vers.1. Electa ad quam scribit Ioan.2. epistolam. Diuites fuerunt Thessalonicenses, vt apparet ex Pauli epistolis ad eos. Denique Psalm.37. Ps.62. vers.1o. Psalm.112. Diuitiae etiam hae, et quidem terrenae promittuntur piis, nedum, vt eas detestari, et abiicere debeamus, si qua legitima ratione nobis obueniant, aut si eas Deus dederit.  \pend\pstart Vetat igitur Paulus ne studium animumque ditescendi habeamus, vel opes nostras retinendi. Vtrumque enim complectitur Paulus siue ditescere velimus opes comparando quas antea non habebamus: vel nostras conseruando ardentiori solicitudine quam oportet. Non vetat, ne familiae curam habeamus, ne bona facultatesque a Deo nobis concessas moderate tueamur, curemus, et iusta ratione conseruemus: vel etiam, si qua ratione legitima fieri possit, eas augeamus. Hoc loco sunt abusi Monachi, qui miseriam et paupertatem pro summo et maxime homine Christiano digno bono vouerunt. Tria sunt enim Monachicarum religionum vota Obedientia, Continenia, Paupertas. Sed etiam hoc eodem loco abusi sunt et abutuntur. Catabaptistae, atque  Anabaptistae, sed et aliis similibus locis scripturae. Nam hos arripiunt, vt doceant homini Christiano non licere quicquam suum aut priuatum habere et possidere, ne φιλοπλεσία alatur inter nos Quod si sit communio bonorum inter nos, cessaturum contendunt hoc omne studium et turpem cupiditatem direscendi, cùm tunc nihil sit pro\pend
\section*{CAPVT  VI. }
\marginpar{[ p.449 ]}\pstart prie nostrum futurum. Sin autem aliquid proprium retinemus, iam hoc damnabile studium in nobis accendi. Sed exemplo Monachorum, qui omnia in communi habent, docemur, ne hanc quidem communionem bonorum extinguere oiλοπλουσίαν, cùm nihil sit Monachis auarius, et ad opes illas communes congerendas ardentius. Nam etiam quod commune est, quodammodo nostrum esse arbitramur, et quanto communitas ipsa erit ditior, nos quoque  ipsos opulentiores, beatiores, ditioresque esse existimamus. Itaque Iesuitarum exemplum et vita siue regula, qui Anabaptisticam illam bonorum communionem admiserunt, sequunturque hodie in suis monasteriis, ostendit euidentissime φιλσπλνσίαν propterea minime extingui. Neque etiam vt nos ab ista habendi cupiditate retrahat Dei spiritus hoc loco, reuocat ad κοπωνίαν bonorum, sed proponit pericula, in quae sese coniiciunt, qui opes sic appetunt, siue priuatas, siue communes. Quae enim mala hic enumerat Paulus tam in communione bonorum, quam in priuata eorum possessione emergunt, ceoriuneuti De qua Irnlavapthtarum sententia plura scriberemus, nisi eam opere elegantissimo vir piae memoriae D. Bullingerus abunde refutasset sex libris contra Catabaptistas editis, in quorum liber 2. Tract.6.hoc delirium reiicitur.  \pend\pstart Gignit igitur et parit hoc studium et affectus varia pericula.  \pend\pstart I Incidunt, qui tales sunt, in laqueos et tentationem, Diaboli nimirum. A quo periculo nobis toto conatu, et velis remisque fugiendum esse docemur quotidie in oratione Dominica his  \pend
\section*{AD I. PAVL. AD TIM. }
\marginpar{[ p.450 ]}\pstart verbis. Ne nos inducas in tentationem, etc. Hoc autem fieri non in vno vel altero tantùm: sed in omnibus φιλοπλουσίοις tum rerum experientia probat: tum etiam docent profanorum hominum sententiae ex naturae sensu et rerum, quas cernebant, vsu prolatae et expressae: qualis illa, Μνξιπολις κακίας απλνςια, item senarius iste vulgatiss.  \pend\pstart Vnde ab Hebraeis appellata est auaritia vra Betsah quod nullum desyderii modum habeat. DeIlque legatur Bernard. in ser.3 in Ps. Qui habitat. et Augustinus liber 1. de Ciuit. Dei cap.1o.epist. IIL.ad Iulia. Comitem. Sed et in Psalm.128.V.36. vna auaritia omnibus Dei mandatis opponitur, quae ex animo nostro delet prorsus, vt Iudae Iscallotae, Itemque Ananiae et Saphirae exempla  \pend\pstart Rationem vero subiicit Paulus, cur tentationes hae obrepant, a quibus vincimur, nimirum quia cupiditates multas affert hic ditescendi animus. Nihil enim insatiabilius istorum hominum animo, nihil concupiscentius: nihil sollicitius nihil cura magis plenum, nihil ardentius et vigilantius  \pend
\section*{CAPVT  VI. }
\marginpar{[ p.451 ]}\pstart nihil curiosius. Omnia circumspectant emissitiis suis illis oculis, omnia spe vorant, omnibus auide inhiant, et Monedularum more, auro et rebus suis perpetuo incumbunt et idem faciunt, quod qui in abysso fodiunt, vt nempe neque fundum reperiant, neque vnquam saccum suum (vt loquitur Aggaeus cap.1.) impleant.  \pend\pstart Appellat autem has cupldstates. Amentes siue stultas, et Exitiosas siue damnosas Pro ἀνοώτες legit vetus interpres ἀνονότος, id est, Inutiles.  \pend\pstart Duplex igitur vitium taxat Paulus in istis auarorum et ditescere cupientium animis. I Quod nullum modum vnquam teneant, nullamque suae personae, status, dignitatis aut sortis rationem habeant. Itaque plane sint isti stultorum et furiosorum hominum more insani, imprudentes, amentes. Itaque  furiosa appellata est cupido opum a Nasone.  \pend\pstart 2 Quod etiam sibi ipsis mala et pericula non tantùm corporis, sed etiam animi accersunt, et creant. Vnde illis et ἀπολεια (quam vocem ad corpus refero) et ὄλεθρος (quam vocem ad animum argumento 2.Thess.1.vers.9. pertinere puto) tandem, et merito quidem, euenit. In quo ipso non tantùm stultitia summa inest: sed etiam miseria, infoelicitas, et horrenda conditio.  \pend\pstart Haec autem omnia effecta opponit Paulus illi laudi et plausui, quem auaris ditescentibus, et opes suas augentibus tribuere solet mundus clamans eos solos beatos, foelices, industrios, et sanos quibus nihil esse et miserius, et furiosius pronuntiat hic spiritus Dei.  \pend\pstart Praeterea eo miserabilior est haec amentia, vt  \pend
\section*{AP I. PAVL. AD TIM. }
\marginpar{[ p.452 ]}\pstart praeclare hoc loco aunotauit Chrysostom. quod quae tantopere optant isti, et quae tanti aestimant, vt dies noctesque propterea crucientur et laborent, opinione tantùm magna sunt et pulchra, non natura. Sola enim opinio illis pretium facit. Nam si quis eos admoneat, doceat que aurum tantùm terram esse: argentum glebam: verum esse apparebit et erubescere cogentur. Praeterea non eodem in pretio ista sunt apud omnes gentes. Nam aurum, vt docet Tertull. De cultu mulierum, et genmae apud Indos, et Æthiopas pro luto habentur, nec nisi in soccis et calceis, idque in contemptum ab iis gestari solebant. Quae autem natura pulchra sunt et pretiosa, ea apud omnes peraeque populos suum pretium suamque aestimationem obtinent, vti Solis istius splendor, coeli pulchritudo, aquae et reliquorum elementorum vtilitas.  \pend\pstart Quod autem idem Chrysost. putat facillimam esse rationem huius flammae et cupiditatis in nobis extinguendae, nimirum si voluerimus (voluisse enimeam extinguere, ait, est eam superasse) Certe de his dici non potest, qui amant diuitias: quia neque ii volunt eam flammam extinguere: neque etiam quandiu illi indulgent, possunt. Praeterea haec ipsa voluntas, per quam huic vitio mederi volumus, non est a nobis: sed a Deo ipso. Sed ita loquitur Chrysost. vt hominibus pudorem incutiat, neque vllam tanto malo medendi desperationem iniiciat, aut excusationem relinquat.  \pend
\phantomsection
\addcontentsline{toc}{subsection}{\textit{10 Siquidem radix omnium malorum est amor pecuniae: quam quidam dum appetunt, aberrarunt a fide, et seipsos transfixerunt }}
\subsection*{\textit{10 Siquidem radix omnium malorum est amor pecuniae: quam quidam dum appetunt, aberrarunt a fide, et seipsos transfixerunt }}
\section*{CAPVT  VI. }
\marginpar{[ p.453 ]}
\phantomsection
\addcontentsline{toc}{subsection}{\textit{xerunt doloribus multis. }}
\subsection*{\textit{xerunt doloribus multis. }}\pstart Alterum argumentum est superioris sententiae confirmandae gratia, ex eodem loco ab effectis ductum. Est autem hoc effectum longe superioribus, quanquam teterrimis, tetrius. Nam docetur φιλσπλεσία et φιλαργυεία esse, non vnius tantùm vitii: sed omnium fons et radix, imprimisque Defectionis a fide, et Cruciatus summi animi. Similis locus Coloss3. v.5. Ephes. 5. vers.3. Matth.13. vers. 22. vbi diuttiae fallaces a Christo nominantur, quod nos a recto tramite abducant, et spirituale semen verbi Dei in nobis extinguant ac suffocent, vt est etiam Ecclesiastici 1o. vers.9. Dissimilis vero huic locus Ecclesiastici cap.1o.v. 15.Initium peccati est superbia.  \pend\pstart Hos locos vt concilient quidam, volunt hoc dicendi genus, quo vsus est Paulus, esse υρερβολικὸν, quodrem nimium augeat, neque sit ad viuum resecandum. Respondeo vero verum dixisse Paulum, neque plus dixisse, quam res est, quamque ipsa expetientia quotidie verum esse docet: sed non esse hanc sententiam in vna tantùm auaritia exploratam: verum in multis quoque aliis vitiis, de quibus nos idem statuere possimus, ea nimirum omne malorum genus gignere et producere. Qanto enim quódque vitium nos longius a Dei metu, et vera fide abducit, eo plura maioraque vitia in nobis parit. Cùm sit enim fides omnium virtutum fons, quanto longius a fide recedimus, eo pluribus grauioribusque peccatis abundamus. At auarus et φιλάργυαςς animus longissime abest a fide, siue fiducia in Deum collo\pend
\section*{AD I. PAVL. AD TIM. }
\marginpar{[ p.454 ]}\pstart canda, longissime a Dei metu, cùm penuriam atque inopiam plus quam Deum ipsum, et omnia Dei iudicia timeat. Idcirco hic affectus merito quaelibet mala et animi vitia secum afferre dicitur. Nam eadem auaritia idololatria dicta est, et Dei obliuio, terrenarumque rerum, in quibus opes sitae sunt, cultus, adeo hominis ingenium auaritia peruertit. Itaque Isai.57.vers.17. pro summo et capitali delicto auaritia recensetur. Praeterea cùm inter vitia alia aliis sint grauiora, hoc maximum est, quod auaritia nominatur. Itaque varia producit. Sunt denique vti virtutes inter se connexae:sic vitia quoque inter se vinculo quodam coniuncta, vt qui vnum habeat plura simul possideat. Ergo haec mens est Pauli: Nullum est malorum et vitiorum genus, quod auaritia secum tandem non attrahat. Quod et de superbia, et de aliis quoque  vitiis dici posse non negat Paulus. Neque enim dixit solam auaritiam radicem esse omnium malorum, vt alia vitia excluderet: sed eam, vt reliqua alia peccata, quae ex eadem infidelitate et diffidentia oriuntur, omnium malorum fontem esse. Neque absurdum est, vt plures sint eorundem riuorum fontes diuersa ratione (vti terra dici potest producere omnes fructus, quibus vescimur, quos eosdem quoque  Sol producit) Sed superbia ratione Dei, in quem sese effert:auaritia hominum respectu quos offendit, dicitur omnium malorum esse causa. P. Lombardus liber 2.distinct.42. cap. Praeterea sciendum et cap.seq. perplexus est in huius loci explicatione. Durandus etiam in eundem locum quaest.3. et 4.magis adhuc quaestionem inuoluit et perturbat. Distinguunt enim isti vitia,  \pend
\section*{CAPVT e v }
\marginpar{[ p.455 ]}\pstart in Capitalia, et Subnascentia, ex illis capitalibus.  \pend\pstart Capitalia constituunt ex Gregorio Magno libus 31. Moralium septem, nempe, Inanem gloriam Iram, lnuidiam, Accidiam, Auaritiam, Gulam, Luxuriam. Itaque quaerunt qua ratione hic a Paulo auaritia radix omnium malorum dicatur, cum praeter cam ana quoque me capdalia vitia. Sed ita hoc a Paulo de auaritia dicitur, vt de aliis quoque non negetur.  \pend\pstart Maior difficultas est, quod quidam censent a Paulo mutatam esse quaestionem. Cùm enim de ριλοπλυσία supra ageret, nunc de φιλαργυρία eum disserere putant, et ad eam argumentum transferre, cuius incommoda et effecta recenset, non autem φίλοπλοσίας. Sed respondeo etsi diuerso vocabulo Paulus vtitur, rem tamen eandem vtroque  significari. Nam et qui ditescere volunt, sunt φιλάργυεα et qui sunt φηλάργυες cupiunt ditescere, sunt ρiλoπλύσον Quod autem quidam quaerunt quomodo Paulus dixerit ἡς τὶς ἐρεγόμενοι referens ad auaritiam vocem ῆς, cùm nemo auaritiam appetat, sed aurum sabet lacnem rcipounonem. Nec enim refert Paulus vocemῆς ad totam φἰλαργυειας significationem: sed ad postremam tantùm compositionis ipsius vocabuli dictionem nempe ἀργυρίας, quae et ἀπὸ τῶ ἀργύρς et ἀργυρίς deducta est. Significat autem et argentum et nummum. Vtrunque autem auari expetunt, et ardenter sitiunt.  \pend\pstart φιλαργυρίας autem voce studium pecuniae manifestissime designat Paulus: sed et omnem quoque rei cuiuslibet turpis et vitiosae immoderatam appetitionem, in vniuersum complectitur. Qanquam enim ἀργύρον vna quaedam earum rerum speF.iiii.  \pend
\section*{AD I. PAVL. AD TIM. }
\marginpar{[ p.40 ]}\pstart cies est, quas auari consectantur maxime: tamen hic omne opum genus complectitur. Itaque oiλάργνεςς hic a Paulo dicitur, qui siue argentum, siue aurum, siue agros, siue frumentum, siue quaslibet alias opes ardenter expetit, et vitiose, id est, et plus quam oportet, et alio modo, atque ad alium finem quam oportet. Nam et auaritia ipsa Latinum vocabulum, quanquam nomen est ab auro deriuatum, immoderatam tamen quarumlibet opum appetitionem vel studium generalicer compicetieur, quemaumiodum docet Augustinus. Alii auaritiam ab auiditate dictam malunt, vnde multa auara appellant, quae auro non egent, vti dicitur mare auarum.  \pend\pstart Duplex autem malum enumerat hoc loco Paulus ex omnibus iis, quae auaritia procreat, nempe Defectionem a fide, et Animi cruciatum. Defectionem a fide ex nimio opum vel congerendarum, vel conseruandarum studio, accidere, infinita hominum exempla docent, et quotidiana experientia. Quot enim hominum millia in nostra, praesertim Gallia, propterea fidem Christi, quam antea professi fuerant., abnegarunt? Quot adhuc hodie propterea desciscunt?  \pend\pstart Quot etiam in media ipsa Ecclesiarum pace propter opes Deum factis negant, quem ore profitentur? Haec vna maxime causa Samosatenum in maledictam illam haeresin, quae postea Arrianismi radix fuit, induxit vt opes a Zenobia Arabum regina haberet, neque ab ea priuaretur iis bonis, quae vt Episcopus Antiochiae possidebat. Et inter causas haereseomn et a Paulo, et ab omnibus auaritia recensetur. Satan ipse ita tentauit Christum Matth.2. Omnia regna mundi tibi dabo, si cadens  \pend
\section*{C ANVTVI. }
\marginpar{[ p.457 ]}\pstart adoraueris me. Animi cruciatum quantum auari sentiant et taciti perferant docet etiam impius Lucianus in Dialogo qui Simon vel Αλεκλρύων inscribitur: Plautus in Euclione, et pulcherrime disputat A. Gollius, ex verbis C. Crisp. Salustii libus 3. Noct. Attic.cap.1. Adeo vt hic animi dolor qui ex diffidentia atque  infidelitate nascitur, sit eius ipsius poenae quae in inferno auaris istis, nisi resi puerint, parata est, praegustamentum quoddam atque  delibatio miserriina e ihiuauissima. Sed non tancumn aiim. verum etiain corporis perpetuus est cruciatus et fatigatio auaritia, quemadmodum praeclare his versibus cecinit Q. Flaccus.  \pend\pstart Impiger extremos currit mercator ad Indos  \pend\pstart Per mare pauperiem fugiens, per saxa, per ignes.  \pend\pstart Reliqua vitia ad pastorum conciones pertinent, quae noc iococumnulaiiin auaritiam possunt.  \pend
\phantomsection
\addcontentsline{toc}{subsection}{\textit{71 Tu vero, homo Dei, ista fuge, sectare autem iustitiam, pietatem, fidem, charitatem, tolerantiam, lenitatem. }}
\subsection*{\textit{71 Tu vero, homo Dei, ista fuge, sectare autem iustitiam, pietatem, fidem, charitatem, tolerantiam, lenitatem. }}\pstart Κατάβασις est. A generali enim praeceptione, quae omnium Christianorum, quemadmodum dictum est a nobis, communis erat, adspecialem descendit vt hanc doctrinam peculiariter pastoribus in persona Timothei applicet. Summa huius applicationis haec est. Si quibus hominibus Christianis scelus illud et teterrimum auaritiae vitium fugiendum et cauendum sit, vnis maxi me pastoribus, Doctoribus, Praefectis, Ministrisque omnibus Ecclesiae id faciendum. Atque per ᚠ-  \pend
\section*{AD I. PAVL. AD TIM. }
\marginpar{[ p.50. ]}\pstart τίθεσιν docet quaenam sint illis ipsis pastoribus sectanda vt habeant et ipsi in quibus se exerceant da oceupent, ionge que menora, quam auari meditentur.  \pend\pstart Ratio huius praeceptionis est ex illo Christi dicto, quod est Matt.5. vers.13. et 14. Vos estis sal terrae, vos estis lux mundi, si sal euanuerit, in quo salietur? Deinde ex eo Petri loco, qui est I.Pet. cap. 5.vers.3. oportet Pastores et Presbyteros Ecclesiae esse exemplar gregis. Qua enim vel fronte vel libertate auaros reprehendent pastores, si sint ipsi auari, et eodem vitio laborent? Nam vt est satis notum.  \pend\pstart Turpe est doctori, cùm culpa redarguit ipsum. Et illud Christi, Medice cura teipsum. Neque  vero noc faeit quo minus paitores et quaedam propria habere possint, ac stipendia ab Ecclesia accipere: et familiae suae honestam moderatamque curam gerere, contra quam tamen sentiunt homines importuni Catabaptistae, vti diximus.  \pend\pstart Dicit autem Paulus φευγο, id est, fuge. Quod vtiliter obseruauit Chrysostom. vt ex vocabulo ipso intelligatur, quanto studio, quantóque animi conatu sit ab hoc auaritiae vitio illis cauendum: fugimus. n. quae periculosissima et maxime exitiosa nobis esse arbitramur, quale certe est hoc scelus. Itaque  non dixit tmn abscede, vt prius ἀφίσα- „n vers.5. Separa te Α ποχώρο: sed fuge, et quidem strenue et diligenter. Videtur autem Paulus voluisse opponere hoc studium nostrum in auaritia fugienda illi diligentiae indefessae atque indefatigabili labori, quem in opibus cumulandis sumunt auari, qui tantus est, vt nullum remittant tempus  \pend
\section*{CAPVT VIa. }
\marginpar{[ p.459 ]}\pstart neque  se respiciant. Breuiter autem addit pondus huic suae exhortationi, cùm eos, quos in Timothei persona alloquitur, nempe Pastores Ecclesiae vocat homines Dei. Ex hac enim ipsorum laude vult eos diligentius officii sui commonefacere. Neque enim solius Timothei propria est et peculiaris haec appellatio. Obiter tamen illud in nostram vtilitatem et expergefactionem notandum est, quod summopere huiusmodi exhortationes sunt nobis necessariae, cùm ne ipsi quidem Timotheo fuerint inutiles: sed proh Deum immortalem, quali viro, ad quem nos collati nihil nisi maculae et turpitudo ipsa videbimur! Nec vero lubens aut iocans Paulus sic eum admonet, sed quod illi ipsi, quanquam multum in Dei metu progresso, vtilem tamen hanc exhortationem fore viderit et sperarit.  \pend\pstart Quaeritur autem, cur Timotheus vocetur homo Dei, et quae sit huius appellationis propria significatio, tum hoc loco: tum etiam aliis pene infinitis. In vniuersum omnes homines dici possunt homines Dei κληςικῶς et γενικῶς, quia et a Deo sunt conditi, et illius sunt. Domini enim est omnis terra et plenitudo eius, et qui in ea habitant. Psalm.24. et Psal.5o. vers.1o. Quod propter Marcionitas, Manichaeos et alios haereticos scitu necessarium est, quorum alii hominem totum opus alterius δημιόργο, quam veri Dei esse docent: alii partem quandam hominis, velut alii corpus, alii animam. Secundo loco dicuntur homines Dei ἐἰδικῶς et ἀγιαςικῶς soli Pii, Fideles, et Electi, quod ab eo sanctificentur, in iis habitet, et eos pro suo populo peculiariter elegerit Deus. Ex quo  \pend
\section*{400 AD LTAVL. ADTIM.. }\pstart Dei pro eodem sumantur. Tertio denique loco et significato dicuntur homines Dei ἰδίως ii tantùm, qui in Ecclesia et domo Dei munus aliquod Ecclesiasticum legitime habent, siue illud ordinarium vt pastores: siue extraordinarium, vt Propheeu. quoupii pius quam renqui fideles speciali quodam iure Dei esse censeantur, et in eius peculio, quemadmodum olim Leuitae.  \pend\pstart Denique  magis se totos Deo addicere hi debeant. Itaque κλυςικο͂ς et proprio quodam possessionis modo dicuntur esse homines Dei. Sic Prophetae homines Dei dicebantur 2. Reg.1.vers.9.4. vers. 16.5.vers.8. Sic pastores omnes a Paulo ratione muneris 2. Flieil.5.veri.rj.non autem rationeelectionis et praedestinationis suae. Sic δῦλον, id est, serui Christi in vniuersum pii omnes dicuntur, imprimis autem iidem pastores et Ministri, vti apparet ex epistola Iudae vers.1.  \pend\pstart Vnum vero hoc loco de Christo specialiter nobis notandum est, eum et hominem Dei et hominem Deum appellari. Hominem quidem Dei, ratione vocationis, ministerii, praedicationis, vti omnes Apostoli, et Ecclesiae pastores appellantur, inter quos ipse est summus. Nec enim nobis verendum est, si Christum, quatenus homo est, et in his terris suo munere fungens, appellemus Dei seruum, quod tamen quidam dicere reformidant, quia ex Isai.61. et Luc.4.vers.18. satis contra sentientium sententia refelli potest. Homo autem Deus, siue, vt Graeci loquuntur, in dictionum conpositione foeliciores, θεάνθρωπος est idem Christus ratione personae et naturarum, quibus con\pend
\section*{CAPVT  VI. }
\marginpar{[ p.461 ]}\pstart stat. Sic enim Athanasius in suo Symbolo appellat. Est enim verus Deus: et verus homo, atque vtrunque quidem vna persona est. Vnde et ab Augustino, et a Bernardo, HomoDeus vnica voce appellatur: et a Graecis Theologis θεάνSpωπις. Nec enim sic Deum suscepit homo Christus vt reliqui homines suscipiunt. Nam alios homines Deus regit tantum, hunc autem etiam gerit, quemadmodum loquitur Augustinus de Agone Christiano cap.20.  \pend\pstart Sed sectare ] Αντιθεσις est, vt et quibus potissimum studiis incumbere: et quas maxime diuitias comparare debeant fideles omnes, imprimis autem Ecclesiae pastores, intelligamus. Ex vtrorumque autem studii, auarorum dico et pastorum, collatione facile iam animaduerti potest, quam longe pastorum et studia et opes, auarorum et studiis et opibus praestent, et sint anteferendae. Varias autem virtutes enumerat Paolus, vt videant pii pastores habere se, in quibus se strenue exerceant, vt et insanas omnes occupationes fugiant, et curiosas eorum quaestiones, quos antea damnauit vers.4.vitent.  \pend\pstart Dicit autem Αιώκες persequere, sectare, id est. summo studio et animi conatu incumbe, quae vox tum opponitur superiori το φεῦγε: tum etiam ostendit non esse nobis fathiscendum neque defatigandum, sed perseuerandum constanter, alacriter et acriter in harum virtutum progressu faciendo. Non omnes autem virtutes recenset Paulus, quae sunt vel pastori necessariae: vel quae auaritiae et curiositati opponuntur (essent enim omnes commemorandae, cùm sit omnium vitiorum ra\pend
\section*{462 AD I. PAVLS AD TIM. }
\section*{CAPVT . VI }
\marginpar{[ p.463 ]}\pstart Neque vero Timotheum, aut alios pastores ad has virtutes exhortatur, tanquam illis catuerit, aut tanquam pastores his vacui eligi in pastores debeant: sed vt in iis quotidie proficiant progressusque faciant maiores, quemadmodum suo exemplo docet faciendum Paulus Philip.3.vers.12. et 14. Nam, vt ait Chrysost. tum neglentiam, tum etiam arrogantiam parit persuasio illa nostrae ὀυταρκείας et perfectionis.  \pend
\phantomsection
\addcontentsline{toc}{subsection}{\textit{12 Decerta praeclarum certamen fidei: apprenende vitam aternam, cuius causa et vocatus es, et professus praeclaram professionem coram multis testibus. }}
\subsection*{\textit{12 Decerta praeclarum certamen fidei: apprenende vitam aternam, cuius causa et vocatus es, et professus praeclaram professionem coram multis testibus. }}\pstart Αʹὐξνσις est. Docet enim non tantùm has virtutes persequendas: sed pro iis etiam nobis esse certandum et saepe quidem durum certamen. Et quia audita certaminis voce, et facta difficultatis mentione saepe homines vel tepent, vel torpent, vel iter reflectunt, et prorsus a bono suscepto desistunt, idcirco calcar addidit Spiritus sanctus vt alacriores ad illud certamen accedamus, illudque ineamus. Est autem triplex consolationis hic propositae locus, nempe a foelici successu. Bonum ait, a fine et premio foelicissimo illi proposito, vitam aeternam: et ab Honesto vel officio, confessus bonam confessionem. Nihil autem hic aliis praescribit Paulus, quod ipse reipsa non praestiterit, vt testatur 2. Timoth. 4. vers.7. et quod reliqui fideles Dei serui non subierint, vt apparet ex eorum Martyrio. Quale illud est Polycarpi, qui se tanquam optimum Domini granum et frumen\pend
\section*{AD I. PAVL. AD TIM. }
\marginpar{[ p.114  ]}\pstart tum mola contundi gaudebat, et illud de Apostolis, ibant gaudentes Apostoli Actor.5.vers.41. Hic locus ostendit non leui manu, vel opera ministris incumbendum in sui muneris functionem: neque ita facilem eam esse, vt plerique sibi pollicentur, et somniant. Cùm enim officium facere volunt, tum in docendo:tum in reprehendendo, multos expenuntur nabentque aduersarios. Itaque ignaui pastores, timidi, somniculosi nunquam vere officio suo recte fungentur.  \pend\pstart Certamen autem hoc loco bonum appellat Paul. cuius foelix, optandus, et prosper exitus est futurus. Nam Deus pollicetur se suissemper affuturum Nunquam te deseram, neque derelinquam ait Dominus Ios.1, v.5. Dominus mihi auxiliator, ait Dauid Psalm.118 vers.9. Quod ad pios omnes transfert Apostolus ad Hebraeos cap.13. vers.6. item Isai.41.vers.13. Est enim Dominus iis omnibus patronus, qui pro se strenue laborant Math.1o.v. 19. Idem certamen praeclarum, id est καλὼν appellat propter rei, de cuius victoria certatur, eminentiam et dignitatem. Est autem Dei gloria, Christi regnum, salus mundi, pro quibus nulla perieula recuianda iunt, cum iis rebus nihil sit praestantius. Denique hic vtamur ea comparatione, quam affert Paulus 1.Corinth.9.vers.25.Si qui milites pro his terrenis rebus certant, tam sunt et constantes et strenui: quanto magis nos ero rebus coelestibus certantes esse oportet et apcres, et inuictos,  \pend\pstart Apprehende vitam) Haec ἀσύνθιτα maiorem vim et δενότητα orationis habent. Argumentatur autem Paulus a praemio huius certaminis. Est  \pend
\section*{CAPVT  VI. }
\marginpar{[ p.46. ]}\pstart autem vita non haec quidem terrena, sed spiritua lis: non caduca, sed aeterna, qua nihil est homini beatius, foelicius, optabilius. Dicit autem apprehende, vt perseuerandum esse intelligamus vsque ad finem, et non desciscendum in medio cursu. Ij enim demum apprehendunt, qui perseuerant. Nec tamen dicit apprehende, quasi vita aeterna sit meritum nostrum: aut quasi illius apprehensio et consecutio sit nostrarum virium (est enim Deï donum Romanor.6.vers. 23.) sed vt doceat non desistendum a coepto itinere et certamine, donec ad eum locum peruenerimus, ad quem a Deo vocati sumus et ire instituimus. Est enim turpe ab itinere semel surcepto co que foelici desistere, et reflecti. Magnitudo autem istius praemii, et quidem gratuiti, satis nos excitare ad quaelibet onera et pericula pro nomine Christi ferenda debebit.  \pend\pstart Oppome dutem heuu praemum Paulus illi exitio et interitui, quem sibi auari accersunt suis curis et cupiditatibus, vt docuit supra vers.9. vt quam foelicior sit noster, id est, piorum in Dei metu progredientium labor intelligatur, quam auaroru, meóque ipso simus alacriores.  \pend\pstart Denique Paulus argumentatur ab officio, et vocatione. Nam pastores ad huiusmodi certamen vocati sunt, non tantum vt pii omnes: sed multo maxime ipsi, quia sunt paitores et caeterorum duces. Similis locus Luc.9. v.2314. v.23. Si quis sequitur Christum tollat suam crucem. Amplificat autem hoc argumentum, a circunstantia et temporis, et rei ipsius Qtod nimirum efficax fuit haec vocatio. et ipse Timotheus illi respon\pend
\section*{AD I. PAVL. AD TIM. }
\marginpar{[ p.466 ]}\pstart derit, Bonam confessionem iam confessus es coram multis, ait, Ergo eo turpius est tibi referre pedem, qui iam tam strenue Christo militasti. Haec confessio, quae facta dicitur a Timotheo, fuit susceptum ab eo Euangelii ministerium quod iam egregie administrarat. Alii tamen ad Baptismum Timothei referunt. Alii quaerunt, vtrùm aliquid passus sit in corpore suo pro Christo Timotheus, vti Sylas et Paulus Acto.16. Sed de Timotheo nihil quicquam simile legimus: tantùm Hebr. cap.13. vers. 23. apparet Timotheum pro nomine Christi vinctum fuisse. Denique summa fuit laus Iimotner in Der ecciesia, adeo vt disputetur, vtrùm ille sit, quem Paulus in 2. Cor. 8.vers.18.19. fratrem appellat, et cuius magnam esse in omnibus Ecclesiis laudem scribit.  \pend
\phantomsection
\addcontentsline{toc}{subsection}{\textit{13 Denuntio tibi coram Deo, qui viuificat omnia, et Iesu Christo, qui testatam fecit coram Pontio Pilato bonam illam professionem. }}
\subsection*{\textit{13 Denuntio tibi coram Deo, qui viuificat omnia, et Iesu Christo, qui testatam fecit coram Pontio Pilato bonam illam professionem. }}\pstart Παράκλησις est cum superioris doctrinae confirmatione coniuncta, quam vtranque cum grauissima obtestatione proponit hoc loco Paulus. Nam et superiora praecepta hac quasi nota fuerunt obsignanda, ne quis humana potius, quam diuina: et Pauli, non autem Spiritus sancti vocem esse existimaret: et fuit in opere et ministerio iniuncto et demandato Timotheus propter muneris ipsius difficultatem, confirmandus. Apparet autem quam sit vtrunque necessarium. Ac primum quidem, quod etiamnon hodie quidam contentiose  \pend
\section*{CAPVT VI. }
\marginpar{[ p.467 ]}\pstart contendunt, quae hic de forma electionis pastorum Ecclesiasticorum a Paulo tradita sunt, temporaria tantùm praecepta esse: quae non vt perpetuo locum in ecclesia haberent sint tradita. Itaq: ea et omitti, et mutari posse. Quam sententiam falsissimam esse docet hic locus, atque vel vna haec tam grauis diligens et sollicita de his obseruandis obtestatio a Paulo facta, praesertim cùm sup. cap 3.vers.21. aliam quoque obtestationem addiderit Paulus, non minus grauem.  \pend\pstart Secundum autem huius versiculi et loci praeceptum non minus quam prius necessarium est, cùm reipsa experiamur quam pauci pastores vel in suo officio perstent: vel si perstant, quam magnis cum difficultatibus luctentur, adeo vt in hoc onere eadem alacritate animi perstare sit omnino arduum, ac rarum, et, nisi Dominus ipse vires et robur suppeditaret, prorsus αδύνατον. Exempla Ioannis, Marci, et Demae, de quibus tum in Actis: tum in 2.ad Timotheum agitur, quod dicimus verissimum esse comprobarunt. Quare neque temere, neque praeter rem, aut in vanum hanc tam grauem, et plenam inuocatione nominis Dei obtestationem addidit Paulus, ne nomen Dei in vanum sumpsisse, et pro re leui, aut etiam in non necessario loco vel negotio interposuisse Paulus censeatur.  \pend\pstart Duplex autem est obtestationis huius atque denuntiationis authoritas, tum ex persona Patris tum Filii. Vtrique  autem suum epitheton adiunctum est ad maiorem ipsorum testium venerationem. Ait Paulus Denuntio. Est autem quiddam aliud δαγsίλ λεν, id est, denuntiare, quamμαρτυρεῖν qua  \pend
\section*{AD I. PAVL. AD TIM. }
\marginpar{[ p.408. ]}\pstart voce etiam sup. cap.5.vers.21.vsus est, etsi maiestatem Dei vtraque phrasis declarat, et designat. Denuntians ea, in quibus omnem negligentiae, ignorationisue excusationem hominibus ademptam esse volumus, quasi eos commonefaciamus prius iam satis edoctos fuisse, et nunc de poena negligentiae serio metuendum cogitandumque illis esse. Obtestamur, cùm nostrum officium et operam fideliter praestitam esse sigiiiieamtis, de quo etiant Deum testem inuocamus. Itaque denuntiatio ad ipsos homines dirigitur, eósque intuetur. Obtestatio proprie Deum quem spectat.  \pend\pstart Coram Deo) Dei nomine Patrem hic intelligi, vti alibi saepe in scriptura, certissimum est, et ex his verbis apparet quae mox sequuntur Jesu Christo Sic passim in epistolis Paulus accipit, praesertim in ipsarum ὀιγραραῖς, vti in hac ipsa, et Tit.1. Ex quo fit, vt recte sit annotatum Dei vocabulum quanquam proprie ad naturam ipsam diuinam significandam repertum est: vt vox haec, homo ad humanam:tamen saepe essentialiter, saepe personaliter sumi Essentialiter igitur acceptum tribus personis sanctae Trinitatis aeque competit, easque generaliter complectitur. Personaliter autem positum eam tantùm Trinitatis personam proprie designat, quam et locus et argumentum et res, de qua agitur, ostendit intelligendam. Saepeigitur Patrem significat, vt hic:saepe filium vt cùm dicitur Deus apparuit Abrahamo Genes. 18.vers.1.item vocatus Abraham auscultauit Deo Hebr.11.vers.8. Saepe Spiritum sanctum, vt cùm dicitur Deus loquutus per Prophetas Hebr.I.v.  \pend
\section*{CAPVT  VI. }
\marginpar{[ p.469 ]}\pstart 1.Actor.28. vers.25. Sic in concilio Nicaeno primo dictum est Deus de Deo, id est, Filius de Patre, non quod diuersa essentia sit filii a Patris essentia. Nam vti Pater est Deus, sic et filius: sed a Patre Deo genitus est Filius Deus, idem tamen et vnus vterque Deus, etsi non vna et eadem vterque persona. Cur autem vna tantùm ex tribus diuinis personis Deus possit appellari, id est, voce totam naturam et essentiam diuinam signiMeunte de compieccentesuenotuiry ratio est, quod in quaque ex illis tribus personis inest tota illa essentia diuina, atque ea etiam ad numerum vna et eadem, non autem pars tantum Deitatis.  \pend\pstart Additur autem epitheton, Qui viuificat omnia. Similis est locus Actor.17.vers.28. Nam quod de hominibus dicitur illic, ad omnia rerum genera pertinet in vniuersum, vt apparet ex Psalm.36.v. 6. et 9. Viuificare autem significat et vitam dare: et datam conseruare. Vtrunque enim praestat Dominus omnibus rebus, quae viuunt, non alterum tantùm. Neque enim deserit res a se creatas Matth.1o.vers.37. Vnde praeclare Bernard. serm. 4. in Cantica. Sane esse omnium dixerim Deum, non quia illa sint, quod est ille: sed quia ex ipso, per ipsum, et in ipso sunt omnia, quia esse omnium est, causale quidem, non autem materiale, quia eorum est factor. Nec enim diuinae essentiae participatione res viuunt: sed qualitatis tantum et virtutis ipsius aliqua communicatione et effectu. Vnde mn Deus appellatur, non autem Iehouah tantùm  \pend\pstart Sic autem magnifice, et, vt res est, de Dei potentia concionatur Paulus, vt in doctrinae Chri\pend
\section*{AD I. PAVL. AD TIM. }
\marginpar{[ p.470 ]}\pstart stianae veritate tuenda, et muneris sui conscientia retinenda Pastores opponant hanc Dei maiestatem, robur, et potentiam omnibus hominum etiam regum, minis, viribus, et conatibus. Nam hominum ipsorum Dei Euangelio inimicorum vires ab vnius Dei potestate dependent, vt est Psalm.59. vers. Io. Itaque non sunt nobis ii metuem di, ne que propterea expauescendum est.  \pend\pstart Duo autem quaeruntur hoc loco. Primum ex eo, quod dixit Paulus, Omnia viuificat. Vtrum omniarerum genera, quae Deus creauit, viuant, id est, vitam habeant. Ac quidem earum rerum quas condidit, aliae sunt Spirituales, vt Angeli, et Aninimae humanae: aliae corporeae, et eae vel Caelestes, vel terrenae. De Angelis siue bonis siue malis: itemque de hominum animabus dubitari non potest quin vtrunque  hoc spiritualium creaturarum genus viuat. De rebus autem corporeis maior quaestio est, inter quas aliae sunt Caelestes, vt coelum ipsum: aliae terrenae, vti Elementa quatuor, Conposita ex elementis corpora, quae dicuntur, vt Mineralia, vt Plantae, vt Animalia, atque  ex iis quaedam perfecta, qualia quae quinque sensus naturales habent ex suae speciei ratione et constitutione. Alia imperfecta animalia dicuntur, vt polypus, ostreorum genus quoddam, et muscarum. De Plantis et Animalibus nemo dubitat, quin hoc vtrunque genus rerum viuat. De Mineralibus, de Coelo ipso, de Terra elemento, de Aqua et reliquis similibus, maior et magis perplexa quaestio est. Nam si haec viuere dixeris, et animam quoque habere concedas necesse est: et sibi simile procreare posse. Quod rerum veritati repu\pend
\section*{enrela va. }
\marginpar{[ p.471 ]}\pstart gnat. Neque enim coelum, aliud coelum sibi simile procreare aut producere potest, vel ex semine vt herba: vel ex suae essentiae profluxu, vt homines et bruta creant. Varro in libris de lingua Latina sentit et scribit, quaedam viuere, quae tamen animam non habeant. Aliis visum est et coelum et terram, suo modo etgenere viuere. Aliis contra haec rerum genera subsistere quidem, sed sine vita et anima, quae in illis interne insit. Vnde ea vitam proprie non habere sentiunt. Respond. Locum hunc Pauli, in quo versamur, esse ita explicandum, vt vocem Omnia ad ea tantùm referamus, quae vitam habent, quia illa Dominus viuificauit. Sic dicit Paulus, vult Deus omnes homines saluos fieri, id est, (vt ait Augustinus) vult Deus omnes homines, qui salui fiunt, saluos fieri. Ergo nihil nobis necesse est in curiosam illam quaestionem, Vtrum omnes creaturae, quae a Deo sunt conditae, viuant, ingredi, etsi eas omnes cuiuscunque sint generis viuere, et vi quidem illis a Deo interne iniecta quidam putant, cùm et sustententur vi formae suae, et motus actionesque suas pulcherrimas habeant, et quemadmodum terra facit multa ex se pariant. Nam volunt spiritum Dei, qui in prima rerum creatione terram, et coelum, et aquas fouebat, et terrae, et aquae, et coelo eam vim internam infudisse, per quam hodie subsistant et durent in sua natura, vti herbae vim internam habent, qua viuere dicuntur, et in genere suo permanent, etsi sibi simile neque coelum, neque elementa generent, vti neque homines, qui sunt steriles ex se vllos procreant, nemo tamen  \pend
\section*{AD I. PAVL. AD TIM. }
\marginpar{[ p.RIE. ]}\pstart eos viuere vi interna a Deo data propterea negauerit, Vti etiam totum mularum genus sterile est ipsa natura:nec tamen propterea viuere quisquam disstebltur. Sed haec disputationis tantùm gratia dicta sunto atque accipiuntor.  \pend\pstart Secundo vero loco quaeritur. Cur hanc operationem tribuat vni tantùm patri Paulus, cùm et Filii et Spiritus sancti sit communis haec actio. Ac de Filio apparet Ioan. 5.vers. 21. Sicut Pater excitat mortuos et viuifieat: ita et Filius, quos vult viuificat. De Spiritu sancto vero item apparet Psalm. 104 vers.3o. Emittis spiritum tuum et creantur, et 1.Samuel.2.vers.6.generaliter de Deo dicitur, Deus est qui viuificat, et morti tradit. Respond. Non assignari a Paulo hic propriam et incommunicabilem Patris operationem, nisi quatenus ille est reliquarum personarum ratione principium, sed communem trium actionem describi: quae tamen in Patris persona hic consideratur, quatenus Pater pro tota diuina natura saepe proponitur, et in eo contrahitur quicquid trium personarum commune est. Quod fit, quia Christum nobis hic subiicit Paulus, considerandum vt hominem. Ergo in persona Patris contrahit quae sunt totius Deitatis propria, id est, trium personarum communia. Coram fesu Christo, qui testatam fecit) Cùm multa sint, quae de Christo et aliis personis proprie incommutabiliter dici possint, veluti quod incarnatus est, passus, mortuus, resurrexit, hoc potissimum delegit Paulus, quod ad praesentis loci et argumenti circunstantiam maxime facere videbatur, et ad superiorem de testimonio fidei  \pend
\section*{CXPVT. VI. }
\marginpar{[ p.473 ]}\pstart perhibendo exhortationem erat aptissimum. Est autem illa confessio Christi, de qua loan.18.vers. 37 mentio fit, et per quam constantissime in ipso mortis periculo comparens coram iudice suo supremo in his terris veritari Dei, ac Euangelio apertissimum atque luculentissimum testimonium reddidit. Itaque ipsius etiam iudicis mentionem fecit Paulus, vt et periculi magnitudo, et mortis terror, qui ante Christi oculos obuersabatur, notaretur a nobis, non tantùm autem vt historiae vetitas atque  rei gestae tempus designaretur. Est autem Christus nobis datus in exemplar, quod imitemur, vt ait Petrus in 1.epist cap.2.vers.21. Huius eiusdem confessionis meminit quoque, quanquam obscurius, Apostolus ad Hebraeos, et ad illius exemplar imitandum nos accendit cap.12. vers.3. Multas enim contradictiones in ea veritate asserenda passus est Christus.. Hoc dictum Pauli explicatur in articulis Symboli Apostolici. Itaque erimus breuiores. Vnum modo notandum est Ethnicos vafre et subdole tempore Imperatoris Iuliani Apostatae commentaria quaedam Pilati supposuisse, ex quibus inepta, blasphema, et prorsus falsa Christo Domino nostro affingebantur, quae et Augustinus liber 1. de Consens. Euangelist. et Cyrillus aduersus Iulia, prolixe refutauit. Appellat autem hanc Christi confessionem bonam, non tantùm propter rem ipsam quam fassus est Christus, quae omnium rerum maxima, optima, fructuosissima fuit. Dei enim voluntatem et veritatem de salute hominum perficienda declarauit constanissimeque asseruit Christus. Sed etiam bonam vocat propter finem, propter quem illud  \pend
\section*{474 AD I. PAVL. AD TIM. }
\marginpar{[ p.]}\pstart testatus ac confessus est Christus, mimirum vt Deus glorificaretur propter misericordiam patefactam et in homines collatam Denique propter effectum vberrimum et amplissimum quem haec confessio protulit. Nam morte ipsa a Christo statim obsignata est et confirmata. Itaque illa confessio fuit verissimum testimonium Christi de voluntate Dei in saluando genere humano, quae multa, imo infinita postea hominum millia Deo lucrata est, quos in ipsa eadem doctrina confirmauit, atque etiam hodie confirmat. Talis quoque est hodie mors piorum Martyrum. Mors autem Socratis, mors haereticorum obstinate suos errores defendentium non est confessio bona, tum quod de haeresi sit ista confessio non de veritate Dei: tum etiam quod non sit obsignatio doctrinae, quae pios fructus proferat. Nam Socratis mors certe non fuit Socratis doctrinae confirmatio, sed potius fuit illius velut extinctio quaedam et suffocatio. Itaque cum morte et cum confessione Christi nullo modo est conferenda, nedum cum ea aequanda.  \pend
\phantomsection
\addcontentsline{toc}{subsection}{\textit{14 Vt serues haec mandata, maculae et culpae omnis expers, vsque ad illustrem illum aduentum Domini nostri Iesu Christi. }}
\subsection*{\textit{14 Vt serues haec mandata, maculae et culpae omnis expers, vsque ad illustrem illum aduentum Domini nostri Iesu Christi. }}\pstart Ε’ξήγηοις est. Docet enim iam de qua re tam grauem obtestationem addiderit et praemiserit. Nimirum vt pastores omnes, imprimisque Timotheus ipse seruet superius mandatum, idque constanter et pure donec Christus illustri suo illo et secundo aduentu appareat in suorum ple\pend
\section*{CAPVT VI }
\marginpar{[ p.475 ]}\pstart nam consolationem, et victoriam. Mandati nomine complectitur Paulus superiora praecepta a se tradita: et ipsum Timothei munus, id est, demandatam illi prouinciam in Ecclelia Ephesina, in qua eum purum et immaculatum versari vult, atque iubet. Quod si tanto viro hac exhortatione fuit opus, quid nobis? Et certe nunquam satis ad officium excitari, aut expergefieri possumus, quia etiam vigilantissimis et perfectissimis obrepunt saepe multae tanquam tepiditates, vel animi et zeli remissiones: denique vitia et errata imprudentibus, propter quae Dei gloria retardatur, Ecclesia offenditur, et nos ipsi coram Deo rei simus. Itaque attente et diligenter omnibus in vniuersum Christianis: imprimis autem pastoribus in suo munere fideliter et constanter peragendo vigilandum est, atque incumbendum.  \pend\pstart Sed haec est magna eorum consolatio, quod munus ipsum quod habent, mandatum Dei appellatur. Ex quo intelligunt primum se Dei ipsius esse voce vocatos, et hanc sibi delegatam essea Deo non ab hominibus tantùm prouinciam, vt in ea mtrcpiue, conidentence aideriterpersentur. Deinde auxilium sibi a Deo semper affuturum, qui eos in ea statione collocauit Actor. 13.vers.3. Coloss.4. vers.17. Galat.1.vers.15.sup.1.vers.1.et 12 denique ex eo ipso cogitent pastores sibi huius muneris demandati rationem esse Deo reddendam, a quo illud habent.  \pend\pstart Expers maculae) Docet iam secundo loco non tantùm quam in rem sit incumbendum pastoribus Paulus: sed etiam quomodo, nempe ea conscientia et fidelitate, eaque diligentia et purita\pend
\section*{AD I. PAVL. AD TIM. }
\marginpar{[ p.310 ]}\pstart te, vt nulla ex parte sibi suóque muneri desint, nec quenquam vel vita vel negligentia offendant demque viuant ireprenensibiles, vt loquutus est sup. 5.vers.22.  \pend\pstart Ergo verba haec sic interpretor, vt ad ipsum Timotheum referantur: non ad mandatum, etsi tamen id patitur ratio Grammatices. Postea vocem ἀπιλον refero ad leuiora delicta: vocem autem ανέπληπιιον ad grauiora peceata, quae et scelera et crimina nominantur, vt apparet ex cap.3.huius epistolae.  \pend\pstart Vsque ad illustrem) Docet tertio loco, quam constanter in tota superiori doctrina ipsóque  suo munere fit et omnibus Christianis et pastoribus perseuerandum. Nimirum vsque ad secundum Christi aduentum, quem Jllustrem appellat. Similis locus Tit.3.vers.13.1.Thess.2. vers.19.  \pend\pstart Est autem huius sententiae et consolationis vsus duplex. Primus quidem aduersus ignauiam nostram, ne fathisca mus, aut deficiamus longi assiduique laboris pertaesi, quasi nimium diu nos exerceat et in certamine retineat Deus. Secundus autem vsus est aduersus impatientiam, ne Deum existimemus vel mentitum: vel diutius, quam par sit, promissorum suorum implementum differre. Nec enim prius certandi finis erit, quam Christus aduenerit iudex viuorum et mortuorum.  \pend\pstart Appellat autem hunc Christi aduentum E’πmiserae illi conditioni, qua nunc premimur Deo vero seruientes, opponamus: tum vt hac cogitatione omnes afflictiones quas pro Christi nomine pati hic necesse est, lubenter deuoremus, in i\pend
\section*{CAPVT  VI. }
\marginpar{[ p.17 ]}\pstart psius coelestem gloriam intentis fixisque penitus oculis intuentes. Atque haec consolatio et admonitio non tantùm illo saeculo fuit opportuna atque  necessaria, quo vera Dei Ecclesia passim ab omni mundo spernebatur atque irridebatur: velut mundi reiectamentum quoddam vt ait Paulus 1. Corint.4 sed etiam hac nostra aetate certe quidem est vtilissima, qua et ab extraneis, id est infidelibus, et ab intraneis, id est, hypocritis non minus uffligimur, spernimur, irridemur, quam olim vere pii, et homines Christiani. Haec enim molesta et dura conditio manet Ecclesiam Dei omni tempore Psalm.123.  \pend
\phantomsection
\addcontentsline{toc}{subsection}{\textit{15 Quem temporibus suis ostendet ille beatus et solus Princeps, Rex regum, et Dominus dominorum. }}
\subsection*{\textit{15 Quem temporibus suis ostendet ille beatus et solus Princeps, Rex regum, et Dominus dominorum. }}\pstart Α'ὔξησις est. Amplificat enim superiorem exhortationem tum a temporis definiti certique circunstantia: tum a Dei ipsius, cuius doctrina seruari praecipitur, et cui veri et fideles Ecclesiae pastores seruiunt, dignitate, et comparatione. Continet quoque hic locus ύποφοραν. Occurrit enim tacitae obiectioni piorum, qui illustrem illum Christi aduentum iamdudum expectantes, nondum tamen eum videbant. Ergo ne vel animis denielanit poit longam expectationem illius aduentus: vel obrepat blasphema ista cogitatio, qua si Christus venturus non sit, et ab eo Ecclesia delu datur, hoc responsum tanquam necessarium animorum fluctuantium pharmacum et remedium subiicit Spiritus sanctus, Nimirum tempora istius aduentus esse certa quidem illa et defjnita: sed soli Deo nota, et constituta.  \pend
\section*{AD I. PAVL. AD TIM. }
\marginpar{[ p.110 ]}\pstart 2 Nostra infirmitas, qui varie tentamur, fluctuamus, et agitamur, vbi tot nobis obiectos aduersarios, in tot tela nos incurrere videmus, dum in Dei metu, in officio, in munere, in professione nominis Christi perseueramus. Itaque sustentandi, contrillandr; corroSorandique sumus aduersus tot hostes, obiecto potentiae Maiestatisque Dei clypeo, et foelici illo beatóque statu, qui nos in aduentu Christi certissime manet, constanter in fide perstantes.  \pend\pstart Magna igitur nostrae temeritatis et impatientiae repressio, magnum lorum est, quod ait Paulus, Tempora illius foelicis et expectati a nobis aduentus esse quidem certa: sed soli Deo cognita neque Deum vel momentum vllum dilaturum, cùm tempus a se decretum impletum fuerit, quin Christum exhibeat redemptorem nostrum, et vitae aeternae donatorem ac assertorem. Eodem vero modo, quo hic de secundo aduentu loquitur, de primo quoque Galat.4.vers.4.Ephes.1. vers. 1o. ait tempora illius certa fuisse a Deo constituta, post quae completa Christus venit, atque homo factus est. Scit enim Dominus non tantùm, quae nobis dona sint salutaria: sed eorum conferendorum et communicandorum tem\pend
\section*{CAPVT  VI. }
\marginpar{[ p.47 ]}\pstart pora et tanquam maturitates, quia solus est Sapiens. Et vt in rebus his terrenis et naturalibus, quibus producendis aliquod a Deo definitum est tempus, patientes esse solemus, veluti anni tempestates expectamus: frugum maturitatem non praecipitamus, dierum incrementa non vno momento esse lubenter ferimus Iac. 5. vers.7. sic in his quoque donis coelestibus et salutaribus accipiendis patienti nos esse animo oportet, et legitimum eorum reuelandorum tempus expectare.  \pend\pstart Tribuit autem Paulus Patri eam secundi et gloriosi Christi aduentus manifestationem: at alibi Christo ipsi videtur vindicari Tit.2.vers.13.1. Thessalonic.4. vers.17. Respond. Quia Christus me nouis nomo de mimus propontus est, Patri tribuuntur, quae sunt propria Dei opera, atque etiam ea ipsa, quae Christo conueniunt quoque, quatenus ipse Deus est. Sic dicitur Pater suscitasse Filium: dicitur etiam Filius seipsum suscitasse.  \pend\pstart Beatus, solus princeps. etc. Est diuinae maiestatis descriptio ex variis illius attributis, vt hanc omnium hominum, quantumuis potentum, viribus et dignitati opponamus, cùm ab illis propter Christi Euangelium, subsannamur vel affligimur, vel de die aduentus Christi, tanquam impedito, ipsi cogitamus. Maxime vero illo seculo fuit necessarium oculos fidelium traducere ab eo splendore, qui in Ecclesiae hostibus apparebat, ad incomparabilem illam vnius Dei maiestarem, thiscerent pii, qui se a potentissimis illis rerum dominis, imperatoribus et Magistratibus Romanis, a fastuosis Iudaeis, denique ab vniuerso mum\pend
\section*{AD I. PAVL. AD TIM. }
\marginpar{[ p.300 ]}\pstart do sperni videbant et se velut mundi sordes et immunditiem abiici, vt docet Paulus I. Corinth. 4.vers.5. qualis est hodie Gallicanarum Ecelesiarum status, quibus Dominus ipse magna sua misericordia aliquando succurret. Praeclare vero et recte annotat Chrysost. hic apparere, quantus sit Pauli in Dei maiestate celebranda zelus, vti alibi quoque saepe idem potest a nobis obseruari. Hic enim sibi temperare non potest, sed pulcherrima variorum epithetorum congerie admirabilem et incomprehensibilem Dei maiestatem esse docet, nósque in eius admirationem rapit. Itaque manum, quod aiunt, non potest tollere de tabula. Sic autem sunt commemorata haec omnia ve quudai gradutio niis obieruata fuisse videatur. Appellat enim Deum primum beatum: deinde solum potentem: tum postea Regem regum et sic deinceps crescente vocum significatione rem ipsam auget.  \pend\pstart Sed quaeritur, cur Deum beatum et potente vocet, qui potius est ipsa beatitudo, et potentia dicendus, quam beatus et potens. Sunt enim hae voces, quibus vtitur Paulus, παρώ υμαι, vt Dialectici loquuntur, quae alterius rei, velut qualitatis, participationem et attributionem denotant fieri in subiecto, de quo vsurpantur. Et tamen Deus nullius beatitudinis communicatione est beatus, nullius potentiae participatione potens: sed est ipsa et beatitudo et potentia, et eo ipso beatus et potens, quia est ipsum illud quod est. Omnia enim quae Deus habet sunt Dei essentia: non qualitates: nec alio quam seipso Deus vel beatus est vel potens. Respond. Nos homines cogi cùm de  \pend
\section*{CAPVT VI }
\marginpar{[ p.481 ]}\pstart Deo loquimur cuius natura nobis est incomprehensibilis, nostro more, nostrisque verbis loqui, vt aliquid de eo pro ingenii nostri captu dicamus et explicemus. Quod etsi non est aequale Dei ipsius maiestati et essentiae: ex nostro tamen more sentiendi et loquendi verum est. Itaque  Deum dicimus et beatum et potentem, non quasi compositam illius naturam statuentes: non quod illi qualitates et accidentia tribuere velimus, et aliud quid praeter simplicissimam ipsius essentiam in eo fingere: sed quod humanus captus atque humanum eloquium aliter explicare commode et satis intelligenter seipsum non potest, cùm de Deo est nobis agendum. Cùm igitur sic loquatur Deriphitus de Deo,neque non ipinde eo sic loqui metuamus, et Deum appellare beatum, potentem, sapientem, et c. Quod fit citra vllam Dei blasphemiam, si, quemadmodum diximus, fuerit acceptum.  \pend\pstart Beatus) Μακάριος. Est Deus omnis foelicitatis fons, et scaturigo, quam et in se, et a se habet, non aliunde, vt angeli, beati: vel etiam vt homines. Sic etiam ab Ethnicis dii ipsorum appellantur μακάριονvt passim ab Homero, quod sit memtibus hominum inscripta a natura haec sententia Nihil Deum esse posse, quin illud idem sit beatum in se, et corruptionis omnis expers et immune. Id quod vox Μακαρ proprie significat, quasi per epenthesin literae μ voce deducta a vocabulo ακα- tùs quod insecabile et imminutum significat. Sic Erasmus ex Grammaticis interpretatur. Secutus sane est Paulus vulgarem et tritam vocis huius significationem, qua beatum et foelicem aliquem  \pend
\section*{AD I. PAVL. AD TIM. }
\marginpar{[ p.402 ]}\pstart ab astrorum quidem influxu loa.19.v.11. Rom. 13.V.I.  \pend\pstart Potens, Rex regum, Dominus dominantium) Haec omnia viribus hominum quantumuis magnis opponuntur, vt confidentius constantiusque  in Dei metu omnes, et quisque in suo munere pergamus, coniectis in vnius Dei potentiam incomprehensibilem oculis, et in ea ipsa retentis. De voce Aυναςuς mouet controuersiam Eras. vtrùm potentem significet: an praefectum et praesidem, qui illius suae potentiae et authoritatis specimen in aliquo populo, vel prouincia regenda exhibeat, nimirum, ne Dei potentiam otiosam fingere videamur. Respondet igitur ille non nudam et otiosam potestatem significari: sed negotiosam, vt ita loquar, et quae in alicuius rei regimine se proferat. Dei enim potentia in totius huius mundi, et singularum eius partium regimine se exerit. Quod quidem verissimum est, tamen tam subtiliter et acute non videtur hoc loco philosophari Paulus. Sed omnem omnium rerum potestatem vni Deo asserere voluisse, omnem item maiestatem, dignitatem, ne potentum regum, imperatorum, aut totius mundi splendore exterreamur et dimoueamur a nostro officio et vera ipsius doctrina. Ergo Rex regum hic appellatur. Sic Apocalyp.17.vers.14.19. vers.16, et cap.1. vers.8. Hlaντοκράτωρ tanquam ὀυτοκράτωρ, vt quidam volunt. Videtur autem hoc sumptum ex Psalm.95. vers.3.96 vers.1o. Ex hoc vero ipso ratio intelligitur cur Dominus dicatur dare regna cui velit: et cui item velit ea adimere Daniel.5. vers.14.item cur dicatur omnis potestas esse superne, id est, a Deo, non  \pend
\section*{CAPVT  VI. }
\marginpar{[ p.483 ]}\pstart ne insanos Astrologos iudiciarios probemus.  \pend\pstart Solus) Hoc ipsum primùm non excludit, neque  tollit regna politica. Sunt enim et ea ipsa a Deo. Itaque quando subalterna sese mutuo non euertunt :recte et Deus solus rex agnosci potest, et reges etiam terreni pro regibus a nobis agnoscuntur, quia eos nobis ipse praeposuit, quasi suos in his terris vicarios et legatos in iis. quae ad vitae huius terrenae rationem pertinent. Deinde, quod dicit Paulus, Solus, etsi ad patrem referri videtur, non excludit tamen filium, vel Spiritum sanctum, sed in persona patris contrahit, et proponit, quae generaliter in tota diuina essentia sunt nobis meditanda et consyderanda, quia de filio mentionem fecit Paulus, tanquam de homine. Nam de Deitate Filii quid sentiat idem Paulus extat luculentissime scriptum Romanor.9. vers.5. Atque haec solutio ad sequentem versiculum etiam est acconmodanda.  \pend
\phantomsection
\addcontentsline{toc}{subsection}{\textit{6 Qui solus habet immortalitatem, lucem habitans inaccessam, quem vidit nemo hominum, neque videre potest: cui bonor et robur aeternum. Amen. }}
\subsection*{\textit{6 Qui solus habet immortalitatem, lucem habitans inaccessam, quem vidit nemo hominum, neque videre potest: cui bonor et robur aeternum. Amen. }}\pstart Εʹ̓πεξεργασία est, et valde necessaria commoratio in imperuestigabili Dei maiestate explicanda, atque supra omnem rerum tum visibilium, tum etiam inuisibilium, dignitatem et eminentiam euehenda. Etsi vero Exod.34.v.6. pulchra quorundam Dei attributorum enumeratio continetur, et Psalm.5o. tamen quae hoc loco recitantur, longe splendidius atque apertius, quam vllus alius scripturae locus, Dei ἰξοχιὼ et praestan\pend
\section*{404 AD I. TAVL. AD TIM. }\pstart tiam nobis describunt. Itaque est hiclocus singularis et egregius, vti et ille, qui est sup.cap.1. vers. 17. qui huie similis est. Quatuor autem attributa referuntur hoc versiculo, nempe. I Quod solus habeat immortalitatem. 2 Quod lucem inhabicermaccemiauenn 5 oueui nemo vnquam viderit. 4 Quod illi soli Honos et gloria debetur.  \pend\pstart Ac primum epitheton excutiamus. Dicitur igitur solus Deus habere immortalitatem, ἀθανα- σίαν dixit Paulus. Immortalitas mortem, tanquam existendi finem, excludere tantum videtur: non etiam principium existendi remouere. Itaque non satis magnifice hoc de Deo videri potest Paulus dixisse. Neque enim sequitur, vt quod immortale est, sit aeternum cùm multa existendi principium habeant, quae tamen minime sunt habitura finem. Itaque ἀἶδότητα potius appellant, quam ἀθανασίαν illam diuinae naturae ὐπαρξν, et existendi modum vt est Hebr.7.vers.3. Vnde Paulus Romanor. 16. vers.26. Deum αιώνιον, id est, aeternum appellat, vt est Isai.57.vers.15. Respond. Immortalitatem proprie hoc loco accipi. Itaque dico eam complecti aet ernitatem, id est, et principium et finem existendi plane secludere. Nam quae immortalia dicuntur, neque sunt aeterna, illa improprie talia sunt, quia alieno tantùm beneficio, non autem sua vi et potestate viuunt, subsistunt, et talia durant, quemadmodum Angeli, et humanae animae. Dei enim sustentatione perdurant et subsistunt, non sua ipsorum. Praeclare enim Irae.lib 4 cap 37. Vita hominis, id est, animae, est visio Dei. Quod ab eo sumptum saepe recitat Augustinus.  \pend
\section*{CAPVT  VI. }
\marginpar{[ p.485 ]}\pstart Itaque nec Angeli, nec humanae animae per se sunt immortales. i. sua vi et potestate propria, sed quatenus a Deo in sua natura sustentantur, et conseruantur, tales permanent. Id quod etiam in Timaeo Plato, quanquam Ethnicus, agneuit, et scripsit. Nec tamen propterea Angelos et animos humanos negamus esse immortales: sed tantùm per se negamus tales esse. August. in liber  de Immortalitate animae, et liber  12. de Ciuit. Dei item liber 1.cap.2. de Origine animae ad Hierony. Tract.in Ioan. Euang. 19. et infinitis pene aliis locis. Vti enim corporis vita est anima: ita animae vita Deus est. Ergo solus Deus immortalis est vere propriaque substantia, id est, υσία, non autem κουωνία quadam, quales Angeli et hominum animi, quia perpetuam durationem et finis expertem tribuit Angelis et hominum animis Deus, vt ait Theodorit. Dialog.3. qui est impatibilis. Ipse autem eam a se habet. Bernard. quoque serm. 31.in Cantic. Canticor. hunc locum praeclare explicat, at que idem quod Theodorit. sentit.  \pend\pstart Cur autem solus Deus sit immortalis ratio duplex affertur I quia est incorporeus. Vti enim Deus solus immortalitatem habet, ita incorporeitatem vt ait Bernard.serm.5.6. et 81.in Can. 2 ratio, Quia est immutabilis omnino Deus Iac. I.vers.17. et haec verissima ratio est, quam liber  1.de Trinit. et liber 6. cap.6.item liber 2. contra Maxim. atque liber 7.cap.6. de Genesi ad liter. Augusti. copiose explicat. In omni enim mutabili natura mutatio est mors, quia facit aliquid in ea non esse quod erat: et aliquid abesse, quod prius inerat Angelos autem esse mutabiles ostendit lapsus  \pend
\section*{400 AD I. TAVL. AD TIM. }\pstart eorum qui similis cum aliis naturae fuerunt Item et humanos animos esse quoque mutabiles, casus primi hominis certissime docet, per quem tanta et animi et corporis in deterius immutatio in nobis facta est. Ergo ii quodammodo mortales, quia sunt mutabiles.  \pend\pstart Secundum Epitheton est, Qui lucem habitat. inaccessam. Hic locus est etiam vberrimus, quem tamen breuiter contrahemus in pauca. Primùm vero videndum est, vnde sumptum sit hoc Pauli dictum. Nec enim haec fuit particulariter facta Paulo reuelatio de natura Dei. Plane vero consentit hic locus atque haec descriptio cum ea, quam de maiestate Dei incomprehensibili Prophetae tradunt. Est igitur similis huic locus Psal. 93.veri.I.Ilai.6.veri.2.57.verf. 15. Ezechiel, 1. et10. Tanta enim est haec lux, qua Deus vndique circumseptus effulget, vt Angeli ipsi beati illius fulgore perstringantur, at que faciem suampropterea obtegant, et operiant.  \pend\pstart Sed dissimiles huic videntur alii quidam esse Scripturae sacrae loci. Imprimis ii omnes, in quibus dicitur Deus habitare in nube, vel ex turbine et caligine loquutus fuisse, vti Exod.14, et 20. Leuitic. 16.Iobus 38.vers.1 Item vbi dicitur Deus fuisse cum leui vento, non cum igne, quemadmodum cùm apparuit Eliae 1. Regum. 19. vers 12. Respond. In nullo iustiusmodi scripturae loco veram Dei sedem describi, quemadmodum hic: sed ea tantùm signa et Symbola declarari, per quae hominum oculis sese parefecit in hoc mundo pro re rum, temporum, personarum, quibuscum illi res erat, conditione, vsu, et diuersa ratione. Ergo ill  \pend
\section*{CAPVT . VI. }
\marginpar{[ p.487 ]}\pstart Praeterea obstat huic loco, quod omnino sequi inde necesse est, si vera sit haec sententia, nimirum Deum, qui est maiestas illa infinita, certo tamen loco contineri, et concludi, quia dicit Paulus eum ὀικεῖν, id est, habitare, at que lucem quidem. Quod sentire plane est blasphemum. Hoc etiam argumentum sensit Chrysost. cuius verissima est haec responsio, nimirum de Deo ista dici ex nostrae cognitionis imbecillitatis sensu, et modo per ἀνθρο ποπάθειάν quandam. Neque enim, quantumuis nitatur humanum ingenium, quicquam satis magnifice et digne de Dei maiestate vel concipere vnquam vel dicere potest. Ergo proposita in superiori epitheto Dei maiestate, nunc ipsius solium lucidissimum et splendidissimun describit, vt vndique effulgeat nobis augusta, ex omni parte sacrosancta, et veneranda Dei illa gloria animis nostris obuersetur, et splendeat. Ex quo fit, vt et lucis et sedendi voces hoc loco fint metaphoricae, non autem natiua et primaeua significatione sumendae. Nec enim, vt praeclare docet etiam hunc locum explicans August.in Psal.33. et in epist. ad Italicam Deus loco comprehendi vllo poroiaueciriniodo, qui locum omnem ipse superat, capit et excedit. Ex eo autem, quod ad literam, vt aiunt, hic locus, ac quidam similes alii accepti sunt, insulsum et plane puerile commentum suum de loco, vbi Deus tanquam in recessu et specu quadam abditissima habitet, finxit atque  excudit impius Mahumetes.  \pend\pstart Neque etiam hic quaerendum est, vtrum lux illa, in qua Deus, qui est aeternus, habitare dici\pend
\section*{D I. TAVL. AD TIM. }
\marginpar{[ p.400 ]}\pstart tur fuerit, antequam lux illa de qua Genes.1.facta mentio est, crearetur et extitisset. Neque enim hic lucis vocabulum natiuo significato sumendum est: sed pro eo splendore, maiestate, statu qui omnem admirationem cogitationemque humanam longe superat, habetque optabilem et laetissimam meditationem, qualis est nobis ipse lucis aipectus, vt ex Ilai.57. vers.15. satis perspici potest et quidem tantam, vt ad eam humanae etiam memtis acies caliget.  \pend\pstart Ergo huius epitheti duplex est sensus. Primus quidem, vt omnem nobis diuinae maiestatis curiosius perscrutandae pruritum ( qui ex magna quadam vanitate nasci solet) adimat, ac tollat. Praeclare enim Bernar. serm. 62.in Cantica opprimitur scrutator maiestatis Dei: non tamen voluntatis. Nam vt ii flammis statim correpti et consumpti sunt, qui tres illos adolescentes Abdenagum, Mizaelum, et Sydracum in fornacem coniecerunt. Daniel. 3. Vti puniti sunt Betsemitae, quod in arcam Domini inspexerint 1. Samuel.6. Sic consumuntur et obruuntur Dei maiestate, qui in eam curiose, et aliter, quam oportet, intueri atque inquirere volunt. Denique  hoc ipsum est, quod ait alibi Paulus, Dei maiestatem, atque iudicia esse imperuestigabilia atque ἀνεξερευιητα. Ergo praeclare Bern.ser.31.in Can. ad Deum accedendum, est, ait, nonirruendum, ne irreuerens scrutator maiestatis opprimatur ab ipsius gloria et splendore. In quo tamen illud est diligenter obseruandum quod idem Bernard. dixit serm.62.opprimi scrutatorem maiestatis Dei:non item voluntatis. Nam eam voluntatem Dei, quam quidem nobis notam  \pend
\section*{CAPVT  VI.. }
\marginpar{[ p.489 ]}\pstart esse voluit, inquirere diligenter et licet et concessum est, vt docet exemplo Patrum Pet.1. epi.1.V.II. Dicitur haec Dei lux inaccessa. Ne proprio marte propriisque  ingenii nostri, quantumuis viuidi, viribus ad eam accedere tentemus. Qui enim sic ad eam adire tentant, excaecantur illius fulgore. Sed prius a Dei spiritu nos illuminari petamus et patiamur. Nam illuminatis patet facilis ad illam Dei maiestatem et lucem intuendam accessus: et Deus ipse prius eos illuminat, a quibus se videri et sua arcana cognosci, atque tanquam conspici vult. Ps34.v.6.Cùm autem ad eam venire volumus, neque  locorummutatione, neque corporis agitatione ad Deum appropinquamus: sed claritatibus, incorporeis et spiritualibus, suiisque muneribus, qualis est fides, vitae sanctitas, zelus verus ad eam accedere nos oportet, vt ait Bernardus, et exBernardo Anselmus.  \pend\pstart Porro secundus huius epitheti vsus est. Quod cùm dicitur Dei domicilium lucidum esse: imo esse lux ipsa, culpa omnis nostrae de Deo ignorationis in nos ipsos transfertur. Nec enim Deus in tenebris habitat, vt agnosci aut videri non possit: sed in luce splendidissima. Ergo quod eum, illiusque admirabilem maiestatem cernere non possumus, id aciei nostrae imbecillitatis, et ingenii tenuitate fit, non ipsius quidem loci, in quo Deus est, obscuritati tribuendum est. Omnia enim quodammodo Dei ipsius maiestatem nobis repraesentant, et illustrant, quocunque conuertamus oculos, vt docet Paulus Rom.1.vers.2o. Sed imprimis verbum Euangelii, in quo se nobis in Christo perspicuum et aspectabilem reddit, eum  \pend
\section*{AD I. PAVL. AD TIM. }
\marginpar{[ p.450 ]}\pstart nobis lucidissime demonstrat. Denique ex hoc ipso loco intelligitur, quam vanum et blasphemum fuerit Manichaeorum delirium, qui deum saum, quem tenebrarum Deum esse et alterum principium fingebant, in tenebris spississimis, non in luce habitare docuerunt.  \pend\pstart Tertium epitheton est, quod ait Paulus, Deum nemo vnquam vidit. Ergo vt lucidustotus est : ita est inuisibilis Αόρατεςvt ait Paulus sup.1. vers.17. Haec epitheta quia sunt diuinae essentiae propria attriouta, iuin leuuio et umgemer obseruanda a piis omnibus, vt ex iis Dei maiestas tanquam aliqua delineatione et rudi pictura in animis nostris informari possit.  \pend\pstart Similes loci huic plures in tota scriptura sup. 1.vers.17. Ioan.1.vers.18. Exod.33. vers.20,21,22. Genes.32.vers.3o.  \pend\pstart Dissimiles autem alii quoque scripturae loci, et in speciem cum hoc pugnantes. Ac primus quidem is, qui est apud Paulum 2. Corinth.3.vers.18. vbi ait nos sideles per lucem Euangelii etiam retecta facie Dei gloriam intueri. Secundus autem quod Deus saepe patribus visus fuisse dicitur, vt Abrahamo, Isaaco, Iacobo, Mosi, et aliis plerisque  Respondeo Omnia illa quae obiiciuntur, κατότι esse accipienda, non eo significato et modo, quo hic loquitur Paulus. Agit enim de ipsius diuinae essentiae intuitiua, vt loquuntur in scholis, visione, quo modo a nemine vnquam mortali visus est, non autem de Symbolica repraesentatione aut explicatione. Ac ne quidem cùm videbatur Christus in his terris, qui Deus conspicuus fuisse dicitur fup3 vers.16. ipsius diuina essentia oculis  \pend
\section*{CAPVT VI.. }
\marginpar{[ p.491 ]}\pstart humanis cernebatur: sed illius tantùm humanitas conspiciebatur, quae velut velum, certe augustissimum, Deitatis Christi templum fuit. Hebr. 10.vers.20. Ergo quod Dei faciem et gloriam in Euangelii luce et praedicatione iam intueri dicimur, id dicitur comparate ad patres, qui sub lege, vmbris et caeremoniis vixerunt, vt ex eo Exodi loco, qui est cap.33.vers.20, et 23. apparet. Patres enim posteriora Dei in legis vmbris viderunt: nos autem in Euangelii luce et implemento videmus faciem ipsam, et ipsam vmbrarum veritatem: non tamen ipsam Dei essentiam, qualis in sese est, propterea videmus. Sed et cùm Abrahamus scribitur vidisse Deum, id certe metonimicôs intelligendum est, quia signa et Symbola ocuns cernebat, per quae Deus se praesentem illic esse testabatur, ipsam autem Dei essentiam non videbat. Sic in columba visus est Spiritus sanctus a Ioanne Baptista Luc.3.vers.22. Iren.libus  4. cap.37.explicat.  \pend\pstart Quod autem haeretici quidam senserunt ipsam Filii Dei θεότητα fuisse per sese semper visibilem, quiu natrious apparuitsetlair aite quamn incarnaretur Christus, nempe quando cum illis est colcutus. Itaque quod hoc loco dicit Paulus, Deum videri non posse, de sola Patris persona et Deitate accipiendum, plane (inquam) hoc falsum est: et praeclare ab Augustino refutatum in liber 8.de Genesi ad literam.  \pend\pstart Quaeritur autem vtrum vox haec ὀδεὶς ἀνθρώπων ad homines tantùm sit referenda, et eos quidem adhuc mortales vel hic degentes, an Angelos etiam, imprimis beatos complectatur: atque pios  \pend
\section*{AD I. PAVL. AD TIM. }
\marginpar{[ p.492 ]}\pstart defunctos et iam in coelum receptos quia de Angelis scriptum est, Matt.17.eos videre faciem Patris, id est, Dei: et de piis beatis Matth.5.vers.8. Mundos corde visuros esse Deum. Hic praetermitto vanas et curiosas Scholasticorum quaestiones, quae de Dei visione fiunt in liber 4. Sentent. Lombard. distinct. 49. Item non afferam ea, quae prolixa epistola ad Paulinam et ad Dardanum de Videndo Deo disserit Augustinus, itemque liber 1. de Trinitate cap.6. Sunt enim ex iis quaedam ἀδηλα, quaedam etiam curiosa, quod tanti viri pace dictum sit. Neque item dicam quod Ansel. Hominum nomine hic carnales duntaxat comprehendi, non etiam pios, renatos, et in Dei obsequium iam formatos homines. Haec enim interpretatio minime vera est. Respondeo igitur superiori quaestioni breuissime ex Theodorito in Polymorpho, itemque ex Bernardo serm. 81. in Cant. Ne ab Angelis quidem ipsis beatis Dei essentiam cerni posse, nisi per quandam συγκτάΒασν, et quasi sui demissionem, vt loquitur Theodorit. Soli enim sibi Deus est et visibilis et comprehensibilis. Ergo multo minus a piis in coelum traductis ipla Deι οσια, qualis quataque est, et ipsa per se videri potest, si nequit ab Angelis. Neque certe tam de hoc nobis est quaerendum, quam via ipsa per quam peruenire ad eum possimus, et scrutanda et tenenda nobis. Est autemilla via vera in Christum fides, Vitae puritas et sanctificatio item pietas charitasque syncera.  \pend\pstart Addidit autem hoc epitheton Paulus tum vt esset superioris tanquam explicatio quaedam, tum etiam ne ex harum rerum, in quibus aliqua Dei  \pend
\section*{CAPVT VI. }
\marginpar{[ p.493 ]}\pstart praesentia cernitur aspectu, diuinae maiestatis emi nentiam metiamur et aestimemus. Denique vt sit haec tanquam hypophora quaedam, et tacita responsio ad eorum obiectiones et dubitationes, qui idcirco Deum esse vel negant prorsus, vel dubitant, quod ab ipsis oculis illius ὀσία minime cernitur. Itaque dixit Paulus neminem vel vidisse Deum antea, vel nunc, vel in posterum etiam videre posse. Caro enim et sanguis neque sunt, vt ait Bernard. diuinae essentiae capaces, neque capabiles.  \pend\pstart Quartum epitheton sequitur in quo totius superioris diuinae naturae descriptionis fructus, scopus, et finis continetur nempe, vt vni Deo laudem omnem, gloriam, et potentia omnium rerum tribuamus, idque  non ad tempus aliquod tantùm: sed in aeternum, et in secula et vltra. Haec particula sup. cap. Medouiquupurocepheutocle, ce lape recurrit, idcirco a plerisque exposita est, et a nobis latiùs non traderur. Vox Amen ipsius Pauli assensom in eas omnes Dei laudes declarat, et eas quasi sialla queaulalandacaa teuiniaes,  \pend
\phantomsection
\addcontentsline{toc}{subsection}{\textit{17 Iis qui diuites sunt in hoc seculo, denuntia ne efferantur animo, neque spemponant in diuitiis incertis, sed in Deo viuo, qui praebet nobis omnia copiose ad fruendum. }}
\subsection*{\textit{17 Iis qui diuites sunt in hoc seculo, denuntia ne efferantur animo, neque spemponant in diuitiis incertis, sed in Deo viuo, qui praebet nobis omnia copiose ad fruendum. }}\pstart Εʹπάνοδοςest et regressio. Intermissam enim de diuitum, et diurtiarum cura, vsu, totóue hoc argumento exhortationem repetit. Idque duplici de causa. Primùm quod, vti ex Iac.2. apparet, facile sint in reliquos iniuriosi diuites, iamque in Ec\pend
\section*{AD I. PAVL. AD TIM. }
\marginpar{[ p.494 ]}\pstart clesia quidam tales existerent. Praesertim autem in amplis et opulentis ciuitatibus saepe id accidere et cernisolet. Hoc autem et ipsis diuitibus exitiosum est: et caeteris in Dei Ecclesia praebet offendiculum. Cui malo remedium afferendum fuit. Atque tale quid iam tunc in Ephesina Eccieila coneigme, veriime mit ex eo sit, quod non frustra has praeceptiones de diuitiis repetit Paulus.  \pend\pstart Secunda vero huius regressionis causa est, ne, cùm superioribus versiculis tam acerbe in diuitiarum appetitum inuectus esset, diuites Christiani omni spe et salutis, et bene vtendi suis diuitiis se destitutos esse statuerent et quererentur. Itaque eas abiicerent, tanquam omnino cum Christiani nominis professione pugnantes et ὀσυζύγος, Deóque exosas. Neque enim supra diuitias perse damnauit Paulus, quemadmodum recte Augustinus scribit: sed diuitiarum morbum: neque praecipit Paulus diuitibus, vt suas diuitias abiiciant, quod insanus ille Antisthenes philosophus et post eum Aristippus fecit: sed vt iis bene vtantur, tum ad Dei gloriam: tum ad suam ipsorum salutem, et proximorum commodum et aedificationem. Haec igitur causa est repetitae huius exhortationis, vt hoc tam necessario praecepto, tanquam praeclaro colophone et fine, totam hanc epistolam tandem clauderet Paulus.  \pend\pstart Ergo quid haec praeceptio contineat sigillatim nobis expendendum est. Duo vero obseruemus hoc loco, et Quibus praecipiat, et Quid praecipiat. Praecipit igitur Paulus diuitibus, et iis quidem, qui in hoc seculo diuites sunt. Quae verba  \pend
\section*{CAPVT  VI. }
\marginpar{[ p.495 ]}\pstart suo momento non carent. Primum enim ostendunt esse duplex diuitum genus. Vnum eorum, qui in hoc seculo et rerum terrenarum ratione solùm diuites et abundantes censeri possunt. Hi sunt, quos vulgus hominum diuites appellat, et de quibus agit hoc loco Paulus. Aliud est diuitum genus, eorum nempe, qui in futuro seculo vel rerum coelestium spiritualiumque  ratione et copia sunt diuites, opulenti, abundantes, de quibus loqui tur Iac .2. v.5. Paul. 2.Cor.8.v 9. Apoc.2.v.9. Qui diuites fide, et bonis operibus diuites, et haeredes in regno Dei dicuntur, quales sunt omnes vere filii Dei, vtcunque in hoc mundo sint egentes, atque etiam Iro ipso, vt est in Prouerbio, pauperiores. Qualis fuit Iobus in mediis afflictionibus: qualis item fuit ille languens Lazarus, et vlceribus plenus et fame oppressus, de quo Luc.16. Praeclare enim Ambrosius epist. 1o.omnis Christianus (vt antea diaiihusjee lapiens est, e diues, rrorum autem spiritualium diuitum diuitiae, vt recte dixit August. in serm. de verbus  Euangel. sunt non quidem pecunia, sed iustitia: quae certe vere sunt diuitiae, neque amitti possunt, quemadmodum etiam docet D. N. Iesus Christus.  \pend\pstart Sed addit quoque alia de causa haec verba Paul.in hoc seculo, vt ipse contemptim de his diuitiis loqueretur, quibus minime fidendum esse demonstrat, quasque res esse vanas, incertas, caducas pronuntiat. Itaque vt diuites ipsos melius deprimeret, diuitias eorum huius seculi appellat, id est, euanidas et fallaces, et instabiles, qualis est huius ipsius seculi et mundi conditio, et eorum omnium, quae in eo sunt. Qua ratione etiam voce  \pend
\section*{AD I. PAVL. AD TIM. }
\marginpar{[ p.490 ]}\pstart σράγ sελλε vsum esse Paulum putant, non autem voce λέγε, vel διδασκο, quod diuitum ingenium subiiciendum sit, et contundendum, cùm nimiùm per se semper insolescat. Certum autem est hoc vocabulo seriam admonitionem a Paulo significari quae diuitibus a pastore, et crebro quidem, subiierdeotc at medicauis  \pend\pstart Ergo duo maxime vult effagi a diuitibus, superbiam animi nimirum, et spem vanam, quae ex opum abundantia ab illis concipitur. Vtrunque sane vitium non modo diuitibus familiare est: sed comnatum et pene a diuitiis inseparabile, cumiisque coalescens et ab incunabulis crescens, vt qui diuitias accipiat, simul quoque eademque  manu haec vitia quodammodo acceptare videatur. Quorum vitiorum vana spes est animi superbiae fons et mater: Ex illa enim haec producitur, tanquam foetus infoelicissimus. Id quod exemplo Ezechiae regis, quanquam alias pii, comprobatur, qui nimia sua foelicitate intumescens sese nimium extulit. Isai.39. Neque id tamen ipsarum opum, aut rerum, in quibus illae diuitiae consistunt, vitio, aut culpa accidit: sed nostra vitiositate dumtaxat, et humani ingenii vanitate, quae per peccatum nobis infusa est nimirum.  \pend\pstart Ne haec vitia rebus ipsis a Deo esse iniecta  \pend
\section*{CAPVT  VI }
\marginpar{[ p.497 ]}\pstart credantur et ita Deo opum ipsarum creatori vitium impie ascribatur, sed vt nostrae corruptionis et reuemonis nae reuitarie en isdamo vis et magnitudo intelligatur, damneturque.  \pend\pstart Superbiam igitur ex animi superbi cogitatione describit Paulus dum eam ὐαηλοφρονίαν esse declarat idem Paulus Rom.12.vers.16. Recte igitur quidam dixit, Vt vter, inflati sunt superbi. Si attendis eos, falleris: si appendis, vanos et vacuos reperies, elatos tamen semper, et magnifice de se sentientes comperies.  \pend\pstart Alterum vitium est vana spes, cuius causa est, quia diuitiae sunt incertae, imo fallaces, et deceptrices, vt ait Christus, Matth.13.vers.22. Mar.4.v. 19. Neque enim cùm haberi creduntur, tenentur sed statim effluunt. Neque cùm tenentur stabiles Iunt, vt eripr nobis ab aliis non pofsint, aut apud nos ipsae per se perire, et consumi nequeant Luc 12. De quo argumento infiniti sunt tum in Prouerbiorum liber  11.vers.28.13.vers.8. tum etiam in Ecclesiaste loci, qui sunt cum hoc conferendi.  \pend\pstart Nec obstat huic loco Prouerbiorum cap. 14. verneie voriten diuitia dicuntur esse diuitum corona. et Ecclesiastic.7.vers.12. In vmbra argenti requiescet homo. Respond. Additum est Sapientum, id est, horum esse coronam qui recte diuitiis vtuntur. Non sunt autem stultorum diuitiae iis ipsis corona, quia vt est Eccles. 4.vers.13. praestat pauper sapiens rege et diuite stulto.  \pend\pstart Caeterum his vitiis contrarias virtutes iam videamus, quas diuitibus consectandas esse docet Paulus. Sunt autem hae tres. I Spes siue fiducia in Deo viuo. 2 Benignitas. 3 Facilitas siue  \pend
\section*{AD I. PAVL. AD TIM. }
\marginpar{[ p.330) ]}\pstart comitas morum.  \pend\pstart Pulchra vero est haec antithesis inter Spem et fiduciam in incerto diuitiarum, et inter Spem siue fiduciam quae ponitur in Deo viuo, et asfatim omnia nobis praebente, non tantùm ad tenuem et necessarium vsum, sed etiam ad liberalem et opiparum. His autem epithetis, quae hic Deo tribuumtur, non tantùm verus Deus ab omni humani ingenii figmento, quod pro Deo habetur, non etiam solùm ab omnibus idolis iam receptis pro diis secernitur: sed etiam verissima nostrae fiduciae in eum collocandae ratio, causaque explicatur, nenpe, quod sit ipse omnis vitae fons, siue tempotariae, siue etiam aeternae. Deinde, quod quaecunque ad vitam hanc etiam tuendam sunt necessaria, abunde et copiose ille vnus omnibus nobis suppeditat, et subministrat, vt ab eo solo pendere debeamus, quacunque ratione vel viuere, vel vitam ipsam tueri et conseruare cupiamus. At qui hae caulae sunt, cur oiuites conndant suis opibus et diuitiis, quod per eas solas se vitam hanc tueri et sustinere posse credant. Quorum ratio non tantùm ex hoc Pauli loco refutatur: sed etiam ex ea parabola, quain amert enriltus Luc 12.vers.16. et sequenti.  \pend\pstart Vnus Deus) appellatur, non tantùm quod et in sese, et per sese viuat, sed quod omnbus, quae viuunt, vitam largiatur. Quaecunque enim viuunt, illius beneficio et dono viuunt, tam Angeli, quam vermes ipsi, vt ait Augustinus.  \pend\pstart Praebet idem Deus omnia, idque πυσίως, id est, affatim et copiose:sed (vt ait Iacobus) etiam be\pend
\section*{CAPVT VI. }
\marginpar{[ p.499 ]}\pstart nigne. In quo pulcherrima est annominatio ad vocem πλεσίων, quos exhortatur Paulus. Non sunt enim illi diuites vere et proprie: sed Deus est dides, qui et omnla ponidet, et abunde cunctis latgitur et opulenter.  \pend\pstart Quod autem additur, Nobis, ad suscepti argumenti rationem est accommodatum, quod ea vitia, quae taxat hic Paulus, in solos homines cadant: quódque nostra causa tantùm institutae sunt hae praeceptiones. Neque vero restringit Paulus admirabilem illam Dei benignitatem vel ad solos pios et fideles homines, cùm etiam impii et reprobi eam sentiant in hac vita, vt est Matt. 5.vers.45. vel etiam ad solum genus humanum, cutir etiaili ortita aminantia pareat Deus, illisque victum suppeditet Psal 104 Matth.6.vers.26. Psa. 36.vers.5. Imo infinita Dei beneficia in omnes mundi huius et creaturas et partes diffunditur, et se exerit.  \pend\pstart Ε’ίς ἀπόλαυσν, id est, ad fruendum. Quia distinguit Augustinus lib 1.de Doctr. Christ. inter frui et vti: et fruitionem ad res diuinas: vsum ad terrenas et caducas refert: idcirco sic hunc locum explicat, et post eum Beda, vt de vita aeterna bonisque coelestibus tantùm accipiendum putet, quia fruendi verbo vsus est Paulus. Sed primum falsa est ista distinctio Augustini (pace tanti viri dictum sit) cùm aeque et proprie, et frui, et vti dicamus rebus siue terrenis siue coelestibus: etsi apud Iurisconsultos aliquod discrimen statuitur inter vsum et fructum: et inter ius vtendi, et ius fruendi re aliqua, quae iurisconsultorum distinctio ad Theologos minime pertinet. Deinde in  \pend
\section*{AD I. PAVL. AD TIM. }
\marginpar{[ p.309. ]}\pstart verbis Graecis, id est, inter χρῆος et Α'πολάυη nulla huiusmodi differentia reperietur obseruata, vt χρη̃͂ος tantùm de bonis terrenis: ἀπολάυον de bonis coelestibus tantùm accipiatur. Denique hic ipse locus ostendit, de bonis terrenis praecipue agi, in quibus diuites huius seculi spem ponunt, quae tamen a solo Deo dari docet Paulus hoc loco. Quid igitur sibi vult haec dicendi ratio? Nimirum, non ea tantum nobis a Deo suppeditari quae necessaria sunt ad stricte, anguste, duriter, et tenuiter vitam hanc tuendam et sustentandam: sed benigne, profuse, largiter et liberaliter, multa etiam nobis a Deo praeter necessarium vsum tribui. Hoc dicunt septuaginta interpretes in veteri Testamento ἀωμεῖν Mautem Tullius liber  1. de Legibus  Prodigere, vt ex hac tam liberali Dei in nos benignitate intelligamus non tantùm quam luitiuo hDeo iperandum sit, quam in diuitiis: sed quam opulentius et liberalius a Deo omnia accipiamus, quam vlli diuites a suis opibus habeit poinn. E mhique ex hoc loco immensae Dei benignitatis in homines commendatio fieri debet.  \pend
\phantomsection
\addcontentsline{toc}{subsection}{\textit{18 Vt aliis benefaciant, vt diuites sint operibus bonis, faciles ad impartiendum, facilis conuictus. }}
\subsection*{\textit{18 Vt aliis benefaciant, vt diuites sint operibus bonis, faciles ad impartiendum, facilis conuictus. }}\pstart Sequuntur aliae virtutes duae, quae superioribus etiam vitiis, quae in diuitibus notauit Paulus, opponuntur. Hae autem sunt, Benignitas, et Comitas. Ac benignitas quidem siue beneficentia est verissimus diuitiarum fructus. Comitas vero velut condimentum ipsius beneficentiae, certe  \pend
\section*{CAPVT . VI. }
\marginpar{[ p.501 ]}\pstart superbiae illi, quam reprehendit supra Paulus, virtus est plane contraria atque repugnans.  \pend\pstart Beneficentiam autem variis verbis designauit tum ad illius virtutis commendationem: tum etiam ad explicationem. Nam eam commendant hae voces Α'γαθοθγεῖν, et Πιλετεῖν ἐν εφγοις καλοῖς. Explicant istx αμεταδότοςτθ).  \pend\pstart Itaque hic bona opera designantur, non tam quidem omnia in genere illa quae scriptura Bona opera vocat, et praecipit: quanquam diuites Christiani iis omnibus animum applicare debent: quam quae Charitatis et Elemosynae opera sunt, et huic argumento conueniunt: quod ea sint praecipua, in quibus diuites sese debent exercere, et quorum praestandorum causa diuitias a Domino acceperunt. Omne autem proximi subleuandi et beneficentiae genus complectitur Paulus, cùm vult ita diuites suas opes collocare et dispensare, vt ditescant abundentque in omni genere bonorum operum, non vt vnius aut alterius tantùm beneficii in proximum exhibitione sibi satisfaciant, quasi sint satis officio functi; sed congerere cumulare, et addere eos oportet beneficia beneficiis minime propterea fathiscentes, quia pluribs subueniant: vel quia vni et eidem saepius iam contulerunt donaruntque aliquid: in quo ipso praeclara est annominatio, et argumen tum a coniugatis a Paulo perstrictum. Oportet enim Πλυσίος πλυτεῖν, id est, diuites congerere: sed in acerui genere, quem fieri iubet Paulus, differentia esta Illi enim vt plurimum congerunt scelera sceleribus, vt ditescant: at Dei spiritus vult eos congerere et cumulare bona opera bonis operibus, vt illis  \pend
\section*{AD I. PAVL. AD TIM. }
\marginpar{[ p.JUa. ]}\pstart ipsis suae diuitiae sint salutares. Pro ἀγαθορργεῖν, legit Ambrosius ὀθελοεργεῖν, quod huic loco et menti Pauli non conuenit. Bona igitur opera vocat Paulus ea, per quae in proximum beneficium confertur, et illi bene fit, quemadmodum idem loquitur Galat.6.vers.9.1o. quanquam verbi Graeci significatio etiam ad alia operum genera pertinet. Bona vero opera sunt eleemosynae, quia a bona radice, id est a fide et vera Christianaque charitate nascuntur. Alias ipsa pecuniae datio quae fit pauperi, per se potest esse opus indifferens, vt docet Paulus 1.Corinth.13.vers.3. vel etiam malum, vt si pauperi dem, quem inescare cupio, vt hac liberalitate postea persuasus a me facilius iuguletur, vel laedatur, exemplum est Matth.6.vers.1.2. 3. eleemosynarum quae damnantur. Ergo non tam externum opus in istis rebus est opus illud quod bonum appellatur, quam fons, origo, finis et scopus ipse, vnde illae eleemosynae proficiscuntur: et ad quem tendunt. Quod idcirco notandum est, quia homines infideles et hypocritae haec opera externa longe potiora melioraque putant, quam veram in verum Deum fiduciam, carnis mortificationem, spiritus renouationem, quoniam speciem maiorem sanctitatis et humanitatis prae se ferre videantur. Bonum enim opus metiuntur ex eo quod est et apparens, et alteri vtile et commodum. Bona sunt quidem haec opera quatenus ex vera fide et charitate fiunt: et ad Dei gloriam et obsequium proximique aedificationem penitus referuntur.  \pend\pstart Quod addit, faciles ad impertiendum, ostendit verissimam liberalitatis, quae in diuitibus esse de\pend
\section*{CAPVT vis }
\marginpar{[ p.2 ]}\pstart bet, vim et naturam, quae in eo est, vt non tantùm velint benefacere: sed reipsa benefaciant, quia dare possunt. Nam in iis, qui possunt praestare aliquid, vera charitas non est in lingua, neque  tantùm in affectu ipso largiendi:sed in ipso opere et praestatione, quemadmodum Ioannes et Iacobus docent I. Ioan.3.vers.18. Iac.2.vers.15.16. Qui nihil habent, quod donent, propter rerum omnium inopiam, qua premuntur, in iis syncerus bene faciendi et largiendi affectus et commiserationis plenus, est sane laudabilis, vti in Petro et Ioanne fuit Actor.3.vers.6.  \pend\pstart Denique sequitur tertia virtus, cùm Christianis omnibus, tum vero diuitibus imprimis consectanda et digna, Nempe comitas, communicabilitasque siue affabilitas. Hanc enim descripsit Paulus voce κοιρωνικές. Haec benignitas nostra est velut condimentum et lepos quidam: et omnis vitae nostrae commendatio, atque conuersationis iucunditas, quae alios homines ad nos allicit, fastumque omnem animi et superbiam deiicit, qua plurimum diuites laborant.  \pend
\phantomsection
\addcontentsline{toc}{subsection}{\textit{19 Recondentes sibiipsis fundamentum bonum in posterum, vt apprehendant aeternam vitam. }}
\subsection*{\textit{19 Recondentes sibiipsis fundamentum bonum in posterum, vt apprehendant aeternam vitam. }}\pstart Α'ιτιολογίά est et commendatio vsus dluitiarum imprimis autem Beneficentiae, quae maxime, vti supra docuit, diuites decet. Est autem ducta haec siue ratio, siue commendatio a fructu siue effecto. Nimirum, quod illa sit vitaẽ aeternae consequendae certissimum testimonium, et tanquam arrhabo, et pignus quoddam a Dei spiritu im\pend
\section*{AD I. PAVL. AD TIM. }
\marginpar{[ p.304. ]}\pstart missum. Cur autem cùm tria quaedam a diuitibus exegerit, in vna praesertim beneficentia commendanda immoretur, haec ratio est, quod cùm natura sint homines et tenaces, et auari, difficillime ad exercendam munificentiam erga alios adducuntur, et de ea virtute persuadentur, adeo vt facilius diuites, et fiduciam in Deum: et erga homines comitatem quandam simulent, quam beneficentiam, quoniam nisi opere ipso et eleemosynae praestatione appareat, solo ore, solaque animi compatientis et miserantis protestatione non potest fingi, aut cuiquam a diuitibus persuaderi. Quanto igitur maior tenacitas cum ditescendi, aut diuitias retinendi animo coniuncta est:tanto magis fuerunt ad liberalitatem et benignitatem diuites exhortandi. Itaqueconsulto et de industria in hac virtute praecipue commendanda et praedicanda insistit Paulus.  \pend\pstart Vtitur autem ea dicendi ratione, ex qua non tantùm optimum et praestantissimum beneficentiae praeconium auditur: sed aculeus quidam ad eam animis iniicitur etiam auarissimis, cùm nempe ait, eos qui benefaciunt, thesaurum sibi recondere, eumque et firmum, et bonum. Itaque  accommodat Paulus suam orationem ad diuitum affectum, qui in eo toti sunt, vt thesaurizent sibi. Id quod eos facturos esse pronuntiat, et longe vberius, vtilius, certius, meliusqs opes suas largiendo pauperibus quam opes congerendo. Nam et in vita aeterna thesaurizant stabilem aceruum. Itaque  praecla ra est haec Augustini sententia et obseruatio ex hoc loco in serm.de Verbus  Euangel.5 Non enim quia dico, vt facile tribuant, eos spoliar  \pend
\section*{CAPVT  VI }
\marginpar{[ p.5 09 ]}\pstart volo, non nudare illos volo, non inanes relinquere. Doceo lucrum, cùm ostendo, vt thesaurizent sibi. Non dico, vt perdant, sed quo migrent ostendo. Haec Augustinus. Ergo qui benefacit et largitur proximo, ille thesaurum coelestem sibi congerit, condit, et coaceruat, omni terreno et numanothelauro ionge praeicantiorem. Similis locus Luc 12.vers.33. Nam qui dat pauperi foeneratur Deo, vt ait Solomon.  \pend\pstart Appellat autem has ex benignitate nostra erga alios comparatas opes thesaurum coelestem, itemque Fundamentum, et Bonum quidem. Fumdamenti vox opponitur tum incertitudini, tum crialiritugnituil terrenarum opum. Sunt enim non tantùm incertae: sed, etiamsi teneantur caducae, quia consumuntur, et ipso vsu pereunt. At coelestes opes neque auferri possunt, neque consumuntur vnquam.  \pend\pstart Bonum autem appellat. Quae vox oponitur et terrenarum opum fallaciae, et miseriis, quas eaedem opes pariunt in animis diuitum, tum vt quaerantur, tum vt conseruentur. Itaque non tam commoditas quaedam, quam incommoditas magna, nisi bene iis vtamur, in iis opibus inesse videtur. Vnde proprie non videntur appellendae Bona quaedam, sed impedimenta, et laqueus.  \pend\pstart Quod addit seipsis et in posterum suaiμρά Cαιnο̃ caret. Nam sua quaerunt diuites. Sibi autem, vt ait Paulus, has coelestes opes congerant. Itaque  eo in benefaciendo largiendóque alacriores et promptiores esse debent. Cùm autem ait, id in posterum, appariturum, omnem eorum cogitatonem ab huius vitae, et status intuitu auertit, et  \pend
\section*{AD I. PAVL. AD TIM. }
\marginpar{[ p.J00. ]}\pstart ad coelestem vitam traducit. Denique tollit adimitque omnem querimoniam, si fortasse conqueruntur et aegreferant se neque in hac vita, neque statim has coelestes opes et thesauros percipere. In futura enim vita promittuntur illis non in hac.  \pend\pstart Denique ipsum pulcherrimum beneficentiae nostrae fructum exponit iam aperte, quem nenpe vocat apprehensionem vitae aeternae. In quo duo quaeruntur, Primum quomodo eleemosyna et liberalitas dicatur esse modus vitae aeternae consequendae. Res. Quia sunt hae virtutes et haec opera bona, verae fidei, veraeque charitatis in nobis inhabitantis certa testimonia, effecta et fructus 2. Pet.1.vers.1o, Iacobus 2. Electionem enim nostram confirmant et ratam faciunt. Sunt igitur signa viae et itineris illius, per quod in coelum nos ituros esse Dominus pronuntiauit: non autem sunt haec opera causa vel pretium vitae aeternae, quoniam illa est gratuitum Dei donum, vt ait Paulus Romanor.6, vers.23.  \pend\pstart Secundo loco quaeritur. Quomodo hoc sit conciliandum cum iis locis scripturae, vbi non ex operibus, sed ex fide nos saluos fieri dicimus Romanor.3.et 4.1.Pet.1.vers.9. Respond. Quod haec bona opera, quae sunt eleemosinae, neque hic constituantur, neque hic commemorentur tanquam causa ipsa vitae aeternae vel per sese, vel in nobis: sed tantum velut signa, testimonia, fundamentaz ex quibus intelligimus nos ex eorum numero esse, quibus vita aeterna a Deo gratis promissa est: quatenus nimirum cùm bona opera in vniuersum, tum haec imprimis sunt adoptionis no\pend
\section*{CAPVT  VI. }
\marginpar{[ p.]}\pstart strae, et gratuitae Dei de nobis electionis signa et consequentia, itemque fidei nostrae certissimi fructus, vt docet Iacobus 2. Per se enim Eleemosinae, quantumuis magnae non possunt esse pretium vitae aeternae, imo nullo modo vllius coram Deo, nisi in Christo, sunt pretii. Sunt enim res nimium exiguae et nullo modo cum regno Dei comparandae, quin etiam extra Christum Deo displicent. In nobis autem esse salutis causae non possunt, quia fiutiqualir eus qu'antum ad communis charitatis officium praestandum tenemur, facimus, vt male ex hoc loco Papistae merita operum confirmare conentur.  \pend
\phantomsection
\addcontentsline{toc}{subsection}{\textit{20 Timothee, depositum serua, et auersare profanos de rebus inanibus clamores et oppositiones falso nominatae scientiae. }}
\subsection*{\textit{20 Timothee, depositum serua, et auersare profanos de rebus inanibus clamores et oppositiones falso nominatae scientiae. }}\pstart Conclusio totius Epistolae, in qua quid maxime sit Pastori fideli et amanti Dei gloriae, gregisque sui aedificationis fugiendum inculcat. Non frustra id quidem, neque temere: sed quod hoc inanis metaeologiae vitium fore tanquam fontem omnis corruptelae, quae in vera Theologia postea inducta est, iam tunc praeuideret, atque praediceret Spiritus sanctus. Ergo et mature et serio huic malo prospici et caueri voluit. Atque haec causa est, cur hanc admonitionem repetat, quae sup.vers.5. et cap.1.vers.6. tam diligenter posita erat. Primum autem breuiter Timotheo, id est, cuilibet legitimo pastori suum munus conmendat, atque ad illud strenue obeundum exhortatur, cùm ait, serua depositum. vti Coloss.4.v. 17. Archippum quoque de suo officio breuiter  \pend
\section*{AD I. PAVL. AD TIM. }
\marginpar{[ p.398 ]}\pstart admonet. Facit enim cùm nostra segnities, tum ipsius muneris onus et magnitudo, vt huiusmodi cohortationes etiam optimis quibusque sint necelariae, quando et pro seipso quoque a tota Ecclesia Deum orari voluit Paulus, vt recte suo munere fungeretur Ephes.6.v19. Depositi vero nomine quid hic Paulus significet quaesitum est. Alii referunt ad dona Timothei peculiaria. Alii ad doctrinam cognitionemque fidei, quae omnium Christianorum communis est. Alii ad munus ipsum pastoris, quo tunc in Ecclesia Ephesina fungebatur ad tempus Timotheus. Ego vero cùm hic in persona Timothei omnium fidelium Ecclesiae pastorum conditionem pingi et describi putem, refero vocem depositum ad ipsius, id est, pastorum munus. Quanquam enim fuit Euangelista Timotheus: non autem vni alicui Ecclesiae addictus pastor, tamen Euangelistae munus tunc in Ecclesia publicum munus fuit et Ecclesiasticum: et quo tempore haec scribebat Paulus, Ephesi velut pastoris partes susceperat ad tempus Timotheus. Est igitur pastoris munus Depositum, sed Dei depositum, non autem hominum, cui propterea illius depositi ratio reddenda est. Quanquam autem hominum suffragio pastores eliguntur, eius tamen electionis vis et authoritas a Deo ipso authore: non ab hominum voce pendet. Itaque non tam homines, quam Deus ipse cuique pastori Ecclesiam, et gregem cui praesicitur commendat, et velut depositum, quod repetiturus est, tradit seruandum ad tempus. Vnde illa est eiusdem Pauli sententia Act. 2o. versa8. Attendite igitur vos ipsos, et totum  \pend
\section*{CAPVT VE }
\marginpar{[ p.509 ]}\pstart gregem, in quo vos Spiritus sanctus constituit Episcopos, etc. Et Apocalyps.cap.1.v.2. et Malac. 2. nomine Angelorum et Nuntiorum Dei appellantur Ecclesiae pastores. Ergo huius depositi rationem reddent. Neque Domini sunt gregis sibi commissi, vt etiam Petrus ait 1.Pet.5. sed Deus iple mus gregis Doinus est. v ti depositarius rei apud se depositae Dominus non est: sed tantum est illius custos ad tempus.  \pend\pstart Seruare autem est non modo gregem Domini non minuere, sed etiam augere. At que ita se gerere, vt se Dei donis indignos et spoliandos non praebeant, qui ea habent, quique Ecclesiae praeficiuntur. Cùm enim Dei donis abutimur, aut ea Ipermime,ne que Minnue grorlam impendimus, tunc iis priuamur a Deo, ita quidem vt etiam auferantur ab eo, qui vel nihil vel pauca habet, quia quod habet talentum non conuertit in lucrum creditoris et Domini, a quo acceperat Matth.15.  \pend\pstart Cur autem depositum sibi datum non seruent homines, duplex causa est, vel quod doctrinae fideique puritas violatur et in haeresin vel mataeologiam conuertitur: vel quod turpis vitae exempla caeteris praebent, Deumque ipsi offendunt. Quo vtroque vitio carere debent, qui depositum et dona sibi a Deo collata habere, retinere et ferdate m idam iarttem cupiunt.  \pend\pstart Atque vt depositum Dei a nobis seruandum est: ita illius corruptela fugienda atque auersanda est. Ea autem est vniuersa mataeologia et quaestionum, oppositionum, contentionum et inanium speculationum plena disputatio, ostentatio, et scientia, quam falsam scientiam appellat Paulus, quae facit vt a vera fide illiusque doctrina lon\pend
\section*{AD I. PAVL. AD TIM. }
\marginpar{[ p.P ]}\pstart gissime aberrmus, et diuertamus.  \pend\pstart Ergo non taxat Paulus hoc loco apertas et manifestas haereses: sed quid minus, nimirum, quaestiones tantum ipsas et disputationes, ex quibnulla aedificatio, nullaque pietatis accessio nobis fit: sed quae totae in curiosa et inani rerum cognitione et inuestigatione versantur. Quales sunt Scholasticorum hodie omnes scholae, et qualis omnis Theologia, quae ab illis appellatur. Pingit igitur graphice Paulus illud malum, quod iam regnat latissime in Papatu, vtinam etiam in nullis Euangelicis Ecclesiis. Istorum methodum docet versari primùm in oppositionibus. Certe de quolibet fidei articulo in vtranque partem disputandum est in eorum scholis. Neque potest optimus et egregius doctor inter eos censeri, qui non sit Quod libetarius, et in vtranque partem, Academicorum more, de quolibet apice Theologico disputare paratus pro et contra, vt ipsi loquuntur, vt nihil sit apud eos tam certum, et deimirum, quod non in dubium reuocari possit contrariasque rationes habeat.  \pend\pstart Finem autem harum quaestionum et oppositionum docet esse vanam cognitionem et prorfus inutilem, imo vero etiam profanam, et in Deum ipsum tandem blasphemam.  \pend\pstart Denique modum et rationem earundem quaestionum habendarum et tractandarum explicat quam docet esse clamorem. Id quod ipsa experientia docet. Inde fit vt eorum alii sint Reales, alii Nominales, alii Scholastici Thomistae, et ita se in varias sectas et haereses quae nunquam tamen inter se placari poterunt, dissecent, atque  \pend
\section*{CAPVT  VI. }
\marginpar{[ p.511 ]}\pstart ambitiosissime partiantur.  \pend\pstart Neque tamen damnat Paulus omnes omnino συζητήσεις et quaestiones, quae de causis et modo salutis nostrae, fieri et agitari possunt: cùm ipse iubeat Episcopum esse ἐλεδάλικὸν, id quod sine disputatione fieri non potest: et cùm Paulus ipse laudet Apollo, quod fuerit in Iudaeis refutandis acutus et solers Tit 1.vers.9. Actor.18. vers.27. Sed eam demum tractandae fidei rationem improbat, ex qua nihil pietatis, nihil aedificationis, mim denque fonur emergie, aeque colligitur, quales sunt pene omnes Scholasticorum contentiosae, et operosae disputationes, quae in ipsorum scholis magna Monachorum corona habentur.  \pend\pstart Nec vero ignorauit Paulus quo potissimum fuco ilti homines le suarque illas inutiles disceptationes tegant, ac commendent, nimirum Scientiae, Sapientiae, Theologiae. Sed hunc fumum discutit, cùm tam praeclaro nomine donari falso hanc cognitionem et studium respondeat.  \pend\pstart Falso enim nominatur γνῶσις, cùm qui ea praediti sunt, semper querant, ad ipsam autem veritatis cognitionem nunquam perueniant, vt ait idem 2.Corinth.3.vers.7. Denique ista est profanatio verbi Dei, et euersio, non autem professio et intelligentia. Veteres, velut Augustinus, Ambrosius, Chrysostomus, Beda. et alii pro κενοφωνίας legunt per diphthongum κανορωνίας, id est, Nouitates verborum. Et certe vtrunque vitium et vanitatis et nouitatis inter se coniunctum est et affine: atque Scholasticis vtrunque accidit, vt dum nubes et inania captant pro solida fidei doctrinanoua quoque  et inaudita rerum vocabula, et pror\pend
\section*{AD I. PAVL. AD TIM. }
\marginpar{[ p.512 ]}\pstart sus inepta et barbara excudant, qualia apud eos sunt pene infinita. Praeclare tamen Augustinus libus 15.contra Faustum cap.3. Non omnis nouitas est fugienda:non omnis vetustas appetenda. Nam et Christus nouumse praeceptum dare dixit, et οκούσιον Patres vocemvt ait Athanasius tunc temporis nouam aduersus Arrianos induxerunt ad noubo nompe morum ertores nigulandos. Sed istorum noua vocabulares ipsas obscurant, et plane inuoluunt.  \pend
\phantomsection
\addcontentsline{toc}{subsection}{\textit{21 Quam nonnulli prae se ferentes, circa fidem aberrarunt a scopo. Gratia sit tecum Amen. }}
\subsection*{\textit{21 Quam nonnulli prae se ferentes, circa fidem aberrarunt a scopo. Gratia sit tecum Amen. }}\pstart Α'ιτιολογέα, confirmat enim superiorem sententiam ab effectu, ne frustra aut temere istam mataeologiam tam seuere damnare videatur. Euentus enim ipse iam tempore Pauli docuit omnes huiusmodi inanes quaestiones et clamoris, et ineptae occupationis plenas verissimam esse doctrinae fidei pestem, et labem. Nam ex iis qui, eas funt consectati, quidam a vera fide iam defecerant Quod si obiicitur. At nondum erat nata scholastica Theologia, quae demum fub Innocentio3. et post Alberti Magni tempora extare coepit, Thoma Aquinate nimirum multum eam promouente et commendante. Respond. Aliae quaestiones iam tempore Pauli pro vera et solida fidei doctrina inuehebantur tum a Iudaeis, tum ab haereticis, de quibus cap 1.et cap.4. et hoc ipso capite facta mentio est, quibus postea depulsis et explosis in earum locum scholastica et inanis Theologia, quae in Papistarum scholis viget, substituta  \pend
\section*{CAPVT VI. }
\marginpar{[ p.513 ]}\pstart est. Itaque iam tunc Scholasticae Theologiae et Mathaeologis quos hic reprehendit Paulus semina iaciebantur. Talis quoque est iuris Canonici ars et scientia, quae et ipsa quoque verissima est verae sacraeque l’heologiae, quae ex vno Dei verbo haurienda est, pernities pestis et corruptela.  \pend\pstart Aruuie Puume precauonemee mus confirmationem ac velut suum de ea votum et sigillum cùm ait Amen.  \pend\pstart Gratia, de qua agit, est Dei gratia, quae etiam nostrae salutis causa est. Est enim ea, per quam a Deo in ipsius filiorum numerum adoptamur, illique per Christum placemus et reconciliamur. Subscriptio vero quae huic Epistolae addita est, tanquam haec Epistola Laodiceae scripta esset (que hic dicitur Metropolis Phrygiae Pacatianae) est falsa, vt verissime variis rationibus euincit D. Th. Beza. Etsi enim Phrygiae Pacatianae facta est mentio in 5. Synodi, quae est Constantinopolita. actis: tamen neque apud Plinium qui post obitum Pauli scripsit, istius Phrygiae Pacatianae facta memtio est: neque apud vllum veterem scriptorem. Deinde ex Epistola Pauli ad Coloss.cap.2. vers.1. quae postrema aetate Pauli ab eo scripta est, et ante hanc ipsam, apparet Paulum nunquam Laodiceam adiisse quae in Phrygia sita erat, et Colossis atque  Hierapoli vicina vt ipse testatur. Laus Deo. FINIS  \pend
\endnumbering
\end{pages}
\end{document}
        
